\documentclass[12pt]{report}

% ---------------------------------------------------------------
% 字型與中文排版
% ---------------------------------------------------------------
\usepackage{fontspec}
\usepackage{xeCJK}
\setmainfont{Times New Roman}
\setCJKmainfont{Songti TC}
\setCJKsansfont{PingFang TC}
% 等寬字改用 Latin Modern Mono，跟 exp2 參考檔案的排版字體一致
% （它不是系統字型，只存在 TeXLive 樹裡，所以要用 Path 明確指定檔案位置）。
\setmonofont{lmmono10-regular}[
    Path = /usr/local/texlive/2024/texmf-dist/fonts/opentype/public/lm/,
    Extension = .otf,
    AutoFakeBold = 2.5,
]

% ---------------------------------------------------------------
% 本書排版套件（packages/ 目錄）
% ---------------------------------------------------------------
\usepackage{sleek}
\usepackage{sleek-theorems}
\usepackage{sleek-code}
\usepackage{sleek-diagrams}
% 注意：sleek-listings.sty 不載入——它是舊版 lstlistings/inconsolata 排版套件，
% 本書實際用的是 sleek-code.sty 的 minted 版 code/example 環境，兩者同時載入會讓
% minted 抓不到 \setmonofont 設定的等寬字（回退成 Latin Modern Roman），
% 且會讓程式碼裡的直引號被轉成全形彎引號。

% 章節標題中文化：「第 N 章」
\usepackage{titlesec}
\titleformat{\chapter}{\normalfont\huge\bfseries}{第~\thechapter~章}{1em}{}
\titlespacing*{\chapter}{0pt}{-20pt}{20pt}

\title{C 程式設計基礎}
\author{}
\date{\today}

\begin{document}

\maketitle
\romantableofcontents

% ---------------------------------------------------------------
% 各章節內容：每個單元的敘述式教材放在對應的
% "chapters/<單元>/chapter.tex"，
% 與該單元的題目（範例練習）放在同一個資料夾方便對照維護。
% 題目本身不在正文中逐題引用——沿用原書「教材與習題分開」的慣例。
% ---------------------------------------------------------------

\chapter{介紹}
\label{ch:introduction}

程式設計，白話說，就是寫出一段段指令，命令計算機替我們完成工作。
計算機本身只看得懂非常底層的機器指令，例如把記憶體某個位置的數字
載入暫存器、再加 $3$ 這類一步一動作的操作。但我們不可能用這種方式
寫出稍具規模的程式。高階程式語言讓我們可以用接近日常思考的方式描述
問題，再由\emph{編譯器}把這些描述翻譯成計算機真正執行的機器碼。本書
要教的高階語言是 C。它夠接近硬體，讓我們能看清程式在記憶體裡實際
發生什麼事，又比機器語言或組合語言好寫得多。對之後想走系統程式、
作業系統、嵌入式或底層開發的讀者，這是一條很實用的起點。

這一章不會一次塞進太多語法細節。我們先把幾件基礎事情講清楚：程式到
底是什麼、機器語言和高階語言差在哪裡、為什麼需要編譯器、寫程式時如
何用流程圖把思路畫出來，以及一份 \texttt{.c} 文字檔怎麼變成螢幕上的
結果。路徑是：先看到一支最小的 \texttt{Hello, world!} 並逐行讀懂，
再用這支程式示範 \texttt{gcc} 與編譯流程，最後集中說明語法、執行期
與邏輯三種錯誤。更完整的變數、輸入輸出與運算，留到下一章基本程式架
構再細講。

\section{程式是什麼}
\label{sec:intro-what-is-programming}

程式的英文是 \emph{program}。如果把計算機想成一位只會照劇本演出的演
員，那程式就是給它看的劇本：劇本裡一行行寫著先做 A、再做 B、如果條
件成立就改做 C。計算機的中央處理器（CPU）其實只做一件很單純的事
。從記憶體讀出一條指令、執行它、再讀下一條。但把大量這種簡單動作
串起來，就能播放音樂、渲染網頁、控制馬達、或是批改作業。

我們平常接觸到的檔案、照片、影片，在計算機裡本質上都是\emph{數字}：
文字每個字元對應一個編號，圖片的每個像素也是一組數字。CPU 能直接處
理的指令同樣是一組組數字，這種計算機原生的指令格式就叫做\emph{機器語
言}（machine language）。機器語言對人類極不友善：同一件把兩個整數
相加的事，可能要拆成把記憶體位址 $A$ 的值載入暫存器、把記憶體位
址 $B$ 的值載入另一個暫存器、把兩個暫存器相加、把結果寫回記憶體位址
$C$ 好幾步，而且每一步的編碼還會隨 CPU 廠牌而不同。

\emph{暫存器}（register）是 CPU 內部極少量、極快的儲存格；\emph{記憶
體}（main memory）容量大得多，但存取比較慢。程式執行時，資料經常在
這兩層之間搬來搬去。現在不必把硬體細節背下來，只要先有這個畫面：我
們寫的每一個變數，最終都會落在某個記憶體位址上；CPU 則靠暫存器做計
算。後面學到指標時，這個畫面會變成核心工具。

\emph{高階程式語言}（high-level programming language）在機器語言和我
們之間加了一層：我們寫 \texttt{total = a + b;}，編譯器負責把它展開成
機器看得懂的指令序列。這樣我們可以把精力放在要解決什麼問題，而
不是這顆 CPU 的第幾號暫存器該放什麼。本書使用的 C 就是這樣一種
高階語言。但它又比許多現代語言更貼近機器：C 程式編譯後通常很精簡、
跑得很快，而且語言本身提供了直接操作記憶體位址、位元組合的能力，所
以作業系統、驅動程式、嵌入式韌體這類必須跟硬體打交道的軟體，長
期以來都以 C 為主要工具之一。

\keypoint{程式是給計算機執行的指令序列。機器語言是 CPU 看得懂的數字
編碼；高階語言讓人用接近思考的方式寫，再由編譯器翻譯成機器語言。}

跟機器語言同一層、但用助記符號（例如 \texttt{ADD}、\texttt{MOV}）來寫
的，叫做\emph{組合語言}（assembly language）。組合語言仍然幾乎一對一
對應某顆 CPU 的指令，換一種處理器通常就要重寫。C 比組合語言高一層：
同一份 C 原始碼，只要有對應的編譯器，就有機會在不同硬體上產生各自的
機器碼。這就是下一節要說的可攜性。

\begin{drills}[label=drills:intro-what-is-programming]
\item 機器語言跟高階語言，哪一種是 CPU 直接看得懂的？哪一種比較貼近
人類原本的思考方式？
\item 暫存器（register）跟記憶體（main memory），哪一個容量比較大、
但存取比較慢？
\item 用一句話說明「編譯器」在這一節裡扮演的角色。
\end{drills}

\section{為什麼是 C}
\label{sec:intro-why-c}

C 語言誕生於 1970 年代初的 Bell Labs。Dennis Ritchie 與 Ken Thompson
在開發 UNIX 作業系統時，需要一種比組合語言更好維護、又比當時其他高
階語言更高效的語言；C 從 Thompson 的 B 語言演化而來，並在 1973 年前後
讓整個 UNIX 核心改用 C 重寫。這件事的意義很大：同一套 C 原始碼，只
要換一個\emph{編譯器}，就有機會在不同的硬體上執行，\emph{可攜性}
（portability）從此成為 C 的招牌之一。

後來 Brian Kernighan 與 Dennis Ritchie 合著的 \emph{The C Programming
Language}（常被稱為 K\&R）在 1978 年出版，幾乎成為當時 C 的事實標準；
1989 年 ANSI 公布正式標準（常稱 C89 或 C90），之後又有 C99、C11 等修
訂。

\textbf{本書採用 C99 標準}。C99 是 1999 年制定的標準，在現代實務中仍
是最常見的選擇：相比 C89，它新增了許多方便的特性（例如 \texttt{for}
迴圈裡可以直接宣告迴圈變數 \texttt{int i}、支持 \texttt{//} 單行註解），
又沒有 C11 以後那些比較複雜的新機制。本書所有範例和說明都依循 C99 標準，
用 \texttt{gcc -std=c99} 或 Clang 等現代編譯器都可以順利編譯執行。如果
你的環境是舊版 C89，會遇到個別語法不相容（例如迴圈變數宣告會失敗），到
時候按照編譯器的錯誤訊息調整即可。

跟 Python、Java 這類記憶體由執行環境代管的語言相比，C 把更多責任
留給程式設計師：陣列有沒有越界、指標指到哪裡、配置到的記憶體有沒有
歸還，編譯器往往不會幫你擋。這聽起來像缺點，但對想理解程式在機器
上實際怎麼跑的人來說，反而是優點。你寫的每一行，比較容易對應到記
憶體裡發生的事。許多語言的執行環境、虛擬機、甚至作業系統核心，底層
仍大量用 C 寫成；學會 C，等於拿到一把能往下挖的鏟子。

當然，C 不是唯一選擇，也不會是終點。C++、Rust、Go 等語言在 C 的基礎
上各自加了不同的抽象。但 C 的語法與思維方式影響深遠：學過 C 之後再
讀其他語言，很多底層概念（堆疊、指標、位元表示）會比較容易對上號。
本書後半會一路教到指標、記憶體配置、位元操作與資料結構。都是走
low-level／systems 方向時避不開的基本功。

\keypoint{C 夠接近硬體、執行效率高、可攜性好；代價是許多錯誤要靠程
式設計師自己避免。本書選擇 C，正是為了把記憶體裡實際發生什麼事
講清楚。}

\subsection{編譯器和直譯器差在哪裡}

把高階語言變成機器能跑的東西，常見有兩條路。\emph{編譯器}把整份程
式一次翻譯成執行檔，之後每次執行都直接跑那份機器碼，不必再看原始
碼。C 走這條路：你改一行，就要重新編譯，換來的是執行時通常比較快、
也比较接近硬體。

\emph{直譯器}（interpreter）則是邊讀邊執行：它讀到一條敘述，立刻做
那一條要求的事，通常不會先產出一份獨立的執行檔。許多教學語言一開
始用直譯，改完馬上試，回饋比較快；代價是執行期多了一層讀程式、再
做事的開銷，而且發佈給別人跑時，對方機器上也得有那個直譯器。

現實世界很少那麼黑白。有些語言先編譯成一種中間碼，再由虛擬機執行；
有些直譯器也會把熱門片段編譯加速。對我們來說，先抓住一件事就夠：
\textbf{C 程式必須先通過編譯，才有東西可以執行}。編譯失敗時，螢幕
上不會出現你預期的 \texttt{Hello, world!}，只會出現編譯器的抱怨。
後面會先給出第一支程式，再示範 \texttt{gcc} 怎麼把它變成可執行檔，
最後集中說明怎麼讀編譯器與執行期的訊息。

標準函式庫也屬於這條翻譯鏈的一部分。C 語言核心並沒有印一行字
這種內建指令；\texttt{printf} 的宣告寫在 \texttt{stdio.h} 裡，真正的
機器碼則在系統提供的函式庫裡。編譯器負責檢查你有沒有用對，連結器
負責把函式庫裡的實作接進來。這就是為什麼最小的程式也要寫
\texttt{\#include <stdio.h>}：它是告訴工具鏈我要用標準
輸入輸出這一套現成的工具。

\begin{drills}[label=drills:intro-why-c]
\item C 語言最早是為了開發哪個作業系統而設計出來的？
\item 本書採用哪一個 C 標準？相比 C89，它新增的其中一項方便特性是什
麼？
\item 編譯器跟直譯器，哪一種是「先把整份程式翻譯成執行檔，之後每次
執行都直接跑那份機器碼」？C 走哪一條路？
\item 為什麼很多作業系統、驅動程式仍然用 C 寫，而不是 Python 這類語
言？
\end{drills}

\section{用流程圖把思路畫出來}
\label{sec:intro-flowchart}

還沒寫程式之前，先問一個問題：這支程式要做哪些步驟、遇到條件時走哪
一條路？ 把答案畫成圖，常常比直接在編輯器裡打字更不容易卡住。
\emph{流程圖}（flowchart）用方塊表示一個動作，用菱形表示一個是非
問題，用箭頭表示下一步往哪裡走。流程圖給\emph{人}看的草圖，編譯
器不會讀它；後面學到判斷與迴圈時，菱形對應 \texttt{if} 的條件，繞
回來的箭頭對應 \texttt{while}。現在只先養成習慣：步驟不清時先畫，
不要一開場就在細節裡迷路。

下面這張圖比 \texttt{Hello, world!} 還小：先印一行字，然後結束。中
間沒有任何判斷。下一節就要看到對應的 C 程式。

\begin{center}
\begin{tikzpicture}[node distance=10mm]
    \node[startend] (s) {開始};
    \node[process, below=of s] (p) {印出 Hello, world!};
    \node[startend, below=of p] (e) {結束};
    \draw[flowarrow] (s.south) -- (p.north);
    \draw[flowarrow] (p.south) -- (e.north);
\end{tikzpicture}
\diagcaption{最小流程圖：沒有判斷，從開始依序執行一個動作後結束。
對應到下一節的第一個 C 程式。}
\label{fig:ch01-hello-world-flow}
\end{center}

\stylepoint{流程圖裡每個判斷只問一件事，每個動作方塊只做一件事。圖
畫得太擠，通常代表思路還沒切乾淨，寫成程式也會擠在同一段裡。}

\begin{drills}[label=drills:intro-flowchart]
\item 流程圖裡，方塊、菱形、箭頭分別代表什麼？
\item 編譯器會不會讀你畫的流程圖？它是畫給誰看的？
\item 圖~\ref{fig:ch01-hello-world-flow} 那張「最小流程圖」裡，一共有
幾個動作方塊？做了什麼事？
\end{drills}

\section{第一個 C 程式}
\label{sec:intro-first-program}

傳統上，學一門新語言的第一個程式會印出 \texttt{Hello, world!}。它夠
短，卻已經包含一個完整 C 程式必備的幾個部分：引入標準標頭、定義
\texttt{main} 函式、在函式裡下達指令、最後把控制權交還作業系統。範
例 \thechapter.1 就是這個程式；我們先把它讀懂，再逐行說明每一部分
在做什麼。

\begin{example}{c}{(hello-world.c) 第一個 C 程式}
#include <stdio.h>

int main(void) {
    printf("Hello, world!\n");
    return 0;
}
\end{example}

\iosample{}{
Hello, world! \\
}

這個程式沒有讀任何輸入，因此下面只示範輸出；執行後終端機只會印出一
行 \texttt{Hello, world!}，然後程式正常結束。

\subsection{逐行閱讀}

\textbf{第 1 行：\texttt{\#include <stdio.h>}}

這是給\emph{前置處理器}（preprocessor）的指令。
前置處理器在編譯器正式分析程式之前先跑一輪；\texttt{\#include} 的意
思是把另一個檔案的內容插入到這裡。\texttt{stdio.h} 是 C 標準函
式庫裡與\emph{標準輸入／輸出}（standard I/O）有關的標頭檔（header），
裡面宣告了 \texttt{printf} 等函式，讓編譯器知道呼叫 \texttt{printf}
時參數該長什麼樣子。沒有這一行，後面的 \texttt{printf} 可能無法通
過編譯。標頭檔通常只放\emph{宣告}；真正的函式機器碼在編譯器附帶的
函式庫裡，連結階段會自動把需要的部分拉進執行檔。

尖括號 \texttt{<stdio.h>} 表示到系統提供的標準目錄裡找這個檔案。
之後若要引入自己寫的標頭，會改用雙引號 \texttt{"myheader.h"}，那是後
面章節的事。現在只要記住：幾乎每個會印東西、讀輸入的程式，開頭都有
這一行。

\textbf{第 2 行：空行}

\texttt{\#include} 之後空一行，再進入 \texttt{main}，是本書（以及許
多 C 專案）的排版習慣：把引入外部標頭和我們自己寫的程式本體
視覺上分開，比較好讀。C 編譯器會忽略多餘的空白行，加不加不影響語意。

\textbf{第 3 行：\texttt{int main(void) \{}}

C 程式是由\emph{函式}（function）組成的；每個可執行程式都必須有一個
名叫 \texttt{main} 的函式。程式開始執行時，作業系統會把控制權交給
\texttt{main}，從這裡開始跑。\texttt{int} 表示 \texttt{main} 結束時
要\emph{回傳}一個整數給作業系統：慣例上回傳 \texttt{0} 代表正常結
束，非零值常用來表示發生錯誤。函式名稱後面的 \texttt{(void)} 明確
表示 \texttt{main} 不接收任何參數——\texttt{void} 在這裡的意思是
「空的」；若需要接收命令列參數，可以寫成
\texttt{int main(int argc, char *argv[])}，那是進階話題，現在先記住
程式從 \texttt{main} 開始就夠了。行尾的 \texttt{\{} 標記函式
\emph{本體}的開始，對應的最後一行 \texttt{\}} 標記結束。

本書之後所有範例都會寫成 \texttt{int main(void)}，不會省略成
\texttt{int main()}：兩者在大多數編譯器裡效果相同，但空括號
\texttt{()} 在 C 語言裡原本代表「參數不確定」，只是幾乎所有編譯器都
把它當空參數處理；\texttt{(void)} 才是 C 標準明確定義的「不接收參
數」寫法，讀起來也更不容易誤會。

\textbf{第 4 行：\texttt{printf("Hello, world!\textbackslash n");}}

\texttt{printf} 是標準函式庫裡負責\emph{格式化輸出}的函式，名字來自
print formatted。括號裡的字串 \texttt{"Hello, world!\textbackslash n"}
是要印的文字；其中的 \texttt{\textbackslash n} 是\emph{跳行}
（newline）字元，代表印完這行後換到下一行。如果拿掉
\texttt{\textbackslash n}，游標會停在同一行末尾，後續輸出會接在後面
。很多初學者的第一個怎麼輸出怪怪的問題，都跟有沒有換行有關。
整行結尾的分號 \texttt{;} 不能少：在 C 裡，大多數敘述都以分號結束，
漏掉分號是編譯階段最常見的語法錯誤之一。

\textbf{第 5 行：\texttt{return 0;}}

把 \texttt{0} 回傳給作業系統，表示程式依預期跑完。對這麼短的程式，
即使省略 \texttt{return 0;}，在 C99 之後的環境裡 \texttt{main} 也常會
被視為隱含回傳 \texttt{0}；但養成明確寫 \texttt{return 0;} 的習慣，之
後寫較大的程式時比較不會搞混到底有沒有正常結束。

\textbf{第 6 行：\texttt{\}}}

對應第 3 行的 \texttt{\{}，函式本體在這裡結束。整個 \texttt{main} 裡
目前只有一條有效敘述（\texttt{printf}）加上 \texttt{return}；之後章
節會在 \texttt{\{} \texttt{\}} 之間放進變數宣告、迴圈、條件判斷等更
多內容，但骨架始終是 \texttt{main} 收到控制權 $\rightarrow$ 依序執
行裡面的敘述 $\rightarrow$ 回傳結束碼。

同一個程式也可以拆成兩次 \texttt{printf}。範例 \thechapter.2 先印
\texttt{Hello,} 並換行，再印 \texttt{world!} 並換行。執行後會看到兩
行，而不是一行裡的 \texttt{Hello, world!}。差在逗號後面有沒有立刻換
行，以及 \texttt{world} 前面有沒有空白。這只是在示範：\texttt{main}
裡的敘述依照寫下的順序一條一條執行，每一次 \texttt{printf} 都接著上
一次結束的位置繼續印。

\begin{example}{c}{(hello-two-lines.c) 用兩次 printf 印出兩行}
#include <stdio.h>

int main(void) {
    printf("Hello,\n");
    printf("world!\n");
    return 0;
}
\end{example}

\iosample{}{
Hello, \\
world! \\
}

\texttt{\textbackslash n} 這種寫在字串裡、以反斜線開頭的組合，叫做
\emph{跳脫序列}（escape sequence）。它不是把兩個字元反斜線和
n 印到螢幕上，而是編譯器把它理解成一個特殊動作——這裡是換行。
常用的還有 \texttt{\textbackslash t}（水平定位，看起來像跳到下一個
定位點）、以及 \texttt{\textbackslash\textbackslash}（真的要印出一
個反斜線時用）。範例 \thechapter.3 把逗號後面的換行改成定位，執行
後 \texttt{Hello,} 與 \texttt{world!} 會出現在同一行，中間隔開一段
空白。

\begin{example}{c}{(hello-tab.c) 用跳脫序列插入定位}
#include <stdio.h>

int main(void) {
    printf("Hello,\tworld!\n");
    return 0;
}
\end{example}

\iosample{}{
Hello,	world! \\
}

若把範例 \thechapter.1 裡的 \texttt{\textbackslash n} 整段拿掉，編譯
仍然成功，螢幕上也看得到 \texttt{Hello, world!}，但下一行的終端機提
示符會黏在同一行後面。這只是程式沒有把游標移到下一行。
養成習慣：每一段你希望單獨成行的輸出，都自己寫上換行，不要依
賴作業系統幫忙。

\keypoint{最小的 C 程式需要 \texttt{\#include <stdio.h>}、一個
\texttt{main} 函式，以及至少一條敘述。\texttt{printf} 字串裡的
\texttt{\textbackslash n} 負責換行，敘述結尾的分號不能省。}

\subsection{註解：寫給人看的說明}

編譯器只在意語法是否合法；人還需要記得這一段為什麼這樣寫。
\emph{註解}（comment）會被編譯器整段忽略，專門留給閱讀程式的人。C 有
兩種常見寫法：\texttt{/* ... */} 可以跨很多行；\texttt{//} 從雙斜線寫
到該行結尾（C99 起成為標準的一部分，現代編譯器都收）。

\begin{code}{c}{兩種註解的寫法}
/* 這一段到右方星號斜線為止，都可以寫說明。*/

// 從雙斜線到這一行結束，都是註解。
printf("Hello, world!\n");  // 印出一行字並換行
\end{code}

註解不要逐行把程式翻譯成中文（這裡把 $i$ 加一這種句子，看程式本
身就知道）。比較有用的是解釋\emph{為什麼}：為什麼緩衝區要多留五格、
為什麼這裡用 \texttt{long long} 而不是 \texttt{int}。範例
\thechapter.1 短到幾乎不需要註解；程式一長，沒有註解的區塊就會變成謎。

\stylepoint{註解寫為什麼，不要寫做了什麼。程式碼已經說了做
什麼；註解負責補上讀程式的人看不到的意圖。}

\begin{drills}[label=drills:intro-first-program]
\item 把範例 \thechapter.1 裡的 \texttt{\textbackslash n} 拿掉，猜猜終
端機會出現什麼狀況？
\item \texttt{int main(void)} 裡的 \texttt{(void)} 是什麼意思？跟省
略成 \texttt{int main()} 差在哪裡？
\item \texttt{//} 註解跟 \texttt{/* */} 註解，哪一種可以跨很多行？
\item 要把 \texttt{hello-world.c} 改成印兩次 \texttt{Hello, world!}，
該怎麼改？
\end{drills}

\section{從原始碼到執行}
\label{sec:intro-build-pipeline}

上一節已經看過 \texttt{hello-world.c} 長什麼樣子。這一節用同一份程式
示範：怎麼把它存成 \texttt{.c} 檔、用 \texttt{gcc} 編譯、在終端機執
行。寫程式的日常很少一次寫對；比較常見的是寫成 \texttt{.c} 檔、編
譯、執行、對照結果、再改。編譯器和機器只會照你寫下的做，不會自動
猜你其實想做什麼。

\subsection{編輯}

用任意文字編輯器把程式寫進一個檔案，慣例上副檔名用 \texttt{.c}（例如
\texttt{hello-world.c}），方便一眼看出這是 C 原始碼。編輯器本身不
懂 C，它只是幫我們存檔；有些編輯器會順便標示語法、自動縮排，但那
是額外的幫手，不是必要條件。

原始碼必須是純文字。用文書處理軟體存成帶格式的檔案，編譯器會看不懂。
檔名盡量只用英文字母、數字、連字號與底線，避免空白與中文。C 語言本身
並沒有禁止這種命名習慣，只是這樣終端機指令比較不容易寫錯。

\subsection{編譯：前置處理、翻譯、連結}

編譯器讀取 \texttt{.c} 檔，檢查語法是否合法，再把合法的部分翻譯成機
器碼，通常產生一個\emph{執行檔}（executable）。表面上我們只打一條
\texttt{gcc} 指令，內部其實常分成好幾段工作，初學者先記住三個名字就
夠了：

\begin{enumerate}
    \item \textbf{前置處理}（preprocessing）：處理以 \texttt{\#} 開頭
          的指令，例如把 \texttt{\#include <stdio.h>} 展開成標頭檔的內
          容。這一步還在文字層級，還沒開始真正翻譯成機器碼。
    \item \textbf{編譯}（compiling，狹義）：把處理過的原始碼翻譯成目
          標碼（object code）。語法錯誤多半在這一步被抓出來。
    \item \textbf{連結}（linking）：把我們的目標碼，跟標準函式庫裡
          \texttt{printf} 等函式的實作接在一起，產生最終執行檔。標頭
          檔只告訴編譯器 \texttt{printf} 長什麼樣子；真正的機器碼
          在函式庫裡，靠連結器找出來。
\end{enumerate}

這三步也可以拆開看。若執行
\texttt{gcc -c hello-world.c}，編譯器會停在目標碼這一層，通常產
生 \texttt{hello-world.o}：裡面已經是機器指令，但還沒接上
\texttt{printf} 的實作，因此還不能直接 \texttt{./} 執行。再對
\texttt{.o} 做連結，才會得到執行檔。初學階段用
\texttt{gcc hello-world.c -o hello-world} 一次做完即可；之後寫到多個
\texttt{.c} 檔組成一個程式時，就會真的需要先分別編成 \texttt{.o} 再
連結。現在只要知道：你在編輯器裡看到的文字，跟作業系統載入進記憶
體執行的位元組，中間隔著好幾層翻譯。

本書使用 \texttt{gcc} 編譯器。在終端機裡的典型用法如下：

\begin{code}{c}{用 gcc 編譯一個 C 原始碼檔}
gcc -Wall 原始碼.c -o 執行檔名稱
\end{code}

例如把 \texttt{hello-world.c} 編譯成名叫 \texttt{hello-world} 的執行檔，
可以輸入 \texttt{gcc -Wall hello-world.c -o hello-world}。\texttt{-o}
用來指定輸出的執行檔名稱；本書範例一律這樣寫，執行時打
\texttt{./hello-world}。若省略 \texttt{-o}，\texttt{gcc} 會用預設檔名
\texttt{a.out}，執行時則打 \texttt{./a.out}。

指令建議加上 \texttt{-Wall}，請編譯器把常見的可疑寫法用警告講出來：

\begin{code}{c}{編譯時打開常見警告}
gcc -Wall hello-world.c -o hello-world
\end{code}

\texttt{-Wall} 並不會很魔法地把所有警告都挑出來。編譯器會選擇
對初學者特別有用的檢查。養成習慣之後，很多執行期才爆出的問題，會在編譯階
段就被點名。

編譯階段若原始碼違反 C 的語法規則（例如少寫分號、括號不配對），編譯
器會報\emph{錯誤}（error），且\textbf{不會}產生執行檔。有時編譯器也會給\emph{警告}（warning）：語法上勉強過
得去，但可能有隱患。警告不一定會阻止產生執行檔，但習慣上仍應認真閱
讀。漏分號這類語法錯誤長什麼樣子，下一節會用剛才那支程式示範。

\keypoint{編譯失敗表示原始碼還不是合法的 C，執行檔不會出現。先把編
譯器指出的那一行（以及它的上一行）看完，再改、再編譯；不要猜測著亂
改。}

\subsection{執行}

編譯成功後，在終端機啟動執行檔。在 Unix 或 macOS 上，若執行檔在
\emph{目前目錄}，通常要寫成 \texttt{./hello-world}。前面的
\texttt{./} 表示就在這個資料夾裡找，避免跟系統 PATH 裡其他同名程
式搞混。程式跑起來後，它可能讀取鍵盤輸入、在螢幕印出結果；執行到一半才暴露的問題，叫做\emph{執行期錯誤}
（runtime error），得靠測試與除錯回頭修正原始碼，再重新編譯。

\subsection{自己動手試一次}

若你的環境已安裝 \texttt{gcc}，可以照下面步驟重現範例 \thechapter.1
（路徑與檔名可自訂，重點是副檔名 \texttt{.c} 與編譯指令）：

\begin{enumerate}
    \item 用編輯器建立 \texttt{hello-world.c}，內容與範例
          \thechapter.1 相同。
    \item 在該檔所在目錄執行
          \texttt{gcc -Wall hello-world.c -o hello-world}。
    \item 執行 \texttt{./hello-world}，應看到一行
          \texttt{Hello, world!}。
    \item （選做）把 \texttt{printf} 那一行的分號刪掉再編譯，對照終端
          機裡的錯誤訊息；細節見下一節。
    \item （選做）把範例 \thechapter.1 改成範例 \thechapter.2 的兩次
          \texttt{printf}，確認終端機出現兩行而不是一行。
    \item （選做）把 \texttt{\textbackslash n} 改成
          \texttt{\textbackslash t}，對照範例 \thechapter.3 的輸出。
\end{enumerate}

\begin{drills}[label=drills:intro-build-pipeline]
\item 編譯的三個步驟依序是什麼？各自大致負責什麼？
\item \texttt{-Wall} 這個編譯選項的作用是什麼？
\item 如果編譯指令省略了 \texttt{-o}，產生的執行檔預設叫什麼名字？
\item 在 Unix 或 macOS 上，為什麼要打 \texttt{./hello-world} 而不是
直接打 \texttt{hello-world}？
\end{drills}

\section{三種錯誤：語法、執行期、邏輯}
\label{sec:intro-error-types}

寫程式幾乎一定會出錯。比較有用的分類是錯在哪一層。因為不同層的錯，
出現的時機不一樣，修法也不一樣。

\textbf{語法錯誤}（syntax error）發生在編譯階段。C 對標點極嚴：敘
述結尾要分號、括號要成對、字串的直引號要關起來。把範例
\thechapter.1 的 \texttt{printf} 那一行分號拿掉，在這台機器上用
\texttt{gcc -Wall hello-world.c -o hello-world} 編譯時，得到的訊息大
意是：第 $4$ 行結尾 \texttt{expected ';' after expression}，並用
\texttt{\^{}} 指出少分號的位置。有錯誤就\textbf{不會}產生執行檔，只能
改原始碼再編譯。編譯器有時也會給警告：語法勉強合法，但寫法可疑。
警告擋不住執行檔，卻常常是下一類錯誤的預告，不要當成裝飾關掉。錯誤
訊息有時會指到下一行。因為編譯器讀完一行，到下一
行才確定上一句還沒結束。除錯時記得也檢查看起來沒報錯的上一行。

\textbf{執行期錯誤}（runtime error）發生在程式已經跑起來之後：語法
過了，但執行到某一刻，行為超出語言保證的範圍，或作業系統把程式停
掉。整數除以零在 C 裡是\emph{未定義行為}（undefined behavior）。
在 x86 的 Linux 或 macOS 上，整數除以零\emph{常見}的結果是 CPU 觸
發硬體例外、作業系統中斷程式，終端機可能出現 \texttt{Floating point
exception}、\texttt{Abort trap} 或 \texttt{core dumped} 這類訊息。
語言標準\emph{不保證}每次都崩潰，也不保證印出同一個數字；有的環境
可能靜默印出一個看起來合理的值。本書只描述觀察到的現象，不把某次
結果當成標準答案。更後面才會碰到的存取不該碰的記憶體，在 Unix 上常以
\emph{segmentation fault} 的形式出現。共同點是：\texttt{gcc} 當下
可能一句話都不說，問題要等執行才爆。

\textbf{邏輯錯誤}（logical error）最陰險：編譯過、執行也正常結束，
螢幕上就是有字，只是那不是題目要的答案。除錯（debug）平常說的，
多半就是在找這種錯。範例 \thechapter.7 看起來跟範例 \thechapter.1
幾乎一樣，差在 \texttt{printf} 被寫了兩次。若規格是只印一行
問候語，這個程式語法完全合法、也不會崩潰，但輸出多了一行。機器
很聽話，它印了你寫下的每一句，沒有義務去猜你其實只想印一次。

\begin{example}{c}{(hello-twice.c) 能跑、但結果不符合預期}
#include <stdio.h>

int main(void) {
    printf("Hello, world!\n");
    printf("Hello, world!\n");
    return 0;
}
\end{example}

\iosample{}{
Hello, world! \\
Hello, world! \\
}

下一章開始讀入數字、做運算之後，邏輯錯誤會變成該加寫成乘、該
比大小寫反這類更難一眼看出的狀況。現在先建立一個習慣：每寫完一
版，都用自己心裡的預期輸出對一下螢幕。流程圖可以畫成下面這樣。
三個菱形分別對應上面三種錯。任何一個答案不滿意，都回到編輯。

\begin{center}
\begin{tikzpicture}[
node distance=8mm and 12mm,
% 標籤文字樣式 (帶白底以防遮擋線條)
label text/.style={
    font=\scriptsize\sffamily, fill=white, inner sep=2pt
}
]

% --- 節點定義 ---
\node[startend] (edit) {編輯原始碼};
\node[process, below=of edit] (cc) {編譯};
\node[decision, below=of cc] (syn) {有語法錯誤？};
\node[process, below=of syn] (run) {執行};
\node[decision, below=of run] (rt) {執行期錯誤？};
\node[decision, below=of rt] (ok) {結果正確？};
\node[startend, below=of ok] (done) {完成};

% --- 主流程箭頭 (向下) ---
\draw[flowarrow] (edit.south) -- (cc.north);
\draw[flowarrow] (cc.south) -- (syn.north);
\draw[flowarrow] (syn.south) -- node[label text] {否} (run.north);
\draw[flowarrow] (run.south) -- (rt.north);
\draw[flowarrow] (rt.south) -- node[label text] {否} (ok.north);
\draw[flowarrow] (ok.south) -- node[label text] {是} (done.north);

% --- 迴圈箭頭 (向右折返) ---
\draw[flowarrow] (syn.east) -- node[above, font=\scriptsize] {是} ++(2.0cm,0) |- (edit.east);
\draw[flowarrow] (rt.east)  -- node[above, font=\scriptsize] {是} ++(3.2cm,0) |- (edit.east);
\draw[flowarrow] (ok.east)  -- node[above, font=\scriptsize] {否} ++(4.4cm,0) |- (edit.east);

\end{tikzpicture}
\diagcaption{寫程式的除錯迴圈：先過語法，再過執行，最後核對結果是
否符合預期。任何一關失敗都回到編輯。}
\label{fig:ch01-debug-loop}
\end{center}

\keypoint{錯誤比警告嚴重：有語法錯誤就沒有執行檔。警告仍應讀完。
能跑也不等於對。邏輯錯誤要靠你拿預期結果去對。}

\stylepoint{先讓程式能編譯，再讓它能跑完，最後才追求答案正確。三
件事混在一次修改裡，出了問題會不知道該從哪一行查起。}

到這裡，我們已經具備幾個起點：程式是高階語言寫成、由編譯器翻譯後執
行的指令序列；流程圖用來整理步驟；C 是貼近系統、可攜又高效的主力語
言之一；一個最小的 C 程式需要 \texttt{\#include}、\texttt{main} 與至
少一條敘述；而寫程式的日常，就是在語法、執行與邏輯三層之間來回修
正。下一章會在這個骨架上，從 \texttt{printf} 輸出格式與變數接下去，
不再重講 \texttt{\#include} 與編譯三步驟。到時候 \texttt{Hello, world!} 裡的
\texttt{\textbackslash n} 和 \texttt{printf} 格式字串，會變成可以靈活
組合的工具，而不只是照抄的一行魔術。

\begin{drills}[label=drills:intro-error-types]
\item 把 \texttt{printf} 那一行的分號拿掉，這個錯誤會在哪個階段被抓
出來？
\item 整數除以零，在 C 的三種錯誤分類裡屬於哪一種？
\item 範例 \thechapter.7（\texttt{hello-twice.c}）屬於哪一種錯誤？它
明明能編譯、能執行，為什麼還是有問題？
\item 「警告」（warning）跟「錯誤」（error）有什麼不同？警告會不會
擋下執行檔？
\end{drills}

\addcontentsline{toc}{section}{練習解答}
\section*{練習解答}

各節練習的答案收在這裡。 短題只給能一句話核對的結果。

\paragraph{第~\ref{sec:intro-what-is-programming}~節}
\begin{enumerate}
\item 機器語言是 CPU 直接看得懂的；高階語言比較貼近人類的思考方式，
要靠編譯器翻譯成機器語言才能執行。
\item 記憶體容量比較大，但存取比暫存器慢；暫存器容量極少，但速度很
快。
\item 編譯器把我們寫的高階語言描述，翻譯成計算機真正執行的機器碼。
\end{enumerate}

\paragraph{第~\ref{sec:intro-why-c}~節}
\begin{enumerate}
\item UNIX。
\item C99；例如 \texttt{for} 迴圈裡可以直接宣告迴圈變數
\texttt{int i}，或支援 \texttt{//} 單行註解（任一個即可）。
\item 編譯器；C 走這條路，先整份翻譯成執行檔，之後每次執行都直接跑
那份機器碼。
\item 因為 C 把記憶體管理的責任交給程式設計師，執行效率高、又貼近硬
體，適合需要直接跟硬體打交道的系統程式。
\end{enumerate}

\paragraph{第~\ref{sec:intro-flowchart}~節}
\begin{enumerate}
\item 方塊表示一個動作，菱形表示一個是非問題，箭頭表示下一步往哪裡
走。
\item 不會。 流程圖是畫給人看的草圖，幫助整理思路，編譯器不會讀它。
\item 兩個：「印出 Hello, world!」，然後結束；中間沒有任何判斷。
\end{enumerate}

\paragraph{第~\ref{sec:intro-first-program}~節}
\begin{enumerate}
\item 螢幕上仍然看得到 \texttt{Hello, world!}，但游標不會換到下一
行，下一行的終端機提示符會黏在同一行後面。
\item \texttt{(void)} 明確表示 \texttt{main} 不接收任何參數；省略成
\texttt{main()} 在大多數編譯器裡效果相同，但語意上代表「參數未指
定」，不如 \texttt{(void)} 清楚。
\item \texttt{/* */} 可以跨很多行；\texttt{//} 只到該行結尾為止。
\item 把 \texttt{printf("Hello, world!\textbackslash n");} 那一行複製
一次，變成連續兩條 \texttt{printf} 敘述。
\end{enumerate}

\paragraph{第~\ref{sec:intro-build-pipeline}~節}
\begin{enumerate}
\item 前置處理（展開 \texttt{\#include} 等指令）、編譯（翻譯成目標
碼，語法錯誤在這一步被抓出來）、連結（把目標碼跟函式庫實作接在一
起，產生執行檔）。
\item 請編譯器把常見的可疑寫法用警告講出來，幫忙抓出很多執行期才會
爆出來的問題。
\item \texttt{a.out}。
\item \texttt{./} 表示到目前目錄裡找執行檔，避免跟系統 PATH 裡其他同
名程式搞混。
\end{enumerate}

\paragraph{第~\ref{sec:intro-error-types}~節}
\begin{enumerate}
\item 編譯階段：這是語法錯誤，編譯器會直接報錯，不會產生執行檔。
\item 執行期錯誤：語法過了，但執行到某一刻，行為超出語言保證的範
圍。
\item 邏輯錯誤：語法合法、執行也正常結束，只是輸出的內容跟預期的規
格不一樣（多印了一行）。
\item 警告不會擋下執行檔的產生，但常常是下一類（執行期或邏輯）錯誤
的預告，不應該當成裝飾忽略；錯誤則會讓編譯直接失敗。
\end{enumerate}

\chapter{基本程式架構}
\label{ch:basicprogramstructure}

第一章已經交代程式骨架、\texttt{gcc} 編譯流程，以及 \texttt{Hello,
world!} 怎麼跑起來。 這一章不再重講 \texttt{\#include} 與前置處理／
編譯／連結，直接從輸出與輸入的實務技巧接下去：資料怎麼透過
\texttt{printf} 和 \texttt{scanf} 在程式與外界之間來回傳遞。

這一章依四個階段推進：先把 \texttt{printf} 與換行講透，再一口氣搞懂變數
的宣告與賦值，接著用 \texttt{scanf} 與回音程式確認輸入輸出的管子是通的，
最後才處理多個數字的讀寫、補零格式與 \texttt{double}。 算式本身（加減
乘除、整數除法怎麼截斷、什麼時候該用 \texttt{double} 做計算）留到下一
章運算再集中處理；這裡的目標是讓讀者能獨立寫出讀一些整數或浮點數、依
題目要求印出結果的程式骨架。

\section{用 printf 輸出文字}
\label{sec:basic-printf}

要印出文字，我們在 \texttt{printf} 的格式字串裡寫想輸出的內容；一行一
個 \texttt{printf}，就會一行一行印到螢幕上。

\begin{code}{c}{連續印出多行文字}
printf("first line\n");
printf("second line\n");
printf("third line\n");
\end{code}

格式字串用雙引號包起來；引號裡的內容會原樣輸出（除了 \texttt{\textbackslash n}
這類跳脫字元）。 如果題目要求印出固定的英文句子、而且順序與換行都已
經規定好，通常不需要變數，也不需要 \texttt{scanf}。 直接把句子寫進
\texttt{printf} 即可。

範例 \thechapter.2 依序印出三句英文；每一句各自一行，中間沒有多餘的
空白行。

\begin{example}{c}{(three-lines.c) 印出三句固定的英文}
#include <stdio.h>

int main(void) {
    printf("This is a book.\n");
    printf("It teaches C programming.\n");
    printf("Read it carefully.\n");
}
\end{example}

\iosample{}{
This is a book. \\
It teaches C programming. \\
Read it carefully. \\
}

這種純輸出、無輸入的題目，檢查的重點往往是：大小寫、標點、以及每
一行末尾有沒有照規定換行。 寫完之後用 \texttt{diff} 跟預期輸出比對，
比用肉眼掃過更可靠。

\texttt{printf} 印完一段文字之後，游標停在那段文字的右邊。 若立刻再
印下一段，兩段會黏成一長串，中間沒有空白、也沒有換行。 範例
\thechapter.3 連續印 \texttt{"17"} 和 \texttt{"29"}，螢幕上看到的是
\texttt{1729}。 這只是兩次輸出拼在一起，跟數學上的 $1729$ 無關。

\begin{example}{c}{(print-glued.c) 兩次 printf 中間沒有換行}
#include <stdio.h>

int main(void) {
    printf("17");
    printf("29");
    return 0;
}
\end{example}

\iosample{}{
1729 \\
}

補上 \texttt{\textbackslash n} 之後，每次印完都會把游標移到下一行，
如範例 \thechapter.4。 這跟上一章印 \texttt{Hello, world!} 是同一
件事。

\begin{example}{c}{(print-two-lines.c) 每段文字印完後換行}
#include <stdio.h>

int main(void) {
    printf("17\n");
    printf("29\n");
    return 0;
}
\end{example}

\iosample{}{
17 \\
29 \\
}

若希望兩段文字出現在同一行、中間用空白隔開，格式字串應寫成
\texttt{"17 29\textbackslash n"}，空白是你自己放進字串裡的字元，編
譯器不會自動插入。 題目範例輸出有幾個空白、最後有沒有換行，都要對
著抄，不要看起來差不多就過關。

\stylepoint{固定字串輸出題不需要宣告變數；把要求的文字原樣放進
\texttt{printf}，並確認 \texttt{\textbackslash n} 的位置跟題目一致。
需要換行就自己寫 \texttt{\textbackslash n}；螢幕上文字黏在一起，多半
是上一次 \texttt{printf} 沒把游標移走。}

\begin{drills}[label=drills:basic-printf]
\item 連續執行 \texttt{printf("17");} 和 \texttt{printf("29");}（中間
沒有 \texttt{\textbackslash n}），螢幕上會看到什麼？
\item 要讓 \texttt{"17"} 和 \texttt{"29"} 印在同一行、中間隔一個空
白，格式字串該怎麼寫？
\item 範例 \thechapter.4（\texttt{print-two-lines.c}）印出兩行
\texttt{17} 與 \texttt{29}，靠的是哪個字元？
\end{drills}

\section{變數與 int 型別}
\label{sec:basic-int-vars}

一旦題目說讀入一個整數，我們就需要地方暫存讀進來的值。
\emph{變數}（variable）就是程式執行期間在記憶體裡保留的一格（或多格）
空間，我們用名字（例如 \texttt{count}、\texttt{year}）存取它，型別
（例如 \texttt{int}）決定這一格有多大、能存什麼範圍的整數。

\begin{code}{c}{宣告整數變數}
int x;
\end{code}

\texttt{int x;} 向編譯器申請一個名為 \texttt{x} 的整數變數。 在
\texttt{main} 裡宣告的變數，生命週期只到 \texttt{main} 結束為止。 它
們通常配置在\emph{堆疊}（stack）上：函式被呼叫時，區域變數依序在堆
疊上開闢空間；函式返回時，這些空間一併回收。 現階段只要記住：變數名
是給人讀的，編譯器真正關心的是在記憶體的哪裡、佔幾個位元組。

每個變數可以想成同時帶著四件事：\emph{型別}（這一格能放什麼）、
\emph{名字}（我們怎麼稱呼它）、\emph{值}（這一格目前存的內容）、以
及\emph{位址}（它在記憶體裡的編號）。 \texttt{scanf} 需要的是位址，
所以才會寫 \texttt{\&x}；\texttt{printf("\%d", x)} 需要的是值，所以
寫 \texttt{x} 就夠。 把這兩種用途混在一起，是這一章後半最常見的錯。

名字不是隨便取的。 C 的\emph{識別字}（identifier）只能由英文字母、
數字和底線組成，而且不能以數字開頭；大小寫視為不同，所以
\texttt{sum} 和 \texttt{Sum} 是兩個變數。 有些單字被語言保留當關鍵
字（\texttt{int}、\texttt{return}、\texttt{if}），不能拿來當變數
名。 取名時優先用看得懂的字：\texttt{score}、\texttt{year} 比
\texttt{a1}、\texttt{tmp2} 更容易在除錯時對上號。 C 程式常見兩種風
格：小寫加底線 \texttt{student\_score}，或駝峰式 \texttt{studentScore}；
兩種都可以，重點是同一支程式裡保持一致，不要讓 \texttt{a}、
\texttt{b}、\texttt{c} 滿天飛。

\begin{center}
\begin{tikzpicture}
    \node[cell] (v) {?};
    \node[idxlbl, left=3mm of v] (name) {\texttt{x}};
    \node[idxlbl, right=3mm of v, align=left] (note) {尚未賦值};
\end{tikzpicture}
\diagcaption{變數 \texttt{x} 在概念上可以想成記憶體裡的一格；宣告後、
\texttt{scanf} 或賦值之前，內容是未定義的，不該拿來做計算。}
\label{fig:ch02-variable-box}
\end{center}

\keypoint{\texttt{int} 用來存整數；變數必須先宣告型別和名字，才能被
\texttt{scanf} 寫入或在運算式中使用。}

\stylepoint{變數名要讓人看得懂用途；\texttt{student\_score} 或
\texttt{studentScore} 都比 \texttt{a} 好維護。同一支程式裡選一種命名
風格並保持一致。}

\begin{drills}[label=drills:basic-int-vars]
\item 一個變數同時帶著哪四件事？ \texttt{scanf} 需要用到其中哪一個？
\item \texttt{int x;} 宣告完、還沒賦值之前，\texttt{x} 裡面是什麼？ 這
時候可以直接拿去 \texttt{printf} 嗎？
\item \texttt{Sum} 和 \texttt{sum} 是不是同一個變數？ 為什麼？
\end{drills}

\section{初始化與賦值}
\label{sec:basic-init-assign}

宣告變數之後，我們得把值寫進去，程式才有明確的狀態可以運作。 一個好
習慣是在宣告的當下就指定初值，讓讀程式的人一眼看出起始狀態，也避免
沒賦值就拿來計算的模糊情況。 在較複雜的程式裡，每個變數一開始都有
確定的值，除錯時也比較少追蹤到莫名其妙的殘留內容。

\begin{code}{c}{宣告並初始化}
int count = 0;
int year = 2024;
\end{code}

這段程式一邊宣告變數 \texttt{count}，一邊把它設為 $0$；\texttt{year}
則一開始就有 $2024$ 這個明確的值。 等號在這裡代表賦值：把右邊的值
寫進左邊的變數，可以想成一支向左的箭頭。 \texttt{int n = 42;} 就是
在宣告的當下就給初值；也可以先宣告、稍後再用賦值敘述改它的內容。

\begin{code}{c}{先宣告，再把值寫進變數}
int n;
n = 42;
\end{code}

如果有幾個用途相近的變數（例如同樣是迴圈用的索引、或同一組座標），
可以在同一行內一起宣告，只要型別一致即可：

\begin{code}{c}{一行宣告多個變數}
int i, j, k;
\end{code}

也可以在宣告時一併給初值：

\begin{code}{c}{宣告並初始化多個變數}
int a = 1, b = 2, c = 3;
\end{code}

範例 \thechapter.10 把宣告與初值合成一個完整程式：宣告 \texttt{answer}、
立刻給初值、再用 \texttt{\%d} 印出來。 \texttt{\%d} 的 \texttt{d} 表
示十進位（decimal）；後面那個參數必須是整數值，不是位址。

\begin{example}{c}{(show-answer.c) 設定整數並印出}
#include <stdio.h>

int main(void) {
    int answer = 42;
    printf("%d\n", answer);
    return 0;
}
\end{example}

\iosample{}{
42 \\
}

同一個意思也可以拆成先空宣告、再賦值，如範例 \thechapter.11。 執行
結果一樣是 $42$。 兩種寫法編譯器都收；初值在宣告時就寫下，比較
不容易忘記。 若宣告了卻既不初始化、也不賦值，就拿去 \texttt{printf}，
那一格裡可能是殘留的任意位元；後面會看到編譯器對這件事的警告。

\begin{example}{c}{(assign-later.c) 宣告與賦值分成兩步}
#include <stdio.h>

int main(void) {
    int n;
    n = 42;
    printf("%d\n", n);
    return 0;
}
\end{example}

\iosample{}{
42 \\
}

在這台機器上，\texttt{sizeof(int)} 的結果是 $4$，也就是一個
\texttt{int} 佔四個位元組（三十二個位元）。 這不是語言保證的唯一可
能；別的平臺上 \texttt{int} 可能更寬或更窄。 現階段只要有這個數量級
的感覺：它不是無限大的整數，太大的數字之後會在運算那一章碰到溢
位。

\keypoint{賦值是把右邊的值寫進左邊的變數。 宣告時給初值，或宣告
後再賦值，都可以；拿尚未寫入的變數去做計算，結果沒有保證。}

\stylepoint{相似用途的變數可以放在一起宣告，增加可讀性。寫作時盡量讓
名字和分組呈現變數之間的邏輯關係。}

\begin{drills}[label=drills:basic-init-assign]
\item \texttt{int n; n = 42;} 跟 \texttt{int n = 42;}，兩者最後
\texttt{n} 的值一樣嗎？
\item 宣告了 \texttt{int x;} 卻沒有初始化、也沒有賦值，就直接拿去
\texttt{printf}，這樣安全嗎？
\item 把 \texttt{int a, b, c;} 改成一行內就給三個初值，該怎麼寫？
\end{drills}

\section{用 scanf 讀入整數}
\label{sec:basic-scanf}

C 語言用 \texttt{scanf} 從標準輸入讀資料。 讀一個整數的典型寫法是：

\begin{code}{c}{讀入一個 int 並存進變數 x}
int x;
scanf("%d", &x);
\end{code}

格式字串 \texttt{"\%d"} 表示期待一個十進位整數；第二個參數必須是
變數的\emph{位址}，所以變數名前面要加 \texttt{\&}（取址運算子）。 可以
把 \texttt{\&x} 理解成請把讀到的值寫進 \texttt{x} 在記憶體裡的那一格。
\texttt{scanf} 需要知道寫到哪裡，光給變數名而不給位址是不夠的。

上一節學 \texttt{printf} 時，我們在格式字串裡用 \texttt{\textbackslash n}
控制換行。 讀入時不要照抄這個習慣：\texttt{scanf("\%d\textbackslash n",
\&n)} 是常見的地雷。 格式字串裡的空白字元（包含 \texttt{\textbackslash n}）
會一直跳過輸入裡的空白，直到看見下一個非空白字元為止；若你在讀完整數
後多寫了 \texttt{\textbackslash n}，程式可能看起來像卡住，等你在鍵盤
上多按一次才繼續。 讀整數用 \texttt{"\%d"} 就夠；換行留給 \texttt{printf}
去處理。

\keypoint{\texttt{scanf} 讀 \texttt{int} 時，格式用 \texttt{\%d}，參
數必須是 \texttt{\&變數名}；漏寫 \texttt{\&} 是初學者最常見、也最危險
的錯誤之一，可能導致程式寫壞記憶體。格式字串裡不要加 \texttt{\textbackslash n}。}

\begin{drills}[label=drills:basic-scanf]
\item 讀一個 \texttt{int} 變數 \texttt{n}，\texttt{scanf} 的格式字串
該寫成什麼？參數該怎麼寫？
\item \texttt{scanf("\%d\textbackslash n", \&n)} 這樣寫，可能會讓程式
看起來像卡住，為什麼？
\item \texttt{scanf} 的第二個參數為什麼一定要是變數的位址（加
\texttt{\&}），不能直接寫變數名？
\end{drills}

\section{回音程式：讀進來再原樣印出}
\label{sec:basic-echo}

真正解題時，數字來自標準輸入。 最小的完整迴路是：宣告、\texttt{scanf}
寫入、\texttt{printf} 讀出。 範例 \thechapter.13 什麼都不算，只把讀
到的整數印回去。 看起來無聊，卻最能確認輸入輸出這條管子是通的。 若
連回音都印錯，問題在格式字串或 \texttt{\&}，不必先懷疑後面的算式。

\begin{example}{c}{(echo-int.c) 讀入一個整數再印出}
#include <stdio.h>

int main(void) {
    int n;
    scanf("%d", &n);
    printf("%d\n", n);
    return 0;
}
\end{example}

\iosample{
100 \\
}{
100 \\
}

執行順序必須是：先有變數、再讀入、最後印出。 顛倒成先 \texttt{printf}
再 \texttt{scanf}，印出來的會是尚未寫入的內容；先 \texttt{scanf} 卻
漏掉 \texttt{\&}，則是請 \texttt{scanf} 把整數寫到一個不像位址的地
方。

管子通了之後，就可以接上一點簡單的計算。 範例 \thechapter.14 讀入兩
個整數，把它們相加後印出。 這是讀入 $\to$ 計算 $\to$ 輸出三段式程
式的最小完整版本：中間的加法本身很簡單，重點在於 \texttt{scanf} 與
\texttt{printf} 的配對是否正確；算式的細節留到下一章再展開。

\begin{example}{c}{(sum-two.c) 讀入兩個整數並輸出總和}
#include <stdio.h>

int main(void) {
    int a, b;
    scanf("%d %d", &a, &b);

    printf("%d\n", a + b);
}
\end{example}

\iosample{
3 7 \\
}{
10 \\
}

以這份輸入為例：\texttt{scanf} 依空白分隔讀到 \texttt{a = 3}、
\texttt{b = 7}，\texttt{printf} 印出 \texttt{a + b = 10}。 格式字串
\texttt{"\%d\textbackslash n"} 裡的 \texttt{\%d} 會被替換成整數值；
我們只放一個 \texttt{\%d}，所以輸出只有數字本身加換行。

我們在這台機器上用 \texttt{gcc -Wall} 編譯過一份漏寫 \texttt{\&}
的程式：\texttt{scanf("\%d", n)}。 編譯器提出兩則警告，大意是格式
期待 \texttt{int *}（指向整數的位址），實際卻收到 \texttt{int}，而
且 \texttt{n} 在呼叫當下還未初始化。 警告沒有阻止產生執行檔，但這
種程式執行期可能寫壞別的記憶體，屬於上一章說的執行期錯誤的近親。
\textbf{看到格式與型別不符的警告，先改到警告消失，再談答案對不對。}

同一個編譯器對宣告了 \texttt{int n;} 卻直接 \texttt{printf("\%d",
n)} 也會警告 \texttt{variable 'n' is uninitialized when used
here}，並建議寫 \texttt{= 0}。 初始化成 $0$ 能讓警告消失，但不表示
$0$ 就是題目要的值。 它只是一個確定的值。 該從鍵盤讀的，還是得
\texttt{scanf}。

\keypoint{\texttt{printf} 用變數的值，\texttt{scanf} 用變數的位址。
\texttt{-Wall} 對漏寫 \texttt{\&}、以及使用未初始化變數，通常會先
喊一聲；那一聲值得當成錯誤來修。}

\begin{drills}[label=drills:basic-echo]
\item 為什麼要先宣告變數、再 \texttt{scanf}、最後才 \texttt{printf}，
順序顛倒會怎樣？
\item 範例 \thechapter.14（\texttt{sum-two.c}）輸入 \texttt{3 7}，會
印出什麼？
\item \texttt{scanf("\%d", n)}（漏寫 \texttt{\&}）編譯時，\texttt{gcc
-Wall} 通常會給什麼樣的警告？
\end{drills}

\section{讀取多個值與調整輸出格式}
\label{sec:basic-multi-read}

題目有時要求分多行讀入，或在同一行讀入多個數字但輸出順序不同。 兩種
常見片語如下：

\begin{code}{c}{同一行讀入兩個整數}
scanf("%d %d", &a, &b);
\end{code}

\begin{code}{c}{分三次各讀一個整數}
scanf("%d", &a);
scanf("%d", &b);
scanf("%d", &c);
\end{code}

同一行的 \texttt{"\%d \%d"} 會跳過中間的空白字元；分三次 \texttt{scanf}
則是每次讀到下一個整數為止。 兩者都能處理一行多個數字或每行一
個數字的輸入，但程式可讀性上，若題目明確說一行兩個整數，用一
次 \texttt{scanf} 讀齊通常比較清楚。

範例 \thechapter.17 讀入兩個整數，但輸出時先印第二個、再印第一個。 
只需要交換 \texttt{printf} 裡的變數順序，不需要真的交換變數的值。

\begin{example}{c}{(swap-two.c) 讀入兩個整數並以相反順序輸出}
#include <stdio.h>

int main(void) {
    int a, b;
    scanf("%d %d", &a, &b);

    printf("%d %d\n", b, a);
}
\end{example}

\iosample{
3 7 \\
}{
7 3 \\
}

範例 \thechapter.18 則示範分三次讀入三個整數：每次 \texttt{scanf} 各處
理一個 \texttt{\%d}，適合題目把三個數字分成三行給定的情況。

\begin{example}{c}{(read-three-sum.c) 分三次讀入三個整數}
#include <stdio.h>

int main(void) {
    int a, b, c;
    scanf("%d", &a);
    scanf("%d", &b);
    scanf("%d", &c);

    printf("%d\n", a + b + c);
}
\end{example}

\iosample{
1 \\
2 \\
3 \\
}{
6 \\
}

\stylepoint{輸出多個以空白分隔的整數時，格式字串寫成
\texttt{"\%d \%d\textbackslash n"}；最後一個數字後面要不要空白，要
對照題目範例輸出，不要自己多加。}

\begin{drills}[label=drills:basic-multi-read]
\item 用一次 \texttt{scanf("\%d \%d", \&a, \&b)} 跟分兩次
\texttt{scanf("\%d", \&a); scanf("\%d", \&b);}，讀到的結果一樣嗎？
\item 範例 \thechapter.17（\texttt{swap-two.c}）輸入 \texttt{3 7}，
輸出是什麼？靠的是交換變數的值，還是交換 \texttt{printf} 裡的順序？
\item 若題目把三個整數分三行給，該用一次 \texttt{scanf} 還是三次？
\end{drills}

\section{printf 的寬度與補零}
\label{sec:basic-width-padding}

除了 \texttt{\%d}，\texttt{printf} 還可以在 \texttt{\%} 與轉換字元之間
插入\emph{寬度}，控制至少印出幾個字元；寬度前面再加 \texttt{0}，不足
的部分會用前導零補齊，而不是空白。

\begin{code}{c}{將整數以至少 5 位數輸出，不足補零}
printf("%05d\n", id);
\end{code}

\texttt{\%05d} 裡的 \texttt{5} 是最小欄寬，\texttt{0} 表示用零補位。
例如 \texttt{id = 42} 只佔兩位，輸出會變成 \texttt{00042}。 若數字本
身已超過五位（例如 \texttt{123456}），不會被截斷，會完整印出六位。

範例 \thechapter.20 讀入一個學號或編號，依題目要求補零成五位。

\begin{example}{c}{(padded-id.c) 將整數補零成五位輸出}
#include <stdio.h>

int main(void) {
    int id;
    scanf("%d", &id);

    printf("%05d\n", id);
}
\end{example}

\iosample{
42 \\
}{
00042 \\
}

這種格式控制在處理固定長度的數字編碼時很常見。 例如日期裡的月
分、流水號。 算式怎麼從整數裡拆出各位數字，則屬於下一章的除法與餘數
技巧；這裡只處理已知一個整數，要怎麼漂亮地印出來。

\keypoint{\texttt{\%05d} 的 \texttt{0} 是補零旗標，\texttt{5} 是
最小寬度；兩者缺一不可，\texttt{\%5d} 只會用空白補位。}

\begin{drills}[label=drills:basic-width-padding]
\item \texttt{printf("\%05d\textbackslash n", 42);} 會印出什麼？
\item 若 \texttt{id = 123456}（超過 5 位數），\texttt{\%05d} 會把它
截斷成 5 位嗎？
\item \texttt{\%5d}（沒有前導 \texttt{0}）跟 \texttt{\%05d}，不足位數
時分別用什麼字元補齊？
\end{drills}

\section{double 型別與浮點輸入輸出}
\label{sec:basic-double}

有些輸出需要小數點（例如金額），這時候整數型別不夠用，要改用
\texttt{double}（雙精度浮點數）。 讀寫 \texttt{double} 時，格式字串
跟 \texttt{int} 不同：

\begin{code}{c}{讀入 double 並印到小數點後兩位}
double x;
scanf("%lf", &x);
printf("%.2f\n", x);
\end{code}

\texttt{scanf} 讀 \texttt{double} 必須用 \texttt{\%lf}（這裡的
\texttt{l} 代表 \texttt{long double} 家族，對 \texttt{double} 就是正
確的格式）；\texttt{printf} 印 \texttt{double} 則用 \texttt{\%f} 就
夠。 這是 C 語言留下來的不對稱寫法，之後浮點數那一章還會再從記憶體
表示的角度細講。

範例 \thechapter.22 讀入一個浮點數，四捨五入到小數點後兩位輸出。 這
裡的重點是 \texttt{\%lf}/\texttt{\%.2f} 的配對；至於整數分鐘換算
成含小數的元這類題目，分子該除以 \texttt{100.0} 還是 \texttt{100}，
則留到下一章討論型別轉換。

\begin{example}{c}{(read-double-print.c) 讀入浮點數並印出兩位小數}
#include <stdio.h>

int main(void) {
    double price;
    scanf("%lf", &price);

    printf("%.2f\n", price);
}
\end{example}

\iosample{
3.14159 \\
}{
3.14 \\
}

\texttt{\%.2f} 裡的 \texttt{.2} 表示小數位數；編譯器會做四捨五入。 若
不指定小數位數，\texttt{\%f} 預設常印出六位小數，對很多題目來說既冗
長又可能暴露浮點誤差。 養成依題目要求寫 \texttt{\%.Nf} 的習慣，可以少
掉不少格式不符的錯誤。

\stylepoint{讀 \texttt{double} 用 \texttt{\%lf}，印 \texttt{double} 用
\texttt{\%f} 或 \texttt{\%.Nf}；不要把 \texttt{scanf} 的格式直接套到
\texttt{printf} 上。}

C 對標點幾乎沒有容錯：少一個分號、括號不成對、字串少關一個直引號，
編譯器就無法確定你想表達什麼。 上一章的 \texttt{Hello, world!} 已經
示範過分號；這一章開始，每一個 \texttt{\%}、\texttt{\&}、逗號、雙引
號都各自有職務，不要把它們當成裝飾。 下一章會把加、減、乘、除、以
及整數除法的截斷，接到這個骨架上。 輸入輸出的管子先通，算式才有
地方可放。

\begin{drills}[label=drills:basic-double]
\item \texttt{scanf} 讀 \texttt{double} 用什麼格式？\texttt{printf}
印 \texttt{double} 用什麼格式？兩者一樣嗎？
\item \texttt{printf("\%.2f\textbackslash n", 3.14159);} 會印出什麼？
\item 若不指定小數位數，直接用 \texttt{\%f}，通常會印出幾位小數？
\end{drills}

\addcontentsline{toc}{section}{練習解答}
\section*{練習解答}

各節練習的答案收在這裡。 短題只給能一句話核對的結果。

\paragraph{第~\ref{sec:basic-printf}~節}
\begin{enumerate}
\item \texttt{1729}：兩次輸出直接黏在一起，中間沒有空白也沒有換行。
\item \texttt{"17 29\textbackslash n"}——空白要自己放進格式字串裡，
編譯器不會自動加。
\item \texttt{\textbackslash n}（換行字元）。
\end{enumerate}

\paragraph{第~\ref{sec:basic-int-vars}~節}
\begin{enumerate}
\item 型別、名字、值、位址；\texttt{scanf} 需要的是位址，所以要寫
\texttt{\&x}。
\item 內容不確定（殘留值）；不該直接拿去 \texttt{printf}，結果沒有保
證。
\item 不是同一個變數。 C 的識別字大小寫視為不同。
\end{enumerate}

\paragraph{第~\ref{sec:basic-init-assign}~節}
\begin{enumerate}
\item 一樣，最後 \texttt{n} 都是 \texttt{42}。
\item 不安全。 未初始化、未賦值的變數內容是不確定的殘留值，不該直接
使用。
\item \texttt{int a = 1, b = 2, c = 3;}
\end{enumerate}

\paragraph{第~\ref{sec:basic-scanf}~節}
\begin{enumerate}
\item \texttt{scanf("\%d", \&n);}
\item 因為 \texttt{scanf} 格式字串裡的空白（含 \texttt{\textbackslash
n}）會一直跳過輸入裡的空白字元，直到遇到下一個非空白字元；多打的
\texttt{\textbackslash n} 會讓程式等使用者多輸入一次才繼續。
\item 因為 \texttt{scanf} 需要知道要把讀到的值\emph{寫到哪裡}（位
址），只給變數名（值）它無從得知該寫進哪一格記憶體。
\end{enumerate}

\paragraph{第~\ref{sec:basic-echo}~節}
\begin{enumerate}
\item 因為要先有變數才能讀入，要先讀入才有值可印；顛倒的話，
\texttt{printf} 會印出尚未寫入的內容，或 \texttt{scanf} 沒有地方可
存。
\item \texttt{10}（\texttt{3 + 7}）。
\item 格式期待 \texttt{int *}（位址）卻收到 \texttt{int}（值），並提
醒 \texttt{n} 在呼叫時尚未初始化。
\end{enumerate}

\paragraph{第~\ref{sec:basic-multi-read}~節}
\begin{enumerate}
\item 一樣：兩種寫法都能把兩個以空白（或換行）分隔的整數分別讀進
\texttt{a}、\texttt{b}。
\item 輸出 \texttt{7 3}；靠的是交換 \texttt{printf} 裡兩個變數的順
序，值本身沒有互換。
\item 分三次 \texttt{scanf("\%d", ...)}，每次讀一個。
\end{enumerate}

\paragraph{第~\ref{sec:basic-width-padding}~節}
\begin{enumerate}
\item \texttt{00042}。
\item 不會截斷，會完整印出 \texttt{123456}（六位）。
\item \texttt{\%5d} 用空白補位；\texttt{\%05d} 用 \texttt{0} 補位。
\end{enumerate}

\paragraph{第~\ref{sec:basic-double}~節}
\begin{enumerate}
\item 不一樣：\texttt{scanf} 讀 \texttt{double} 要用 \texttt{\%lf}，
\texttt{printf} 印則用 \texttt{\%f}（或 \texttt{\%.Nf}）。
\item \texttt{3.14}。
\item 預設通常印出六位小數。
\end{enumerate}

\chapter{運算}
\label{ch:arithmetic}

上一章把程式的骨架、變數、以及 \texttt{printf}/\texttt{scanf} 的用法
建立起來了。 這一章專門處理\emph{數值運算}：怎麼用 \texttt{+}、\texttt{-}、
\texttt{*}、\texttt{/}、\texttt{\%} 組合出題目要的結果，賦值運算子
\texttt{=} 怎麼改變變數裡的值，運算式裡多個運算子先算哪個，以及整數
跟浮點數混在一起時，編譯器會怎麼自動轉型。

很多初學者以為「會寫 \texttt{scanf} 和 \texttt{printf} 就能解題了」，
卻在 \texttt{9 / 5} 變成 \texttt{1}、或 \texttt{7 / 2} 變成 \texttt{3}
而不是 \texttt{3.5} 的地方卡住。 這不是 \texttt{printf} 寫錯，而是對
C 的整數除法與型別規則不熟。 這一章會把這些規則講清楚，並用題庫難度
曲線上常見的應用（時間換算、位數拆解、簡單幾何公式）串起來。 至於
「條件成立才做某件事」的判斷與邏輯運算子，留給下一章。

\section{算術運算子與優先順序}
\label{sec:arith-operators}

C 提供的算術運算子，一次列出來如下：\texttt{+} 加、\texttt{-} 減、
\texttt{*} 乘、\texttt{/} 除、\texttt{\%} 取餘數，以及只需要一個運算
元的一元 \texttt{-}（取相反數）。 這六種是本章、乃至於之後幾乎所有
章節都會不斷用到的基本積木。

\begin{code}{c}{六種基本算術運算}
sum  = i + j;   // 加
diff = i - j;   // 減
prod = i * j;   // 乘
quot = i / j;   // 除
rem  = i % j;   // 餘數
neg  = -i;      // 相反數（一元負號，只有一個運算元）
\end{code}

\texttt{+}、\texttt{-}、\texttt{*} 的行為跟數學課本一致：兩個
\texttt{int} 相加、相減、相乘，結果仍是 \texttt{int}。 \texttt{/} 和
\texttt{\%} 比較特別：當兩個運算元都是整數時，\texttt{/} 執行\emph{整
數除法}，小數部分直接捨去（例如 \texttt{7 / 2} 是 \texttt{3}，不是
\texttt{3.5}）；\texttt{\%} 則是這次除法\emph{剩下的餘數}，\texttt{/}
跟 \texttt{\%} 通常成對出現，一個取「除得盡的部分」，一個取「除不盡
剩下的部分」。 下一節會把整數除法與餘數的規則講透。

這裡要先提醒兩個容易忽略、卻會直接讓程式當掉或算錯的狀況：

\begin{itemize}
  \item \textbf{除數為 \texttt{0}}：如果 \texttt{j} 是 \texttt{0}，
  \texttt{i / j} 和 \texttt{i \% j} 都是\emph{未定義行為}
  （undefined behavior），不是「得到一個奇怪但看起來合理的數字」，
  而是這台機器上通常會讓程式直接收到 \texttt{SIGFPE} 訊號而中止執
  行。 跟很多「靜默算錯」的陷阱不同，除以 \texttt{0} 通常一測就會發
  現；但只要除數是從讀入或計算得來、而不是寫死的常數，寫程式時就要
  留意它有沒有可能是 \texttt{0}。
  \item \textbf{\texttt{\%} 遇到負數}：\texttt{\%} 在負數時的行為，
  常常跟直覺（或國小學的「餘數不能是負的」）不一樣——C 的餘數符號跟
  著被除數走，不是恆為非負。 下一節會用具體數字把規則講清楚。
\end{itemize}

範例 \thechapter.2 讀入兩個整數，輸出第一個減第二個的差。 這是「讀
入 $\to$ 單一運算 $\to$ 輸出」最直接的模板。

\begin{example}{c}{(subtract-two.c) 輸出兩整數的差}
#include <stdio.h>

int main(void) {
    int a, b;
    scanf("%d %d", &a, &b);

    printf("%d\n", a - b);
}
\end{example}

\iosample{
10 3 \\
}{
7 \\
}

範例 \thechapter.3 把乘法用在幾何題：長方形面積是長乘寬。 題目本身
沒有新語法，只是把 \texttt{*} 放進 \texttt{printf} 的運算式裡。

\begin{example}{c}{(rect-area.c) 長方形面積}
#include <stdio.h>

int main(void) {
    int length, width;
    scanf("%d %d", &length, &width);

    printf("%d\n", length * width);
}
\end{example}

\iosample{
4 5 \\
}{
20 \\
}

以輸入 \texttt{4 5} 為例：\texttt{length = 4}、\texttt{width = 5}，
\texttt{length * width = 20}，程式印出 \texttt{20}。 乘法在 C 裡沒有
位元溢位的自動保護。 若兩個很大的 \texttt{int} 相乘超過 \texttt{int}
能表示的範圍，結果會\emph{溢位}（overflow）而變成錯誤的值；題庫目
前給的測資都在安全範圍內，但這個限制之後遇到 \texttt{long long} 時
會再回來談。

同一條運算式裡有多個運算子時，C 語言依\emph{優先順序}（precedence）
決定先算哪一部分。 跟數學課本類似：乘法比加法先算，括號可以強制改
變順序。

\begin{code}{c}{乘法優先於加法；括號可改變順序}
result = a + b * c;      // 先算 b * c，再加 a
result = (a + b) * c;    // 先算 a + b，再乘 c
\end{code}

若不確定，就加括號。 編譯器不會抱怨多餘的括號，但少括號導致順序錯
誤，程式卻可能靜默地算出錯誤答案。 減法和除法還有\emph{結合方向}
（associativity）的問題：同樣優先順序的運算子由左向右算，所以
\texttt{a - b - c} 其實是 \texttt{(a - b) - c}，不是
\texttt{a - (b - c)}。 例如長方體表面積公式
$2(ab + bc + ca)$，若寫成 \texttt{2 * a * b + b * c + c * a}，乘法
會先結合，整個式子的意義就偏離公式了；正確做法是把加法的部分括起
來：\texttt{2 * (a * b + b * c + c * a)}。

本章涉及的主要運算子按優先順序由高到低如下：

\begin{center}
\begin{tabular}{cl}
優先順序 & 運算子 \\
\hline
1（最高） & \texttt{++ --}（後綴、單獨成一行時） \\
2 & \texttt{++ --}（前綴）、一元 \texttt{-} \\
3 & \texttt{* / \%}（乘、除、餘） \\
4 & \texttt{+ -}（加、減） \\
5 & \texttt{=}、\texttt{+= -= *= /= \%=}（賦值、複合指定） \\
\end{tabular}
\end{center}

例如 \texttt{a + b * c} 優先算 \texttt{b * c}（優先順序 3）再加 \texttt{a}
（優先順序 4）。 同一個優先順序的多個運算子按\emph{結合方向}決定：加減
乘除由左向右，賦值與複合指定由右向左。

\stylepoint{寫多項加乘混合的公式時，寧可多寫括號對照題目原文，也不要
依賴「我記得優先順序」。 這在考試時間壓力下特別容易失手。}

範例 \thechapter.4 則是純整數加法：民國年加 \texttt{1911} 換算成西
元。 常數跟變數一樣可以出現在運算式裡。

\begin{example}{c}{(roc-year.c) 民國年換算成西元年}
#include <stdio.h>

int main(void) {
    int roc;
    scanf("%d", &roc);

    printf("%d\n", roc + 1911);
}
\end{example}

\iosample{
100 \\
}{
2011 \\
}

\begin{drills}[label=drills:arith-operators]
\item \texttt{a + b * c} 跟 \texttt{(a + b) * c}，兩者算出來的結果一
定一樣嗎？ 為什麼？
\item 寫出 \texttt{2 * a * b + b * c + c * a} 這個式子哪裡沒有正確表
達出 $2(ab+bc+ca)$，該怎麼加括號修正？
\item \texttt{a - b - c} 是先算 \texttt{a - b} 還是先算 \texttt{b - c}？
這是靠優先順序還是結合方向決定的？
\end{drills}

\section{整數除法與餘數}
\label{sec:arith-intdiv}

C 的 \texttt{/} 在兩個運算元都是整數時，執行\emph{整數除法}：結果
的小數部分直接捨去（對正數而言等同無條件截斷 toward zero）。

\begin{code}{c}{整數除法：商的小數部分被捨去}
int q = 7 / 2;        // q 是 3，不是 3.5
int avg = (a + b + c) / 3;
\end{code}

範例 \thechapter.5 讀入三個整數，輸出算術平均的\emph{整數部分}。 
先加總再除以 \texttt{3}，小數同樣被截斷。

\begin{example}{c}{(three-average.c) 三個整數的整數平均}
#include <stdio.h>

int main(void) {
    int a, b, c;
    scanf("%d %d %d", &a, &b, &c);

    printf("%d\n", (a + b + c) / 3);
}
\end{example}

\iosample{
1 2 3 \\
}{
2 \\
}

\texttt{(1 + 2 + 3) / 3 = 6 / 3 = 2}，剛好整除。 若輸入是 \texttt{1 2 4}，
和為 \texttt{7}，\texttt{7 / 3 = 2}（餘 \texttt{1} 被捨去），不是
\texttt{2.333...}。

範例 \thechapter.6 用攝氏轉華氏再反推的整數版公式：$C = (F - 32) \times
5 / 9$。 括號確保先算 \texttt{F - 32}，再乘 \texttt{5}，最後除以
\texttt{9}；每一步都是整數運算，最終結果同樣截斷。

\begin{example}{c}{(f-to-c.c) 華氏轉攝氏（整數版）}
#include <stdio.h>

int main(void) {
    int f;
    scanf("%d", &f);

    printf("%d\n", (f - 32) * 5 / 9);
}
\end{example}

\iosample{
32 \\
212 \\
}{
0 \\
100 \\
}

\texttt{F = 32} 時，\texttt{(32 - 32) * 5 / 9 = 0}；\texttt{F = 212}
時，\texttt{(212 - 32) * 5 / 9 = 180 * 5 / 9 = 900 / 9 = 100}。 若
想要四捨五入而不是截斷，通常要在分子上加調整（例如先加 \texttt{4}
再除）。 那是較進階的技巧，現階段題庫多半接受截斷版。

\texttt{\%} 是\emph{餘數}（remainder）運算子。 在 C99 標準中（本書採用
的標準），\texttt{a \% b} 遵循「向 0 取整」的規則：

$$a = b \times (a / b) + (a \bmod{b})$$

當都是正數時，結果符合日常的「餘數」概念：\texttt{7 \% 3 = 1}（7 除以 3
商 2 餘 1）。 但當涉及負數時，\textbf{餘數的符號跟著被除數（除號前那個數
字），不跟著除數}。 例如：
\begin{itemize}
  \item \texttt{-7 \% 3 = -1}（不是 2！）
  \item \texttt{7 \% -3 = 1}（不是 -2！）
\end{itemize}

簡單記住：C 語言的 \texttt{\%} 與 \texttt{/} 配對，一起向 0 取整。

\begin{code}{c}{商與餘數：總分鐘換算成幾小時又幾分}
int hours = total / 60;         // 向 0 取整，得完整小時數
int minutes = total % 60;       // 剩餘分鐘
\end{code}

範例 \thechapter.7 讀入總分鐘數，輸出「幾小時又幾分鐘」。 \texttt{/}
得到完整的小時數，\texttt{\%} 得到不足一小時的剩餘分鐘。 這是時間換
算題最經典的片語。

\begin{example}{c}{(split-minutes.c) 總分鐘換成小時與分鐘}
#include <stdio.h>

int main(void) {
    int total;
    scanf("%d", &total);

    printf("%d %d\n", total / 60, total % 60);
}
\end{example}

\iosample{
135 \\
}{
2 15 \\
}

\texttt{135 / 60 = 2}、\texttt{135 \% 60 = 15}，輸出 \texttt{2 15}。
反過來「小時加分鐘換總分鐘」則用 \texttt{hours * 60 + minutes}，綜
合應用一節的範例 \thechapter.17 示範這個方向。

\texttt{-} 除了當減法，也可以當\emph{一元}運算子：\texttt{-n} 表示
「\texttt{n} 的相反數」，只需要一個運算元。 它跟減法長得一樣，所以
更要靠空白與括號讓人一眼看出你要的是哪一種。

範例 \thechapter.8 讀入一個整數（測資用 \texttt{-17}），依序印出相
反數、除以 $5$ 的商、以及除以 $5$ 的餘數。 在這台機器上觀察到：
\texttt{-(-17) = 17}，\texttt{-17 / 5 = -3}，\texttt{-17 \% 5 = -2}。
也就是說，C99 的整數除法往零的方向截斷（$-3.4$ 變成 $-3$，不是
$-4$），餘數的符號跟被除數相同。 同一個程式若輸入 \texttt{17}，則
商是 $3$、餘數是 $2$。 正數那一側比較符合小學的「有餘數除法」。

\begin{example}{c}{(unary-and-neg.c) 相反數、負數除法與餘數}
#include <stdio.h>

int main(void) {
    int n;
    scanf("%d", &n);
    printf("%d\n", -n);
    printf("%d\n", n / 5);
    printf("%d\n", n % 5);
    return 0;
}
\end{example}

\iosample{
-17 \\
}{
17 \\
-3 \\
-2 \\
}

\textbf{奇偶與餘數}：對正整數，\texttt{n \% 2} 只會是 \texttt{0} 或
\texttt{1}，分別對應偶數與奇數。 但負奇數的餘數是 \texttt{-1}，例
如 \texttt{-3 \% 2 = -1}，\texttt{-5 \% 2 = -1}。 因此「餘數等於
\texttt{1}」並不能當成奇數的判準；下一章教判斷時會用 \texttt{n \% 2 != 0}
來涵蓋負數。 題庫目前多半只在非負整數範圍出題，但語言本身允許負數
運算，自己寫程式時不要假設餘數一定是非負整數。

\keypoint{兩個 \texttt{int} 做 \texttt{/}，結果仍是 \texttt{int}，小
數部分直接捨去；\texttt{a / b} 與 \texttt{a \% b} 配對使用，可以從
一個總量拆出「完整單位數」與「剩餘部分」。 在 C99 中，兩者都向 0 取
整；當有負數時，餘數的符號跟著被除數。}

\begin{drills}[label=drills:arith-intdiv]
\item \texttt{-7 \% 3} 等於多少？ 為什麼不是 \texttt{2}？
\item 讀入總分鐘數 \texttt{125}，用 \texttt{/ 60} 和 \texttt{\% 60} 算
出來的「幾小時幾分」分別是多少？
\item \texttt{n \% 2 == 1} 能不能拿來判斷 \texttt{n} 是奇數？ 對所有
整數都成立嗎？
\end{drills}

\section{賦值、更新與交換}
\label{sec:arith-assign-update-swap}

上一章已經介紹過 \texttt{=} 是\emph{賦值}：把右邊運算式的值送進左邊
的變數，方向由右向左。 這裡補充幾個在運算章節裡常碰到的用法。

指定會\emph{蓋掉}左邊變數的舊值。 範例 \thechapter.9 先印出 \texttt{left}
讀進來的值，再執行 \texttt{left = right}，然後再印一次。 第二次印出
的是 \texttt{right} 的拷貝；\texttt{right} 自己並沒有被改掉，但
\texttt{left} 原本的數字已經不在了。

\begin{example}{c}{(copy-over.c) 指定之後左邊的舊值消失}
#include <stdio.h>

int main(void) {
    int left, right;
    scanf("%d %d", &left, &right);
    printf("%d\n", left);
    left = right;
    printf("%d\n", left);
    return 0;
}
\end{example}

\iosample{
8 21 \\
}{
8 \\
21 \\
}

若目標是讓兩個變數\emph{互換}內容，直接寫 \texttt{a = b;} 再
\texttt{b = a;} 會失敗：第一句已經把 \texttt{a} 蓋掉，第二句只是把同
一個值再抄回去。 正確作法是先把其中一個值存進暫存變數。

\begin{code}{c}{用暫存變數交換 a 與 b}
temp = a;
a = b;
b = temp;
\end{code}

範例 \thechapter.10 用暫存變數交換兩個整數。

\begin{example}{c}{(swap-temp.c) 交換兩個整數}
#include <stdio.h>

int main(void) {
    int a, b;
    scanf("%d %d", &a, &b);
    int temp = a;
    a = b;
    b = temp;
    printf("%d %d\n", a, b);
    return 0;
}
\end{example}

\iosample{
8 21 \\
}{
21 8 \\
}

指定敘述本身也可以當成運算式，值就是指定完成後左邊那個變數的新值。
因為 \texttt{=} 由右向左結合，\texttt{c = b = a} 等價於
\texttt{c = (b = a)}：先把 \texttt{a} 抄進 \texttt{b}，再把 \texttt{b}
的新值抄進 \texttt{c}。

\begin{code}{c}{連鎖指定，由右向左}
c = b = a;
\end{code}

範例 \thechapter.11 讀入一個數字，三個變數最後印出來會一樣。

\begin{example}{c}{(chain-assign.c) 一個值指定給三個變數}
#include <stdio.h>

int main(void) {
    int a = 1, b = 2, c = 3;
    scanf("%d", &a);
    c = b = a;
    printf("%d %d %d\n", a, b, c);
    return 0;
}
\end{example}

\iosample{
14 \\
}{
14 14 14 \\
}

「把變數加上（或減去）一個量再存回去」太常見，C 提供了較短的寫法。
\texttt{n++} 讓 \texttt{n} 的值加一；\texttt{n--} 讓 \texttt{n} 的值減一；
\texttt{n += 10} 等價於 \texttt{n = n + 10}。 還有 \texttt{-=}、\texttt{*=}、
\texttt{/=}、\texttt{\%=}，規則都是「先用舊值做運算，再寫回同一個變數」。

\begin{code}{c}{加一、減一，以及複合指定}
n++;
n--;
n += 10;
n -= 5;
n *= 2;
n /= 3;
n %= 7;
\end{code}

範例 \thechapter.12 讀入一個整數，先加一印出，再減 $3$ 印出，再加
$10$ 印出。 輸入 \texttt{5} 時，先看到 \texttt{6}，再看到 \texttt{3}，
最後看到 \texttt{13}。

\begin{example}{c}{(count-up.c) 加一、減一與複合指定}
#include <stdio.h>

int main(void) {
    int n;
    scanf("%d", &n);
    n++;
    printf("%d\n", n);
    n -= 3;
    printf("%d\n", n);
    n += 10;
    printf("%d\n", n);
    return 0;
}
\end{example}

\iosample{
5 \\
}{
6 \\
3 \\
13 \\
}

\texttt{++} 與 \texttt{--} 還有寫在變數前面（\texttt{++n} 或 \texttt{--n}）
與寫在後面（\texttt{n++} 或 \texttt{n--}）的差別：單獨成一行時兩者都只
是加一或減一；若把它嵌進更大的運算式，前綴是「先加/減再取值」，後綴
是「先取值再加/減」。 現階段請先養成「加一或減一自己寫成一行」的習慣，
不要把 \texttt{++} 或 \texttt{--} 塞進 \texttt{printf} 參數裡。

下面這段是\emph{不良示範}：同一個運算式裡對 \texttt{a} 做兩次更新，
結果是\emph{未定義行為}，印出來的數字因編譯器而異，甚至可能編譯時
就收到警告。

\begin{code}{c}{不良示範：同一式子裡多次更新 a（未定義行為）}
int a = 1;
printf("%d\n", a++ + ++a);   // 不要這樣寫
\end{code}

範例 \thechapter.13 展示複合指定運算子的行為。 \texttt{a += 5} 等於
\texttt{a = a + 5}；\texttt{+=} 的優先順序最低，所以 \texttt{n += 10}
會先把 \texttt{n} 加 $10$ 的結果再指定回 \texttt{n}。

\begin{example}{c}{(compound-ops.c) 複合指定運算子的優先序}
#include <stdio.h>

int main(void) {
    int a = 10, b = 20;
    printf("a = %d, b = %d\n", a, b);

    a += 5;     // 等於 a = a + 5
    printf("a += 5: a = %d\n", a);

    b -= 3;     // 等於 b = b - 3
    printf("b -= 3: b = %d\n", b);

    a *= 2;     // 等於 a = a * 2
    printf("a *= 2: a = %d\n", a);

    b /= 4;     // 等於 b = b / 4
    printf("b /= 4: b = %d\n", b);

    return 0;
}
\end{example}

\iosample{}{
a = 10, b = 20 \\
a += 5: a = 15 \\
b -= 3: b = 17 \\
a *= 2: a = 30 \\
b /= 4: b = 4 \\
}

\stylepoint{讓變數加一、減一或累加時，優先寫成單獨敘述。 複合指定能少
打一次變數名，但不要為了短而把三個運算擠在同一行。 二元運算子前後各
留一個空白，\texttt{=} 左右也是。}

\begin{drills}[label=drills:arith-assign-update-swap]
\item 直接寫 \texttt{a = b; b = a;} 能不能交換 \texttt{a} 跟 \texttt{b}
的值？ 為什麼？
\item \texttt{c = b = a;} 是先把 \texttt{a} 存進哪個變數，再存進哪個
變數？
\item \texttt{n += 10} 等價於哪一句用 \texttt{=} 寫出來的敘述？
\end{drills}

\section{浮點轉型}
\label{sec:arith-float-cast}

當整數運算式裡出現浮點常數（例如 \texttt{100.0}），C 會把整數運算元
\emph{提升}（promote）成浮點再計算，整個結果變成 \texttt{double}。

\begin{code}{c}{用 100.0 讓整數除法變成浮點除法}
double dollars = cents / 100.0;
\end{code}

若寫成 \texttt{cents / 100}，兩邊都是 \texttt{int}，先做整數除法，
\texttt{199 / 100 = 1}，再轉成 \texttt{1.00}。 跟預期的 \texttt{1.99}
差很遠。 分母寫成 \texttt{100.0}，\texttt{cents} 會被提升成
\texttt{double}，除法在浮點域進行，才得到正確的小數。

範例 \thechapter.14 讀入以「分」為單位的整數，輸出「元」並保留兩位
小數。

\begin{example}{c}{(cents-dollars.c) 分換算成元}
#include <stdio.h>

int main(void) {
    int cents;
    scanf("%d", &cents);

    printf("%.2f\n", cents / 100.0);
}
\end{example}

\iosample{
199 \\
}{
1.99 \\
}

\keypoint{需要小數結果時，檢查算式裡是否至少有一個浮點運算元；在整數
常數後面加 \texttt{.0} 是最簡單、也最不容易忘記的作法。 兩個 \texttt{int}
做 \texttt{/}，結果仍是 \texttt{int}，小數部分直接捨去。}

\begin{drills}[label=drills:arith-float-cast]
\item \texttt{cents / 100} 跟 \texttt{cents / 100.0}，哪一個會做浮點
除法？
\item 若 \texttt{cents = 199}，\texttt{cents / 100} 的結果是多少？
\texttt{cents / 100.0} 呢？
\end{drills}

\section{綜合應用}
\label{sec:arith-synthesis}

許多題目把前面所有運算子串在同一個 \texttt{main} 裡：先讀入多個量，
依公式算出中間結果，再 \texttt{printf}。 寫這類程式時，把「讀輸入 /
算公式 / 印輸出」分成三個區塊，中間必要時用額外變數存中間值，可讀
性會比把所有式子塞進一個 \texttt{printf} 好得多。

整數除法與餘數還可以處理\emph{位數}：從一個整數抽出百位、十位、個
位，或把兩個一位數合併成兩位數。

\begin{code}{c}{取出百位、十位、個位}
int hundreds = n / 100;
int tens = (n / 10) % 10;
int ones = n % 10;
\end{code}

\begin{code}{c}{將兩個一位數 d1、d2 合併成兩位數}
int combined = 10 * d1 + d2;
\end{code}

範例 \thechapter.15 讀入一個三位數，依序印出百、十、個位。

\begin{example}{c}{(digit-parts.c) 輸出百位、十位、個位}
#include <stdio.h>

int main(void) {
    int n;
    scanf("%d", &n);

    printf("%d %d %d\n", n / 100, (n / 10) % 10, n % 10);
}
\end{example}

\iosample{
482 \\
}{
4 8 2 \\
}

範例 \thechapter.16 做反向操作：讀入兩個 \texttt{0} 到 \texttt{9} 的
數字，合併成一個兩位數 \texttt{10 * d1 + d2}。

\begin{example}{c}{(combine-digits.c) 合併兩個一位數}
#include <stdio.h>

int main(void) {
    int d1, d2;
    scanf("%d %d", &d1, &d2);

    printf("%d\n", 10 * d1 + d2);
}
\end{example}

\iosample{
4 8 \\
}{
48 \\
}

範例 \thechapter.17 把「幾小時幾分鐘」換成總分鐘數，公式是
\texttt{hours * 60 + minutes}。 跟範例 \thechapter.7 互為逆運算。

\begin{example}{c}{(minutes-total.c) 小時與分鐘換算成總分鐘}
#include <stdio.h>

int main(void) {
    int hours, minutes;
    scanf("%d %d", &hours, &minutes);

    printf("%d\n", hours * 60 + minutes);
}
\end{example}

\iosample{
2 15 \\
}{
135 \\
}

範例 \thechapter.18 一次算出長方體的表面積與體積：表面積
$2(ab + bc + ca)$、體積 $abc$。 這裡用兩個變數 \texttt{surface} 和
\texttt{volume} 分別存結果，比直接在 \texttt{printf} 裡重複寫長公式
更容易對照題目、也更容易除錯。

\begin{example}{c}{(cuboid.c) 長方體表面積與體積}
#include <stdio.h>

int main(void) {
    int a, b, c;
    scanf("%d %d %d", &a, &b, &c);

    int surface = 2 * (a * b + b * c + c * a);
    int volume = a * b * c;
    printf("%d %d\n", surface, volume);
}
\end{example}

\iosample{
2 3 4 \\
}{
52 24 \\
}

「只要最後幾位」與「買十送一」都是把總量拆成完整組數與剩餘。 範例
範例 \thechapter.19 用 \texttt{n \% 10000} 取出最右邊四位，再用上一章的
\texttt{\%04d} 補零。 輸入 \texttt{123456} 得到 \texttt{3456}；輸入
\texttt{7} 得到 \texttt{0007}。

\begin{example}{c}{(last-four.c) 取整數的末四位並補零}
#include <stdio.h>

int main(void) {
    int n;
    scanf("%d", &n);
    printf("%04d\n", n % 10000);
}
\end{example}

\iosample{
123456 \\
}{
3456 \\
}

範例 \thechapter.20 模擬「每買十件送一件」：每 $11$ 件裡有 $10$ 件要
付錢、$1$ 件免費。 讀入想拿到的件數 \texttt{n}，完整組數是
\texttt{n / 11}，剩下 \texttt{n \% 11} 件沒湊成一組、仍要原價。 應付
件數因此是 \texttt{(n / 11) * 10 + (n \% 11)}。

\begin{example}{c}{(pay-packs.c) 買十送一的應付件數}
#include <stdio.h>

int main(void) {
    int n;
    scanf("%d", &n);
    int packs = n / 11;
    int extra = n % 11;
    int pay = packs * 10 + extra;
    printf("%d\n", pay);
}
\end{example}

\iosample{
11 \\
}{
10 \\
}

題庫後段還有更複雜的應用，例如從年齡反推並重組十位與個位、或解多
未知數的線性方程組。 那些題目建立在本章所有片語之上：整數四則、截
斷除法、餘數、位數拆解、以及必要時的 \texttt{long long} 避免溢位。
把這一章的範例與片語練熟，下一章就可以專心處理判斷與分支，而不會
被算式細節絆住。

\keypoint{複雜公式先拆成多個有名字的中間變數，再分別輸出；末 $k$ 位
用對 $10^k$ 取餘；「每 $m$ 個一組」用 \texttt{/ m} 與 \texttt{\% m}。
先用兩、三組小數字心算核對，再交給程式跑大測資。}

\begin{drills}[label=drills:arith-synthesis]
\item 要從一個三位數 \texttt{n} 取出十位數字，該怎麼寫？
\item 「每買 11 件送 1 件」，讀入想拿到的件數 \texttt{n}，應付件數的
公式是什麼？
\item 為什麼把長公式拆成幾個有名字的中間變數（例如 \texttt{surface}、
\texttt{volume}），會比全部塞進同一個 \texttt{printf} 好？
\end{drills}

\addcontentsline{toc}{section}{練習解答}
\section*{練習解答}

各節練習的答案收在這裡。 短題只給能一句話核對的結果。

\paragraph{第~\ref{sec:arith-operators}~節}
\begin{enumerate}
\item 不一定。 \texttt{a + b * c} 先算 \texttt{b * c} 再加 \texttt{a}；
\texttt{(a + b) * c} 先算 \texttt{a + b} 再乘 \texttt{c}，兩式只有在
特定數字下才會剛好相等。
\item 乘法會先跟緊鄰的項結合，寫成 \texttt{2 * a * b + b * c + c * a}
只有第一項被乘 $2$，跟 $2(ab+bc+ca)$ 不同；要改成
\texttt{2 * (a * b + b * c + c * a)}。
\item 都不是靠優先順序（減法只有一種優先順序），是靠結合方向：加減
乘除由左向右結合，所以 \texttt{a - b - c} 是先算 \texttt{a - b}。
\end{enumerate}

\paragraph{第~\ref{sec:arith-intdiv}~節}
\begin{enumerate}
\item \texttt{-1}。 C99 的 \texttt{\%} 與 \texttt{/} 配對向 $0$ 取整，
餘數的符號跟著被除數（\texttt{-7}），不是恆為非負。
\item \texttt{125 / 60 = 2} 小時，\texttt{125 \% 60 = 5} 分。
\item 不能。 對非負整數可以，但負奇數的餘數是 \texttt{-1}（例如
\texttt{-3 \% 2 = -1}），不是 \texttt{1}。
\end{enumerate}

\paragraph{第~\ref{sec:arith-assign-update-swap}~節}
\begin{enumerate}
\item 不能。 第一句 \texttt{a = b;} 已經把 \texttt{a} 原本的值蓋掉，
第二句 \texttt{b = a;} 只是把同一個值抄回去，\texttt{b} 原本的值已經
遺失。
\item 先存進 \texttt{b}（\texttt{b = a}），再把 \texttt{b} 的新值存進
\texttt{c}（因為 \texttt{=} 由右向左結合）。
\item \texttt{n = n + 10;}。
\end{enumerate}

\paragraph{第~\ref{sec:arith-float-cast}~節}
\begin{enumerate}
\item \texttt{cents / 100.0}：分母是浮點常數，\texttt{cents} 會被提
升成 \texttt{double}，除法在浮點域進行。
\item \texttt{cents / 100} 是整數除法，結果是 \texttt{1}；
\texttt{cents / 100.0} 是浮點除法，結果是 \texttt{1.99}。
\end{enumerate}

\paragraph{第~\ref{sec:arith-synthesis}~節}
\begin{enumerate}
\item \texttt{(n / 10) \% 10}。
\item \texttt{(n / 11) * 10 + (n \% 11)}。
\item 每個中間變數都有名字，對照題目公式時比較容易確認每一步算對了
沒有；擠在同一個 \texttt{printf} 裡難以分段檢查，也更容易在括號和運
算子優先順序上出錯。
\end{enumerate}

\chapter{判斷}
\label{ch:conditionals}

到目前為止，我們寫的程式大多是「讀入資料、計算、印出結果」，不管輸入是什麼，
執行的步驟都一樣。 但實際問題裡，程式常常需要\textbf{依條件選擇不同的
做法}：分數夠就及格、不夠就不及格；年份符合閏年規則就印 \texttt{YES}，否
則印 \texttt{NO}；輸入的字元是 \texttt{R} 就停、是 \texttt{G} 就走⋯⋯
這種「先檢查條件，再決定走哪條路」的結構，統稱為\textbf{判斷}或
\textbf{分支}（branching）。

這一章依語法複雜度推進：先認識關係式與 $0$／$1$，再學單路 \texttt{if}、
雙路 \texttt{if-else}、多路 \texttt{else if} 與獨立 \texttt{if} 的差別，
然後才組合 \texttt{\&\&}、\texttt{||} 與短路求值，最後用 \texttt{?:} 與
\texttt{switch} 處理特定情境。 在硬體層面，CPU 遇到條件分支時，會依條件
的真偽改變\textbf{程式計數器}（program counter）指向的位址，也就是「下一
條要執行的指令」；現代處理器還會做分支預測，猜錯了就得把已經投機執行的工
作丟掉重來。

沿路上我們也會特別留意幾個 C 語言裡容易踩到的坑：把 \texttt{==} 寫成
\texttt{=} 的賦值錯誤，以及 \texttt{else if} 與獨立 \texttt{if} 該怎麼
選。

\section{關係式與真假值}
\label{sec:cond-relational}

要做判斷，第一步是把「我們想問的問題」寫成一個可以求值的\textbf{關係式}
（relational expression）。 C 的關係運算子跟數學裡的 $<$、$>$、$=$ 很像，
但結果是\textbf{整數}：條件成立時結果是 $1$，不成立時是 $0$。 這個
$0$／$1$ 的慣例會一直跟著我們；之後 \texttt{if} 括號裡的條件，其實就
是在問「這個運算式的值是不是非零」。

\begin{code}{c}{六種關係運算子}
a < b    // a 是否小於 b
a > b    // a 是否大於 b
a <= b   // a 是否小於或等於 b
a >= b   // a 是否大於或等於 b
a == b   // a 是否等於 b（兩個等號！）
a != b   // a 是否不等於 b
\end{code}

範例 \thechapter.2 讀入兩個整數，直接把 \texttt{a == b} 與 \texttt{a > b}
的結果印出來。 輸入 \texttt{8 8} 時得到 \texttt{1} 再 \texttt{0}：相等
為真、大於為假。 輸入 \texttt{9 4} 則是 \texttt{0} 再 \texttt{1}。 這
時還沒用到 \texttt{if}，但已經看見比較式的值就是 $0$ 或 $1$。

\begin{example}{c}{(print-cmp.c) 把比較結果當成整數印出}
#include <stdio.h>

int main(void) {
    int a, b;
    scanf("%d %d", &a, &b);
    printf("%d\n", a == b);
    printf("%d\n", a > b);
    return 0;
}
\end{example}

\iosample{
8 8 \\
}{
1 \\
0 \\
}

C99 起可以用 \texttt{<stdbool.h>} 裡的 \texttt{bool}、\texttt{true}、
\texttt{false} 來強調邏輯值；印出來通常仍是 $1$ 與 $0$。 比較的優先順序
低於加減乘除。 \texttt{i > j + n} 其實是 \texttt{i > (j + n)}；仍建議
把算術部分括起來，讀的時候才不必先在腦子裡排優先順序。

C 語言裡最常見的初學者錯誤之一，是把「比較是否相等」的 \texttt{==} 寫成
「指派」用的 \texttt{=}：

\begin{code}{c}{把 == 誤寫成 = 的後果}
if (i == 0) {      // 正確：測試 i 是否等於 0
    ...
}

if (i = 0) {       // 錯誤：先把 0 指派給 i，再測試結果是否非零
    ...            // 這個 if 永遠不會執行（因為 0 是假）
}
\end{code}

編譯器通常不會把 \texttt{if (i = 0)} 當成語法錯誤。 它是一個合法的指派
運算式。 所以這類 bug 只能靠仔細閱讀程式碼或靠編譯器的警告訊息來抓。
養成習慣：想問「等不等於」時，一律寫兩個等號 \texttt{==}。

\keypoint{關係式的值是整數 $0$ 或 $1$。 \texttt{if (expression)} 把「非
零」當成真、「零」當成假。}

\begin{drills}[label=drills:cond-relational]
\item \texttt{a == b} 這個關係式，結果會是哪兩個可能的整數值？
\item 已知 \texttt{a = 3}、\texttt{b = 5}，\texttt{a < b} 的值是多少？
\item \texttt{if (i = 0)} 這行為什麼不會被編譯器擋下來？它實際做的是
比較還是賦值？
\end{drills}

\section{只有 if：條件成立才做事}
\label{sec:cond-if-only}

\texttt{if} 敘述的結構很簡單：括號裡放一個運算式，非零就執行後面那一
段，否則整段跳過。 不是每次判斷都有 \texttt{else}；有時條件成立才
做事，否則這一段直接跳過。

\begin{code}{c}{if 敘述}
if (condition) {
    statement1;
}
\end{code}

\begin{center}
\begin{tikzpicture}[node distance=1.3cm and 2.2cm]
    \node[startend] (start) {Start};
    \node[decision, below=of start] (cond) {condition?};
    \node[process, below left=1.3cm and 2.2cm of cond] (stmt) {statement1};
    \node[startend, below=4.4cm of cond] (end) {End};

    \draw[flowarrow] (start.south) -- (cond.north);
    \draw[flowarrow] (cond.west) -| node[above,pos=0.25,font=\scriptsize] {true} (stmt.north);
    \draw[flowarrow] (cond.south) -- node[right,font=\scriptsize] {false} (end.north);
    \draw[flowarrow] (stmt.south) -- ++(0,-0.5) |- (end.west);
\end{tikzpicture}
\diagcaption{只有 \texttt{if}：\texttt{condition} 為真時執行
\texttt{statement1}，為假時跳過，直接到 \texttt{End}。}
\end{center}

條件為假時，控制流沿主幹往下，不進入 \texttt{statement1}；為真時
向左執行動作後再匯入 \texttt{End}。

範例 \thechapter.5 讀入一個整數，只有正數才印出來。 輸入 \texttt{7}
會看到一行 \texttt{7}；輸入 \texttt{-2} 則什麼都不印。 不是當機，
只是那一個 \texttt{printf} 沒被執行。 這是本章第一個含 \texttt{if} 的
完整程式。

\begin{example}{c}{(print-if-positive.c) 只印出正數}
#include <stdio.h>

int main(void) {
    int n;
    scanf("%d", &n);
    if (n > 0) {
        printf("%d\n", n);
    }
    return 0;
}
\end{example}

\iosample{
7 \\
}{
7 \\
}

\begin{drills}[label=drills:cond-if-only]
\item 範例 print-if-positive.c 輸入 \texttt{-2} 時，程式會印出什麼？
\item 只有 \texttt{if}、沒有 \texttt{else} 時，條件為假會發生什麼事？
\item \texttt{if} 的括號裡放的是什麼型別的值？非零與零分別對應什麼？
\end{drills}

\section{if-else：二選一分支}
\label{sec:cond-if-else}

如果「條件不成立時也要做別的事」，就加上 \texttt{else}，形成\textbf{二
選一}：

\begin{code}{c}{if-else 敘述}
if (condition) {
    statement1;
} else {
    statement2;
}
\end{code}

\begin{center}
\begin{tikzpicture}[node distance=1.3cm and 2.2cm]
    \node[startend] (start) {Start};
    \node[decision, below=of start] (cond) {condition?};
    \node[process, below left=1.3cm and 2.2cm of cond] (stmt1) {statement1};
    \node[process, below right=1.3cm and 2.2cm of cond] (stmt2) {statement2};
    \node[startend, below=3.2cm of cond] (end) {End};

    \draw[flowarrow] (start.south) -- (cond.north);
    \draw[flowarrow] (cond.west) -| node[above,pos=0.25,font=\scriptsize] {true} (stmt1.north);
    \draw[flowarrow] (cond.east) -| node[above,pos=0.25,font=\scriptsize] {false} (stmt2.north);
    \draw[flowarrow] (stmt1.south) |- (end.west);
    \draw[flowarrow] (stmt2.south) |- (end.east);
\end{tikzpicture}
\diagcaption{\texttt{if-else}：\texttt{condition} 為真走 \texttt{statement1}，
為假走 \texttt{statement2}，兩路最後匯入 \texttt{End}。}
\end{center}

\texttt{condition} 為真時向左，為假時向右；兩條路各自執行完再往下
匯流。

範例 \thechapter.7 用一個分數門檻做最經典的 \texttt{if-else}：$60$ 分以
上及格，否則不及格。

\begin{example}{c}{(pass-or-fail.c) 依分數判斷及格或不及格}
#include <stdio.h>

int main(void) {
    int score;
    scanf("%d", &score);

    if (score >= 60) {
        printf("PASS\n");
    } else {
        printf("FAIL\n");
    }
}
\end{example}

\iosample{
55 \\
}{
FAIL \\
}

只有兩種互斥的結果時，\texttt{if-else} 很合適。

\begin{drills}[label=drills:cond-if-else]
\item 範例 pass-or-fail.c 輸入 \texttt{60} 時，會印出 \texttt{PASS}
還是 \texttt{FAIL}？
\item \texttt{if-else} 適合互斥的兩種結果，還是允許兩段都執行？
\end{drills}

\section{懸空的 else}
\label{sec:cond-dangling-else}

\texttt{else} 會配給\textbf{最近、還沒配對}的那個 \texttt{if}。 若內
層 \texttt{if} 沒加大括號，外層的 \texttt{else} 可能配錯對象。 這就
是懸空的 \texttt{else}（dangling else）。 本書一律加大括號，這個坑
會少很多。

\begin{code}{c}{危險：內層 if 沒加大括號時，else 可能配錯}
if (outer) {
    if (inner)
        statement1;
else
    statement2;
}
\end{code}

正確作法是讓每一個 \texttt{if}／\texttt{else} 的本體都是複合敘述，
縮排也跟配對一致。

\keypoint{\texttt{if} 和 \texttt{else} 各自控制\textbf{一段}敘述；如果
本體超過一行，一定要用 \texttt{\{\}} 包成複合敘述。 只寫一行不加括號、
下一行又寫別的敘述，編譯器會把 \texttt{else} 配錯對象，產生語法錯誤或
懸空的 \texttt{else}。}

\begin{drills}[label=drills:cond-dangling-else]
\item \texttt{else} 會配對給哪一個 \texttt{if}：離它最近的，還是最外
層的？
\item 為什麼本書建議每個 \texttt{if}／\texttt{else} 本體都用大括號包
起來？
\end{drills}

\section{用取餘數做二元分類}
\label{sec:cond-modulo}

很多判斷題目其實是在問「某個整數屬不屬於某一類」。 偶數還是奇數、能不能
被某個數整除、餘數是多少⋯⋯ 這類問題幾乎都離不開\textbf{取餘數}
（modulo）運算子 \texttt{\%}：\texttt{x \% d} 的值是 $x$ 除以 $d$ 的餘
數。 學會 \texttt{if-else} 之後，這是最常見的第一批實戰應用。

\begin{code}{c}{用 x \% d == 0 測試整除}
if (x % d == 0) {
    // x 能被 d 整除
} else {
    // x 不能被 d 整除
}
\end{code}

範例 \thechapter.10 是最常見的特例：\texttt{d = 2}，用餘數是 $0$ 還是
$1$ 來區分偶數與奇數。

\begin{example}{c}{(even-or-odd.c) 判斷偶數或奇數}
#include <stdio.h>

int main(void) {
    int n;
    scanf("%d", &n);

    if (n % 2 == 0) {
        printf("EVEN\n");
    } else {
        printf("ODD\n");
    }
}
\end{example}

\iosample{
4 \\
}{
EVEN \\
}

把 $d$ 換成任意正整數，就得到更一般的整除測試。 範例 \thechapter.11 讀入
\texttt{X} 和 \texttt{D}，判斷 \texttt{X} 能不能被 \texttt{D} 整除。 邏輯
跟偶奇判斷完全一樣，只是把 \texttt{2} 換成變數 \texttt{D}。

\begin{example}{c}{(divisible-by.c) 判斷 X 能否被 D 整除}
#include <stdio.h>

int main(void) {
    int X, D;
    scanf("%d %d", &X, &D);

    if (X % D == 0) {
        printf("YES\n");
    } else {
        printf("NO\n");
    }
}
\end{example}

\iosample{
12 3 \\
}{
YES \\
}

\stylepoint{「$n$ 是偶數」寫成 \texttt{n \% 2 == 0} 比寫成
\texttt{n / 2 * 2 == n} 更直接，也更容易一眼看出意圖。 若測資可能
含負數，奇數判斷應寫 \texttt{n \% 2 != 0}，因為負奇數的餘數是
\texttt{-1}，不是 \texttt{1}。}

\begin{drills}[label=drills:cond-modulo]
\item 判斷偶數要寫 \texttt{n \% 2 == 0} 還是 \texttt{n \% 2 == 1}？
\item \texttt{-3 \% 2} 的值是多少？為什麼奇數判斷建議寫成
\texttt{n \% 2 != 0} 而不是 \texttt{n \% 2 == 1}？
\item \texttt{X \% D == 0} 這個條件在問什麼？
\end{drills}

\section{用判斷改寫數值}
\label{sec:cond-rewrite-values}

判斷也可以拿來改變數，而不只是選哪一句 \texttt{printf}。 範例
\thechapter.12 讀入兩個整數，印出較大的那一個：把 \texttt{a >= b} 當成
問題「$a$ 有沒有至少跟 $b$ 一樣大？」，是就印 \texttt{a}，否則印
\texttt{b}。 輸入 \texttt{3 7} 時，\texttt{a >= b} 的值是 $0$，改走
\texttt{else} 印 $7$。

\begin{example}{c}{(maximum-of-two.c) 印出兩個整數中較大者}
#include <stdio.h>

int main(void) {
    int a, b;
    scanf("%d %d", &a, &b);

    if (a >= b) {
        printf("%d\n", a);
    } else {
        printf("%d\n", b);
    }
}
\end{example}

\iosample{
3 7 \\
}{
7 \\
}

範例 \thechapter.13 求絕對值：若 \texttt{n} 為負，就把 \texttt{n} 改成
\texttt{-n}，最後一律印出非負的結果。

\begin{example}{c}{(abs-int.c) 用 if 求絕對值}
#include <stdio.h>

int main(void) {
    int n;
    scanf("%d", &n);
    if (n < 0) {
        n = -n;
    }
    printf("%d\n", n);
    return 0;
}
\end{example}

\iosample{
-9 \\
}{
9 \\
}

「三個數裡最大的那個」不必一次寫出六道兩兩比較。 先假設第一個就是
目前最大，再跟第二個比、跟第三個比，較大的就更新。 這是一串\textbf{
獨立}的 \texttt{if}，因為每一個都可能改寫 \texttt{max}。

\begin{example}{c}{(max-three.c) 用獨立 if 更新目前最大值}
#include <stdio.h>

int main(void) {
    int a, b, c;
    scanf("%d %d %d", &a, &b, &c);
    int max = a;
    if (b > max) {
        max = b;
    }
    if (c > max) {
        max = c;
    }
    printf("%d\n", max);
    return 0;
}
\end{example}

\iosample{
4 9 1 \\
}{
9 \\
}

範例 \thechapter.15 則是「必要時才交換」：若 \texttt{a > b}，用上一
章的暫存變數對調，最後由小到大印出。 輸入 \texttt{9 2} 得到
\texttt{2 9}。

\begin{example}{c}{(print-ordered.c) 兩數由小到大印出}
#include <stdio.h>

int main(void) {
    int a, b;
    scanf("%d %d", &a, &b);
    if (a > b) {
        int temp = a;
        a = b;
        b = temp;
    }
    printf("%d %d\n", a, b);
    return 0;
}
\end{example}

\iosample{
9 2 \\
}{
2 9 \\
}

\begin{drills}[label=drills:cond-rewrite-values]
\item 求絕對值時，\texttt{n < 0} 成立要對 \texttt{n} 做什麼動作？
\item 找三個數的最大值，為什麼要用三個獨立的 \texttt{if}，而不是
\texttt{else if} 階梯？
\item 範例 print-ordered.c 輸入 \texttt{9 2} 時，輸出的兩數順序是什
麼？
\end{drills}

\section{else if 階梯：多路選擇}
\label{sec:cond-else-if-ladder}

當有\textbf{好幾個互斥的區間}要對應不同結果。 例如分數換等第。 與其寫一
串互相嵌套的 \texttt{if-else}，不如排成\textbf{階梯}：從上往下依序測試，
第一個成立的條件決定結果，後面全部跳過。

\begin{code}{c}{else if 階梯}
if (condition1) {
    ...
} else if (condition2) {
    ...
} else if (condition3) {
    ...
} else {
    ...
}
\end{code}

\begin{center}
\begin{tikzpicture}[node distance=1.3cm and 2.0cm]
    \node[startend] (start) {Start};
    \node[decision, below=of start] (cond1) {cond1?};
    \node[decision, below=1.6cm of cond1] (cond2) {cond2?};
    \node[decision, below=1.6cm of cond2] (cond3) {cond3?};
    \node[process, below=1.6cm of cond3] (default) {default};
    \node[startend, below=1.3cm of default] (end) {End};

    \node[process, left=3.2cm of cond1] (stmt1) {stmt1};
    \node[process, left=3.2cm of cond2] (stmt2) {stmt2};
    \node[process, left=3.2cm of cond3] (stmt3) {stmt3};

    \draw[flowarrow] (start.south) -- (cond1.north);
    \draw[flowarrow] (cond1.south) -- node[right,font=\scriptsize] {F} (cond2.north);
    \draw[flowarrow] (cond2.south) -- node[right,font=\scriptsize] {F} (cond3.north);
    \draw[flowarrow] (cond3.south) -- node[right,font=\scriptsize] {F} (default.north);
    \draw[flowarrow] (default.south) -- (end.north);

    \draw[flowarrow] (cond1.west) -- node[above,font=\scriptsize] {T} (stmt1.east);
    \draw[flowarrow] (cond2.west) -- node[above,font=\scriptsize] {T} (stmt2.east);
    \draw[flowarrow] (cond3.west) -- node[above,font=\scriptsize] {T} (stmt3.east);

    % 匯流線畫在 stmt 左側，避免從 stmt1 往下穿過 stmt2、stmt3
    \coordinate (bus) at ([xshift=-0.8cm]stmt1.west);
    \coordinate (b1) at (bus |- stmt1.west);
    \coordinate (b2) at (bus |- stmt2.west);
    \coordinate (b3) at (bus |- stmt3.west);
    \coordinate (bend) at (bus |- end.west);

    \draw[flowarrow] (stmt1.west) -- (b1);
    \draw[flowarrow] (stmt2.west) -- (b2);
    \draw[flowarrow] (stmt3.west) -- (b3);
    \draw[black!70, thick] (b1) -- (bend);
    \draw[flowarrow] (bend) -- (end.west);
\end{tikzpicture}
\diagcaption{沿右側主幹依序測試 \texttt{cond1}、\texttt{cond2}、\texttt{cond3}；
第一個成立的條件執行對應 \texttt{stmt} 後整條結束，全部不成立才走
\texttt{default}。}
\end{center}

右側主幹上，條件不成立就往下測下一個；成立時向左執行對應
\texttt{stmt}，再匯入 \texttt{End}。

範例 \thechapter.17 把分數換成等第字母：$A \ge 90$、$B \ge 80$、依此類
推。 注意條件必須\textbf{從高到低}排。 如果先測 \texttt{score >= 60}，
那 $95$ 分也會被歸成 \texttt{D}。

\begin{example}{c}{(grade-level.c) 分數換等第}
#include <stdio.h>

int main(void) {
    int score;
    scanf("%d", &score);

    if (score >= 90) {
        printf("A\n");
    } else if (score >= 80) {
        printf("B\n");
    } else if (score >= 70) {
        printf("C\n");
    } else if (score >= 60) {
        printf("D\n");
    } else {
        printf("E\n");
    }
}
\end{example}

\iosample{
72 \\
}{
C \\
}

三種互斥情形（正、負、零）也可以用同一套階梯來寫。 範例 \thechapter.18
依正負零印出 \texttt{POSITIVE}、\texttt{NEGATIVE}、\texttt{ZERO}。

\begin{example}{c}{(sign-of-number.c) 判斷正數、負數或零}
#include <stdio.h>

int main(void) {
    int n;
    scanf("%d", &n);

    if (n > 0) {
        printf("POSITIVE\n");
    } else if (n < 0) {
        printf("NEGATIVE\n");
    } else {
        printf("ZERO\n");
    }
}
\end{example}

\iosample{
-3 \\
}{
NEGATIVE \\
}

範例 \thechapter.19 用二次方程的判別式判斷根的類型。 判別式
$\Delta = b^2 - 4ac$ 決定根的個數：$\Delta > 0$ 時有兩個實根，
$\Delta = 0$ 時有一個根，$\Delta < 0$ 時無實根。

\begin{example}{c}{(quadratic-roots.c) 判斷二次方程的根}
#include <stdio.h>

int main(void) {
    int a, b, c;
    scanf("%d %d %d", &a, &b, &c);

    int delta = b * b - 4 * a * c;

    if (delta > 0) {
        printf("TWO REAL ROOTS\n");
    } else if (delta == 0) {
        printf("ONE ROOT\n");
    } else {
        printf("NO REAL ROOTS\n");
    }

    return 0;
}
\end{example}

\iosample{
1 -3 2 \\
1 -2 1 \\
1 0 1 \\
}{
TWO REAL ROOTS \\
ONE ROOT \\
NO REAL ROOTS \\
}

如果有三種（或更多）\textbf{互斥}的情形，例如兩數相等、前者較大、前者
較小，也要用連鎖的 \texttt{if-else if-else}。 範例 \thechapter.20 讀入
\texttt{A} 和 \texttt{B}，依大小關係印出 \texttt{A=B}、\texttt{A<B}、或
\texttt{A>B}。 順序很重要：必須先測 \texttt{A == B}，再測 \texttt{A < B}，
最後的 \texttt{else} 才代表 \texttt{A > B}。

\begin{example}{c}{(size-relation.c) 三種大小關係}
#include <stdio.h>

int main(void) {
    int A, B;
    scanf("%d %d", &A, &B);

    if (A == B) {
        printf("A=B\n");
    } else if (A < B) {
        printf("A<B\n");
    } else {
        printf("A>B\n");
    }
}
\end{example}

\iosample{
5 5 \\
}{
A=B \\
}

\keypoint{\texttt{else if} 階梯裡，一旦某個 \texttt{if} 或 \texttt{else if}
成立並執行完，\textbf{整條階梯就結束}。 不會再往下測試。 這正是「區間
互斥」題目該用階梯、而不是一串獨立 \texttt{if} 的原因。}

\begin{drills}[label=drills:cond-else-if-ladder]
\item 把分數換成等第的 \texttt{else if} 階梯，條件為什麼必須從高分排
到低分？
\item 一旦某個 \texttt{else if} 成立並執行完，階梯剩下的條件還會被
測試嗎？
\end{drills}

\section{獨立 if 與決策表：可以同時符合多條規則}
\label{sec:cond-independent-if}

\texttt{else if} 階梯適合「互斥選一」；但有些題目允許\textbf{多條規則同
時成立}，只是最後要依題意取「加總」或「最大值」。 這時候應該寫\textbf{一
串獨立的 \texttt{if}}，每條規則各自檢查、各自更新結果，而不是用
\texttt{else if} 把後面的規則擋掉。

範例 \thechapter.21 模擬刮刮卡：三條獎金規則可能同時符合，但只領\textbf{
最高}的那一項。 程式用三個獨立的 \texttt{if} 依序檢查，每次用「新獎金是
否比目前累計更高」來更新 \texttt{prize}。

\begin{example}{c}{(scratch-card-max.c) 刮刮卡取最高獎金}
#include <stdio.h>

int main(void) {
    int d1, d2;
    scanf("%d %d", &d1, &d2);

    int prize = 0;

    // 規則 1：d1 為奇數，100 元
    if (d1 % 2 == 1) {
        prize = 100;
    }

    // 規則 2：d2 是 2、5 或 8，200 元
    if (d2 == 2 || d2 == 5 || d2 == 8) {
        if (prize < 200) {
            prize = 200;
        }
    }

    // 規則 3：兩數相同，50 元
    if (d1 == d2) {
        if (prize < 50) {
            prize = 50;
        }
    }

    printf("%d\n", prize);
}
\end{example}

\iosample{
1 3 \\
}{
100 \\
}

\iosample{
2 5 \\
}{
200 \\
}

如果把三條規則改成 \texttt{else if} 階梯，第一條成立後後面永遠不會再
檢查。 對「取最高獎金」這種\textbf{非互斥}的題目就會答錯。 寫之前先問：
「這些條件是互斥的，還是可能同時成立？」

\keypoint{\texttt{else if} 表示「前面都不成立才看這條」；獨立的 \texttt{if}
表示「不管前面如何，這條都要檢查」。 決策表題目如果允許多條規則同時命中，
幾乎一定該用獨立 \texttt{if}。}

\stylepoint{取最大、取最小這類「目前最佳解」用獨立 \texttt{if} 逐一
更新；只有互斥選一時才用 \texttt{else if}。}

\begin{drills}[label=drills:cond-independent-if]
\item 什麼情況該用一串獨立的 \texttt{if}，而不是 \texttt{else if}
階梯？
\item 若把刮刮卡的三條獎金規則改寫成 \texttt{else if} 階梯，會有什麼
後果？
\end{drills}

\section{邏輯運算子：組合多個條件}
\label{sec:cond-logical-ops}

單一條件有時不夠用。 「$X$ 在 $L$ 到 $H$ 之間（含端點）」需要\textbf{
兩個}條件同時成立；「$i$ 在範圍外」則是\textbf{任一}條件成立即可。 C 提
供三個邏輯運算子來組合條件：

\begin{code}{c}{邏輯運算子 !、\&\&、||}
!expr          // 邏輯非：expr 為 0 時結果為 1，否則為 0
expr1 && expr2 // 邏輯且：兩邊都非 0 時結果為 1
expr1 || expr2 // 邏輯或：任一邊非 0 時結果為 1
\end{code}

比較運算子的結果是 $0$ 或 $1$，可以直接當成邏輯運算子的運算元。 C 把
\textbf{任何非零}都當成真，所以 \texttt{3 \&\& 0} 是假（$0$），
\texttt{3 || 0} 是真（$1$），\texttt{!3} 是假。 能寫清楚的關係式就
寫清楚，不要靠「非零就算真」讓人猜。

範例 \thechapter.23 讀取下界 \texttt{L}、上界 \texttt{H} 和待測值
\texttt{X}，用 \texttt{L <= X \&\& X <= H} 一次表達「$L \le X \le H$」。

\begin{example}{c}{(in-range.c) 判斷 X 是否在 L 到 H 之間（含端點）}
#include <stdio.h>

int main(void) {
    int L, H, X;
    scanf("%d %d %d", &L, &H, &X);

    if (L <= X && X <= H) {
        printf("YES\n");
    } else {
        printf("NO\n");
    }
}
\end{example}

\iosample{
1 10 5 \\
}{
YES \\
}

\begin{drills}[label=drills:cond-logical-ops]
\item \texttt{!}、\texttt{\&\&}、\texttt{||} 分別是什麼意思？
\item \texttt{L <= X \&\& X <= H} 這個複合條件在表達什麼？
\item \texttt{3 \&\& 0} 的值是多少？
\end{drills}

\section{短路求值}
\label{sec:cond-short-circuit}

\texttt{\&\&} 和 \texttt{||} 有一個重要性質：只要左邊已經足夠決定整
個運算式的真假，就不會去算右邊。 這叫短路求值（short-circuit
evaluation）：\texttt{\&\&} 左邊為假（$0$）時，右邊完全不會被求值；
\texttt{||} 左邊為真（非 $0$）時，右邊也不會被求值。

\begin{code}{c}{短路求值的規則}
expr1 && expr2   // 左邊為 0 時，expr2 不會被算
expr1 || expr2   // 左邊非 0 時，expr2 不會被算
\end{code}

這不只是效能上的小聰明，還能拿來擋掉本來會出錯的運算。 整數除以零是
未定義行為，把除數是否為零放在 \texttt{\&\&} 左邊，除法就只在確定安全
時才會被算。

\begin{example}{c}{(guard-divide.c) 用 \&\& 短路擋掉除以零}
#include <stdio.h>

int main(void) {
    int i;
    scanf("%d", &i);

    if (i != 0 && 100 / i > 5) {
        printf("YES\n");
    } else {
        printf("NO\n");
    }
    return 0;
}
\end{example}

\iosample{
0 \\
}{
NO \\
}

\iosample{
10 \\
}{
YES \\
}

輸入 \texttt{0} 時，\texttt{i != 0} 已經是假，\texttt{\&\&} 不會再去
算 \texttt{100 / i}，程式直接印出 \texttt{NO}，不會真的去執行那個除
法。 條件的順序不是隨意的：如果寫反成
\texttt{100 / i > 5 \&\& i != 0}，輸入 \texttt{0} 就會先算
\texttt{100 / i}，真的去做這個未定義行為的運算。 除以零在不同平台的
下場不一樣：有些會讓程式收到訊號直接中止，有些會安靜地給出一個沒有
意義的結果，不保證每次、每台機器都相同。 把會出錯的那段放在保護它的
條件後面，才能保證它永遠不會真的被執行。

\texttt{||} 的短路則用來跳過用不到的計算。 下面的
\texttt{is\_multiple\_of\_seven} 會在被呼叫時印一行訊息，用來證明它
到底有沒有被呼叫。

\begin{example}{c}{(skip-expensive-check.c) 用 || 短路跳過昂貴檢查}
#include <stdio.h>

int is_multiple_of_seven(int x) {
    printf("(checking %d)\n", x);
    return x % 7 == 0;
}

int main(void) {
    int x;
    scanf("%d", &x);

    if (x < 0 || is_multiple_of_seven(x)) {
        printf("YES\n");
    } else {
        printf("NO\n");
    }
    return 0;
}
\end{example}

\iosample{
-3 \\
}{
YES \\
}

\iosample{
14 \\
}{
(checking 14) \\
YES \\
}

輸入 \texttt{-3} 時，\texttt{x < 0} 已經是真，\texttt{||} 不會再呼叫
\texttt{is\_multiple\_of\_seven}，所以沒有印出 \texttt{(checking ...)}
這一行。 輸入 \texttt{14} 時左邊是假，右邊才真正被求值，
\texttt{(checking 14)} 這行證明函式確實被呼叫了一次。 沒有這行訊息，
讀者只能相信程式碼裡的註解；有了它，短路有沒有發生就是印出來的事實。

短路也表示放在右邊的副作用不保證每次都發生。 把遞增、遞減或函式呼叫
寫進 \texttt{\&\&}、\texttt{||} 右邊，副作用能不能發生會隨左邊的真假
改變，讀程式的人很容易漏看。 後面短路碰上指定那一節會示範更極端的情
形：右邊藏著一次指定，短路一發生，那個變數就悄悄維持舊值。

\keypoint{\texttt{\&\&} 左邊為假時不算右邊；\texttt{||} 左邊為真時
不算右邊。 拿短路擋掉會出錯的運算（例如除以零）很安全；拿它決定副作
用（遞增、指定、呼叫函式）發不發生，讀的人很容易漏看，應該避免。}

\begin{drills}[label=drills:cond-short-circuit]
\item \texttt{i != 0 \&\& 100 / i > 5}，輸入 \texttt{i = 0} 時，
\texttt{100 / i} 會不會被求值？為什麼這樣寫是安全的？
\item 若把上一題的條件寫反成 \texttt{100 / i > 5 \&\& i != 0}，輸入
\texttt{i = 0} 會發生什麼事？
\item \texttt{||} 左邊為真時，右邊還會不會被求值？
\end{drills}

\section{連鎖比較：別把數學式直接搬進 C}
\label{sec:cond-chained-compare}

數學上我們習慣寫 $L \le X \le H$，但 C 語言沒有連鎖比較這種語法。
\texttt{L <= X <= H} 其實會被解讀成 \texttt{(L <= X) <= H}：先算出
\texttt{L <= X} 的結果（$0$ 或 $1$），再拿這個 $0$／$1$ 去跟
\texttt{H} 比，幾乎永遠不是我們想要三者之間的大小關係。 正確寫法是拆
成兩個比較，中間用 \texttt{\&\&} 接起來；上一節的短路在這裡也用得上
，只要前段比較已經為假，\texttt{\&\&} 就不會再看後段。

「$A$、$B$、$C$ 嚴格遞增」在數學裡常寫成 $A < B < C$，但 C 不支援
這種寫法。 下面這段看起來像數學，其實完全不是：

\begin{code}{c}{錯誤：以為 a < b < c 代表三者嚴格遞增}
if (a < b < c) {   // 實際上是 (a < b) < c
    ...
}
\end{code}

\texttt{a < b} 的結果只有 $0$ 或 $1$。 如果 \texttt{a < b} 成立，左半部
是 $1$，接著 \texttt{1 < c} 只在 \texttt{c > 1} 時才成立。 這跟「三者
嚴格遞增」是兩回事。 正確寫法是把兩段比較拆開：

\begin{code}{c}{正確：用 \&\& 串起兩個嚴格不等式}
if (a < b && b < c) {
    ...
}
\end{code}

範例 \thechapter.29 讀入三個整數，依這個正確寫法判斷是否嚴格遞增。

\begin{example}{c}{(strictly-increasing.c) 判斷三數是否嚴格遞增}
#include <stdio.h>

int main(void) {
    int a, b, c;
    scanf("%d %d %d", &a, &b, &c);

    if (a < b && b < c) {
        printf("YES\n");
    } else {
        printf("NO\n");
    }
}
\end{example}

\iosample{
1 2 3 \\
}{
YES \\
}

\iosample{
2 1 3 \\
}{
NO \\
}

第二組測資 \texttt{2 1 3}：\texttt{a < b} 是 \texttt{2 < 1}，結果
$0$，整個 \texttt{\&\&} 因短路不再看 \texttt{b < c}，條件為假，印
\texttt{NO}。

\stylepoint{看到題目裡的「$L \le X \le H$」「$A < B < C$」這類數學連鎖比
較，在 C 裡一律改寫成兩個（或更多）關係式，用 \texttt{\&\&} 串起來。}

\begin{drills}[label=drills:cond-chained-compare]
\item \texttt{a < b < c} 在 C 裡實際上會被解讀成什麼？
\item 正確表達「\texttt{a}、\texttt{b}、\texttt{c} 嚴格遞增」該怎麼
寫？
\end{drills}

\section{複合條件：閏年規則}
\label{sec:cond-leap-year}

有些條件沒辦法用單一關係式表達，需要把多個整除測試用 \texttt{\&\&} 和
\texttt{||} 組合起來。 格里曆的閏年規則是：能被 $4$ 整除但不能被 $100$
整除，\textbf{或者}能被 $400$ 整除。

範例 \thechapter.30 直接把這個規則寫進 \texttt{if} 的條件裡。 $2000$ 年
能被 $400$ 整除，是閏年；$1900$ 年能被 $100$ 整除但不能被 $400$ 整除，
不是閏年。

\begin{example}{c}{(leap-year.c) 判斷是否為閏年}
#include <stdio.h>

int main(void) {
    int Y;
    scanf("%d", &Y);

    if ((Y % 4 == 0 && Y % 100 != 0) || Y % 400 == 0) {
        printf("YES\n");
    } else {
        printf("NO\n");
    }
}
\end{example}

\iosample{
2000 \\
}{
YES \\
}

\iosample{
1900 \\
}{
NO \\
}

用 $2000$ 走一遍：\texttt{Y \% 4 == 0} 成立，但 \texttt{Y \% 100 != 0}
不成立，前半段 \texttt{\&\&} 整體為假；接著 \texttt{Y \% 400 == 0} 成
立，\texttt{||} 右邊為真，整個條件為真，印 \texttt{YES}。

外層括號 \texttt{(Y \% 4 == 0 \&\& Y \% 100 != 0)} 不是多餘的裝飾。 
\texttt{\&\&} 比 \texttt{||} 優先序高，但加上括號可以讓「能被 $4$ 整除且
不是世紀年」這一整段意義一眼可讀，也跟題目文字描述的結構對齊。

\begin{drills}[label=drills:cond-leap-year]
\item 閏年規則的完整條件式怎麼寫（用 \texttt{\%4}、\texttt{\%100}、
\texttt{\%400}）？
\item 西元 \texttt{1900} 年是不是閏年？為什麼？
\end{drills}

\section{更多複合條件}
\label{sec:cond-more-compound}

「三邊能否圍成三角形」要兩邊都顧到：每邊為正，且任兩邊之和大於第
三邊。 少寫一條，就會把 \texttt{8 4 3} 這種退化情形誤判成合法。

\begin{example}{c}{(triangle-ok.c) 判斷三邊能否圍成三角形}
#include <stdio.h>

int main(void) {
    int a, b, c;
    scanf("%d %d %d", &a, &b, &c);
    if (a > 0 && b > 0 && c > 0 &&
        a + b > c && b + c > a && a + c > b) {
        printf("YES\n");
    } else {
        printf("NO\n");
    }
}
\end{example}

\iosample{
3 4 5 \\
}{
YES \\
}

\iosample{
8 4 3 \\
}{
NO \\
}

直角三角形則是畢氏定理：把「斜邊是 $c$」寫成
\texttt{a * a + b * b == c * c}。 範例 \thechapter.32 假設輸入已經依
「兩股、一斜」的順序給定。

\begin{example}{c}{(right-triangle.c) 判斷是否為以 c 為斜邊的直角三角形}
#include <stdio.h>

int main(void) {
    int a, b, c;
    scanf("%d %d %d", &a, &b, &c);
    if (a * a + b * b == c * c) {
        printf("YES\n");
    } else {
        printf("NO\n");
    }
}
\end{example}

\iosample{
3 4 5 \\
}{
YES \\
}

三個數是否都為正，是單純的 \texttt{\&\&}：任一個 $\le 0$ 就整段為
假。 「恰好一個條件成立」則比較像互斥或：兩邊的真假不能相同。 範
例 \thechapter.34 把「非零當成真」，用 \texttt{(p != 0) != (q != 0)}
表達「恰好一邊為真」。

\begin{example}{c}{(all-positive.c) 三數是否皆為正}
#include <stdio.h>

int main(void) {
    int a, b, c;
    scanf("%d %d %d", &a, &b, &c);
    if (a > 0 && b > 0 && c > 0) {
        printf("YES\n");
    } else {
        printf("NO\n");
    }
}
\end{example}

\iosample{
1 2 3 \\
}{
YES \\
}

\begin{example}{c}{(exactly-one.c) 恰好一個為真}
#include <stdio.h>

int main(void) {
    int p, q;
    scanf("%d %d", &p, &q);
    if ((p != 0) != (q != 0)) {
        printf("ONE\n");
    } else {
        printf("NOT_ONE\n");
    }
}
\end{example}

\iosample{
1 0 \\
}{
ONE \\
}

\texttt{\&\&} 與 \texttt{!} 也可以同時出現。 範例 \thechapter.35 要
求能被 $3$ 與 $5$ 整除、但不能被 $7$ 整除：\texttt{15} 符合，
\texttt{105}（$3\times5\times7$）不符合。

\begin{example}{c}{(mult-35-not7.c) 3 與 5 的倍數，但不是 7 的倍數}
#include <stdio.h>

int main(void) {
    int n;
    scanf("%d", &n);
    if (n % 3 == 0 && n % 5 == 0 && n % 7 != 0) {
        printf("YES\n");
    } else {
        printf("NO\n");
    }
}
\end{example}

\iosample{
15 \\
}{
YES \\
}

\begin{drills}[label=drills:cond-more-compound]
\item 判斷三邊能否圍成三角形，除了三邊都要為正，還要檢查什麼？
\item \texttt{(p != 0) != (q != 0)} 這個寫法在表達「恰好一個為真」，
它的邏輯是什麼？
\end{drills}

時間區間常常先換成「從午夜算起的總分鐘」，再當成普通整數比較。 範
例 \thechapter.36 判斷是否落在下午三點到五點半（含兩端）。

\begin{example}{c}{(tea-window.c) 把時與分換成總分鐘再比區間}
#include <stdio.h>

int main(void) {
    int h, m;
    scanf("%d %d", &h, &m);
    int t = h * 60 + m;
    if (t >= 15 * 60 && t <= 17 * 60 + 30) {
        printf("YES\n");
    } else {
        printf("NO\n");
    }
}
\end{example}

\iosample{
16 0 \\
}{
YES \\
}

三數取中位數可以用巢狀的「是否夾在另外兩個之間」來寫：先問 \texttt{a}
是不是中間那個，再問 \texttt{b}，否則就是 \texttt{c}。 相等時「夾在中
間」仍然成立，所以 \texttt{2 2 5} 的中位數是 $2$。 範例 \thechapter.37
示範這種複合條件的寫法。

\begin{example}{c}{(median-three.c) 用巢狀區間判斷三數的中位數}
#include <stdio.h>

int main(void) {
    int a, b, c;
    scanf("%d %d %d", &a, &b, &c);
    int mid;
    if ((a >= b && a <= c) || (a >= c && a <= b)) {
        mid = a;
    } else if ((b >= a && b <= c) || (b >= c && b <= a)) {
        mid = b;
    } else {
        mid = c;
    }
    printf("%d\n", mid);
    return 0;
}
\end{example}

\iosample{
9 1 4 \\
}{
4 \\
}

\section{短路碰上指定}
\label{sec:cond-short-circuit-assign}

前面一節說過 \texttt{||} 左邊為真就不再算右邊。 這裡看進階陷阱：若右邊
藏著指定，那一次「沒算」會讓變數維持舊值。 範例 \thechapter.38 把指定
寫進條件裡（這不是推薦寫法，只為了把陷阱攤開）：
\texttt{(k = i) > 0 || (ell = j) > 0}。 輸入 \texttt{7 2} 時左邊已經為
真，右邊的指定被跳過，\texttt{ell} 仍是 $0$；輸入 \texttt{0 2} 時左邊
為假，右邊才執行，\texttt{ell} 變成 $2$。

\begin{example}{c}{(short-side.c) 短路讓右邊的指定不一定發生}
#include <stdio.h>

int main(void) {
    int i, j;
    int k = 0;
    int ell = 0;
    scanf("%d %d", &i, &j);
    if ((k = i) > 0 || (ell = j) > 0) {
        printf("TRUE\n");
    } else {
        printf("FALSE\n");
    }
    printf("%d %d\n", k, ell);
    return 0;
}
\end{example}

\iosample{
7 2 \\
}{
TRUE \\
7 0 \\
}

\keypoint{不要把指定寫進 \texttt{\&\&}／\texttt{||} 的運算元裡。 短
路是為了少做不必做的測試，不是為了選擇性賦值。}

\begin{drills}[label=drills:cond-short-circuit-assign]
\item \texttt{(k = i) > 0 || (ell = j) > 0}，輸入 \texttt{7 2} 時，
\texttt{ell} 最後是多少？為什麼？
\item 為什麼不建議把賦值寫進 \texttt{\&\&}／\texttt{||} 的運算元裡？
\end{drills}

\section{條件運算子：寫在運算式裡的二選一}
\label{sec:cond-ternary}

有時你要的不是選擇執行哪一段敘述，只是依條件選兩個值裡的一個。
條件運算子 \texttt{?:} 的寫法是 \texttt{condition ? a : b}：條件非
零就取 \texttt{a}，否則取 \texttt{b}。 它是運算式，可以直接出現在
賦值右邊，不需要額外的變數先存起來再判斷。

\begin{code}{c}{條件運算子}
mag = (n >= 0) ? n : -n;
\end{code}

\texttt{?:} 的優先順序很低，條件兩邊建議加括號，讀的時候才不必先在
腦子裡排優先序。

\subsection{巢狀 ?:：一行做多選一}

\texttt{?:} 的第二個分支可以再放一個 \texttt{?:}，而且是\textbf{右
結合}：\texttt{a ? b : c ? d : e} 會被讀成
\texttt{a ? b : (c ? d : e)}。 這讓好幾個互斥的區間可以縮成一行，效果
跟 \texttt{else if} 階梯一樣，只是階梯選的是要執行的敘述，這裡選的是
要放進變數的值。

\begin{code}{c}{巢狀條件運算子}
grade = (score >= 90) ? 'A' :
        (score >= 80) ? 'B' :
        (score >= 70) ? 'C' : 'D';
\end{code}

範例 \thechapter.41 把前面等第字母那個 \texttt{else if} 階梯改寫成一
串巢狀 \texttt{?:}，答案完全一樣：輸入 \texttt{72} 一樣得到
\texttt{C}。

\begin{example}{c}{(grade-ternary.c) 用巢狀 ?: 選出等第字母}
#include <stdio.h>

int main(void) {
    int score;
    scanf("%d", &score);

    char grade = (score >= 90) ? 'A' :
                 (score >= 80) ? 'B' :
                 (score >= 70) ? 'C' :
                 (score >= 60) ? 'D' : 'E';

    printf("%c\n", grade);
    return 0;
}
\end{example}

\iosample{
72 \\
}{
C \\
}

每個分支都只是一個字元，沒有其他動作，這正是巢狀 \texttt{?:} 適用的
場合：每條規則都只是在選一個值。 分支一旦要做多件事，或條件超過三、
四層，讀起來會比對應的 \texttt{else if} 階梯更累；那時就該換回
\texttt{if-else}，不要硬把敘述塞進問號冒號兩側。

\subsection{陷阱：兩個分支的型別可能不一樣}

\texttt{?:} 的兩個分支不用同型別，但運算式最終只有\textbf{一個}型
別。 C 會先套用一般算術轉換，把兩個分支的型別統一，再決定整個
\texttt{?:} 表達式的型別。 就算實際選中的是 \texttt{int} 那一支，只
要另一支是 \texttt{double}，結果型別就是 \texttt{double}。

\begin{example}{c}{(ternary-type.c) 沒被選中的分支也決定了型別}
#include <stdio.h>

int main(void) {
    int n = 5;
    printf("%zu\n", sizeof(n > 0 ? n : 0));
    printf("%zu\n", sizeof(n > 0 ? n : 0.0));

    double result = (n > 0) ? n : 0.0;
    printf("%g\n", result);
    return 0;
}
\end{example}

\iosample{
}{
4 \\
8 \\
5 \\
}

\texttt{n > 0} 為真，兩次都選中 \texttt{n} 這一支。 第一行兩個分支都
是 \texttt{int}，\texttt{sizeof} 量到 $4$。 第二行另一支是
\texttt{0.0}，整個 \texttt{?:} 表達式因此變成 \texttt{double}，
\texttt{sizeof} 量到 $8$；\texttt{n} 本身沒有變，是選中它之後又被轉成
\texttt{double} 才交出去。 把這個結果直接印成 \texttt{\%d} 會讀到錯的
位元組，跟第 13 章格式指定符不對齊的陷阱同一個道理。

\keypoint{\texttt{?:} 的型別由兩個分支一起決定，不是由實際選中的那一
支單獨決定。 分支型別不同時，較窄的那一支會被轉換，就算它沒被選中。}

\stylepoint{\texttt{?:} 只適合選一個值；分支需要做動作、或巢狀超過兩
三層，就改回 \texttt{if-else}，不要把敘述硬塞進冒號兩側。}

\begin{drills}[label=drills:cond-ternary]
\item \texttt{mag = (n >= 0) ? n : -n;} 這行在做什麼？
\item \texttt{a ? b : c ? d : e} 的結合方向是左結合還是右結合？
\item \texttt{?:} 的兩個分支型別不同（一個 \texttt{int}、一個
\texttt{double}）時，整個 \texttt{?:} 表達式的型別由誰決定？
\end{drills}

\section{switch：一個值對應多種固定選項}
\label{sec:cond-switch}

當程式要依\textbf{一個整數或字元}的值，從\textbf{很多固定選項}中挑一個
來做，\texttt{if-else if} 階梯雖然可行，但 \texttt{switch} 通常更清楚，
編譯器也較容易生成有效率的跳轉表（jump table）。

\begin{code}{c}{switch 的基本結構}
switch (expression) {
    case 常數1:
        ...
        break;
    case 常數2:
        ...
        break;
    default:
        ...
        break;
}
\end{code}

\begin{center}
\begin{tikzpicture}[node distance=1.3cm and 2.0cm]
    \node[startend] (start) {Start};
    \node[process, below=of start] (expr) {Evaluate expr};
    \node[decision, below=1.3cm of expr] (case1) {expr == c1?};
    \node[decision, below=1.6cm of case1] (case2) {expr == c2?};
    \node[decision, below=1.6cm of case2] (case3) {expr == c3?};
    \node[process, below=1.6cm of case3] (default) {default};
    \node[startend, below=1.3cm of default] (end) {End};

    \node[process, left=3.2cm of case1] (stmt1) {stmt1};
    \node[process, left=3.2cm of case2] (stmt2) {stmt2};
    \node[process, left=3.2cm of case3] (stmt3) {stmt3};

    \draw[flowarrow] (start.south) -- (expr.north);
    \draw[flowarrow] (expr.south) -- (case1.north);
    \draw[flowarrow] (case1.south) -- node[right,font=\scriptsize] {N} (case2.north);
    \draw[flowarrow] (case2.south) -- node[right,font=\scriptsize] {N} (case3.north);
    \draw[flowarrow] (case3.south) -- node[right,font=\scriptsize] {N} (default.north);
    \draw[flowarrow] (default.south) -- (end.north);

    \draw[flowarrow] (case1.west) -- node[above,font=\scriptsize] {Y} (stmt1.east);
    \draw[flowarrow] (case2.west) -- node[above,font=\scriptsize] {Y} (stmt2.east);
    \draw[flowarrow] (case3.west) -- node[above,font=\scriptsize] {Y} (stmt3.east);

    % 匯流線畫在 stmt 左側，避免從 stmt1 往下穿過 stmt2、stmt3
    \coordinate (bus) at ([xshift=-0.8cm]stmt1.west);
    \coordinate (b1) at (bus |- stmt1.west);
    \coordinate (b2) at (bus |- stmt2.west);
    \coordinate (b3) at (bus |- stmt3.west);
    \coordinate (bend) at (bus |- end.west);

    \draw[flowarrow] (stmt1.west) -- (b1);
    \draw[flowarrow] (stmt2.west) -- (b2);
    \draw[flowarrow] (stmt3.west) -- (b3);
    \draw[black!70, thick] (b1) -- (bend);
    \draw[flowarrow] (bend) -- (end.west);
\end{tikzpicture}
\diagcaption{先求出 \texttt{expr}，再沿主幹依序問 \texttt{expr == c1?}、
\texttt{c2?}、\texttt{c3?}；相等時向左執行對應 \texttt{stmt}，都不符才
走 \texttt{default}。}
\end{center}

跟 \texttt{else if} 階梯同一套視覺語言：\texttt{switch} 是對同一個值
連續問相等，\texttt{else if} 則是對不同條件依序提問。

\texttt{case} 後面必須是\textbf{編譯期可知的常數}（字面常數或 \texttt{enum}
常數），不能是變數。 每個 \texttt{case} 結束通常要寫 \texttt{break}，否
則程式會\textbf{穿透}（fall through）到下一個 \texttt{case} 繼續執行。

範例 \thechapter.44 用 \texttt{enum} 定義一週七天，再依輸入的整數（$0$
到 $6$）印出對應的縮寫名稱。

\begin{example}{c}{(weekday-name.c) 用 enum 與 switch 印出星期名稱}
#include <stdio.h>

enum Weekday {
    SUN, MON, TUE, WED, THU, FRI, SAT
};

int main(void) {
    int n;
    scanf("%d", &n);

    enum Weekday day = n;
    switch (day) {
        case SUN:
            printf("SUN\n");
            break;
        case MON:
            printf("MON\n");
            break;
        case TUE:
            printf("TUE\n");
            break;
        case WED:
            printf("WED\n");
            break;
        case THU:
            printf("THU\n");
            break;
        case FRI:
            printf("FRI\n");
            break;
        case SAT:
            printf("SAT\n");
            break;
        default:
            printf("UNKNOWN\n");
            break;
    }
}
\end{example}

\iosample{
3 \\
}{
WED \\
}

多個 \texttt{case} 可以\textbf{刻意}不加 \texttt{break}，讓它們共用同一段
程式碼。 這叫做\textbf{分組 case}（grouped cases）。 範例 \thechapter.45
依月份輸出季節。

\begin{example}{c}{(season-from-month.c) 用分組 case 依月份輸出季節}
#include <stdio.h>

enum Season { SPRING, SUMMER, AUTUMN, WINTER };

int main(void) {
    int M;
    scanf("%d", &M);

    enum Season season;
    switch (M) {
        case 3:
        case 4:
        case 5:
            season = SPRING;
            break;
        case 6:
        case 7:
        case 8:
            season = SUMMER;
            break;
        case 9:
        case 10:
        case 11:
            season = AUTUMN;
            break;
        case 12:
        case 1:
        case 2:
            season = WINTER;
            break;
        default:
            season = WINTER;
            break;
    }

    switch (season) {
        case SPRING:
            printf("SPRING\n");
            break;
        case SUMMER:
            printf("SUMMER\n");
            break;
        case AUTUMN:
            printf("AUTUMN\n");
            break;
        case WINTER:
            printf("WINTER\n");
            break;
    }
}
\end{example}

\iosample{
4 \\
}{
SPRING \\
}

\stylepoint{只有當多個 \texttt{case} 要\textbf{刻意}共用同一段程式碼時，才
省略 \texttt{break} 讓它穿透；其他情況每個 \texttt{case} 結尾都應寫
\texttt{break}，避免執行流滑進下一個標籤。}

\texttt{switch} 的控制運算式也可以是字元：字元在 C 裡就是一種整數。
讀字元時，格式前面加一個空白 \texttt{" \%c"}，讓 \texttt{scanf} 先
跳過換行或空白，否則很容易讀到上一行留下來的 \texttt{\textbackslash n}。
範例 \thechapter.46 依 \texttt{R}／\texttt{Y}／\texttt{G} 印出停、等、
走。

\begin{example}{c}{(traffic.c) 用字元 switch 模擬號誌}
#include <stdio.h>

int main(void) {
    char light;
    scanf(" %c", &light);
    switch (light) {
        case 'R':
            printf("STOP\n");
            break;
        case 'Y':
            printf("WAIT\n");
            break;
        case 'G':
            printf("GO\n");
            break;
        default:
            printf("?\n");
            break;
    }
    return 0;
}
\end{example}

\iosample{
G \\
}{
GO \\
}

四則運算選單也很適合 \texttt{switch}。 範例 \thechapter.47 讀
\texttt{a}、運算子、\texttt{b}；除法時若除數為 $0$ 就印 \texttt{ERR}，
避免未定義行為。

\begin{example}{c}{(mini-calc.c) 依運算子字元做整數四則}
#include <stdio.h>

int main(void) {
    int a, b;
    char op;
    scanf("%d %c %d", &a, &op, &b);
    switch (op) {
        case '+':
            printf("%d\n", a + b);
            break;
        case '-':
            printf("%d\n", a - b);
            break;
        case '*':
            printf("%d\n", a * b);
            break;
        case '/':
            if (b == 0) {
                printf("ERR\n");
            } else {
                printf("%d\n", a / b);
            }
            break;
        default:
            printf("ERR\n");
            break;
    }
    return 0;
}
\end{example}

\iosample{
8 * 3 \\
}{
24 \\
}

\begin{drills}[label=drills:cond-switch]
\item \texttt{case} 後面必須是什麼樣的值？可以是變數嗎？
\item 忘記在 \texttt{case} 結尾寫 \texttt{break} 會發生什麼事？這種
現象叫做什麼？
\item 什麼情況下「刻意」不寫 \texttt{break} 是合理的？
\end{drills}

到這裡，關係式、邏輯組合、\texttt{if}／\texttt{else if}、獨立
\texttt{if}、以及 \texttt{switch} 都已經能處理題庫裡大多數「先分類、
再輸出」的題目。 下一章「迴圈」會把同一段判斷重複很多次。 條件的
寫法不變，只是它會被問過一遍又一遍。

\addcontentsline{toc}{section}{練習解答}
\section*{練習解答}

各節練習的答案收在這裡。 短題只給能一句話核對的結果。

\paragraph{第~\ref{sec:cond-relational}~節}
\begin{enumerate}
\item $0$ 或 $1$。
\item $1$（真）。
\item 它是一個合法的賦值運算式：先把 \texttt{0} 存進 \texttt{i}，再
把賦值後的結果（\texttt{0}，也就是假）當成條件值，所以這個 \texttt{if}
永遠不會執行。
\end{enumerate}

\paragraph{第~\ref{sec:cond-if-only}~節}
\begin{enumerate}
\item 什麼都不印：\texttt{n > 0} 不成立，\texttt{printf} 那段被跳
過，但程式不會當機。
\item 直接跳過整個本體，繼續往下執行 \texttt{if} 之後的敘述。
\item 整數（$0$ 或非零）；非零視為真，零視為假。
\end{enumerate}

\paragraph{第~\ref{sec:cond-if-else}~節}
\begin{enumerate}
\item \texttt{PASS}（\texttt{score >= 60} 成立）。
\item 互斥的兩種結果——同一時間只會執行其中一段。
\end{enumerate}

\paragraph{第~\ref{sec:cond-dangling-else}~節}
\begin{enumerate}
\item 離它最近、還沒配對的那個 \texttt{if}。
\item 因為只寫一行不加括號時，\texttt{else} 可能配到內層的 \texttt{if}
而不是外層，跟縮排暗示的結構不一致，容易產生邏輯錯誤。
\end{enumerate}

\paragraph{第~\ref{sec:cond-modulo}~節}
\begin{enumerate}
\item \texttt{n \% 2 == 0}。
\item $-1$；因為負奇數的餘數是 $-1$ 而不是 $1$，用 \texttt{== 1} 會
把負奇數誤判成偶數。
\item \texttt{X} 能不能被 \texttt{D} 整除。
\end{enumerate}

\paragraph{第~\ref{sec:cond-rewrite-values}~節}
\begin{enumerate}
\item 把 \texttt{n} 改成 \texttt{-n}。
\item 因為每個數都可能改寫目前最大值，用 \texttt{else if} 的話，一旦
第一條成立，後面的數就不會再被拿來比較，可能漏掉真正的最大值。
\item \texttt{2 9}（由小到大）。
\end{enumerate}

\paragraph{第~\ref{sec:cond-else-if-ladder}~節}
\begin{enumerate}
\item 因為階梯由上往下、第一個成立的條件就決定結果；如果先測低分的
門檻（例如 \texttt{score >= 60}），高分也會被這條擋下來，永遠測不
到後面更高的門檻。
\item 不會，階梯在第一個成立的條件執行完後就結束。
\end{enumerate}

\paragraph{第~\ref{sec:cond-independent-if}~節}
\begin{enumerate}
\item 多條規則允許同時成立、且結果要取加總或最大值時。
\item 第一條規則成立後，後面的規則就不會再被檢查，可能選不到真正該
給的最高獎金。
\end{enumerate}

\paragraph{第~\ref{sec:cond-logical-ops}~節}
\begin{enumerate}
\item \texttt{!} 是邏輯非；\texttt{\&\&} 是邏輯且（兩邊都非零才為
真）；\texttt{||} 是邏輯或（任一邊非零就為真）。
\item \texttt{X} 是否介於 \texttt{L} 和 \texttt{H} 之間（含兩端）。
\item $0$（假）。
\end{enumerate}

\paragraph{第~\ref{sec:cond-short-circuit}~節}
\begin{enumerate}
\item 不會；\texttt{\&\&} 左邊已經為假，右邊直接跳過，所以除以零不
會真的被執行。
\item \texttt{100 / i} 會先被算，輸入 \texttt{0} 時觸發除以零的未定
義行為。
\item 不會。
\end{enumerate}

\paragraph{第~\ref{sec:cond-chained-compare}~節}
\begin{enumerate}
\item \texttt{(a < b) < c}：先算出 \texttt{a < b} 的 $0$／$1$，再拿
這個值跟 \texttt{c} 比較，不是三者遞增關係。
\item \texttt{a < b \&\& b < c}。
\end{enumerate}

\paragraph{第~\ref{sec:cond-leap-year}~節}
\begin{enumerate}
\item \texttt{(Y \% 4 == 0 \&\& Y \% 100 != 0) || Y \% 400 == 0}。
\item 不是；\texttt{1900} 能被 \texttt{100} 整除，但不能被 \texttt{400}
整除。
\end{enumerate}

\paragraph{第~\ref{sec:cond-more-compound}~節}
\begin{enumerate}
\item 任兩邊之和要大於第三邊。
\item 兩者的「非零狀態」不同（一個為真、一個為假）時，結果才為真。
\end{enumerate}

\paragraph{第~\ref{sec:cond-short-circuit-assign}~節}
\begin{enumerate}
\item $0$；因為左邊 \texttt{(k = i) > 0} 已經為真，短路讓右邊的
\texttt{(ell = j) > 0} 完全不會被求值，\texttt{ell} 維持原本的初值。
\item 因為短路可能讓賦值不被執行，副作用發不發生要看左邊的真假，容
易讓讀程式的人漏看或誤判。
\end{enumerate}

\paragraph{第~\ref{sec:cond-ternary}~節}
\begin{enumerate}
\item 若 \texttt{n} 非負就把 \texttt{n} 存進 \texttt{mag}，否則存
\texttt{-n}——這是在算絕對值。
\item 右結合。
\item 由兩個分支型別一起決定：較窄的一支會被轉換成較寬的型別，就算
它沒被選中。
\end{enumerate}

\paragraph{第~\ref{sec:cond-switch}~節}
\begin{enumerate}
\item 編譯期可知的常數（字面常數或 \texttt{enum} 常數）；不能是變
數。
\item 執行流會「穿透」到下一個 \texttt{case} 繼續執行，這現象叫做
fall through。
\item 想讓多個 \texttt{case} 刻意共用同一段程式碼時（分組 case）。
\end{enumerate}

\chapter{迴圈}
\label{ch:loop}

到目前為止，我們寫的程式多半只執行一次就結束。 讀一個數字、算一個結果、印出
來。 但實際上很多問題需要\textbf{重複}做同一件事：印出 $N$ 行相同的文字、把
一串數字逐一讀進來、一直問使用者輸入直到格式正確、或是模擬每一天細菌怎麼繁殖。
如果每一輪都手動複製貼上同一段程式碼，程式不但冗長，改一個地方還得改很多
處，幾乎無法維護。

\textbf{迴圈}（loop）就是為了解決這個問題：把「要重複做的事」寫成一段，再告
訴電腦「重複幾次」或「重複到什麼條件為止」。 C 語言提供三種迴圈語法。 
\texttt{for}、\texttt{while}、\texttt{do-while}。 語法不同，但背後都是同一
套機制：每次迭代（iteration）更新迴圈變數或狀態，直到條件不再成立為止。 從
硬體的角度看，CPU 其實只是在反覆執行「比較條件 → 跳回或繼續往下」這兩步；
迴圈讓我們用高階語法描述這個過程，而不必自己寫跳轉指令。

這一章依語法複雜度推進：先認識三種迴圈與 \texttt{for}/\texttt{while} 的等價
關係，再學 \texttt{break} 與 \texttt{continue}，接著用巢狀迴圈輸出圖形，然
後練解題平台必備的輸入模板，最後才進入數論算法與需要追蹤多個狀態的模擬題。
題庫裡這個單元有八十多題，難度跨度很大；我們沿著整條難度曲線走，把核心語法
與經典解題模型分開講清楚。

\section{固定次數的重複：for 迴圈}
\label{sec:loop-for-basics}

\texttt{for} 迴圈最適合「次數事先知道」的情形：從 $0$ 數到 $n-1$、從 $1$ 數
到 $n$、或是倒數 $n$ 次。 它的結構把「初始化」「條件」「每次迭代後的更新」
集中寫在一行，一眼就能看出迴圈跑幾輪。

\begin{code}{c}{for 迴圈的基本形狀}
for (初始化; 條件; 更新) {
    重複執行的敘述;
}
\end{code}

範例 \thechapter.2 讀入正整數 $n$，把同一行文字印 $n$ 次。 這是最單純的
「已知次數、每次做同一件事」的模型。

\begin{example}{c}{(print-fixed-line.c) 把固定文字印 N 次}
#include <stdio.h>

int main(void) {
    int n;
    scanf("%d", &n);

    for (int i = 0; i < n; i++) {
        printf("Hello\n");
    }
    return 0;
}
\end{example}

\iosample{
3 \\
}{
Hello \\
Hello \\
Hello \\
}

\texttt{for (int i = 0; i < n; i++)} 裡，\texttt{i} 從 $0$ 開始，每輪加
$1$，當 \texttt{i == n} 時條件不成立、迴圈結束，所以總共執行 $n$ 次。
這跟陣列索引從 $0$ 開始的慣例一致。 之後用 \texttt{for} 走訪陣列時，幾
乎都會看到 \texttt{i < n} 這種寫法。


\keypoint{\texttt{for (int i = 0; ...)} 裡宣告的 \texttt{i}，作用域就是這
一句 \texttt{for}；迴圈結束後在 \texttt{main} 裡不能再寫 \texttt{printf("\%d", i)}。
若要在迴圈外使用計數器，把 \texttt{int i} 提到 \texttt{for} 前面。}

\keypoint{\texttt{for} 迴圈的三個部分（初始化、條件、更新）都在迴圈\textbf{開
始前}就決定好；條件在每一輪\textbf{一開始}先檢查，不成立就整個跳過，一次都
不執行。}

\begin{drills}[label=drills:loop-for-basics]
\item \texttt{for (int i = 0; i < n; i++)} 印 \texttt{Hello} $n$ 次，
若輸入 $n = 0$，迴圈本體會執行幾次？
\item \texttt{for} 裡宣告的 \texttt{int i}，離開 \texttt{for} 這一句之後
還能不能在 \texttt{main} 裡繼續用？
\item \texttt{i < n} 與 \texttt{i <= n} 兩種條件，各自讓迴圈跑幾次？
\end{drills}

\section{改變計數方向：倒數與反向列印}
\label{sec:loop-countdown}

次數固定，但「每一輪要做的事」跟計數器的值有關時，只要改初始化、條件、更新
的方向就好。 範例 \thechapter.3 從 $n$ 倒數印到 $1$，每一行印一個數字。

\begin{example}{c}{(print-backward.c) 從 n 倒數印到 1}
#include <stdio.h>

int main(void) {
    int n;
    scanf("%d", &n);

    for (int i = n; i >= 1; i--) {
        printf("%d\n", i);
    }
    return 0;
}
\end{example}

\iosample{
5 \\
}{
5 \\
4 \\
3 \\
2 \\
1 \\
}

範例 \thechapter.4 則把 $n$ 到 $0$ 全部印在同一行，數字之間以空白分隔。 
這裡多了一個常見的輸出細節：最後一個數字後面不要多印空白。

\begin{example}{c}{(countdown.c) 倒數 n 到 0，空白分隔}
#include <stdio.h>

int main(void) {
    int n;
    scanf("%d", &n);

    for (int i = n; i >= 0; i--) {
        printf("%d", i);
        if (i > 0) {
            printf(" ");
        }
    }
    printf("\n");
    return 0;
}
\end{example}

\iosample{
3 \\
}{
3 2 1 0 \\
}

以 $n = 3$ 走一遍：\texttt{i} 依序是 $3, 2, 1, 0$；印完 $3$、$2$、$1$ 時
\texttt{i > 0} 成立，後面補空白；印到 $0$ 時不再補空白，最後換行。

\stylepoint{輸出一排以空白分隔的值時，用 \texttt{if (i > 0)}（或類似判
斷）避免行尾多餘空白，比事後再處理字串乾淨得多。}

\begin{drills}[label=drills:loop-countdown]
\item \texttt{for (int i = n; i >= 1; i--)} 跟
\texttt{for (int i = n; i >= 0; i--)}，哪一個會多印一個數字？多印的是哪
一個？
\item \texttt{countdown.c} 為什麼要用 \texttt{if (i > 0)} 才印空白，而
不是每個數字後面都印？
\item 若把 \texttt{if (i > 0)} 誤寫成 \texttt{if (i >= 0)}，輸出會有什
麼問題？
\end{drills}

\section{while 迴圈：條件成立就繼續}
\label{sec:loop-while-basics}

\texttt{while} 迴圈適合「不知道要跑幾次，只知道什麼時候該停」的情形。 它只
有條件，沒有把初始化跟更新寫在括號裡。 這些通常要自己在迴圈本體裡或迴圈
前後處理。

\begin{code}{c}{while 迴圈的基本形狀}
while (條件) {
    重複執行的敘述;
}
\end{code}

範例 \thechapter.6 讀入正整數 $n$，只要 $n$ 還是偶數就除以 $2$，直到變
成奇數為止。 這種「一直做某個操作直到狀態改變」的模式，用 \texttt{while}
比 \texttt{for} 自然。

\begin{example}{c}{(divide-until-odd.c) 一直除以 2 直到變成奇數}
#include <stdio.h>

int main(void) {
    int n;
    scanf("%d", &n);

    while (n % 2 == 0) {
        n /= 2;
    }
    printf("%d\n", n);
    return 0;
}
\end{example}

\iosample{
24 \\
}{
3 \\
}

$n = 24$ 時：$24 \to 12 \to 6 \to 3$，三次除以 $2$ 後 $3$ 是奇數，迴圈
結束，印出 $3$。


若本體裡忘了更新條件相關的變數，迴圈可能永遠停不下來：

\begin{code}{c}{危險：忘記更新，形成無窮迴圈}
int i = 0;
while (i < n) {
    printf("%d\n", i);
    // 漏寫 i++：i 永遠是 0，條件永遠成立
}
\end{code}

終端機若卡住不動、風扇開始狂轉，先檢查迴圈裡有沒有把計數器或狀態往終止
條件推進。 開發時用 \texttt{Ctrl-C} 強制結束程式，再回去補上漏掉的更新。

\keypoint{\texttt{while} 迴圈在每一輪\textbf{一開始}先檢查條件；如果一開始
條件就不成立，迴圈本體一次都不會執行。 \texttt{do-while} 則把條件
檢查放在本體之後，至少會跑一次；後面會專門對比。}

\begin{drills}[label=drills:loop-while-basics]
\item \texttt{divide-until-odd.c} 輸入 $24$，迴圈總共跑幾輪？每一輪
\texttt{n} 依序變成什麼？
\item 範例裡若漏寫 \texttt{n /= 2;}，會發生什麼事？
\item \texttt{while} 迴圈在每一輪的\emph{開始}還是\emph{結束}檢查條
件？
\end{drills}

\section{for 跟 while 其實是同一件事}
\label{sec:loop-for-while-equiv}

\texttt{for (初始化; 條件; 更新)} 可以拆成等價的 \texttt{while}：初始
化只做一次，然後「先測條件、再執行本體、最後更新」，再回頭測條件。
條件一開始就不成立時，本體一次都不會跑。 這點 \texttt{for} 與
\texttt{while} 相同，跟 \texttt{do-while} 不同。

\begin{code}{c}{把 for 拆成等價的 while}
初始化;
while (條件) {
    本體;
    更新;
}
\end{code}

範例 \thechapter.9 用 \texttt{while} 印出 $0$ 到 $n-1$，效果等同
\texttt{for (int i = 0; i < n; i++)}。 輸入 \texttt{3} 得到三行
\texttt{0}、\texttt{1}、\texttt{2}。

\begin{example}{c}{(for-as-while.c) 用 while 模擬固定次數的 for}
#include <stdio.h>

int main(void) {
    int n;
    scanf("%d", &n);
    int i = 0;
    while (i < n) {
        printf("%d\n", i);
        i++;
    }
    return 0;
}
\end{example}

\iosample{
3 \\
}{
0 \\
1 \\
2 \\
}

\begin{center}
\begin{tikzpicture}[node distance=8mm and 20mm]
    \node[process] (init) {初始化};
    \node[decision, below=of init] (cond) {條件？};
    \node[process, below=of cond] (body) {本體};
    \node[process, right=of body] (upd) {更新};
    \node[process, below=of body] (done) {離開迴圈};
    \draw[flowarrow] (init.south) -- (cond.north);
    \draw[flowarrow] (cond.south) -- node[left,font=\scriptsize] {是} (body.north);
    \draw[flowarrow] (body.east) -- (upd.west);
    \draw[flowarrow] (upd.north) |- (cond.east);
    \draw[flowarrow] (cond.west) -- ++(-2cm,0) |- node[below,font=\scriptsize,pos=0.25] {否} (done.west);
\end{tikzpicture}
\diagcaption{\texttt{for} 與 \texttt{while} 的共同流程：先測條件。 條
件為假時本體一次也不執行。}
\label{fig:ch05-for-while-flow}
\end{center}

計數器作用域已在第一節說明；這裡專心看 \texttt{for} 怎麼改寫成
\texttt{while}，以及兩者共用的流程。

差一錯誤（off-by-one）幾乎都出在條件用 \texttt{<} 還是 \texttt{<=}。
\texttt{i = 0; i < n} 跑 $n$ 次；\texttt{i = 1; i <= n} 也跑 $n$
次。 兩種都對，混用就容易多跑或少跑一輪。 先在紙上寫出 $i$ 會取哪些
值，再寫條件。

逗號運算子可以把兩個運算式合成一個，值是右邊那個。 偶爾會看到
\texttt{for (i = 0, sum = 0; i < n; i++)} 這種「初始化兩個變數」的寫
法；能讀懂即可，不必為了短而把更新也擠進逗號裡。

\begin{code}{c}{for 的初始化裡用逗號隔開兩個指定}
for (i = 0, sum = 0; i < n; i++) {
    sum += i;
}
\end{code}

\begin{drills}[label=drills:loop-for-while-equiv]
\item 把 \texttt{for (初始化; 條件; 更新)} 拆成等價的 \texttt{while}
時，「更新」要放在迴圈本體的最前面還是最後面？
\item \texttt{for} 與 \texttt{while}，當條件一開始就不成立時，本體各
會執行幾次？
\item \texttt{for-as-while.c} 輸入 $3$，會印出哪三個數字？
\end{drills}

\section{do-while：至少執行一次}
\label{sec:loop-do-while}

\texttt{do-while} 把條件檢查放在迴圈\textbf{本體之後}，所以不管條件一開始
成不成立，本體至少會跑一次。 這在「一定要先問一次、再決定要不要繼續問」
的輸入驗證特別有用。

\begin{code}{c}{do-while 迴圈的基本形狀}
do {
    重複執行的敘述;
} while (條件);
\end{code}

範例 \thechapter.12 一直讀整數，直到讀到嚴格正數為止。 使用者可能第一次就
輸入錯的，但我們仍然要先讀進來才能判斷。

\begin{example}{c}{(read-until-positive.c) 讀到嚴格正數為止}
#include <stdio.h>

int main(void) {
    int x;

    do {
        scanf("%d", &x);
    } while (x <= 0);

    printf("%d\n", x);
    return 0;
}
\end{example}

\iosample{
-1 \\
0 \\
5 \\
}{
5 \\
}

第一次讀到 $-1$、第二次 $0$，條件 \texttt{x <= 0} 都成立，繼續讀；第三
次讀到 $5$，條件不成立，跳出迴圈。

\keypoint{\texttt{do-while} 保證迴圈本體至少執行一次；輸入驗證、選單、
「重試直到成功」這類題目，幾乎都適合用 \texttt{do-while} 而不是
\texttt{while}。}

\texttt{while} 與 \texttt{for} 都可能一次都不進本體。 \texttt{do-while}
相反：先做，再問。 範例 \thechapter.13 把 \texttt{n} 設成 $5$，迴圈條
件卻是 \texttt{n < 0}。 條件從頭到尾都是假，但仍會印出一行 \texttt{5}，
因為本體在測試之前就跑過了。

\begin{example}{c}{(do-once.c) 條件為假時，do-while 仍執行一次}
#include <stdio.h>

int main(void) {
    int n = 5;
    do {
        printf("%d\n", n);
        n++;
    } while (n < 0);
    return 0;
}
\end{example}

\iosample{}{
5 \\
}

這就是為什麼輸入驗證愛用 \texttt{do-while}：你總得先讀一次，才知道合
不合法。

\begin{drills}[label=drills:loop-do-while]
\item \texttt{do-while} 跟 \texttt{while} 最大的差別是什麼？
\item \texttt{do-once.c} 裡 \texttt{n = 5}，迴圈條件是
\texttt{n < 0}（從頭到尾都是假），但程式仍會印出什麼？為什麼？
\item \texttt{read-until-positive.c} 若第一個輸入就是正數，迴圈本體
會執行幾次？
\end{drills}

\section{break、continue 與提早離開迴圈}
\label{sec:loop-break-continue}

有時候迴圈不必跑完所有輪次。 找到答案就可以 \texttt{break} 跳出；某些
輪次要跳過則用 \texttt{continue} 直接進下一輪。 後面的質數判定與
Collatz 序列都會再用到 \texttt{break}。

\begin{code}{c}{在 for 迴圈裡用 break 提早結束}
for (int i = 0; i < n; i++) {
    if (找到答案) {
        break;
    }
}
\end{code}


範例 \thechapter.15 讀入 $n$，從 $1$ 掃到 $n$，找到第一個 $7$ 的倍數就
\texttt{break}，印出那個數；若都沒有就印 \texttt{-1}。

\begin{example}{c}{(first-multiple-of-7.c) 找到第一個 7 的倍數就停}
#include <stdio.h>

int main(void) {
    int n;
    scanf("%d", &n);

    int found = -1;
    for (int i = 1; i <= n; i++) {
        if (i % 7 == 0) {
            found = i;
            break;
        }
    }
    printf("%d\n", found);
    return 0;
}
\end{example}

\iosample{
10 \\
}{
7 \\
}

\iosample{
5 \\
}{
-1 \\
}

\begin{code}{c}{在 for 迴圈裡用 continue 跳過本輪}
for (int i = 0; i < n; i++) {
    if (這一輪要跳過) {
        continue;
    }
    處理本輪;
}
\end{code}

\stylepoint{一個迴圈盡量只有一個入口、一個出口，可讀性較高；但「找到就
停」這類明確的 \texttt{break}，通常比設一堆旗標變數更直覺。 不要為了
避免 \texttt{break} 而把邏輯寫繞。}

\texttt{continue} 適合「這一輪剩下的事可以整段不理」。 範例
\thechapter.17 讀 $n$ 個整數，只把正數加進總和；非正數一碰到就
\texttt{continue}，後面的 \texttt{sum += x} 不會執行。 輸入
\texttt{4} 筆資料 \texttt{7 -2 3 8}，略過 $-2$，得到 $18$。

\begin{example}{c}{(sum-positives.c) 用 continue 跳過非正數}
#include <stdio.h>

int main(void) {
    int n;
    scanf("%d", &n);
    int sum = 0;
    for (int count = 0; count < n; count++) {
        int x;
        scanf("%d", &x);
        if (x <= 0) {
            continue;
        }
        sum += x;
    }
    printf("%d\n", sum);
    return 0;
}
\end{example}

\iosample{
4 7 -2 3 8 \\
}{
18 \\
}

同一題也可以寫成「只有正數才加」，如範例 \thechapter.18。 結果一樣，
但讀程式的人看到 \texttt{if (x > 0)} 往往會往下找 \texttt{else}。 其
實沒有 else，只是「其餘情況什麼都不做」。 \texttt{continue} 把「不理
這一輪」寫在條件成立的當下，後面的加總邏輯就專心處理正數。

\begin{example}{c}{(sum-positives-if.c) 不用 continue 的等價寫法}
#include <stdio.h>

int main(void) {
    int n;
    scanf("%d", &n);
    int sum = 0;
    for (int count = 0; count < n; count++) {
        int x;
        scanf("%d", &x);
        if (x > 0) {
            sum += x;
        }
    }
    printf("%d\n", sum);
    return 0;
}
\end{example}

\iosample{
4 7 -2 3 8 \\
}{
18 \\
}

迴圈本體裡宣告的 \texttt{x}，離開那一對大括號就不能再用。 這正好符合
「它只是這一輪的暫存」。 宣告和使用靠在一起，比把 \texttt{x} 提到
\texttt{main} 開頭更清楚。

\keypoint{\texttt{continue} 結束的是\textbf{這一輪}，不是整個迴圈；
\texttt{break} 才是整段離開。 兩者都不要拿來代替本來該寫清楚的條件。}

\begin{drills}[label=drills:loop-break-continue]
\item \texttt{break} 跟 \texttt{continue}，哪一個是整段離開迴圈，哪
一個只結束這一輪？
\item \texttt{first-multiple-of-7.c} 輸入 $5$（找 $1$ 到 $5$ 裡第一個
$7$ 的倍數），會印出什麼？
\item \texttt{sum-positives.c} 輸入 \texttt{4 7 -2 3 8}，\texttt{continue}
跳過的是哪一個數字？最後 \texttt{sum} 是多少？
\end{drills}

\section{兩段式輸出：先上後下}
\label{sec:loop-two-stage}

有些題目要依序做兩段不同但結構類似的輸出。 與其寫一個複雜的迴圈，不如老老
實實寫兩個 \texttt{for}：第一段 $1$ 到 $n$，第二段 $n-1$ 到 $1$。 可讀性
通常比硬塞在一個迴圈裡好。

\begin{example}{c}{(up-and-down.c) 先印 1 到 n，再印 n-1 到 1}
#include <stdio.h>

int main(void) {
    int n;
    scanf("%d", &n);

    for (int i = 1; i <= n; i++) {
        printf("%d\n", i);
    }
    for (int i = n - 1; i >= 1; i--) {
        printf("%d\n", i);
    }
    return 0;
}
\end{example}

\iosample{
3 \\
}{
1 \\
2 \\
3 \\
2 \\
1 \\
}

兩個迴圈各自獨立，第二段從 \texttt{n - 1} 開始而不是 \texttt{n}，避免
$n$ 被印兩次。 這種「邊界少印或多印一次」是迴圈題目最常見的失誤之一。

\stylepoint{兩段結構相似的輸出，用兩個清楚的迴圈分開寫，通常比一個迴圈夾
雜 \texttt{if} 判斷「現在是上升還是下降段」更容易讀、也比較不容易算錯邊
界。}

\begin{drills}[label=drills:loop-two-stage]
\item 範例 \texttt{up-and-down.c} 裡，第二個 \texttt{for} 為什麼要從
\texttt{n - 1} 開始，而不是從 \texttt{n} 開始？
\item 如果兩段都從 \texttt{n} 開始遞減／遞增，$n = 3$ 的輸出會多印出
什麼？
\end{drills}

\section{巢狀迴圈：矩形與數字圖形}
\label{sec:loop-nested-shapes}

當外層迴圈控制「第幾列」、內層控制「這一列印幾次或印什麼」，就形成\textbf{
巢狀迴圈}（nested loop）。 內層每跑完一輪，外層才前進一格。 總執行次數
大約是外層次數乘以內層次數。

\begin{code}{c}{巢狀 for 迴圈印矩形}
for (int row = 0; row < height; row++) {
    for (int col = 0; col < width; col++) {
        printf("#");
    }
    printf("\n");
}
\end{code}

範例 \thechapter.21 讀入寬 $W$、高 $H$，印出 $H$ 行、每行 $W$ 個
\texttt{\#}。

\begin{example}{c}{(print-rectangle.c) 印 W x H 的井號矩形}
#include <stdio.h>

int main(void) {
    int w, h;
    scanf("%d %d", &w, &h);

    for (int row = 0; row < h; row++) {
        for (int col = 0; col < w; col++) {
            printf("#");
        }
        printf("\n");
    }
    return 0;
}
\end{example}

\iosample{
4 3 \\
}{
\#\#\#\# \\
\#\#\#\# \\
\#\#\#\# \\
}

範例 \thechapter.22 第 $i$ 列印數字 $i$ 共 $i$ 次。 內層次數跟外層變數
綁在一起，形成三角形。

\begin{example}{c}{(number-triangle.c) 第 i 列重複數字 i 共 i 次}
#include <stdio.h>

int main(void) {
    int n;
    scanf("%d", &n);

    for (int i = 1; i <= n; i++) {
        for (int j = 0; j < i; j++) {
            printf("%d", i);
        }
        printf("\n");
    }
    return 0;
}
\end{example}

\iosample{
4 \\
}{
1 \\
22 \\
333 \\
4444 \\
}

範例 \thechapter.23 第 $i$ 列印 $1$ 到 $i$。 同一個外層 \texttt{i}，
內層從 $1$ 數到 \texttt{i}，圖形從「重複同一數字」變成「遞增序列」。

\begin{example}{c}{(repeated-pattern.c) 第 i 列印 1 到 i}
#include <stdio.h>

int main(void) {
    int n;
    scanf("%d", &n);

    for (int i = 1; i <= n; i++) {
        for (int j = 1; j <= i; j++) {
            printf("%d", j);
        }
        printf("\n");
    }
    return 0;
}
\end{example}

\iosample{
4 \\
}{
1 \\
12 \\
123 \\
1234 \\
}

\stylepoint{巢狀迴圈寫圖形時，先想清楚「外層控制列還是行、內層控制這一列要
印幾個字元或什麼內容」，在紙上畫小表格對一下，比直接寫程式瞎試省時。}

\begin{drills}[label=drills:loop-nested-shapes]
\item 巢狀迴圈裡，外層通常控制什麼、內層通常控制什麼？
\item \texttt{number-triangle.c} 第 $i$ 列印數字 $i$ 共幾次？$n = 4$
時第 $3$ 列印出什麼？
\item \texttt{repeated-pattern.c} 跟 \texttt{number-triangle.c} 印出的
圖形有什麼不同？
\end{drills}

\section{乘法表與週期圖樣}
\label{sec:loop-mult-table}

$N \times N$ 乘法表是巢狀迴圈的經典練習：外層 \texttt{i}、內層 \texttt{j}，
印 \texttt{i * j}，並用 \texttt{\%2d} 對齊欄寬。

\begin{example}{c}{(multiplication-table.c) N x N 乘法表}
#include <stdio.h>

int main(void) {
    int n;
    scanf("%d", &n);

    for (int i = 1; i <= n; i++) {
        for (int j = 1; j <= n; j++) {
            printf("%2d", i * j);
        }
        printf("\n");
    }
    return 0;
}
\end{example}

\iosample{
3 \\
}{
 1 2 3 \\
 2 4 6 \\
 3 6 9 \\
}

範例 \thechapter.25 在 $N \times N$ 網格上，行列索引之和能被 $3$ 整除的
格子印 \texttt{@}，其餘印 \texttt{-}。 這類「週期圖樣」題目，關鍵往往是
找到 \texttt{(i + j) \% k} 或類似規律。

\begin{example}{c}{(diagonal-tile.c) 週期對角線圖樣}
#include <stdio.h>

int main(void) {
    int n;
    scanf("%d", &n);

    for (int i = 0; i < n; i++) {
        for (int j = 0; j < n; j++) {
            if ((i + j) % 3 == 0) {
                printf("@");
            } else {
                printf("-");
            }
        }
        printf("\n");
    }
    return 0;
}
\end{example}

\iosample{
4 \\
}{
@--@ \\
--@- \\
-@-- \\
@--@ \\
}

\begin{drills}[label=drills:loop-mult-table]
\item \texttt{multiplication-table.c} 裡 \texttt{printf("\%2d", ...)}
的 \texttt{\%2d} 是做什麼用的？
\item \texttt{diagonal-tile.c} 用 \texttt{(i + j) \% 3 == 0} 判斷要印
\texttt{@} 還是 \texttt{-}；$n = 4$ 時第一列（\texttt{i = 0}）有哪幾欄
是 \texttt{@}？
\end{drills}

\section{置中的數字金字塔}
\label{sec:loop-pyramid}

巢狀迴圈不一定印實心矩形。 有時內層要先印「左邊的空白（或底線）」，
再印這一列的內容。 範例 \thechapter.26 第 $i$ 列先印 $n-i$ 個底線，再
把數字 $i$ 重複 $i$ 次。 $n = 3$ 時外形會像金字塔。

\begin{example}{c}{(number-pyramid.c) 用底線對齊的數字金字塔}
#include <stdio.h>

int main(void) {
    int n;
    scanf("%d", &n);
    for (int i = 1; i <= n; i++) {
        for (int s = 0; s < n - i; s++) {
            printf("_");
        }
        for (int k = 0; k < i; k++) {
            printf("%d", i);
        }
        printf("\n");
    }
    return 0;
}
\end{example}

\iosample{
3 \\
}{
\_\_1 \\
\_22 \\
333 \\
}

外層仍管「第幾列」，內層拆成兩段：一段專管左邊界、一段專管內容。 先
在紙上把 $n=3$ 的三列畫出來，再數底線個數，比直接在終端機裡試縮排快。

\begin{drills}[label=drills:loop-pyramid]
\item \texttt{number-pyramid.c} 第 $i$ 列先印幾個底線、再印幾個數
字？
\item 內層的兩段 \texttt{for} 分別負責圖形的哪一部分？
\end{drills}

\section{do-while 驗證輸入：密碼、分數與字元}
\label{sec:loop-input-validation}

同一套「讀入 → 檢查 → 不對就再讀」可以套在不同型別的輸入上。 範例
\thechapter.27 模擬 PIN 碼驗證：不是 \texttt{1234} 就一直重試。

\begin{example}{c}{(pin-until-match.c) 重試直到 PIN 正確}
#include <stdio.h>

int main(void) {
    int pin;

    do {
        scanf("%d", &pin);
    } while (pin != 1234);

    printf("OK\n");
    return 0;
}
\end{example}

\iosample{
9999 \\
1234 \\
}{
OK \\
}

範例 \thechapter.28 限制分數在 $0$ 到 $100$ 之間。 條件變成區間檢查，
結構完全一樣。

\begin{example}{c}{(valid-grade.c) 讀到 0 到 100 的分數為止}
#include <stdio.h>

int main(void) {
    int score;

    do {
        scanf("%d", &score);
    } while (score < 0 || score > 100);

    printf("%d\n", score);
    return 0;
}
\end{example}

\iosample{
120 \\
-5 \\
85 \\
}{
85 \\
}

讀字元時，空白字元（空格、換行）常會留在輸入緩衝區裡干擾下一次
\texttt{scanf}。 在 \texttt{\%c} 前面加一個空格。 寫成 \texttt{" \%c"}。 
會讓 \texttt{scanf} 先跳過前面的空白再讀一個字元。 範例 \thechapter.29
用這個技巧讀 \texttt{Y} 或 \texttt{N}。

\begin{example}{c}{(yes-or-no.c) 讀到 Y 或 N 為止}
#include <stdio.h>

int main(void) {
    char c;

    do {
        scanf(" %c", &c);
    } while (c != 'Y' && c != 'N');

    printf("%c\n", c);
    return 0;
}
\end{example}

\iosample{
A \\
M \\
Y \\
}{
Y \\
}

\stylepoint{用 \texttt{scanf(" \%c", \&c)} 讀單一字元時，格式字串開頭的空
格會跳過換行與空白，避免上一次 \texttt{scanf} 留下的換行被當成輸入讀進
來。}

範例 \thechapter.30 讀星期（$1$ 到 $7$），不在範圍就再讀，直
到合法為止。 輸入 \texttt{0 8 3} 時前兩次被丟掉，最後印 \texttt{3}。

\begin{example}{c}{(valid-weekday.c) 讀到 1 到 7 的星期為止}
#include <stdio.h>

int main(void) {
    int d;
    do {
        scanf("%d", &d);
    } while (d < 1 || d > 7);
    printf("%d\n", d);
    return 0;
}
\end{example}

\iosample{
0 8 3 \\
}{
3 \\
}
\begin{drills}[label=drills:loop-input-validation]
\item 這一節的每個例子都用 \texttt{do-while} 而不是 \texttt{while}
驗證輸入，為什麼？
\item \texttt{scanf(" \%c", \&c)} 格式字串開頭的空格，作用是什麼？
\item \texttt{valid-grade.c} 輸入依序是 \texttt{120}、\texttt{-5}、
\texttt{85}，迴圈本體執行幾次？最後印出什麼？
\end{drills}

\section{哨兵值、選單與累加}
\label{sec:loop-sentinel}

\textbf{哨兵值}（sentinel）是一個約定好的特殊輸入，代表「資料結束了」。 
例如 $-1$ 或 $0$。 搭配 \texttt{do-while} 或 \texttt{while}，可以讀一
串未知長度的資料並做統計。

範例 \thechapter.31 模擬選單：選 $0$ 離開，其他選項印 \texttt{AGAIN}。

\begin{example}{c}{(menu-until-quit.c) 選單迴圈直到選 0}
#include <stdio.h>

int main(void) {
    int choice;

    do {
        scanf("%d", &choice);
        if (choice != 0) {
            printf("AGAIN\n");
        }
    } while (choice != 0);

    printf("BYE\n");
    return 0;
}
\end{example}

\iosample{
1 \\
2 \\
0 \\
}{
AGAIN \\
AGAIN \\
BYE \\
}

範例 \thechapter.32 在讀到 $-1$ 之前，數有幾筆有效資料（$-1$ 本身不計）。

\begin{example}{c}{(sum-until-sentinel.c) 以 -1 為哨兵值計數}
#include <stdio.h>

int main(void) {
    int x, count = 0;

    do {
        scanf("%d", &x);
        if (x != -1) {
            count++;
        }
    } while (x != -1);

    printf("%d\n", count);
    return 0;
}
\end{example}

\iosample{
1 \\
2 \\
3 \\
-1 \\
}{
3 \\
}

範例 \thechapter.33 則把哨兵改成 $0$，讀到的所有數字（包含最後那個 $0$）
加總。 注意這裡 $0$ 既參與累加，又觸發結束。

\begin{example}{c}{(sum-until-zero.c) 累加直到讀到 0}
#include <stdio.h>

int main(void) {
    int x, sum = 0;

    do {
        scanf("%d", &x);
        sum += x;
    } while (x != 0);

    printf("%d\n", sum);
    return 0;
}
\end{example}

\iosample{
1 \\
2 \\
3 \\
0 \\
}{
6 \\
}

\keypoint{哨兵值要不要參與計算，完全看題目規定；寫程式前先確認「結束符號
本身算不算一筆資料」，比寫完再改邏輯省時。}

\begin{drills}[label=drills:loop-sentinel]
\item 什麼是「哨兵值」？
\item \texttt{sum-until-sentinel.c}（以 $-1$ 為哨兵）跟
\texttt{sum-until-zero.c}（以 $0$ 為哨兵）最大的差別是什麼？
\item \texttt{menu-until-quit.c} 輸入 \texttt{1 2 0}，會依序印出哪幾
行？
\end{drills}

\section{多筆測資：已知筆數與 Case 前綴}
\label{sec:loop-multi-case}

競賽與題庫常見的輸入格式是：第一行給測資筆數 $N$，後面接 $N$ 組資料。
這時外層用 \texttt{for (int i = 0; i < n; i++)} 包住內層處理即可。

\begin{example}{c}{(multi-case-compare.c) N 組資料比較兩數大小}
#include <stdio.h>

int main(void) {
    int n;
    scanf("%d", &n);

    for (int i = 0; i < n; i++) {
        int a, b;
        scanf("%d %d", &a, &b);
        if (a > b) {
            printf("A\n");
        } else if (a < b) {
            printf("B\n");
        } else {
            printf("T\n");
        }
    }
    return 0;
}
\end{example}

\iosample{
3 \\
1 2 \\
2 2 \\
3 1 \\
}{
B \\
T \\
A \\
}

有些題目要求每組輸出前面加 \texttt{Case k:} 前綴。 範例 \thechapter.35
把迴圈計數器從 $1$ 開始（\texttt{k = 1; k <= n; k++}），直接當 Case 編
號。

\begin{example}{c}{(case-prefix.c) 多筆測資加上 Case 前綴}
#include <stdio.h>

int main(void) {
    int n;
    scanf("%d", &n);

    for (int k = 1; k <= n; k++) {
        int a, b;
        scanf("%d %d", &a, &b);
        printf("Case %d: ", k);
        if (a > b) {
            printf("A\n");
        } else if (a < b) {
            printf("B\n");
        } else {
            printf("T\n");
        }
    }
    return 0;
}
\end{example}

\iosample{
2 \\
1 2 \\
2 2 \\
}{
Case 1: B \\
Case 2: T \\
}

測資筆數未知、直到 EOF（檔案結束）為止時，可以用 \texttt{while (scanf(...)
!= EOF)} 這種寫法。 檔案處理那一章會更完整地讲，這裡先記住「外層迴圈控
制還有沒有下一組資料」這個概念就好。

\begin{code}{c}{一直讀到 EOF 為止}
while (scanf("%d %d", &a, &b) != EOF) {
    處理這一組資料;
}
\end{code}

另一種常見終止條件是「某個特殊值組合」。 例如兩個 \texttt{0} 代表結束。 
邏輯跟哨兵值相同，只是形式從單一數字變成多個欄位。


範例 \thechapter.37 是解題平台上最常見的「測資筆數未知、讀到 EOF 為止」模板：
每讀到兩個整數就印它們的和，沒有下一行時 \texttt{scanf} 回傳 \texttt{EOF}、
迴圈結束。 檔案與標準輸入重定向的細節留給後面章節；這裡先記住這個外層
\texttt{while} 骨架。

\begin{example}{c}{(sum-until-eof.c) 一直讀兩整數相加，直到 EOF}
#include <stdio.h>

int main(void) {
    int a, b;
    while (scanf("%d %d", &a, &b) != EOF) {
        printf("%d\n", a + b);
    }
    return 0;
}
\end{example}

\iosample{
1 2 \\
3 4 \\
}{
3 \\
7 \\
}


\begin{drills}[label=drills:loop-multi-case]
\item 「先給測資筆數 $N$，再接 $N$ 組資料」這種輸入格式，通常用什麼
迴圈結構處理？
\item \texttt{case-prefix.c} 跟 \texttt{multi-case-compare.c} 比起
來，多做了什麼事？
\item 測資筆數\emph{未知}、要讀到檔案結束為止時，可以用哪種
\texttt{while} 寫法？
\end{drills}

\section{連續整數求和}
\label{sec:loop-range-sum}

求 $A + (A+1) + \cdots + B$ 這類問題，用迴圈從 $A$ 加到 $B$ 最直接。
如果輸入可能 $A > B$，先交換再算。 這比寫兩個方向各算一遍更清楚。

\begin{code}{c}{若 a 大於 b 就先交換}
if (a > b) {
    int t = a;
    a = b;
    b = t;
}
for (int i = a; i <= b; i++) {
    sum += i;
}
\end{code}

範例 \thechapter.39 除了印總和，還把算式 $a + (a+1) + \cdots = s$ 完整
輸出。 迴圈裡用旗標 \texttt{first} 控制加號，避免算式開頭多一個
\texttt{+}。

\begin{example}{c}{(sum-equation.c) 連續整數求和並印出算式}
#include <stdio.h>

int main(void) {
    int a, b;
    scanf("%d %d", &a, &b);

    if (a > b) {
        int t = a;
        a = b;
        b = t;
    }

    long long sum = 0;
    int first = 1;
    for (int i = a; i <= b; i++) {
        if (!first) {
            printf("+");
        }
        printf("%d", i);
        sum += i;
        first = 0;
    }
    printf("=%lld\n", sum);
    return 0;
}
\end{example}

\iosample{
1 5 \\
}{
1+2+3+4+5=15 \\
}

\keypoint{連續整數求和雖然有公式 $(A+B)(B-A+1)/2$，但用迴圈實作更直覺，
也為之後「不是等差數列、不能套公式」的累加題打好基礎。}

\begin{drills}[label=drills:loop-range-sum]
\item 如果輸入的 $a$ 大於 $b$，程式在開始累加前會先做什麼？
\item \texttt{sum-equation.c} 裡的旗標 \texttt{first} 是用來避免什麼
問題？
\end{drills}

\section{累加的變奏：次方和、下一個倍數、最小公倍數}
\label{sec:loop-accum-variants}

同一層 \texttt{for} 可以同時累加好幾種量。 範例 \thechapter.40 讀
\texttt{n}，一次算出 $1$ 到 $n$ 的和、平方和、立方和。 $n = 4$ 時是
$10$、$30$、$100$。

\begin{example}{c}{(sum-powers.c) 連續整數的和、平方和、立方和}
#include <stdio.h>

int main(void) {
    int n;
    scanf("%d", &n);
    int sum = 0;
    int sq = 0;
    int cu = 0;
    for (int i = 1; i <= n; i++) {
        sum += i;
        sq += i * i;
        cu += i * i * i;
    }
    printf("%d %d %d\n", sum, sq, cu);
    return 0;
}
\end{example}

\iosample{
4 \\
}{
10 30 100 \\
}

有些 \texttt{while} 不是在「數次數」，而是在「把狀態推到下一個合格
點」。 範例 \thechapter.41 從 $0$ 每次加 $5$，直到不小於輸入的
\texttt{j}。 也就是「$5$ 的倍數裡，不小於 $j$ 的最小那個」。
\texttt{j = 17} 得到 $20$；\texttt{j = 15} 得到 $15$。

\begin{example}{c}{(next-five.c) 不小於 j 的最小 5 的倍數}
#include <stdio.h>

int main(void) {
    int j;
    scanf("%d", &j);
    int i = 0;
    while (i < j) {
        i += 5;
    }
    printf("%d\n", i);
    return 0;
}
\end{example}

\iosample{
17 \\
}{
20 \\
}

最小公倍數可以從已經會的輾轉相除再走一步：$\mathrm{lcm}(a,b) =
ab / \gcd(a,b)$。 為了少溢位，先除後乘：\texttt{a / gcd * b}。 範例
\thechapter.42 對 \texttt{12 18} 印出 \texttt{36}。

\begin{example}{c}{(lcm-two.c) 用 gcd 迴圈求最小公倍數}
#include <stdio.h>

int main(void) {
    int a, b;
    scanf("%d %d", &a, &b);
    int x = a;
    int y = b;
    while (y != 0) {
        int t = x % y;
        x = y;
        y = t;
    }
    printf("%d\n", a / x * b);
    return 0;
}
\end{example}

\iosample{
12 18 \\
}{
36 \\
}

\begin{drills}[label=drills:loop-accum-variants]
\item \texttt{next-five.c} 輸入 \texttt{j = 17}，\texttt{i} 從 $0$
開始每次加 $5$，最後印出多少？
\item \texttt{lcm-two.c} 怎麼從 \texttt{gcd} 算出最小公倍數？
\item 為什麼要寫成 \texttt{a / gcd * b} 而不是 \texttt{a * b / gcd}？
\end{drills}

\section{位數處理：剝除與反轉}
\label{sec:loop-digits}

整數在記憶體裡是以二進位存的，但我們常需要按\textbf{十進位}的個位、十位
來處理。 取餘數 \texttt{n \% 10} 得到最後一位，再 \texttt{n /= 10} 把
最後一位「剝掉」。 這個 \texttt{while (n > 0)} 模式在數論題目裡非常常見。

\begin{example}{c}{(digit-sum.c) 各位數字相加}
#include <stdio.h>

int main(void) {
    int n;
    scanf("%d", &n);

    int s = 0;
    while (n > 0) {
        s += n % 10;
        n /= 10;
    }
    printf("%d\n", s);
    return 0;
}
\end{example}

\iosample{
12345 \\
}{
15 \\
}

範例 \thechapter.44 用同一套剝除邏輯，但把每一位「接」到新整數的末尾，
達成數字反轉。

\begin{example}{c}{(reverse-digits.c) 反轉整數的十進位數字}
#include <stdio.h>

int main(void) {
    int n;
    scanf("%d", &n);

    int rev = 0;
    while (n > 0) {
        rev = rev * 10 + n % 10;
        n /= 10;
    }
    printf("%d\n", rev);
    return 0;
}
\end{example}

\iosample{
12345 \\
}{
54321 \\
}

以 $12345$ 走一遍反轉：\texttt{rev} 依序變成 $5 \to 54 \to 543 \to
5432 \to 54321$。 每次 \texttt{rev * 10} 把舊結果左移一位，再加上新的
個位。

\stylepoint{位數處理的 \texttt{while (n > 0)} 迴圈，每一輪只做兩件事：
\texttt{n \% 10} 取個位、\texttt{n /= 10} 剝掉個位。 把這個節奏記熟，
很多題目都能套同一個骨架。}

\begin{drills}[label=drills:loop-digits]
\item \texttt{digit-sum.c} 用哪兩個運算子，把整數的個位剝出來、再把
它去掉？
\item \texttt{reverse-digits.c} 裡 \texttt{rev = rev * 10 + n \% 10;}
這一行在做什麼？
\end{drills}

\section{累乘、階乘型乘積與狀態更新}
\label{sec:loop-product}

連加對應的是 \texttt{sum += ...}，連乘則是 \texttt{product *= ...}。
排列數 $P(N,R) = N \times (N-1) \times \cdots \times (N-R+1)$ 需要
\texttt{long long}，因為乘積很快超過 \texttt{int} 範圍。

\begin{example}{c}{(falling-factorial.c) 計算 P(N,R)}
#include <stdio.h>

int main(void) {
    int n, r;
    scanf("%d %d", &n, &r);

    long long ans = 1;
    for (int i = 0; i < r; i++) {
        ans *= (n - i);
    }
    printf("%lld\n", ans);
    return 0;
}
\end{example}

\iosample{
5 3 \\
}{
60 \\
}

有些 \texttt{while} 迴圈每一輪要同時更新好幾個變數。 範例 \thechapter.46
模擬細菌：每天族群乘以 $3$，直到達到目標。

\begin{example}{c}{(bacteria-growth.c) 三倍成長直到達標}
#include <stdio.h>

int main(void) {
    int pop, target;
    scanf("%d %d", &pop, &target);

    int days = 0;
    while (pop < target) {
        pop *= 3;
        days++;
    }
    printf("%d\n", days);
    return 0;
}
\end{example}

\iosample{
1 27 \\
}{
3 \\
}

範例 \thechapter.47 則是「兌換迴圈」：每 $5$ 根棍子可換 $1$ 根，換來的
又能繼續湊。 迴圈裡同時更新「目前手上」與「累計總數」兩個狀態。

\begin{example}{c}{(exchange-sticks.c) 每 5 根兌換 1 根}
#include <stdio.h>

int main(void) {
    int sticks;
    scanf("%d", &sticks);

    int total = sticks;
    while (sticks >= 5) {
        int gained = sticks / 5;
        total += gained;
        sticks = sticks % 5 + gained;
    }
    printf("%d\n", total);
    return 0;
}
\end{example}

\iosample{
23 \\
}{
28 \\
}

$23$ 根：先換 $4$ 根（總 $27$，剩 $7$），再換 $1$ 根（總 $28$，剩 $3$），
$3 < 5$ 結束。

\keypoint{累乘、累加一樣要選夠大的型別；乘積比加法更容易溢出，看到排列、
組合、階乘這類關鍵字，優先考慮 \texttt{long long}。}

\begin{drills}[label=drills:loop-product]
\item \texttt{falling-factorial.c} 為什麼用 \texttt{long long} 存
\texttt{ans}，而不是 \texttt{int}？
\item \texttt{bacteria-growth.c} 的 \texttt{while} 迴圈裡，每一輪同時
更新哪兩個變數？
\end{drills}

\section{歐幾里得算法與質數判定}
\label{sec:loop-gcd-prime}

\textbf{輾轉相除法}（Euclid 算法）用「大數除以小數、取餘數、把小數當新
的大數」反覆進行，直到餘數為 $0$，此時的小數就是最大公因數。 整個過程
就是一個 \texttt{while (b != 0)} 迴圈。

\begin{example}{c}{(gcd-euclid.c) 輾轉相除法求最大公因數}
#include <stdio.h>

int main(void) {
    int a, b;
    scanf("%d %d", &a, &b);

    while (b != 0) {
        int t = a % b;
        a = b;
        b = t;
    }
    printf("%d\n", a);
    return 0;
}
\end{example}

\iosample{
48 18 \\
}{
6 \\
}

判定質數則用 \texttt{for} 從 $2$ 試除到 $\sqrt{n}$。 只要找到任何一個
因數就可以用前面學過的 \texttt{break} 提早結束，不必試完所有數字。

\begin{example}{c}{(prime-check.c) 試除法判定質數}
#include <stdio.h>

int main(void) {
    int n;
    scanf("%d", &n);

    if (n < 2) {
        printf("NO\n");
        return 0;
    }

    int isPrime = 1;
    for (int d = 2; (long long)d * d <= n; d++) {
        if (n % d == 0) {
            isPrime = 0;
            break;
        }
    }
    printf("%s\n", isPrime ? "YES" : "NO");
    return 0;
}
\end{example}

\iosample{
17 \\
}{
YES \\
}

\keypoint{試除到 $\sqrt{n}$ 就夠，因為若 $n = ab$ 且 $a \le b$，則
$a \le \sqrt{n}$；用 \texttt{(long long)d * d <= n} 避免 \texttt{d * d}
在 \texttt{d} 很大時溢出。}

\begin{drills}[label=drills:loop-gcd-prime]
\item 輾轉相除法裡，\texttt{while (b != 0)} 迴圈結束時，哪個變數存的
就是最大公因數？
\item \texttt{prime-check.c} 為什麼只試除到 $\sqrt{n}$ 就夠，不用試到
$n$？
\item 條件寫成 \texttt{(long long)d * d <= n} 而不是 \texttt{d * d <= n}
的目的是什麼？
\end{drills}

\section{Collatz 序列}
\label{sec:loop-collatz}

$3n+1$ 問題（Collatz 猜想）：若 $n$ 是偶數就除以 $2$，奇數就變成
$3n+1$，重複直到變成 $1$。 沒人證明是否一定會到 $1$，但程式可以忠實
模擬這個過程。

範例 \thechapter.50 印出整段序列；範例 \thechapter.51 只數長度（包含起
點）。 兩者都適合用 \texttt{while (n != 1)}，差別在要不要把中間每一步
都印出來。

\begin{example}{c}{(collatz-sequence.c) 印出 Collatz 序列}
#include <stdio.h>

int main(void) {
    int n;
    scanf("%d", &n);

    int first = 1;
    while (1) {
        if (!first) {
            printf(" ");
        }
        printf("%d", n);
        first = 0;
        if (n == 1) {
            break;
        }
        if (n % 2 == 0) {
            n /= 2;
        } else {
            n = 3 * n + 1;
        }
    }
    printf("\n");
    return 0;
}
\end{example}

\iosample{
6 \\
}{
6 3 10 5 16 8 4 2 1 \\
}

\begin{example}{c}{(collatz-length.c) 計算 Collatz 序列長度}
#include <stdio.h>

int main(void) {
    int n;
    scanf("%d", &n);

    int len = 1;
    while (n != 1) {
        if (n % 2 == 0) {
            n /= 2;
        } else {
            n = 3 * n + 1;
        }
        len++;
    }
    printf("%d\n", len);
    return 0;
}
\end{example}

\iosample{
6 \\
}{
9 \\
}

\keypoint{Collatz 迴圈裡，奇數步 \texttt{3 * n + 1} 可能讓 \texttt{n}
短時間內變很大；若題目上限很大，要確認 \texttt{int} 是否夠用，必要時
改用 \texttt{long long}。}

\begin{drills}[label=drills:loop-collatz]
\item Collatz 規則：$n$ 是偶數時怎麼變？是奇數時怎麼變？
\item \texttt{collatz-length.c} 跟 \texttt{collatz-sequence.c} 用的是
同一套規則，兩支程式的輸出差在哪裡？
\end{drills}

\section{因數和、餘數分類與同號連續段}
\label{sec:loop-divisor-classify}

把 $1$ 到 $n-1$ 中所有整除 $n$ 的因數加起來，可以判斷 $n$ 是完數
（Perfect）、盈數（Abundant）還是虧數（Deficient）。 範例 \thechapter.52
用一個 \texttt{for} 累加真因數和。

\begin{example}{c}{(classify-divisors.c) 依真因數和分類整數}
#include <stdio.h>

int main(void) {
    int n;
    scanf("%d", &n);

    int sum = 0;
    for (int d = 1; d < n; d++) {
        if (n % d == 0) {
            sum += d;
        }
    }

    if (sum == n) {
        printf("Perfect\n");
    } else if (sum > n) {
        printf("Abundant\n");
    } else {
        printf("Deficient\n");
    }
    return 0;
}
\end{example}

\iosample{
28 \\
}{
Perfect \\
}

範例 \thechapter.53 讀 $n$ 個整數，依除以 $3$ 的餘數分三類計數。 這
跟陣列章的「分桶」概念類似，但這裡只用三個計數器，不必開陣列。

\begin{example}{c}{(residue-count.c) 依 mod 3 餘數分類計數}
#include <stdio.h>

int main(void) {
    int n;
    scanf("%d", &n);

    int c0 = 0, c1 = 0, c2 = 0;
    for (int i = 0; i < n; i++) {
        int x;
        scanf("%d", &x);
        if (x % 3 == 0) {
            c0++;
        } else if (x % 3 == 1) {
            c1++;
        } else {
            c2++;
        }
    }
    printf("%d %d %d\n", c0, c1, c2);
    return 0;
}
\end{example}

\iosample{
5 \\
3 4 5 6 7 \\
}{
2 2 1 \\
}

範例 \thechapter.54 處理「連續同號整數的長度」。 掃描時若符號跟上一個
相同就延長計數，否則輸出上一段並重新開始。 這是\textbf{run-length}（游
程編碼）的雛形，之後字串、序列題目還會再遇到。

\begin{example}{c}{(run-length-sign.c) 同號連續段長度}
#include <stdio.h>

int main(void) {
    int n;
    scanf("%d", &n);

    int prev, count = 0, first = 1;
    for (int i = 0; i < n; i++) {
        int x;
        scanf("%d", &x);
        if (i == 0) {
            prev = x;
            count = 1;
            continue;
        }
        if ((prev > 0 && x > 0) || (prev < 0 && x < 0)) {
            count++;
        } else {
            if (!first) {
                printf(" ");
            }
            printf("%d", (prev > 0 ? 1 : -1) * count);
            first = 0;
            prev = x;
            count = 1;
        }
    }

    if (n > 0) {
        if (!first) {
            printf(" ");
        }
        printf("%d\n", (prev > 0 ? 1 : -1) * count);
    }
    return 0;
}
\end{example}

\iosample{
6 \\
1 2 3 -1 -2 5 \\
}{
3 -2 1 \\
}

正數段 $1,2,3$ 長度 $3$ 輸出 \texttt{3}；負數段 $-1,-2$ 長度 $2$ 輸出
\texttt{-2}（符號代表正負）；最後單獨的 $5$ 輸出 \texttt{1}。

\begin{drills}[label=drills:loop-divisor-classify]
\item \texttt{classify-divisors.c} 怎麼判斷一個數是不是「完數」
（Perfect）？
\item \texttt{residue-count.c} 用了幾個計數器？分別對應除以 $3$ 餘幾
的情況？
\item \texttt{run-length-sign.c} 在處理的是哪一種常見模式（跟後面字
串章節會再遇到的哪個概念一樣）？
\end{drills}

\section{日期計算：從年月日推算星期}
\label{sec:loop-day-of-week}

日期題目通常要處理：各月天數不同、閏年二月多一天、從某個基準日累加總天
數再取餘數。 範例 \thechapter.55 簡化版：假設該年 1 月 1 日是星期五
（編號 $6$，若星期日為 $0$），讀入月日後累加從年初到該日前的天數。

\begin{example}{c}{(day-of-week.c) 由月日推算星期（簡化基準）}
#include <stdio.h>

int isLeap(int y) {
    return (y % 400 == 0) || (y % 4 == 0 && y % 100 != 0);
}

int main(void) {
    int y, m, d;
    scanf("%d %d %d", &y, &m, &d);

    int days[] = {0, 31, 28, 31, 30, 31, 30, 31, 31, 30, 31, 30, 31};
    if (isLeap(y)) {
        days[2] = 29;
    }

    int total = 0;
    for (int i = 1; i < m; i++) {
        total += days[i];
    }
    total += d - 1;

    int w = (total + 6) % 7;
    printf("%d\n", w);
    return 0;
}
\end{example}

\iosample{
2024 3 1 \\
}{
3 \\
}

閏年判斷 \texttt{isLeap} 抽出成函式，主程式用 \texttt{for} 累加整月天
數。 這種「查表 + 累加 + 取餘」的組合，是日期類題目的共通骨架。 更完整
的月曆輸出（格子排版、 weekday offset）只是把同一套天數計算嵌進巢狀
迴圈裡控制列與欄。

\keypoint{日期題目裡，閏年規則（四年一閏、百年不閏、四百年再閏）跟「各
月天數表」一定要先寫對；邏輯錯在這裡，後面再漂亮的迴圈也救不了。}

\begin{drills}[label=drills:loop-day-of-week]
\item \texttt{isLeap} 函式判斷閏年的三個條件是什麼？
\item 為什麼程式要在確定是不是閏年之後，才把 \texttt{days[2]} 改成
$29$？
\end{drills}

\section{進階模擬：多狀態與複合迴圈}
\label{sec:loop-advanced-sim}

題庫後段（分數 $80$ 以上）的題目，往往把多種技巧疊在一起：在網格上模擬
機器人移動、光線在鏡面間反彈、多個容器之間反覆倒水直到平衡、或是在狀態
空間裡用 \texttt{while} 迭代直到收斂。 這些題目沒有單一「萬用片語」，
但共通結構是：

\begin{enumerate}
    \item 用變數（或陣列）完整描述「系統現狀」；
    \item 外層 \texttt{while} 控制「模擬還要不要繼續」；
    \item 內層 \texttt{for} 處理「這一輪裡每個物件各做一步」；
    \item 每步更新後檢查終止條件（碰撞、填滿、空池、超過步數上限等）。
\end{enumerate}

練習時建議先在紙上對小測資手工模擬幾輪，確認「每一輪哪些變數會變、什麼
情況該停」，再把它翻成程式。 若單純用 \texttt{for} 掃一次無法表達「要
重複調整直到穩定」，就該懷疑需要外層 \texttt{while}。 這是從基礎迴圈
過渡到模擬題的分水嶺。

題庫裡像「兩球在格子裡彈跳並交換位置」「多個杯子與無限水池之間倒水」
「機器人找燃料並在環狀地圖上相遇」這類高分題，都建立在本章前面所有
技巧之上：輸入驗證、哨兵、多筆測資、巢狀迴圈、\texttt{break}、以及多
變數狀態更新。 把這一章的基礎練熟，再回頭挑戰那些題目，會比一開始就
硬啃模擬題有效率得多。

\begin{drills}[label=drills:loop-advanced-sim]
\item 進階模擬題的共通結構列出的四個步驟依序是什麼？
\item 當「單純用 \texttt{for} 掃一次無法表達」時，通常代表需要哪一種
額外的迴圈結構？
\end{drills}

\addcontentsline{toc}{section}{練習解答}
\section*{練習解答}

各節練習的答案收在這裡。 短題只給能一句話核對的結果。

\paragraph{第~\ref{sec:loop-for-basics}~節}
\begin{enumerate}
\item 一次都不會執行：條件 \texttt{0 < 0} 一開始就不成立。
\item 不能，\texttt{i} 的作用域只到這一句 \texttt{for} 結束為止。
\item \texttt{i < n} 跑 $n$ 次；\texttt{i <= n} 跑 $n + 1$ 次。
\end{enumerate}

\paragraph{第~\ref{sec:loop-countdown}~節}
\begin{enumerate}
\item \texttt{i >= 0} 那個會多印一個；多印的是 $0$。
\item 因為最後一個數字（$0$）後面不該有空白；\texttt{if (i > 0)} 讓
只有非最後一個數字才補空白。
\item 每個數字後面都會多印一個空白，行尾會多出一個空格。
\end{enumerate}

\paragraph{第~\ref{sec:loop-while-basics}~節}
\begin{enumerate}
\item 跑三輪：$24 \to 12 \to 6 \to 3$，$3$ 是奇數時結束。
\item 條件永遠成立（\texttt{n} 不會變），形成無窮迴圈。
\item 開始。
\end{enumerate}

\paragraph{第~\ref{sec:loop-for-while-equiv}~節}
\begin{enumerate}
\item 放在本體的最後面（先執行本體，再更新，才回頭測條件）。
\item 都是 $0$ 次。
\item $0$、$1$、$2$。
\end{enumerate}

\paragraph{第~\ref{sec:loop-do-while}~節}
\begin{enumerate}
\item \texttt{do-while} 保證本體至少執行一次；\texttt{while} 若一開始
條件就不成立，可能一次都不執行。
\item 印出 \texttt{5}：因為本體在條件被檢查之前就先跑過一次。
\item 一次：讀到正數後條件 \texttt{x <= 0} 不成立，迴圈立刻結束。
\end{enumerate}

\paragraph{第~\ref{sec:loop-break-continue}~節}
\begin{enumerate}
\item \texttt{break} 整段離開迴圈；\texttt{continue} 只結束這一輪，
接著跑下一輪。
\item 印出 \texttt{-1}：$1$ 到 $5$ 裡沒有 $7$ 的倍數。
\item 跳過的是 \texttt{-2}；\texttt{sum = 7 + 3 + 8 = 18}。
\end{enumerate}

\paragraph{第~\ref{sec:loop-two-stage}~節}
\begin{enumerate}
\item 避免 $n$ 被印兩次——第一段已經印過 $n$ 了。
\item $n$ 會被連續印兩次（一次在第一段結尾，一次在第二段開頭）。
\end{enumerate}

\paragraph{第~\ref{sec:loop-nested-shapes}~節}
\begin{enumerate}
\item 外層通常控制「第幾列」，內層控制「這一列印幾次或印什麼」。
\item 印 $i$ 次；第 $3$ 列印出 \texttt{333}。
\item \texttt{number-triangle.c} 每列重複同一個數字；
\texttt{repeated-pattern.c} 每列印遞增序列 $1$ 到 $i$。
\end{enumerate}

\paragraph{第~\ref{sec:loop-mult-table}~節}
\begin{enumerate}
\item 把每個數字用固定寬度（$2$ 格）對齊輸出，讓欄位排整齊。
\item $j = 0, 3$ 時 $(0+j) \% 3 == 0$ 成立，也就是第 $0$、$3$ 欄印
\texttt{@}。
\end{enumerate}

\paragraph{第~\ref{sec:loop-pyramid}~節}
\begin{enumerate}
\item 先印 $n - i$ 個底線，再印 $i$ 個數字。
\item 第一段負責左邊的留白（底線），第二段負責這一列的內容（數字）。
\end{enumerate}

\paragraph{第~\ref{sec:loop-input-validation}~節}
\begin{enumerate}
\item 因為一定要先讀一次才能判斷合不合法，\texttt{do-while} 保證本體
至少執行一次。
\item 跳過輸入裡殘留的空白字元（包含前一次讀取留下的換行）。
\item 三次；最後印出 \texttt{85}。
\end{enumerate}

\paragraph{第~\ref{sec:loop-sentinel}~節}
\begin{enumerate}
\item 一個約定好、代表「資料結束了」的特殊輸入值。
\item \texttt{-1} 那版不把哨兵本身計入結果；\texttt{0} 那版連最後讀到
的 $0$ 都一起加進總和。
\item 依序印出 \texttt{AGAIN}、\texttt{AGAIN}、\texttt{BYE}。
\end{enumerate}

\paragraph{第~\ref{sec:loop-multi-case}~節}
\begin{enumerate}
\item 外層用 \texttt{for (int i = 0; i < n; i++)} 包住處理每一組資料
的邏輯。
\item \texttt{case-prefix.c} 多印了 \texttt{Case k:} 這個前綴。
\item \texttt{while (scanf(...) != EOF)}。
\end{enumerate}

\paragraph{第~\ref{sec:loop-range-sum}~節}
\begin{enumerate}
\item 先把 $a$、$b$ 交換，確保 $a \le b$ 再累加。
\item 避免算式開頭多印一個 \texttt{+} 號。
\end{enumerate}

\paragraph{第~\ref{sec:loop-accum-variants}~節}
\begin{enumerate}
\item $20$（$5$ 的倍數裡不小於 $17$ 的最小值）。
\item 先用 \texttt{while} 迴圈算出 \texttt{gcd}，再用
\texttt{a / gcd * b}。
\item 避免 \texttt{a * b} 在中間步驟就超過型別能表示的範圍（先除可以
讓中間值小一點）。
\end{enumerate}

\paragraph{第~\ref{sec:loop-digits}~節}
\begin{enumerate}
\item \texttt{\% 10} 取出個位，\texttt{/= 10} 把個位剝掉。
\item 把新讀到的個位「接」到 \texttt{rev} 的末尾，達成反轉的效果。
\end{enumerate}

\paragraph{第~\ref{sec:loop-product}~節}
\begin{enumerate}
\item 因為排列數的乘積很快就會超過 \texttt{int} 能表示的範圍。
\item 同時更新族群數量 \texttt{pop}（乘以 $3$）與天數 \texttt{days}
（加 $1$）。
\end{enumerate}

\paragraph{第~\ref{sec:loop-gcd-prime}~節}
\begin{enumerate}
\item \texttt{a}：迴圈結束時 \texttt{a} 就是最大公因數。
\item 因為若 $n = a \times b$ 且 $a \le b$，則必有 $a \le \sqrt{n}$，
超過 $\sqrt{n}$ 不會再找到新的因數。
\item 避免 \texttt{d * d} 在 \texttt{d} 很大時，用 \texttt{int} 運算會
先溢位。
\end{enumerate}

\paragraph{第~\ref{sec:loop-collatz}~節}
\begin{enumerate}
\item 偶數除以 $2$；奇數變成 $3n + 1$。
\item \texttt{collatz-length.c} 只印最終的步數；\texttt{collatz-sequence.c}
把過程中每一步的值都印出來。
\end{enumerate}

\paragraph{第~\ref{sec:loop-divisor-classify}~節}
\begin{enumerate}
\item 真因數和 \texttt{sum} 剛好等於 $n$ 本身時就是完數。
\item 三個計數器 \texttt{c0}、\texttt{c1}、\texttt{c2}，分別對應除以
$3$ 餘 $0$、$1$、$2$ 的情況。
\item 游程編碼（run-length encoding）：把連續相同（這裡是同號）的一
段資料記成「值＋長度」。
\end{enumerate}

\paragraph{第~\ref{sec:loop-day-of-week}~節}
\begin{enumerate}
\item 能被 $400$ 整除；或者能被 $4$ 整除但不能被 $100$ 整除。
\item 因為閏年二月是 $29$ 天、平年是 $28$ 天，要先判斷完閏年才知道
\texttt{days[2]} 該填哪個值。
\end{enumerate}

\paragraph{第~\ref{sec:loop-advanced-sim}~節}
\begin{enumerate}
\item 描述系統現狀、外層 \texttt{while} 控制模擬要不要繼續、內層
\texttt{for} 處理每個物件各走一步、每步後檢查終止條件。
\item 通常代表需要外層 \texttt{while}（因為要「重複調整直到穩定」，
次數不是事先就知道的）。
\end{enumerate}

\chapter{陣列}
\label{ch:array}

陣列（array）讓我們可以用一個名字，儲存一整排相同型別的資料，並用索引
（index）存取其中的每一個元素。 在這之前，我們只能宣告一個一個獨立的變數，
如果要處理 $100$ 個數字，就得寫 $100$ 個變數名稱，非常不切實際。 更麻煩
的是，變數的數量在寫程式的當下就要決定好；如果題目說「$n$ 個數字，$n$
由輸入決定」，我們根本沒辦法事先知道要宣告幾個變數。

有了陣列，我們只需要一個名字，再搭配迴圈，就能有系統地處理一整排資料，
不管 $n$ 實際上是多少，程式的寫法都一樣。 這也是陣列跟迴圈幾乎總是成對
出現的原因：陣列負責「存」，迴圈負責「一個一個處理」。

這一章分三段推進：先把宣告、初始化、越界與「什麼時候其實不需要陣列」講
清楚，再進入一維陣列的核心解題模式（兩次掃描、雙指標、分流、查表、分桶、
滑動視窗），最後才談二維陣列在記憶體裡怎麼排列、怎麼走訪與建構。 沿路上
我們也會開始接觸一些跟記憶體本身有關的細節。 陣列的大小為什麼要先決定、
索引為什麼從 $0$ 開始、資料實際上是怎麼一塊一塊排列的。 這些細節在陣列
這一章可能只是小提醒，但到了指標那一章會變成核心觀念。

\section{宣告與讀入陣列}
\label{sec:array-declare-read}

要使用陣列，我們必須先宣告陣列的型別與大小，再用迴圈把資料一個一個讀進去。
我們先看一個很小、很具體的例子：宣告一個只能放 $5$ 個整數的陣列，把 $5$
個數字讀進去，再依序印出來。

\begin{code}{c}{宣告一個大小為 5 的整數陣列並讀入、印出資料}
int a[5];
for (int i = 0; i < 5; i++) {
    scanf("%d", &a[i]);
}
for (int i = 0; i < 5; i++) {
    printf("%d\n", a[i]);
}
\end{code}

\texttt{int a[5]} 跟系統要了一段連續的記憶體，剛好放得下 $5$ 個
\texttt{int}。 這 $5$ 個位置各自有一個編號，叫做\emph{索引}（index），
分別是 \texttt{a[0]}、\texttt{a[1]}、\texttt{a[2]}、\texttt{a[3]}、
\texttt{a[4]}。 編號是從 $0$ 開始，不是從 $1$，所以最後一個位置是
\texttt{a[4]} 而不是 \texttt{a[5]}（這件事很重要，我們馬上會再細講）。
上面的程式先用一個迴圈把 $5$ 個數字依序存進 \texttt{a[0]} 到
\texttt{a[4]}，再用另一個迴圈依序印出來；兩個迴圈裡的變數 \texttt{i}
每一輪剛好對應到一個索引，範圍都是 $0$ 到 $4$。

\keypoint{陣列的大小要在宣告時用常數決定，宣告之後就固定不變；
\texttt{int a[5]} 有 $5$ 個位置，索引是 \texttt{a[0]} 到 \texttt{a[4]}。}

但實際解題時，陣列的大小往往不能像上面這樣直接寫死成 $5$。 題目通常只
會說「有 $n$ 筆資料」，而 $n$ 是執行期才從輸入讀進來的，寫程式的當下並
不知道確切數字，也就沒辦法像 \texttt{a[5]} 這樣把大小寫成一個固定的字面
值。 這時候常見的做法，是先看題目給的上限（例如「$n \le 1000$」），宣告
一個「夠大」的陣列，再用實際讀到的 $n$ 決定要用到前面幾格。 片語
\thechapter.3 就是這種寫法：宣告 \texttt{a[1005]}（比 $1000$ 多留一點空
間），但實際只使用前 $n$ 格；跟剛才 \texttt{a[5]} 的例子比起來，只是把
「讀 5 次」換成「先讀 $n$，再讀 $n$ 次」，陣列本身的用法完全一樣。

\begin{code}{c}{宣告並讀入一個大小為 n 的整數陣列}
int a[1005], n;
scanf("%d", &n);
for (int i = 0; i < n; i++) {
    scanf("%d", &a[i]);
}
\end{code}

\keypoint{陣列的大小要在宣告時用常數決定，通常會比題目上限多留一點空間
（例如 $n \le 1000$ 就宣告 \texttt{a[1005]}）。}

這裡有一個從第一天用陣列就會遇到、之後也會一直跟著我們的規則：C 語言的
陣列索引從 $0$ 開始，不是從 $1$ 開始。 一個宣告成 \texttt{a[1005]} 的陣
列，合法的索引是 $0$ 到 $1004$，總共 $1005$ 個位置，\texttt{a[1005]} 本
身反而是不合法的（超出了我們要來的那塊記憶體）。 這個規則不是 C 語言隨
意的選擇：索引其實代表的是「跟陣列開頭差了幾格」的\emph{位移量}
（offset），而不是「這是第幾個元素」這種日常語感裡從 $1$ 算起的排名。
陣列的第一個元素跟開頭的位移量是 $0$（因為它就是開頭本身），所以索引自
然從 $0$ 開始；陣列的最後一個元素，跟開頭差了 $n - 1$ 格，索引就是
$n - 1$。 這件事現在只需要記住結論，指標那一章會看到 \texttt{a[i]} 實
際上就是「從 \texttt{a} 的位址往後移動 \texttt{i} 個元素大小」，到時候
「索引 = 位移量」這件事會變得非常具體。

\keypoint{陣列索引從 $0$ 開始，是因為索引代表的是離開頭的位移量，不是
「第幾個」的排名；大小為 $n$ 的陣列合法索引是 $0$ 到 $n - 1$，存取
\texttt{a[n]} 是超出範圍的錯誤，即使它「感覺」應該是最後一個元素。}

\begin{center}
\begin{tikzpicture}
    \def\vals{1,4,8,9,10}
    \foreach \v [count=\i from 0] in \vals {
        \pgfmathtruncatemacro{\iprev}{\i-1}
        \ifnum\i=0
            \node[cell] (c\i) at (0,0) {\v};
        \else
            \node[cell, right=0pt of c\iprev] (c\i) {\v};
        \fi
        \node[idxlbl, below=2pt of c\i] {\i};
    }
\end{tikzpicture}
\diagcaption{陣列 \texttt{a} 的內容與索引；灰色小字是索引，索引從 $0$
開始，最後一個元素 \texttt{10} 的索引是 $n - 1 = 4$，不是 $5$。}
\label{fig:ch06-array-basic}
\end{center}

有了陣列之後，我們不只可以照原順序把資料印出來，也可以用任何順序走訪它
。 陣列存在記憶體裡的順序，跟我們讀取／使用它的順序，兩者是獨立的。
範例 \thechapter.3 讀入 $n$ 個整數後，反向印出整個序列：讀進來的順序是
從第一個到最後一個，但輸出的迴圈刻意從索引 $n - 1$ 開始往回跑到
$0$，資料本身完全沒有搬動位置，改變的只是我們走訪它的方向。

\begin{example}{c}{(print-reverse.c) 反向輸出序列}
#include <stdio.h>

int main(void) {
    int a[1005], n;

    // 讀入 n 個整數
    scanf("%d", &n);
    for (int i = 0; i < n; i++) {
        scanf("%d", &a[i]);
    }

    // 反向印出，元素之間留一個空白
    for (int i = n - 1; i >= 0; i--) {
        printf("%d", a[i]);
        if (i) {
            printf(" ");
        }
    }
}
\end{example}

\iosample{
5 \\
1 4 8 9 10 \\
}{
10 9 8 4 1 \\
}

以範例 \thechapter.3 的輸入為例：$n = 5$，陣列依序存了
$a[0]=1, a[1]=4, a[2]=8, a[3]=9, a[4]=10$。 反向迴圈第一次執行時
$i = 4$，印出 $a[4] = 10$；接著 $i = 3$，印出 $a[3] = 9$；依此類推，直
到 $i = 0$ 印出 $a[0] = 1$ 之後迴圈結束。 每次印完一個數字，程式都要判
斷「這是不是最後一個要印的元素」，決定要不要在後面補一個空白。

\stylepoint{輸出一整排以空白分隔的數字時，最後一個元素後面不要印多餘的
空白；範例 \thechapter.3 用 \texttt{if (i)} 判斷還沒印到最後一個元素才印
空白。}

\begin{drills}[label=drills:array-declare-read]
\item \texttt{int a[1005];} 宣告後，合法的索引範圍是多少？
\texttt{a[1005]} 本身合法嗎？
\item 陣列索引為什麼從 $0$ 開始，而不是從 $1$ 開始？
\item 範例 print-reverse.c 反向印出陣列時，陣列本身存的資料有沒有被
搬動或改變？
\item 把 $5$ 個數字讀進陣列、再依序印出來，需要幾個迴圈？各自的用途
是什麼？
\end{drills}

\section{陣列的初始化與限制}
\label{sec:array-init-limits}

到目前為止，我們的陣列都是先宣告、再用迴圈把值一個一個讀進去。 但如果
陣列的內容在寫程式的當下就已經知道（例如一份固定的對照表），可以直接用
\emph{初始化}（initializer）在宣告的同時填進去，不用等到執行期才填。

\begin{code}{c}{陣列初始化的幾種寫法}
int a[5] = {1, 2, 3, 4, 5};   // 完整初始化

int b[5] = {1, 2};            // 只給前兩個值，其餘自動補 0
                               // b 的內容是 {1, 2, 0, 0, 0}

int c[5] = {0};                // 把整個陣列歸零的慣用寫法

int d[] = {1, 2, 3, 4, 5};    // 省略長度，編譯器依初始值數量算出長度 5
\end{code}

\keypoint{\texttt{int c[5] = \{0\};} 是把整個陣列歸零最常見的寫法：初
始化值的數量比陣列小時，剩下的元素會自動補 $0$，不用一個一個列出來。}

範例 \thechapter.5 用初始化建立一個「每個月天數」的對照表，並示範另一個
好用的技巧：用 \texttt{sizeof} 自動算出陣列長度，不用另外記或手動數。

\begin{example}{c}{(days-in-month.c) 用初始化建立對照表}
#include <stdio.h>

// 平年每個月的天數（1 月到 12 月）
const int daysInMonth[] = {31, 28, 31, 30, 31, 30, 31, 31, 30, 31, 30, 31};

int main(void) {
    int month;
    scanf("%d", &month);
    printf("%d\n", daysInMonth[month - 1]);

    // 用 sizeof 算出陣列長度，不用手動數有幾個月
    int numMonths = sizeof(daysInMonth) / sizeof(daysInMonth[0]);
    int totalDays = 0;
    for (int i = 0; i < numMonths; i++) {
        totalDays += daysInMonth[i];
    }
    printf("%d\n", totalDays);
}
\end{example}

\iosample{
2 \\
}{
28 \\
365 \\
}

\texttt{sizeof(daysInMonth)} 是整個陣列佔用的位元組數，
\texttt{sizeof(daysInMonth[0])} 是單一個元素佔用的位元組數，兩者相除就
是元素個數。 這個寫法的好處是：就算之後把對照表的內容改長或改短，這一
行完全不用跟著修改。

\keypoint{\texttt{sizeof(a) / sizeof(a[0])} 是計算陣列長度的經典寫法；
陣列前面加上 \texttt{const} 則是告訴編譯器「這個陣列的內容不會被修
改」，編譯器會在我們不小心寫壞它時直接報錯，對像對照表這種本來就不該
被更動的資料很有幫助。}

陣列雖然可以用初始化「一次填好」，但\textbf{不能}在宣告之後用
\texttt{=} 把一整個陣列指定給另一個陣列。 這是宣告時初始化用的
\texttt{=} 和一般賦值運算子看起來一樣、但規則完全不同的地方。

\begin{code}{c}{陣列不能整個賦值}
int a[5] = {1, 2, 3, 4, 5};
int b[5];

b = a;   // 錯誤：無法通過編譯，陣列不能整個賦值

for (int i = 0; i < 5; i++) {
    b[i] = a[i];   // 正確：一個元素一個元素複製
}
\end{code}

\keypoint{要把一個陣列的內容複製到另一個陣列，只能用迴圈逐一複製元素
（或之後會遇到的 \texttt{memcpy} 函式），不能直接寫
\texttt{b = a;}。 這跟一般變數的賦值行為不一樣，是陣列作為一個「整
體」時少數不能直接操作的地方之一。}

\begin{drills}[label=drills:array-init-limits]
\item \texttt{int c[5] = \{0\};} 之後，陣列裡每個元素分別是多少？
\item \texttt{int b[5] = \{1, 2\};} 之後，\texttt{b[2]}、\texttt{b[3]}、
\texttt{b[4]} 分別是多少？
\item 為什麼不能寫 \texttt{b = a;} 來複製整個陣列？該怎麼做才對？
\item \texttt{sizeof(a) / sizeof(a[0])} 算出來的是什麼？
\end{drills}

\section{越界存取：一個真實的當機示範}
\label{sec:array-out-of-bounds}

C 語言完全不會在執行期檢查陣列索引有沒有超出範圍。 這不是「編譯器比較
寬鬆」而已。 存取一個超出範圍的索引，屬於\textbf{未定義行為}
（undefined behavior）：語言標準完全沒有規定會發生什麼事，程式可能正常
跑完、可能印出亂七八糟的數字、可能當掉，結果完全取決於編譯器、作業系統，
甚至那一次執行當下的記憶體狀態。

範例 \thechapter.7 故意把迴圈條件寫錯（\texttt{i <= 5} 而不是
\texttt{i < 5}），對一個宣告成 \texttt{a[5]} 的陣列多存取了一格。

\begin{example}{c}{(out-of-bounds-demo.c) 索引寫錯範圍的示範（請勿模仿）}
#include <stdio.h>

int main(void) {
    int a[5];

    // 錯誤：迴圈多跑了一次，i 會等於 5，超出 a[5] 的合法範圍
    for (int i = 0; i <= 5; i++) {
        a[i] = i * i;
    }

    for (int i = 0; i < 5; i++) {
        printf("%d\n", a[i]);
    }
}
\end{example}

我們實際編譯並執行了這段程式：\texttt{gcc} 用預設選項編譯完全沒有任何
警告，但一執行，程式就非正常終止（收到 SIGABRT 訊號），\emph{連原本應
該印出來的五個數字都沒有顯示}。 因為程式在把緩衝區裡的輸出真正寫到畫面
上之前就已經被中止了。 這正是未定義行為可怕的地方：它不會給你一個乾淨
明確的錯誤訊息。 K.N. King 的教科書裡示範同一種錯誤時，觀察到的是另一種
症狀。 變數 \texttt{i} 剛好被安排在陣列後面的記憶體，多寫的那一格意外
把 \texttt{i} 重設成 $0$，造成迴圈整個變成無窮迴圈。 兩種現象、兩種系統，
但共同的教訓只有一個：陣列的索引範圍必須自己顧好，編譯器不會替你檢查。

\keypoint{陣列索引超出範圍是未定義行為，不是「執行時會跳出一個錯誤訊
息」這種可預期的結果；不同編譯器、不同系統上，同一段越界的程式碼可能
表現出完全不同、甚至看起來「正常」的行為，這也是為什麼不能靠「反正沒
出事」來判斷程式是否正確。}

除了索引本身不能超出範圍，索引\emph{運算式}裡如果帶有副作用（例如
\texttt{++}、\texttt{--}），也要特別小心：

\begin{code}{c}{危險寫法：下標運算式裡帶有副作用}
int i = 0;
while (i < n) {
    a[i] = b[i++];   // 未定義行為：i 在同一運算式裡被讀取又被修改
}
\end{code}

\texttt{a[i] = b[i++]} 這一行，C 語言沒有規定「先算左邊的 \texttt{i}
還是先算右邊的 \texttt{i++}」，同一個變數在同一運算式裡被讀取又被修改，
結果同樣是未定義行為。 想避免這個問題，把遞增獨立成一行就好：

\begin{code}{c}{安全寫法：把遞增拆成獨立的一行，或直接用 for 迴圈}
for (int i = 0; i < n; i++) {
    a[i] = b[i];
}
\end{code}

\stylepoint{下標運算式裡不要同時做「讀取」跟「有副作用的修改」；把遞增
或遞減獨立成一行，或直接改用 \texttt{for} 迴圈管理索引，兩者都能避免
這類難以察覺的未定義行為。}

\begin{drills}[label=drills:array-out-of-bounds]
\item 陣列索引超出範圍時，C 語言會不會在執行期幫忙檢查、報錯？
\item 為什麼「這次執行沒有出事」不能證明一段越界存取的程式碼是對的？
\item \texttt{a[i] = b[i++];} 這一行為什麼是未定義行為？該怎麼改寫比
較安全？
\end{drills}

\section{串流處理：不是每次都需要陣列}
\label{sec:array-stream-no-array}

前面幾個範例都需要先把整排資料存進陣列，才能重複走訪它。 但不是每個
「看起來需要處理一整排數字」的題目都真的需要陣列。 如果題目只要一個
最終的統計值（例如總和、最大值），而不需要事後回頭看某一筆舊資料，那麼
邊讀邊用幾個變數累加，可能就已經足夠，完全不用開陣列。

範例 \thechapter.10 讀入 $n$ 個整數，只需要印出總和與最大值。 這類題目完全不需要事後回頭比較任何一筆舊資料，所以只用兩個變數 \texttt{sum} 和 \texttt{maxVal} 邊讀邊更新就
夠了。

\begin{example}{c}{(stream-sum-max.c) 不存陣列，邊讀邊算總和與最大值}
#include <stdio.h>

int main(void) {
    int n;
    scanf("%d", &n);

    // 不用陣列，只靠兩個變數邊讀邊更新
    long long sum = 0;
    int maxVal = -2000000000;
    for (int i = 0; i < n; i++) {
        int x;
        scanf("%d", &x);
        sum += x;
        if (x > maxVal) {
            maxVal = x;
        }
    }

    printf("%lld %d\n", sum, maxVal);
}
\end{example}

\iosample{
5 \\
1 4 8 9 10 \\
}{
32 10 \\
}

這個判斷值得養成習慣：拿到一個處理一整排數字的題目時，先問自己「這題到
底需不需要事後回頭看某一筆舊資料？」如果答案是否定的，通常代表不需要陣
列，改用幾個變數邊讀邊算，程式更短、用的記憶體也更少。 當 $n$ 大到陣列
可能裝不下時（題庫裡常見的提示是 MLE，記憶體超出限制），這往往是唯一
可行的做法。

\keypoint{陣列不是處理「一整排數字」的必要條件；只有在需要事後回頭走訪
或重複使用舊資料時，才真的需要把資料存起來。}

\begin{drills}[label=drills:array-stream-no-array]
\item 判斷「這題到底需不需要陣列」的關鍵問題是什麼？
\item 範例 stream-sum-max.c 只用兩個變數就算出總和跟最大值，為什麼
不需要陣列？
\item 當 $n$ 大到陣列可能裝不下時（常見的提示是 MLE），為什麼「不用
陣列」往往是唯一可行的做法？
\end{drills}

\section{用陣列做兩次計算：找出第二大的值}
\label{sec:array-second-largest}

上一節的總和與最大值，邊讀邊用兩個變數就算完了。 但有些題目沒
辦法只掃一次就得到答案，需要先掃一遍拿到某個資訊（例如整排資料的最大值），
再掃第二遍利用這個資訊算出真正要的答案。 如果一邊讀輸入一邊處理、又沒
有把資料存下來，讀到第二遍的時候，原始資料早就不見了。 資料已經存在
陣列裡，我們可以照需要的順序、甚至重複走訪它，這是陣列比起單一變數好用
的地方。

範例 \thechapter.11 讀入 $n$ 個整數（保證至少有兩個相異值），印出其中「第二
大」的值。 也就是比最大值小、但在所有數值中最大的那一個。 這個定義聽
起來直覺，但直接動手寫的時候有一個陷阱：如果資料裡有重複的最大值，「第
二大」指的還是要比最大值\emph{嚴格小}的那個值，不是「陣列裡由大到小數
過來的第二個位置」。

\begin{example}{c}{(second-largest.c) 找出第二大的相異值}
#include <stdio.h>

int main(void) {
    int a[1005], n;

    // 讀入 n 個整數
    scanf("%d", &n);
    for (int i = 0; i < n; i++) {
        scanf("%d", &a[i]);
    }

    // 第一遍：找出最大值
    int maxVal = a[0];
    for (int i = 1; i < n; i++) {
        if (a[i] > maxVal) {
            maxVal = a[i];
        }
    }

    // 第二遍：在比最大值小的數之中找最大的一個
    int secondVal = -2000000000;
    for (int i = 0; i < n; i++) {
        if (a[i] > secondVal && a[i] < maxVal) {
            secondVal = a[i];
        }
    }

    printf("%d\n", secondVal);
}
\end{example}

\iosample{
6 \\
5 5 5 3 3 1 \\
}{
3 \\
}

用範例的測資 \texttt{5 5 5 3 3 1} 走一遍邏輯：第一遍掃描後
\texttt{maxVal} 是 $5$。 第二遍掃描時，三個 $5$ 都會因為
\texttt{a[i] < maxVal} 不成立而被跳過，兩個 $3$ 會先後被記錄成目前的
\texttt{secondVal}，最後的 $1$ 因為比 $3$ 小而不會取代它，答案是 $3$。
如果只是單純去找「第二個位置」而不做這個判斷，很容易在有重複最大值時
答錯。

\begin{center}
\begin{tikzpicture}
    \node[hlcell2] (c0) {5};
    \node[hlcell2, right=0pt of c0] (c1) {5};
    \node[hlcell2, right=0pt of c1] (c2) {5};
    \node[hlcell, right=0pt of c2] (c3) {3};
    \node[hlcell, right=0pt of c3] (c4) {3};
    \node[cell, right=0pt of c4] (c5) {1};
    \node[idxlbl, yshift=-0.25cm] at (c1.south) {maxVal = 5};
    \node[idxlbl, yshift=-0.25cm] at ($(c3.south)!0.5!(c4.south)$) {secondVal = 3};
\end{tikzpicture}
\diagcaption{粉色是等於 \texttt{maxVal} 的元素（第二遍會被跳過），橘色
是 \texttt{secondVal} 最後鎖定的元素。}
\label{fig:ch06-second-largest-scan}
\end{center}

\keypoint{「第二大」不是索引為 $1$ 的元素，也不是單純「第二次出現的最大
值」；範例 \thechapter.11 特別排除等於 \texttt{maxVal} 的情形，才能在有
重複值時正確找到 $3$ 而不是 $5$。}

\begin{drills}[label=drills:array-second-largest]
\item 為什麼求「第二大」需要掃描陣列兩次，而不能只掃一次？
\item 陣列 \texttt{5 5 5 3 3 1} 裡，第二大的值是多少？為什麼不是 $5$？
\item 如果只是取「陣列裡第二個位置的值」而不做額外判斷，在有重複最
大值時會答錯什麼？
\end{drills}

\section{原地反轉陣列：交換元素}
\label{sec:array-reverse-swap}

範例 \thechapter.3 用改變走訪方向來「看起來」反向輸出，但陣列本身的內容
完全沒變。 如果題目要求的是把陣列\emph{本身}反轉過來（之後可能還要拿反
轉後的陣列做別的計算，不是印出來就結束），就需要真的把元素互相交換，而
不只是換個方向讀。

一個常見的手法是同時從陣列的頭尾各放一個索引，往中間逼近，每一步把兩邊
的元素互換：左邊的索引從 $0$ 開始，右邊的索引從 $n - 1$ 開始（正好用到
剛剛提過的「最後一個元素索引是 $n - 1$」），直到兩個索引相遇或錯身。

\begin{center}
\begin{tikzpicture}
    \def\vals{1,4,8,9,10}
    \foreach \v [count=\i from 0] in \vals {
        \pgfmathtruncatemacro{\iprev}{\i-1}
        \ifnum\i=0
            \node[hlcell] (c\i) at (0,0) {\v};
        \else\ifnum\i=4
            \node[hlcell, right=0pt of c\iprev] (c\i) {\v};
        \else
            \node[cell, right=0pt of c\iprev] (c\i) {\v};
        \fi\fi
        \node[idxlbl, below=2pt of c\i] {\i};
    }
    \draw[diagarrow] (c0.north) to[bend left=45] node[above,font=\scriptsize] {交換} (c4.north);
\end{tikzpicture}
\diagcaption{\texttt{left}（索引 $0$）與 \texttt{right}（索引
$n-1=4$）指向的元素互換，接著兩個索引各往中間靠一格，重複同樣的動作。}
\label{fig:ch06-reverse-two-pointer}
\end{center}

\begin{example}{c}{(reverse-in-place.c) 用頭尾交換原地反轉陣列}
#include <stdio.h>

int main(void) {
    int a[1005], n;
    scanf("%d", &n);
    for (int i = 0; i < n; i++) {
        scanf("%d", &a[i]);
    }

    // 兩個指標從頭尾往中間逼近，互換元素
    int left = 0, right = n - 1;
    while (left < right) {
        int temp = a[left];
        a[left] = a[right];
        a[right] = temp;
        left++;
        right--;
    }

    for (int i = 0; i < n; i++) {
        if (i) {
            printf(" ");
        }
        printf("%d", a[i]);
    }
    printf("\n");
}
\end{example}

\iosample{
5 \\
1 4 8 9 10 \\
}{
10 9 8 4 1 \\
}

跟同樣的輸入走一遍：一開始 \texttt{left = 0}、\texttt{right = 4}，交換
\texttt{a[0]} 和 \texttt{a[4]}（$1$ 和 $10$互換）；接著
\texttt{left = 1}、\texttt{right = 3}，交換 \texttt{a[1]} 和
\texttt{a[3]}（$4$ 和 $9$ 互換）；此時 \texttt{left = 2 = right}，迴圈
結束（正中間的元素不用跟自己交換）。 陣列變成
$\{10, 9, 8, 4, 1\}$，這次是真的被改寫了，不是只有印出來的順序不同。

\keypoint{交換兩個變數的值一定要借助第三個暫存變數（範例裡的
\texttt{temp}）：如果直接寫 \texttt{a[left] = a[right]; a[right] =
a[left];}，第二行的 \texttt{a[left]} 其實已經被第一行蓋掉了，換到的會
是同一個值。}

\begin{drills}[label=drills:array-reverse-swap]
\item 交換兩個變數的值時，為什麼需要第三個暫存變數？
\item \texttt{left}、\texttt{right} 兩個索引分別從哪裡開始？往哪個方
向逼近？迴圈什麼時候該停止？
\item 範例 reverse-in-place.c 跟範例 print-reverse.c（反向印出）比起
來，本質上差在哪裡？
\end{drills}

\section{把資料依條件分成兩個陣列}
\label{sec:array-partition}

有時候我們要的不是一個值，而是把原始資料\emph{分流}成好幾組，各自存進不
同的陣列。 這時候陣列本身就是輸出的一部分，不只是計算過程中暫存資料
的地方。

範例 \thechapter.13 讀入 $n$ 個整數，依照奇偶把它們分別存進兩個陣列，再
依序印出偶數陣列、奇數陣列。 兩個陣列各自有自己的計數器，記錄目前已經
存了幾個元素。

\begin{center}
\begin{tikzpicture}
    \foreach \v [count=\i from 0] in {1,2,3,4,5,6} {
        \pgfmathtruncatemacro{\iprev}{\i-1}
        \ifnum\i=0
            \node[cell] (in0) {\v};
        \else
            \node[cell, right=0pt of in\iprev] (in\i) {\v};
        \fi
    }
    \node[idxlbl, left=8pt of in0] {輸入};

    \node[hlcell, below=0.2cm of in0] (ev0) {2};
    \node[hlcell, right=0pt of ev0] (ev1) {4};
    \node[hlcell, right=0pt of ev1] (ev2) {6};
    \node[idxlbl, left=8pt of ev0] {evens};

    \node[hlcell2, below=0.2cm of ev0] (od0) {1};
    \node[hlcell2, right=0pt of od0] (od1) {3};
    \node[hlcell2, right=0pt of od1] (od2) {5};
    \node[idxlbl, left=8pt of od0] {oddsArr};
\end{tikzpicture}
\diagcaption{輸入依序掃過一次，偶數依序存進 \texttt{evens}，奇數依序存
進 \texttt{oddsArr}，兩個陣列各自保留原本的先後順序。}
\label{fig:ch06-partition-even-odd}
\end{center}

\begin{example}{c}{(split-by-parity.c) 依奇偶把資料分成兩個陣列}
#include <stdio.h>

int main(void) {
    int n;
    scanf("%d", &n);

    // 依奇偶分別存進兩個陣列
    int evens[1005], oddsArr[1005];
    int evenCount = 0, oddCount = 0;
    for (int i = 0; i < n; i++) {
        int x;
        scanf("%d", &x);
        if (x % 2 == 0) {
            evens[evenCount] = x;
            evenCount++;
        } else {
            oddsArr[oddCount] = x;
            oddCount++;
        }
    }

    for (int i = 0; i < evenCount; i++) {
        if (i) {
            printf(" ");
        }
        printf("%d", evens[i]);
    }
    printf("\n");

    for (int i = 0; i < oddCount; i++) {
        if (i) {
            printf(" ");
        }
        printf("%d", oddsArr[i]);
    }
    printf("\n");
}
\end{example}

\iosample{
6 \\
1 2 3 4 5 6 \\
}{
2 4 6 \\
1 3 5 \\
}

用範例的測資走一遍：讀到 $1$（奇）存進 \texttt{oddsArr[0]}；讀到 $2$
（偶）存進 \texttt{evens[0]}；依此類推，掃完之後
\texttt{evens = \{2, 4, 6\}}、\texttt{oddsArr = \{1, 3, 5\}}，各自保留
了原本在輸入裡出現的先後順序。

\keypoint{兩個陣列各自需要自己的計數器（\texttt{evenCount}、
\texttt{oddCount}），不能共用同一個索引變數，否則其中一個陣列會出現沒
填到值的空洞。}

\begin{drills}[label=drills:array-partition]
\item 範例 split-by-parity.c 為什麼需要兩個獨立的計數器，而不是共用同
一個索引變數？
\item 若輸入依序是 \texttt{1 2 3 4 5 6}，\texttt{evens} 和
\texttt{oddsArr} 掃完之後的內容依序分別是什麼？
\end{drills}

\section{用陣列記錄有沒有看過：以值為索引的查表}
\label{sec:array-lookup}

到目前為止，陣列索引多半代表「第幾筆輸入」。 下一個關鍵跳躍是
讓索引代表「數值本身」：宣告一個跟數值範圍一樣大的陣列，索引就是數值，
內容只記「看過了嗎」的 0 或 1。 這種「查表」（lookup table）手法能把
「這個值之前出現過嗎」的判斷，從一筆一筆比對，變成直接用索引存取，是陣
列在演算法裡最常見的加速手法之一。

範例 \thechapter.14 讀入 $n$ 張卡片的數值（每個數值最多到 $v$），每當某
個數值\emph{第二次}出現，就印出它。 模擬翻牌配對遊戲裡「配到相同的牌」
的瞬間。

\begin{center}
\begin{tikzpicture}
    \node[hlcell] (c0) {3};
    \node[hlcell, right=0pt of c0] (c1) {5};
    \node[hlcell2, right=0pt of c1] (c2) {3};
    \node[hlcell, right=0pt of c2] (c3) {7};
    \node[hlcell2, right=0pt of c3] (c4) {5};
    \node[hlcell2, right=0pt of c4] (c5) {7};
    \node[idxlbl, align=left, anchor=west, above right=0pt and 8pt of c5.east] {橘色：第一次看到};
    \node[idxlbl, align=left, anchor=west, below right=0pt and 8pt of c5.east] {粉色：第二次看到，印出};
\end{tikzpicture}
\diagcaption{依序讀入的六個數值；一個數值第二次出現時，\texttt{seen}
陣列裡它的位置已經是 $1$，程式就把它印出來。}
\label{fig:ch06-seen-array-lookup}
\end{center}

\begin{example}{c}{(seen-lookup.c) 找出第二次出現的數值}
#include <stdio.h>

int main(void) {
    int n, v;
    scanf("%d %d", &n, &v);

    // seen[x] 記錄數值 x 有沒有出現過，索引就是數值本身
    int seen[1005] = {0};
    for (int i = 0; i < n; i++) {
        int x;
        scanf("%d", &x);
        if (seen[x]) {
            printf("%d\n", x);
        }
        seen[x] = 1;
    }
}
\end{example}

\iosample{
6 10 \\
3 5 3 7 5 7 \\
}{
3 \\
5 \\
7 \\
}

依序讀入 $3, 5, 3, 7, 5, 7$：前兩個新數字 $3$、$5$ 第一次出現，只標記
\texttt{seen}、不印出；第三個數字又是 $3$，這時 \texttt{seen[3]} 已經
是 $1$，所以印出 $3$；接下來 $7$ 第一次出現不印，$5$ 第二次出現印出，
最後 $7$ 第二次出現印出，跟輸出的 $3, 5, 7$ 一致。

\keypoint{查表陣列的大小要開到「數值可能的最大範圍」，而不是「資料筆數
$n$」。 範例裡陣列大小取決於數值上限 $v$，如果 $v$ 遠大於 $n$，陣列可
能比輸入資料本身還大很多。}

\begin{drills}[label=drills:array-lookup]
\item \texttt{seen[x]} 這種寫法裡，陣列索引 \texttt{x} 代表的是什麼？
跟前面幾節「索引代表第幾筆輸入」有什麼不同？
\item 查表陣列的大小該取決於資料筆數 $n$，還是數值的可能範圍 $v$？
\item 依序讀入 \texttt{3 5 3 7 5 7}，程式會依序印出哪些數字？
\end{drills}

\section{陣列當作分類桶：用索引記錄統計資料}
\label{sec:array-bucket}

查表陣列裡存的是 0 或 1。 同一套「索引代表類別」的想法還可以
再往前推：索引仍代表某個類別（例如除以 $m$ 的餘數），但格子裡不再只記
有沒有看過，而是把屬於這個類別的數值累加起來。 這種用法常被稱為「桶」
（bucket）或「計數陣列」，是查表手法在統計題上的自然擴充。

範例 \thechapter.15 讀入 $n$ 個整數與一個模數 $m$，把每個數依照除以 $m$
的餘數分類，累加進對應的桶子，最後印出每個桶子的總和。 陣列索引在這裡
不是「第幾個輸入」，而是「餘數是多少」。

\begin{center}
\begin{tikzpicture}
    \foreach \v/\r [count=\i from 0] in {1/1,2/2,3/0,4/1,5/2} {
        \pgfmathtruncatemacro{\iprev}{\i-1}
        \ifnum\i=0
            \node[cell] (v0) {\v};
        \else
            \node[cell, right=0pt of v\iprev] (v\i) {\v};
        \fi
        \node[idxlbl, above=2pt of v\i] {餘\r};
    }
    \node[hlcell, below=0.5cm of v1] (b0) {3};
    \node[hlcell, right=0pt of b0] (b1) {5};
    \node[hlcell, right=0pt of b1] (b2) {7};
    \node[idxlbl, below=2pt of b0] {b[0]};
    \node[idxlbl, below=2pt of b1] {b[1]};
    \node[idxlbl, below=2pt of b2] {b[2]};
\end{tikzpicture}
\diagcaption{每個輸入依照除以 $3$ 的餘數，累加進對應的桶子；例如 $1$ 和
$4$ 餘數都是 $1$，兩者都加進 \texttt{bucketSum[1]}。}
\label{fig:ch06-bucket-count-mod3}
\end{center}

\begin{example}{c}{(remainder-bucket-sum.c) 依餘數分桶並累加}
#include <stdio.h>

int main(void) {
    int n, m;
    scanf("%d %d", &n, &m);

    // 用餘數當作陣列索引，直接累加進對應的桶子
    int bucketSum[1005] = {0};
    for (int i = 0; i < n; i++) {
        int x;
        scanf("%d", &x);
        int b = x % m;
        bucketSum[b] += x;
    }

    for (int b = 0; b < m; b++) {
        if (b) {
            printf(" ");
        }
        printf("%d", bucketSum[b]);
    }
    printf("\n");
}
\end{example}

\iosample{
5 3 \\
1 2 3 4 5 \\
}{
3 5 7 \\
}

$m = 3$ 時，餘數只有 $0, 1, 2$ 三種可能：$1, 4$ 餘 $1$，$2, 5$ 餘
$2$，$3$ 餘 $0$。 所以 \texttt{bucketSum[0] = 3}、
\texttt{bucketSum[1] = 1 + 4 = 5}、\texttt{bucketSum[2] = 2 + 5 = 7}，
跟輸出的 \texttt{3 5 7} 一致。

\keypoint{\texttt{int bucketSum[1005] = \{0\};} 在宣告的同時把整個陣列
初始化成 $0$；如果只宣告不初始化，陣列裡一開始的值是不確定的，直接拿
來累加會得到錯誤的結果。}

\begin{drills}[label=drills:array-bucket]
\item 「查表」跟「分類桶」這兩種用法，陣列格子裡存的東西有什麼不同？
\item \texttt{int bucketSum[1005] = \{0\};} 如果拿掉 \texttt{= \{0\}}
只留宣告，會有什麼風險？
\item $m = 3$、輸入 \texttt{1 2 3 4 5} 時，三個桶 \texttt{bucketSum[0]}、
\texttt{bucketSum[1]}、\texttt{bucketSum[2]} 的總和分別是多少？
\end{drills}

\section{固定範圍的滑動視窗：區塊最大值}
\label{sec:array-sliding-window}

有些題目要求把陣列切成很多個固定大小的區塊，分別對每個區塊做統計。 這
跟前面「桶」的差別在於：桶是依\emph{值}分類，區塊則是依\emph{位置}切
段。 用索引的範圍（而不是索引本身）來決定資料屬於哪一組。

範例 \thechapter.16 讀入 $n$ 個整數與區塊大小 $k$，從頭開始每 $k$ 個分成
一組，對每個「完整」的區塊印出最大值；如果 $n$ 不是 $k$ 的倍數，最後會
剩下一個不足 $k$ 個元素的區塊，這個不完整的區塊改成印最小值，藉此練習
「最後一組要不要特別處理」這個陣列題目常見的邊界情形。

\begin{center}
\begin{tikzpicture}
    \foreach \v [count=\i from 0] in {5,2,9,1,8,3,4,7} {
        \pgfmathtruncatemacro{\iprev}{\i-1}
        \ifnum\i=0
            \node[cell] (c0) {\v};
        \else
            \node[cell, right=0pt of c\iprev] (c\i) {\v};
        \fi
    }
    \draw[decorate,decoration={brace,amplitude=4pt,mirror}] (c0.south west) -- (c2.south east)
        node[midway,below,font=\scriptsize,yshift=-2pt] {最大值 9};
    \draw[decorate,decoration={brace,amplitude=4pt,mirror}] (c3.south west) -- (c5.south east)
        node[midway,below,font=\scriptsize,yshift=-2pt] {最大值 8};
    \draw[decorate,decoration={brace,amplitude=4pt,mirror}] (c6.south west) -- (c7.south east)
        node[midway,below,font=\scriptsize,yshift=-2pt] {不足額，最小值 4};
\end{tikzpicture}
\diagcaption{$n=8$、$k=3$：前兩個區塊剛好 $3$ 個元素，各自報最大值；最
後一個區塊只剩 $2$ 個元素，改報最小值。}
\label{fig:ch06-sliding-window-block-max}
\end{center}

\begin{example}{c}{(block-extremes.c) 逐區塊統計最大值，最後不足額區塊統計最小值}
#include <stdio.h>

int main(void) {
    int n, k;
    scanf("%d %d", &n, &k);

    int a[100005];
    for (int i = 0; i < n; i++) {
        scanf("%d", &a[i]);
    }

    // 每 k 個一組，逐區塊掃描
    for (int start = 0; start < n; start += k) {
        int end = start + k;
        if (end > n) {
            end = n;
        }

        int blockSize = end - start;
        int best = a[start];
        for (int i = start + 1; i < end; i++) {
            if (blockSize == k) {
                if (a[i] > best) {
                    best = a[i];
                }
            } else {
                if (a[i] < best) {
                    best = a[i];
                }
            }
        }
        printf("%d\n", best);
    }
}
\end{example}

\iosample{
8 3 \\
5 2 9 1 8 3 4 7 \\
}{
9 \\
8 \\
4 \\
}

$n = 8$、$k = 3$ 時，區塊依序是 $\{5,2,9\}$（完整，最大值 $9$）、
$\{1,8,3\}$（完整，最大值 $8$）、$\{4,7\}$（只剩 $2$ 個，不完整，改報
最小值 $4$）。 \texttt{start} 每次跳 $k$ 格，\texttt{end} 用
\texttt{if (end > n)} 修正成不超出陣列範圍，這正是區塊型陣列題目最容易
出錯、也最需要小心處理的邊界。

\stylepoint{迴圈裡先算出 \texttt{end} 再修正上限，比每次存取元素時都判
斷「是不是超出範圍」更清楚：邊界只需要處理一次，後面的迴圈本體完全不用
再操心。}

\begin{drills}[label=drills:array-sliding-window]
\item 「桶」跟這一節的「區塊」，兩種分組方式的依據有什麼不同？
\item $n = 8$、$k = 3$ 時，一共分成幾組？最後一組有幾個元素？
\item 為什麼最後一個不足額的區塊，程式改成統計最小值而不是最大值？
\end{drills}

\section{二維陣列與矩陣：從記憶體的真實排列說起}
\label{sec:array-2d-memory}

有些問題本身的資料就是二維的，例如棋盤、表格、或是數學上的矩陣。 這時候
可以宣告二維陣列，用兩個索引 \texttt{grid[row][col]} 分別代表
「第幾列」與「第幾行」。 可以把二維陣列想成「陣列的陣列」：外層索引先
選出第幾列，內層索引再選出那一列裡的第幾格；宣告時兩個維度的大小都要
先用常數決定，跟一維陣列的道理一樣。

\begin{code}{c}{宣告並填滿一個 NxN 的二維整數陣列}
int grid[1000][1000];
for (int r = 0; r < N; r++) {
    for (int c = 0; c < N; c++) {
        grid[r][c] = 0;
    }
}
\end{code}

我們習慣把二維陣列畫成一張表格，但電腦的記憶體本質上是一長條、一維的
空間，沒有「第二維」這種東西。 C 語言處理這個落差的方式，是把二維陣列
按照「先排完第 $0$ 列，再排第 $1$ 列⋯」的順序，攤平成一長條連續的記憶
體，稱為\textbf{列優先}（row-major）排列。

\begin{center}
\begin{tikzpicture}
    \node[hlcell] (r0c0) {1};
    \node[hlcell, right=0pt of r0c0] (r0c1) {2};
    \node[hlcell, right=0pt of r0c1] (r0c2) {3};
    \node[hlcell2, below=0pt of r0c0] (r1c0) {4};
    \node[hlcell2, right=0pt of r1c0] (r1c1) {5};
    \node[hlcell2, right=0pt of r1c1] (r1c2) {6};
    \node[idxlbl, left=8pt of r0c0] {列 0};
    \node[idxlbl, left=8pt of r1c0] {列 1};

    \node[hlcell, below=0.5cm of r1c0] (mem0) {1};
    \node[hlcell, right=0pt of mem0] (mem1) {2};
    \node[hlcell, right=0pt of mem1] (mem2) {3};
    \node[hlcell2, right=0pt of mem2] (mem3) {4};
    \node[hlcell2, right=0pt of mem3] (mem4) {5};
    \node[hlcell2, right=0pt of mem4] (mem5) {6};
    \foreach \i in {0,1,2,3,4,5} {
        \node[idxlbl, below=2pt of mem\i] {\i};
    }
    \node[idxlbl, left=8pt of mem0] {記憶體};

    \draw[diagarrow, black!45] (r0c1.south) to[out=270,in=90] (mem1.north);
    \draw[diagarrow, black!45] (r1c1.south) to[out=270,in=90] (mem4.north);
\end{tikzpicture}
\diagcaption{同一份資料的兩種畫法：上面是我們習慣的二維表格畫法，下面是
它在記憶體裡實際的樣子。 列 $0$ 的三格排完，緊接著才是列 $1$。 顏色對應
同一列：橘色是列 $0$，粉色是列 $1$。}
\label{fig:ch06-2d-array-two-views}
\end{center}

知道這件事之後，我們可以精確算出 \texttt{grid[i][j]} 在這條攤平記憶體
裡的位移量：先跳過前面 $i$ 整列（每列 $C$ 格），再往右走 $j$ 格，也就是
$i \times C + j$。

\begin{center}
\begin{tikzpicture}
    \node[hlcell] (g0c0) {1};
    \node[cell, right=0pt of g0c0] (g0c1) {2};
    \node[cell, right=0pt of g0c1] (g0c2) {3};
    \node[cell, below=0pt of g0c0] (g1c0) {4};
    \node[cell, right=0pt of g1c0] (g1c1) {5};
    \node[hlcell2, right=0pt of g1c1] (g1c2) {6};

    \node[hlcell, below=0.5cm of g1c0] (m0) {1};
    \node[cell, right=0pt of m0] (m1) {2};
    \node[cell, right=0pt of m1] (m2) {3};
    \node[cell, right=0pt of m2] (m3) {4};
    \node[cell, right=0pt of m3] (m4) {5};
    \node[hlcell2, right=0pt of m4] (m5) {6};
    \node[idxlbl, below=2pt of m0] {0};
    \node[idxlbl, below=2pt of m5] {5};

    \node[idxlbl, align=left, right=4.3cm of g0c2] {grid[0][0]\\位移量 $= 0 \times 3 + 0 = 0$};
    \node[idxlbl, align=left, right=4.3cm of g1c2] {grid[1][2]\\位移量 $= 1 \times 3 + 2 = 5$};
\end{tikzpicture}
\diagcaption{一個 $R \times C$ 的二維陣列，\texttt{grid[i][j]} 在記憶體
裡的位移量是 $i \times C + j$；這裡 $C = 3$，所以 \texttt{grid[1][2]}
落在位移量 $5$ 的位置，正好是整個陣列的最後一格。}
\label{fig:ch06-2d-array-memory-layout}
\end{center}

這個「列優先」的事實，會實際影響程式跑起來的快慢。 存取記憶體裡\emph{
緊鄰}的位置，通常比大幅度跳動快得多（這跟 CPU 的快取記憶體 cache 有
關：緊鄰的資料常常會一起被載入，跳著存取則常常要重新載入）。 走訪二維
陣列時，如果內層迴圈跟著「列優先」的方向走（固定列、讓行的索引連續變
化），每一步都剛好移到記憶體裡緊鄰的下一格；反過來，如果內層迴圈固定
行、讓列的索引變化，每一步都要在記憶體裡大幅度跳動。

\begin{center}
\begin{tikzpicture}
    \node[cell] (c0) {1};
    \node[cell, right=0pt of c0] (c1) {2};
    \node[cell, right=0pt of c1] (c2) {3};
    \node[cell, right=0pt of c2] (c3) {4};
    \node[cell, right=0pt of c3] (c4) {5};
    \node[cell, right=0pt of c4] (c5) {6};

    \draw[diagarrow] ([yshift=-0.25cm]c0.south) -- ([yshift=-0.25cm]c1.south) --
        ([yshift=-0.25cm]c2.south) -- ([yshift=-0.25cm]c3.south) --
        ([yshift=-0.25cm]c4.south) -- ([yshift=-0.25cm]c5.south);
    \node[idxlbl, anchor=west] at ([xshift=6pt,yshift=-0.25cm]c5.east) {列優先：依序移動};

    \draw[diagarrow, black!45] ([yshift=0.35cm]c0.north) to[out=45,in=135] ([yshift=0.35cm]c3.north);
    \draw[diagarrow, black!45] ([yshift=0.45cm]c3.north) to[out=45,in=135] ([yshift=0.45cm]c1.north);
    \draw[diagarrow, black!45] ([yshift=0.55cm]c1.north) to[out=45,in=135] ([yshift=0.55cm]c4.north);
    \node[idxlbl, anchor=west] at ([xshift=6pt,yshift=0.4cm]c5.east) {行優先：來回跳著移動};
\end{tikzpicture}
\diagcaption{對同一段記憶體，列優先走訪（下方箭頭）每一步都移到緊鄰的下
一格；行優先走訪（上方箭頭）卻要在記憶體裡大幅度來回跳動。 前者對
cache 更友善，通常也跑得比較快。}
\label{fig:ch06-row-major-traversal}
\end{center}

\keypoint{二維陣列在記憶體裡是列優先排列的：同一列的元素彼此緊鄰，不同
列之間則相隔一整列的距離。 走訪二維陣列時，讓內層迴圈跟著「列優先」的
方向（也就是讓行的索引 \texttt{j} 在內層變化）走訪，通常比反過來更有效
率。}

存取二維陣列元素時，另一個常見的手誤是把 \texttt{grid[i][j]} 誤寫成
\texttt{grid[i,j]}。 這行仍然可以編譯通過，但意思完全不對：

\begin{code}{c}{常見手誤：把 [i][j] 誤寫成 [i,j]}
grid[i][j] = 0;   // 正確：存取第 i 列第 j 行

grid[i, j] = 0;   // 可以編譯，但等同 grid[j] = 0
                  // 逗號運算子會先算 i、丟棄結果，再算 j
\end{code}

\stylepoint{\texttt{grid[i, j]} 在 C 裡會被逗號運算子解讀成「先算
\texttt{i}、丟棄結果，再算 \texttt{j}」，整句等同 \texttt{grid[j]}。 
這不是語法錯誤，編譯器通常也不會警告，是二維陣列最容易被忽略的手誤，一
定要記得寫成兩個獨立的中括號 \texttt{grid[i][j]}。}

二維陣列的初始化可以用巢狀大括號，每一層對應一個維度：

\begin{code}{c}{二維陣列的初始化}
int grid[2][3] = {
    {1, 2, 3},
    {4, 5, 6},
};
\end{code}

前面幾個範例裡，我們都把二維陣列開成固定的最大尺寸（例如
\texttt{grid[105][105]}），即使實際的 $R$、$C$ 遠小於這個上限。 C99 其
實也允許用執行期才知道的變數當作陣列大小（稱為\emph{可變長度陣列}，
variable-length array，簡稱 VLA），例如直接宣告
\texttt{int grid[R][C];}；但 VLA 不能是全域變數、也不能在宣告時給初始
值，限制不少，所以題庫裡的範例多半還是選擇開一個夠大的固定陣列，用實際
的 $R$、$C$ 決定要用到多少。

\texttt{grid[i][j]} 這種「列、行」兩個索引的寫法，剛好跟數學上矩陣的記
號 $a_{ij}$（第 $i$ 列第 $j$ 行）幾乎一模一樣，所以二維陣列很自然地拿
來表示矩陣。 範例 \thechapter.20 計算一個 $R \times C$ 矩陣的\emph{轉
置}（transpose）：轉置後的矩陣是 $C \times R$，原本第 $i$ 列第 $j$ 行
的值，變成新矩陣的第 $j$ 列第 $i$ 行。

\begin{example}{c}{(matrix-transpose.c) 計算矩陣的轉置}
#include <stdio.h>

int main(void) {
    int R, C;
    scanf("%d %d", &R, &C);

    int grid[105][105];
    for (int i = 0; i < R; i++) {
        for (int j = 0; j < C; j++) {
            scanf("%d", &grid[i][j]);
        }
    }

    // 轉置：輸出的第 j 列第 i 行 = 原本的第 i 列第 j 行
    for (int j = 0; j < C; j++) {
        for (int i = 0; i < R; i++) {
            if (i) {
                printf(" ");
            }
            printf("%d", grid[i][j]);
        }
        printf("\n");
    }
}
\end{example}

\iosample{
2 3 \\
1 2 3 \\
4 5 6 \\
}{
1 4 \\
2 5 \\
3 6 \\
}

這段程式跟本章稍早「逐行加總」的迴圈順序完全一樣。 外層固定
\texttt{j}、內層才走訪 \texttt{i}。 只是這次不是拿來加總，而是直接把
讀到的值印到新的位置。 兩個例子背後是同一件事：只要想清楚「要用什麼順
序走訪二維陣列」，很多看似不同的題目其實只是同一套走訪順序的不同應用。

\keypoint{\texttt{grid[i][j]} 的兩個索引對應數學矩陣的「列、行」記號；
一旦題目描述用到「矩陣」「表格」這類詞彙，通常代表可以直接用二維陣列
搭配列優先／行優先的走訪順序來解題，不需要另外設計新的資料結構。}

\begin{drills}[label=drills:array-2d-memory]
\item 二維陣列 \texttt{grid[R][C]} 在記憶體裡實際上是怎麼排列的？
\item 一個 $R \times C$ 的二維陣列，\texttt{grid[i][j]} 在攤平記憶體裡
的位移量公式是什麼？
\item \texttt{grid[i, j]} 這樣寫會發生什麼事？跟 \texttt{grid[i][j]}
是同一件事嗎？
\item 走訪二維陣列時，為什麼讓內層迴圈跟著列優先的方向走，通常比反過
來更快？
\end{drills}

\section{二維陣列的暖身：逐行加總}
\label{sec:array-2d-column-sum}

前面已經看過二維陣列在記憶體裡的真實排列方式，這裡用一個簡單的例子練習
最基本的二維陣列讀取與逐行加總，為下一節稍微複雜一點的鄰居檢查暖身。

範例 \thechapter.21 讀入一個 $R$ 列 $C$ 行的整數表格，對每一行（column）
算出該行所有數字的整數平均。 這裡刻意選「行」而不是「列」，因為列平均
只需要走訪陣列裡本來就連續的一段，但行平均需要在兩個維度之間跳躍讀取，
更能練習到二維索引。 也剛好呼應上一節「列優先 vs. 行優先」的討論。

\begin{center}
\begin{tikzpicture}
    \node[cell] (g0c0) {1};
    \node[hlcell, right=0pt of g0c0] (g0c1) {2};
    \node[cell, right=0pt of g0c1] (g0c2) {3};
    \node[cell, below=0pt of g0c0] (g1c0) {5};
    \node[hlcell, right=0pt of g1c0] (g1c1) {6};
    \node[cell, right=0pt of g1c1] (g1c2) {7};
    \node[hlcell, below=0.55cm of g1c1] (res) {4};
    \draw[diagarrow] (g1c1.south) -- (res.north);
    \node[idxlbl, anchor=west, align=left] at ([xshift=8pt]$(g1c1.south)!0.5!(res.north)$) {column 1\\加總取平均};
\end{tikzpicture}
\diagcaption{中間這一行（column 1）的兩個值 $2$ 和 $6$ 加總後除以列數
$R=2$，得到平均 $4$。}
\label{fig:ch06-2d-column-sum}
\end{center}

\begin{example}{c}{(column-average.c) 計算表格每一行的整數平均}
#include <stdio.h>

int main(void) {
    int R, C;
    scanf("%d %d", &R, &C);

    int grid[105][105];
    for (int i = 0; i < R; i++) {
        for (int j = 0; j < C; j++) {
            scanf("%d", &grid[i][j]);
        }
    }

    // 對每一行（column）累加各列的值，再取整數平均
    for (int j = 0; j < C; j++) {
        int sum = 0;
        for (int i = 0; i < R; i++) {
            sum += grid[i][j];
        }
        if (j) {
            printf(" ");
        }
        printf("%d", sum / R);
    }
    printf("\n");
}
\end{example}

\iosample{
2 3 \\
1 2 3 \\
5 6 7 \\
}{
3 4 5 \\
}

注意這裡兩層迴圈的順序：外層固定「第幾行」\texttt{j}，內層才走訪「第幾
列」\texttt{i}，這跟前面填滿陣列、逐列印出的順序（外層 \texttt{i}、內層
\texttt{j}）剛好相反。 依照上一節的討論，這個順序在記憶體裡其實是屬於
「行優先」的跳躍走訪，只是資料量小的時候感覺不出差異。

\stylepoint{兩層迴圈要處理二維陣列的哪個維度在外、哪個在內，取決於我們
想要的走訪順序，而不是陣列宣告時兩個維度的先後順序；寫之前先想清楚「
外層代表什麼、內層代表什麼」，比邊寫邊試更不容易搞混。}

\begin{drills}[label=drills:array-2d-column-sum]
\item 範例 column-average.c 為什麼外層固定 \texttt{j}（行）、內層走訪
\texttt{i}（列），而不是反過來？
\item 依上一節的討論，這種走訪順序是列優先還是行優先？
\end{drills}

\section{二維陣列的鄰居檢查}
\label{sec:array-2d-neighbor}

二維陣列除了逐列或逐行整體處理之外，很多題目要對\emph{每一格}分別檢查
它上下左右的鄰居。 例如判斷這一格是不是區域最大值、是不是被特定的圖案
包圍。 這類題目的共通困難點在於邊界：格子在網格邊緣時，某些方向的鄰居
根本不存在，程式必須先判斷鄰居的座標有沒有落在網格範圍內，才能安全地去
讀它的值。

範例 \thechapter.22 讀入一個 $R \times C$ 的網格，對每一格檢查它是不是
嚴格大於所有\emph{存在}的上下左右鄰居（邊緣或角落的格子鄰居比較少，只
需要跟存在的比），統計這樣的「區域最大值」總共有幾格。

\begin{center}
\begin{tikzpicture}
    \node[cell] (r0c0) {1};
    \node[hlcell, right=0pt of r0c0] (up) {2};
    \node[cell, right=0pt of up] (r0c2) {1};
    \node[hlcell, below=0pt of r0c0] (left) {2};
    \node[hlcell2, right=0pt of left] (mid) {9};
    \node[hlcell, right=0pt of mid] (right) {2};
    \node[cell, below=0pt of left] (r2c0) {1};
    \node[hlcell, right=0pt of r2c0] (down) {2};
    \node[cell, right=0pt of down] (r2c2) {1};
    \draw[diagarrow] (mid.north) -- (up.south);
    \draw[diagarrow] (mid.south) -- (down.north);
    \draw[diagarrow] (mid.west) -- (left.east);
    \draw[diagarrow] (mid.east) -- (right.west);
\end{tikzpicture}
\diagcaption{中央的 $9$（粉色）對照上下左右四個鄰居（橘色），全部嚴格
比它小，所以中央是區域最大值；角落與邊緣格子鄰居較少，逐一判斷是否存
在再比較。}
\label{fig:ch06-2d-neighbor-check}
\end{center}

\begin{example}{c}{(local-max-2d.c) 統計網格中區域最大值的格數}
#include <stdio.h>

int main(void) {
    int R, C;
    scanf("%d %d", &R, &C);

    int grid[105][105];
    for (int i = 0; i < R; i++) {
        for (int j = 0; j < C; j++) {
            scanf("%d", &grid[i][j]);
        }
    }

    // 上、下、左、右四個方向的座標偏移
    int dr[4] = {-1, 1, 0, 0};
    int dc[4] = {0, 0, -1, 1};

    int count = 0;
    for (int i = 0; i < R; i++) {
        for (int j = 0; j < C; j++) {
            int isLocalMax = 1;
            for (int d = 0; d < 4; d++) {
                int ni = i + dr[d], nj = j + dc[d];
                if (ni >= 0 && ni < R && nj >= 0 && nj < C) {
                    if (grid[ni][nj] >= grid[i][j]) {
                        isLocalMax = 0;
                    }
                }
            }
            if (isLocalMax) {
                count++;
            }
        }
    }

    printf("%d\n", count);
}
\end{example}

\iosample{
3 3 \\
1 2 1 \\
2 9 2 \\
1 2 1 \\
}{
1 \\
}

範例輸入是一個十字形排列、中央特別高的網格。 中央的 $9$ 上下左右都是
$2$，四個方向都嚴格比它小，是區域最大值；四個角落（例如左上角的 $1$）
只有右邊和下面兩個鄰居存在，都是 $2 \ge 1$，所以角落都不是區域最大值；
邊緣中點（例如上方中間的 $2$）有三個鄰居存在，其中下方的 $9 \ge 2$，也
不是區域最大值。 整個網格只有中央這一格符合，答案是 $1$。

\keypoint{用 \texttt{dr}、\texttt{dc} 兩個小陣列存四個方向的座標偏移量，
再用一個迴圈依序套用，比為上下左右各寫一次幾乎一樣的判斷式更精簡，之後
遇到八個方向（含斜角）的題目時，只需要把這兩個陣列换成八個方向就好，邏
輯完全不用重寫。}

\begin{drills}[label=drills:array-2d-neighbor]
\item \texttt{dr}、\texttt{dc} 這兩個小陣列的用途是什麼？
\item 角落格子跟中央格子相比，鄰居的數量有什麼不同？程式怎麼處理這個
差異？
\item 範例輸入的十字形網格裡，區域最大值總共有幾格？
\end{drills}

\section{二維陣列與座標式問題：幻方}
\label{sec:array-2d-magic-square}

有些問題需要一邊走訪二維陣列、一邊依規則\emph{建構}它的內容，而不只是
讀取既有的資料。 這一節用一個經典的例子，把陣列的走訪、環繞索引、以及
先前學過的化簡策略全部串起來：幻方。 幻方是一個 $N \times N$ 的矩陣，
依下列步驟填入正整數 $1, 2, \cdots, N^2$：首先將 $1$ 填在起始位置。 之
後每次向右上對角線方向移一格並填入下一個數字；如果那一格已經填過，就改
為向正下方移一格，再繼續往右上填。 「向右上」或「向下」的移動超出方格
範圍時，會環繞到最後一列、第一行或第一列。 這正是陣列索引很適合搭配
\texttt{\%} 運算子做環繞（wrap-around）的地方，因為
\texttt{(r + 1) \% N} 會在 $r$ 超出 $N - 1$ 時自動繞回 $0$。

\begin{center}
\begin{tikzpicture}
    \foreach \r in {0,1,2} {
        \foreach \c in {0,1,2} {
            \ifnum\c=0
                \ifnum\r=0
                    \node[cell] (g-\r-\c) {};
                \else
                    \pgfmathtruncatemacro{\rprev}{\r-1}
                    \node[cell, below=0pt of g-\rprev-0] (g-\r-\c) {};
                \fi
            \else
                \pgfmathtruncatemacro{\cprev}{\c-1}
                \node[cell, right=0pt of g-\r-\cprev] (g-\r-\c) {};
            \fi
        }
    }
    \node[hlcell] at (g-2-0.center) {1};
    \node[hlcell] at (g-1-1.center) {2};
    \node[hlcell] at (g-0-2.center) {3};
    \draw[diagarrow] (g-2-0.north east) -- (g-1-1.south west);
    \draw[diagarrow] (g-1-1.north east) -- (g-0-2.south west);
    \node[idxlbl, yshift=0.1cm] at ($(g-0-0.north)!0.5!(g-0-1.north)$) {向右上填下一格};
    \coordinate (wrap) at ($(g-1-2.east)+(1.25,0)$);
    \draw[diagarrow] (g-0-2.east) to[out=15,in=90] (wrap) to[out=270,in=345] (g-2-0.south east);
    \node[idxlbl, align=left, right=8pt of wrap] {超出範圍\\就繞回};
\end{tikzpicture}
\diagcaption{正常情況下每次向右上移一格；移動超出方格範圍時，環繞到對
面繼續（這裡示範從右上角繞回左下角）。}
\label{fig:ch06-magic-square-move}
\end{center}

\begin{center}
\begin{tikzpicture}
    \node[hlcell2] (cur) {6};
    \node[hlcell, anchor=south west] (blocked) at (cur.north east) {4};
    \node[hlcell, below=0cm of cur] (down) {7};
    \draw[diagarrow, black!35] (cur.north east) -- (blocked.south west);
    \node[font=\large] at ($(cur.north east)!0.5!(blocked.south west)$) {$\times$};
    \node[idxlbl, align=left, right=8pt of blocked] {右上已填過};
    \draw[diagarrow] (cur.south) -- (down.north);
    \node[idxlbl, align=left, right=8pt of down] {改往下填};
\end{tikzpicture}
\diagcaption{向右上那一格已經填過數字時（這裡是已經填了 $4$ 的格
子），改成往正下方移一格繼續填。}
\label{fig:ch06-magic-square-collision}
\end{center}

把兩條規則反覆套用在一個 $3 \times 3$、從左下角開始（$k=1$）的小幻方
上，就會得到下圖的結果：

\begin{center}
\begin{tikzpicture}
    \node[cell] (g0c0) {5};
    \node[cell, right=0pt of g0c0] (g0c1) {7};
    \node[cell, right=0pt of g0c1] (g0c2) {3};
    \node[cell, below=0pt of g0c0] (g1c0) {9};
    \node[cell, right=0pt of g1c0] (g1c1) {2};
    \node[cell, right=0pt of g1c1] (g1c2) {4};
    \node[cell, below=0pt of g1c0] (g2c0) {1};
    \node[cell, right=0pt of g2c0] (g2c1) {6};
    \node[cell, right=0pt of g2c1] (g2c2) {8};
\end{tikzpicture}
\diagcaption{$3 \times 3$、從左下角開始填 $1$ 的完整結果；沿右上對角線
走兩步（$1 \to 2 \to 3$）就碰到邊界繞回，之後遇到已填格子時改往下
（$3 \to 4$、$6 \to 7$）。 正是前面兩張圖示範的規則。}
\label{fig:ch06-magic-square-result}
\end{center}

範例 \thechapter.23 給定矩陣大小 $N$、以及整數 $k$（$1 \le k \le N$）所在的
起始位置 $(x, y)$，建構整個幻方並印出來。 由於已知的是「$k$ 所在的位置」
而不是「$1$ 所在的位置」，我們先沿著對角線反向走 $k - 1$ 步，推算出
$1$ 應該在的起始位置，再從那裡開始，用同一套規則把整個矩陣填滿。

\begin{example}{c}{(magic-square.c) 建構幻方}
#include <stdio.h>

int grid[1000][1000];

int main(void) {
    int N, k, x, y;
    scanf("%d %d %d %d", &N, &k, &x, &y);

    // 從「k 的位置」沿對角線反推「1 應該在的位置」
    int r = x, c = y;
    for (int i = 0; i < k - 1; i++) {
        r = (r + 1) % N;
        c = (c - 1 + N) % N;
    }

    // 從 1 開始依規則把整個矩陣填滿
    for (int v = 1; v <= N * N; v++) {
        grid[r][c] = v;

        int nr = (r - 1 + N) % N, nc = (c + 1) % N;
        if (grid[nr][nc] != 0) {
            nr = (r + 1) % N;
            nc = c;
        }
        r = nr;
        c = nc;
    }

    // 輸出整個矩陣
    for (int i = 0; i < N; i++) {
        for (int j = 0; j < N; j++) {
            if (j) {
                printf(" ");
            }
            printf("%d", grid[i][j]);
        }
        printf("\n");
    }
}
\end{example}

\iosample{
5 1 1 2 \\
}{
11 18 25 2 9 \\
17 24 1 8 15 \\
23 5 7 14 16 \\
4 6 13 20 22 \\
10 12 19 21 3 \\
}

程式先用一個獨立的迴圈把起始位置反推回去，得到「$1$ 應該在的位置」之後，
才進入真正填數字的主迴圈；這個主迴圈每一步都會檢查「向右上那一格是不
是已經填過」，填過的話才改成向下移動，這正好對應到題目描述的規則。

\keypoint{用 \texttt{(r - 1 + N) \% N} 而不是 \texttt{(r - 1) \% N} 處
理環繞，是因為 C 裡負數取餘數可能是負的（例如 \texttt{-1 \% 5} 是
\texttt{-1}，不是 \texttt{4}）。}

範例 \thechapter.23 示範的是一個常見的解題策略：與其直接處理「已知 $k$ 的
位置」這個比較複雜的情形，不如先想清楚「已知 $1$ 的位置」這個最簡單的情
形要怎麼解，再把複雜的情形\emph{化簡}成簡單的情形。 這裡能這樣做的
關鍵觀察是：只要 $k \le N$，起始位置到 $1$ 的位置之間就只會發生「向右上」
的移動、不會發生「向下」的移動，所以反向走 $k - 1$ 步一定能精確地推回
$1$ 的位置。 如果 $k$ 可以大於 $N$，這個反推的技巧就不一定成立了，因為
「向右上」的路徑上有可能已經發生過填滿而改向下的情形，反推時就不能只
單純地走對角線。 這也是為什麼許多這類題目會把 $k \le N$ 設計成一個難度
較低的子題。

\keypoint{先確認一個問題的簡化情形要怎麼解，再想辦法把一般情形化簡成簡化
情形，往往比一開始就想著同時處理所有情形更容易寫對。}

\begin{drills}[label=drills:array-2d-magic-square]
\item 為什麼要用 \texttt{(r + 1) \% N} 而不是直接寫 \texttt{r + 1}？
\item 為什麼要寫成 \texttt{(r - 1 + N) \% N} 而不是 \texttt{(r - 1) \% N}？
\item 幻方規則裡，「向右上填」跟「改往下填」是根據什麼條件切換的？
\end{drills}

\addcontentsline{toc}{section}{練習解答}
\section*{練習解答}

各節練習的答案收在這裡。 短題只給能一句話核對的結果。

\paragraph{第~\ref{sec:array-declare-read}~節}
\begin{enumerate}
\item 合法索引是 $0$ 到 $1004$，總共 $1005$ 個位置；\texttt{a[1005]}
本身不合法，是超出範圍的存取。
\item 因為索引代表的是離陣列開頭的\emph{位移量}，不是「第幾個」的排
名；第一個元素跟開頭的位移量是 $0$，所以索引從 $0$ 開始。
\item 沒有。 反向印出只是改變走訪的方向，陣列裡的資料完全沒有被搬動
或修改。
\item 兩個迴圈：第一個負責把 $5$ 個數字讀進陣列，第二個負責依序把陣
列內容印出來。
\end{enumerate}

\paragraph{第~\ref{sec:array-init-limits}~節}
\begin{enumerate}
\item 全部是 $0$：\texttt{\{0, 0, 0, 0, 0\}}。
\item 三者都是 $0$（初始化值不足時，剩下的元素自動補 $0$）。
\item 陣列不能整個賦值；必須用迴圈逐一複製每個元素（\texttt{b[i] =
a[i];}），或使用 \texttt{memcpy}。
\item 陣列的元素個數（整個陣列的位元組數除以單一元素的位元組數）。
\end{enumerate}

\paragraph{第~\ref{sec:array-out-of-bounds}~節}
\begin{enumerate}
\item 不會。 C 語言完全不在執行期檢查陣列索引，越界存取是未定義行為。
\item 因為未定義行為的表現因編譯器、系統、甚至該次執行當下的記憶體
狀態而異；這次沒出事不代表下次、換一台機器也不會出事。
\item 因為 C 沒有規定「先算左邊被讀取的 \texttt{i} 還是先算右邊
\texttt{i++} 的副作用」；把遞增獨立成一行，或直接用 \texttt{for} 管
理索引，就能避免。
\end{enumerate}

\paragraph{第~\ref{sec:array-stream-no-array}~節}
\begin{enumerate}
\item 這題到底需不需要事後回頭看某一筆舊資料。
\item 因為總和與最大值都只需要邊讀邊更新，不需要回頭比較任何一筆舊
資料，所以兩個變數就夠了。
\item 當 $n$ 大到陣列可能裝不下（MLE）時，邊讀邊算往往是唯一可行的
做法。
\end{enumerate}

\paragraph{第~\ref{sec:array-second-largest}~節}
\begin{enumerate}
\item 因為要先確定最大值是誰，才能在第二遍裡正確排除它，找出比它嚴
格小的最大值；只掃一次沒辦法同時做到這兩件事。
\item 是 $3$。 三個 $5$ 都等於最大值本身，必須排除，剩下最大的是 $3$。
\item 會誤把「第二次出現的最大值」（例如另一個 $5$）當成答案，而不是
真正嚴格小於最大值的那個值。
\end{enumerate}

\paragraph{第~\ref{sec:array-reverse-swap}~節}
\begin{enumerate}
\item 因為 \texttt{a[left] = a[right];} 執行後，\texttt{a[left]} 原本
的值已經被蓋掉，若沒先存起來，第二行就換不回正確的值。
\item \texttt{left} 從 $0$ 開始，\texttt{right} 從 $n - 1$ 開始，往中
間逼近；當 \texttt{left >= right} 時迴圈結束。
\item print-reverse.c 只是改變印出的順序，陣列內容不變；
reverse-in-place.c 是真的把陣列元素互換，陣列本身被改寫了。
\end{enumerate}

\paragraph{第~\ref{sec:array-partition}~節}
\begin{enumerate}
\item 因為兩個陣列各自累積的資料筆數不同步；共用一個索引會讓其中一
個陣列出現沒填到值的空洞。
\item \texttt{evens = \{2, 4, 6\}}、\texttt{oddsArr = \{1, 3, 5\}}，
各自保留原本在輸入裡出現的先後順序。
\end{enumerate}

\paragraph{第~\ref{sec:array-lookup}~節}
\begin{enumerate}
\item \texttt{x} 代表的是數值本身；前面幾節的索引多半代表「第幾筆輸
入」，這裡索引直接就是要查詢的那個值。
\item 數值的可能範圍 $v$；陣列要開到能涵蓋所有可能出現的數值。
\item 依序印出 $3, 5, 7$（各自第二次出現時被印出）。
\end{enumerate}

\paragraph{第~\ref{sec:array-bucket}~節}
\begin{enumerate}
\item 查表存的是「有沒有看過」的 $0$／$1$；分類桶存的是累加起來的統
計值。
\item 陣列初始內容不確定，直接拿來累加會得到無法預期的錯誤結果。
\item \texttt{bucketSum[0] = 3}、\texttt{bucketSum[1] = 5}、
\texttt{bucketSum[2] = 7}。
\end{enumerate}

\paragraph{第~\ref{sec:array-sliding-window}~節}
\begin{enumerate}
\item 桶依\emph{值}分類；區塊依\emph{位置}切段，用索引的範圍決定資料
屬於哪一組。
\item 分成 $3$ 組；最後一組只剩 $2$ 個元素。
\item 因為題目刻意設計成「最後一組要不要特別處理」的邊界練習：不足額
的區塊改統計最小值，藉此練習邊界情形跟完整區塊分開處理。
\end{enumerate}

\paragraph{第~\ref{sec:array-2d-memory}~節}
\begin{enumerate}
\item 列優先排列：先排完第 $0$ 列，再排第 $1$ 列，依此類推，攤平成一
長條連續記憶體。
\item $i \times C + j$（$C$ 是每列的行數）。
\item \texttt{grid[i, j]} 會被逗號運算子解讀成「先算 \texttt{i}、丟棄
結果，再算 \texttt{j}」，等同 \texttt{grid[j]}，跟 \texttt{grid[i][j]}
完全不是同一件事，但可以編譯通過。
\item 因為列優先走訪時，內層迴圈每一步都移到記憶體裡緊鄰的下一格，對
CPU 快取更友善；反過來則要在記憶體裡大幅度跳動。
\end{enumerate}

\paragraph{第~\ref{sec:array-2d-column-sum}~節}
\begin{enumerate}
\item 因為要對每一\emph{行}分別加總，所以外層固定行 \texttt{j}，內層
才走訪列 \texttt{i} 把同一行的各列值加起來。
\item 是行優先：內層迴圈的 \texttt{i} 每次都要在記憶體裡跳一整列的距
離，不是緊鄰的下一格。
\end{enumerate}

\paragraph{第~\ref{sec:array-2d-neighbor}~節}
\begin{enumerate}
\item 存四個方向（上下左右）相對於目前格子的座標偏移量，搭配一個迴
圈依序套用，取代四段幾乎一樣的判斷式。
\item 角落格子只有兩個鄰居存在，邊緣格子有三個，中央格子有四個；程式
用座標範圍判斷（\texttt{ni}、\texttt{nj} 是否落在 $0$ 到 $R-1$／
$C-1$ 之間）決定鄰居是否存在，只跟存在的鄰居比較。
\item $1$ 格（十字中央的 $9$）。
\end{enumerate}

\paragraph{第~\ref{sec:array-2d-magic-square}~節}
\begin{enumerate}
\item 因為 \texttt{r + 1} 可能等於 $N$、超出合法的列索引範圍；
\texttt{\% N} 讓它在超出時自動環繞回 $0$。
\item 因為 C 語言裡負數取餘數可能是負的（例如 \texttt{-1 \% 5} 是
\texttt{-1}，不是 \texttt{4}）；先加上 $N$ 再取餘數，才能保證結果落
在 $0$ 到 $N-1$ 之間。
\item 根據「向右上那一格是否已經填過數字」切換：還沒填過就繼續往右
上，已經填過就改往正下方移動。
\end{enumerate}

\chapter{浮點數}
\label{ch:floatingpointnumbers}

到目前為止我們處理的資料大多是整數。 但很多真實世界的數量。 溫度、距
離、比例。 本質上是連續的，需要小數點。 C 語言用 \texttt{float} 與
\texttt{double} 儲存這類「浮點數」（floating-point number）。

跟整數不同，浮點數在記憶體裡\textbf{不是}十進位小數的精確表示，而是用固
定數量的位元，以二進位、類似科學記號的方式去\emph{逼近}一個值。 我們熟
悉的十進位小數（例如 $0.1$）換成二進位之後，很多時候變成無限循環小數，
電腦只能在有限的位元裡把它捨入成一個「很接近但不完全相等」的值。 這件
事聽起來抽象，但它會直接影響我們怎麼輸出浮點數、怎麼判斷兩個浮點數是不
是「相等」，貫穿這一整章；甚至會影響我們該不該一開始就用浮點數來解一個
問題。

這一章分五段推進：先把宣告、格式化、整數除法陷阱與型別轉型講完，再一口
氣建立記憶體近似、輸出假象、比較陷阱與 epsilon 的心智模型，接著教能
用整數就不用浮點的實戰技巧，然後才是 \texttt{math.h} 的幾何應用，最後用
累加穩定度與迭代方法收尾。 內容對齊本單元課綱（\texttt{float}／\texttt{double}
型別與 \texttt{scanf}/\texttt{printf} 格式化、精度與位元表示、\texttt{math.h}
數學函式），難度曲線對應題庫從溫度轉換到多邊形周長的整條路徑。

\section{宣告、讀寫與格式化輸出}
\label{sec:float-declare-io}

浮點數的宣告跟整數很像，差別在型別關鍵字：\texttt{float} 是單精度（通
常占 $4$ 個位元組、約 $7$ 位有效數字），\texttt{double} 是雙精度（通
常占 $8$ 個位元組、約 $15$ 位有效數字）。 沒有特別理由的話，一般都用
\texttt{double}，因為多數 C 函式庫（包括接下來會用到的數學函式）預設就
是用 \texttt{double} 在運算，混用 \texttt{float} 反而容易在型別轉換
的地方出錯。

\begin{code}{c}{宣告並讀入一個 double，輸出到小數點後 N 位}
double x;
scanf("%lf", &x);
printf("%.1f\n", x);
\end{code}

\texttt{printf} 裡的 \texttt{\%.1f} 表示四捨五入到小數點後 $1$ 位。 浮
點數輸出幾乎都要指定小數位數，不指定的話預設印 $6$ 位，容易印出一堆沒
意義的尾數。 這些尾數往往正是浮點數用二進位逼近十進位小數時產生的誤
差，指定小數位數等於是把這些誤差四捨五入掉，讓輸出看起來乾淨。

\keypoint{\texttt{scanf} 讀 \texttt{double} 要用 \texttt{\%lf}，但
\texttt{printf} 印 \texttt{double} 只需要 \texttt{\%f}，這是 C 語言留下
來的不對稱寫法。}

範例 \thechapter.2 把攝氏溫度轉換成華氏溫度，公式是 $F = C \times
\frac{9}{5} + 32$。 這是本章最簡單的範例，重點不在演算法，而在浮點數運
算最基本、也最容易忽略的一個地雷：整數除法。

\begin{example}{c}{(celsius-to-fahrenheit.c) 攝氏轉華氏}
#include <stdio.h>

int main(void) {
    double c;
    scanf("%lf", &c);

    // F = C * 9/5 + 32
    double f = c * 9.0 / 5.0 + 32.0;

    printf("%.1f\n", f);
}
\end{example}

\iosample{
100 \\
}{
212.0 \\
}

\stylepoint{公式裡的 \texttt{9.0}、\texttt{5.0} 特意寫成浮點數常數；如
果寫成 \texttt{c * 9 / 5}，\texttt{9 / 5} 會先做整數除法變成
\texttt{1}，答案就錯了。}

這個陷阱值得再多說一句：C 語言的運算子沒辦法「預知」我們接下來想拿這個
值去跟浮點數相乘，它只看眼前這一個運算子兩側的型別。 \texttt{9} 和
\texttt{5} 都是整數常數，\texttt{9 / 5} 這個運算單獨看就是整數除法，結
果被無條件捨去成 \texttt{1}，接下來不管乘上什麼都救不回來。 只要算式裡
有任何一個運算元是浮點數，那一步運算才會被拉去用浮點數計算，所以在常
數後面補上 \texttt{.0} 是最簡單、也最保險的做法。 輸入 $100$ 攝氏，算
出 $212.0$ 華氏，跟範例輸出一致。

\begin{drills}[label=drills:float-declare-io]
\item \texttt{scanf} 讀 \texttt{double} 該用哪個格式碼？ \texttt{printf}
印 \texttt{double} 呢？
\item 為什麼公式裡要把 \texttt{9} 和 \texttt{5} 寫成 \texttt{9.0} 和
\texttt{5.0}？ 若寫成 \texttt{c * 9 / 5} 會發生什麼事？
\item \texttt{printf("\%f", x)} 不指定小數位數時，預設會印幾位？
\end{drills}

\section{科學記號、e 格式，以及整數除法再看一次}
\label{sec:float-sci-notation}

很大或很小的浮點數，用小數點寫會一長串零。 C 允許科學記號常數（例如
\texttt{1.2e-4} 表示 $1.2 \times 10^{-4}$），\texttt{printf} 的
\texttt{\%e} 則把值印成「尾數 $\times$ 十的次方」。 範例
\thechapter.3 把 $12345.678$ 先用三位科學記號印出，再用兩位
\texttt{\%f} 印出。 四捨五入後分別是 \texttt{1.235e+04} 與
\texttt{12345.68}。

\begin{example}{c}{(sci-format.c) 科學記號與固定小數位}
#include <stdio.h>

int main(void) {
    double x = 12345.678;
    printf("%.3e\n", x);
    printf("%.2f\n", x);
    return 0;
}
\end{example}

\iosample{}{
1.235e+04 \\
12345.68 \\
}

讀 \texttt{float} 用 \texttt{\%f}（或 \texttt{\%e}），讀
\texttt{double} 用 \texttt{\%lf}；印的時候 \texttt{float} 會被提升成
\texttt{double}，所以 \texttt{printf} 一律用 \texttt{\%f} 或
\texttt{\%e} 即可。 \texttt{long double} 更寬，讀寫常用 \texttt{\%Lf}，
本書範例仍以 \texttt{double} 為主。

整數除法的陷阱在運算那一章講過，這裡再用浮點輸出對一次。 範例
\thechapter.4 印出三種寫法：\texttt{9 / 5} 仍是整數 $1$；
\texttt{9.0 / 5} 與 \texttt{(double)9 / 5} 才是 $1.8$。

\begin{example}{c}{(mix-div.c) 整數除法與浮點除法並排}
#include <stdio.h>

int main(void) {
    printf("%d\n", 9 / 5);
    printf("%.1f\n", 9.0 / 5);
    printf("%.1f\n", (double)9 / 5);
    return 0;
}
\end{example}

\iosample{}{
1 \\
1.8 \\
1.8 \\
}

第一行必須用 \texttt{\%d}：\texttt{9 / 5} 的型別是 \texttt{int}，若
硬用 \texttt{\%.1f} 去印，\texttt{printf} 會把整數位元樣式誤讀成浮
點數，印出來的東西沒有意義。 第二、三行才是浮點除法，所以用
\texttt{\%.1f}。 攝氏轉華氏那一節把常數寫成 \texttt{9.0 / 5.0}，就
是為了讓除法從一開始就走這條路。

\begin{drills}[label=drills:float-sci-notation]
\item \texttt{1.235e+04} 這個 \texttt{\%e} 輸出，換算成一般小數是多少？
\item 為什麼 \texttt{printf("\%d\textbackslash n", 9 / 5)} 不能改成
\texttt{printf("\%.1f\textbackslash n", 9 / 5)}？
\item \texttt{9.0 / 5} 和 \texttt{(double)9 / 5} 這兩種寫法，哪一個運
算元被轉成浮點數的方式不一樣？
\end{drills}

\section{強制轉型，以及由左向右的減法}
\label{sec:float-cast-assoc}

除了在常數後面加 \texttt{.0}，也可以用\emph{強制轉型}（cast）把運算
元暫時看成另一種型別。 \texttt{(double)(a + b) / 2} 先把整數和轉成
\texttt{double}，再除以 $2$，商才會帶小數。 括號必須包住要轉的那一
段：若寫成 \texttt{(double) a + b / 2}，除法仍可能先在整數裡做完。

\begin{example}{c}{(half-cast.c) 用轉型得到帶小數的平均}
#include <stdio.h>

int main(void) {
    int a, b;
    scanf("%d %d", &a, &b);
    printf("%.1f\n", (double)(a + b) / 2);
    return 0;
}
\end{example}

\iosample{
1 2 \\
}{
1.5 \\
}

最後用減法把結合方向講死。 範例用對 \texttt{10 3 2}
分別印出 \texttt{a - b - c} 與 \texttt{a - (b - c)}。 前者是
\texttt{(10 - 3) - 2 = 5}；後者是 \texttt{10 - (3 - 2) = 9}。 差 $4$
不是溢位，只是括號放錯位置。

\begin{example}{c}{(assoc-minus.c) 減法由左結合，括號會改答案}
#include <stdio.h>

int main(void) {
    int a, b, c;
    scanf("%d %d %d", &a, &b, &c);
    printf("%d\n", a - b - c);
    printf("%d\n", a - (b - c));
    return 0;
}
\end{example}

\iosample{
10 3 2 \\
}{
5 \\
9 \\
}

整數加到超出 \texttt{int} 能表示的範圍時，行為在語言裡沒有保證。 我
們在這台機器上把當時的最大值再加 $1$，觀察到數字從
\texttt{2147483647} 繞成 \texttt{-2147483648}。 這是這次編譯器與硬
體的現象，不是「加一之後一定變負數」的定理。 題目數字可能變大時，
應改用 \texttt{long long}，不要賭溢位剛好繞成你要的答案。

\keypoint{同樣優先順序的加減由左向右算。 需要小數時，用 \texttt{.0}
或 \texttt{(double)} 把除法送進浮點；需要更大範圍時，提早換
\texttt{long long}。}

\begin{drills}[label=drills:float-cast-assoc]
\item 為什麼 \texttt{(double) a + b / 2} 不能保證得到帶小數的平均？
正確寫法該把括號放在哪裡？
\item 輸入 \texttt{10 3 2}，\texttt{a - b - c} 跟 \texttt{a - (b - c)}
分別是多少？ 差在哪裡？
\item 減法的結合方向是由左向右還是由右向左？
\end{drills}

\section{記憶體中的浮點數：有限位元下的近似}
\label{sec:float-bit-repr}

整數在記憶體裡通常就是它的二進位值本身：\texttt{int} 占 $4$ 個位元組，
能表示的範圍固定、每個可表示的整數都精確無誤。 浮點數則不同。 它占用的
位元數也是固定的（\texttt{float} 通常 $4$ 個位元組、\texttt{double}
通常 $8$ 個），但這些位元\textbf{不是}直接存十進位小數，而是存一組「符
號 + 指數 + 尾數」，用類似科學記號 $m \times 2^e$ 的方式去逼近一個
值。 這套格式在幾乎所有現代處理器上都遵循 IEEE 754 標準。

\begin{center}
\begin{tikzpicture}
    \node[idxlbl] (lbl) {\texttt{float}：};
    \node[cell, minimum width=0.5cm, right=0.15cm of lbl] (s) {S};
    \node[hlcell, minimum width=2.7cm, right=0pt of s] (e) {指數 $E$（8 位）};
    \node[hlcell2, minimum width=5.4cm, right=0pt of e] (m) {尾數 $M$（23 位）};
    \node[idxlbl, below=0.15cm of s] {1 位};
    \node[idxlbl, below=0.15cm of e] {8 位};
    \node[idxlbl, below=0.15cm of m] {23 位};
\end{tikzpicture}
\diagcaption{\texttt{float} 在記憶體裡占 $32$ 位（$4$ 個位元組），分成
符號、指數、尾數三段；\texttt{double} 占 $64$ 位，指數 $11$ 位、尾數
$52$ 位，精度更高。}
\label{fig:ch07-float-bit-layout}
\end{center}

這代表什麼？ 代表電腦能「精確表示」的數，其實只是所有實數裡的一個有限
子集。 很多十進位看起來很簡單的小數。 $0.1$、$0.2$、$0.3$。 在二進位
裡是無限循環小數，跟 $\frac{1}{3} = 0.333\ldots$ 在十進位裡循環一樣；
電腦只能存前 $23$（或 $52$）位尾數，其餘全部捨入。 所以 \texttt{double
x = 0.1;} 這行寫起來像是把 $0.1$ 存進 \texttt{x}，但 \texttt{x} 裡實
際上存的是「最接近 $0.1$ 的那個可表示二進位浮點數」，兩者不完全相等。

範例 \thechapter.7 把這件事具體化：先印出 \texttt{float} 和
\texttt{double} 各占幾個位元組，再把 \texttt{0.1} 用完整精度印出來，
最後用 \texttt{union} 把 \texttt{0.1f} 的 $32$ 位元原封不動當成整數
印成十六進位。 這是最後一行的用途，不是要你背十六進位值，而是讓你看
到「同一塊記憶體，解讀成浮點數跟解讀成整數，看到的是完全不同的東
西」。

\begin{example}{c}{(float-storage-demo.c) 觀察浮點數的儲存方式}
#include <stdio.h>
#include <stdint.h>

int main(void) {
    printf("sizeof(float) = %zu, sizeof(double) = %zu\n",
           sizeof(float), sizeof(double));

    // 0.1 在記憶體裡不是精確的 0.1
    double x = 0.1;
    printf("%.17f\n", x);

    // 把 float 的 32 位元原樣當成整數印出（十六進位）
    union {
        float f;
        uint32_t u;
    } bits;
    bits.f = 0.1f;
    printf("%08x\n", bits.u);
}
\end{example}

\iosample{}{
sizeof(float) = 4, sizeof(double) = 8 \\
0.10000000000000001 \\
3dcccccd \\
}

第一行確認 \texttt{float} 占 $4$ 個位元組、\texttt{double} 占 $8$ 個
位元組。 第二行用 \texttt{\%.17f} 印出 \texttt{0.1} 的完整精度。 注意
最後一位是 \texttt{1} 而不是 \texttt{0}，這就是捨入後的近似值。 第三
行 \texttt{3dcccccd} 是 \texttt{0.1f} 在記憶體裡的 $32$ 位元樣子；同
一個 \texttt{union} 裡，改寫 \texttt{bits.f} 再讀 \texttt{bits.u}，就
是在看「這塊記憶體的 raw bits」，之後指標那一章會更常碰到這種「同一塊
記憶體、不同型別解讀」的觀念。

指數那一段不是直接存數學上的 $e$，而是存一個平移過的整數，讓負指數
也能用無號數表示。 以 \texttt{float} 來說，8 位指數常見的做法是把實
際指數加上 $127$ 再寫進去。 這個平移量叫做偏移（bias）。 我們現在
不必手算 bias，只要建立這個直覺：同樣是 $32$ 位，浮點數把多數位元拿
去記「小數有多長」，所以能表示的範圍比 \texttt{int} 大很多，但相鄰
兩個可表示值的間距不是固定的。 靠近 $0$ 的地方比較密，數字越大，間
距越寬。 這也是後面「相對誤差」會出現的硬體原因：大數字本來就沒辦
法分辨「差 $0.5$」這種在小數點附近看起來很大的差距。

\keypoint{浮點數變數存的是「在有限位元下最接近的那個可表示值」，不是
數學上精確的十進位小數；這是硬體層面的限制，不是某個編譯器的 bug。}

\begin{drills}[label=drills:float-bit-repr]
\item \texttt{float} 跟 \texttt{double} 各占幾個位元組？ 各自的指數欄
位有幾位？
\item \texttt{double x = 0.1;} 之後，\texttt{x} 裡存的真的是數學上精確
的 $0.1$ 嗎？
\item 浮點數格式裡，靠近 $0$ 的地方跟數字很大的地方，哪一邊相鄰可表示
值的間距比較密？
\end{drills}

\section{輸出位數會把誤差藏起來}
\label{sec:float-hidden-error}

指定小數位數不只是「讓輸出好看」，它實際上是在做四捨五入。 同一個
\texttt{0.1}，用 \texttt{\%.1f} 印成 \texttt{0.1}，用 \texttt{\%.2f}
印成 \texttt{0.10}，看起來都像精確的十分之一；只有把精度拉到
\texttt{\%.17f}，才看得到記憶體裡真正存的
\texttt{0.10000000000000001}。 這跟上一節範例 \thechapter.7 用完整
精度印 \texttt{0.1} 時看到的尾數是同一回事。 範例 \thechapter.8 把這三行並排印
出。 解題時若只用眼睛對輸出的短格式，很容易以為浮點數「算準了」，
等到後面用 \texttt{==} 比較才突然失敗。

\begin{example}{c}{(prec-mask.c) 同一筆 0.1，三種小數位數}
#include <stdio.h>

int main(void) {
    double x = 0.1;
    printf("%.1f\n", x);
    printf("%.2f\n", x);
    printf("%.17f\n", x);
    return 0;
}
\end{example}

\iosample{}{
0.1 \\
0.10 \\
0.10000000000000001 \\
}

\keypoint{題目要求印到小數點後 $k$ 位，是在規定輸出怎麼捨入，不是
保證變數裡的值精確到那一位。 比較相等仍然要用容許誤差，不能靠
\texttt{printf} 的外觀。}

\begin{drills}[label=drills:float-hidden-error]
\item \texttt{0.1} 用 \texttt{\%.1f}、\texttt{\%.2f}、\texttt{\%.17f}
分別印出來，結果一樣「乾淨」嗎？ 要看到真正的近似值該用哪一個？
\item 「印出來看起來是整數」跟「變數裡的值精確等於那個整數」是同一件
事嗎？
\end{drills}

\section{為什麼不能用 \texttt{==} 比較浮點數}
\label{sec:float-no-equality}

上一節用 \texttt{\%.1f} 把 \texttt{0.1} 印成乾乾淨淨的小數，但
\texttt{==} 比的是記憶體裡的位元樣式，不是輸出長相。 這一節看這件事在
「比較相等」上會造成什麼後果。 範例 \thechapter.9 把 $0.1$ 加十次，理論上應該等於
$1.0$，但印出完整精度後可以看到誤差。 而且用 \texttt{==} 直接比較會得到
「不相等」。

\begin{example}{c}{(float-precision-demo.c) 觀察浮點數的累積誤差}
#include <stdio.h>
#include <math.h>

int main(void) {
    // 把 0.1 加十次，理論上應該等於 1.0
    double sum = 0.0;
    for (int i = 0; i < 10; i++) {
        sum += 0.1;
    }

    printf("%.17f\n", sum);
    printf("%d\n", sum == 1.0);
    printf("%d\n", fabs(sum - 1.0) < 1e-9);
}
\end{example}

\iosample{}{
0.99999999999999989 \\
0 \\
1 \\
}

第一行印出的完整精度顯示，\texttt{sum} 其實是 $0.99999999999999989$，
跟數學上的 $1.0$ 差了一點點；第二行用 \texttt{==} 直接比較，因為兩者不
是真的相等，得到 \texttt{0}（假）；第三行改用容許誤差比較，因為差距遠
小於 \texttt{1e-9}，得到 \texttt{1}（真）。

為什麼 \texttt{==} 會失敗？ 因為 \texttt{==} 比的是「記憶體裡那兩個
位元樣式是否完全相同」，不是「數學上是否相等」。 每次 \texttt{sum +=
0.1}，加進去的 \texttt{0.1} 本身就已經是近似值，十次累加後誤差會一點一
點堆上去；而 \texttt{1.0} 在記憶體裡是另一個可精確表示的值，兩者的位
元樣式自然對不上。 這不是 C 語言的特例。 幾乎所有遵循 IEEE 754 的語言
都有同樣的問題。

\begin{code}{c}{用容許誤差 eps 比較兩個浮點數是否「足夠接近」}
#include <math.h>
double eps = 1e-9;
if (fabs(a - b) < eps) {
    // 視為相等
}
\end{code}

片語 \thechapter.10 是浮點數比較的標準寫法：先算兩者的差距，再跟一個很
小的容許值 \texttt{eps} 比。 \texttt{eps} 取多少要看題目要求的精度。 
輸出到小數點後 $3$ 位，\texttt{1e-6} 通常就夠；需要更高精度時再調小。
\keypoint{永遠不要用 \texttt{==} 直接比較兩個浮點數；應該判斷兩者的差
有沒有小於一個很小的容許誤差 \texttt{eps}，也就是
\texttt{fabs(a - b) < eps}。}

\begin{drills}[label=drills:float-no-equality]
\item 把 \texttt{0.1} 加十次，數學上應該等於 \texttt{1.0}；這台機器上
用 \texttt{==} 比較會得到真還是假？ 為什麼？
\item 正確比較兩個浮點數是否「足夠接近」該怎麼寫？
\item \texttt{==} 比的到底是「數學上是否相等」還是「記憶體裡的位元樣
式是否相同」？
\end{drills}

\section{0.1 加 0.2，以及絕對誤差對相對誤差}
\label{sec:float-abs-rel-error}

把 $0.1$ 加十次會偏離 $1.0$，更短的式子也一樣。 範例
\thechapter.11 計算 \texttt{0.1 + 0.2}：這台機器上印出
\texttt{0.30000000000000004}，跟常數 \texttt{0.3} 做 \texttt{==} 得到
$0$。 不是加法算錯，是三個字面常數各自被捨入成不同的可表示值。 
\texttt{0.1}、\texttt{0.2}、\texttt{0.3} 在二進位裡都是循環小數，
兩段捨入加起來，不必剛好等於第三段捨入。

\begin{example}{c}{(point-one-two.c) 0.1 加 0.2 並不恰好是 0.3}
#include <stdio.h>

int main(void) {
    double s = 0.1 + 0.2;
    printf("%.17f\n", s);
    printf("%d\n", s == 0.3);
    return 0;
}
\end{example}

\iosample{}{
0.30000000000000004 \\
0 \\
}

\texttt{fabs(a - b) < 1e-9} 是絕對誤差：不管數字本身有多大，都只允
許差 $10^{-9}$。 兩個接近 $10^{12}$ 的數，差 $0.5$ 在數學上相對很小
（大約十兆分之零點五），但絕對誤差測試會判「不相等」，因為 $0.5$
遠大於 $10^{-9}$。 這不是比較函式寫錯，是門檻跟數字的數量級不匹配：
大數字在浮點格式裡本來就比較「疏」，$10^{12}$ 附近相鄰兩個可表示值
的間距已經遠大於 $10^{-9}$，要求兩個大數差不到 $10^{-9}$，等於在要
求硬體做不到的事。

範例 \thechapter.12 對 \texttt{1e12} 與 \texttt{1e12 + 0.5} 分別用
絕對與相對門檻：前者失敗，後者
\texttt{fabs(a - b) < 1e-9 * fabs(a)} 成功。 相對誤差把容許差距
按 $|a|$ 放大，數量級大的時候門檻也跟著變寬。 題目若沒指定，先用
絕對誤差（對齊輸出位數最直覺）；數值範圍可能差很多個數量級時，才
考慮相對誤差。 比較前也可以先把兩個數都除掉某個典型大小，變成「相
對 $1$ 附近」再比，觀念上跟相對誤差是同一件事。

\begin{drills}[label=drills:float-abs-rel-error]
\item \texttt{0.1 + 0.2 == 0.3} 這台機器上印出真還是假？
\item 兩個接近 $10^{12}$ 的數字差 $0.5$，用絕對誤差 \texttt{1e-9} 測試
會判斷「相等」還是「不相等」？ 為什麼跟直覺不一樣？
\item 相對誤差的門檻該怎麼寫？（提示：把絕對誤差按哪個值放大？）
\end{drills}

\begin{example}{c}{(rel-eps.c) 絕對誤差與相對誤差}
#include <stdio.h>
#include <math.h>

int main(void) {
    double a = 1e12;
    double b = 1e12 + 0.5;
    printf("%d\n", fabs(a - b) < 1e-9);
    printf("%d\n", fabs(a - b) < 1e-9 * fabs(a));
    return 0;
}
\end{example}

\iosample{}{
0 \\
1 \\
}

\section{float.h 裡的精度常數}
\label{sec:float-limits}

編譯器在 \texttt{<float.h>} 提供這台機器上浮點型別的範圍與精度，不
必背數字。 範例 \thechapter.13 印出 \texttt{DBL\_DIG}（十進位有效位數
的保證值）、\texttt{DBL\_EPSILON}（$1$ 與下一個可表示值的間距），以
及 \texttt{DBL\_MAX}。 這台機器上分別觀察到 $15$、$2.22045\times
10^{-16}$、以及約 $1.80\times 10^{308}$。 與 IEEE 754 雙精度的常見數
字一致，但這是觀察結果，不是語言保證每台電腦都印同一個字串。

\begin{example}{c}{(float-limits.c) 印出 double 的精度與範圍}
#include <stdio.h>
#include <float.h>

int main(void) {
    printf("%d\n", DBL_DIG);
    printf("%.5e\n", DBL_EPSILON);
    printf("%.5e\n", DBL_MAX);
    return 0;
}
\end{example}

\iosample{}{
15 \\
2.22045e-16 \\
1.79769e+308 \\
}

\texttt{DBL\_EPSILON} 量的是「$1$ 附近」的間距。 前面說過，浮點數越
大，相鄰可表示值差得越遠，所以在 $10^{12}$ 附近用
\texttt{1e-9} 當絕對門檻會失敗，並不是 \texttt{DBL\_EPSILON} 變了，
而是那個間距已經被指數放大。 \texttt{eps} 若小過
\texttt{DBL\_EPSILON} 很多，對接近 $1$ 的數根本分辨不了「下一個可表
示值」；若大到 $0.01$，又等於放棄小數點後兩位。 門檻要對齊題目要求
的輸出位數，而不是抄一個永遠正確的魔術數字。

\begin{drills}[label=drills:float-limits]
\item \texttt{<float.h>} 要 include 才能用；\texttt{DBL\_EPSILON} 量
的是「哪個數附近」的間距？
\item 為什麼在 $10^{12}$ 附近用 \texttt{1e-9} 當絕對誤差門檻會失敗？
是 \texttt{DBL\_EPSILON} 變了嗎？
\end{drills}

\section{無窮大與 NaN}
\label{sec:float-inf-nan}

IEEE 754 把「超出範圍」與「不是一個數」編進位元樣式裡。 浮點除以零
在多數實作上不會像整數那樣未定義：\texttt{1.0 / 0.0} 得到正無窮，
\texttt{0.0 / 0.0} 得到 NaN（Not a Number）。 範例 \thechapter.14
用 \texttt{isinf}、\texttt{isnan} 檢查這台機器上的結果，並確認正無
窮比一個很大的有限數還大。

\begin{example}{c}{(inf-nan.c) 觀察無窮大與 NaN}
#include <stdio.h>
#include <math.h>

int main(void) {
    double inf = 1.0 / 0.0;
    double nanv = 0.0 / 0.0;
    printf("%d\n", isinf(inf));
    printf("%d\n", isnan(nanv));
    printf("%d\n", inf > 1e300);
    return 0;
}
\end{example}

\iosample{}{
1 \\
1 \\
1 \\
}

NaN 最容易讓人踩的坑是：它跟任何值比較（包括它自己）都不會得到
「小於」或「等於」那種直覺結果。 這台機器上
\texttt{nanv == nanv}、\texttt{nanv < 0.0}、\texttt{nanv > 0.0} 三
個比較全部印出 $0$。 所以一旦算式可能產生 NaN（例如對負數開平方
根、對非正數取對數），後面的 \texttt{if} 可能整條走錯。 條件看起來
寫了，但永遠不會成立。 發現輸出變成 \texttt{nan} 或 \texttt{inf}，
先回頭查有沒有除以零、對數的真數非正、或平方根的參數為負，再用
\texttt{isnan}／\texttt{isinf} 把異常值擋下來，不要假裝它還是普通
的有限數字。

\keypoint{看到 \texttt{nan} 或 \texttt{inf}，先查定義域（除數、對
數、開根），不要繼續拿這個值去做比較或累加。}

\begin{drills}[label=drills:float-inf-nan]
\item \texttt{1.0 / 0.0} 跟 \texttt{0.0 / 0.0} 這台機器上分別得到什麼？
\item \texttt{nanv == nanv} 這種「NaN 跟自己比較」會是真還是假？
\item 看到輸出變成 \texttt{nan} 或 \texttt{inf}，第一步該檢查什麼？
\end{drills}

\section{能用整數算的錢，就不要用浮點}
\label{sec:float-money-int}

金額看起來像小數（$1.99$ 元），本質上卻是離散的「分」。 若用
\texttt{double} 去累加一長串價錢，每一筆的 $0.1$、$0.01$ 都會帶進
二進位捨入，加到最後可能連「該不該剛好是整數元」都說不準。 更穩的
做法是全程用整數記「分」，只在最後輸出時除以 $100$ 變成小數點。

\begin{code}{c}{金額用「分」累加，最後再換成元}
long long cents = 0;
cents += price_in_cents;
printf("%lld.%02lld\n", cents / 100, cents % 100);
\end{code}

範例 \thechapter.16 讀 $n$ 筆以分為單位的價錢，加總後印成「元.分」。
輸入 $199$、$250$、$51$ 分，總和 $500$ 分，正好是 $5.00$ 元。 全程
沒有浮點加法。 \texttt{\%02lld} 讓分的部分固定兩位，所以 $5$ 元 $0$
分仍會印成 \texttt{5.00}，不會變成 \texttt{5.0}。

\begin{example}{c}{(cents-sum.c) 用整數分累加金額}
#include <stdio.h>

int main(void) {
    int n;
    scanf("%d", &n);

    long long cents = 0;
    for (int i = 0; i < n; i++) {
        int p;
        scanf("%d", &p);
        cents += p;
    }

    printf("%lld.%02lld\n", cents / 100, cents % 100);
    return 0;
}
\end{example}

\iosample{
3 \\
199 \\
250 \\
51 \\
}{
5.00 \\
}

問題本身是有理數或離散量時，全程用整數往往比較乾淨。 溫度、距離、
統計才是浮點數真正該上場的地方。

\begin{drills}[label=drills:float-money-int]
\item 為什麼金額用 \texttt{double} 直接累加，加多了可能連整數元都對不
準？
\item \texttt{\%02lld} 印分的部分，作用是什麼？ 少了它 \texttt{5.00}
可能印成什麼？
\item 199、250、51 分加起來，換算成「元.分」該怎麼印？
\end{drills}

\section{避免比較浮點數：用整數交叉相乘}
\label{sec:float-cross-mult}

上一節用整數記「分」累加金額；這一節換一個角度：比較兩個分數的大小，
也可以全程只用整數乘法，搜尋過程完全不出現浮點數。 假設要在好幾組「折扣金額
$d$、原價 $p$」中，找出折扣率 $d / p$ 最高的一組。 如果 $d$ 跟 $p$ 都是
整數，與其每組都算一次除法比較，不如用交叉相乘比較兩個分數的大小：
$\dfrac{d_1}{p_1} > \dfrac{d_2}{p_2} \iff d_1 \times p_2 > d_2 \times
p_1$（假設 $p_1, p_2 > 0$）。 右式全部都是整數運算，完全不會有浮點誤差，
也不用擔心「比較兩個很接近的浮點數」的問題，因為整個搜尋過程根本不曾出
現浮點數。

\begin{example}{c}{(best-discount-rate.c) 用整數交叉相乘找出最高折扣率}
#include <stdio.h>

int main(void) {
    int n;
    scanf("%d", &n);

    // 用交叉相乘比較 d/p，搜尋過程完全不用浮點數
    long long bestD = 0, bestP = 1;
    for (int i = 0; i < n; i++) {
        long long d, p;
        scanf("%lld %lld", &d, &p);
        if (d * bestP > bestD * p) {
            bestD = d;
            bestP = p;
        }
    }

    printf("%.4f\n", (double)bestD / bestP);
}
\end{example}

\iosample{
3 \\
3 10 \\
5 20 \\
9 20 \\
}{
0.4500 \\
}

三組折扣分別是 $0.3$、$0.25$、$0.45$，程式全程只用整數乘法比較大小，
最後印出結果時才轉成浮點數：\texttt{9/20 = 0.4500} 確實是三者之中最高的。
比較 \texttt{9 * 1 > 0 * 20} 時，左邊 $9$、右邊 $0$，更新最佳答案；下一
輪 \texttt{5 * 20 > 9 * 20} 不成立（$100 \not> 180$）；第三輪
\texttt{9 * 20 > 5 * 20} 成立，最終得到 \texttt{bestD=9, bestP=20}。

\keypoint{搜尋過程全部用整數運算完成，只在最後印出結果時才做一次除法；
能用整數解決的比較，就不要提早換成浮點數。}

\begin{drills}[label=drills:float-cross-mult]
\item 比較 $\frac{d_1}{p_1}$ 跟 $\frac{d_2}{p_2}$ 哪個大，交叉相乘要怎
麼寫成不含除法的整數比較？
\item 這個技巧的前提是 $p_1, p_2$ 要符合什麼條件？
\item 為什麼交叉相乘完全不用擔心浮點數比較的誤差問題？
\end{drills}

\section{數學函式庫：以計算兩點距離為例}
\label{sec:float-math-lib}

C 的標準函式庫把平方根、三角函數、對數等數學函式集中放在 \texttt{math.h}
裡，使用前要 \texttt{\#include <math.h>}。 範例 \thechapter.18 用兩點座
標算歐幾里得距離 $\sqrt{(x_2-x_1)^2 + (y_2-y_1)^2}$，會用到 \texttt{sqrt}
函式；這是我們第一次在程式裡使用「數學函式庫」而不是自己手刻演算法，之
後遇到平方根、對數、三角函數，都可以直接呼叫對應的函式，不用重新發明。

\begin{example}{c}{(euclidean-distance.c) 計算兩點距離}
#include <stdio.h>
#include <math.h>

int main(void) {
    double x1, y1, x2, y2;
    scanf("%lf %lf %lf %lf", &x1, &y1, &x2, &y2);

    // 歐幾里得距離公式
    double dx = x2 - x1, dy = y2 - y1;
    printf("%.3f\n", sqrt(dx * dx + dy * dy));
}
\end{example}

\iosample{
0 0 3 4 \\
}{
5.000 \\
}

範例輸入的兩點是 $(0, 0)$ 與 $(3, 4)$，$dx = 3$、$dy = 4$，
$\sqrt{3^2 + 4^2} = \sqrt{25} = 5$，剛好是最經典的一組整數邊長直角三角
形，方便我們用心算驗證程式對不對。 \texttt{sqrt} 回傳的是
\texttt{double}，我們用 \texttt{\%.3f} 四捨五入到小數點後三位輸出。

\keypoint{使用 \texttt{math.h} 裡的函式（如 \texttt{sqrt}）時，編譯常
需要額外加上 \texttt{-lm} 連結數學函式庫，例如
\texttt{gcc program.c -o program -lm}。}

\begin{drills}[label=drills:float-math-lib]
\item 兩點 $(0,0)$ 與 $(3,4)$ 之間的歐幾里得距離是多少？
\item 用 \texttt{math.h} 的函式時，編譯常常需要額外加上什麼選項？
\end{drills}

\section{floor、ceil、round 與 fmod}
\label{sec:float-round-fns}

\texttt{math.h} 除了剛才用到的 \texttt{sqrt}，還有往不同方向
取整的函式。 \texttt{floor} 往 $-\infty$ 取整，\texttt{ceil} 往
$+\infty$ 取整，\texttt{round} 對「恰好一半」的情況往遠離零的整數
靠（C99）。 範例 \thechapter.19 對輸入 \texttt{3.5} 印出
\texttt{3 4 4}：$\lfloor 3.5\rfloor = 3$，$\lceil 3.5\rceil = 4$，
四捨五入也是 $4$。

這些函式對負數的方向最容易搞反。 對 $-1.5$，這台機器上
\texttt{floor} 得到 $-2$（往更小的整數）、\texttt{ceil} 得到 $-1$
（往比較大的整數）、\texttt{round} 得到 $-2$（一半時遠離零）。 如
果你心裡想的是「拿掉小數部分」，那比較接近往零取整，C99 另有
\texttt{trunc}；\texttt{floor} 不是那個意思。

\begin{example}{c}{(round-fns.c) 下取整、上取整與四捨五入}
#include <stdio.h>
#include <math.h>

int main(void) {
    double x;
    scanf("%lf", &x);
    printf("%.0f %.0f %.0f\n", floor(x), ceil(x), round(x));
    return 0;
}
\end{example}

\iosample{
3.5 \\
}{
3 4 4 \\
}

浮點數沒有 \texttt{\%}。 要「實數除法的餘數」用 \texttt{fmod}。 範例
\thechapter.20：\texttt{7.5} 除以 \texttt{2.0}，商是 $3$、餘
\texttt{1.50}，因為 $3 \times 2.0 + 1.50 = 7.5$。 跟整數 \texttt{\%}
一樣，它給的是「扣掉整數倍之後剩下的」，不是四捨五入。 週期現象
（角度模 $2\pi$、時間模一天的秒數）常用這一招，但別忘了浮點餘數本
身也帶捨入誤差，累加很多次之後仍可能慢慢漂離。

\begin{example}{c}{(fmod-demo.c) 浮點數的餘數}
#include <stdio.h>
#include <math.h>

int main(void) {
    double a, b;
    scanf("%lf %lf", &a, &b);
    printf("%.2f\n", fmod(a, b));
    return 0;
}
\end{example}

\iosample{
7.5 2.0 \\
}{
1.50 \\
}

距離公式若寫成 \texttt{sqrt(x * x + y * y)}，當 $x$、$y$ 都很大時，
平方可能先溢成無窮，再開根就得到 \texttt{inf}，即使真正的距離其實
還在 \texttt{double} 能表示的範圍內。 \texttt{hypot(x, y)} 是函式庫
提供的較穩寫法，內部會先縮放再算，減少這種中間溢位。 範例
\thechapter.21 對 \texttt{3 4} 仍得到 \texttt{5.0000}，跟前面的兩點
距離一致；日常測資兩者沒差，極端座標時差就出現了。

\begin{example}{c}{(hypot-demo.c) 用 hypot 算歐氏長度}
#include <stdio.h>
#include <math.h>

int main(void) {
    double x, y;
    scanf("%lf %lf", &x, &y);
    printf("%.4f\n", hypot(x, y));
    return 0;
}
\end{example}

\iosample{
3 4 \\
}{
5.0000 \\
}

\stylepoint{座標可能很大時，優先寫 \texttt{hypot(dx, dy)}，而不是
自己平方再開根；正確性通常一樣，數值穩定性比較好。}

\begin{drills}[label=drills:float-round-fns]
\item 對 $-1.5$，\texttt{floor}、\texttt{ceil}、\texttt{round} 這台機
器上分別得到什麼？
\item \texttt{fmod(7.5, 2.0)} 的結果是多少？ 它跟整數的 \texttt{\%} 概
念上做的是同一件事嗎？
\item 為什麼 \texttt{sqrt(x*x + y*y)} 在 $x$、$y$ 很大時可能不如
\texttt{hypot(x, y)} 穩定？
\end{drills}

\section{用對數計算位數，避免處理超大整數}
\label{sec:float-log-digits}

一個正整數 $N$ 的十進位位數是 $\lfloor \log_{10} N \rfloor + 1$，而
$\log_{10}(a^b) = b \times \log_{10} a$。 這代表我們可以直接用對數算出
$a^b$ 有幾位數，不需要真的算出 $a^b$ 本身。 畢竟 $a^b$ 可能大到連
\texttt{long long} 都裝不下，而位數只是一個很小的整數，用對數就能一步
算出來，完全不用面對「數字太大存不下」的問題。

\begin{example}{c}{(power-digit-count.c) 計算 a 的 b 次方有幾位數}
#include <stdio.h>
#include <math.h>

int main(void) {
    int a, b;
    scanf("%d %d", &a, &b);

    // log10(a^b) = b * log10(a)，位數 = floor(log10(a^b)) + 1
    double logVal = b * log10((double)a);
    int digits = (int)floor(logVal + 1e-9) + 1;

    printf("%d\n", digits);
}
\end{example}

\iosample{
2 10 \\
}{
4 \\
}

驗證一下：$2^{10} = 1024$，確實是 $4$ 位數，跟 $\lfloor \log_{10} 1024
\rfloor + 1 = \lfloor 3.01 \rfloor + 1 = 3 + 1 = 4$ 吻合。 這裡我們
把 \texttt{log10} 的結果存在 \texttt{double} 裡。 對數本身也是浮點
數，所以最後取 \texttt{floor} 轉成整數時，要記得上一節的 \texttt{floor} 用法，以及前面容許誤差比較裡 \texttt{1e-9} 的防禦性寫法。

\stylepoint{\texttt{floor} 前先加一個很小的 \texttt{1e-9} 是防禦性寫
法：\texttt{log10} 的結果理論上可能因浮點誤差落在整數邊界下方，先加一
點容許值可以避免少算一位。}

\begin{drills}[label=drills:float-log-digits]
\item $2^{10} = 1024$ 是幾位數？ 用 $\log_{10}$ 怎麼算出這個位數，不必
真的算出 $2^{10}$ 本身？
\item 為什麼這個技巧特別適合 $a^b$ 大到 \texttt{long long} 裝不下的情況？
\item \texttt{floor(logVal + 1e-9)} 裡加 \texttt{1e-9} 是為了防止什麼
情況？
\end{drills}

\section{綜合應用：凸多邊形的周長}
\label{sec:float-polygon-perimeter}

先用一個把距離累加起來的綜合範例：給定一個凸多邊形依序排列的頂點座
標，求周長。
做法是把本章學到的技巧串起來。 用陣列存下所有頂點、用迴圈依序走訪相鄰
的兩個頂點、用 \texttt{sqrt} 算兩點距離、把每一段的距離加總。 最後一個
頂點要繞回第一個頂點，這裡沿用陣列那一章介紹過的 \texttt{\% n} 環繞寫法：
當 \texttt{i} 是最後一個索引 $n - 1$ 時，\texttt{(i + 1) \% n} 會繞回
$0$，正好是第一個頂點。

\begin{example}{c}{(polygon-perimeter.c) 計算凸多邊形周長}
#include <stdio.h>
#include <math.h>

int main(void) {
    int n;
    scanf("%d", &n);

    double x[1005], y[1005];
    for (int i = 0; i < n; i++) {
        scanf("%lf %lf", &x[i], &y[i]);
    }

    // 累加每一段相鄰頂點的距離，最後一段用 % n 繞回第一個頂點
    double perimeter = 0;
    for (int i = 0; i < n; i++) {
        int j = (i + 1) % n;
        double dx = x[j] - x[i], dy = y[j] - y[i];
        perimeter += sqrt(dx * dx + dy * dy);
    }

    printf("%.3f\n", perimeter);
}
\end{example}

\iosample{
4 \\
0 0 \\
4 0 \\
4 4 \\
0 4 \\
}{
16.000 \\
}

範例輸入是一個邊長 $4$ 的正方形，四段邊長分別是 $4, 4, 4, 4$，加總後周
長是 $16$，跟輸出一致。 這個範例也再次用到本章強調的重點：每一段距離都
是用 \texttt{sqrt} 算出來的浮點數，而 \texttt{perimeter} 是好幾個浮點
數累加的結果，如果題目要求判斷周長是不是恰好等於某個整數，就要用容許誤
差比較（片語 \thechapter.10 的寫法），而不是 \texttt{==}。
題庫後半還會碰到「頂點不是直接給座標、而是兩條直線的交點」。 那個
寫法留到本章更後面，先把「有座標就能累加邊長」練熟。

\begin{drills}[label=drills:float-polygon-perimeter]
\item 邊長 $4$ 的正方形四個頂點依序給定時，周長怎麼算？
\item \texttt{(i + 1) \% n} 在 \texttt{i} 是最後一個索引時算出什麼？
這個寫法的用途是什麼？
\item 若要判斷周長是不是恰好等於某個整數，該用 \texttt{==} 還是容許誤
差比較？
\end{drills}

\section{弧長、給定的 pi，以及繞過障礙的面積}
\label{sec:float-arc-area}

周長不一定是直線。 若路徑是半徑 $r$ 的圓的四分之三弧，長度是
$\frac{3}{4} \times 2\pi r$。 題目若給定 $\pi$，就用那個常數，不要
自己另記一串數字。 兩個人各背一組不同位數的 $\pi$，答案會在小數點
後幾位對不起來。 範例 \thechapter.24 用
\texttt{3.141592653589793} 計算 $r=10$ 的弧長：
$1.5 \times \pi \times 10 \approx 47.123889\ldots$，印到小數點後三
位是 \texttt{47.124}。

\begin{example}{c}{(three-quarter-arc.c) 四分之三圓周}
#include <stdio.h>

int main(void) {
    double r;
    scanf("%lf", &r);
    double pi = 3.141592653589793;
    double len = 0.75 * 2.0 * pi * r;
    printf("%.3f\n", len);
    return 0;
}
\end{example}

\iosample{
10 \\
}{
47.124 \\
}

面積也常是「整圓的幾分之幾」。 想像一條繩子拴在矩形房屋的牆角，房
屋擋住一個象限，繩子夠短、還繞不過牆的時候，動物能走的就是其餘
$\frac{3}{4}$ 圓。 繩子比某一邊牆還長時，多出來的長度可以沿著那一面
牆再繞出一個半徑更短的 $\frac{1}{4}$ 圓。 兩邊牆都可能發生，所以最
多要加兩次。

\begin{center}
\begin{tikzpicture}[scale=0.7]
    \fill[black!12] (0,0) rectangle (2.2,1.6);
    \draw (0,0) rectangle (2.2,1.6);
    \node at (1.1,0.8) {\small 房屋};
    \fill (0,0) circle (1.6pt);
    \draw[thick] (0,2.4) arc[start angle=90, end angle=360, radius=2.4];
    \node[idxlbl] at (0.15,-0.35) {拴點};
    \node[idxlbl] at (-1.55,1.1) {$\frac{3}{4}$ 圓};
\end{tikzpicture}
\diagcaption{房屋占掉一個象限時，短繩子掃過的是其餘四分之三圓；繩
子長過牆邊才會再繞出更小的四分之一圓。}
\label{fig:ch07-rope-corner-arc}
\end{center}

範例 \thechapter.25 用題目給定的 $\pi = 3.1415926$。 輸入邊長
$10$、$10$ 與繩長 $5$：繩子比兩邊都短，面積就是
$\frac{3}{4}\pi\cdot 5^2 = \frac{3}{4}\times 3.1415926\times 25 =
58.90486125$，用 \texttt{\%f} 印六位小數得到 \texttt{58.904861}。
若繩長改成 $15$，兩邊都會繞出半徑 $5$ 的四分之一圓，面積變成
$\frac{3}{4}\pi\cdot 15^2 + 2\times\frac{1}{4}\pi\cdot 5^2$，這台
機器印出 \texttt{569.413659}。 先畫圖判斷「會不會繞過牆」，再決定
要不要加上那兩個小扇形，比一開始就堆公式不容易漏情況。

\begin{example}{c}{(cow-graze.c) 四分之三圓加上繞牆的扇形}
#include <stdio.h>

int main(void) {
    double a, b, c;
    scanf("%lf %lf %lf", &a, &b, &c);
    double pi = 3.1415926;
    double area = 0.75 * pi * c * c;
    if (c > a) {
        double r = c - a;
        area += 0.25 * pi * r * r;
    }
    if (c > b) {
        double r = c - b;
        area += 0.25 * pi * r * r;
    }
    printf("%f\n", area);
    return 0;
}
\end{example}

\iosample{
10.0 10.0 5.0 \\
}{
58.904861 \\
}

\stylepoint{題目寫「假設 $\pi = \cdots$」時，把那一串數字原樣寫進程
式，不要改用 \texttt{acos(-1.0)} 另算一個 $\pi$。 評測比的是輸出字
串，常數不一致就會錯在最後幾位。}

\begin{drills}[label=drills:float-arc-area]
\item 半徑 $r$ 的圓，四分之三弧長的公式是什麼？
\item 題目給定一個 $\pi$ 的近似值時，該用題目給的，還是自己用
\texttt{acos(-1.0)} 另外算一個？ 為什麼？
\item 繩子拴在牆角，什麼情況下動物能走的範圍不只是四分之三圓？
\end{drills}

\section{由直線交點圍成多邊形}
\label{sec:float-line-intersect}

稍早的周長範例直接讀頂點座標。 另一種常見輸入是：$n$ 條直線
$ax+by+c=0$，相鄰兩條的交點才是頂點，再依序把邊長加起來。 兩條直線
$a_1x+b_1y+c_1=0$、$a_2x+b_2y+c_2=0$ 的交點可以用二階克拉瑪公式：

\[
x = \frac{b_1 c_2 - b_2 c_1}{a_1 b_2 - a_2 b_1},\quad
y = \frac{c_1 a_2 - c_2 a_1}{a_1 b_2 - a_2 b_1}.
\]

分母是兩個法向量的行列式，平行時為 $0$。 凸多邊形的相鄰邊不應該平
行，測資通常保證交點存在。 第 $i$ 條與第 $i+1$ 條交出第 $i$ 個頂
點，最後一條再跟第一條交出最後一個頂點，然後跟前面一樣用
\texttt{\% n} 繞一圈累加 \texttt{hypot}。

\begin{center}
\begin{tikzpicture}[scale=0.55]
    \draw[thick] (-0.5,0) -- (5.5,0);
    \draw[thick] (0,-0.5) -- (0,5.5);
    \draw[thick] (4,-0.5) -- (4,5.5);
    \draw[thick] (-0.5,4) -- (5.5,4);
    \fill (0,0) circle (2pt);
    \fill (4,0) circle (2pt);
    \fill (4,4) circle (2pt);
    \fill (0,4) circle (2pt);
    \node[idxlbl, below left] at (0,0) {$(0,0)$};
    \node[idxlbl, below right] at (4,0) {$(4,0)$};
    \node[idxlbl, above right] at (4,4) {$(4,4)$};
    \node[idxlbl, above left] at (0,4) {$(0,4)$};
\end{tikzpicture}
\diagcaption{四條直線 $x=0$、$y=0$、$x=4$、$y=4$ 兩兩相鄰相交，頂點
正好是邊長 $4$ 的正方形。}
\label{fig:ch07-polygon-intersections}
\end{center}

範例 \thechapter.26 就用這四條直線（寫成
$1\cdot x+0\cdot y+0=0$ 這種一般式）。 四個交點是
$(0,0)$、$(4,0)$、$(4,4)$、$(0,4)$，周長 $16$，跟直接給頂點的正方形
例子一致。 寫的時候每一條直線都用 \texttt{double} 讀 $a,b,c$，交點
也存 \texttt{double}；最後用 \texttt{\%.3f} 輸出，對齊常見的周長精
度要求。

\begin{example}{c}{(line-polygon.c) 相鄰直線交點圍成的周長}
#include <stdio.h>
#include <math.h>

int main(void) {
    int n;
    scanf("%d", &n);

    double a[250], b[250], c[250];
    for (int i = 0; i < n; i++) {
        scanf("%lf %lf %lf", &a[i], &b[i], &c[i]);
    }

    double px[250], py[250];
    for (int i = 0; i < n; i++) {
        int j = (i + 1) % n;
        double det = a[i] * b[j] - a[j] * b[i];
        px[i] = (b[i] * c[j] - b[j] * c[i]) / det;
        py[i] = (c[i] * a[j] - c[j] * a[i]) / det;
    }

    double peri = 0.0;
    for (int i = 0; i < n; i++) {
        int j = (i + 1) % n;
        peri += hypot(px[j] - px[i], py[j] - py[i]);
    }
    printf("%.3f\n", peri);
    return 0;
}
\end{example}

\iosample{
4 \\
1 0 0 \\
0 1 0 \\
1 0 -4 \\
0 1 -4 \\
}{
16.000 \\
}

\begin{drills}[label=drills:float-line-intersect]
\item 兩條直線 $a_1x+b_1y+c_1=0$ 與 $a_2x+b_2y+c_2=0$ 的交點公式裡，
分母是什麼？ 它等於 $0$ 代表什麼幾何情況？
\item 求出所有交點之後，怎麼把它們串成周長？
\end{drills}

\section{累加與數值穩定性：平均與變異數}
\label{sec:float-mean-variance}

有些題目需要對一整排浮點數做統計，例如平均值與變異數。 變異數的公式是
「每個值與平均值的差的平方，取平均」，所以要先知道平均值才能算每個值的
差。 得先掃一遍陣列算出平均值，才能再掃第二遍算變異數。 這跟陣列那一
章「先掃一遍拿到資訊、再掃第二遍利用它」的節奏完全一樣，差別只在於這裡
處理的是浮點數，多了一層數值精確度的考量。 前面也看過：浮點數
累加會一點一點堆誤差，統計題尤其要留意。

\begin{example}{c}{(mean-and-variance.c) 兩階段計算平均與變異數}
#include <stdio.h>

int main(void) {
    int n;
    scanf("%d", &n);

    // 第一遍：讀入資料並累加，算出平均值
    double a[100005];
    double sum = 0;
    for (int i = 0; i < n; i++) {
        scanf("%lf", &a[i]);
        sum += a[i];
    }
    double mean = sum / n;

    // 第二遍：用平均值算出變異數
    double sq = 0;
    for (int i = 0; i < n; i++) {
        double d = a[i] - mean;
        sq += d * d;
    }
    double variance = sq / n;

    printf("%.4f %.4f\n", mean, variance);
}
\end{example}

\iosample{
5 \\
1 2 3 4 5 \\
}{
3.0000 2.0000 \\
}

用 $1, 2, 3, 4, 5$ 驗證：平均值是 $\frac{1+2+3+4+5}{5} = 3$；每個值與
平均值的差分別是 $-2, -1, 0, 1, 2$，平方後是
$4, 1, 0, 1, 4$，總和是 $10$，除以 $5$ 得到變異數 $2$，跟範例輸出一致。

變異數也可以用單次掃描公式 $\overline{x^2} - \overline{x}^2$ 邊讀邊算，
不用存整個陣列，也不用掃兩遍。 只要邊讀邊累加 $\sum x_i$ 跟
$\sum x_i^2$，最後一次算完就好，寫起來甚至比兩階段版本更短。 但兩個很接
近的大數相減會放大浮點誤差：當 $\overline{x^2}$ 跟 $\overline{x}^2$ 本
身都很大、而兩者的差相對很小時，這個差值的有效位數會被嚴重侵蝕，資料量
大或數值範圍大時甚至可能算出負的變異數（理論上變異數永遠不會是負的）。

\keypoint{兩階段的做法比單次掃描公式慢，但數值上更穩定，這是精確度與效
率之間常見的取捨。}

\begin{drills}[label=drills:float-mean-variance]
\item $1, 2, 3, 4, 5$ 這組資料的平均值與變異數各是多少？
\item 為什麼算變異數要先掃一遍拿到平均值，才能掃第二遍？
\item 單次掃描公式 $\overline{x^2} - \overline{x}^2$ 比兩階段法快，但
可能有什麼數值上的風險？
\end{drills}

\section{迭代逼近：牛頓法}
\label{sec:float-newton}

有些題目不給閉式公式，而要你把同一個更新做 $k$ 次。 求 $\sqrt{a}$ 可
以看成解 $f(x) = x^2 - a = 0$。 牛頓法的一步是沿著切線走到 $x$ 軸：
$x \leftarrow x - f(x)/f'(x)$。 對這個 $f$，化簡之後正好是
$x \leftarrow \dfrac{1}{2}\bigl(x + a/x\bigr)$。 每一步都把目前的猜測
跟「$a$ 除以猜測」平均一次，直觀上就是在逼近幾何平均。

範例 \thechapter.28 從 $x=1$ 開始，對 $a=9$ 做 $6$ 次。 我們逐步把每
一輪的值寫出來（這台機器、中間用更高精度印）：

\begin{center}
\begin{tabular}{cr}
輪次 & $x$ \\
1 & $5.000000$ \\
2 & $3.400000$ \\
3 & $3.023529$ \\
4 & $3.000092$ \\
5 & $3.000000$ \\
6 & $3.000000$ \\
\end{tabular}
\end{center}

第一輪從 $1$ 跳到 $5$，還差很遠；第三輪已經到 $3.02$，第四輪以後在
\texttt{\%.4f} 底下就跟精確的 $3$ 分不出來了。 次數不夠會停在近似
值；題目若規定做 $k$ 次再輸出，就真的只做 $k$ 次，不要自己改成「誤
差夠小就停」。 那是另一種寫法，答案對測資不一定符合。

\begin{example}{c}{(newton-sqrt.c) 牛頓迭代求平方根}
#include <stdio.h>

int main(void) {
    double a;
    int k;
    scanf("%lf %d", &a, &k);

    double x = 1.0;
    for (int i = 0; i < k; i++) {
        x = 0.5 * (x + a / x);
    }
    printf("%.4f\n", x);
    return 0;
}
\end{example}

\iosample{
9 6 \\
}{
3.0000 \\
}

推廣到任意多項式時，每一輪都要同時算出 $f(x)$ 與 $f'(x)$。 把係數
從高次寫到常數項，霍納法則（Horner）用一次掃描同時求出兩者：每讀進
一個係數，先用「目前的 $f$」去更新導數（乘上 $x$ 再累加），再用這個
係數更新 $f$。 片語 \thechapter.30 就是這一層；外面再套 $k$ 次
$x \leftarrow x - f/f'$。

\begin{code}{c}{霍納法則同時求多項式值與導數}
double f = 0.0, df = 0.0;
for (int i = 0; i <= degree; i++) {
    df = df * x + f;
    f = f * x + coef[i];
}
x = x - f / df;
\end{code}

範例 \thechapter.30 對 $f(x) = x^2 - 2$（係數 \texttt{1 0 -2}）從
$x=1$ 做 $5$ 次，每一步先印目前的 $x$ 與 $f(x)$，再更新。 第一行
\texttt{1.0000 -1.0000} 就是起點：$1^2-2=-1$。 下一步
$x \leftarrow 1 - (-1)/2 = 1.5$，對應第二行 $1.5^2-2=0.25$。 之後
很快靠近 $\sqrt{2}\approx 1.4142$，後兩行的函數值在四位小數下已經
是 $0$。 導數若變成 $0$，這一除會爆炸。 測資通常會保證切線不水平，
寫程式時仍該知道這個失敗模式。

\begin{drills}[label=drills:float-newton]
\item 牛頓法求 $\sqrt{a}$ 的更新公式是什麼？
\item 題目規定做 $k$ 次迭代時，該做滿 $k$ 次，還是「誤差夠小就提早停」？
\item 霍納法則在同一次掃描裡，同時算出多項式的哪兩個量？
\end{drills}

\begin{example}{c}{(newton-poly.c) 對多項式做牛頓迭代}
#include <stdio.h>

int main(void) {
    int d, k;
    scanf("%d", &d);

    double coef[12];
    for (int i = 0; i <= d; i++) {
        scanf("%lf", &coef[i]);
    }
    scanf("%d", &k);
    double x;
    scanf("%lf", &x);

    for (int t = 0; t < k; t++) {
        double f = 0.0, df = 0.0;
        for (int i = 0; i <= d; i++) {
            df = df * x + f;
            f = f * x + coef[i];
        }
        printf("%.4f %.4f\n", x, f);
        x = x - f / df;
    }
    return 0;
}
\end{example}

\iosample{
2 \\
1 0 -2 \\
5 \\
1 \\
}{
1.0000 -1.0000 \\
1.5000 0.2500 \\
1.4167 0.0069 \\
1.4142 0.0000 \\
1.4142 0.0000 \\
}

\section{上三角方程組的回代}
\label{sec:float-back-sub}

解一般的線性方程組要用高斯消去，但若係數矩陣已經是上三角（對角線
以下都是 $0$），從最後一個未知數往回代就夠了。 兩條式子
$a_{11}x + a_{12}y = b_1$、$a_{22}y = b_2$ 時，先
$y = b_2/a_{22}$，再 $x = (b_1 - a_{12} y)/a_{11}$。 範例
\thechapter.31 讀入 \texttt{2 1 5} 與 \texttt{1 1}，對應
$2x+y=5$、$y=1$，得到 $y=1$、$x=(5-1)/2=2$。

\begin{example}{c}{(back-sub.c) 二階上三角方程組的回代}
#include <stdio.h>

int main(void) {
    double a11, a12, b1;
    double a22, b2;
    scanf("%lf %lf %lf", &a11, &a12, &b1);
    scanf("%lf %lf", &a22, &b2);
    double y = b2 / a22;
    double x = (b1 - a12 * y) / a11;
    printf("%.4f %.4f\n", x, y);
    return 0;
}
\end{example}

\iosample{
2 1 5 \\
1 1 \\
}{
2.0000 1.0000 \\
}

$n$ 階時，第 $i$ 個未知數（從 $0$ 數）等於「右端減掉右邊已經求出的
那些項，再除以對角線」。 所以外層迴圈從 $n-1$ 走到 $0$，內層把
$j = i+1,\ldots,n-1$ 的 $a_{ij}x_j$ 扣掉。

\begin{center}
\begin{tikzpicture}
    \node[cell] (r0c0) {\texttt{a00}};
    \node[cell, right=0pt of r0c0] (r0c1) {\texttt{a01}};
    \node[cell, right=0pt of r0c1] (r0c2) {\texttt{a02}};

    \node[cell, below=0pt of r0c0] (r1c0) {\texttt{0}};
    \node[cell, right=0pt of r1c0] (r1c1) {\texttt{a11}};
    \node[cell, right=0pt of r1c1] (r1c2) {\texttt{a12}};

    \node[cell, below=0pt of r1c0] (r2c0) {\texttt{0}};
    \node[cell, right=0pt of r2c0] (r2c1) {\texttt{0}};
    \node[hlcell, right=0pt of r2c1] (r2c2) {\texttt{a22}};

    \node[idxlbl, right=0.8cm of r0c2] {$x_0$（最後算）};
    \node[idxlbl, right=0.8cm of r1c2] {$x_1$};
    \node[idxlbl, right=0.8cm of r2c2] {$x_2$（最先算）};
    \draw[diagarrow] (r2c2.east) -- ++(0.6,0);
\end{tikzpicture}
\diagcaption{上三角矩陣從右下角的對角線除起，每求出一個 $x_i$，就
拿去消掉上面那些列裡還沒處理的項。}
\label{fig:ch07-back-substitution}
\end{center}

範例 \thechapter.32 讀 $n=3$ 與

\[
\begin{bmatrix}
1 & 1 & 1 \\
0 & 1 & 1 \\
0 & 0 & 1
\end{bmatrix}
\begin{bmatrix} x_0 \\ x_1 \\ x_2 \end{bmatrix}
=
\begin{bmatrix} 6 \\ 3 \\ 1 \end{bmatrix}.
\]

先 $x_2 = 1/1 = 1$，再 $x_1 = (3-1\cdot 1)/1 = 2$，最後
$x_0 = (6-1\cdot 2-1\cdot 1)/1 = 3$。 對角線接近 $0$ 時，這一層會
把上面的捨入誤差放大。 這是數值線性代數的起點，現階段先保證測資
的樞紐不是零，也別在過程中把已經是 $0$ 的下三角元素拿來除。

\begin{example}{c}{(back-sub-n.c) n 階上三角回代}
#include <stdio.h>

int main(void) {
    int n;
    scanf("%d", &n);

    double a[16][16], y[16], x[16];
    for (int i = 0; i < n; i++) {
        for (int j = 0; j < n; j++) {
            scanf("%lf", &a[i][j]);
        }
    }
    for (int i = 0; i < n; i++) {
        scanf("%lf", &y[i]);
    }

    for (int i = n - 1; i >= 0; i--) {
        double s = y[i];
        for (int j = i + 1; j < n; j++) {
            s -= a[i][j] * x[j];
        }
        x[i] = s / a[i][i];
    }

    for (int i = 0; i < n; i++) {
        printf("%.4f\n", x[i]);
    }
    return 0;
}
\end{example}

\iosample{
3 \\
1 1 1 \\
0 1 1 \\
0 0 1 \\
6 \\
3 \\
1 \\
}{
3.0000 \\
2.0000 \\
1.0000 \\
}

\keypoint{浮點數適合連續量的近似計算。 輸出位數、容許誤差、以及能不
能改成整數交叉相乘，這三件事要想清楚，再開始寫公式。}

\begin{drills}[label=drills:float-back-sub]
\item 上三角方程組回代時，先解出哪一個未知數？ 最後解出哪一個？
\item $n$ 階回代時，外層迴圈從哪個索引走到哪個索引？
\item 若對角線元素非常接近 $0$，回代過程可能出現什麼問題？
\end{drills}

\addcontentsline{toc}{section}{練習解答}
\section*{練習解答}

各節練習的答案收在這裡。 短題只給能一句話核對的結果。

\paragraph{第~\ref{sec:float-declare-io}~節}
\begin{enumerate}
\item \texttt{scanf} 用 \texttt{\%lf}；\texttt{printf} 用 \texttt{\%f}。
\item 因為 \texttt{9}、\texttt{5} 都是整數常數，\texttt{9 / 5} 會先做
整數除法變成 \texttt{1}，之後乘什麼都救不回來。
\item 預設印 $6$ 位小數。
\end{enumerate}

\paragraph{第~\ref{sec:float-sci-notation}~節}
\begin{enumerate}
\item $1.235 \times 10^4 = 12350$。
\item 因為 \texttt{9 / 5} 的型別是 \texttt{int}，用 \texttt{\%.1f} 印
會把整數位元樣式誤讀成浮點數，印出無意義的內容。
\item \texttt{9.0 / 5} 是常數 \texttt{9.0} 本身就是浮點數；
\texttt{(double)9 / 5} 是用強制轉型把 \texttt{9} 轉成浮點數。
\end{enumerate}

\paragraph{第~\ref{sec:float-cast-assoc}~節}
\begin{enumerate}
\item 因為除法可能先在括號外的整數運算裡做完；正確寫法要把
\texttt{a + b} 整個括起來轉型：\texttt{(double)(a + b) / 2}。
\item \texttt{a - b - c = 5}，\texttt{a - (b - c) = 9}，差在括號位置
改變了運算順序。
\item 由左向右。
\end{enumerate}

\paragraph{第~\ref{sec:float-bit-repr}~節}
\begin{enumerate}
\item \texttt{float} 占 $4$ 個位元組（指數 $8$ 位）；\texttt{double}
占 $8$ 個位元組（指數 $11$ 位）。
\item 不是，\texttt{x} 裡存的是「最接近 $0.1$ 的可表示二進位浮點數」。
\item 靠近 $0$ 的地方比較密。
\end{enumerate}

\paragraph{第~\ref{sec:float-hidden-error}~節}
\begin{enumerate}
\item 三者印出來的字面不同（\texttt{0.1}、\texttt{0.10}、
\texttt{0.10000000000000001}）；要看到真正的近似值要用
\texttt{\%.17f}。
\item 不是同一件事；印出來像整數只是輸出位數把誤差藏起來了。
\end{enumerate}

\paragraph{第~\ref{sec:float-no-equality}~節}
\begin{enumerate}
\item 得到假（\texttt{0}），因為十次累加的誤差讓 \texttt{sum} 跟
\texttt{1.0} 的位元樣式不同。
\item \texttt{fabs(a - b) < eps}，用一個很小的容許誤差判斷是否「足夠
接近」。
\item 比的是記憶體裡的位元樣式是否相同，不是數學上是否相等。
\end{enumerate}

\paragraph{第~\ref{sec:float-abs-rel-error}~節}
\begin{enumerate}
\item 印出假（\texttt{0}）。
\item 判斷「不相等」，因為 $0.5$ 遠大於絕對門檻 \texttt{1e-9}，即使相
對這兩個大數字而言差距很小。
\item \texttt{fabs(a - b) < eps * fabs(a)}，把容許誤差按 $|a|$ 放大。
\end{enumerate}

\paragraph{第~\ref{sec:float-limits}~節}
\begin{enumerate}
\item 量的是 $1$ 附近的間距。
\item 因為 $10^{12}$ 附近相鄰可表示值的間距已經遠大於 \texttt{1e-9}，
不是 \texttt{DBL\_EPSILON} 本身變了，而是間距被指數放大了。
\end{enumerate}

\paragraph{第~\ref{sec:float-inf-nan}~節}
\begin{enumerate}
\item \texttt{1.0 / 0.0} 得到正無窮；\texttt{0.0 / 0.0} 得到 NaN。
\item 假（\texttt{0}）；NaN 跟任何值比較，包括自己，都不會得到「相
等」。
\item 檢查有沒有除以零、對數的真數非正、或平方根的參數為負。
\end{enumerate}

\paragraph{第~\ref{sec:float-money-int}~節}
\begin{enumerate}
\item 因為每一筆金額的小數部分累加時都帶著二進位捨入誤差，加久了可能
連該不該是整數元都說不準。
\item 讓分的部分固定印兩位；少了它，\texttt{5.00} 可能印成
\texttt{5.0}。
\item $199 + 250 + 51 = 500$ 分，印成 \texttt{5.00} 元。
\end{enumerate}

\paragraph{第~\ref{sec:float-cross-mult}~節}
\begin{enumerate}
\item $d_1 \times p_2 > d_2 \times p_1$。
\item $p_1, p_2 > 0$。
\item 因為整個搜尋過程只用整數乘法，完全不出現浮點數，所以沒有浮點誤
差的問題。
\end{enumerate}

\paragraph{第~\ref{sec:float-math-lib}~節}
\begin{enumerate}
\item $\sqrt{3^2+4^2} = 5$。
\item 常要加上 \texttt{-lm} 連結數學函式庫。
\end{enumerate}

\paragraph{第~\ref{sec:float-round-fns}~節}
\begin{enumerate}
\item \texttt{floor} 得 $-2$，\texttt{ceil} 得 $-1$，\texttt{round}
得 $-2$。
\item \texttt{fmod(7.5, 2.0) = 1.50}；概念上是同一件事，只是換成浮點數
版本。
\item 因為 $x$、$y$ 平方後可能先溢成無窮，\texttt{hypot} 內部會先縮放
再算，減少中間溢位。
\end{enumerate}

\paragraph{第~\ref{sec:float-log-digits}~節}
\begin{enumerate}
\item $4$ 位數；用 $\lfloor \log_{10}(2^{10}) \rfloor + 1$ 直接算出來，
不必真的算出 $2^{10}$。
\item 因為 $a^b$ 可能大到任何整數型別都裝不下，而位數本身只是一個很小
的整數。
\item 防止 \texttt{log10} 的結果因浮點誤差落在整數邊界下方，導致少算
一位。
\end{enumerate}

\paragraph{第~\ref{sec:float-polygon-perimeter}~節}
\begin{enumerate}
\item 把相鄰頂點的距離依序用 \texttt{sqrt} 算出來再加總，得到 $16$。
\item 算出 $0$，讓最後一段繞回第一個頂點。
\item 容許誤差比較，不要用 \texttt{==}。
\end{enumerate}

\paragraph{第~\ref{sec:float-arc-area}~節}
\begin{enumerate}
\item $\frac{3}{4} \times 2\pi r$。
\item 用題目給的；不同位數的 $\pi$ 常數會讓答案在小數點後幾位對不起
來。
\item 繩子比某一邊牆還長時，會沿那面牆再繞出一個更小的四分之一圓。
\end{enumerate}

\paragraph{第~\ref{sec:float-line-intersect}~節}
\begin{enumerate}
\item 分母是兩個法向量的行列式 $a_1 b_2 - a_2 b_1$；等於 $0$ 代表兩條
直線平行，沒有交點。
\item 依序求出相鄰兩條直線的交點當頂點，再用 \texttt{\% n} 繞一圈累加
兩點距離。
\end{enumerate}

\paragraph{第~\ref{sec:float-mean-variance}~節}
\begin{enumerate}
\item 平均值 $3$，變異數 $2$。
\item 因為變異數公式需要用到平均值，必須先算完平均值才能算每個值跟它
的差。
\item 兩個很接近的大數相減，可能讓有效位數被嚴重侵蝕，資料量大時甚至
可能算出負的變異數。
\end{enumerate}

\paragraph{第~\ref{sec:float-newton}~節}
\begin{enumerate}
\item $x \leftarrow \frac{1}{2}(x + a/x)$。
\item 做滿規定的 $k$ 次，不要自己改成誤差夠小就提早停。
\item 同時算出多項式的值 $f(x)$ 與導數 $f'(x)$。
\end{enumerate}

\paragraph{第~\ref{sec:float-back-sub}~節}
\begin{enumerate}
\item 先解出最後一個未知數（對角線最下面那個），最後解出第一個。
\item 從 $n-1$ 走到 $0$。
\item 這一層會把上面已經有的捨入誤差放大。
\end{enumerate}

\chapter{函式}
\label{ch:function}

到目前為止，我們寫過的程式幾乎都只有一個 \texttt{main} 函式：所有變數、
所有邏輯都擠在同一層裡。 當程式變長、或某段邏輯需要重複使用時，把它
獨立寫成一個\textbf{函式}（function）會讓程式更容易閱讀、修改、測試。
一個函式就像一個小工具：給它固定形式的輸入（\textbf{參數}，parameter），
它做完一件明確的事，然後用 \texttt{return} 交回一個結果；中間怎麼實作，
呼叫它的人完全不需要知道。

這一章分四段推進：先學會寫有回傳值、無回傳值與多參數函式，並用純函式
對上 \texttt{void} 的設計差異；再從呼叫堆疊解釋整數與陣列參數為什麼行
為不同；接著談函式原型、標頭檔與多檔編譯；最後才看變參 \texttt{printf}
／\texttt{scanf}，以及遞迴的進階用法。 函式指標會在指標那一章正式展開。
堆疊與生命週期會在這一章建立直覺，到指標與遞迴章再串成完整圖像。

\section{定義有回傳值的函式}
\label{sec:func-return-value}

一個有回傳值的函式由幾個部分組成：\textbf{回傳型別}、\textbf{函式名}、
\textbf{參數列表}、\textbf{函式本體}。 本體內用 \texttt{return} 交回一個
值給呼叫者。

\begin{code}{c}{有回傳值的函式}
回傳型別 函式名(參數型別 參數名) {
    ...
    return 型別匹配的值;
}
\end{code}

範例 \thechapter.2 把「計算平方」獨立成 \texttt{square}；雖然
\texttt{n * n} 直接寫在 \texttt{main} 裡也可以，但包成函式之後，只要
看到 \texttt{square(x)} 這個名字，就不用重新想一遍它在做什麼。 呼叫
\texttt{square(n)} 時，\texttt{n} 的值被複製一份給 \texttt{square} 的
參數，計算完畢後回傳值再交給 \texttt{printf}。

\begin{example}{c}{(square.c) 計算平方的函式}
#include <stdio.h>

int square(int n) {
    return n * n;
}

int main(void) {
    int n;
    scanf("%d", &n);
    printf("%d\n", square(n));
}
\end{example}

\iosample{
7 \\
}{
49 \\
}

\subsection{main 函式}

\texttt{main} 也是一個函式，它的特殊之處是作為程式的\textbf{入口點}
（entry point）——作業系統啟動程式時，會呼叫 \texttt{main}。 \texttt{main}
回傳的整數值會交給作業系統，告知程式是否成功執行。 慣例上 \texttt{0}
表示成功，非 \texttt{0} 表示失敗。

\begin{code}{c}{main 的簽名}
int main(void) {
    // 程式邏輯
    return 0;  // 告知作業系統程式成功
}

// 或者接受命令行引數（進階用法，這裡先不詳述）
int main(int argc, char *argv[]) {
    // ...
}
\end{code}

\subsection{判斷閏年}

實際應用中，我們常把「某個明確條件判斷」包成一個函式。 例如判斷是否
為閏年：年份能被 4 整除但不能被 100 整除，或者能被 400 整除。

\begin{example}{c}{(leap-year.c) 判斷閏年的函式}
#include <stdio.h>

// 回傳 1 表示閏年，0 表示平年
int isLeapYear(int year) {
    if (year % 400 == 0) {
        return 1;  // 能被 400 整除，是閏年
    }
    if (year % 100 == 0) {
        return 0;  // 能被 100 整除但不能被 400 整除，是平年
    }
    if (year % 4 == 0) {
        return 1;  // 能被 4 整除但不能被 100 整除，是閏年
    }
    return 0;  // 其他都是平年
}

int main(void) {
    int year;
    scanf("%d", &year);
    printf("%d\n", isLeapYear(year));
}
\end{example}

\iosample{
2024 \\
}{
1 \\
}

輸入 \texttt{2024}，按第三個條件 \texttt{2024 \% 4 == 0}（且不被 100
整除），回傳 \texttt{1}。 輸入 \texttt{2000} 則按第一個條件回傳 \texttt{1}。

\keypoint{\texttt{return} 一執行就立刻結束這次函式呼叫，後面的程式碼
不會再跑。 用連續的 \texttt{if} 搭配提早 \texttt{return}，通常比複雜的
\texttt{if-else} 鏈更平坦好讀。}

\begin{drills}[label=drills:func-return-value]
\item \texttt{square} 函式呼叫 \texttt{square(n)} 時，\texttt{n} 的值
是直接傳進去，還是複製一份給參數？
\item \texttt{isLeapYear(2000)} 會先命中哪一個 \texttt{if}？回傳值是
多少？
\item 把 \texttt{isLeapYear} 裡三個 \texttt{if} 的判斷順序打亂（先判斷
\% 4，再判斷 \% 400），對 2000 這個輸入會不會算錯？為什麼？
\end{drills}

\section{定義無回傳值的函式}
\label{sec:func-void}

\texttt{void} 型別表示函式不回傳任何值。 這類函式通常用來執行某個動作
或修改某些狀態——比如印出結果、修改陣列、更新全域變數。

\begin{code}{c}{無回傳值的函式}
void 函式名(參數型別 參數名) {
    // 執行動作
    // 沒有 return 值（或 return; 後面不跟值）
}
\end{code}

範例~\ref{ex:print-digits} 寫一個 \texttt{void} 函式，把一個整數的各位數字
分別印在不同行。 這個函式不回傳任何計算結果，只是有副作用（把值印出來）。

\begin{example}[label=ex:print-digits]{c}{(print-digits.c) 將各位數印在不同行}
#include <stdio.h>

// 無回傳值的函式：只負責輸出，不計算結果
void printDigits(int n) {
    if (n < 0) n = -n;  // 處理負數

    if (n == 0) {
        printf("0\n");
        return;
    }

    // 遞迴式地從最高位開始印
    if (n >= 10) {
        printDigits(n / 10);
    }
    printf("%d\n", n % 10);
}

int main(void) {
    int num;
    scanf("%d", &num);
    printDigits(num);
}
\end{example}

\iosample{
12345 \\
}{
1 \\
2 \\
3 \\
4 \\
5 \\
}

\texttt{printDigits(12345)} 遞迴呼叫自己，依序把 \texttt{12345/10 = 1234}
傳進去，最後回到 \texttt{1}；再回溯時逐一印出 \texttt{2, 3, 4, 5}。

另一個 \texttt{void} 函式的例子是初始化或填充陣列：

\begin{code}{c}{用 void 函式初始化陣列}
void fillZeros(int arr[], int n) {
    for (int i = 0; i < n; i++) {
        arr[i] = 0;  // 副作用：修改陣列
    }
}
\end{code}

這類函式的副作用（修改陣列）是預期的，所以函式名往往暗示「會改東西」
（如 \texttt{fillZeros}、\texttt{reverse}、\texttt{sort}）。

\keypoint{\texttt{void} 函式沒有回傳值，呼叫它是為了它的副作用（印出
結果、修改陣列、改變全域狀態等）。 在函式名上要講清楚是否會有副作用。}

\begin{drills}[label=drills:func-void]
\item \texttt{printDigits} 沒有 \texttt{return} 值，那它是靠什麼讓
數字出現在螢幕上？
\item \texttt{fillZeros(arr, n)} 呼叫結束後，\texttt{main} 裡原本的
陣列內容會不會被改變？
\item 一個函式如果同時「有副作用」又「有回傳值」，合理嗎？舉一個你覺
得可以接受的例子。
\end{drills}

\section{多個參數的函式}
\label{sec:func-params}

函式可以接多個參數。 呼叫時，每個引數各複製一份值給對應的參數，彼此
獨立。

\begin{code}{c}{多參數函式}
回傳型別 函式名(型別1 參數1, 型別2 參數2, ...) {
    return 根據參數計算的結果;
}
\end{code}

\subsection{計算天數}

一個常見的例子是把「年、月、日」轉換成該年的第幾天。

\begin{example}{c}{(day-of-year.c) 計算日期是一年中的第幾天}
#include <stdio.h>

// 給定年月日，回傳該年第幾天
int dayOfYear(int year, int month, int day) {
    // 各月的天數
    int daysInMonth[] = {31, 28, 31, 30, 31, 30, 31, 31, 30, 31, 30, 31};

    // 處理閏年
    if (year % 400 == 0 || (year % 4 == 0 && year % 100 != 0)) {
        daysInMonth[1] = 29;
    }

    // 加總前面各月的天數
    int result = 0;
    for (int m = 0; m < month - 1; m++) {
        result += daysInMonth[m];
    }

    // 加上這個月的日期
    result += day;
    return result;
}

int main(void) {
    int y, m, d;
    scanf("%d %d %d", &y, &m, &d);
    printf("%d\n", dayOfYear(y, m, d));
}
\end{example}

\iosample{
2024 3 1 \\
}{
61 \\
}

\subsection{印出連續整數序列}

另一個例子是印出從 \texttt{start} 到 \texttt{end} 的連續整數。

\begin{code}{c}{印出一個數字範圍}
void printRange(int start, int end) {
    for (int i = start; i <= end; i++) {
        if (i > start) printf(" ");
        printf("%d", i);
    }
    printf("\n");
}
\end{code}

呼叫 \texttt{printRange(3, 7)} 會印出 \texttt{3 4 5 6 7}。

\subsection{算式判斷}

一些題目要求判斷一個算式是否成立。 例如判斷三個邊長是否能組成直角三角形。

\begin{example}{c}{(right-triangle.c) 判斷是否為直角三角形}
#include <stdio.h>

// 判斷三邊能否組成直角三角形
int isRightTriangle(int a, int b, int c) {
    // 找最長邊
    int max = a;
    if (b > max) max = b;
    if (c > max) max = c;

    // 檢查勾股定理：最長邊的平方 = 其他兩邊平方和
    if (a == max) {
        return a * a == b * b + c * c;
    }
    if (b == max) {
        return b * b == a * a + c * c;
    }
    return c * c == a * a + b * b;
}

int main(void) {
    int a, b, c;
    scanf("%d %d %d", &a, &b, &c);
    printf("%d\n", isRightTriangle(a, b, c));
}
\end{example}

\iosample{
3 4 5 \\
}{
1 \\
}

\subsection{比較判斷}

判斷兩個值的大小或相等關係也很常見。

\begin{code}{c}{比較函式}
int compare(int a, int b) {
    if (a < b) return -1;
    if (a > b) return 1;
    return 0;
}

// 或者直接回傳布林結果
int isGreater(int a, int b) {
    return a > b;
}
\end{code}

\keypoint{多參數函式的每個參數在呼叫時各拿到一份獨立的複本；參
數間互不干擾。}

\begin{drills}[label=drills:func-params]
\item \texttt{dayOfYear(2024, 3, 1)} 為什麼要先處理閏年，才能算出正確
的答案？
\item \texttt{printRange(3, 7)} 印出的最後一個數字是幾？
\item \texttt{isRightTriangle} 為什麼要先找出最長邊，才能檢查勾股定理？
\end{drills}

\section{函式的優點}
\label{sec:func-benefits}

到目前為止，我們已經寫過有回傳值的 \texttt{square}、\texttt{dayOfYear}，
也有只負責輸出的 \texttt{printDigits}。 這一節整理函式為什麼值得寫，以及
\texttt{void} 與有回傳值函式在設計上該怎麼分工。

\subsection{純函式}

一個\textbf{純函式}（pure function）給定相同的輸入，\textbf{永遠會回傳相同
的結果}，而且不會改變程式中其他地方的狀態（如全域變數、輸入輸出等）。
例如 \texttt{square(5)} 在哪裡呼叫、呼叫幾次，都一定回傳 \texttt{25}。

相反地，一個\textbf{有副作用的函式}（function with side effects）不只回傳
結果，還會「順便」改變其他東西——修改陣列、改寫全域變數、印出畫面等。 例如
\texttt{printDigits} 只用於印出；本節後面的 \texttt{incrementAndReturn} 還會改全域變數。

\begin{example}{c}{(pure-vs-impure.c) 純函式 vs 有副作用的函式}
#include <stdio.h>

// 純函式：給定 n，永遠回傳同一個值
int square(int n) {
    return n * n;
}

// 純函式：判斷是否偶數，不改變任何狀態
int isEven(int n) {
    return n % 2 == 0;
}

// 有副作用：修改全域變數
int count = 0;
int incrementAndReturn(int n) {
    count++;  // 副作用：修改全域變數
    return n + 1;
}

int main(void) {
    printf("%d\n", square(5));    // 純函式，多次結果同
    printf("%d\n", square(5));    // 輸出：25, 25

    printf("%d %d\n", incrementAndReturn(5), count);  // 輸出：6 1
    printf("%d %d\n", incrementAndReturn(5), count);  // 輸出：6 2
}
\end{example}

\iosample{
}{
25 \\
25 \\
6 1 \\
6 2 \\
}

同樣呼叫 \texttt{incrementAndReturn(5)} 兩次，第二次執行時全域變數
\texttt{count} 已經從 \texttt{1} 變成 \texttt{2}，結果不再「純粹」。

\subsection{純函式的優點}

\textbf{易於測試}：純函式不依賴外部狀態，測試時只需檢查輸入和輸出是否
對應。 有副作用的函式往往需要檢查「執行前後，哪些全域變數改了」，測試
更複雜。

\textbf{易於組合}：寫 \texttt{minOfThree(a, b, c)} 時，可以安心呼叫
\texttt{minOfTwo} 兩次，因為知道它不會改變 \texttt{a}、\texttt{b}、
\texttt{c} 的值。

\textbf{易於並行執行}：純函式的呼叫之間沒有隱藏的依賴，多個呼叫可以
同時進行而不會衝突。

\begin{example}{c}{(compose-functions.c) 組合純函式}
#include <stdio.h>

int min(int a, int b) {
    return a < b ? a : b;
}

int minOfThree(int a, int b, int c) {
    return min(min(a, b), c);
}

int main(void) {
    printf("%d\n", minOfThree(5, 2, 9));  // 輸出 2
}
\end{example}

\iosample{
}{
2 \\
}

\subsection{模組化與重用}

把邏輯獨立成函式的另一個好處是\textbf{重用}。 一旦寫好 \texttt{gcd}（最大
公因數），就能在求最小公倍數、約分、題目檢查等多個地方使用，不用重複寫。

\begin{example}{c}{(lcm.c) 用輔助函式組合邏輯}
#include <stdio.h>

// 歐幾里得演算法
int gcd(int a, int b) {
    while (b != 0) {
        int t = b;
        b = a % b;
        a = t;
    }
    return a;
}

// 用 gcd 組合出最小公倍數
int lcm(int a, int b) {
    return a / gcd(a, b) * b;
}

int main(void) {
    int n;
    scanf("%d", &n);

    int result;
    scanf("%d", &result);
    for (int i = 1; i < n; i++) {
        int x;
        scanf("%d", &x);
        result = lcm(result, x);
    }

    printf("%d\n", result);
}
\end{example}

\iosample{
3 \\
4 6 8 \\
}{
24 \\
}

\texttt{main} 自己完全沒有寫任何輾轉相除的邏輯，只是把資料交給
\texttt{lcm}；\texttt{lcm} 也不需要知道歐幾里得演算法的細節，只要相信
\texttt{gcd} 正確即可。

\subsection{控制複雜度}

當函式複雜時，可以再把它拆成更小的輔助函式。 這樣每層都只負責一個明確
的任務，整個程式讀起來就像在讀高階的敘述，而非深陷細節。

\subsection{修改與除錯}

有 bug 時，如果邏輯集中在一個函式裡，修改一個地方就能解決所有問題；如果
邏輯散落在 \texttt{main} 的各個角落，bug 修改起來就容易遺漏。

\keypoint{純函式好測試、好組合；有副作用的函式要在名字和註解裡講清楚。
優先寫純函式，只在必要時才接受副作用。}

\begin{drills}[label=drills:func-benefits]
\item \texttt{incrementAndReturn} 是純函式嗎？為什麼？
\item 為什麼「純函式易於組合」？用 \texttt{minOfThree} 呼叫
\texttt{min} 兩次來說明。
\item \texttt{lcm} 函式裡完全沒寫輾轉相除的邏輯，這是靠什麼原則做到的？
\end{drills}

\section{呼叫堆疊與區域變數的生命週期}
\label{sec:func-callstack}

上一節區分純函式與有副作用的函式時，我們強調好的設計盡量不改呼叫者
的狀態。 但就算函式看起來想改參數，在 C 裡通常也改不到 \texttt{main}
裡原本的整數變數。 原因在\textbf{呼叫堆疊}（call stack）：函式裡宣告
的變數（包含參數）都是\textbf{區域變數}（local variable），只存在於這
一次函式呼叫所佔用的 stack frame 裡；函式 \texttt{return} 後，這些變
數占用的記憶體就釋放，名字也隨之失效。 \texttt{main} 裡的 \texttt{int x}
跟某個函式裡的 \texttt{int x} 即使名字相同，也是兩塊完全不相干的記憶體。

\begin{center}
\begin{tikzpicture}
    \node[cell, minimum width=2.8cm, minimum height=0.75cm] (squareTop) {返回位址};
    \node[hlcell, minimum width=2.8cm, minimum height=0.75cm, below=0pt of squareTop] (squareLocals) {區域變數 \texttt{n}};
    \node[cell, minimum width=2.8cm, minimum height=0.75cm, below=0.75cm of squareLocals] (mainTop) {返回位址};
    \node[hlcell, minimum width=2.8cm, minimum height=0.75cm, below=0pt of mainTop] (mainLocals) {區域變數 \texttt{x}, \texttt{y}};

    \node[idxlbl, anchor=east, xshift=-0.15cm] at ($(squareTop.west)!0.5!(squareLocals.west)$) {\texttt{square}};
    \node[idxlbl, anchor=east, xshift=-0.15cm] at ($(mainTop.west)!0.5!(mainLocals.west)$) {\texttt{main}};

    \coordinate (squareArrowStart) at ($(squareLocals.east)+(0.4,0)$);
    \coordinate (squareArrowEnd) at ($(squareLocals.east)+(1.2,0)$);
    \draw[diagarrow] (squareArrowStart) -- (squareArrowEnd);
    \node[idxlbl, anchor=west, align=left, right=0.05cm of squareArrowEnd] {目前執行\\這一層};

    \node[idxlbl, anchor=west, align=left] at ($(mainTop.east)+(1.25,0)$) {暫停，等上面\\做完再回來};
\end{tikzpicture}
\diagcaption{呼叫 \texttt{square(n)} 時，堆疊頂端多一層 \texttt{square}
的 frame（淺灰是這次呼叫專屬的區域變數）；\texttt{square} 返回後，這一
層被整個彈出，\texttt{main} 的 \texttt{x}、\texttt{y} 仍在下面那一層
等待繼續執行。}
\label{fig:ch08-call-stack-square}
\end{center}

片語 \thechapter.16 用佔位符概括「一次函式呼叫」在堆疊上發生的事：編譯
器為這次呼叫準備一層 frame，把引數的值複製進參數，跳去執行函式本體，最
後把回傳值交給呼叫者，並銷毀這一層 frame。

\begin{code}{c}{一次函式呼叫在堆疊上的生命週期（概念）}
// 呼叫者這一層的 frame 仍在堆疊上，只是暫停
結果 = 被呼叫函式(引數的值);   // 引數「值」被複製一份給參數

// 被呼叫函式本體執行期間：
//   1. 在堆疊頂端建立新的 frame
//   2. 參數與函式內宣告的區域變數都活在這層 frame 裡
//   3. return 時銷毀這層 frame，把回傳值留給呼叫者
\end{code}

這個「參數是複製一份值進去」的規則，在 C 裡叫做\textbf{傳值呼叫}
（call by value）。 範例 \thechapter.17 寫了一個看起來像要交換兩個整數
的 \texttt{swap} 函式，但執行後 \texttt{main} 裡的 \texttt{x}、
\texttt{y} 完全沒變。 因為 \texttt{swap} 交換的是它自己 frame 裡那兩
份複本，跟 \texttt{main} 的 \texttt{x}、\texttt{y} 無關。

\begin{example}{c}{(swap-params.c) 在函式裡改參數，改不到呼叫者的變數}
#include <stdio.h>

void swap(int a, int b) {
    int t = a;
    a = b;
    b = t;
}

int main(void) {
    int x = 3, y = 7;
    swap(x, y);
    printf("%d %d\n", x, y);
}
\end{example}

\iosample{}{
3 7 \\
}

用 \texttt{x = 3}、\texttt{y = 7} 走一遍：\texttt{swap(x, y)} 把
\texttt{3} 複製給參數 \texttt{a}、\texttt{7} 複製給 \texttt{b}，在
\texttt{swap} 的 frame 裡 \texttt{a} 與 \texttt{b} 確實互換了，但
\texttt{main} 的 \texttt{x}、\texttt{y} 從頭到尾都在另一層 frame，值
分別仍是 \texttt{3} 和 \texttt{7}。 如果題目真的要在函式裡改到呼叫者
的變數，需要指標（下一章才會正式教）；現階段只要記住：一般整數參
數傳進去的是「值的複本」。

\begin{center}
\begin{tikzpicture}
    \node[hlcell2, minimum width=1.1cm] (mx) {3};
    \node[hlcell2, minimum width=1.1cm, right=0.2cm of mx] (my) {7};
    \node[idxlbl, left=0.3cm of mx] {\texttt{main}};
    \node[hlcell, minimum width=1.1cm, right=1.1cm of my] (pa) {3};
    \node[hlcell, minimum width=1.1cm, right=0.2cm of pa] (pb) {7};
    \node[idxlbl, left=0.3cm of pa] {\texttt{swap}};
    \draw[diagarrow, black!45] (mx.east) -- (pa.west);
    \draw[diagarrow, black!45] (my.east) -- (pb.west);
    \node[idxlbl, align=left, right=0.6cm of pb] {呼叫當下\\各複製一份};
\end{tikzpicture}
\diagcaption{\texttt{main} 的 \texttt{x}、\texttt{y}（深灰）與
\texttt{swap} 的參數 \texttt{a}、\texttt{b}（淺灰）是四個獨立位置；
在 \texttt{swap} 裡改 \texttt{a}、\texttt{b} 不會回寫到 \texttt{x}、
\texttt{y}。}
\label{fig:ch08-stack-frame-scope}
\end{center}

\keypoint{區域變數（含參數）只活在這次函式呼叫的 stack frame 裡；
函式 \texttt{return} 後，這些變數占用的記憶體就釋放，名字也隨之失效。
不同函式裡同名變數互不影響。}

\begin{drills}[label=drills:func-callstack]
\item \texttt{swap(x, y)} 執行完後，\texttt{main} 裡的 \texttt{x}、
\texttt{y} 為什麼沒有互換？
\item 函式裡宣告的區域變數，什麼時候會被釋放？
\item \texttt{main} 裡有一個 \texttt{int x}，某個函式裡也宣告了
\texttt{int x}，這是同一塊記憶體嗎？
\end{drills}

\section{參數傳遞}
\label{sec:func-passing}

上一節的 \texttt{swap} 已經示範：整數參數傳進去的是複本。 這一節把
整數、陣列、二維陣列分開講清楚：整數型別傳的是值的複本；陣列和結構體
有特殊的處理方式。

\subsection{整數參數傳遞}

整數參數遵循\textbf{傳值呼叫}（call by value）：引數的值被複製一份進
stack frame，函式內對參數的修改改不到呼叫者的變數。

\begin{code}{c}{整數傳值示例}
void increment(int n) {
    n++;  // 只改這個 frame 裡的副本
}

int main(void) {
    int x = 5;
    increment(x);
    printf("%d\n", x);  // 仍是 5，不是 6
}
\end{code}

這是一種\textbf{保護}：寫函式的人不用擔心參數會洩漏回去。 但如果題目要
求在函式內真的改到呼叫者的整數，就得用指標（下一章的內容）。

\subsection{一維陣列參數傳遞}

陣列是例外：傳給函式的\textbf{不是整份陣列的複本}，而是陣列的起始位址
（即第一個元素的指標）。 函式對陣列元素做的修改會直接反映到呼叫者的
原陣列上。

\begin{example}{c}{(modify-array.c) 陣列參數會被修改}
#include <stdio.h>

// 計算陣列裡偶數的個數，並把偶數改成 0
int countEvensAndNullify(int arr[], int n) {
    int count = 0;
    for (int i = 0; i < n; i++) {
        if (arr[i] % 2 == 0) {
            count++;
            arr[i] = 0;  // 副作用：直接改陣列
        }
    }
    return count;
}

int main(void) {
    int nums[] = {1, 2, 3, 4, 5, 6};
    int n = 6;

    int evenCount = countEvensAndNullify(nums, n);

    printf("偶數個數：%d\n", evenCount);
    printf("修改後的陣列：");
    for (int i = 0; i < n; i++) {
        if (i) printf(" ");
        printf("%d", nums[i]);
    }
    printf("\n");
}
\end{example}

\iosample{
}{
偶數個數：3 \\
修改後的陣列：1 0 3 0 5 0 \\
}

陣列參數的注意事項：
\begin{itemize}
\item 函式簽名中 \texttt{int arr[]} 看起來像陣列，但實際上只是一個指標。
\item 函式無法直接知道陣列大小，所以通常需要另外傳一個 \texttt{int n}。
\item 函式對陣列的修改是真實的、會反映到呼叫者的原陣列。
\end{itemize}

\subsection{多維陣列參數傳遞}

多維陣列傳進函式時，第一維可以省略，但之後的維度必須明確指定。

\begin{example}{c}{(matrix-operation.c) 二維陣列參數}
#include <stdio.h>

// 把 src 順時針旋轉 90 度，寫進另一塊記憶體 dst
void rotate90(int n, int src[n][n], int dst[n][n]) {
    for (int i = 0; i < n; i++) {
        for (int j = 0; j < n; j++) {
            dst[j][n - 1 - i] = src[i][j];
        }
    }
}

int main(void) {
    int n;
    scanf("%d", &n);

    int a[n][n], b[n][n];
    for (int i = 0; i < n; i++) {
        for (int j = 0; j < n; j++) {
            scanf("%d", &a[i][j]);
        }
    }

    rotate90(n, a, b);

    for (int i = 0; i < n; i++) {
        for (int j = 0; j < n; j++) {
            if (j) printf(" ");
            printf("%d", b[i][j]);
        }
        printf("\n");
    }
}
\end{example}

\iosample{
2 \\
1 2 \\
3 4 \\
}{
3 1 \\
4 2 \\
}

在 \texttt{rotate90(n, a, b)} 中，\texttt{a} 和 \texttt{b} 傳的是位址，
函式裡對 \texttt{dst[i][j]} 的寫入會直接改到 \texttt{main} 的 \texttt{b}
陣列。 這是多維陣列參數和一維陣列參數一樣的行為。

\keypoint{陣列參數傳的是位址而非複本；函式對陣列的修改會直接影響原陣列。
傳入多維陣列時，第一維可省略，但後續維度必須明確。}

\begin{drills}[label=drills:func-passing]
\item 整數參數與陣列參數傳給函式的分別是什麼？
\item \texttt{countEvensAndNullify(nums, n)} 執行完後，\texttt{main}
裡的 \texttt{nums} 陣列內容會不會改變？
\item 為什麼陣列參數通常還要另外傳一個 \texttt{int n}？
\end{drills}

\section{常見陷阱與最佳實踐}
\label{sec:func-pitfalls}

前兩節講過整數參數改不到原變數、陣列參數卻能改到元素、以及函式本身不知道
陣列有多大。 這一節把新手最常踩的三個坑集中列出，方便對照。

\subsection{陷阱 1：整數參數傳值，改不到原變數}

\begin{code}{c}{傳值的限制}
void tryToIncrement(int n) {
    n++;  // 只改副本
}

int main(void) {
    int x = 5;
    tryToIncrement(x);
    printf("%d\n", x);  // 仍是 5
}
\end{code}

這是設計特性，不是 bug。 要改呼叫者的變數，得用指標。

\subsection{陷阱 2：陣列參數與整數參數行為不同}

\begin{code}{c}{對比：整數 vs 陣列}
void modifyInt(int n) {
    n = 999;  // 改不到
}

void modifyArray(int arr[], int n) {
    arr[0] = 999;  // 能改到
}
\end{code}

新手往往忽略這個區別。 同一個道理——陣列傳的是位址——但機制與整數完全不同。

\subsection{陷阱 3：忘記陣列參數需要傳大小}

\begin{code}{c}{陷阱：函式不知道陣列有多大}
void sum(int arr[]) {
    // 無法知道 arr 有多少個元素！
    // int size = sizeof(arr);  // 這會給出指標的大小，不是陣列大小
}
\end{code}

總是一起傳陣列和它的大小：

\begin{code}{c}{正確做法}
int sum(int arr[], int n) {
    int total = 0;
    for (int i = 0; i < n; i++) {
        total += arr[i];
    }
    return total;
}
\end{code}

\subsection{最佳實踐}

\begin{itemize}
\item \textbf{優先寫純函式}：如果不需要副作用，就回傳結果；不要修改參數或
改全域變數。

\item \textbf{講清楚副作用}：若函式必須修改陣列，函式名要暗示這一點
（\texttt{reverse}、\texttt{fillWithZeros}），在註解裡也要說明。

\item \textbf{檢查回傳值}：\texttt{scanf}、\texttt{fopen} 等函式回傳狀態碼，
呼叫時要檢查。

\item \textbf{用函式原型}：大檔案裡，在前面集中寫原型；小檔案可以把定義
寫在前面。

\item \textbf{避免全域變數}：如果可以用參數和回傳值，就不用全域變數；
全域狀態容易導致副作用難預測。

\item \textbf{陣列總是要傳大小}：沒有例外。
\end{itemize}

\begin{drills}[label=drills:func-pitfalls]
\item 為什麼 \texttt{void modifyInt(int n) \{ n = 999; \}} 改不到呼叫者
的變數？
\item \texttt{sizeof(arr)} 在函式參數 \texttt{int arr[]} 裡量到的是什
麼？為什麼不能拿它當陣列長度？
\item 陣列參數需要修改時，函式名該不該暗示這件事？舉一個好名字和一個
壞名字。
\end{drills}

\section{呼叫函式}
\label{sec:func-declare}

到目前為止，函式定義都寫在同一個 \texttt{.c} 檔裡，而且多半放在
\texttt{main} 前面或後面。 程式變大、拆成多個檔案時，編譯器仍然必須在
每一次呼叫前知道函式「長什麼樣子」：它叫什麼名字、接什麼參數、回傳什麼
型別。 這就是 C 語言\textbf{宣告}與\textbf{定義}分離的由來。 有三種常見
方式告訴編譯器這些資訊。

\subsection{定義在前}

最簡單的方式是把函式定義寫在 \texttt{main} 之前。 這樣編譯器掃描到
\texttt{main} 時，已經看過函式的完整定義，知道 \texttt{square(n)} 存在
且回傳 \texttt{int}。

\begin{code}{c}{函式定義在 main 前面}
#include <stdio.h>

// 定義在前
int square(int n) {
    return n * n;
}

// 呼叫在後
int main(void) {
    int x;
    scanf("%d", &x);
    printf("%d\n", square(x));
}
\end{code}

\subsection{函式原型}

當檔案很大、或想把主要邏輯寫在 \texttt{main} 時，可以在前面先寫一個
\textbf{函式原型}（function prototype）——只告訴編譯器函式的簽名，不寫本體。

\begin{code}{c}{用函式原型告訴編譯器函式簽名}
#include <stdio.h>

// 函式原型（只有簽名，沒有本體）
int square(int n);

int main(void) {
    int x;
    scanf("%d", &x);
    printf("%d\n", square(x));
}

// 函式定義放在後面
int square(int n) {
    return n * n;
}
\end{code}

函式原型必須放在函式被呼叫之前，通常在檔案頂部集中寫。 編譯器見到
\texttt{square(x)} 時，會查詢前面宣告的原型，確認簽名正確。

\subsection{標頭檔}

當程式擴展到多個檔案時，用\textbf{標頭檔}（header file）是標準做法。
標頭檔（\texttt{.h}）裡放的是函式原型和常用的型別定義，源檔（\texttt{.c}）
裡放的是實作。 例如標準庫的 \texttt{stdio.h} 就是一個標頭檔，裡面宣告了
\texttt{printf}、\texttt{scanf} 等函式的原型；實際的實作在系統的標準庫
文件裡。

\begin{code}{c}{標頭檔的結構}
// math.h - 公開介面
int square(int n);
int add(int a, int b);

// math.c - 實作
#include "math.h"
int square(int n) {
    return n * n;
}
int add(int a, int b) {
    return a + b;
}

// main.c - 呼叫
#include "math.h"
int main(void) {
    printf("%d\n", square(5));
    printf("%d\n", add(3, 4));
}
\end{code}

編譯時用 \texttt{gcc main.c math.c -o prog} 把兩個源檔一起編譯。

\subsection{檢查函式回傳值}

當函式回傳一個值時，呼叫者常需要檢查這個值是否代表成功或失敗。 很多
標準庫函式約定用 \texttt{0} 代表成功、非 \texttt{0} 代表各種失敗情況。

\begin{example}{c}{(check-return.c) 檢查 scanf 的回傳值}
#include <stdio.h>

int main(void) {
    int x, y;

    // scanf 回傳成功讀入的項目數
    // 成功讀入兩個整數時回傳 2
    if (scanf("%d %d", &x, &y) != 2) {
        printf("輸入錯誤\n");
        return 1;  // 表示程式失敗
    }

    printf("讀到 %d 和 %d\n", x, y);
    return 0;  // 表示程式成功
}
\end{example}

\iosample{
3 5 \\
}{
讀到 3 和 5 \\
}

\texttt{scanf} 回傳它成功讀入的值的個數。 上例中，若輸入 \texttt{3 5}
就回傳 \texttt{2}（兩個都讀到）；若只輸入 \texttt{3}（或輸入非整數的
字串），就回傳較小的數字或 \texttt{0}。 檢查回傳值是寫穩健程式的習慣。

\begin{drills}[label=drills:func-declare]
\item 函式原型跟函式定義差在哪裡？
\item 標頭檔（\texttt{.h}）裡通常放什麼？源檔（\texttt{.c}）裡放什麼？
\item \texttt{scanf("\%d \%d", \&x, \&y)} 若只成功讀到一個整數，回傳
值會是多少？
\end{drills}

\section{不定個數參數：printf 與 scanf}
\label{sec:func-variadic}

上一節的 \texttt{stdio.h} 只宣告了 \texttt{printf}、\texttt{scanf}
等函式的原型；這一節來看其中一個特殊之處：它們的參數個數在編譯時不固定
——程式設計者寫多少個 \texttt{\%d} 就對應多少個值。 這稱為\textbf{變
參函式}（variadic function），C 標準庫通過格式字串（format string）來支
持這個特性。

\subsection{\texttt{printf} 用法}

\begin{code}{c}{printf 的簽名（簡化版）}
int printf(const char *format, ...);
\end{code}

\texttt{const char *format} 是格式字串，\texttt{...} 代表任意多個參數。
\texttt{printf} 解析格式字串，遇到 \texttt{\%} 轉義序列就從剩餘參數裡
取一個值。

\begin{example}{c}{(printf-demo.c) printf 的不同用法}
#include <stdio.h>

int main(void) {
    // 無參數，只印字面量
    printf("Hello\n");

    // 一個參數
    printf("整數：%d\n", 42);

    // 多個參數
    printf("名字：%s，年齡：%d，分數：%.2f\n", "Alice", 25, 95.7);

    // 常見格式指定
    printf("十進制：%d，十六進制：%x，八進制：%o\n", 255, 255, 255);

    // 字元和字串
    printf("字元：%c，字串：%s\n", 'A', "world");

    return 0;
}
\end{example}

\iosample{
}{
Hello \\
整數：42 \\
名字：Alice，年齡：25，分數：95.70 \\
十進制：255，十六進制：ff，八進制：377 \\
字元：A，字串：world \\
}

常見的格式指定符：
\begin{itemize}
\item \texttt{\%d}：十進制整數
\item \texttt{\%x}、\texttt{\%X}：十六進制（小寫/大寫）
\item \texttt{\%o}：八進制
\item \texttt{\%f}、\texttt{\%.2f}：浮點數（\.2f 表示兩位小數）
\item \texttt{\%s}：字串
\item \texttt{\%c}：單個字元
\item \texttt{\%\%}：印出一個 \texttt{\%}
\end{itemize}

\subsection{\texttt{scanf} 用法}

\begin{code}{c}{scanf 的簽名（簡化版）}
int scanf(const char *format, ...);
\end{code}

\texttt{scanf} 的工作類似 \texttt{printf}，但方向相反：它按格式字串
解析\textbf{輸入}，把讀到的值填進後面參數指向的位置（所以參數要用
\texttt{\&} 取位址）。

\begin{example}{c}{(scanf-demo.c) scanf 的基本用法}
#include <stdio.h>

int main(void) {
    int x;
    float f;
    char c;
    char str[100];

    // 讀一個整數
    printf("輸入整數：");
    scanf("%d", &x);
    printf("讀到 %d\n", x);

    // 讀一個浮點數
    printf("輸入浮點數：");
    scanf("%f", &f);
    printf("讀到 %f\n", f);

    // 讀一個字元（跳過空白）
    printf("輸入字元：");
    scanf(" %c", &c);  // 注意前面的空白
    printf("讀到 %c\n", c);

    // 讀一個字串（不含空白）
    printf("輸入字串：");
    scanf("%s", str);
    printf("讀到 %s\n", str);

    return 0;
}
\end{example}

\iosample{
輸入整數：42 \\
}{
讀到 42 \\
}

\texttt{scanf} 的陷阱：
\begin{itemize}
\item \texttt{\%s} 不檢查長度，容易緩衝區溢出。 用 \texttt{\%99s} 限制長
度比較安全。
\item 在格式字串前加空白（如 \texttt{" \%c"}）能跳過前面的換行符；否則讀
字元時容易誤讀輸入的尾端 \texttt{\textbackslash n}。
\item \texttt{scanf} 的回傳值是成功讀入的項目數，可用來檢查輸入是否正確。
\end{itemize}

\keypoint{\texttt{printf} 和 \texttt{scanf} 是變參函式，通過格式字串決定
參數個數；掌握常見格式指定符對正確處理輸入輸出至關重要。}

\begin{drills}[label=drills:func-variadic]
\item \texttt{printf} 怎麼知道 \texttt{...} 裡實際傳了幾個參數？
\item 讀一個字元前，格式字串為什麼常寫成 \texttt{" \%c"} 而不是
\texttt{"\%c"}？
\item \texttt{scanf("\%s", str)} 有什麼安全疑慮？該怎麼改比較安全？
\end{drills}

\section{高級應用與設計模式}
\label{sec:func-advanced}

\subsection{函式指標回調}

函式本身也有位址，所以我們可以宣告指向函式的指標，在執行期才決定要呼
叫哪一個。 這種寫法常用來做回調（callback），或動態切換演算法。 宣告
語法、跟 \texttt{int *f(int)} 的差異、以及用函式指標陣列切換策略，都
會在指標那一章完整展開；這裡先知道這件事存在即可。

\subsection{遞迴最佳實踐}

遞迴優雅但需謹慎。 底線案例和終止條件很關鍵。

\begin{code}{c}{安全的遞迴寫法}
// 好：有明確的終止情況，且每次呼叫都在逼近它
int factorial(int n) {
    if (n <= 1) return 1;  // 終止情況
    return n * factorial(n - 1);
}

// 差：可能無窮遞迴
// int bad_recursion(int n) {
//     return bad_recursion(n + 1);  // 永遠不會停
// }

// 更好：用記憶化避免重複計算
int fib_memo(int n, int memo[]) {
    if (n <= 1) return n;
    if (memo[n] != -1) return memo[n];
    memo[n] = fib_memo(n-1, memo) + fib_memo(n-2, memo);
    return memo[n];
}
\end{code}

\keypoint{遞迴需明確終止。 複雜遞迴用記憶化最佳化。}

\begin{drills}[label=drills:func-advanced]
\item 函式指標跟一般變數指標，指向的東西有什麼不一樣？
\item \texttt{bad\_recursion} 為什麼永遠不會停？
\item \texttt{fib\_memo} 比沒有記憶化的費氏遞迴版本快在哪裡？
\end{drills}

\addcontentsline{toc}{section}{練習解答}
\section*{練習解答}

各節練習的答案收在這裡。 短題只給能一句話核對的結果。

\paragraph{第~\ref{sec:func-return-value}~節}
\begin{enumerate}
\item 複製一份給參數；\texttt{square} 裡的 \texttt{n} 是呼叫端那個值
的複本，跟呼叫端的變數無關。
\item 先命中 \texttt{year \% 400 == 0}（2000 能被 400 整除），回傳
\texttt{1}。
\item 不會算錯。 2000 能被 4 整除也能被 400 整除，不管先判斷哪一個，
最後都會落在「是閏年」的分支；但若判斷順序寫反、又忘記排除「被 100
整除但不被 400 整除」的情況，其他年份（如 1900）就會算錯。
\end{enumerate}

\paragraph{第~\ref{sec:func-void}~節}
\begin{enumerate}
\item 靠函式內部呼叫 \texttt{printf} 的副作用；它沒有 \texttt{return}
任何值。
\item 會。 陣列參數傳的是位址，\texttt{fillZeros} 對 \texttt{arr[i]}
的寫入會直接反映到呼叫者的原陣列。
\item 合理，例如 \texttt{scanf} 本身：它把值寫進參數指到的位址（副作
用），也回傳成功讀入的個數（回傳值）。
\end{enumerate}

\paragraph{第~\ref{sec:func-params}~節}
\begin{enumerate}
\item 因為閏年的 2 月有 29 天，會影響 3 月以後每個月累加的天數基準。
\item 7。
\item 因為勾股定理要用「最長邊的平方」對「另外兩邊平方和」比較，先找
出哪一邊最長，才知道該把哪個變數放在等號左邊。
\end{enumerate}

\paragraph{第~\ref{sec:func-benefits}~節}
\begin{enumerate}
\item 不是。 它會修改全域變數 \texttt{count}，同樣的參數 \texttt{n}
在不同次呼叫會讓程式的其他狀態（\texttt{count}）改變。
\item 因為 \texttt{min} 不會偷改 \texttt{a}、\texttt{b}，所以
\texttt{minOfThree} 可以放心呼叫兩次 \texttt{min} 而不用擔心中途結果
被污染。
\item 抽象／重用原則：\texttt{lcm} 只負責呼叫 \texttt{gcd} 並組合結
果，相信 \texttt{gcd} 已經正確實作，不需要重新展開它的邏輯。
\end{enumerate}

\paragraph{第~\ref{sec:func-callstack}~節}
\begin{enumerate}
\item 因為整數參數是傳值呼叫：\texttt{swap} 交換的是它自己 stack
frame 裡的複本 \texttt{a}、\texttt{b}，跟 \texttt{main} 的 \texttt{x}、
\texttt{y} 是不同的記憶體。
\item 函式 \texttt{return} 的時候，它的 stack frame 被整個銷毀。
\item 不是。 兩個 \texttt{x} 分別在各自函式的 stack frame 裡，是兩塊
互不相干的記憶體，只是剛好同名。
\end{enumerate}

\paragraph{第~\ref{sec:func-passing}~節}
\begin{enumerate}
\item 整數參數傳的是值的複本；陣列參數傳的是陣列起始位址（第一個元素
的指標）。
\item 會改變；因為函式拿到的是原陣列的位址，對 \texttt{arr[i]} 的寫入
就是在寫原陣列。
\item 因為陣列參數退化成指標後，函式內部無法從 \texttt{arr} 本身得知
陣列有幾個元素，必須額外告知。
\end{enumerate}

\paragraph{第~\ref{sec:func-pitfalls}~節}
\begin{enumerate}
\item 因為 \texttt{n} 是整數參數，只拿到呼叫端那個值的複本；改
\texttt{n} 只改到這份複本，不會回寫到呼叫端的變數。
\item 量到的是指標的大小（通常是 4 或 8 個位元組），不是陣列實際的元
素個數，因為陣列參數在函式簽名裡已經退化成指標。
\item 該。 例如 \texttt{fillWithZeros}／\texttt{reverse} 是好名字，
清楚暗示會修改參數；\texttt{getMax} 這種名字如果偷偷改了陣列，就會誤
導呼叫者。
\end{enumerate}

\paragraph{第~\ref{sec:func-declare}~節}
\begin{enumerate}
\item 函式原型只宣告簽名（名字、參數型別、回傳型別），沒有本體；函式
定義則包含完整的程式邏輯本體。
\item 標頭檔通常放函式原型和型別定義；源檔放實際的函式實作。
\item 1（只成功讀到 \texttt{x} 那一個整數）。
\end{enumerate}

\paragraph{第~\ref{sec:func-variadic}~節}
\begin{enumerate}
\item 靠格式字串：\texttt{printf} 解析格式字串裡每個 \texttt{\%} 轉義
序列，依序從 \texttt{...} 裡取一個對應的參數，並不會另外檢查參數個數
是否吻合。
\item 因為前面的輸入（例如上一個 \texttt{scanf("\%d", ...)}）常留下一
個換行符在輸入緩衝區裡；格式字串開頭的空白會先跳過所有空白字元，避免
直接讀到那個殘留的換行。
\item \texttt{\%s} 不檢查目的地緩衝區的長度，輸入太長會造成緩衝區溢
出；應改用 \texttt{\%99s} 這種帶寬度限制的寫法。
\end{enumerate}

\paragraph{第~\ref{sec:func-advanced}~節}
\begin{enumerate}
\item 函式指標指向一段可執行的程式碼（某個函式），一般變數指標指向資
料（例如一個整數或陣列元素）。
\item 因為它沒有終止情況：每次呼叫都把 \texttt{n} 加大再繼續呼叫自
己，沒有任何一步會停下來回傳。
\item 快在避免重複計算：一般遞迴費氏數列會重複展開同樣的子問題很多
次，記憶化把算過的結果存進 \texttt{memo}，之後遇到同樣的 \texttt{n}
直接查表回傳。
\end{enumerate}


\chapter{指標}
\label{ch:pointer}

到目前為止，我們操作的變數多是值本身：把 $7$ 存進 \texttt{int x}，
\texttt{x} 就代表 $7$。 但每個變數在記憶體裡都占一個\textbf{位址}。
\textbf{指標}（pointer）是一種變數，存的是另一個物件的位址；用
\texttt{\&} 取位址、用 \texttt{*} 沿位址讀寫，是這一章的核心操作。

這一章分五段推進：先建立左值／右值與 wild pointer 的防禦習慣，再整
合指標參數、純函式與副作用的取捨、\texttt{const} 保護、回傳指標的生
命週期；接著釐清陣列與指標是不同東西、退化與指標算術如何運作；然後
看多維攤平、列指標與函式指標；最後集中越界、型別不匹配與懸空指標等
陷阱。 函式章已預告傳值改不到原變數；這一章用位址把那一限制補完整。

\section{左值、右值與指標基礎}
\label{sec:ptr-basics}

左值（l-value）與右值（r-value）是理解指標的起點。 寫下
\texttt{int x = 7} 時，編譯器會在記憶體裡分配一塊剛好放得下
\texttt{int} 的空間，並把名字 \texttt{x} 綁到這塊空間的位址。 這個
名字代表一塊\textbf{有固定位址}的記憶體物件，可以出現在賦值左邊，
也可以對它取 \texttt{\&}；這類物件稱為左值，l-value 的 l 來自
locator，意思是它能指出物件住在哪。 字面常數 \texttt{42}、算式
\texttt{x + 1} 的結果則只是\textbf{內容值}，沒有自己的持久位址，稱
為右值。 \texttt{\&} 只能加在左值前面：我們取的是這個物件住在哪。

指標變數本身也是左值：編譯器同樣給它分配一塊空間，只是這塊空間裡
存的內容剛好是另一個物件的位址。 所以指標有兩層：\texttt{p} 自己
住在某個位址，\texttt{p} 裡又記著別人的位址。

\subsection{宣告}

指標變數的宣告格式是 \textbf{型別 *變數名}。 例如 \texttt{int *p}
代表 p 是一個指向 \texttt{int} 的指標。 星號 \texttt{*} 在宣告時是
\textbf{型別修飾符}，讀法是指向 \ldots 的指標。

\begin{code}{c}{指標宣告的多種形式}
int *p;           // p 是指向 int 的指標
double *q;        // q 是指向 double 的指標
int *p1, *p2;     // 多個指標：都要加 *
int *p3, x;       // 陷阱：p3 是指標，但 x 只是普通 int
\end{code}

宣告指標後必須立即指派一個有效的位址，否則對它解參考是未定義行為。

\subsection{指定（賦值）位址}

將一個變數或陣列元素的位址指派給指標，使用 \texttt{\&} 運算子（取
位址）。

\begin{code}{c}{用 \& 取位址，指派給指標}
int x = 42;
int *p = &x;      // p 現在指向 x（存著 x 的位址）

int arr[5];
int *pa = &arr[0];  // 指向陣列第一個元素
int *pa2 = arr;     // 陣列名在此當成首元素位址（退化，後面會細講）
\end{code}

\subsection{取值（解參考）}

用 \texttt{*} 運算子沿著指標指向的位址讀出（或修改）那個位置的值。
這個動作叫做\textbf{解參考}（dereference）。

\begin{code}{c}{用 * 取值}
int x = 42;
int *p = &x;
printf("%d\n", *p);   // 讀：印出 42
*p = 100;             // 寫：現在 x 變成 100
\end{code}

\texttt{*p} 與 \texttt{x} 代表同一個整數，只是走不同的路徑：
\begin{itemize}
  \item \texttt{x} 是直接存取變數
  \item \texttt{*p} 是沿著指標指向的位址讀取
\end{itemize}

\begin{center}
\begin{tikzpicture}
    \node[cell] (xcell) {42};
    \node[idxlbl, below=0.15cm of xcell] {\texttt{x}};

    \node[cell, right=1.3cm of xcell] (pcell) {};
    \node[idxlbl, font=\scriptsize] at (pcell.center) {$\rightarrow$~\texttt{x}};
    \node[idxlbl, below=0.15cm of pcell] {\texttt{p}};

    \draw[diagarrow] (pcell.west) .. controls ($(pcell.west)!0.5!(xcell.east)+(0,0.35)$) .. (xcell.east);
\end{tikzpicture}
\diagcaption{變數 \texttt{x} 存值 $42$；指標 \texttt{p} 存的是 \texttt{x}
的位址（箭頭表示指向）。 \texttt{*p} 沿箭頭讀到的就是 $42$。}
\label{fig:ch09-pointer-basic}
\end{center}

\subsection{Wild pointer 與 NULL}

若宣告指標卻不初始化，它會含當時那塊記憶體碰巧留下的垃圾值。 這種
指標叫做 \textbf{wild pointer}（野指標）。 對它做 \texttt{*} 等於往
一個隨機位址讀寫：輕則立刻 Segmentation Fault，重則默默改掉別的變
數，之後的 bug 極難追。 宣告當下就要給它一個有效位址，或先設成
\texttt{NULL}。

\texttt{NULL} 是刻意表示目前不指向任何有效物件的哨兵值。 對
\texttt{NULL} 做 \texttt{*} 幾乎一定立刻 crash（Segmentation
Fault），這反而比較好查：錯就停在解參考那一步，不會默默改到一塊看
起來還能跑的記憶體。 因此在使用指標前，一定要確認它不是
\texttt{NULL}。

\begin{code}{c}{安全地使用指標}
int *p = NULL;           // 宣告時初始化成 NULL
int x = 5;
p = &x;                  // 現在指向有效變數

if (p != NULL) {
    printf("%d\n", *p);  // 安全：先檢查再使用
}
\end{code}

習慣上，宣告指標時立刻初始化成 \texttt{NULL} 或指向某個有效物件，
避免指標含垃圾值。

\subsection{取位址運算子 \&}

\texttt{\&} 運算子得到變數在記憶體中的位址。 對於陣列，陣列名本身
在很多上下文中就代表起始位址。

\begin{code}{c}{\& 運算子的使用}
int x = 42;
printf("%p\n", &x);      // 印出 x 的位址（格式 %p）

int arr[5] = {1, 2, 3, 4, 5};
printf("%p\n", arr);      // 陣列名 = 首元素位址
printf("%p\n", &arr[0]);  // 明確取首元素位址（相同）
printf("%p\n", &arr[2]);  // 取第三個元素的位址
\end{code}

範例 \thechapter.6 把取位址、解參考、透過指標改寫原變數串在一起。
跑起來會看到 \texttt{\&x} 與 \texttt{p} 印出同一個位址；對
\texttt{*p} 賦值之後，\texttt{x} 的值跟著變成 $100$。

\begin{example}{c}{(pointer-basics.c) 指標的基本操作}
#include <stdio.h>

int main(void) {
    int x = 42;
    int *p = &x;           // p 指向 x

    printf("x 的值：%d\n", x);
    printf("*p 的值：%d\n", *p);    // 沿 p 讀 x
    printf("x 的位址：%p\n", &x);
    printf("p 存的位址：%p\n", p);

    *p = 100;              // 透過指標改 x
    printf("改變後 x：%d\n", x);

    return 0;
}
\end{example}

\iosample{}{
x 的值：42 \\
*p 的值：42 \\
x 的位址：0x7ffc8b7c4abc \\
p 存的位址：0x7ffc8b7c4abc \\
改變後 x：100 \\
}

\keypoint{\texttt{\&} 從變數得到位址，\texttt{*} 從位址得到值；指標
變數本身存的是位址。 型別必須一致：\texttt{int *p} 只能指向
\texttt{int}。}

\begin{drills}[label=drills:ptr-basics]
\item 字面常數 \texttt{42} 跟運算式 \texttt{x + 1}，是左值還是右值？
\texttt{\&} 運算子只能加在哪一種前面？
\item \texttt{int *p3, x;} 這樣宣告，\texttt{x} 是指標還是普通
\texttt{int}？
\item 為什麼指標宣告後不能放著不管，要立刻指派有效位址或
\texttt{NULL}？
\item 對 \texttt{NULL} 指標解參考通常會發生什麼事？為什麼這種
crash 反而「比較好查」？
\end{drills}

\section{指標參數與多重回傳}
\label{sec:ptr-params}

函式那一章說過：一般變數當參數傳進函式時，傳的是\textbf{值的複製}，
函式裡改參數不會影響呼叫者原本的變數。 片語 \thechapter.2 已經讓指
標握著別人的位址；把這個位址傳進函式，函式就能沿著位址改寫呼叫者
那邊的記憶體。

\subsection{交換變數值}

\begin{code}{c}{用指標參數交換兩個整數}
void swap(int *pa, int *pb) {
    int temp = *pa;
    *pa = *pb;
    *pb = temp;
}

int main(void) {
    int a = 3, b = 7;
    swap(&a, &b);          // 傳位址
    printf("%d %d\n", a, b);  // 輸出 7 3
}
\end{code}

\texttt{swap(\&a, \&b)} 傳進去的是 \texttt{a}、\texttt{b} 的位址。
函式裡的 \texttt{*pa}、\texttt{*pb} 實際上在改呼叫者那邊的
\texttt{a}、\texttt{b}。 對照函式章那個改不到原變數的 \texttt{swap}：
那一版交換的是自己 stack frame 裡的複本；這一版沿位址改的是呼叫者
的那兩格。

\begin{center}
\begin{tikzpicture}
    \node[hlcell2, minimum width=1.1cm] (a1) {3};
    \node[idxlbl, below=0.15cm of a1] {\texttt{a}};
    \node[hlcell2, minimum width=1.1cm, right=0.4cm of a1] (b1) {7};
    \node[idxlbl, below=0.15cm of b1] {\texttt{b}};

    \node[hlcell2, minimum width=1.1cm, below=1.3cm of a1] (a2) {7};
    \node[idxlbl, below=0.15cm of a2] {\texttt{a}};
    \node[hlcell2, minimum width=1.1cm, right=0.4cm of a2] (b2) {3};
    \node[idxlbl, below=0.15cm of b2] {\texttt{b}};

    \draw[diagarrow] ($(a1)!0.5!(b1)$ |- a1.south) -- ($(a2)!0.5!(b2)$ |- a2.north);
    \coordinate (rowmid) at ($(a1)!0.5!(a2)$);
    \node[idxlbl, right=1.7cm of rowmid] {交換後};
\end{tikzpicture}
\diagcaption{函式沿位址改寫兩格記憶體，交換完成。}
\label{fig:ch09-pointer-swap}
\end{center}

\begin{example}{c}{(swap-with-pointers.c) 用指標交換兩個整數}
#include <stdio.h>

void swap(int *pa, int *pb) {
    int temp = *pa;
    *pa = *pb;
    *pb = temp;
}

int main(void) {
    int a, b;
    scanf("%d %d", &a, &b);
    swap(&a, &b);
    printf("%d %d\n", a, b);
}
\end{example}

\iosample{
3 7 \\
}{
7 3 \\
}

用輸入 $3$ 與 $7$ 走一遍：\texttt{swap} 收到的是這兩個變數的位址，
\texttt{temp} 先接住 $3$，再把 $7$ 寫進 \texttt{a}、把 $3$ 寫進
\texttt{b}。 \texttt{main} 印出的就是 $7$ $3$。

\subsection{輸出參數}

除了交換，指標參數還常拿來當\textbf{輸出參數}（output parameter）：
函式除了 \texttt{return} 一個值，還透過指標把額外結果寫回給呼叫者。

\begin{example}{c}{(output-parameter.c) 用指標參數回傳最大值與索引}
#include <stdio.h>

void findMax(int a[], int n, int *maxVal, int *maxIdx) {
    *maxVal = a[0];
    *maxIdx = 0;
    for (int i = 1; i < n; i++) {
        if (a[i] > *maxVal) {
            *maxVal = a[i];
            *maxIdx = i;
        }
    }
}

int main(void) {
    int a[1005], n;
    scanf("%d", &n);
    for (int i = 0; i < n; i++) {
        scanf("%d", &a[i]);
    }

    int maxVal, maxIdx;
    findMax(a, n, &maxVal, &maxIdx);
    printf("%d %d\n", maxVal, maxIdx);
}
\end{example}

\iosample{
5 \\
1 4 8 9 10 \\
}{
10 4 \\
}

呼叫端準備 \texttt{maxVal}、\texttt{maxIdx} 兩個變數，把位址傳給函
式；函式執行完後，這兩個變數裡就寫好了結果。 這跟
\texttt{scanf("\%d", \&n)} 是同一個模式：我們請函式把答案寫進我們
準備好的格子。 測資裡最大值是 $10$、落在索引 $4$，所以印出
\texttt{10 4}。

\subsection{純函式與副作用的取捨}

指標參數讓函式可以改寫呼叫者的變數，這正是一種副作用。 函式那一章
已經把\textbf{純函式}（pure function）講清楚了：給定相同輸入永遠得
到相同輸出，而且不改外部狀態。 純函式好測試、好組合、也好並行；指
標把改寫外部狀態這條路打開了，所以更要克制。 需要改寫呼叫者、或一
次交出兩個結果時，才用指標當輸出參數；其餘情況仍優先寫只靠
\texttt{return} 交回結果的純函式。

\keypoint{指標參數用於改寫呼叫者的變數，或一次回傳多個結果。 呼叫
時必須傳 \texttt{\&變數}。 能寫純函式就寫純函式，只在必要時才用指
標製造副作用。}

\begin{drills}[label=drills:ptr-params]
\item \texttt{swap(\&a, \&b)} 裡的 \texttt{\&a}、\texttt{\&b} 在傳什
麼？ 如果寫成 \texttt{swap(a, b)}（漏掉 \texttt{\&}）會發生什麼事？
\item \texttt{findMax} 為什麼要用兩個指標參數 \texttt{maxVal}、
\texttt{maxIdx}，而不是只 \texttt{return} 一個值？
\item 純函式跟用指標製造副作用的函式，什麼情況該選前者、什麼情況才
用後者？
\end{drills}

\section{用 \texttt{const} 保護唯讀參數}
\label{sec:ptr-const}

\texttt{const} 加在指標宣告上，可以保護兩件不同的事：指向的資料，
或指標自己存的位址。 兩種寫法長得很像，語意差很多，讀的時候從右往
左拆。

\texttt{const int *p}：\texttt{p} 是指向 \texttt{const int} 的指標。
可以改 \texttt{p} 指向哪裡，不能透過 \texttt{p} 改那格的值。
\texttt{int * const p}：\texttt{p} 是指向 \texttt{int} 的常數指標。
可以改 \texttt{*p}，不能讓 \texttt{p} 改指向別的物件。

\begin{code}{c}{const int * 與 int * const}
int x = 10, y = 20;

const int *p_data = &x;   // 保護資料：不能寫 *p_data
p_data = &y;              // 可以改指向
// *p_data = 0;           // 編譯錯誤

int * const p_addr = &x;  // 保護指標本身：不能改 p_addr
*p_addr = 30;             // 可以改資料
// p_addr = &y;           // 編譯錯誤
\end{code}

範例 \thechapter.11 把這兩種組合都跑一遍：先讓指向資料的指標改指向
\texttt{y}，再讓常數指標把 \texttt{x} 改成 $30$。

\begin{example}{c}{(const-pointer-kinds.c) 兩種 const 指標組合}
#include <stdio.h>

int main(void) {
    int x = 10;
    int y = 20;

    const int *p_data = &x;   // 可改指向，不可改資料
    p_data = &y;
    printf("%d\n", *p_data);

    int * const p_addr = &x;  // 可改資料，不可改指向
    *p_addr = 30;
    printf("%d\n", x);

    return 0;
}
\end{example}

\iosample{}{
20 \\
30 \\
}

還有第三種寫法 \texttt{const int * const p}，資料與指向都不能改。
那一種、以及 \texttt{static} 等修飾詞，會在宣告修飾那一章系統整理。
\texttt{*}、\texttt{[]}、\texttt{()} 疊在一起的讀法，也在那一章後段練。
這一章先把最常用的兩種 \texttt{const} 指標分開即可。

函式若只需要\textbf{讀}陣列、不需要改元素，就在參數上加
\texttt{const}，向編譯器保證本函式不會透過這個指標寫入。 這對唯讀
走訪、加總、查找特別有用，也能避免手滑改到呼叫者的資料。 參數寫成
\texttt{const int a[]} 與 \texttt{const int *a} 在這裡語意相近，都
是保護資料的那種組合。

\begin{code}{c}{唯讀陣列參數}
int sumReadonly(const int a[], int n) {
    int total = 0;
    for (int i = 0; i < n; i++) {
        total += a[i];
    }
    // a[i] = 0;  // 編譯錯誤：不能改 const 資料
    return total;
}
\end{code}

\begin{example}{c}{(sum-readonly.c) 用 const 保護唯讀陣列參數}
#include <stdio.h>

int sumReadonly(const int a[], int n) {
    int total = 0;
    for (int i = 0; i < n; i++) {
        total += a[i];
    }
    return total;
}

int main(void) {
    int a[1005], n;
    scanf("%d", &n);
    for (int i = 0; i < n; i++) {
        scanf("%d", &a[i]);
    }
    printf("%d\n", sumReadonly(a, n));
}
\end{example}

\iosample{
5 \\
1 4 8 9 10 \\
}{
32 \\
}

\texttt{sumReadonly} 只把元素加起來，完全不需要改陣列，所以參數加
\texttt{const}。 測資 $1+4+8+9+10 = 32$。 需要改陣列時才拿掉
\texttt{const}，像前面的 \texttt{swap} 就不能加。

\keypoint{\texttt{const int *p} 保護指向的資料；\texttt{int *
const p} 保護指標本身。 唯讀走訪的函式參數用前一種，編譯器會擋掉
意外的賦值。}

\begin{drills}[label=drills:ptr-const]
\item \texttt{const int *p} 跟 \texttt{int * const p}，哪一個不能改
\texttt{p} 指向哪裡？哪一個不能改 \texttt{*p} 的值？
\item 一個只需要讀陣列、不需要改元素的函式，參數該不該加
\texttt{const}？ 這樣做編譯器會幫忙擋掉什麼？
\item \texttt{const int * const p} 同時保護了什麼？
\end{drills}

\section{回傳指標與生命週期}
\label{sec:ptr-return-lifetime}

上一節用指標參數把結果\textbf{寫回}呼叫者；這一節談函式\textbf{回
傳}指標本身。 回傳指標安全的前提是：那個物件的生命週期必須超過這
次函式呼叫。 區域變數活在這次呼叫的 stack frame 裡，函式
\texttt{return} 後這一層就被銷毀，位址立刻失效。

\subsection{安全的回傳}

範例 \thechapter.14 示範兩種安全情況：回傳全域陣列的位址，以及回傳
呼叫者傳進來的陣列裡某個元素的位址。 這兩塊記憶體在函式結束後都還
在。

\begin{example}{c}{(return-pointer-safe.c) 安全的指標回傳}
#include <stdio.h>

int global_arr[10];

int *getGlobalArray() {
    return global_arr;     // 全域物件，函式結束後仍在
}

int *findInArray(int arr[], int n, int target) {
    for (int i = 0; i < n; i++) {
        if (arr[i] == target) {
            return &arr[i];  // 指向呼叫者的陣列，安全
        }
    }
    return NULL;
}

int main(void) {
    int nums[] = {1, 2, 3, 4, 5};
    int *found = findInArray(nums, 5, 3);
    if (found != NULL) {
        printf("找到：%d\n", *found);
    }
}
\end{example}

\iosample{}{
找到：3 \\
}

\texttt{findInArray} 在 \texttt{nums} 裡找到 $3$，回傳
\texttt{\&nums[2]}。 \texttt{nums} 是 \texttt{main} 的陣列，函式結
束後仍在，所以 \texttt{*found} 讀到 $3$。 找不到時回傳
\texttt{NULL}，呼叫端先檢查再解參考。

\subsection{危險的做法}

\begin{code}{c}{危險：回傳局部變數指標}
int *badFunction() {
    int x = 42;
    return &x;      // x 是區域變數，函式結束後會被銷毀
}
\end{code}

\texttt{badFunction} 回傳的位址指向一塊已經歸還給系統的記憶體。 之
後讀這個指標是未定義行為：有時碰巧還看得到 $42$，有時會讀到被蓋掉
的垃圾，有時直接 crash。 動態配置的記憶體生命週期由我們自己用
\texttt{malloc}／\texttt{free} 決定，可以活過函式呼叫，那一種會在
高等技術那一章展開；現階段能安全回傳的，就是參數陣列與全域變數。

\keypoint{回傳指標安全的條件是：指向的物件活得比這次函式呼叫久。
參數陣列、全域變數可以；區域變數不行。}

\begin{drills}[label=drills:ptr-return-lifetime]
\item \texttt{badFunction} 回傳 \texttt{\&x} 為什麼是危險的？
\texttt{x} 是哪一種生命週期的變數？
\item 回傳指標安全的條件是什麼？參數陣列和全域變數為什麼符合這個條
件？
\item \texttt{findInArray} 找不到 \texttt{target} 時回傳什麼？呼叫端
在解參考回傳值之前應該先做什麼？
\end{drills}

\section{陣列與指標：兩種不同的東西}
\label{sec:ptr-array-vs-pointer}

指標參數常常搭配陣列一起出現，\texttt{int *p = a} 看起來也很順，
所以初學者很容易把兩者混成同一個東西。 編譯器對它們的處理差很多。
\texttt{int a[5]} 在宣告當下就配置連續 $5$ 個 \texttt{int} 的記憶
體，\texttt{a} 這個名字在編譯期就綁在這塊空間的起始位址上。
\texttt{int *p} 只是\textbf{一個}存位址的變數，本身不占那 $5$ 格；
執行時 CPU 得先把 \texttt{p} 裡的位址載出來，再拿那個位址當起點。

最直覺的證據是 \texttt{sizeof}。 範例 \thechapter.16 在同一個程式
裡印出陣列與指標的大小。

\begin{example}{c}{(array-vs-pointer.c) 陣列與指標的 sizeof 不同}
#include <stdio.h>

int main(void) {
    int a[5];
    int *p = a;

    printf("%zu %zu\n", sizeof(a), sizeof(p));
    return 0;
}
\end{example}

\iosample{}{
20 8 \\
}

在常見的 $64$ 位元環境上，\texttt{sizeof(a)} 是 $5$ 個
\texttt{int} 的總位元組（例如 $20$），\texttt{sizeof(p)} 只是指標
變數本身的大小（例如 $8$）。 \texttt{p = a} 之後兩者都能用來走訪
同一排元素，型別系統仍把它們當成不同東西：陣列是一塊固定大小的記
憶體，指標是一個存放位址的變數。

\keypoint{陣列是連續配置的固定大小記憶體；指標是存放位址的變數。
語法上常一起出現，\texttt{sizeof} 會立刻拆穿兩者的差別。}

\begin{drills}[label=drills:ptr-array-vs-pointer]
\item \texttt{int a[5]; int *p = a;} 之後，\texttt{sizeof(a)} 跟
\texttt{sizeof(p)} 會一樣嗎？
\item 陣列在宣告當下做了什麼、指標又做了什麼，導致上一題兩者的
\texttt{sizeof} 不同？
\end{drills}

\section{退化與指標算術}
\label{sec:ptr-decay-arith}

上一節強調陣列與指標是不同型別；但在大多數\textbf{運算式}裡，陣列
名會暫時\textbf{退化}（decay）成指向第一個元素的指標。 傳進函式當
參數時也會退化，所以 \texttt{void f(int a[])} 與
\texttt{void f(int *a)} 對編譯器來說幾乎是同一種參數。 這讓我們能
用指標語法走訪陣列。

\texttt{sizeof(a)} 與 \texttt{\&a} 是少數不會退化的情況：前者量的
是整塊陣列，後者得到的是指向整塊陣列的指標，型別是
\texttt{int (*)[N]}，跟指向首元素的 \texttt{int *} 不同。 下一節談
多維陣列時還會再碰到這種列指標。

\begin{code}{c}{陣列名退化為指標}
int a[5] = {1, 4, 8, 9, 10};

int *p = a;           // a 退化成 &a[0]
// a[i]   等同 *(a + i)
// &a[i]  等同 (a + i)
\end{code}

\begin{example}{c}{(walk-array-pointer.c) 用指標走訪陣列並累加}
#include <stdio.h>

int main(void) {
    int a[1005], n;
    scanf("%d", &n);
    for (int i = 0; i < n; i++) {
        scanf("%d", &a[i]);
    }

    int sum = 0;
    int *p = a;              // p 指向 a[0]
    for (int i = 0; i < n; i++) {
        sum += *p;
        p++;                 // 移到下一個元素
    }
    printf("%d\n", sum);
}
\end{example}

\iosample{
5 \\
1 4 8 9 10 \\
}{
32 \\
}

\texttt{p++} 移動一個 \texttt{int} 的寬度（通常 $4$ bytes），不是
$1$ byte。 這種寫法與 \texttt{for (int i = 0; i < n; i++) sum +=
a[i];} 完全等價，測資加總仍是 $32$。

\subsection{指標算術}

\texttt{p + 1} 的意義是：記憶體位址增加 \texttt{sizeof(*p)} 個位元
組，也就是往後一個元素。 編譯器依指標的型別決定跨多遠，我們不用自
己乘上元素大小。

指標可以做有限的算術：
\begin{itemize}
  \item \texttt{p + n}：移動 n 個元素的距離
  \item \texttt{p - q}（兩個指標都指向同一個陣列時）：算出相隔幾
        個元素
\end{itemize}

\begin{code}{c}{指標算術示例}
int a[5] = {10, 20, 30, 40, 50};
int *p = &a[1];
int *q = &a[3];

printf("%d\n", *(p + 2));    // 等同 a[3]，印出 40
printf("%d\n", q - p);        // 相隔 2 個元素
\end{code}

\begin{example}{c}{(pointer-arithmetic.c) 指標加減}
#include <stdio.h>

int main(void) {
    int arr[5] = {10, 20, 30, 40, 50};
    int *p = arr;

    printf("%d\n", *p);           // 10
    printf("%d\n", *(p + 1));     // 20
    printf("%d\n", *(p + 4));     // 50

    int *q = &arr[4];
    printf("%ld\n", q - p);       // 4 個元素的差距
}
\end{example}

\iosample{}{
10 \\
20 \\
50 \\
4 \\
}

\texttt{p} 指向 \texttt{arr[0]}，\texttt{*(p + 4)} 就是
\texttt{arr[4]}。 \texttt{q - p} 得到 $4$，單位是元素個數，已經把
位元組差距除以 \texttt{sizeof(int)} 了。

\keypoint{陣列名在多數運算式裡會退化成指向首元素的指標；\texttt{p
+ 1} 移動的是一個元素的位元組數，兩個指向同一陣列的指標相減得到元
素間距。}

\begin{drills}[label=drills:ptr-decay-arith]
\item \texttt{sizeof(a)} 跟 \texttt{\&a} 為什麼是陣列名「不會退化」
的兩個例外？
\item \texttt{p + 1} 實際上移動幾個位元組？這跟 \texttt{p} 的型別有
沒有關係？
\item 兩個指向同一陣列的指標相減，得到的單位是位元組還是元素個數？
\end{drills}

\section{多維陣列與記憶體攤平}
\label{sec:ptr-multidim}

一維走訪弄清楚之後，二維陣列在記憶體裡仍然是一長條。 陣列章已經把
列優先（row-major）講過：\texttt{mat[r][c]} 對應的偏移是
\texttt{r * C + c}。 若手上有兩個指向同一塊矩陣的指標，相減就能還
原出行、列座標。

\begin{code}{c}{由指標相減還原列與欄索引}
int mat[10][10], R = 8, C = 6;
int *start = &mat[0][0];           // 矩陣第一格
int *ptr = &mat[r][c];             // 指向 mat[r][c]

int offset = (int)(ptr - start);
int row = offset / C;
int col = offset % C;
\end{code}

範例 \thechapter.22 用一維陣列模擬攤平後的矩陣，驗證這組公式。

\begin{example}{c}{(matrix-pointer.c) 二維陣列的指標操作}
#include <stdio.h>

int main(void) {
    int R, C, r, c;
    scanf("%d %d %d %d", &R, &C, &r, &c);

    int mat[10005];
    for (int i = 0; i < R * C; i++) {
        mat[i] = i + 1;
    }

    int *start = mat;
    int *ptr = mat + r * C + c;
    int row = (int)((ptr - start) / C);
    int col = (int)((ptr - start) % C);
    printf("%d %d\n", row, col);
}
\end{example}

\iosample{
8 6 4 2 \\
}{
4 2 \\
}

輸入 $R=8$、$C=6$、$r=4$、$c=2$，偏移量是 $4 \times 6 + 2 = 26$，還
原回來仍是列 $4$、欄 $2$。

\subsection{雙重解參考與列指標}

\texttt{a[i][j]} 在編譯器裡會展開成 \texttt{*(*(a + i) + j)}。 先
看型別就順了。 \texttt{int a[M][N]} 的元素是長度為 $N$ 的一維陣列；
\texttt{a} 退化之後得到的是指向第一列的指標，型別是
\texttt{int (*)[N]}。 \texttt{a + i} 一次跳一整列（$N$ 個
\texttt{int}），\texttt{*(a + i)} 是第 $i$ 列這條陣列，再退化成
\texttt{int *}，加上 \texttt{j} 才指到格子本身。

所以指向二維陣列的指標必須宣告成 \texttt{int (*p)[N]}：括號不能省，
\texttt{int *p[N]} 會變成指標陣列，下一節再對照。 \texttt{p + 1}
移動的是 \texttt{sizeof(*p)}，也就是一列的位元組數。

\begin{code}{c}{a[i][j] 展開與列指標}
int a[2][3] = {{1, 2, 3}, {4, 5, 6}};
int (*p)[3] = a;          // p 指向一列（3 個 int）

// a[i][j]  展開成  *(*(a + i) + j)
// p[1][2]  與      *(*(p + 1) + 2)  都是 6
// sizeof(*p) 是 3 * sizeof(int)
\end{code}

範例 \thechapter.24 印出 \texttt{sizeof(*p)} 與右下角那一格，用來
核對一次跳一列、再解兩次參考的路徑。

\begin{example}{c}{(row-pointer.c) 用列指標走訪二維陣列}
#include <stdio.h>

int main(void) {
    int a[2][3] = {
        {1, 2, 3},
        {4, 5, 6}
    };
    int (*p)[3] = a;

    printf("%zu\n", sizeof(*p));
    printf("%d\n", *(*(p + 1) + 2));
    return 0;
}
\end{example}

\iosample{}{
12 \\
6 \\
}

\texttt{sizeof(*p)} 是 $12$，因為一列有 $3$ 個 \texttt{int}。
\texttt{p + 1} 跳到第二列，再加 $2$ 走到該列第三格，值是 $6$，跟
\texttt{a[1][2]} 相同。

\keypoint{\texttt{a[i][j]} 會展開成雙重解參考。 指向二維陣列的指標
要寫成 \texttt{int (*p)[N]}，\texttt{p + 1} 一次跳一整列。}

\begin{drills}[label=drills:ptr-multidim]
\item \texttt{int (*p)[3] = a;} 之後，\texttt{p + 1} 移動的是多少個
\texttt{int}？
\item \texttt{a[i][j]} 展開成哪一種雙重解參考的寫法？
\item \texttt{int (*p)[N]} 跟 \texttt{int *p[N]}，型別上差在哪裡？
\end{drills}

\section{陣列的陣列與指標陣列}
\label{sec:ptr-array-of-pointers}

真正的二維陣列 \texttt{int grid[2][3]} 在記憶體裡是連續 $6$ 個
\texttt{int}；用 \texttt{grid[r][c]} 存取時，編譯器用列優先公式一
次算好偏移。 另一種常見寫法是 \texttt{int *rows[2]}：這是\textbf{
指標陣列}（array of pointers），每個元素是一個 \texttt{int *}，可
以指向長度不同的列。 這種各列分開、用指標串起來的結構，有時叫做
Iliffe vector。 兩者語法都有 \texttt{[行][列]}，底層記憶體布局完
全不同。

\begin{code}{c}{連續二維陣列 vs 指標陣列}
int grid[2][3] = {{1, 2, 3}, {4, 5, 6}};  // 6 個 int 連續存放
int row0[] = {1, 2, 3};
int row1[] = {4, 5, 6};
int *rows[2] = {row0, row1};              // 兩個指標，各指向一列

// grid[1][2] 靠偏移一次算完；rows[1][2] 先取指標再偏移
\end{code}

\begin{center}
\begin{tikzpicture}
    \node[cell] (c0) {1};
    \node[cell, right=0pt of c0] (c1) {2};
    \node[cell, right=0pt of c1] (c2) {3};
    \node[cell, right=0pt of c2] (c3) {4};
    \node[cell, right=0pt of c3] (c4) {5};
    \node[cell, right=0pt of c4] (c5) {6};
    \node[idxlbl, left=0.2cm of c0] {\texttt{grid}};

    \node[cell, below=1.5cm of c0] (p0) {};
    \node[idxlbl, font=\scriptsize] at (p0.center) {$\rightarrow$};
    \node[idxlbl, below=0.15cm of p0] {\texttt{rows[0]}};
    \node[cell, right=0.35cm of p0] (p1) {};
    \node[idxlbl, font=\scriptsize] at (p1.center) {$\rightarrow$};
    \node[idxlbl, below=0.15cm of p1] {\texttt{rows[1]}};

    \node[cell, right=1.6cm of p1] (r00) {1};
    \node[cell, right=0pt of r00] (r01) {2};
    \node[cell, right=0pt of r01] (r02) {3};
    \node[cell, below=0pt of r00] (r10) {4};
    \node[cell, right=0pt of r10] (r11) {5};
    \node[cell, right=0pt of r11] (r12) {6};

    \draw[diagarrow] (p0) -- (r00.west);
    \draw[diagarrow] (p1) -- (r10.west);
\end{tikzpicture}
\diagcaption{上面是連續的 \texttt{int grid[2][3]}，六格排成一長條。
下面是指標陣列：\texttt{rows[0]}、\texttt{rows[1]} 各握一列的起始
位址，兩列不必相鄰。}
\label{fig:ch09-iliffe}
\end{center}

\texttt{grid[1][2]} 用 $1 \times 3 + 2$ 一次算出位移；
\texttt{rows[1][2]} 先從 \texttt{rows[1]} 取出第二列的位址，再往後
走 $2$ 格。 步數不同，列與列之間也不保證連續，所以不能把其中一種
的指標直接當成另一種來用。 後面做到動態配置列、二層表時，多半就是
指標陣列這條路。

\keypoint{\texttt{int a[M][N]} 是陣列的陣列，記憶體連續；
\texttt{int *a[M]} 是指標陣列，每一列可以分開存放。 存取步數與布
局不同，不可混用。}

\begin{drills}[label=drills:ptr-array-of-pointers]
\item \texttt{int grid[2][3]} 跟 \texttt{int *rows[2]}，記憶體布局上
差在哪裡？
\item \texttt{rows[1][2]} 要經過幾個步驟才能找到那個值？跟
\texttt{grid[1][2]} 比起來呢？
\end{drills}

\section{函式指標}
\label{sec:ptr-funcptr}

變數有位址，函式也有。 \textbf{函式指標}（function pointer）存的
是某個函式的進入點，讓我們在執行期才決定要呼叫哪一個函式。 宣告時
括號的位置差一個，型別就完全不同。

\texttt{int *f(int)} 是函式：接受 \texttt{int}，回傳
\texttt{int *}。 因為函式宣告的 \texttt{()} 綁得比 \texttt{*} 緊，
星號先貼上了回傳型別。 \texttt{int (*f)(int)} 才是指標：指向一個
接受 \texttt{int}、回傳 \texttt{int} 的函式。 要宣告函式指標，括
號必須把 \texttt{*f} 包起來。

\begin{code}{c}{函式指標的宣告與回調}
int *g(int n);            // 函式：回傳 int *
int (*f)(int);            // 指標：指向  int(int) 這種函式

typedef int (*BinOp)(int, int);

int apply(BinOp op, int a, int b) {
    return op(a, b);      // 執行期才決定呼叫哪一個
}

int add(int x, int y) { return x + y; }
int mul(int x, int y) { return x * y; }
// apply(add, 5, 3) 得到 8
// apply(mul, 5, 3) 得到 15
\end{code}

\texttt{apply} 把函式當參數傳進來，這就是回調（callback）：呼叫者
決定策略，被呼叫者在適當時機執行。 把多個函式放進陣列，還能依輸入
動態切換演算法，這種寫法常被叫做策略模式（strategy pattern）。

範例 \thechapter.27 用函式指標陣列存放加、減、乘，再用一個整數指令
選要用哪一個。

\begin{example}{c}{(function-pointer-ops.c) 用函式指標陣列切換運算}
#include <stdio.h>

int add(int a, int b) {
    return a + b;
}

int sub(int a, int b) {
    return a - b;
}

int mul(int a, int b) {
    return a * b;
}

int main(void) {
    int (*ops[3])(int, int) = {add, sub, mul};
    int a, b, cmd;

    scanf("%d %d %d", &a, &b, &cmd);
    if (cmd >= 0 && cmd < 3) {
        printf("%d\n", ops[cmd](a, b));
    }
    return 0;
}
\end{example}

\iosample{
10 3 0 \\
}{
13 \\
}

\texttt{ops[0]} 是 \texttt{add}、\texttt{ops[1]} 是 \texttt{sub}、
\texttt{ops[2]} 是 \texttt{mul}。 輸入 \texttt{10 3 0} 時呼叫
\texttt{add(10, 3)}，得到 $13$；指令改成 $1$ 會得到 $7$，改成 $2$
會得到 $30$。 \texttt{main} 裡沒有一串 \texttt{if}-\texttt{else}
去分運算，只是查表之後呼叫。

函式指標也可以當成狀態機的目前狀態：每個狀態對應一個處理函式，處
理完再把指標改成下一個狀態。 剖析輸入、協定處理常用這種寫法；觀念
上就是把目前該做的事存在一個函式指標裡，跟上面的策略表同一種工具。

\keypoint{\texttt{int (*f)(int)} 是函式指標，\texttt{int *f(int)}
是回傳指標的函式。 函式指標用來做回調、動態切換策略，以及把狀態機
的目前狀態存在指標裡。}

\begin{drills}[label=drills:ptr-funcptr]
\item \texttt{int *f(int)} 跟 \texttt{int (*f)(int)}，哪一個才是函式
指標？
\item \texttt{apply(BinOp op, int a, int b)} 這種把函式當參數傳進去
的寫法，叫做什麼？
\item 範例 \texttt{function-pointer-ops.c} 裡，若輸入指令改成
\texttt{2}，會呼叫哪一個函式？
\end{drills}

\section{常見陷阱與最佳實踐}
\label{sec:ptr-pitfalls}

前幾節已經談過 wild pointer、回傳區域變數指標、陣列與指標的型別差
異。 指標讓改寫呼叫者、走訪陣列、多維布局與函式策略都成為可能，也
把越界、型別、生命週期的責任交給我們。 這一節集中列出 C 不會幫你
擋的幾類實戰陷阱。

\subsection{越界與 buffer overflow}

C 語言不檢查陣列邊界；指標算術多走一步，可能讀寫到別的變數，或直
接 Segmentation Fault。 這就是 buffer overflow：寫過了這塊緩衝區
的盡頭，把後面的記憶體破壞掉。

\begin{code}{c}{越界走訪（錯誤示範）}
int arr[5];
int *p = arr;
for (int i = 0; i <= 5; i++) {   // i=5 時已越界
    printf("%d\n", *p++);
}
\end{code}

比較穩的走法是先算出尾端指標，迴圈條件用 \texttt{p < end}，讓合法
範圍寫在一個地方。

\begin{code}{c}{用 end 指標限定範圍}
int *end = arr + 5;
for (int *p = arr; p < end; p++) {
    printf("%d\n", *p);
}
\end{code}

\subsection{型別不匹配}

\texttt{int *} 只能指向 \texttt{int}；把 \texttt{double *} 硬指到
\texttt{int} 會觸發警告或未定義行為。 就算強制轉型編譯過了，解參
考時仍會用錯的寬度去讀記憶體。

在 little-endian 機器上，較低的位元組排在較低的位址。 例如
\texttt{int x = 0x01020304} 在記憶體裡是 \texttt{04 03 02 01} 這四
個位元組。 若把 \texttt{\&x} 轉成 \texttt{short *} 再解參考，讀到
的是前兩個位元組組成的 \texttt{0x0304}，也就是 $772$，後面兩個位
元組被截掉了。 數值看起來像個普通整數，其實已經跟原本的
\texttt{int} 對不上。

\begin{code}{c}{型別必須一致}
double x = 3.14;
double *p = &x;              // 正確
// int *q = &x;              // 編譯警告：型別不符

int n = 0x01020304;
short *ps = (short *)&n;     // 強制轉型，little-endian 上 *ps 常是 772
\end{code}

\subsection{懸空指標}

動態配置的完整規則會在高等技術那一章系統地講；這裡只要先認識一種
跟生命週期有關的陷阱。 \texttt{malloc} 從 heap 要來的記憶體，生命
週期由我們自己決定：\texttt{free} 之後那塊空間已經不屬於程式，但
指標變數裡往往還握著舊位址。 這種指向已釋放記憶體的指標叫做
\textbf{懸空指標}（dangling pointer）。 再對它讀寫是未定義行為。

好習慣是 \texttt{free} 之後立刻把指標設成 \texttt{NULL}。 之後若
不小心再解參考，會立刻以對 \texttt{NULL} 解參考的形式爆掉，比默默
用到已經被別人回收的記憶體好查得多。

\begin{code}{c}{free 之後把指標設成 NULL}
#include <stdlib.h>

int *p = (int *)malloc(sizeof(int));
if (p != NULL) {
    *p = 42;
    free(p);
    p = NULL;     // 避免懸空指標
    // *p = 1;    // 現在會立刻出問題，比較好查
}
\end{code}

區域變數回傳造成的失效指標，本質上也是懸空指標：物件先消失，指標
還留著舊位址。 兩種情況的防禦相同：不要讓指標活得比它指向的物件久，
失效後就把它設成 \texttt{NULL}。

\subsection{跨檔案：陣列 vs 指標宣告}

前面強調陣列與指標型別不同，跨檔用 \texttt{extern} 宣告時也要一致。
一邊定義成陣列、另一邊卻宣告成指標，連結之後行為會異常。

\begin{code}{c}{file1.c：定義陣列}
int arr[100];
\end{code}

\begin{code}{c}{file2.c：extern 必須與定義一致}
extern int arr[];    // 正確：仍是陣列型別
// extern int *arr;  // 錯誤：型別變成指標，連結後行為異常
\end{code}

\subsection{最佳實踐}

\begin{itemize}
\item \textbf{宣告即初始化}：指標立刻指向有效物件或 \texttt{NULL}
\item \textbf{使用前檢查}：解參考前確認 \texttt{p != NULL}
\item \textbf{確認範圍}：走訪時用 \texttt{end} 指標或索引邊界
\item \textbf{型別匹配}：指標型別與指向物件一致，不要靠強制轉型
      混用寬度不同的型別
\item \textbf{陣列參數帶 \texttt{n}}：函式無法從 \texttt{int a[]}
      推得長度
\item \textbf{唯讀走訪加 \texttt{const}}：避免意外修改
\item \textbf{失效即清空}：\texttt{free} 之後、或物件生命週期結
      束之後，把指標設成 \texttt{NULL}
\end{itemize}

\keypoint{指標讓改寫呼叫者、走訪陣列、動態配置、資料結構與函式指
標成為可能，也是 bug 與安全問題的來源。 把位址、退化、生命週期與
\texttt{const} 記清楚，大多數陷阱都能避開。}

\begin{drills}[label=drills:ptr-pitfalls]
\item 什麼是 buffer overflow？範例裡 \texttt{for (int i = 0; i <= 5;
i++)} 走訪 \texttt{int arr[5]} 錯在哪？
\item \texttt{free(p)} 之後為什麼建議立刻把 \texttt{p} 設成
\texttt{NULL}？
\item \texttt{extern int arr[];} 跟 \texttt{extern int *arr;} 為什
麼不能拿來宣告同一個陣列？
\end{drills}

\addcontentsline{toc}{section}{練習解答}
\section*{練習解答}

各節練習的答案收在這裡。 短題只給能一句話核對的結果。

\paragraph{第~\ref{sec:ptr-basics}~節}
\begin{enumerate}
\item \texttt{42} 是右值，\texttt{x + 1} 也是右值；\texttt{\&} 只能
加在左值前面。
\item \texttt{x} 只是普通 \texttt{int}；只有 \texttt{p3} 前面有
\texttt{*}，才是指標。
\item 因為宣告完但未指派的指標是 wild pointer，內容是垃圾值，對它解
參考是未定義行為。
\item 幾乎立刻 Segmentation Fault；因為錯誤會停在解參考那一步，不會
默默改到別的記憶體，比較容易定位。
\end{enumerate}

\paragraph{第~\ref{sec:ptr-params}~節}
\begin{enumerate}
\item 傳的是 \texttt{a}、\texttt{b} 的位址；漏寫 \texttt{\&} 會變成
傳值，\texttt{swap} 內部改的只是自己的參數複本，\texttt{main} 裡的
\texttt{a}、\texttt{b} 不會交換。
\item 因為一次要交出兩個結果（最大值與其索引），\texttt{return} 只
能回傳一個值，所以用指標參數把兩個結果都寫回呼叫者。
\item 沒有必要改呼叫者狀態、或只需要一個回傳值時，優先寫純函式；需
要改寫呼叫者的變數或一次回傳多個結果時，才用指標參數。
\end{enumerate}

\paragraph{第~\ref{sec:ptr-const}~節}
\begin{enumerate}
\item \texttt{int * const p} 不能改 \texttt{p} 指向哪裡；\texttt{const
int *p} 不能改 \texttt{*p} 的值。
\item 該加；編譯器會擋掉函式內任何想寫入該資料的敘述，避免手滑改到
呼叫者的陣列。
\item 同時保護了指標本身（不能改指向）與指向的資料（不能改值）。
\end{enumerate}

\paragraph{第~\ref{sec:ptr-return-lifetime}~節}
\begin{enumerate}
\item \texttt{x} 是區域變數，函式結束後它的 stack frame 就被銷毀，
\texttt{\&x} 變成指向已回收記憶體的位址。
\item 指向的物件生命週期要長過這次函式呼叫；參數陣列與全域變數都不
是這次呼叫的區域變數，函式結束後仍然存在。
\item 回傳 \texttt{NULL}；呼叫端要先檢查回傳值是否為 \texttt{NULL}，
再解參考。
\end{enumerate}

\paragraph{第~\ref{sec:ptr-array-vs-pointer}~節}
\begin{enumerate}
\item 不一樣；\texttt{sizeof(a)} 是整個陣列的位元組數，\texttt{sizeof(p)}
只是一個指標變數的大小。
\item 陣列宣告時配置了連續 $5$ 個 \texttt{int} 的空間；指標宣告時只
配置一個存位址的變數，本身不占那 $5$ 格。
\end{enumerate}

\paragraph{第~\ref{sec:ptr-decay-arith}~節}
\begin{enumerate}
\item \texttt{sizeof(a)} 量的是整塊陣列的大小，\texttt{\&a} 得到的是
指向整個陣列的指標，兩者都沒有退化成指向首元素的 \texttt{int *}。
\item 移動 \texttt{sizeof(*p)} 個位元組；有關係，編譯器依指標指向的
型別決定跨多遠。
\item 元素個數；位元組差距已經被除以該型別的大小。
\end{enumerate}

\paragraph{第~\ref{sec:ptr-multidim}~節}
\begin{enumerate}
\item 移動 $3$ 個 \texttt{int}（一整列）。
\item \texttt{*(*(a + i) + j)}。
\item \texttt{int (*p)[N]} 是指向一整列（$N$ 個 \texttt{int}）的指
標；\texttt{int *p[N]} 是有 $N$ 個 \texttt{int *} 元素的指標陣列。
\end{enumerate}

\paragraph{第~\ref{sec:ptr-array-of-pointers}~節}
\begin{enumerate}
\item \texttt{grid} 是連續 $6$ 個 \texttt{int}；\texttt{rows} 的每一
列可以分別存放在記憶體裡不相鄰的地方。
\item \texttt{rows[1][2]} 要先讀 \texttt{rows[1]} 拿到該列位址，再偏
移 $2$；\texttt{grid[1][2]} 只需用列優先公式一次算出偏移。
\end{enumerate}

\paragraph{第~\ref{sec:ptr-funcptr}~節}
\begin{enumerate}
\item \texttt{int (*f)(int)} 是函式指標；\texttt{int *f(int)} 是回
傳 \texttt{int *} 的函式。
\item 回調（callback），也是策略模式的一種寫法。
\item \texttt{mul}，因為 \texttt{ops[2]} 對應 \texttt{mul}。
\end{enumerate}

\paragraph{第~\ref{sec:ptr-pitfalls}~節}
\begin{enumerate}
\item buffer overflow 是寫或讀超過緩衝區邊界；\texttt{i <= 5} 讓
\texttt{i = 5} 時仍執行，但 \texttt{arr[5]} 已經越界。
\item 避免懸空指標：之後若不小心再解參考 \texttt{p}，會立刻對
\texttt{NULL} 解參考而明顯出錯，比默默用到已釋放的記憶體好查。
\item 兩者型別不同：一個仍是陣列，一個變成指標，跨檔案宣告與定義的
型別不一致，連結後行為會異常。
\end{enumerate}

\chapter{字元}
\label{ch:character}

C 語言用 \texttt{char} 型別表示一個字元。 \texttt{char} 其實是一個位元
組（byte）大小的整數型別，字元只是我們選擇賦予它的意義。 在電腦裡，
\texttt{'A'} 就是某個固定的整數（$65$），\texttt{'7'} 也是另一個固定的整
數（$55$），只是我們平常用 \texttt{\%c} 把它們印成看得見的符號而已。

很多看起來像「文字處理」的技巧，本質上都只是整數的讀取、比較、累加；這
一章會從這個角度出發，逐步看到字元怎麼跟標準函式庫配合、怎麼夾在數字
中間被讀取、以及怎麼用來處理長到裝不進任何整數型別的數字字串。

\section{字元的本質：一個位元組的整數}
\label{sec:char-nature}

每個字元對應一個固定的整數編碼。 英文環境最常用的是 ASCII：例如
\texttt{'A'} 是 $65$、\texttt{'a'} 是 $97$、\texttt{'0'} 是 $48$。
這份對照表是硬體、作業系統、程式語言共同遵守的標準。

從記憶體的角度看，\texttt{char} 佔一個 byte（通常是 $8$ 個 bit），所
以它的值域是 $0$ 到 $255$（若視為 \texttt{unsigned char}）或 $-128$
到 $127$（若視為 \texttt{signed char}）。 我們可以隨時在兩種角度切換：
一種是「這是哪個符號」，一種是「這是哪個數字」。

\begin{center}
\begin{tikzpicture}[node distance=0.2cm]
    \node[cell] (c0) {\texttt{0}};
    \node[cell, right=of c0] (c1) {\texttt{A}};
    \node[cell, right=of c1] (c2) {\texttt{a}};
    \node[idxlbl, below=0.15cm of c0] {$48$};
    \node[idxlbl, below=0.15cm of c1] {$65$};
    \node[idxlbl, below=0.15cm of c2] {$97$};
\end{tikzpicture}
\diagcaption{三個字元與它們常見的 ASCII 整數值；同一個 \texttt{char}
變數，用 \texttt{\%c} 印是符號，用 \texttt{\%d} 印是整數。}
\label{fig:ch10-char-ascii}
\end{center}

\keypoint{同一個變數用 \texttt{\%c} 印出來是字元本身，用 \texttt{\%d}
印出來是整數編碼；資料本身沒有改變，只是解讀方式不同。}

\begin{drills}[label=drills:char-nature]
\item \texttt{'A'} 的 ASCII 值是多少？ \texttt{'a'} 呢？
\item 同一個 \texttt{char} 變數，用 \texttt{\%c} 印跟用 \texttt{\%d}
印，資料本身有沒有改變？ 改變的是什麼？
\item 若 \texttt{char} 被視為 \texttt{unsigned char}，值域是多少？ 視
為 \texttt{signed char} 呢？
\end{drills}

\section{讀寫字元}
\label{sec:char-io}

最基本的操作是用 \texttt{\%c} 讀進一個字元，再用不同的格式化字串印出
來。 \texttt{scanf} 的 \texttt{\%c} 不會跳過前面的空白，它會讀下緊接著
的下一個字元。

\begin{example}{c}{(char-to-ascii.c) 印出字元的 ASCII 碼}
#include <stdio.h>

int main(void) {
    char c;
    scanf("%c", &c);
    printf("%c\n", c);     // 印成字元
    printf("%d\n", c);     // 印成整數
}
\end{example}

\iosample{
A \\
}{
A \\
65 \\
}

許多題目會逐字元讀取一整段輸入直到 EOF。 用 \texttt{getchar()} 比
\texttt{scanf("\%c")} 更直觀：

\begin{code}{c}{逐字元讀直到 EOF}
int ch;
while ((ch = getchar()) != EOF) {
    // 處理 ch
}
\end{code}

\texttt{getchar()} 在檔案結束時回傳 \texttt{EOF}（通常是 $-1$），適合
串流處理。

\texttt{getchar()} 的回傳型別是 \texttt{int}。 \texttt{EOF} 這個額外的
訊號值通常是 $-1$，超出 \texttt{unsigned char} 能表示的 $0$ 到 $255$，
有號 \texttt{char} 的 $-128$ 到 $127$ 也裝不穩。 若把回傳值存進
\texttt{char}，某些系統上迴圈比不到真正的 \texttt{EOF}，就不會正常結束。

\begin{code}{c}{危險：用 char 接 getchar 的回傳值}
char ch;
while ((ch = getchar()) != EOF) {
    // 某些系統上比不到真正的 EOF
}
\end{code}

\begin{code}{c}{正確：用 int 接 getchar 的回傳值}
int ch;
while ((ch = getchar()) != EOF) {
    // int 裝得下 EOF
}
\end{code}

\keypoint{迴圈終止條件若跟 \texttt{EOF} 比較，接收變數就該宣告成
\texttt{int}。 這條規則對 \texttt{getchar()} 和後面章節會看到的
\texttt{getc()}／\texttt{fgetc()} 都成立。}

\begin{drills}[label=drills:char-io]
\item \texttt{scanf("\%c")} 會不會跳過輸入前面的空白字元？
\item 用 \texttt{char} 而不是 \texttt{int} 接 \texttt{getchar()} 的回
傳值，可能導致什麼問題？
\item \texttt{while ((ch = getchar()) != EOF)} 這段迴圈，什麼時候會終
止？
\end{drills}

\section{轉義序列}
\label{sec:char-escape}

有些字元無法直接寫在程式碼裡。 例如換行符本身會斷行、雙引號會結束字
串、反斜線會啟動特殊意義。 C 用\textbf{轉義序列}在反斜線後接字母或符號
來表示這些字元。

\begin{center}
\begin{tabular}{|c|c|c|}
\hline
\textbf{轉義序列} & \textbf{名稱} & \textbf{效果} \\
\hline
\texttt{\textbackslash n} & 換行 & 光標移到下一行開頭 \\
\texttt{\textbackslash t} & 製表符 & 向右跳到下個製表位 \\
\texttt{\textbackslash\textbackslash} & 反斜線 & 印出一個 \texttt{\textbackslash} \\
\texttt{\"} & 雙引號 & 印出一個 \texttt{"} \\
\texttt{\'} & 單引號 & 印出一個 \texttt{'} \\
\texttt{\textbackslash 0} & 空字元 & 字串終止符 \\
\hline
\end{tabular}
\end{center}

\begin{example}{c}{(escape-demo.c) 常見的轉義序列}
#include <stdio.h>

int main(void) {
    printf("第一行\n");
    printf("Name\tAge\n");
    printf("路徑: C:\\\\Program Files\n");
    printf("他說\"你好\"\n");
}
\end{example}

\iosample{
}{
第一行 \\
Name\quad Age \\
路徑: C:\textbackslash Program Files \\
他說"你好" \\
}

\begin{drills}[label=drills:char-escape]
\item \texttt{\textbackslash t} 和 \texttt{\textbackslash n} 分別代表
什麼效果？
\item 要在字串裡印出一個雙引號字元，該怎麼寫？
\item 要印出一個反斜線字元，轉義序列該怎麼寫？
\end{drills}

\section{字元分類}
\label{sec:char-classify}

判斷字元的性質是很常見的需求。 C 的標準庫在 \texttt{ctype.h} 裡提供了
分類函式。

\subsection{內建分類函式}

\begin{code}{c}{ctype.h 的常用函式}
#include <ctype.h>

isdigit(c)     // '0'..'9' ？
isalpha(c)     // a..z, A..Z ？
isalnum(c)     // 字母或數字 ？
isspace(c)     // 空白（空格、tab、換行等）？
islower(c)     // 小寫字母 ？
isupper(c)     // 大寫字母 ？
\end{code}

上面六個最常用。 \texttt{ctype.h} 還有幾個同樣好用、容易漏掉的分類函式。

\begin{center}
\begin{tabular}{|c|c|c|}
\hline
\textbf{函式} & \textbf{判斷內容} & \textbf{範例} \\
\hline
\texttt{isxdigit} & 十六進位數字 & \texttt{isxdigit('f')} 真，\texttt{isxdigit('g')} 假 \\
\texttt{ispunct} & 標點符號 & \texttt{ispunct('!')} 真，\texttt{ispunct('a')} 假 \\
\texttt{isgraph} & 可看見、不含空格 & \texttt{isgraph('A')} 真，\texttt{isgraph(' ')} 假 \\
\texttt{isprint} & 可看見、含空格 & \texttt{isprint('A')} 真，\texttt{isprint(' ')} 真 \\
\texttt{iscntrl} & 控制字元 & \texttt{iscntrl('\textbackslash n')} 真，\texttt{iscntrl('A')} 假 \\
\texttt{isblank} & C99：只認空格與 tab & \texttt{isblank(' ')} 真，\texttt{isblank('\textbackslash n')} 假 \\
\hline
\end{tabular}
\end{center}

最容易搞混的是 \texttt{isgraph} 與 \texttt{isprint}：\texttt{isprint}
把空格算進去，\texttt{isgraph} 不算。 \texttt{isblank} 比
\texttt{isspace} 窄，換行不算空白。

\begin{example}{c}{(classify-stream.c) 統計文字中各類字元的個數}
#include <stdio.h>
#include <ctype.h>

int main(void) {
    char c;
    int digits = 0, letters = 0, others = 0;

    while (scanf("%c", &c) != EOF) {
        if (isdigit(c)) {
            digits++;
        } else if (isalpha(c)) {
            letters++;
        } else {
            others++;
        }
    }

    printf("%d %d %d\n", digits, letters, others);
}
\end{example}

\iosample{
abc123 XYZ! \\
}{
3 6 2 \\
}

\subsection{自訂判斷式}

\texttt{ctype.h} 沒有「母音」、「子音」之類的分類，需要自己寫條件。

\begin{example}{c}{(is-vowel.c) 判斷字元是否為母音}
#include <stdio.h>

int isVowel(char c) {
    return c == 'A' || c == 'E' || c == 'I' || c == 'O' || c == 'U' ||
           c == 'a' || c == 'e' || c == 'i' || c == 'o' || c == 'u';
}

int main(void) {
    char c;
    scanf("%c", &c);
    printf("%d\n", isVowel(c));
}
\end{example}

\iosample{
a \\
}{
1 \\
}

\begin{drills}[label=drills:char-classify]
\item \texttt{isdigit('7')} 的結果是真還是假？ \texttt{isalpha('7')}
呢？
\item \texttt{isprint(' ')} 和 \texttt{isgraph(' ')} 的結果一樣嗎？
差在哪裡？
\item \texttt{ctype.h} 裡有沒有現成的函式可以判斷一個字元是不是母
音？ 要怎麼自己寫這個判斷式？
\end{drills}

\section{字元轉換}
\label{sec:char-convert}

\texttt{ctype.h} 提供轉換函式：\texttt{toupper} 轉大寫、\texttt{tolower}
轉小寫。 若輸入不是字母，這兩個函式會原樣回傳。

\begin{example}{c}{(to-case.c) 轉換字元大小寫}
#include <stdio.h>
#include <ctype.h>

int main(void) {
    char c;
    scanf("%c", &c);
    printf("%c %c\n", toupper(c), tolower(c));
}
\end{example}

\iosample{
A \\
}{
A a \\
}

\begin{drills}[label=drills:char-convert]
\item \texttt{toupper('a')} 的結果是什麼？ \texttt{toupper('5')} 呢
（輸入不是字母的情況）？
\item \texttt{to-case.c} 範例讀入字元 \texttt{A}，輸出是什麼？
\end{drills}

\section{字元在運算中的應用}
\label{sec:char-arithmetic}

\subsection{數字字元與整數的轉換}

因為 \texttt{'0'} 到 \texttt{'9'} 的 ASCII 碼連續遞增（48 到 57），
一個數字字元減去 \texttt{'0'} 就得到它的整數值。

\begin{example}{c}{(digit-convert.c) 把數字字元轉成整數}
#include <stdio.h>

int main(void) {
    char c;
    scanf("%c", &c);
    printf("%d\n", c - '0');
}
\end{example}

\iosample{
7 \\
}{
7 \\
}

反過來，整數加上 \texttt{'0'} 就成了字元：\texttt{(7 + '0')} 是
\texttt{'7'}。

\subsection{混合讀入字元與數字}

\texttt{scanf} 可以在同一次呼叫裡混合讀入數字和字元。 注意 \texttt{\%d}
會跳過前面空白，但 \texttt{\%c} 不會。

\begin{example}{c}{(operator-calc.c) 讀運算式並計算}
#include <stdio.h>

int main(void) {
    int a, b;
    char op;

    scanf("%d%c%d", &a, &op, &b);

    switch (op) {
        case '+':
            printf("%d\n", a + b);
            break;
        case '-':
            printf("%d\n", a - b);
            break;
        case '*':
            printf("%d\n", a * b);
            break;
    }
}
\end{example}

\iosample{
3+5 \\
}{
8 \\
}

輸入中間不能有空白，因為 \texttt{\%c} 會讀到空格而不是運算子。

\subsection{逐字元處理超大整數}

超大整數無法存進任何變數，但整除性只需要逐字元累加。 例如判斷能否被 3
整除只需看各位數字之和。

\begin{example}{c}{(divisibility.c) 判斷超大整數的整除性}
#include <stdio.h>
#include <ctype.h>

int main(void) {
    int sum = 0, alt = 0, sign = 1, lastDigit = 0;
    int ch;

    while ((ch = getchar()) != EOF && ch != '\n') {
        if (!isdigit(ch)) continue;

        int d = ch - '0';
        lastDigit = d;
        sum += d;
        alt += sign * d;
        sign = -sign;
    }

    // 檢查整除性
    printf("%s ", (lastDigit % 2 == 0) ? "yes" : "no");    // 被 2 整除
    printf("%s\n", (sum % 3 == 0) ? "yes" : "no");         // 被 3 整除
}
\end{example}

\iosample{
190949 \\
}{
no yes \\
}

\keypoint{流式處理讓我們處理任意長度的數字，只用固定的幾個小變數。}

\begin{drills}[label=drills:char-arithmetic]
\item \texttt{'7' - '0'} 的結果是什麼？ 這背後靠的是 ASCII 碼的什麼性
質？
\item \texttt{scanf("\%d\%c\%d", \&a, \&op, \&b)} 讀入 \texttt{3+5}
時，\texttt{op} 會存到什麼字元？
\item 判斷一個超大整數能不能被 \texttt{3} 整除，為什麼不需要把整個數
字存進一個 \texttt{int} 變數？
\end{drills}

\section{進階應用}
\label{sec:char-advanced}

\subsection{從字元流組出單字}

單個字符之外，常需要將連續的字母累積成一個單詞。 遇到空白或標點就把目
前的單詞結案、開始下一個。

\begin{code}{c}{逐字元累積單字}
char word[100];
int len = 0;

if (isalpha(c)) {
    word[len++] = c;
} else if (len > 0) {
    word[len] = '\0';
    // 處理這個單詞
    len = 0;
}
\end{code}

\subsection{前瞻一格：句號與句子邊界}

某些問題需要查看「下一個字元」才能決定如何處理目前字元。 例如判斷句號是
否結束句子，需要看句號後面是什麼。

\begin{code}{c}{前瞻一格的做法}
for (int i = 0; i < n; i++) {
    if (text[i] == '.') {
        char next = (i + 1 < n) ? text[i + 1] : '\0';
        if (next == ' ' || next == '\n' || next == '\0') {
            // 這個句號算數
        }
    }
}
\end{code}

\subsection{滑動視窗：三字模式}

固定大小的窗口逐個字符滑動，可以提取連續的模式。 重點是窗口遇到邊界
時要重設。

\begin{code}{c}{長度為 3 的滑動窗口}
char window[3];
int wlen = 0;

if (isalpha(c)) {
    if (wlen < 3) {
        window[wlen++] = c;
    } else {
        window[0] = window[1];
        window[1] = window[2];
        window[2] = c;
    }
} else {
    wlen = 0;  // 重設窗口
}
\end{code}

\subsection{驗證使用者輸入的單一字元}

使用者常常會多打空白或整行輸入。 如果只想要第一個字元，需要把剩下的
字元丟掉，不然下一次讀取會讀到殘留的內容。

範例 \thechapter.16 問要不要繼續，只接受 \texttt{y} 或 \texttt{n}；其他
輸入就印提示、再讀一次。 讀到第一個字元之後，把該行讀到換行符為止，
殘留就清掉了。

\begin{example}{c}{(get-yes-no.c) 只接受 y 或 n 的驗證迴圈}
#include <stdio.h>

char getYesNo() {
    int ch;
    while (1) {
        ch = getchar();
        while (getchar() != '\n') {
            continue;              // 丟掉這一行剩下的字元
        }
        if (ch == 'y' || ch == 'n') {
            return ch;
        }
        printf("請輸入 y 或 n：\n");
    }
}

int main(void) {
    printf("要繼續嗎？(y/n)\n");
    char ans = getYesNo();
    printf("%c\n", ans);
}
\end{example}

\iosample{
hello \\
y \\
}{
要繼續嗎？(y/n) \\
請輸入 y 或 n： \\
y \\
}

\keypoint{讀取單一字元時，把該行剩下的字元清掉（讀到 \texttt{'\textbackslash n'}
為止），才不會讓下一次讀取吃到殘留輸入。 這跟陷阱 1 裡 \texttt{\%c}
讀到換行符是同一個原因。}

\begin{drills}[label=drills:char-advanced]
\item 逐字元累積單字時，遇到什麼樣的字元代表一個單字結束了？
\item \texttt{get-yes-no.c} 範例裡，為什麼讀到第一個字元後，還要用一
個內層 \texttt{while} 迴圈把該行剩下的字元讀掉？
\item 滑動視窗的大小固定為 \texttt{3}，當視窗已經滿了、又進來一個新
字元時，程式怎麼處理？
\end{drills}

\section{常見陷阱與最佳實踐}
\label{sec:char-pitfalls}

\subsection{陷阱 1：\texttt{\%c} 讀到換行符}

\texttt{scanf("\%d")} 讀完後，換行符留在緩衝區。 下一個 \texttt{\%c}
會誤讀這個換行。

\begin{example}{c}{(newline-trap.c) \%c 吞掉換行的陷阱}
#include <stdio.h>

int main(void) {
    int x;
    char c;

    scanf("%d", &x);           // 讀完 x，\n 留在緩衝區
    scanf("%c", &c);           // \%c 讀到 \n 而不是下一個字元
    printf("x=%d, c=%c\n", x, c);
}
\end{example}

\iosample{
42 \\
A \\
}{
x=42, c= \\
}

修正：在第二個 \texttt{scanf} 前用 \texttt{scanf("\%c", \&dummy)} 消
耗換行，或改用 \texttt{getchar()}。

\subsection{陷阱 2：假設字元在某範圍}

\texttt{c - '0'} 只在 \texttt{c} 確實是數字時有意義。

\begin{code}{c}{不安全的假設}
char c;
scanf("%c", &c);
int digit = c - '0';  // 若 c 是 'a'，結果無意義
\end{code}

安全的做法：先檢查
\begin{code}{c}{先驗證再使用}
if (isdigit(c)) {
    int digit = c - '0';
} else {
    // 處理非數字
}
\end{code}

\subsection{陷阱 3：\texttt{isspace} vs 只檢查空格}

\texttt{isspace} 涵蓋空格、tab、換行等。 若題目說「空白」，用
\texttt{isspace} 才對。

\begin{code}{c}{不夠全面}
if (c == ' ')  // 只檢查空格，漏掉 tab 和換行
\end{code}

\begin{code}{c}{正確做法}
if (isspace(c))  // 涵蓋所有空白
\end{code}

\subsection{陷阱 4：大小寫轉換不用手寫公式}

\texttt{toupper}/\texttt{tolower} 已經封裝好，不要手寫 \texttt{c - 32}
之類的黑魔法。

\subsection[陷阱 5：ctype.h 參數要轉 unsigned char]{陷阱 5：\texttt{ctype.h} 的參數先轉成 unsigned char}

\texttt{ctype.h} 的函式都吃 \texttt{int} 參數。 若系統上 \texttt{char}
是有號的，處理中文或其他非 ASCII 位元組時，某個 \texttt{char} 的值會
超過 $127$，自動轉成 \texttt{int} 就變成負數。 把負數丟進
\texttt{isalpha} 這類函式，行為沒有保證，不同編譯器可能不一樣。

保險的做法是先轉成 \texttt{unsigned char} 再傳進去。

\begin{code}{c}{有風險：有號 char 直接丟進 ctype}
char c;
isalpha(c);   // c 超過 127 時可能變成負數
\end{code}

\begin{code}{c}{安全：先轉成 unsigned char}
char c;
isalpha((unsigned char)c);
\end{code}

\subsection{最佳實踐}

\begin{itemize}
\item \textbf{用 \texttt{getchar()} 處理串流}：逐字元讀到 EOF 時比
\texttt{scanf} 更直觀
\item \textbf{字元判斷用 \texttt{ctype.h}}：\texttt{isdigit} 比手寫比較式
更清楚、更不易出錯
\item \textbf{先驗證後使用}：在 \texttt{c - '0'} 前先檢查 \texttt{isdigit(c)}
\item \textbf{混合讀取要小心空白}：\texttt{\%d} 跳過空白，\texttt{\%c} 不跳
\item \textbf{轉換用標準函式}：\texttt{toupper}/\texttt{tolower} 優於手寫
公式
\end{itemize}

\keypoint{字元處理的核心是理解「字元是整數」以及「逐字元流式處理」的思
想；掌握這兩點，從簡單的分類判斷到複雜的狀態機就都能應對。}

\begin{drills}[label=drills:char-pitfalls]
\item \texttt{scanf("\%d", \&x)} 之後緊接著 \texttt{scanf("\%c", \&c)}，
\texttt{c} 常常會讀到什麼字元？ 該怎麼修正？
\item 為什麼把 \texttt{char c} 直接丟進 \texttt{isalpha(c)} 有風險？
正確做法是什麼？
\item \texttt{toupper}/\texttt{tolower} 已經處理好大小寫轉換，為什麼
不建議自己手寫 \texttt{c - 32} 這類算式？
\end{drills}

\addcontentsline{toc}{section}{練習解答}
\section*{練習解答}

各節練習的答案收在這裡。 短題只給能一句話核對的結果。

\paragraph{第~\ref{sec:char-nature}~節}
\begin{enumerate}
\item \texttt{'A'} 是 $65$；\texttt{'a'} 是 $97$。
\item 資料本身沒有改變，改變的只是\emph{解讀方式}：同一組位元，
\texttt{\%c} 當成符號印，\texttt{\%d} 當成整數印。
\item 視為 \texttt{unsigned char} 時是 $0$ 到 $255$；視為
\texttt{signed char} 時是 $-128$ 到 $127$。
\end{enumerate}

\paragraph{第~\ref{sec:char-io}~節}
\begin{enumerate}
\item 不會。 \texttt{\%c} 會直接讀下一個字元，包括空白字元。
\item 某些系統上 \texttt{char} 裝不下 \texttt{EOF}（通常是 $-1$），迴
圈可能永遠比不到真正的 \texttt{EOF}，變成無窮迴圈。
\item 讀到 \texttt{getchar()} 回傳 \texttt{EOF} 時終止，也就是輸入串
流結束的時候。
\end{enumerate}

\paragraph{第~\ref{sec:char-escape}~節}
\begin{enumerate}
\item \texttt{\textbackslash t} 讓游標跳到下一個定位點；
\texttt{\textbackslash n} 讓游標換到下一行開頭。
\item \texttt{\textbackslash"}。
\item \texttt{\textbackslash\textbackslash}。
\end{enumerate}

\paragraph{第~\ref{sec:char-classify}~節}
\begin{enumerate}
\item \texttt{isdigit('7')} 是真；\texttt{isalpha('7')} 是假（\texttt{'7'}
是數字不是字母）。
\item 不一樣。 \texttt{isprint(' ')} 是真（空格算可印字元），
\texttt{isgraph(' ')} 是假（\texttt{isgraph} 不把空格算進去）。
\item 沒有現成函式；要自己寫條件式，例如逐一比對
\texttt{'a'/'e'/'i'/'o'/'u'} 及其大寫。
\end{enumerate}

\paragraph{第~\ref{sec:char-convert}~節}
\begin{enumerate}
\item \texttt{toupper('a')} 是 \texttt{'A'}；\texttt{toupper('5')} 原
樣回傳 \texttt{'5'}（不是字母就不轉換）。
\item 輸出 \texttt{A a}：先印大寫版本，再印小寫版本。
\end{enumerate}

\paragraph{第~\ref{sec:char-arithmetic}~節}
\begin{enumerate}
\item 結果是 $7$；靠的是 \texttt{'0'} 到 \texttt{'9'} 的 ASCII 碼連續
遞增這個性質。
\item \texttt{op} 會存到 \texttt{'+'}。
\item 因為超大整數裝不進任何整數型別；只需要逐字元累加各位數字之
和，就能判斷整除性，不需要把整個數字存起來。
\end{enumerate}

\paragraph{第~\ref{sec:char-advanced}~節}
\begin{enumerate}
\item 遇到非字母字元（例如空白或標點）時，代表目前累積的單字結束
了。
\item 因為使用者可能多打了字元或整行輸入；不清掉的話，這些殘留字元
會被下一次讀取誤讀，讓程式跳過應有的提示或輸入。
\item 把視窗最前面的字元丟掉，其餘往前移一位，新字元補進最後一個位
置（\texttt{window[0]=window[1]; window[1]=window[2]; window[2]=c;}）。
\end{enumerate}

\paragraph{第~\ref{sec:char-pitfalls}~節}
\begin{enumerate}
\item 常常讀到上一次 \texttt{scanf("\%d")} 留在緩衝區裡的換行符
\texttt{'\textbackslash n'}；修正方式是在第二個 \texttt{scanf} 前先用
\texttt{scanf("\%c", \&dummy)} 或 \texttt{getchar()} 把換行消耗掉。
\item 若系統上 \texttt{char} 是有號的，超過 $127$ 的位元組值轉成
\texttt{int} 後會變成負數，丟進 \texttt{isalpha} 這類函式行為沒有保
證；正確做法是先轉型成 \texttt{(unsigned char)} 再傳入。
\item 因為 \texttt{c - 32} 這類算式只對特定編碼（如 ASCII）成立，且
沒有處理輸入不是字母的情況；標準函式已經封裝好這些細節，比手寫公式
更可靠。
\end{enumerate}

\chapter{字串}
\label{ch:string}

字串是一個約定：一塊 \texttt{char} 陣列，最後一格放 NUL 終止符
\texttt{\textbackslash 0}。 C 沒有獨立的 \texttt{String} 型別，也不會
檢查這個約定有沒有被遵守。 後面每一節都在看同一個約定在不同情境下怎
麼被遵守、又怎麼被打破：宣告時、輸出入時、呼叫函式時、自己動手寫函式
時。 最後用一個簡易 BASIC 直譯器當總驗收，那支程式同時用到矩形陣列存
字串、指標陣列做查表，剛好是本章要建立的核心對比。

\section{字串是一種約定，不是一種型別}
\label{sec:str-convention}

寫下 \texttt{char s[] = "hi"} 時，記憶體裡實際是三個位元組
\texttt{'h'}、\texttt{'i'}，以及最後的 \texttt{'\textbackslash 0'}。
\texttt{printf("\%s", s)}、\texttt{strlen(s)}、\texttt{strcpy} 都靠那
個值為 $0$ 的字元判斷結尾。 沒有人在編譯期或執行期強制檢查它在不在；
忘記放、放錯位置、寫過頭蓋掉它，程式仍會編譯通過，只是印出垃圾、讀過
邊界，或直接當掉。

\begin{code}{c}{忘記加 NUL，printf 會印出垃圾}
char s[3];
s[0] = 'h';
s[1] = 'i';
// 沒有 s[2] = '\0'
printf("%s\n", s);   // 從 s[0] 一路印到碰巧遇到 0 為止
\end{code}

這段程式完全合法。 \texttt{\%s} 不知道陣列有多長，只會往後走，直到看
見一個位元組的值是 $0$。 那個 $0$ 可能在隔壁變數裡，也可能在更遠的地
方。 第 1 節之後談到的 \texttt{sizeof} 對不上 \texttt{strlen}、
\texttt{strcpy} 不檢查邊界、改字串常值會出事，都是同一個原因的不同展
現：語言把結尾交給我們自己守。

順便把兩個長得很像的名字分開。 \textbf{NUL} 是值為 $0$ 的字元，寫成
\texttt{'\textbackslash 0'}，用來標字串結尾。 \textbf{NULL} 是指標的
哨兵值，表示目前不指向任何有效物件。 一個是字元、一個是指標，混用會
讓分類與比較都默默錯下去。

\keypoint{字串 $=$ \texttt{char} 陣列 $+$ NUL 終止符。 C 不檢查這個約
定；每次讀寫都要自己確保結尾在、空間夠。}

\begin{drills}[label=drills:str-convention]
\item \texttt{char s[3] = \{'h', 'i'\};} 少了最後補 \texttt{s[2] = '\textbackslash 0'}，
拿去 \texttt{printf("\%s", s)} 會發生什麼事？
\item NUL 跟 NULL 有什麼不同？ 一個標記什麼結束，一個表示什麼？
\item 為什麼 \texttt{strlen} 讀過邊界的錯誤，編譯器通常抓不出來？
\end{drills}

\section{宣告與儲存位置}
\label{sec:str-declare-storage}

同一個約定，放進不同種類的記憶體，規則就不一樣。 常見有三種：指標指
向字串常值、陣列自己擁有一份可改的複本、\texttt{malloc} 從 heap 要一
塊空間。

\subsection{指標宣告}

\begin{code}{c}{字串指標宣告}
char *p = "hello";    // p 指向唯讀的字串常值
\end{code}

\texttt{p} 自己是一個位址變數，存的是常值 \texttt{"hello"} 的起始位址。
那塊常值通常住在唯讀區，不能改內容。 硬寫 \texttt{p[0] = 'H'} 行為沒
有保證，常常直接 Segmentation Fault。

編譯器還可能把相同的字面值去重、只留一份。 兩個指標都寫
\texttt{"hello"}，執行時可能指向同一塊記憶體。 若其中一個被改到（例如
先轉成非 \texttt{const} 再寫入），另一個看起來毫不相干的字面值也會跟
著變，除錯時非常難查。

\subsection{陣列宣告}

\begin{code}{c}{字串陣列宣告}
char str[100];              // 未初始化，內容為垃圾
char str[] = "hello";       // 編譯器自動計算大小（6）
char str[6] = "hello";      // 明確指定大小
char str[100] = "hello";    // 初始化，剩餘位置為 NUL
\end{code}

陣列在宣告當下就配置連續的 \texttt{char}。 用字面值初始化時，編譯器
會把字元連同最後的 NUL 一併複製進這塊空間，所以陣列內容可以改。 大小
必須至少是字元數加 $1$，留給終止符。

\subsection{堆上分配}

\begin{code}{c}{用 malloc 分配字串}
#include <stdlib.h>
#include <string.h>

char *str = (char *)malloc(100);
if (str == NULL) {
    // malloc 失敗
    return;
}
strcpy(str, "hello");
free(str);
\end{code}

\texttt{malloc} 要來的空間生命週期由我們決定，內容可改，用完必須
\texttt{free}。 完整的動態配置規則在高等技術那一章；這裡只要先記住：
指標指向的那塊記憶體，必須活得比我們使用它的時間久。

\subsection[陣列名字 vs 指標：常數與變數]{陣列名字是位址常數，指標是位址變數}

這是兩種宣告最容易混的地方。 \texttt{char arr[] = "hello"} 的
\texttt{arr} 在編譯期就綁在這塊陣列的起始位址上，不能改指向別處。
\texttt{char *p = "hello"} 的 \texttt{p} 是變數，可以再指向別的字串。

\begin{code}{c}{陣列不能重新指向，指標可以}
char arr[] = "hello";
char *p = "hello";

// arr = "world";   // 編譯錯誤：陣列不是可賦值的左值
p = "world";        // 合法：p 改指向另一個常值
\end{code}

\subsection[陷阱：別用 = 或 == 處理字串]{陷阱：不能用 \texttt{=} 賦值、不能用 \texttt{==} 比較}

初學者寫壞第一支字串程式，多半卡在這兩個直覺。 整數可以 \texttt{x =
3}、\texttt{x == y}；字串看起來也像一份資料，手就會寫出同樣的符號。

\begin{code}{c}{❌ 危險：用 = 把字面值賦給陣列}
char s[10];
s = "abc";    // 編譯錯誤
\end{code}

\begin{code}{c}{✓ 安全：用 strcpy 把字元一個一個複製過去}
char s[10];
strcpy(s, "abc");
\end{code}

\texttt{=} 對陣列做的是改指向，陣列名字改不了指向，所以編譯失敗。 要
的是把字元複製進已有的格子，那是 \texttt{strcpy} 的工作。

\begin{code}{c}{❌ 危險：用 == 比較兩個字串}
char s1[] = "abc";
char s2[] = "abc";
if (s1 == s2) {       // 比的是兩個陣列的起始位址
    // 幾乎一定是假，即使內容相同
}
\end{code}

\begin{code}{c}{✓ 安全：用 strcmp 比較內容}
char s1[] = "abc";
char s2[] = "abc";
if (strcmp(s1, s2) == 0) {
    // 內容相同
}
\end{code}

\texttt{==} 比較的是位址。 兩個陣列各有自己的空間，位址不同，內容再
怎麼一樣也會得到假。 要比內容，用 \texttt{strcmp}：回傳 $0$ 表示相等，
負數表示左邊字典序較小，正數表示較大。

\subsection{空字串}

\begin{example}{c}{(empty-string.c) 空字串的長度是 0}
#include <stdio.h>
#include <string.h>

int main(void) {
    char empty[] = "";         // 陣列大小為 1，只有 NUL
    printf("%zu\n", strlen(empty));
}
\end{example}

\iosample{}{
0 \\
}

空字串長度為 $0$，但仍佔 $1$ 個位元組給 NUL。 指標也可以寫
\texttt{char *emptyPtr = ""}，指向同一種只有終止符的常值。 沒有終止
符的空陣列，跟空字串是兩件事。

\subsection{\texttt{sizeof} 與 \texttt{strlen}}

宣告完立刻就會撞到這組對比。

\begin{example}{c}{(sizeof-strlen.c) 比較 sizeof 與 strlen}
#include <stdio.h>
#include <string.h>

int main(void) {
    char arr[] = "hello";
    printf("%zu\n", sizeof(arr));   // 6，含 NUL
    printf("%zu\n", strlen(arr));   // 5，走到 NUL 之前
}
\end{example}

\iosample{}{
6 \\
5 \\
}

\texttt{sizeof(arr)} 在編譯期就算好整塊陣列的位元組數，含結尾的 NUL。
\texttt{strlen} 執行時從開頭走到第一個 NUL，回傳途中經過幾個字元，不
含終止符。 若改成指標 \texttt{char *p = "hello"}，\texttt{sizeof(p)}
只會得到指標本身的寬度，跟字串多長無關。 複製時預留的空間必須是
\texttt{strlen(src) + 1}。

\begin{code}{c}{計算複製空間}
char dest[100];
char src[] = "hello";

strcpy(dest, src);                    // dest 至少要 6 格
strncpy(dest, src, sizeof(dest) - 1); // 留下最後一格給 NUL
\end{code}

\keypoint{\texttt{strlen} 回傳字元數（不含 NUL），\texttt{sizeof} 回傳
陣列大小（含 NUL）。 複製時預留 $1$ 位元組給終止符。}

\begin{drills}[label=drills:str-declare-storage]
\item \texttt{char *p = "hello";} 跟 \texttt{char arr[] = "hello";}，哪
一個之後可以改指向別的字串？ 哪一個之後可以改內容？
\item \texttt{sizeof(arr)} 跟 \texttt{strlen(arr)} 對同一個
\texttt{char arr[] = "hello"}，各會得到多少？ 差在哪裡？
\item 用 \texttt{malloc} 要來的字串記憶體，用完之後該做什麼？
\end{drills}

\section{字串的輸出與輸入}
\label{sec:str-io}

輸出入函式同樣只認 NUL，不認陣列長度。 選哪一個函式，差在它讀到哪裡
停、會不會自己加換行、會不會檢查空間。

\subsection{輸出：\texttt{printf} 與 \texttt{puts}}

\begin{example}{c}{(printf-puts.c) 用 printf 與 puts 印字串}
#include <stdio.h>

int main(void) {
    char str[] = "hello";
    printf("%s\n", str);
    printf("Name: %s\n", str);
    puts(str);                 // 自動再印一個換行
}
\end{example}

\iosample{}{
hello \\
Name: hello \\
hello \\
}

\texttt{\%s} 從起始位址一路印到 NUL。 字串沒正確終止，就會印過邊界。
\texttt{puts(str)} 等價於 \texttt{printf("\%s\textbackslash n", str)}，
少打幾個字，但一樣靠 NUL 停下來；遇不到 \texttt{'\textbackslash 0'}，
它會一直印到記憶體某處碰巧找到 $0$ 為止。 這跟第 1 節那個忘記加終止
符的例子是同一件事。

\subsection{用 \texttt{scanf} 讀入}

\begin{example}{c}{(scanf-word.c) scanf 只讀到空白為止}
#include <stdio.h>

int main(void) {
    char str[100];
    scanf("%99s", str);        // 最多 99 個字元，留下 1 格給 NUL
    printf("%s\n", str);
}
\end{example}

\iosample{
hello world \\
}{
hello \\
}

\texttt{\%s} 會跳過前導空白，讀到下一個空白為止，所以
\texttt{hello world} 只留下 \texttt{hello}。 \textbf{務必指定長度限
制}；若寫成 \texttt{scanf("\%s", str)}，輸入太長就會寫過陣列盡頭。

早期還有一個看起來更直覺的 \texttt{gets(str)}：讀整行、不檢查長度。
緩衝區溢位的漏洞大到 C11 已經把它從標準裡移除。 現在讀整行，用下面
的 \texttt{fgets}。

\subsection{用 \texttt{fgets} 讀入}

\begin{example}{c}{(fgets-line.c) fgets 讀整行再去掉換行}
#include <stdio.h>
#include <string.h>

int main(void) {
    char str[100];
    if (fgets(str, 100, stdin) != NULL) {
        char *nl = strchr(str, '\n');
        if (nl != NULL) {
            *nl = '\0';        // 把換行換成終止符
        }
        printf("%s\n", str);
    }
}
\end{example}

\iosample{
hello world \\
}{
hello world \\
}

\texttt{fgets} 最多讀 \texttt{size - 1} 個字元，自己補 NUL，所以比
\texttt{scanf("\%s")} 安全。 它會把行尾的 \texttt{'\textbackslash n'}
一併讀進來（緩衝區夠的話）。 後面要比對內容、算長度時，通常先把換行
拿掉。 上面用 \texttt{strchr} 找第一個換行，找到就改成
\texttt{'\textbackslash 0'}；\texttt{strchr} 的完整用法第 5 節再展開。

\keypoint{輸出入都靠 NUL 停下來。 讀入時指定長度；讀整行用
\texttt{fgets}，再用 \texttt{strchr} 清掉換行。}

\begin{drills}[label=drills:str-io]
\item \texttt{scanf("\%s", str)} 讀 \texttt{"hello world"} 時，\texttt{str}
裡最後會是什麼？
\item 為什麼本書建議用 \texttt{scanf("\%99s", str)} 而不是
\texttt{scanf("\%s", str)}？
\item \texttt{fgets} 讀進來的字串，行尾通常還帶著什麼字元？ 要怎麼把它
去掉？
\end{drills}

\section{動手做：自己寫字串函式}
\label{sec:str-diy-functions}

標準函式庫看起來像黑盒子，打開來看就是：從某個位址出發，一格一格走到
NUL。 我們先用下標數某個字元出現幾次，再改成指標，最後推導出兩個到處
都會碰到的寫法。

\subsection{下標版}

\begin{example}{c}{(count-char.c) 用下標計算字元次數}
#include <stdio.h>

int count_char(const char s[], char target) {
    int n = 0;
    for (int i = 0; s[i] != '\0'; i++) {
        if (s[i] == target) {
            n++;
        }
    }
    return n;
}

int main(void) {
    printf("%d\n", count_char("hello", 'l'));
}
\end{example}

\iosample{}{
2 \\
}

\texttt{s[i] != '\textbackslash 0'} 就是在守那個約定：走到終止符就停。
對 \texttt{"hello"} 數 \texttt{'l'}，\texttt{i} 從 $0$ 走到 $4$，命中
兩次，回傳 $2$。

\subsection{指標版}

同一個函式可以消掉索引 \texttt{i}。 \texttt{s[i]} 等價於
\texttt{*(s + i)}；每次 \texttt{i++} 等價於讓指標本身往前一格。

\begin{code}{c}{用指標計算某個字元出現次數}
int count_char(const char *s, char target) {
    int n = 0;
    while (*s) {
        if (*s == target) {
            n++;
        }
        s++;
    }
    return n;
}
\end{code}

\texttt{*s} 在條件裡被當成整數看：NUL 的值是 $0$，條件為假，迴圈停。
非 NUL 的字元值都不是 $0$，條件為真。 這就是下面第一個 idiom 的來源。

\subsection{Idiom 1：找到字串結尾}

\begin{code}{c}{走到 NUL 上停下}
while (*s) {
    s++;
}
// 離開迴圈時 *s 是 '\0'，s 正指在終止符上
\end{code}

每一步先看目前格子有沒有值，有就前進。 停下來時 \texttt{s} 指著
NUL，所以 \texttt{s - start} 正好是字元數，也就是 \texttt{strlen}。

\subsection{Idiom 2：邊複製邊找結尾}

\texttt{strcpy} 的核心是一行：

\begin{code}{c}{strcpy 的核心邏輯}
while (*p++ = *s2++) {
    ;
}
\end{code}

拆開看，一次迴圈做完四件事：

\begin{enumerate}
\item \texttt{*s2++} 取出 \texttt{s2} 目前的字元，然後 \texttt{s2}
      往前一格。
\item 把這個字元寫進 \texttt{*p}，然後 \texttt{p} 往前一格。
\item 賦值運算式本身的值，就是剛被寫進去的那個字元。
\item \texttt{while} 拿這個值當條件：不是 $0$ 就再跑一輪。
\end{enumerate}

複製到 NUL 那一刻，賦值已經完成（目的地也有終止符了），條件剛好為假，
迴圈結束。 不需要再補一行 \texttt{*p = '\textbackslash 0'}，因為
\texttt{'\textbackslash 0'} 自己已經被複製過去。

\subsection{手刻 \texttt{my\_strlen} 與 \texttt{my\_strcpy}}

把兩個 idiom 收成函式，標準庫裡的對應版本做的就是這些事。

\begin{example}{c}{(my-strlen-strcpy.c) 手刻 strlen 與 strcpy}
#include <stdio.h>
#include <stddef.h>

size_t my_strlen(const char *s) {
    const char *start = s;
    while (*s) {
        s++;
    }
    return (size_t)(s - start);
}

void my_strcpy(char *p, const char *s2) {
    while ((*p++ = *s2++)) {   // 連 NUL 一起複製過去
        ;
    }
}

int main(void) {
    char dest[20];
    my_strcpy(dest, "hello");
    printf("%zu\n", my_strlen(dest));
    printf("%s\n", dest);
}
\end{example}

\iosample{}{
5 \\
hello \\
}

\texttt{my\_strlen("hello")} 得到 $5$。 \texttt{my\_strcpy(dest,
"hello")} 之後 \texttt{dest} 裡是 \texttt{h e l l o \textbackslash 0}，
可以交給 \texttt{\%s} 印。 參數寫成 \texttt{const char *} 的，是來源、
函式不會改它；沒有 \texttt{const} 的是目的地、準備被寫入。 下一節標準
函式庫的原型，用的是同一套標記。

\keypoint{\texttt{while (*s)} 走到結尾；\texttt{while (*p++ = *s2++)}
邊複製邊停在 NUL。 標準字串函式就是這兩種走法的包裝。}

\begin{drills}[label=drills:str-diy-functions]
\item \texttt{s[i]} 跟 \texttt{*(s + i)} 是同一件事嗎？ 兩種寫法哪個比
較適合改寫成指標版迴圈？
\item \texttt{while (*p++ = *s2++)} 這一行，複製完最後一個字元之後，為
什麼迴圈會自動停下來，不用另外判斷？
\item \texttt{my\_strlen} 回傳的長度，含不含最後的 NUL？
\end{drills}

\section{標準函式庫}
\label{sec:str-stdlib}

\texttt{string.h} 把上一節那種走訪包裝成一組現成函式。 原型裡標
\texttt{const char *} 的參數是來源，函式保證不改它；沒標
\texttt{const} 的是目的地，會被寫入。 例如
\texttt{strcpy(char *dest, const char *src)}：\texttt{src} 只讀，
\texttt{dest} 要夠大、由呼叫者負責。

\subsection{\texttt{strlen}：字串長度}

\begin{example}{c}{(strlen-demo.c) strlen 用法}
#include <stdio.h>
#include <string.h>

int main(void) {
    char str[] = "hello";
    printf("%zu\n", strlen(str));
}
\end{example}

\iosample{}{
5 \\
}

回傳字元數（不含 NUL）。

\subsection{\texttt{strcpy}：複製字串}

\begin{code}{c}{strcpy 用法（不安全）}
char dest[100];
strcpy(dest, "hello");  // 不檢查 dest 大小
\end{code}

\subsection{\texttt{strncpy}：有長度限制的複製}

\begin{code}{c}{strncpy 用法（較安全）}
char dest[100];
strncpy(dest, "hello", sizeof(dest) - 1);
dest[sizeof(dest) - 1] = '\0';  // 確保 NUL 終止
\end{code}

\texttt{strncpy} 在來源較長時可能不會自己寫 NUL，所以最後一格要手動補上。

\subsection{\texttt{strcat}：字串連接}

\begin{example}{c}{(strcat-demo.c) strcat 用法}
#include <stdio.h>
#include <string.h>

int main(void) {
    char dest[100] = "Hello ";
    strcat(dest, "World");     // 無長度檢查
    printf("%s\n", dest);
}
\end{example}

\iosample{}{
Hello World \\
}

\subsection{\texttt{strncat}：有長度限制的連接}

\begin{code}{c}{strncat 用法}
char dest[100] = "Hello ";
strncat(dest, "World", sizeof(dest) - strlen(dest) - 1);
\end{code}

可追加的長度是目的地剩餘空間再減 $1$，留給終止符。

\subsection{\texttt{strcmp}：字串比較}

\begin{example}{c}{(strcmp-demo.c) strcmp 用法}
#include <stdio.h>
#include <string.h>

int main(void) {
    char s1[] = "apple";
    char s2[] = "banana";

    if (strcmp(s1, s2) < 0) {
        printf("apple < banana\n");
    }
}
\end{example}

\iosample{}{
apple < banana \\
}

回傳負數（s1 較小）、$0$（相等）、正數（s1 較大）。

\subsection{\texttt{strncmp}：有長度限制的比較}

\begin{code}{c}{strncmp 用法}
char s1[] = "hello";
char s2[] = "help";

if (strncmp(s1, s2, 3) == 0) {
    printf("前 3 字元相同\n");
}
\end{code}

\subsection{\texttt{strchr}：查找字元}

第 3 節去掉 \texttt{fgets} 的換行，用的就是它。

\begin{example}{c}{(strchr-demo.c) strchr 用法}
#include <stdio.h>
#include <string.h>

int main(void) {
    char str[] = "hello world";
    char *p = strchr(str, 'o');
    if (p != NULL) {
        printf("找到 o 在位置 %ld\n", p - str);
    }
}
\end{example}

\iosample{}{
找到 o 在位置 4 \\
}

回傳第一次出現的位址；找不到則回傳 \texttt{NULL}。
\texttt{p - str} 是從開頭算起的位移，單位是字元個數。

\subsection{\texttt{strrchr}：查找最後一個字元}

\begin{code}{c}{strrchr 用法}
char str[] = "hello.txt";
char *p = strrchr(str, '.');  // 找最後的點
\end{code}

\subsection{\texttt{strstr}：查找子字串}

\begin{example}{c}{(strstr-demo.c) strstr 用法}
#include <stdio.h>
#include <string.h>

int main(void) {
    char str[] = "hello world";
    char *p = strstr(str, "wor");
    if (p != NULL) {
        printf("找到 wor 在位置 %ld\n", p - str);
    }
}
\end{example}

\iosample{}{
找到 wor 在位置 6 \\
}

找不到同樣回傳 \texttt{NULL}。

\subsection{\texttt{strtok}：分詞}

\begin{example}{c}{(strtok-demo.c) strtok 用法}
#include <stdio.h>
#include <string.h>

int main(void) {
    char str[] = "apple,banana,orange";
    char *token = strtok(str, ",");   // 第一個詞
    while (token != NULL) {
        printf("%s\n", token);
        token = strtok(NULL, ",");    // 後續詞
    }
}
\end{example}

\iosample{}{
apple \\
banana \\
orange \\
}

\texttt{strtok} 會修改原字串，把分隔符改成 NUL。 若需保留原文，先複製
一份再切。

\subsection{\texttt{strspn} 與 \texttt{strcspn}}

\begin{code}{c}{strspn 與 strcspn}
char str[] = "123abc";
int n = strspn(str, "0123456789");  // 3（前三個都是數字）
n = strcspn(str, "abc");            // 3（前三個不含 a、b、c）
\end{code}

\subsection{\texttt{sprintf} 與 \texttt{snprintf}：組成字串}

有時要把好幾個值拼成一行訊息，存在陣列裡以後再用。
\texttt{sprintf} 跟 \texttt{printf} 格式相同，只是寫進字串而不是螢幕；
它不檢查目的地有多長。 習慣上改用 \texttt{snprintf}，把陣列大小一併傳
進去。

\begin{example}{c}{(snprintf-demo.c) 用 snprintf 組成訊息}
#include <stdio.h>

int main(void) {
    char name[] = "Ann";
    int score = 95;
    char msg[80];

    snprintf(msg, sizeof(msg), "%s 得了 %d 分", name, score);
    puts(msg);
}
\end{example}

\iosample{}{
Ann 得了 95 分 \\
}

\keypoint{來源參數標 \texttt{const}，目的地由呼叫者準備夠大的緩衝區。
找不到時 \texttt{strchr}／\texttt{strstr} 回傳 \texttt{NULL}，先檢查再
解參考。}

\begin{drills}[label=drills:str-stdlib]
\item \texttt{strcmp(s1, s2)} 回傳 $0$ 代表什麼？ 回傳負數呢？
\item \texttt{strchr} 找不到字元時，回傳值是什麼？ 用它的回傳值之前該
先做什麼檢查？
\item \texttt{strtok} 分詞時，會不會動到原本傳進去的字串？ 若要保留原
文，該怎麼做？
\item \texttt{sprintf} 跟 \texttt{snprintf} 的差別是什麼？ 為什麼本節建
議優先用 \texttt{snprintf}？
\end{drills}

\section{字串陣列：矩形 vs 鋸齒}
\label{sec:str-array-rect-jagged}

一組字串可以有兩種排法。 一種是每列固定長度的二維陣列，一種是指標陣
列、各指向長短不一的字串。 最後一節的直譯器兩種都用。

\subsection{矩形陣列}

\begin{example}{c}{(names-rect.c) 矩形陣列存放多個名字}
#include <stdio.h>
#include <string.h>

int main(void) {
    char names[5][40];
    strcpy(names[0], "Ann");
    strcpy(names[1], "Bob");
    names[0][0] = 'a';     // 內容可改
    printf("%s\n", names[0]);
    printf("%s\n", names[1]);
}
\end{example}

\iosample{}{
ann \\
Bob \\
}

\texttt{names[5][40]} 在記憶體裡是連續 $5 \times 40$ 個
\texttt{char}。 每列都有自己的空間，可以修改，所以第一個名字從
\texttt{Ann} 變成 \texttt{ann}。 短名字後面會墊一堆 NUL，比較浪費。

\begin{code}{text}{矩形：每列一樣長}
names[0]  A n n \0  ...（補到 40）
names[1]  B o b \0  ...（補到 40）
\end{code}

\subsection{鋸齒陣列}

\begin{code}{c}{指標陣列，各指向不同長度的字串}
char *table[5] = {"Ann", "Bob", "Cara", "Dan", "Eve"};
// table[0][0] = 'a';   // 若指向字面值，改內容行為沒有保證
table[0] = "Anna";      // 可以改指向
\end{code}

每個元素是一個 \texttt{char *}，指向的字串可長可短，省空間。 若指向
字面值，內容通常不能改；若指向自己的陣列或 \texttt{malloc} 出來的空
間，就可以改。

\begin{code}{text}{鋸齒：指標指向各自的結尾}
table[0] --> A n n \0
table[1] --> B o b \0
table[2] --> C a r a \0
\end{code}

第 9 節的直譯器正好並存這兩種：\texttt{varName} 是
\texttt{char varName[VAR][NAMELENGTH]}，用來真正存放變數名字元；
\texttt{varNameTable} 是 \texttt{char *varNameTable[VAR]}，每個格子指
向 \texttt{varName} 的某一列，查表時交給 \texttt{strcmp}。 矩形負責擁
有記憶體，鋸齒負責快速指來指去。

排序一組字串時，鋸齒陣列可以只交換指標、不必搬動字元本身；完整的排序
範例留給之後的演算法章節。

\keypoint{\texttt{char a[M][N]} 是矩形、列定長、內容可改。
\texttt{char *a[M]} 是鋸齒、省空間，指向常值時不能改內容。}

\begin{drills}[label=drills:str-array-rect-jagged]
\item \texttt{char names[5][40]} 這種矩形陣列，每一列的長度是固定的還
是各自不同？
\item \texttt{char *table[5] = \{"Ann", "Bob", ...\}} 這種鋸齒陣列，
\texttt{table[0] = "Anna";} 合法嗎？ \texttt{table[0][0] = 'a';} 呢？
\item 矩形陣列跟鋸齒陣列，哪一種比較省記憶體？ 為什麼？
\end{drills}

\section{傳遞與回傳字串}
\label{sec:str-pass-return}

函式邊界上，字串參數看起來像陣列，傳進去的其實是起始位址。 不必再傳
長度，除非函式裡要寫入、必須知道緩衝區有多大。 這就是直譯器裡
\texttt{char *varNameTable[]} 能在多個函式之間傳來傳去的前提。

回傳字串時，指向的記憶體必須在函式結束後仍有效。

\begin{itemize}
\item 回傳字面值：常值活到程式結束，安全。
\item 回傳 \texttt{static} 陣列：安全，但下次呼叫會覆蓋上次的內容。
\item 回傳 \texttt{malloc} 的區塊：安全，呼叫者負責 \texttt{free}。
\item 回傳區域陣列的位址：函式一結束那塊 stack 就失效，之後再讀行為
      沒有保證。
\end{itemize}

\begin{code}{c}{❌ 危險：回傳區域陣列的位址}
char *makeTag(int n) {
    char buf[16];
    sprintf(buf, "TAG-%d", n);
    return buf;   // buf 在函式結束後已經無效
}
\end{code}

\begin{example}{c}{(make-tag.c) 呼叫者提供緩衝區寫入標籤}
#include <stdio.h>
#include <stddef.h>

void makeTag(int n, char *buf, size_t size) {
    snprintf(buf, size, "TAG-%d", n);
}

int main(void) {
    char tag[16];
    makeTag(7, tag, sizeof(tag));
    printf("%s\n", tag);
}
\end{example}

\iosample{}{
TAG-7 \\
}

呼叫端自己準備 \texttt{char tag[16]}，把位址與大小傳進去。 函式寫完
就結束，不需要回傳指標。 第 9 節的 \texttt{readVariable}、
\texttt{processCode} 都是這種模式：緩衝區由 \texttt{main} 準備，函式
只沿著傳進來的位址讀寫。

\keypoint{傳字串就是傳位址。 回傳時，指向的物件必須活得比這次呼叫久；
區域陣列不行。}

\begin{drills}[label=drills:str-pass-return]
\item 為什麼 \texttt{char *makeTag(int n) \{ char buf[16]; ...; return
buf; \}} 是危險的寫法？ 函式結束後 \texttt{buf} 發生了什麼事？
\item 要讓函式把結果寫回呼叫端，範例 \texttt{make-tag.c} 用的是哪一種
做法：回傳指標，還是把緩衝區位址傳進去？
\item 回傳 \texttt{static} 陣列安不安全？ 連續呼叫兩次會發生什麼事？
\end{drills}

\section{常見陷阱與最佳實踐}
\label{sec:str-pitfalls}

前幾節談過的約定、賦值、生命週期，這裡收成一張對照表。

\subsection{陷阱 1：緩衝區溢出}

\begin{code}{c}{❌ 危險}
char dest[10];
strcpy(dest, "very long string");  // 溢出
\end{code}

\begin{code}{c}{✓ 安全}
char dest[10];
strncpy(dest, "very long string", sizeof(dest) - 1);
dest[sizeof(dest) - 1] = '\0';
\end{code}

\subsection{陷阱 2：忘記 NUL 終止符}

\begin{code}{c}{❌ 危險}
char dest[10];
for (int i = 0; i < 10; i++) {
    dest[i] = 'a';
}
// 沒有 NUL，strlen 與 %s 都會越界
\end{code}

\subsection{陷阱 3：\texttt{strtok} 修改原字串}

\begin{code}{c}{❌ 危險}
char str[] = "a,b,c";
strtok(str, ",");
printf("%s\n", str);  // 只剩 a
\end{code}

\begin{code}{c}{✓ 安全}
char str[] = "a,b,c";
char copy[100];
strcpy(copy, str);
strtok(copy, ",");     // 修改複本
\end{code}

\subsection{陷阱 4：用 \texttt{=} 賦值、用 \texttt{==} 比較}

第 2 節的核心陷阱，這裡再列一次當複習。

\begin{code}{c}{❌ 危險}
char s[10];
s = "abc";            // 編譯錯誤
if (s == "abc") {     // 比位址
}
\end{code}

\begin{code}{c}{✓ 安全}
char s[10];
strcpy(s, "abc");
if (strcmp(s, "abc") == 0) {
}
\end{code}

\subsection{陷阱 5：回傳區域字串位址}

第 7 節的地雷，再看一眼簽名就知道。

\begin{code}{c}{❌ 危險}
char *label(void) {
    char buf[8] = "ok";
    return buf;
}
\end{code}

\begin{code}{c}{✓ 安全}
void label(char *buf, size_t size) {
    snprintf(buf, size, "ok");
}
\end{code}

\subsection{最佳實踐}

\begin{itemize}
\item \textbf{用 strncpy 和 strncat}，永不用 strcpy 和 strcat
\item \textbf{檢查 malloc 回傳值}，防止 NULL 指標
\item \textbf{字串賦值前複製}，若需保留原文
\item \textbf{明確檢查邊界}，尤其是陣列大小計算
\item \textbf{用 sizeof 陣列得大小}，但記得 sizeof 指標只給指標大小
\end{itemize}

\keypoint{字串處理的核心是記住：字串 $=$ 陣列 $+$ NUL 終止符。 每次操
作都要確保足夠空間與正確終止。}

\begin{drills}[label=drills:str-pitfalls]
\item 陷阱 1 裡 \texttt{strcpy(dest, "very long string")} 若 \texttt{dest}
只有 10 格，會發生什麼事？ 換成 \texttt{strncpy} 版本要注意補上什麼？
\item 陷阱 3 為什麼說 \texttt{strtok} 會「修改原字串」？ 要保留原文該
怎麼做？
\item 陷阱 4 的 \texttt{if (s == "abc")} 比較的是內容還是位址？ 正確寫
法該用哪個函式？
\end{drills}

\section{簡易 BASIC 直譯器}
\label{sec:str-basic-interp}

作為字串處理的綜合應用，我們實作一個簡易的 BASIC 直譯器。 它展示如何
用字串函式解析並執行一個小語言：讀 token、比對指令名稱、查變數表、再
依操作碼執行。

\subsection{設計}

BASIC 直譯器支援以下語句：
\begin{itemize}
\item \texttt{VAR ... END}：變數聲明（輸入時寫成一連串
      \texttt{名稱 = 值}，以 \texttt{END} 結束）
\item \texttt{ASSIGN}：變數賦值
\item \texttt{IF}：條件跳轉
\item \texttt{GOTO}：無條件跳轉
\item \texttt{PRINT}：輸出變數
\item \texttt{STOP}：停止程式
\end{itemize}

\subsection{核心函式說明}

\textbf{readVariable}：解析變數聲明區段

\begin{code}{c}{readVariable 簽名}
void readVariable(char varName[][NAMELENGTH],
                  int varValue[],
                  char *varNameTable[]);
\end{code}

讀入變數名和初始值，建立查詢表。

\textbf{processCode}：解析程式碼

\begin{code}{c}{processCode 簽名}
void processCode(char *varNameTable[],
                 int code[], int opd1[],
                 int opd2[], int compopr[],
                 int gotoLine[],
                 int target[]);
\end{code}

逐行解析 BASIC 語句，提取操作碼和操作數。

\textbf{runCode}：執行程式

\begin{code}{c}{runCode 簽名}
void runCode(char varName[][NAMELENGTH],
             int varValue[], int code[],
             int opd1[], int opd2[],
             int compopr[], int gotoLine[],
             int target[]);
\end{code}

根據解析結果逐行執行，維護變數值和程式計數器。

完整實作包括以下函式：
\begin{itemize}
\item 查表：\texttt{find}、\texttt{findOprCode}、\texttt{findVariable}
\item 解析：\texttt{processIF}、\texttt{processASSIGN}、\texttt{processCode}
\item 執行：\texttt{compare}、\texttt{arithmetic}、\texttt{runCode}
\item 入口：\texttt{readVariable}、\texttt{main}
\end{itemize}

\begin{code}{c}{BASIC 直譯器完整原始碼}
#include <stdio.h>
#include <string.h>
#include <assert.h> /* for assert */
#include <stdlib.h> /* for exit */
/* define */
/* max # of var */
#define VAR 100
/* max length of var name */
#define NAMELENGTH 80
/* max # of lines of code */
#define LINE 200

enum Operator
  {ADD, SUB, MUL, DIV, MOD, EQL, NOTEQL,
   SMALLER, LARGER, SMALLEREQL, LARGEREQL};

enum StatementType
  {STOP, IF, GOTO, PRINT, ASSIGN};

/* find */
int find(char *string, char *table[])
{
  int index = 0;
  while (table[index] != NULL) {
    if (strcmp(string, table[index]) == 0)
      return index;
    index++;
  }
  return(-1); /* not found */
}

/* findOprCode */
int findOprCode(char *name)
{
  char *operators[] = {"+", "-", "*", "/", "%",
                       "==", "!=", "<", ">",
                       "<=", ">=", NULL};
  int code = find(name, operators);
  if (code == -1)
    exit(-1);
  return(code);
}

/* findVariable */
int findVariable (char *name,
                  char *varNameTable[])
{
  int index = find(name, varNameTable);
  if (index == -1)
    exit(-1);
  return(index);
}

/* processIF */
void processIF(int lineNum, char *varNameTable[],
               int code[], int opd1[],
               int opd2[], int compopr[],
               int gotoLine[])
{
  char gotoStr[NAMELENGTH];
  char opd[NAMELENGTH], opr[NAMELENGTH];

  code[lineNum] = IF;
  scanf("%s", opd);
  opd1[lineNum] = 
    findVariable(opd, varNameTable);
  scanf("%s", opr);
  compopr[lineNum] = findOprCode(opr);
  scanf("%s", opd);
  opd2[lineNum] = 
    findVariable(opd, varNameTable);
  scanf("%s", gotoStr);
  assert(strcmp(gotoStr, "GOTO") == 0);
  scanf("%d", &(gotoLine[lineNum]));
}

/* processASSIGN */
void processASSIGN(int lineNum, char *first,
                   char *varNameTable[],
                   int code[], int opd1[],
                   int opd2[], int compopr[],
                   int target[])
{
  char opd[NAMELENGTH], opr[NAMELENGTH];
  code[lineNum] = ASSIGN;
  target[lineNum] = 
    findVariable(first, varNameTable);
  char assignStr[NAMELENGTH];
  scanf("%s", assignStr);
  assert(strcmp(assignStr, "=") == 0);
  scanf("%s", opd);
  opd1[lineNum] = 
    findVariable(opd, varNameTable);
  scanf("%s", opr);
  compopr[lineNum] = findOprCode (opr);
  scanf("%s", opd);
  opd2[lineNum] = 
    findVariable(opd, varNameTable);
}

/* processCode */
void processCode(char *varNameTable[],
                 int code[], int opd1[],
                 int opd2[], int compopr[],
                 int gotoLine[],
                 int target[])
{
  char first[NAMELENGTH];
  char opd[NAMELENGTH];
  int lineNum = 2; /* 1 is for var */
  char *firstTable[] = 
    {"STOP", "IF", "GOTO", "PRINT", NULL};
  int firstTableSize = 
    sizeof(firstTable) / sizeof(char *);
  while (scanf("%s", first) != EOF) {
    int firstType = find(first, firstTable);
    switch (firstType) {
    case GOTO:
      code[lineNum] = GOTO;
      scanf("%d", &(gotoLine[lineNum]));
      break;
    case STOP:
      code[lineNum] = STOP;
      break;
    case IF:
      code[lineNum] = IF;
      processIF(lineNum, varNameTable, code,
                opd1, opd2, compopr, gotoLine);
      break;
    case PRINT:
      code[lineNum] = PRINT;
      scanf("%s", opd);
      target[lineNum] = 
        findVariable(opd, varNameTable);
      break;
    default: /* must be assignment */
      code[lineNum] = ASSIGN;
      processASSIGN(lineNum, first, varNameTable,
                    code, opd1, opd2, compopr,
                    target);
      break;
    }
    lineNum++;
  }
}

/* compare */
int compare(int operand1, int operand2,
            int operator)
{
  switch (operator) {
  case EQL:
    return (operand1 == operand2);
  case NOTEQL:
    return (operand1 != operand2);
  case SMALLER:
    return (operand1 < operand2);
  case LARGER:
    return (operand1 > operand2);
  case SMALLEREQL:
    return (operand1 <= operand2);
  case LARGEREQL:
    return (operand1 >= operand2);
  default:
    printf("invalid comp %d\n", operator);
    exit(-1);
  }
}

/* arithmetic */
int arithmetic(int operand1, int operand2,
               int operator)
{
  switch (operator) {
  case ADD:
    return (operand1 + operand2);
  case SUB:
    return (operand1 - operand2);
  case MUL:
    return (operand1 * operand2);
  case DIV:
    return (operand1 / operand2);
  case MOD:
    return (operand1 % operand2);
  default:
    printf("invalid comp %d\n", operator);
    exit(-1);
  }
}

/* runCode */
void runCode(char varName[][NAMELENGTH],
             int varValue[], int code[],
             int opd1[], int opd2[],
             int compopr[], int gotoLine[],
             int target[])
{
  int line = 2; /* starts at 2 */
  while (code[line] != STOP) {
    switch (code[line]) {
    case IF:
      if (compare(varValue[opd1[line]],
                  varValue[opd2[line]],
                  compopr[line]))
        line = gotoLine[line];
      else
        line++;
      break;
    case GOTO:
      line = gotoLine[line];
      break;
    case ASSIGN:
      varValue[target[line]] = 
        arithmetic(varValue[opd1[line]],
                   varValue[opd2[line]],
                   compopr[line]);
      line++;
      break;
    case PRINT:
      printf("%d\n", varValue[target[line]]);
      line++;
      break;
    default:
      printf("invalid code %d\n", code[line]);
      exit(-1);
    }
  }
}

/* readVariable */
void readVariable(char varName[][NAMELENGTH],
                  int varValue[],
                  char *varNameTable[])
{
  int count = 0;
  char name [NAMELENGTH];

  scanf("%s", name);
  while (strcmp(name, "END") != 0) {
    char assignStr[NAMELENGTH];
    strcpy(varName[count], name);
    scanf("%s", assignStr);
    assert(strcmp(assignStr, "=") == 0);
    scanf("%d", &(varValue[count]));
    varNameTable[count] = varName[count];
    count++;
    scanf("%s", name);
  }
  varNameTable[count] = NULL;
}

/* main */
int main ()
{
  char varName [VAR] [NAMELENGTH];
  int varValue [VAR];
  char *varNameTable [VAR];
  /* get variables */
  readVariable(varName, varValue, varNameTable);
  int code [LINE];
  int opd1 [LINE], compopr[LINE], opd2[LINE];
  int gotoLine[LINE]; /* IF */
  int target[LINE]; /* both PRINT and ASSIGN */
  /* process code */
  processCode(varNameTable, code, opd1,
              opd2, compopr, gotoLine, target);
  /* run code */
  runCode(varName, varValue, code, opd1,
          opd2, compopr, gotoLine, target);
  return 0;
}
/* end */
\end{code}

\subsection{用 BASIC 求質數的程式範例}

以下 BASIC 程式找出 $2$ 到 $100$ 之間的質數。 直譯器要處理變數宣告、
算術、比較、條件跳轉。 輸入時第 1 行是變數區（以 \texttt{END} 結束），
後面各行對應程式行號 $2$ 起：

\begin{code}{c}{質數判斷程式（BASIC 語法）}
1 N = 2 F = 2 ONE = 1 TWO = 2 R = 0 S = 0 ZERO = 0 HUNDRED = 100 END
2 IF N > HUNDRED GOTO 12
3 F = TWO + ZERO
4 R = N % F
5 IF R == ZERO GOTO 10
6 F = F + ONE
7 S = F * F
8 IF S <= N GOTO 4
9 PRINT N
10 N = N + ONE
11 GOTO 2
12 STOP
\end{code}

執行流程：
\begin{enumerate}
\item 第 1 行：初始化變數（N$=2$，F$=2$，ONE$=1$ 等）
\item 第 2 行：檢查 N 是否超過 $100$，若是則跳至停止
\item 第 3--8 行：逐一檢查 $2$ 到 $\sqrt{N}$ 的數字是否能整除 N
\item 第 9 行：若沒有因子，印出 N（質數）
\item 第 10--11 行：N 加 $1$，繼續迴圈
\end{enumerate}

直譯器讀入時沒有行號 token，變數區與各行指令用空白分開即可：

\iosample{
N = 2 F = 2 ONE = 1 TWO = 2 R = 0 S = 0 ZERO = 0 HUNDRED = 100 END \\
IF N > HUNDRED GOTO 12 \\
F = TWO + ZERO \\
R = N \% F \\
IF R == ZERO GOTO 10 \\
F = F + ONE \\
S = F * F \\
IF S <= N GOTO 4 \\
PRINT N \\
N = N + ONE \\
GOTO 2 \\
STOP \\
}{
3 \\
5 \\
7 \\
11 \\
13 \\
17 \\
19 \\
23 \\
29 \\
31 \\
37 \\
41 \\
43 \\
47 \\
53 \\
59 \\
61 \\
67 \\
71 \\
73 \\
79 \\
83 \\
89 \\
97 \\
}

起始時 \texttt{N} 與 \texttt{F} 都是 $2$，第 4 行算出
\texttt{N \% F} 為 $0$，第 5 行就跳到第 10 行把 \texttt{N} 加 $1$，
所以 $2$ 不會被印出來。 從 $3$ 開始，\texttt{F} 每次重設成 $2$，小
於 \texttt{N} 的因子才會觸發那條跳轉。 輸出是 $3$ 到 $97$ 的質數。

這個程式展示了直譯器需要處理的典型複雜度：
\begin{itemize}
\item \textbf{變數管理}：宣告、查詢、修改多個變數
\item \textbf{字串解析}：用空白把一行拆成 token
\item \textbf{字元查詢}：用 \texttt{strcmp} 匹配指令名稱
\item \textbf{運算子解析}：識別 \texttt{+}、\texttt{-}、\texttt{*}、
      \texttt{\%}、\texttt{>}、\texttt{<=}、\texttt{==} 等
\item \textbf{控制流}：管理程式計數器、條件跳轉
\end{itemize}

\subsection[回頭指認：對照前面學過的內容]{回頭指認：這支程式用到前面哪幾節}

讀完整支直譯器之後，可以對回本章的主軸：每一處都在遵守或利用字串那條
約定。

\begin{itemize}
\item \texttt{varName} 是 \texttt{char varName[VAR][NAMELENGTH]}，矩
      形陣列，真正存變數名字元；\texttt{varNameTable} 是
      \texttt{char *varNameTable[VAR]}，鋸齒式指標陣列，每個元素指向
      \texttt{varName} 的一列，最後用 \texttt{NULL} 收尾。 這就是第 6
      節的兩種模式並存：矩形負責擁有記憶體，指標陣列負責查表。
\item \texttt{find}、\texttt{findVariable}、\texttt{findOprCode} 都是
      對一個 \texttt{NULL} 結尾的指標陣列做 \texttt{strcmp}，一個一個
      比到命中為止。 這是第 5 節 \texttt{strcmp} 的手刻查表，沒有用雜
      湊或更複雜的結構。
\item 變數區與每一句指令都靠 \texttt{scanf("\%s", ...)} 讀下一個
      token，空白當分隔。 這是第 3 節的讀法：\texttt{\%s} 讀到空白為
      止。
\item 所有函式簽名裡的 \texttt{char varName[][NAMELENGTH]}、
      \texttt{char *varNameTable[]}，傳的都是起始位址。 \texttt{main}
      準備好兩塊陣列，其餘函式只沿著指標讀寫，沒有回傳區域字串。 這
      是第 7 節的傳遞方式。
\end{itemize}

直譯器示範了字串處理在編譯器與直譯器裡的角色：從源碼讀入、匹配指令、
分割 token、查表，每一步都靠同一條約定：走到 NUL 就停，比內容用
\texttt{strcmp}，傳的是位址而不是複本。

\keypoint{矩形陣列存字元、指標陣列做查表、\texttt{strcmp} 逐一比對、
\texttt{scanf("\%s")} 讀 token、函式之間只傳位址。 直譯器把本章每一節
串在同一支程式裡。}

\begin{drills}[label=drills:str-basic-interp]
\item 質數程式跑到 \texttt{N = 2} 時，第 9 行的 \texttt{PRINT N} 有沒有
執行？ 為什麼 $2$ 不會出現在輸出裡？
\item 直譯器裡 \texttt{varName} 跟 \texttt{varNameTable} 各是矩形陣列還
是鋸齒陣列？ 誰真正擁有字元的記憶體？
\item \texttt{find}、\texttt{findVariable} 這類查表函式，靠哪個字串函
式判斷有沒有找到？
\end{drills}

\addcontentsline{toc}{section}{練習解答}
\section*{練習解答}

各節練習的答案收在這裡。 短題只給能一句話核對的結果。

\paragraph{第~\ref{sec:str-convention}~節}
\begin{enumerate}
\item \texttt{printf} 會從 \texttt{s[0]} 開始一路印，直到記憶體某處碰
巧出現值為 $0$ 的位元組為止——可能印出垃圾字元，甚至讀過陣列邊界。
\item NUL（\texttt{'\textbackslash 0'}）是值為 $0$ 的字元，標記字串結
尾；NULL 是指標的哨兵值，表示不指向任何有效物件。
\item 因為語言本身不檢查這個約定；\texttt{\%s} 只是往後走找 $0$，忘記
放終止符仍然合法編譯，只是執行時行為不對。
\end{enumerate}

\paragraph{第~\ref{sec:str-declare-storage}~節}
\begin{enumerate}
\item \texttt{p} 之後可以改指向別的字串（它是一個指標變數）；
\texttt{arr} 之後可以改內容（它是自己擁有的、可寫入的陣列），但不能
改指向。
\item \texttt{sizeof(arr)} 得到 $6$（含最後的 NUL），\texttt{strlen(arr)}
得到 $5$（不含 NUL，只算走到終止符前的字元數）。
\item 用 \texttt{free} 把它還給系統，避免記憶體洩漏。
\end{enumerate}

\paragraph{第~\ref{sec:str-io}~節}
\begin{enumerate}
\item 只會存到 \texttt{hello}：\texttt{\%s} 讀到空白就停。
\item 因為沒有長度限制的 \texttt{scanf("\%s", str)}，若輸入超過緩衝區
大小，會寫過陣列盡頭；加上長度限制（如 \texttt{\%99s}）能確保最多寫入
99 個字元、留 1 格給 NUL。
\item 通常還帶著換行字元 \texttt{'\textbackslash n'}；可以用
\texttt{strchr} 找到它，改成 \texttt{'\textbackslash 0'} 去掉。
\end{enumerate}

\paragraph{第~\ref{sec:str-diy-functions}~節}
\begin{enumerate}
\item 是同一件事；指標版迴圈用 \texttt{*s} 搭配 \texttt{s++}，改寫起來
更直接。
\item 因為賦值運算式的值就是被寫入的那個字元；複製到 NUL 時，NUL 本
身（值為 $0$）被賦值、也被當成條件判斷，條件為假，迴圈自然結束。
\item 不含最後的 NUL，只回傳走過的字元數。
\end{enumerate}

\paragraph{第~\ref{sec:str-stdlib}~節}
\begin{enumerate}
\item $0$ 代表兩字串內容相同；負數代表 \texttt{s1} 字典序比 \texttt{s2}
小。
\item 回傳 \texttt{NULL}；用它之前要先檢查是否為 \texttt{NULL}，再解參
考，否則會存取無效位址。
\item 會：\texttt{strtok} 把分隔符位置改成 NUL，直接修改原字串；要保留
原文，得先複製一份再對複本呼叫 \texttt{strtok}。
\item \texttt{sprintf} 不檢查目的地大小，\texttt{snprintf} 多一個大小參
數、不會寫過邊界；本節建議優先用 \texttt{snprintf} 是為了避免緩衝區溢
出。
\end{enumerate}

\paragraph{第~\ref{sec:str-array-rect-jagged}~節}
\begin{enumerate}
\item 固定的：每一列都是同樣的長度（這裡是 40）。
\item \texttt{table[0] = "Anna";} 合法（改指向）；\texttt{table[0][0] =
'a';} 若指向的是字面值，行為沒有保證，通常不安全。
\item 鋸齒陣列通常比較省空間，因為每個字串只佔剛好夠用的長度，不會像
矩形陣列那樣把短字串墊滿到固定寬度。
\end{enumerate}

\paragraph{第~\ref{sec:str-pass-return}~節}
\begin{enumerate}
\item \texttt{buf} 是函式裡的區域陣列，函式一結束它在堆疊上的空間就失
效；回傳的位址指向一塊已經失效的記憶體，之後讀取行為沒有保證。
\item 把緩衝區位址（連同大小）傳進去，由呼叫端 \texttt{main} 準備空間，
函式只負責寫入，不回傳指標。
\item 安全，但下次呼叫會覆蓋上次的內容；連續呼叫兩次、把兩次的回傳指
標都留著用，會發現兩個指標最後指向同一份、且是第二次呼叫的內容。
\end{enumerate}

\paragraph{第~\ref{sec:str-pitfalls}~節}
\begin{enumerate}
\item 會寫過 \texttt{dest} 的邊界，造成緩衝區溢位；換成 \texttt{strncpy}
之後，若來源字串跟緩衝區一樣長或更長，\texttt{strncpy} 可能不會自動補
NUL，要自己在最後一格補上 \texttt{'\textbackslash 0'}。
\item 因為 \texttt{strtok} 會把分隔符原地改成 NUL，直接修改傳進去的字
串本身；要保留原文，得先用 \texttt{strcpy} 複製一份，再對複本呼叫
\texttt{strtok}。
\item 比較的是位址（兩個陣列的起始位址），不是內容；正確寫法要用
\texttt{strcmp(s, "abc") == 0} 比較內容。
\end{enumerate}

\paragraph{第~\ref{sec:str-basic-interp}~節}
\begin{enumerate}
\item 沒有執行：\texttt{N = 2} 時第 4 行算出 \texttt{N \% F} 為 $0$，
第 5 行的 \texttt{IF R == ZERO GOTO 10} 成立，直接跳到第 10 行把 N 加
1，第 9 行的 \texttt{PRINT} 被跳過。
\item \texttt{varName} 是矩形陣列（\texttt{char varName[VAR][NAMELENGTH]}），
真正擁有字元的記憶體；\texttt{varNameTable} 是鋸齒式指標陣列
（\texttt{char *varNameTable[VAR]}），只是指向 \texttt{varName} 裡的
某一列，不擁有記憶體。
\item 靠 \texttt{strcmp}：對 \texttt{NULL} 結尾的指標陣列逐一比對，找
到內容相同的字串就算命中。
\end{enumerate}

\chapter{位元操作}
\label{ch:bitoperations}

位元運算是 C 語言最直接、最高效的操作。 從檢查單一位的狀態，到用位移實
現快速乘除，再到用位元操作交換兩數：這些都是硬體層、密碼學、遊戲引擎常
用的技能。 這章從最基礎的二進位表示開始，逐步講到 5 個核心習語：位元
且、位元或、位元互補、位元互斥或，以及移位。

\section{基礎：進制與二進位表示}
\label{sec:bit-basics}

電腦用二進位儲存資料。 十進位的每一位代表 $10^n$，二進位的每一位代表
$2^n$。 寫下 \texttt{int x = 5;} 時，記憶體裡那一格是
\texttt{00000101}（8 位二進位），也可以寫成 \texttt{0x05}（十六進位）。
後面每一個位元運算子，都是在這一排 $0$ 與 $1$ 上面做事。

\subsection{進制轉換}

C 的整數字面值預設是十進位。 八進位前面加 $0$，十六進位前面加
\texttt{0x}。 印出來時用不同的格式化字串，同一個值可以換成三種進位。

\begin{code}{c}{進制轉換}
int x = 13;      // decimal
// 13 in binary:  1101
// 13 in octal:   015
// 13 in hex:     0x0D

printf("%d\n", x);    // decimal: 13
printf("%o\n", x);    // octal: 15
printf("%x\n", x);    // hex: d
\end{code}

十進位換成二進位，就是把這個數拆成若干個 $2$ 的次方再相加。 $13 = 8 +
4 + 1 = 2^3 + 2^2 + 2^0$，所以二進位是 \texttt{1101}。 反過來，看到
\texttt{1101}，就把亮著的權重加起來：$8 + 4 + 1 = 13$。

後面每一節都要講第 $n$ 位：檢查它、設成 $1$、清成 $0$、切換它。 $n$
從哪一端數、哪一格叫 bit $0$，若不先釘死，\texttt{1u << n} 就對不準
那一格。 下面這張圖把編號、權重、MSB／LSB 一次講完；後文說檢查第
$n$ 位、設置第 $n$ 位，指的都是這個編號。

\subsection{位元的編號、MSB 與權重}

一個位元組有 8 個位，由左到右編號 $7$ 到 $0$。 最左邊是 MSB（most
significant bit，最高有效位），權重最大；最右邊是 LSB（least
significant bit，最低有效位），權重最小。 第 $n$ 位為 $1$ 時，貢獻的
值是 $2^n$。

\begin{center}
\begin{tikzpicture}
    \node[cell] (b7) {0};
    \node[hlcell, right=0pt of b7] (b6) {1};
    \node[cell, right=0pt of b6] (b5) {0};
    \node[cell, right=0pt of b5] (b4) {0};
    \node[hlcell, right=0pt of b4] (b3) {1};
    \node[cell, right=0pt of b3] (b2) {0};
    \node[cell, right=0pt of b2] (b1) {0};
    \node[hlcell, right=0pt of b1] (b0) {1};
    \node[idxlbl, above=2pt of b7] {7};
    \node[idxlbl, above=2pt of b6] {6};
    \node[idxlbl, above=2pt of b5] {5};
    \node[idxlbl, above=2pt of b4] {4};
    \node[idxlbl, above=2pt of b3] {3};
    \node[idxlbl, above=2pt of b2] {2};
    \node[idxlbl, above=2pt of b1] {1};
    \node[idxlbl, above=2pt of b0] {0};
    \node[idxlbl, below=2pt of b7] {$128$};
    \node[idxlbl, below=2pt of b6] {$64$};
    \node[idxlbl, below=2pt of b5] {$32$};
    \node[idxlbl, below=2pt of b4] {$16$};
    \node[idxlbl, below=2pt of b3] {$8$};
    \node[idxlbl, below=2pt of b2] {$4$};
    \node[idxlbl, below=2pt of b1] {$2$};
    \node[idxlbl, below=2pt of b0] {$1$};
\end{tikzpicture}
\diagcaption{上方是位編號，下方是權重 $2^n$。 最左 bit $7$ 是 MSB，最
右 bit $0$ 是 LSB。 橘色格子是值為 $1$ 的位：bit $6$、$3$、$0$，加總
$64 + 8 + 1 = 73$。}
\label{fig:ch12-bit-index}
\end{center}

\texttt{1u << n} 會在第 $n$ 位放一個 $1$，其餘為 $0$，正好對準這一格。
印點陣圖、拆封包時，題目若寫從 MSB 開始，就是從最左邊、編號最大的那
一位先處理。

把 $73$ 拆成權重，過程是：最大不超過 $73$ 的權重是 $64$（bit $6$），
剩下 $9$；再取 $8$（bit $3$），剩下 $1$；最後取 $1$（bit $0$）。 其他
權重都是 $0$。

\keypoint{第 $n$ 位的權重是 $2^n$。 MSB 是最高位，LSB 是最低位；一個
位元組裡 MSB 是 bit $7$，LSB 是 bit $0$。}

\subsection{十六進位：四個位元一組}

二進位寫長了不好對。 十六進位的每一位剛好對應 $4$ 個位元（一個
nibble）：$0$ 是 \texttt{0000}，$9$ 是 \texttt{1001}，$A$ 到 $F$ 是
$10$ 到 $15$。 圖 \ref{fig:ch12-bit-index} 那一列可以切成
\texttt{0100} 與 \texttt{1001}，也就是 $4$ 與 $9$，寫成
\texttt{0x49}。

\begin{example}{c}{(print-73.c) 十進位與十六進位}
#include <stdio.h>

int main(void) {
    unsigned char x = 73;
    printf("dec %u\n", x);
    printf("hex 0x%02x\n", x);
}
\end{example}

\iosample{
}{
dec 73 \\
hex 0x49 \\
}

做位元運算時，遮罩習慣寫十六進位：\texttt{0x0F} 是最低 $4$ 位、
\texttt{0xFF} 是一個位元組、\texttt{0x80} 是 8 位裡的 MSB。 腦子裡再
把每個十六進位數字拆回 $4$ 個位即可。

\subsection{複合賦值運算子}

位元運算常常是拿變數自己跟一個遮罩再運算，再寫回自己。 C 為此準備了簡
寫：\texttt{x op= y} 等價於 \texttt{x = x op y}。 後面的片語會直接用這
種寫法。

\begin{code}{c}{複合賦值}
x &= mask;    // same as x = x & mask
x |= bit;     // same as x = x | bit
x ^= bit;     // same as x = x ^ bit
x <<= n;      // same as x = x << n
x >>= n;      // same as x = x >> n
\end{code}

五種位元運算子都有對應的複合賦值。 寫成 \texttt{x |= (1u << 2)} 的時
候，心裡仍然要把它展開成：先算出右邊，再跟左邊做 OR，結果放回
\texttt{x}。

\subsection{有符號數的二補數}

最高位（MSB）在有號整數裡同時是符號位：$0$ 表示非負，$1$ 表示負。 現
代 C 用二補數（two's complement）表示負數。

\begin{example}{c}{(twos-complement.c) 補碼表示}
#include <stdio.h>

int main(void) {
    signed char x = 5;     // binary: 00000101
    signed char neg_x = -x; // binary: 11111011

    printf("5 in hex: 0x%02x\n", (unsigned char)x);
    printf("-5 in hex: 0x%02x\n", (unsigned char)neg_x);
}
\end{example}

\iosample{
}{
5 in hex: 0x05 \\
-5 in hex: 0xfb \\
}

取負的規則可以拆成兩步：先把每一位翻反（\texttt{\~{}x}），再加 $1$。
\texttt{5} 是 \texttt{00000101}，翻反得 \texttt{11111010}，加 $1$ 得
\texttt{11111011}，也就是範例印出的 \texttt{0xfb}。 後面講位元互補時，
會再用到這個翻反的動作。

二補數讓加法本身就能處理正負：$5 + (-5)$ 在 8 位裡是
\texttt{00000101} $+$ \texttt{11111011}，加上進位後低 8 位是
\texttt{00000000}。 $-1$ 的每一位都是 $1$（\texttt{0xff}）：把
\texttt{00000000} 翻反再加 $1$，進位走出格外，留下 8 個 $1$。

\begin{example}{c}{(minus-one.c) -1 與相加歸零}
#include <stdio.h>

int main(void) {
    signed char n = -1;
    signed char a = 5;
    signed char b = -5;
    printf("-1 in hex: 0x%02x\n", (unsigned char)n);
    printf("5 + (-5) = %d\n", a + b);
}
\end{example}

\iosample{
}{
-1 in hex: 0xff \\
5 + (-5) = 0 \\
}

8 位有號的範圍是 $-128$ 到 $127$：\texttt{10000000} 是 $-128$，沒有對
應的正數可再取負。 無號的 \texttt{\~{}0u} 與有號的 $-1$ 位元模式相同
，混在一起時正負會看錯。

\keypoint{二補數取負就是 \texttt{\~{}x + 1}。 $-1$ 的每一位都是 $1$；
有號整數的 MSB 同時是符號位。}

\stylepoint{做位元運算時，變數請宣告成 \texttt{unsigned}。 有號數的右
移、把 $1$ 左移進符號位，行為要嘛未定義、要嘛依實作而定；
\texttt{unsigned} 讓每一位都只是 $0$ 或 $1$，沒有符號位延伸這條後路。}

\subsection{五個運算子怎麼算、怎麼用}

位元運算是對每一個位各自計算，左右鄰居互不干擾，也沒有進位。 這點跟
二補數加法不同：加法會把進位傳到下一格，AND／OR／XOR 不會。

\begin{center}
\begin{tabular}{|l|l|l|}
\hline
\textbf{運算子} & \textbf{每一位} & \textbf{典型用途} \\
\hline
\texttt{\&} & 兩位都是 $1$ 才得 $1$ & 檢查、遮罩 \\
\texttt{|} & 有一位是 $1$ 就得 $1$ & 設置成 $1$ \\
\texttt{\~{}} & $0$ 變 $1$、$1$ 變 $0$ & 做清除用的擋板 \\
\texttt{\textasciicircum{}} & 不同得 $1$、相同得 $0$ & 切換、交換 \\
\texttt{<<}／\texttt{>>} & 整排左右滑 & 對準某一格、乘除 $2^n$ \\
\hline
\end{tabular}
\end{center}

手算時把兩個運算元按 bit $7$ 到 bit $0$ 上下對齊，逐欄套真值表。 後面
五節各對準表裡的一行：且用來檢查、或用來設置、互補用來清除、互斥或用
來交換、移位用來對準再測試或計數。 遮罩這塊擋板會一路跟著走：AND 用
它留下洞、OR 往洞裡塞 $1$、NOT 把洞翻成擋板再清位、XOR 只翻洞裡那一
格。 章末再用同一組旗標與一組檔案權限，把五個運算子串成一支小程式。

\begin{drills}[label=drills:bit-basics]
\item 把十進位 $13$ 拆成 $2$ 的次方相加，寫出對應的二進位。
\item 一個位元組的 MSB 是第幾位？ 它的權重是多少？
\item \texttt{0x49} 換成二進位是多少？ 換成十進位呢？
\item \texttt{x |= (1u << 2)} 這種複合賦值寫法，展開後等價於哪一行完
整算式？
\end{drills}

\section{位元且：檢查位元}
\label{sec:bit-and}

\texttt{x \& y} 對每一位做邏輯 AND。 最常用的是檢查特定位是否為 1。

\keypoint{位元且的真值表：0 AND 0 $= 0$，0 AND 1 $= 0$，1 AND 0 $= 0$，
1 AND 1 $= 1$。 只有兩位都是 1 才得 1。}

\subsection[逐位元計算：AND]{逐位元計算}

AND 是對每一個位各自運算，左右鄰居互不干擾。 把兩個運算元上下對齊，逐
欄看過一次，就能用手算出結果。 橘色格子是這一列為 $1$ 的位；結果列只
在兩排都是橘色的欄留下 $1$。

\begin{center}
\begin{tikzpicture}
    \foreach \col/\bn in {0/7,1/6,2/5,3/4,4/3,5/2,6/1,7/0} {
        \node[idxlbl] at (0.75*\col, 1.35) {\bn};
    }
    \node[idxlbl, anchor=east] at (-0.45, 0.75) {147};
    \foreach \col/\v in {0/1,1/0,2/0,3/1,4/0,5/0,6/1,7/1} {
        \ifnum\v=1
            \node[hlcell, minimum width=0.75cm, minimum height=0.75cm]
                at (0.75*\col, 0.75) {\v};
        \else
            \node[cell, minimum width=0.75cm, minimum height=0.75cm]
                at (0.75*\col, 0.75) {\v};
        \fi
    }
    \node[idxlbl, anchor=east] at (-0.45, 0) {\& 61};
    \foreach \col/\v in {0/0,1/0,2/1,3/1,4/1,5/1,6/0,7/1} {
        \ifnum\v=1
            \node[hlcell, minimum width=0.75cm, minimum height=0.75cm]
                at (0.75*\col, 0) {\v};
        \else
            \node[cell, minimum width=0.75cm, minimum height=0.75cm]
                at (0.75*\col, 0) {\v};
        \fi
    }
    \node[idxlbl, anchor=east] at (-0.45, -0.75) {17};
    \foreach \col/\v in {0/0,1/0,2/0,3/1,4/0,5/0,6/0,7/1} {
        \ifnum\v=1
            \node[hlcell, minimum width=0.75cm, minimum height=0.75cm]
                at (0.75*\col, -0.75) {\v};
        \else
            \node[cell, minimum width=0.75cm, minimum height=0.75cm]
                at (0.75*\col, -0.75) {\v};
        \fi
    }
\end{tikzpicture}
\diagcaption{$147 \& 61$。 上下對齊後逐欄做 AND，只有 bit $4$ 與
bit $0$ 兩排都是 $1$，結果是 $16 + 1 = 17$。}
\label{fig:ch12-and-align}
\end{center}

\begin{code}{c}{逐位元手算 AND}
// bit:  7 6 5 4 3 2 1 0
//       1 0 0 1 0 0 1 1   (147)
//     & 0 0 1 1 1 1 0 1   (61)
//       ---------------
//       0 0 0 1 0 0 0 1   (17)  // bits 4 and 0 survive
\end{code}

十進位驗算：$147 \& 61 = 17$。 後面每次看到 \texttt{\&}，都可以在腦子
裡排成這樣的兩行再對一次。 遮罩就是下面那一排：為 $1$ 的欄是擋板上的
洞，為 $0$ 的欄把上面那一排蓋掉。

\subsection{檢查第 n 位}

要檢查第 $n$ 位是否為 $1$，先把 \texttt{1u} 左移 $n$ 格做出遮罩，再跟
原數做 AND。

\begin{code}{c}{檢查第 n 位}
unsigned int x = 12;  // binary: 1100
int bit2 = (x >> 2) & 1;  // shift right 2, check lowest bit
printf("%d\n", bit2);  // output: 1

// or: AND with a 1u mask
int bit3 = (x & (1u << 3)) != 0;  // check bit 3
printf("%d\n", bit3);  // output: 1
\end{code}

最低位（bit $0$）為 $1$ 就是奇數。 \texttt{x \& 1} 只開這一格的洞，結
果不是 $0$ 就是 $1$。

\begin{code}{c}{用最低位檢查奇偶}
int x = 7;
if (x & 1) {
    printf("odd\n");   // 7 AND 1 = 1
} else {
    printf("even\n");
}
\end{code}

\subsection{應用：檢查多個位}

遮罩（mask）是預先設置好的位元組合，用來篩選特定位。

\begin{code}{c}{多位元遮罩}
unsigned int status = 0b1101;  // 13
unsigned int mask = 0b1010;    // 10

int result = status & mask;  // AND: only bit 3 and bit 1 pass
printf("0x%x\n", result);  // output: 0x8
\end{code}

\subsection{遮罩的意義}

遮罩可以想成一塊擋板，上面打了幾個洞。 洞的位置（位元是 $1$）讓原來的
位看得見、原樣通過；擋板擋住的位置（位元是 $0$）強迫結果變成 $0$。 所
以 \texttt{x \& mask} 只留下 mask 為 $1$ 的那些位，其餘一律清掉。

同一個想法也可以用來截斷。 \texttt{ch \&= 0xFF;} 只開最低 8 個洞，高
於一個位元組的位全部被擋住，結果保證落在 $0$ 到 $255$。 從較寬的整數
裡取出最低一個 byte 時，這行很常見。

\subsection[陷阱：flags == MASK 比較法錯了]{陷阱：不要直接比較 \texttt{flags == MASK}}

要確認 MASK 裡的每一位都亮著，條件寫成
\texttt{(flags \& MASK) == MASK}。 這個寫法只看 MASK 覆蓋到的那些位；
\texttt{flags} 上其他位是 $0$ 或 $1$ 都不影響結果。

\begin{code}{c}{用遮罩檢查旗標}
unsigned int flags = 0x7;   // bits 0, 1, 2
unsigned int MASK  = 0x5;   // bits 0 and 2

if (flags == MASK) {              // 7 == 5, false
    printf("exact match\n");
}
if ((flags & MASK) == MASK) {     // (7 & 5) == 5, true
    printf("mask bits on\n");
}
\end{code}

\texttt{flags} 是 $7$，\texttt{MASK} 是 $5$。 直接比較兩個整數，會要求
每一位都相同，多開的 bit $1$ 讓比較失敗。 先做 AND 再比，只問 bit $0$
與 bit $2$ 有沒有亮，答案是有。

括號一定要寫。 \texttt{\&} 的優先序低於 \texttt{==}，若寫成
\texttt{flags \& MASK == MASK}，C 會先算 \texttt{MASK == MASK}（結果是
$1$），再跟 \texttt{flags} 做 AND，整個條件的意義就變了。

\subsection{具名旗標：檢查}

下一節會用 \texttt{FLAG\_BLUE}、\texttt{FLAG\_GREEN}、
\texttt{FLAG\_RED} 來設置旗標。 檢查走的就是這一節的 AND：洞開在要問
的那幾位，其餘擋住。

\begin{example}{c}{(flag-check.c) 檢查具名旗標}
#include <stdio.h>

#define FLAG_BLUE  1u
#define FLAG_GREEN 2u
#define FLAG_RED   4u

int main(void) {
    unsigned int flags = FLAG_BLUE | FLAG_RED;

    if ((flags & FLAG_BLUE) != 0) {
        printf("blue on\n");
    }
    if ((flags & FLAG_GREEN) != 0) {
        printf("green on\n");
    } else {
        printf("green off\n");
    }
    if ((flags & (FLAG_BLUE | FLAG_RED)) == (FLAG_BLUE | FLAG_RED)) {
        printf("blue and red on\n");
    }
}
\end{example}

\iosample{
}{
blue on \\
green off \\
blue and red on \\
}

\texttt{FLAG\_BLUE | FLAG\_RED} 先合成要問的那兩格。 跟
\texttt{flags == MASK} 那個陷阱同一句：先 AND 再比 MASK，多開的位不影
響。

\subsection{練習：手算 AND}

把 $182$（\texttt{10110110}）跟 $205$（\texttt{11001101}）上下對齊，
逐欄做 AND。 寫完再對解答：$10000100$，也就是 $132$。 只有 bit $7$
與 bit $2$ 兩排都是 $1$。

AND 留下擋板上的洞。 下一節的 OR 則是往指定的格子裡塞 $1$，已經是
$1$ 的位維持 $1$。

\begin{drills}[label=drills:bit-and]
\item 手算 $147 \& 61$（bit 7 到 bit 0 上下對齊逐欄算），結果是多少？
\item 要測試變數 \texttt{x} 的第 3 位是否為 1，該怎麼寫？
\item 為什麼 \texttt{(flags \& MASK) == MASK} 是對的，而
\texttt{flags == MASK} 通常是錯的？
\item \texttt{flags \& MASK == MASK}（沒加括號）會被剖析成什麼？ 為什
麼跟原意不同？
\end{drills}

\section{位元或：設置位元}
\label{sec:bit-or}

\texttt{x | y} 對每一位做邏輯 OR。 最常用的是把特定位設為 1。 手算時
一樣上下對齊：任一欄有 $1$，結果就是 $1$。 跟 AND 那張圖比，這裡只要
有一排是橘色，結果列就是橘色。

\begin{center}
\begin{tikzpicture}
    \foreach \col/\bn in {0/7,1/6,2/5,3/4,4/3,5/2,6/1,7/0} {
        \node[idxlbl] at (0.75*\col, 1.35) {\bn};
    }
    \node[idxlbl, anchor=east] at (-0.45, 0.75) {128};
    \foreach \col/\v in {0/1,1/0,2/0,3/0,4/0,5/0,6/0,7/0} {
        \ifnum\v=1
            \node[hlcell, minimum width=0.75cm, minimum height=0.75cm]
                at (0.75*\col, 0.75) {\v};
        \else
            \node[cell, minimum width=0.75cm, minimum height=0.75cm]
                at (0.75*\col, 0.75) {\v};
        \fi
    }
    \node[idxlbl, anchor=east] at (-0.45, 0) {| 4};
    \foreach \col/\v in {0/0,1/0,2/0,3/0,4/0,5/1,6/0,7/0} {
        \ifnum\v=1
            \node[hlcell, minimum width=0.75cm, minimum height=0.75cm]
                at (0.75*\col, 0) {\v};
        \else
            \node[cell, minimum width=0.75cm, minimum height=0.75cm]
                at (0.75*\col, 0) {\v};
        \fi
    }
    \node[idxlbl, anchor=east] at (-0.45, -0.75) {132};
    \foreach \col/\v in {0/1,1/0,2/0,3/0,4/0,5/1,6/0,7/0} {
        \ifnum\v=1
            \node[hlcell, minimum width=0.75cm, minimum height=0.75cm]
                at (0.75*\col, -0.75) {\v};
        \else
            \node[cell, minimum width=0.75cm, minimum height=0.75cm]
                at (0.75*\col, -0.75) {\v};
        \fi
    }
\end{tikzpicture}
\diagcaption{$128 | 4$。 bit $7$ 與 bit $2$ 各有一排是 $1$，結果兩格
都亮，得到 $132$。}
\label{fig:ch12-or-align}
\end{center}

\begin{code}{c}{逐位元手算 OR}
// bit:  7 6 5 4 3 2 1 0
//       1 0 0 0 0 0 0 0   (128)
//     | 0 0 0 0 0 1 0 0   (4)
//       ---------------
//       1 0 0 0 0 1 0 0   (132)
\end{code}

\subsection{設置第 n 位為 1}

遮罩寫成 \texttt{1u << n}。 字面值 \texttt{1} 的型別是有號
\texttt{int}；在 $32$ 位系統上，\texttt{1 << 31} 會把 $1$ 推進符號位，
行為沒有定義。 \texttt{1u} 是 \texttt{unsigned}，\texttt{1u << 31} 得
到 \texttt{0x80000000u}，每一位的意義都穩定。

\begin{code}{c}{設置第 n 位}
unsigned int x = 8;   // binary: 1000
x |= (1u << 2);      // set bit 2: OR with 0100
printf("0x%x\n", x);  // output: 0xc (12 in decimal)
\end{code}

\subsection{應用：開啟功能}

在嵌入式系統中，硬體暫存器用位表示開關狀態。

\begin{code}{c}{硬體暫存器}
unsigned int flags = 0;
flags |= (1u << 0);  // enable interrupt
flags |= (1u << 2);  // enable timer
printf("0x%x\n", flags);  // output: 0x5
\end{code}

\subsection{具名旗標}

把每個位的意義寫成巨集，讀程式的人就不必心算第幾位代表什麼。 習慣上每
個旗標只占一個位，值是 $1$、$2$、$4$、$8$ 這種 $2$ 的次方，互相不會重疊。

\begin{code}{c}{具名旗標}
#define FLAG_BLUE  1    // bit 0
#define FLAG_GREEN 2    // bit 1
#define FLAG_RED   4    // bit 2

unsigned int flags = 0;
flags |= FLAG_BLUE | FLAG_RED;   // turn on blue and red at once
// flags is now 5 (101)
\end{code}

\texttt{FLAG\_BLUE | FLAG\_RED} 先合成遮罩 $5$（bit $0$ 與 bit $2$），
再 OR 進 \texttt{flags}。 已經亮著的位維持亮著，原本是 $0$ 的位被設成
$1$。 上一節用同一組名字做檢查，下一節用它們做清除。

\subsection{陷阱：\texttt{=} 會蓋掉其他旗標}

要把藍色打開，正確是 \texttt{flags |= FLAG\_BLUE}。 若寫成
\texttt{flags = FLAG\_BLUE}，整格被換成 $1$，原本亮著的綠、紅全部熄
掉。 OR 只把遮罩為 $1$ 的位點亮，其餘位原樣留下；指派則是整份覆寫。

\begin{code}{c}{設置用 OR，不要整份覆寫}
unsigned int flags = FLAG_GREEN | FLAG_RED;  // 6
flags = FLAG_BLUE;     // now 1; green and red are gone
flags = FLAG_GREEN | FLAG_RED;
flags |= FLAG_BLUE;    // now 7; green and red stay
\end{code}

\subsection{練習：手算 OR}

同一對數字：$182$（\texttt{10110110}）OR $205$（\texttt{11001101}）。
解答是 $11111111$，也就是 $255$。 每一欄至少有一個 $1$。

OR 把指定的位設成 $1$。 要關回去，得先做一塊只擋那幾格的板，再 AND
回去，就是下一節的互補。

\keypoint{位元或的真值表：任何位 OR 0 不變，任何位 OR 1 變 1。}

\begin{drills}[label=drills:bit-or]
\item 手算 $182 | 205$（二進位對齊逐欄算），結果是多少？
\item 為什麼設置某個旗標要用 \texttt{|=}，直接寫 \texttt{flags =
FLAG\_BLUE} 會有什麼後果？
\item 為什麼遮罩要寫成 \texttt{1u << n} 而不是 \texttt{1 << n}？
\end{drills}

\section{位元互補：清除位元}
\label{sec:bit-complement}

\texttt{\~{} x} 翻反 \texttt{x} 的每一位：0 變 1，1 變 0。 最常用的是把
特定位清為 0。

\subsection{清除第 n 位為 0}

\begin{code}{c}{清除第 n 位}
unsigned int x = 12;   // binary: 1100
x &= ~(1u << 2);      // clear bit 2: AND with ~0100 = 1011
printf("0x%x\n", x);  // output: 0x8 (8 in decimal)
\end{code}

注意：\texttt{\~{}(1u << 2)} 先左移再翻反。 括號很重要。

\subsection{應用：關閉功能}

\begin{code}{c}{關閉功能}
unsigned int flags = 0x7;  // binary: 0111
flags &= ~(1u << 1);      // disable bit 1
printf("0x%x\n", flags);  // output: 0x5 (0101)
\end{code}

\subsection{清除具名旗標}

延續上一節的 \texttt{FLAG\_*}。 要同時關掉好幾個旗標，先用 OR 把它們
合成一塊，再取互補，最後 AND 回去。

\begin{code}{c}{清除具名旗標}
#define FLAG_BLUE  1
#define FLAG_GREEN 2
#define FLAG_RED   4

unsigned int flags = FLAG_BLUE | FLAG_GREEN | FLAG_RED;  // 7
flags &= ~(FLAG_GREEN | FLAG_RED);   // turn off green and red
// flags is now FLAG_BLUE (1)
\end{code}

先算 \texttt{FLAG\_GREEN | FLAG\_RED}，得到 $6$（二進位 \texttt{110}）。
\texttt{\~{}6} 把這兩位翻成 $0$、其餘翻成 $1$。 再跟原本的 $7$
（\texttt{111}）做 AND，只留下 bit $0$，也就是 \texttt{FLAG\_BLUE}。
補碼那一節說過取負是 \texttt{\~{}x + 1}；這裡只用到前半段的翻反，用來
做擋板。

\subsection[可攜寫法：全部是 1 用 ~0]{可攜寫法：\texttt{\~{}0} 代表全部是 1}

把遮罩寫成 \texttt{0xFF} 或 \texttt{0xFFFF}，等於把位元寬度綁死在 $8$
位或 $16$ 位。 \texttt{int} 若是 $32$ 位，\texttt{0xFF} 只開最低 8 個
洞，高位全是 $0$。 換成 \texttt{\~{}0}，每一位都會是 $1$，寬度跟著當
前型別走。

unsigned 語境請寫 \texttt{\~{}0u}。 有號的 \texttt{\~{}0} 在補碼裡就是
$-1$，跟整數運算混在一起時容易看錯正負。

\begin{code}{c}{可攜的全 1 遮罩}
unsigned int x = 0x3F;     // 00111111
x &= ~0x1Fu;               // clear the lowest 5 bits
// x is now 0x20 (00100000)
\end{code}

\texttt{0x1F} 是最低 5 位為 $1$。 \texttt{\~{}0x1Fu} 把這 5 位翻成 $0$
、其餘維持 $1$，再 AND 回去就清掉最低 5 位。
\texttt{0x3F \& \~{}0x1F} 得 \texttt{0x20}，跟手算一致。 這裡沒有寫死
\texttt{0xFFFFFFE0} 這種跟 $32$ 位綁在一起的常數。

清除是 AND 加上一塊翻反過的擋板：洞變成牆，牆變成洞。 若只要翻某一格
、其他位不動，擋板的想法仍然適用，運算子改成 XOR，見下一節。

\keypoint{位元互補的真值表：NOT 0=1，NOT 1=0。 翻反所有位。
常與 AND 搭配使用來清除特定位。}

\begin{drills}[label=drills:bit-complement]
\item \texttt{x = 12}（\texttt{1100}），執行 \texttt{x \&= \~{}(1u <<
2)} 之後 \texttt{x} 是多少？
\item 為什麼要寫 \texttt{\~{}0u} 而不是寫死 \texttt{0xFF} 或
\texttt{0xFFFF} 當全 1 遮罩？
\item 要同時清除 \texttt{FLAG\_GREEN} 與 \texttt{FLAG\_RED} 兩個旗標，
該怎麼寫成一行？
\end{drills}

\section{位元互斥或：切換與交換}
\label{sec:bit-xor}

\texttt{x \textasciicircum{} y} 對每一位做邏輯 XOR。 相同為 $0$，不同
為 $1$。 最常用的是切換某一格，以及不用臨時變數交換兩個整數。 這兩件
事都靠同一條性質：同一個值 XOR 兩次，會回到原值。

\subsection[逐位元計算：XOR]{逐位元計算}

XOR 一樣上下對齊、逐欄算。 兩排相同的欄得 $0$，不同的欄得 $1$。

\begin{center}
\begin{tikzpicture}
    \foreach \col/\bn in {0/7,1/6,2/5,3/4,4/3,5/2,6/1,7/0} {
        \node[idxlbl] at (0.75*\col, 1.35) {\bn};
    }
    \node[idxlbl, anchor=east] at (-0.45, 0.75) {5};
    \foreach \col/\v in {0/0,1/0,2/0,3/0,4/0,5/1,6/0,7/1} {
        \ifnum\v=1
            \node[hlcell, minimum width=0.75cm, minimum height=0.75cm]
                at (0.75*\col, 0.75) {\v};
        \else
            \node[cell, minimum width=0.75cm, minimum height=0.75cm]
                at (0.75*\col, 0.75) {\v};
        \fi
    }
    \node[idxlbl, anchor=east] at (-0.45, 0) {\textasciicircum{} 12};
    \foreach \col/\v in {0/0,1/0,2/0,3/0,4/1,5/1,6/0,7/0} {
        \ifnum\v=1
            \node[hlcell, minimum width=0.75cm, minimum height=0.75cm]
                at (0.75*\col, 0) {\v};
        \else
            \node[cell, minimum width=0.75cm, minimum height=0.75cm]
                at (0.75*\col, 0) {\v};
        \fi
    }
    \node[idxlbl, anchor=east] at (-0.45, -0.75) {9};
    \foreach \col/\v in {0/0,1/0,2/0,3/0,4/1,5/0,6/0,7/1} {
        \ifnum\v=1
            \node[hlcell, minimum width=0.75cm, minimum height=0.75cm]
                at (0.75*\col, -0.75) {\v};
        \else
            \node[cell, minimum width=0.75cm, minimum height=0.75cm]
                at (0.75*\col, -0.75) {\v};
        \fi
    }
\end{tikzpicture}
\diagcaption{$5 \oplus 12 = 9$。 bit $3$ 與 bit $0$ 兩排不同，結果這
兩格是 $1$；其餘欄相同，結果是 $0$。}
\label{fig:ch12-xor-align}
\end{center}

\keypoint{位元互斥或的真值表：0 XOR 0=0，0 XOR 1=1，1 XOR 0=1，1 XOR 1=0。
相同變 0，不同變 1。 任何數 XOR 自己都是 $0$；任何數 XOR $0$ 都不變。}

\subsection{同一值 XOR 兩次會抵銷}

對每一位來說，把同一個 $k$ 的位 XOR 進去兩次，等於沒做。 寫成式子：
\texttt{(x \textasciicircum{} k) \textasciicircum{} k} 等於
\texttt{x}。 先記住這條，下面的交換與加密都是它的應用。

\begin{code}{c}{XOR 兩次回到原值}
unsigned int x = 0x48;   // 'H'
unsigned int k = 0x2A;
unsigned int y = x ^ k;
unsigned int back = y ^ k;   // back == x
\end{code}

\subsection{為什麼 XOR 交換行得通}

令一開始 $a$ 存 $a_0$、$b$ 存 $b_0$。 三步之後，$a$ 應變成 $b_0$，
$b$ 應變成 $a_0$。 每一步都用上一小節的抵銷。

\begin{center}
\begin{tabular}{|l|l|l|l|}
\hline
\textbf{步驟} & \textbf{$a$} & \textbf{$b$} & \textbf{這一步做了什麼} \\
\hline
開始 & $a_0$（$5$，\texttt{0101}） & $b_0$（$12$，\texttt{1100}） &
兩格各自的原值 \\
\texttt{a \textasciicircum{}= b} & $a_0 \oplus b_0$（$9$，\texttt{1001}） &
$b_0$ & $a$ 暫時變成兩者的 XOR \\
\texttt{b \textasciicircum{}= a} & $a_0 \oplus b_0$ &
$b_0 \oplus (a_0 \oplus b_0) = a_0$ & $b_0$ 被 XOR 兩次，抵銷，留下
$a_0$ \\
\texttt{a \textasciicircum{}= b} &
$(a_0 \oplus b_0) \oplus a_0 = b_0$ & $a_0$ & $a_0$ 被 XOR 兩次，抵
銷，留下 $b_0$ \\
\hline
\end{tabular}
\end{center}

第三步的 $b$ 已經是 $a_0$，所以最後一次 XOR 等於把 $a$ 裡的 $a_0$ 消
掉，只剩 $b_0$。

\begin{code}{c}{用 XOR 交換}
int a = 5;   // binary: 0101
int b = 12;  // binary: 1100

a = a ^ b;  // a: 0101 XOR 1100 = 1001 (9)
b = b ^ a;  // b: 1100 XOR 1001 = 0101 (5)
a = a ^ b;  // a: 1001 XOR 0101 = 1100 (12)

printf("a=%d, b=%d\n", a, b);  // output: a=12, b=5
\end{code}

這個技巧只對整數有效，不會用到額外記憶體。

\subsection[陷阱：交換的兩變數不能同位址]{陷阱：交換的兩個變數不能是同一個位址}

XOR 交換法要求兩個運算元是不同的變數。 任何數 XOR 自己都是 $0$；若對
同一個變數連續做三次 \texttt{a = a \textasciicircum{} a}，第一步就已經
把資料清成 $0$，後面兩步只是 $0$ XOR $0$。

\begin{code}{c}{對同一變數做 XOR 交換}
int a = 5;
a = a ^ a;   // 5 XOR 5 = 0
a = a ^ a;   // 0 XOR 0 = 0
a = a ^ a;   // still 0
printf("%d\n", a);   // output: 0
\end{code}

呼叫函式時若兩個參數其實指向同一格記憶體，例如 \texttt{swap(\&x, \&x)}
，也會踩進這個陷阱。 日常交換兩個整數，用一個臨時變數比較不容易寫錯。

\subsection{具名旗標：切換}

OR 把位設成 $1$，AND 加互補把位清成 $0$。 XOR 則是亮的變暗、暗的變亮
，做一次切一次，做兩次回到原狀，正好是抵銷性質。

\begin{example}{c}{(flag-toggle.c) 切換具名旗標}
#include <stdio.h>

#define FLAG_BLUE  1u
#define FLAG_GREEN 2u
#define FLAG_RED   4u

int main(void) {
    unsigned int flags = FLAG_BLUE | FLAG_RED;
    printf("start: %u\n", flags);
    flags ^= FLAG_GREEN;
    printf("toggle green on: %u\n", flags);
    flags ^= FLAG_GREEN;
    printf("toggle green off: %u\n", flags);
}
\end{example}

\iosample{
}{
start: 5 \\
toggle green on: 7 \\
toggle green off: 5 \\
}

\texttt{FLAG\_GREEN} 只占 bit $1$。 第一次 XOR 把它點亮，$5$ 變成
$7$；第二次再 XOR 同一格，綠燈熄掉，回到 $5$。 藍與紅沒被碰到。

\subsection{練習：手算 XOR}

同一對數字：$182$（\texttt{10110110}）XOR $205$（\texttt{11001101}）。
解答是 $01111011$，也就是 $123$。 跟前面 AND、OR 的練習比：AND 只留兩
排都是 $1$ 的欄，OR 留至少一排是 $1$ 的欄，XOR 留剛好一排是 $1$ 的欄。

切換與交換都靠 XOR 的可逆性。 要對準某一格、或把幾個位元組拆開再組回
去，還需要把整排左右滑，見下一節。

\subsection{應用：XOR 加密}

交換法用過的同一條性質也可以拿來加密。 明文跟一把固定的 key 做 XOR，
得到密文；密文再跟同一把 key 做 XOR，明文就回來了。

\begin{example}{c}{(xor-cipher.c) XOR 加密}
#include <stdio.h>

int main(void) {
    unsigned char msg0 = 'H';
    unsigned char msg1 = 'I';
    unsigned char key = 0x2A;

    unsigned char enc0 = msg0 ^ key;
    unsigned char enc1 = msg1 ^ key;
    unsigned char dec0 = enc0 ^ key;
    unsigned char dec1 = enc1 ^ key;

    printf("plain: %c%c\n", msg0, msg1);
    printf("cipher: 0x%02x 0x%02x\n", enc0, enc1);
    printf("decoded: %c%c\n", dec0, dec1);
}
\end{example}

\iosample{
}{
plain: HI \\
cipher: 0x62 0x63 \\
decoded: HI \\
}

\texttt{'H'} 是 \texttt{0x48}，\texttt{'I'} 是 \texttt{0x49}，key 是
\texttt{0x2A}。 \texttt{0x48 \textasciicircum{} 0x2A} 得 \texttt{0x62}，
\texttt{0x49 \textasciicircum{} 0x2A} 得 \texttt{0x63}。 密文再 XOR 一
次 \texttt{0x2A}，兩個位元組都回到明文。 這把固定 key 的作法只適合用
來理解 XOR 的可逆性；真正的加密還需要更長、不重複使用的金鑰。

\begin{drills}[label=drills:bit-xor]
\item 手算 $182 \oplus 205$，結果是多少？
\item 為什麼 \texttt{(x \textasciicircum{} k) \textasciicircum{} k}
會等於 \texttt{x}？
\item XOR 交換法為什麼不能用在 \texttt{swap(\&x, \&x)} 這種情況？
\item \texttt{flags \textasciicircum{}= FLAG\_GREEN} 連續執行兩次，
\texttt{flags} 最後會變成什麼？
\end{drills}

\section{移位：數字移位、測試與計數}
\label{sec:bit-shift}

前面四節都假設遮罩已經對準要動的那一格。 對準的動作就是移位：
\texttt{1u << n} 做出第 $n$ 位的擋板，\texttt{x >> n} 把第 $n$ 位滑到
最低位再檢查。 這一節把移位本身講清楚，並用它拆、組多位元組的整數。

\subsection{左移與右移}

\texttt{x << n} 左移 n 格，等同乘以 $2^n$。 \texttt{x >> n} 右移 n 格，
等同除以 $2^n$（整數除法）。

\begin{code}{c}{左移與右移}
unsigned int x = 3;   // binary: 11
unsigned int y = x << 2; // left shift: 1100 = 12
unsigned int z = y >> 1; // right shift: 0110 = 6

printf("%u << 2 = %u\n", x, y);  // 3 << 2 = 12
printf("%u >> 1 = %u\n", y, z);  // 12 >> 1 = 6
\end{code}

\keypoint{十進位裡把小數點右移兩位，等於乘以 $100$。 二進位的左移 $n$
位，同樣是乘以 $2^n$。 這個類比在非負數與 \texttt{unsigned} 上完全成
立。}

\subsection{有號數右移：算術位移與邏輯位移}

右移空出來的高位要補什麼，得看型別。 \texttt{unsigned} 永遠補 $0$，這
叫邏輯位移。 有號負數的右移在 C 裡是 implementation-defined：編譯器可
以補 $0$，也可以複製符號位（算術位移，負數會一直保持負的）。

同一份程式換一台機器、換一個編譯器，有號右移的結果可能跟著變。 可攜的
寫法是：有號數若只是要除以 $2$ 的次方，直接用 \texttt{/}，不要用
\texttt{>>}。

就算編譯器對負數做算術位移，結果也跟 \texttt{/} 不同。
\texttt{-3 / 2} 得到 $-1$（向零取整）；算術右移一位通常得到 $-2$（向
負無窮取整）。 兩條路連取整的方向都不一樣，混用會悄悄算錯。

\subsection{測試位元（用右移）}

最直接的方法：右移目標位到最低位，再用 AND 檢查。

\begin{code}{c}{用右移測試位元}
unsigned int x = 12;  // binary: 1100
int bit2 = (x >> 2) & 1;  // shift right 2, check bit
printf("bit 2: %d\n", bit2);  // output: 1

int bit0 = (x >> 0) & 1;  // check bit 0
printf("bit 0: %d\n", bit0);  // output: 0
\end{code}

\subsection{應用：從整數中取出色彩分量}

一個 \texttt{unsigned long} 可以一次塞進三個 byte。 常見的包裝是把紅、
綠、藍依序放進不同的 8 位區間：最高那個 byte 是 R，中間是 G，最低是 B。
要拆開來用，先右移對準，再跟 \texttt{0xFF} 做 AND，只留下一個位元組。

\begin{example}{c}{(extract-color.c) 取出 RGB 分量}
#include <stdio.h>

int main(void) {
    unsigned long color = 0xC04080UL;
    unsigned int r = (color >> 16) & 0xFF;
    unsigned int g = (color >> 8) & 0xFF;
    unsigned int b = color & 0xFF;

    printf("R = 0x%02x\n", r);
    printf("G = 0x%02x\n", g);
    printf("B = 0x%02x\n", b);
}
\end{example}

\iosample{
}{
R = 0xc0 \\
G = 0x40 \\
B = 0x80 \\
}

\texttt{0xC04080} 的三個 byte 分別是 \texttt{0xC0}、\texttt{0x40}、
\texttt{0x80}。 \texttt{(color >> 16) \& 0xFF} 把 R 滑到最低 8 位再截
斷；右移 $8$ 位取出 G；不移位、直接 AND，取出本來就在最右邊的 B。 這
個位移加遮罩的組合，拆封包、拆機器字時會一再出現。

反過來更常寫：手上有三個 $0$ 到 $255$ 的分量，要組回一個整數。 每個
分量先左移到自己的區間，再用 OR 併起來。 AND 負責截斷，OR 負責拼接，
移位負責對準。

\begin{example}{c}{(pack-color.c) 組回 RGB}
#include <stdio.h>

int main(void) {
    unsigned int r = 0xC0;
    unsigned int g = 0x40;
    unsigned int b = 0x80;
    unsigned long color =
        ((unsigned long)r << 16) | ((unsigned long)g << 8) | b;

    printf("color = 0x%06lx\n", color);

    unsigned int r2 = (color >> 16) & 0xFF;
    unsigned int g2 = (color >> 8) & 0xFF;
    unsigned int b2 = color & 0xFF;
    printf("R = 0x%02x\n", r2);
    printf("G = 0x%02x\n", g2);
    printf("B = 0x%02x\n", b2);
}
\end{example}

\iosample{
}{
color = 0xc04080 \\
R = 0xc0 \\
G = 0x40 \\
B = 0x80 \\
}

左移之前先轉成 \texttt{unsigned long}，避免 $32$ 位 \texttt{int} 把高
位切掉。 組回去的值跟上一例拆開前相同，再拆一次仍得原來的三個分量。

\subsection{計數 1 的個數}

給定整數，數一數二進位表示中有多少個 1。

\begin{example}{c}{(count-bits.c) 計算 1 的個數}
#include <stdio.h>

int count_ones(unsigned int x) {
    int count = 0;
    while (x) {
        count += (x & 1);  // check lowest bit
        x >>= 1;           // right shift
    }
    return count;
}

int main(void) {
    printf("%d\n", count_ones(7));   // 111 = 3 ones
    printf("%d\n", count_ones(12));  // 1100 = 2 ones
}
\end{example}

\iosample{
}{
3 \\
2 \\
}

\subsection{進階：更快的計數技巧（選讀）}

上面的迴圈每一步只看最低位，再右移一格。 $1$ 若出現在很高的位，就得一
直移到那一格為止；最壞情況接近型別的位寬。 有一個常見的加速：
\texttt{n \&= (n - 1);} 會一次清掉最低那個 $1$，其餘位維持原樣。

\begin{code}{c}{一次清掉最低位的 1}
int count_ones_fast(unsigned int n) {
    int count = 0;
    while (n) {
        n &= (n - 1);   // clear the lowest 1-bit
        count++;
    }
    return count;
}
\end{code}

以 $12$（\texttt{1100}）為例。 $12 - 1 = 11$（\texttt{1011}），AND 之
後得 $8$（\texttt{1000}），最低的那個 $1$ 消失了。 下一步
$8 - 1 = 7$（\texttt{0111}），AND 得 $0$，迴圈結束。 總共走 $2$ 次，
剛好等於 $1$ 的個數。 前面那個右移版本對 $12$ 要走 $4$ 步才移完。

一般情況差在迴圈次數跟什麼成正比。 右移版每一步只看最低位，一定走到
最高那個 $1$ 為止，步數接近型別位寬；最壞是 $32$ 位裡只有 MSB 為 $1$
，仍要走滿 $32$ 步。 加速版每一步清掉一個 $1$，步數等於 $1$ 的個數。
\texttt{0x80000001} 只有兩個 $1$，一在 bit $0$、一在 bit $31$：右移版
走 $32$ 步，加速版走 $2$ 步。 高位含很多 $1$、中間又夾著長串 $0$ 時，
兩種寫法的差距就是位寬對上 popcount。

產品程式通常不必手寫這兩種迴圈。 GCC 與 Clang 提供
\texttt{\_\_builtin\_popcount}（64 位用 \texttt{\_\_builtin\_popcountll}
），編譯器會把它編成硬體的 popcnt 指令。 作業、講原理、或編譯器沒有
內建時，用上面的迴圈；其餘時候優先用內建。

\begin{code}{c}{編譯器內建的 popcount}
unsigned int x = 0x80000001u;
int n = __builtin_popcount(x);   // 2; GCC / Clang
\end{code}

\keypoint{移位運算在二進位層面上動作，所以超過型別位寬是未定義行為。
例如 32 位整數右移 32 格結果不可預測。}

\begin{drills}[label=drills:bit-shift]
\item \texttt{unsigned int x = 3;}，\texttt{x << 2} 的結果是多少？
\item 有號負數右移的結果，C 標準保證了什麼、又沒保證什麼？
\item 要從 \texttt{0xC04080} 取出中間那個位元組（G），該怎麼寫？
\item \texttt{n \&= (n - 1);} 這一行對 \texttt{n} 的位元模式做了什麼事？
\end{drills}

\section{綜合應用}
\label{sec:bit-applications}

切換第 $n$ 位就是 XOR 加上移位做出的遮罩，上一節的
\texttt{FLAG\_GREEN} 已經做過一次。 這裡用 Unix 風格的檔案權限，把五
個運算子放進同一支程式：OR 設置、AND 檢查、AND 加互補清除、XOR 切換、
移位負責把 owner／group／other 三組 $3$ 位對準。

每一組權限占 $3$ 位，意義與十進位 $4$、$2$、$1$ 相同：讀、寫、執行。
owner 放在 bit $8$ 到 $6$，group 在 bit $5$ 到 $3$，other 在 bit $2$
到 $0$。 \texttt{0755} 這種八進位字面值，剛好一位一位對準這 $9$ 格。

\begin{example}{c}{(file-mode.c) 檔案權限的位元}
#include <stdio.h>

#define R 4u
#define W 2u
#define X 1u

unsigned int pack(unsigned int owner, unsigned int group,
                  unsigned int other) {
    return (owner << 6) | (group << 3) | other;
}

int main(void) {
    unsigned int mode = pack(R | W | X, R | X, R | X);
    printf("mode = 0%o\n", mode);

    unsigned int owner = (mode >> 6) & 7u;
    printf("owner = %u\n", owner);

    if ((owner & W) != 0) {
        printf("owner can write\n");
    }

    mode &= ~(W << 6);
    printf("clear owner write: 0%o\n", mode);

    mode ^= X;
    printf("toggle other exec: 0%o\n", mode);
}
\end{example}

\iosample{
}{
mode = 0755 \\
owner = 7 \\
owner can write \\
clear owner write: 0555 \\
toggle other exec: 0554 \\
}

\texttt{pack} 裡的左移加 OR，跟組 RGB 是同一招，只是每段 $3$ 位而不是
$8$ 位。 \texttt{(mode >> 6) \& 7u} 取出 owner：右移對準，AND
\texttt{7u}（二進位 \texttt{111}）當擋板，只留 $3$ 位。
\texttt{(owner \& W) != 0} 是 AND 節的檢查。 \texttt{\~{}(W << 6)} 先
把寫權限對準 owner 那一段再翻反，AND 回去就清掉它，$0755$ 變成
$0555$。 最後 \texttt{mode \textasciicircum{}= X} 切換 other 的執行位
，變成 $0554$。

結構章會看到 \texttt{unsigned x : 3} 這種由編譯器切位的寫法。 本章先
練手動遮罩與移位；欄位較短、又對齊機器字時，手寫位元運算比較能掌握每
一格在哪。 欄位很多、只想依名字讀寫時，再讓編譯器切。

\subsection{優先級陷阱}

位元運算的優先級\textbf{低於}關係運算子。 括號很重要。

\begin{code}{c}{運算子優先序}
// wrong: a | b == c parsed as a | (b == c)
// correct
if ((a | b) == c) {
    printf("OK\n");
}

// wrong: x & 0x0F = 0 uses = not ==
// correct
if ((x & 0x0F) == 0) {
    printf("OK\n");
}
\end{code}

\keypoint{位元運算子順序：\texttt{\&} 高於 \texttt{\textasciicircum{}}
高於 \texttt{|}。 複雜表達式一定要括號。}

\subsection{慣用語總表}

這一章教過的核心寫法可以收成下面這張小抄。 各行都假設被操作的變數是
\texttt{unsigned}，遮罩寫成 \texttt{1u << n}。 從設置、清除到切換都守
這兩條，算術位移與符號位延伸就進不來。

\begin{code}{c}{位元慣用語小抄}
x |=  (1u << n);              // set bit n
x &= ~(1u << n);              // clear bit n
x ^=  (1u << n);              // toggle bit n
on  = (x & (1u << n)) != 0;   // test bit n
bit = (x >> n) & 1;           // read bit n
a ^= b; b ^= a; a ^= b;       // XOR swap (a and b distinct)
\end{code}

設置、清除、切換分別對應 OR、AND 加互補、XOR；測試與讀取對應 AND，差
在前者只問亮不亮、後者把那一位抽出來變成 $0$ 或 $1$。 交換那一行記得
避開同一個位址。 檔案權限那支程式把這張小抄全用過一遍：
\texttt{pack} 是移位加 OR，取出 owner 是移位加 AND，關寫權限是 AND 加
互補，切換執行位是 XOR。

\begin{drills}[label=drills:bit-applications]
\item Unix 檔案權限裡，owner 的讀寫執行三位放在哪個 bit 範圍？
\item \texttt{a | b == c}（沒加括號）會被剖析成什麼？ 該怎麼加括號修
正？
\item \texttt{mode \&= \~{}(W << 6);} 這一行在做什麼？
\end{drills}

\addcontentsline{toc}{section}{練習解答}
\section*{練習解答}

各節練習的答案收在這裡。 手算題直接給數值結果，觀念題給能一句話核對
的簡答。

\paragraph{第~\ref{sec:bit-basics}~節}
\begin{enumerate}
\item $13 = 8 + 4 + 1 = 2^3 + 2^2 + 2^0$，二進位是 \texttt{1101}。
\item MSB 是第 7 位（一個位元組的最高位），權重是 $2^7 = 128$。
\item \texttt{0x49} 是 \texttt{0100 1001}，十進位是 $73$。
\item 等價於 \texttt{x = x | (1u << 2)}。
\end{enumerate}

\paragraph{第~\ref{sec:bit-and}~節}
\begin{enumerate}
\item $147 \& 61 = 17$（\texttt{10010011 \& 00111101 = 00010001}）。
\item \texttt{(x \& (1u << 3)) != 0}，或等價地 \texttt{(x >> 3) \& 1}。
\item 因為 \texttt{flags} 上其他沒被 \texttt{MASK} 覆蓋到的位是 0 或
1 都不該影響判斷；先 AND 再比，只問 \texttt{MASK} 指定的那幾位。直接
比較兩個整數則要求每一位都相同。
\item 因為 \texttt{\&} 優先序低於 \texttt{==}，會先算
\texttt{MASK == MASK}（結果 1），再跟 \texttt{flags} 做 AND，跟原意完
全不同。
\end{enumerate}

\paragraph{第~\ref{sec:bit-or}~節}
\begin{enumerate}
\item $182 | 205 = 255$（\texttt{10110110 | 11001101 = 11111111}）。
\item \texttt{|=} 只把指定的位設成 1，其餘位不動；直接賦值
\texttt{flags = FLAG\_BLUE} 會把整個 \texttt{flags} 換成新值，原本亮
著的其他旗標全部被蓋掉熄滅。
\item 字面值 \texttt{1} 是有號 \texttt{int}，\texttt{1 << 31} 會把 1
推進符號位、行為未定義；\texttt{1u} 是 \texttt{unsigned}，位移結果穩
定明確。
\end{enumerate}

\paragraph{第~\ref{sec:bit-complement}~節}
\begin{enumerate}
\item $x = 8$（\texttt{1100 \& \~{}0100 = 1100 \& 1011 = 1000}）。
\item \texttt{\~{}0u} 的位寬會自動跟著變數的型別走；寫死
\texttt{0xFF}／\texttt{0xFFFF} 等於把遮罩寬度綁死在 8 位或 16 位，型
別更寬時高位會留下不該有的 0。
\item \texttt{flags \&= \~{}(FLAG\_GREEN | FLAG\_RED);}
\end{enumerate}

\paragraph{第~\ref{sec:bit-xor}~節}
\begin{enumerate}
\item $182 \oplus 205 = 123$（\texttt{10110110 \textasciicircum{}
11001101 = 01111011}）。
\item 因為同一位 XOR 兩次會抵銷：先翻一次變成不同，再翻一次翻回原
值，等於沒做。
\item 因為任何數 XOR 自己都是 0；若兩個運算元其實是同一格記憶體，第
一步就把資料清成 0，後面兩步只是 0 XOR 0，資料整個消失。
\item 第一次切換後 \texttt{FLAG\_GREEN} 的位被點亮，第二次切換同一位
再翻一次，回到原來的值——等於沒變。
\end{enumerate}

\paragraph{第~\ref{sec:bit-shift}~節}
\begin{enumerate}
\item $3 << 2 = 12$（\texttt{0011} 變 \texttt{1100}）。
\item C 標準只保證 \texttt{unsigned} 右移補 0（邏輯位移）；有號負數右
移是 implementation-defined，可能補 0 也可能複製符號位，不保證跨編譯
器一致。
\item \texttt{(color >> 8) \& 0xFF}。
\item 一次清掉 \texttt{n} 裡最低的那個 1 位元，其餘位維持不變。
\end{enumerate}

\paragraph{第~\ref{sec:bit-applications}~節}
\begin{enumerate}
\item bit 8 到 bit 6。
\item 會被剖析成 \texttt{a | (b == c)}，因為 \texttt{|} 優先序低於
\texttt{==}；要改成 \texttt{(a | b) == c}。
\item 把寫權限（\texttt{W}）對準 owner 那一段（左移 6 位）再翻反當擋
板，AND 回去清掉 owner 的寫入位，其餘位不動。
\end{enumerate}

% 變更摘要
% - 第四輪：解讀複雜宣告整節（含原陷阱 6／現陷阱 5）併回本章
% - 格式指定符只留 signed/unsigned 節一張總表
% - const：三種指標之後立刻接跨檔兩種寫法；volatile 等改選讀
% - static：用 print-stars.c 開場，再定義 scope / linkage / duration
% - 陷阱：1 格式、2 有號無號、3 靜默溢位、4 未初始化、5 巢狀 const
% - UB 統一：未定義行為：…這裡印出的數字只是這次編譯的結果…
% TODO(future-chapter): malloc/free/calloc/VLA
% TODO(future-chapter): struct/union/enum 語法
% TODO(future-chapter): 標頭檔保護、makefile、static/dynamic linking、ABI
% TODO(future-chapter): inline 函式

\chapter{宣告修飾}
\label{ch:declarationmodifiers}

整數型別不只有 \texttt{int}。 C 提供了 \texttt{short}、\texttt{long}、
\texttt{long long}、\texttt{signed}、\texttt{unsigned} 等修飾符，讓我們
依範圍與正負選擇型別。 同一套宣告語法也決定變數能不能改
（\texttt{const}）、能活多久、誰看得到（\texttt{static}），以及一串
\texttt{*}、\texttt{[]}、\texttt{()} 該從哪裡讀起。 這一章先把整數家
族、限定符與儲存類別講清楚，再用同一套規則讀複雜宣告；讀完可以去結
構。

\section{short、long 及 long long}
\label{sec:decl-int-sizes}

\subsection{位元組數與範圍}

\begin{code}{c}{整數型別的大小}
#include <stdio.h>
#include <limits.h>

int main(void) {
    printf("short: %zu bytes, range [%d, %d]\n",
           sizeof(short), SHRT_MIN, SHRT_MAX);
    printf("int: %zu bytes, range [%d, %d]\n",
           sizeof(int), INT_MIN, INT_MAX);
    printf("long: %zu bytes, range [%ld, %ld]\n",
           sizeof(long), LONG_MIN, LONG_MAX);
    printf("long long: %zu bytes, range [%lld, %lld]\n",
           sizeof(long long), LLONG_MIN, LLONG_MAX);
}
\end{code}

典型的大小：
\begin{itemize}
\item \texttt{short}：2 位元組（$-32768$ 到 $32767$）
\item \texttt{int}：4 位元組（約 $-2.1 \times 10^9$ 到 $2.1 \times 10^9$）
\item \texttt{long}：4 或 8 位元組（取決於架構）
\item \texttt{long long}：8 位元組（約 $-9.2 \times 10^{18}$ 到 $9.2 \times 10^{18}$）
\end{itemize}

對應的 \texttt{printf} 格式指定符（含無號與八進位）見下一節總表。

\subsection{溢位}

當變數超過其型別的最大值時發生溢位。 帶符號整數加過上限，在許多機器上
會繞回最小負數，但標準把這件事標成未定義行為：編譯器可以環繞、可以刪
掉檢查、也可以做別的事。 無符號整數溢位才有保證，結果是模 $2^n$ 的環
繞。 完整的檢查寫法見後文 \texttt{stdint.h} 那一節。

上面是概念描述，實際跑起來長什麼樣子，看下面這個對照：

\begin{example}{c}{(int-overflow.c) 帶符號與無符號溢位的差別}
#include <stdio.h>
#include <limits.h>

int main(void) {
    int i = INT_MAX;
    printf("INT_MAX = %d\n", i);
    i = i + 1;
    printf("INT_MAX + 1 = %d\n", i);

    unsigned int u = UINT_MAX;
    printf("UINT_MAX = %u\n", u);
    u = u + 1;
    printf("UINT_MAX + 1 = %u\n", u);
}
\end{example}

\iosample{
}{
INT\_MAX = 2147483647 \\
INT\_MAX + 1 = -2147483648 \\
UINT\_MAX = 4294967295 \\
UINT\_MAX + 1 = 0 \\
}

\keypoint{未定義行為：帶符號整數溢位。 這裡印出的數字只是這次編譯的結
果，換機器或最佳化層級不必相同。}

前兩行是帶符號溢位：依標準，\texttt{i + 1} 的結果不受保護。 這裡印出的
\texttt{-2147483648} 只是這台機器、這個編譯器當下剛好呈現的一種可能。
後兩行是無符號溢位：\texttt{u + 1} 保證環繞成 $0$，這是標準明文規定的
行為，換編譯器也一樣。

\begin{drills}[label=drills:decl-int-sizes]
\item \texttt{short} 通常佔幾個位元組？範圍大約是多少？
\item \texttt{INT\_MAX + 1} 在這台機器上實際印出什麼？這個結果是標準
保證的行為嗎？
\item \texttt{UINT\_MAX + 1} 標準怎麼規定？跟帶符號溢位的差別在哪？
\end{drills}

\section{signed 及 unsigned}
\label{sec:decl-signed-unsigned}

\texttt{signed} 和 \texttt{unsigned} 修飾符決定整數是否表示負數。

\subsection{帶符號整數}

\begin{code}{c}{帶符號整數（預設）}
signed int si = -42;
int i = -42;             // default is signed
\end{code}

範圍：$-2^{n-1}$ 到 $2^{n-1} - 1$（其中 $n$ 是位元數）。

\subsection{無符號整數}

\begin{code}{c}{無符號整數}
unsigned int ui = 42;
unsigned short us = 1000;
unsigned long ul = 100000;
unsigned long long ull = 10000000000ULL;
\end{code}

範圍：$0$ 到 $2^n - 1$。

\subsection{格式指定符總表}

把短／長／有號／無號放進同一張表，寫 \texttt{printf} 時可以直接查：

\begin{center}
\begin{tabular}{|l|l|l|l|}
\hline
\textbf{型別} & \textbf{十進位} & \textbf{十六進位} & \textbf{八進位} \\
\hline
\texttt{short} & \texttt{\%hd} & \texttt{\%hx} & \texttt{\%ho} \\
\texttt{int} & \texttt{\%d} & \texttt{\%x} & \texttt{\%o} \\
\texttt{long} & \texttt{\%ld} & \texttt{\%lx} & \texttt{\%lo} \\
\texttt{long long} & \texttt{\%lld} & \texttt{\%llx} & \texttt{\%llo} \\
\texttt{unsigned short} & \texttt{\%hu} & \texttt{\%hx} & \texttt{\%ho} \\
\texttt{unsigned int} & \texttt{\%u} & \texttt{\%x} & \texttt{\%o} \\
\texttt{unsigned long} & \texttt{\%lu} & \texttt{\%lx} & \texttt{\%lo} \\
\texttt{unsigned long long} & \texttt{\%llu} & \texttt{\%llx} & \texttt{\%llo} \\
\hline
\end{tabular}
\end{center}

\texttt{\%x} 與 \texttt{\%o} 不分有號無號，一律把底下的位元當無號數排
版；印一個負數的十六進位，看到的是它的補數表示法。

\keypoint{\texttt{unsigned} 的 \texttt{printf} 格式是 \texttt{\%u}，
寫成 \texttt{\%d} 會依有號解讀同一串位元，印出來的數字就錯了。 格式與
引數型別不符本身是未定義行為，細節見後文陷阱 1。}

\begin{drills}[label=drills:decl-signed-unsigned]
\item \texttt{unsigned int} 的範圍是多少？ 跟 \texttt{signed int} 的範
圍差在哪裡？
\item 用 \texttt{\%d} 印一個 \texttt{unsigned int} 變數，可能出什麼問
題？ 該用哪個格式指定符？
\item \texttt{\%x} 和 \texttt{\%o} 會不會區分有號無號？ 印一個負數的十
六進位，看到的其實是什麼？
\end{drills}

% 合併原 C99 與 stdint 兩節；完整溢位檢查只留 overflow-check.c
\section{C99 與 stdint.h}
\label{sec:decl-stdint}

C99 用 \texttt{stdint.h} 給出固定寬度的整數型別，跨平臺時大小與範圍相
同。

\subsection{固定寬度型別}

\begin{code}{c}{stdint.h 的固定寬度型別}
#include <stdint.h>

int8_t i8 = 127;
int16_t i16 = 32767;
int32_t i32 = 2147483647;
int64_t i64 = 9223372036854775807LL;

uint8_t u8 = 255;
uint16_t u16 = 65535;
uint32_t u32 = 4294967295U;
uint64_t u64 = 18446744073709551615ULL;
\end{code}

同一個標頭還提供至少 $N$ 位的 \texttt{int\_leastN\_t}、實作上較快的
\texttt{int\_fastN\_t}，以及裝得下指標的 \texttt{intptr\_t}／
\texttt{uintptr\_t}。 對應的上下界是 \texttt{INT32\_MAX}、
\texttt{UINT32\_MAX} 這類巨集；\texttt{printf} 可用
\texttt{PRId32}、\texttt{PRIu32} 對準格式。

\subsection{整數提升}

較窄的整數一旦走進運算式，會先被提升成較寬的型別，這叫 integral
promotion。 引數傳進沒有原型、或傳進 \texttt{printf} 這種可變引數函式
時，同樣會提升：所以 \texttt{printf("\%d", short\_var)} 能動，實際送進
去的是 \texttt{int}。

\begin{center}
\begin{tabular}{|l|l|}
\hline
\textbf{出現在運算式裡的型別} & \textbf{提升後} \\
\hline
\texttt{char}、\texttt{signed char}、\texttt{unsigned char} &
\texttt{int} \\
\texttt{short} & \texttt{int} \\
\texttt{unsigned short} & \texttt{int}（裝得下時） \\
\texttt{float}（舊式／無原型引數） & \texttt{double} \\
\hline
\end{tabular}
\end{center}

\texttt{unsigned short} 若裝不進 \texttt{int}，才升成
\texttt{unsigned int}。 \texttt{float} 那一行只適用舊式引數傳遞；兩個
\texttt{float} 相加，在有原型的現代 C 裡仍是 \texttt{float}。 提升之後
若兩邊型別還不同，再走 usual arithmetic conversions：有號與無號
\texttt{int} 一起比較時，有號那邊會轉成無號。 這就是後文陷阱 2 的根
源。

\keypoint{較窄的整數先提升，再跟另一個運算元對齊。
\texttt{int} 碰上 \texttt{unsigned int} 時，是有號轉成無號。}

提升只發生在運算的當下，存回窄型別變數時又是另一回事：

\begin{example}{c}{(promotion-surprise.c) 提升與窄化是兩個不同的步驟}
#include <stdio.h>

int main(void) {
    unsigned char a = 200;
    unsigned char b = 100;
    unsigned char c = a + b;

    printf("a + b 存回 unsigned char = %u\n", c);
    printf("a + b 這個運算式的大小 = %zu bytes\n", sizeof(a + b));
}
\end{example}

\iosample{
}{
a + b 存回 unsigned char = 44 \\
a + b 這個運算式的大小 = 4 bytes \\
}

\texttt{a} 和 \texttt{b} 相加前都先提升成 \texttt{int}，所以
\texttt{a + b} 這個運算式本身的型別是 \texttt{int}（\texttt{sizeof} 量
到 $4$），運算過程中的 $300$ 完全不會溢位。 溢位其實發生在最後一步：
把 \texttt{int} 的 $300$ 窄化存回 \texttt{unsigned char} 時，依無符號
窄化規則模 $256$ 環繞，變成 $44$。 這一步是有定義的。 兩個
\texttt{unsigned char} 相加會不會溢位，因此要拆兩層回答：運算當下不
會，存回窄型別時才可能。

\subsection{溢位的安全檢查}

用固定寬度型別時，上下界是常數，可以在運算前先問會不會跨過邊界。

\begin{example}{c}{(overflow-check.c) 檢查整數溢位}
#include <stdio.h>
#include <stdint.h>
#include <limits.h>

int main(void) {
    int32_t a = INT32_MAX;
    int32_t b = 1;

    if (a > INT32_MAX - b) {
        printf("加法會溢位\n");
    } else {
        printf("結果：%d\n", a + b);
    }

    uint32_t x = UINT32_MAX / 2 + 1;
    uint32_t y = 3;
    if (x > UINT32_MAX / y) {
        printf("乘法會溢位\n");
    } else {
        printf("結果：%u\n", x * y);
    }
}
\end{example}

\iosample{
}{
加法會溢位 \\
乘法會溢位 \\
}

\texttt{INT32\_MAX - b} 先算出還能再加多少，再跟 \texttt{a} 比；乘法則
用除法避開先乘爆掉。 帶符號溢位是未定義行為，所以要在真正算出
\texttt{a + b} 之前攔下來。

\begin{drills}[label=drills:decl-stdint]
\item \texttt{promotion-surprise.c} 裡，\texttt{a + b} 這個運算式本身
的 \texttt{sizeof} 是多少？ 存回 \texttt{unsigned char c} 之後的值又是
多少？ 兩者為什麼不同？
\item 較窄的整數型別（例如 \texttt{short}）走進運算式時，會先被提升成
什麼型別？
\item \texttt{overflow-check.c} 用什麼技巧，在加法真的溢位之前就先攔下
來？
\end{drills}

\section{const 型別限定符}
\label{sec:decl-const}

\texttt{const} 表示這份物件不能透過這個名字改寫。 它可以標在變數、陣列
元素和指標上。

\subsection{const 參數}

\begin{code}{c}{const 函式參數}
void printString(const char *str) {
    printf("%s\n", str);
    // str[0] = 'X';   // compiler error
}

int sumArray(const int arr[], int n) {
    int sum = 0;
    for (int i = 0; i < n; i++) {
        sum += arr[i];
    }
    // arr[0] = 0;     // compiler error
    return sum;
}
\end{code}

參數標 \texttt{const}，是向呼叫者保證函式只讀、不改內容。

\subsection{const 陣列}

\begin{code}{c}{const 陣列}
const int days[] = {31, 28, 31, 30, 31, 30, 31, 31, 30, 31, 30, 31};
// days[0] = 30;       // compiler error

const char *months[] = {"Jan", "Feb", "Mar"};
\end{code}

\subsection{const 指標}

\texttt{const} 可以修飾指標本身，也可以修飾它指向的資料，因此有三種組
合。

\begin{code}{c}{三種 const 指標}
int x = 5;
int y = 8;

const int *p1 = &x;    // pointer to const data
p1 = &y;               // OK: pointer can move
// *p1 = 10;           // error: data is const

int * const p2 = &x;   // const pointer
// p2 = &y;            // error: pointer is const
*p2 = 10;              // OK: data can change

const int * const p3 = &x;  // both const
// p3 = &y;
// *p3 = 10;
\end{code}

這裡先從識別字往外走，系統練習見後文解讀複雜宣告。 \texttt{p1} 左邊是
\texttt{*}，所以它是指標；再往左是 \texttt{const int}，指向的資料不能
改。 \texttt{p2} 的 \texttt{const} 緊貼識別字，改不了指標本身。

\keypoint{從識別字往外讀。 另一種掃法是由右向左：\texttt{*} 左邊的
\texttt{const} 修飾指向的資料，緊貼識別字的 \texttt{const} 修飾指標本
身；兩套走完得到同一個意思。}

\subsection{跨檔案共享 const}

跨檔案共享一份 \texttt{const} 全域資料，有兩種常見寫法。

第一種：一個檔案負責定義，其他檔案用 \texttt{extern} 宣告。 整支程式
只有一份 \texttt{PI}。

\begin{code}{c}{extern 宣告一份全域常數}
/* file1.c: the single definition */
const double PI = 3.14159;

/* file2.c: a reference, no new object */
extern const double PI;
\end{code}

第二種：把定義寫進標頭，並加上 \texttt{static}。 每個包含這份標頭的
\texttt{.c} 各自有一份相同的副本。 資料是常數、內容又一樣，副本之間不
用通信，這樣就夠用。

\begin{code}{c}{標頭裡的 static const}
/* constants.h */
static const double PI = 3.14159;

/* file1.c and file2.c both include constants.h */
\end{code}

若省略 \texttt{static}，兩個檔案各自做出一份可連結的定義，標準不保證
能過。 要讓兩個檔案改同一格、用來通信，用第一種 \texttt{extern}；只是
共用讀得到的常數、不怕複製，用第二種。

\subsection{選讀：\texttt{volatile}、\texttt{restrict}、\texttt{\_Atomic}}

這三個限定符這一章認得出名字即可，主線用不到。

\texttt{volatile} 告訴編譯器：這個值可能被程式以外的東西改掉，例如記憶
體對映的硬體暫存器，或訊號處理函式寫入的旗標。 編譯器因此不會把讀取
優化成暫存複本。

\texttt{restrict} 標在指標上，意思是：在這段生命週期裡，這塊物件只從這
一條指標進出。 編譯器可以更大膽地重排讀寫。 兩個會重疊的緩衝區不要互
相標 \texttt{restrict}。

\texttt{\_Atomic}（C11）讓讀寫在執行緒之間不可切割，用來做旗標或計數
器這類小物件的同步。 完整的並行程式還需要記憶體順序與函式庫，這裡只
先認得這個限定符。

% TODO(future-chapter): inline 函式最多在此帶過——C99 的 function
% specifier，語意與多檔展開留給之後的章節

\begin{drills}[label=drills:decl-const]
\item \texttt{const int *p1} 跟 \texttt{int * const p2}，哪一個是「指
標本身不能改」，哪一個是「指向的資料不能改」？
\item 跨檔案共享一份 \texttt{const} 全域常數，若只是要各檔案各自讀取、
不需要通信，用 \texttt{extern} 那一種寫法還是 \texttt{static const} 那
一種？
\item \texttt{volatile} 通常用在什麼情境？
\end{drills}

\section{static 儲存類別}
\label{sec:decl-static}

幾個函式不小心共用同一個檔案層變數當迴圈計數器，會互相踩值。 下面這
段本意是印出 $10 \times 10$ 的星號。

\begin{example}{c}{(print-stars.c) 共用 i 只印出一列}
#include <stdio.h>

int i;

void print_one_row(void) {
    for (i = 1; i <= 10; i++) {
        printf("*");
    }
}

void print_all_rows(void) {
    for (i = 1; i <= 10; i++) {
        print_one_row();
        printf("\n");
    }
}

int main(void) {
    print_all_rows();
}
\end{example}

\iosample{
}{
********** \\
}

\texttt{print\_one\_row} 返回時 \texttt{i} 已經是 $11$。 外層
\texttt{for} 接著做 \texttt{i++} 再測 $i \le 10$，條件立刻失敗，所以
只印一列。 迴圈變數應宣告在各自的函式裡。

這個例子一次踩到識別字的三件事：\textbf{scope}（哪些程式碼看得到）、
\textbf{linkage}（能不能讓別的翻譯單元用同一個名字）、\textbf{storage
duration}（記憶體何時配置、何時釋放）。 \texttt{i} 寫在所有函式外面，
所以兩個函式都看得到它（file scope）；沒加 \texttt{static}，別的
\texttt{.c} 也可以用 \texttt{extern} 連過來（external linkage）；程式
一開始就有這一格，結束才釋放（static duration）。
\texttt{static} 出現的位置不同，改到的是其中不同的一項。

\subsection{三個詞}

區塊裡宣告的變數，預設是 block scope、沒有 linkage、automatic storage
duration：進入區塊才配置，離開就釋放。 寫在所有函式外面的變數，預設是
file scope、external linkage、static duration：程式一開始就在，整個檔
案都看得到，別的 \texttt{.c} 也可以用 \texttt{extern} 連過來。

\subsection{區塊裡的 static}

在區塊裡加上 \texttt{static}，scope 仍停在這個函式，但 duration 變成
static：只在程式啟動時初始化一次，之後每次呼叫都沿用同一格。

\begin{code}{c}{static 區域變數}
void counter(void) {
    static int count = 0;
    count++;
    printf("%d\n", count);
}

// 1st call: 1
// 2nd call: 2
// 3rd call: 3
\end{code}

遞迴時各層也共用這一格，不會每進一層就長出新的 \texttt{count}。 函式
可以回傳這格的位址，呼叫結束後位址仍然有效；一般的區域變數就不能這樣
做。

\subsection{檔案裡的 static}

寫在函式外面的 \texttt{static} 改的是 linkage：從 external 變成
internal，這個名字只在目前這個 \texttt{.c} 裡看得見。 duration 本來就
是 static，不必再靠關鍵字延長壽命。 開頭那個星號程式若把 \texttt{int
i;} 改成 \texttt{static int i;}，別的檔案就連不過來，但兩個函式仍然共
用同一格，一列星號的 bug 還在。 真正要在多個函式之間共享的狀態，用
file-scope \texttt{static} 可以，並給它一個說得出口的名字，不要拿迴圈
計數器來當共享變數。

幾個函式要共享狀態、又不想讓別的檔案碰到時，用 file-scope
\texttt{static} 正好。 下面是精簡版的堆疊：\texttt{contents} 與
\texttt{top} 對這個檔案裡的 \texttt{push}／\texttt{pop} 可見，對其他
\texttt{.c} 隱藏。

\begin{code}{c}{file-scope static 的堆疊}
#define STACK_SIZE 100

static int contents[STACK_SIZE];
static int top = 0;

void push(int i) {
    contents[top++] = i;
}

int pop(void) {
    return contents[--top];
}
\end{code}

正式程式還要檢查滿與空。 重點是：共享的狀態放在 file-scope
\texttt{static}，比把 \texttt{top} 做成誰都能改的 external 變數安全。

\texttt{static} 不只能標變數，也能標函式，效果一樣是把 linkage 限制在
目前這個檔案：

\begin{code}{c}{static 函式}
static int clamp(int x, int lo, int hi) {
    if (x < lo) {
        return lo;
    }
    if (x > hi) {
        return hi;
    }
    return x;
}
\end{code}

只給這個檔案內部用的輔助函式，習慣上都標 \texttt{static}：一來限制連
結，避免跟別的檔案裡同名的函式衝突；二來也是對讀者的提示，這個函式只
供本檔使用。 標頭檔裡通常只宣告不加 \texttt{static} 的函式，那些才是
這個模組真正想讓外界呼叫的介面。

\keypoint{\texttt{static} 改的是 linkage。 兩個檔案各自宣告同名的
file-scope \texttt{static} 變數或函式，是兩個獨立的東西，連結時不會衝
突。 拿掉 \texttt{static}，兩份 external linkage 的同名定義才會在連結
期打架。}

\texttt{extern} 如何讓多個檔案共用同一個變數或函式，將於之後的多檔／
連結章節詳述。

% TODO(future-chapter): static/dynamic linking、ABI、連結器報表、
% makefile 與標頭檔保護

\subsection{auto 與 register}

\texttt{auto} 是區塊變數的預設，寫出來與不寫相同，現在幾乎沒人用。
\texttt{register} 當年是請編譯器把變數放進 CPU 暫存器；現在編譯器自己
會選，而且標了 \texttt{register} 就不能取址。 當歷史遺跡認得即可。

\keypoint{\texttt{static} 在區塊裡延長壽命，在檔案層改成內部連結。
同一個關鍵字，看它寫在函式內還是函式外。}

\begin{drills}[label=drills:decl-static]
\item \texttt{print-stars.c} 那個例子為什麼只印出一列星星，而不是十
列？
\item 函式內宣告 \texttt{static int count = 0;}，跟不加 \texttt{static}
的區域變數比起來，差在哪裡？
\item 檔案層的變數前面加 \texttt{static}，改變的是 scope、linkage、還
是 storage duration？
\item 兩個不同的 \texttt{.c} 檔案各自宣告同名的 file-scope
\texttt{static} 變數，連結時會不會衝突？
\end{drills}

\section{解讀複雜宣告}
\label{sec:decl-complex-declarations}

陣列、函式、指標、\texttt{const} 前面都見過了。 這裡專門練
\texttt{*}、\texttt{[]}、\texttt{()} 疊在一起時該從哪裡讀起。

宣告裡的 \texttt{*}、\texttt{[]}、\texttt{()} 可以疊在一起。 規則只有
兩條：從被宣告的識別字往外讀；有選擇時，\texttt{[]} 與 \texttt{()} 優
先於 \texttt{*}。 括號可以改掉這個優先序。

\subsection{兩種等價的讀法}

由識別字往外讀，是上面給的規則。 Expert C Programming 給了另一個名字：
\textbf{順時針／螺旋法則}（clockwise/spiral rule）。 兩者是同一套規則
的兩種敘述方式，選一種練熟即可：

\begin{enumerate}
\item 從識別字出發。
\item 沿順時針方向繞：先看右邊，遇到 \texttt{[]} 就唸成：\ldots 的陣
列；遇到 \texttt{()} 就唸成：回傳 \ldots 的函式。 被括號擋住就先處理
括號裡面。
\item 繞不過去了（右邊已經讀完或被括號擋住）就轉看左邊：遇到
\texttt{*} 唸成：指向 \ldots 的指標；遇到 \texttt{const}／
\texttt{volatile} 就接上不可變這類限定。
\item 一路唸到最左邊的基底型別為止。
\end{enumerate}

這其實就是 \texttt{[]}、\texttt{()} 優先於 \texttt{*} 的具體操作方式
：優先權高的東西在螺旋裡先被繞到。

\subsection{暖身：陣列與指標}

\begin{code}{c}{最基本的兩個}
int arr[5];      // arr 是 5 個 int 的陣列
int *p;          // p 是指向 int 的指標
\end{code}

識別字往右先遇到什麼就是什麼，這兩個不用繞。

\subsection{指標陣列與陣列指標}

\begin{code}{c}{指標陣列}
int *ap[10];
\end{code}

識別字是 \texttt{ap}。 右邊是 \texttt{[10]}，左邊是 \texttt{*}；
\texttt{[]} 優先，所以 \texttt{ap} 先是長度 $10$ 的陣列，元素型別才是
指向 \texttt{int} 的指標。 \texttt{ap} 是十個 \texttt{int *}。

\begin{center}
\begin{tikzpicture}
    \node[cell] (t0) {\texttt{int}};
    \node[cell, right=0pt of t0] (t1) {\texttt{*}};
    \node[hlcell, right=0pt of t1] (t2) {\texttt{ap}};
    \node[cell, right=0pt of t2] (t3) {\texttt{[10]}};
    \node[idxlbl, below=0.25cm of t2] {起點};
    \draw[diagarrow] ($(t2.south)+(0.15,-0.35)$)
        to[out=-20,in=-160] node[below, font=\scriptsize] {1 右} ($(t3.south)+(0,-0.15)$);
    \draw[diagarrow] ($(t3.north)+(0,0.12)$)
        to[out=70,in=70] node[above, font=\scriptsize] {2 左} ($(t1.north)+(0,0.12)$);
\end{tikzpicture}
\diagcaption{\texttt{int *ap[10]} 的螺旋：從 \texttt{ap} 先往右讀
\texttt{[10]}（陣列），再往左讀 \texttt{*}（元素是指標），最後剩下
\texttt{int}。}
\label{fig:ch13-spiral-ap}
\end{center}

\begin{code}{c}{陣列指標}
int (*pa)[10];
\end{code}

括號把 \texttt{*} 圈起來，逼它先跟 \texttt{pa} 結合：\texttt{pa} 是指
標，指向一個長度 $10$ 的 \texttt{int} 陣列。 跟前一個對照著看：拿掉括
號、字元順序沒變，型別卻完全不同，這是這一節要練的核心分辨力。

\subsection{函式指標}

\begin{code}{c}{函式指標對照}
int *fp(int n);
int (*fp)(int n);
\end{code}

第一行：識別字 \texttt{fp} 右邊先遇到 \texttt{()}，\texttt{fp} 是函
式，接收一個 \texttt{int}；再往左是 \texttt{*}，回傳值是指標。 這行是
回傳 \texttt{int *} 的函式。

第二行：括號把 \texttt{*fp} 圈住，\texttt{fp} 先是指標；指標右邊才是
\texttt{()}，所以是指向函式的指標，該函式接收 \texttt{int}、回傳
\texttt{int}。 這一組是宣告複雜型別時最常見的分岔點：要不要加括號，差
的是函式回傳指標，或者指標指向函式。

\subsection[參數列：() 與 (void) 不一樣]{函式的參數列：\texttt{()} 與 \texttt{(void)} 不一樣}

函式宣告子裡的括號本身也有語意，不是可有可無的裝飾。 空括號
\texttt{()} 在 C 裡的原始意思是「參數的個數與型別未指定」——這是舊式
（K\&R 風格）C 遺留下來的寫法，編譯器看到 \texttt{()} 不會替你檢查呼
叫端傳的引數對不對。 \texttt{(void)} 才是明確宣告「這個函式不接收任
何參數」，呼叫端一旦傳了引數，編譯器會直接報錯。

\begin{code}{c}{() 不檢查參數，(void) 才會}
int f();       // 舊式寫法：參數未指定，呼叫端傳什麼都不會被擋
int g(void);   // 明確寫法：不接收參數，傳引數會編譯錯誤

int main(void) {
    f(1, 2, 3);   // 能通過編譯，即使 f 的定義沒有用到這些引數
    // g(1, 2, 3);  // 編譯錯誤：g 明明不吃參數
    return 0;
}
\end{code}

這就是本書從第一章開始，所有範例都寫 \texttt{int main(void)} 而不是
\texttt{int main()} 的原因：\texttt{main} 本身不接收任何命令列引數
時，\texttt{(void)} 讓這件事在宣告上就講清楚，而不是靠「反正沒人會
亂傳」這種默契。 其他自訂函式若不需要參數，也適用同一個原則。

\keypoint{函式括號裡的 \texttt{()} 代表「參數未指定」，\texttt{(void)}
代表「明確不接收參數」；兩者對編譯器的檢查力度不同，寫函式時一律用
\texttt{(void)} 取代空括號。}

\subsection{疊到最複雜}

\begin{code}{c}{指標、函式與 const}
char * const *(*next)();
\end{code}

依序走：

\begin{enumerate}
\item 識別字是 \texttt{next}。
\item \texttt{(*next)} 的括號讓 \texttt{*} 先綁住，\texttt{next} 是指
標。
\item 右邊是 \texttt{()}，所以它指向一個函式。
\item 函式左邊還有一個 \texttt{*}，回傳值是指標。
\item 最左邊是 \texttt{char * const}：回傳的是指向字元的 const 指標，
指標本身不能改，指向的字元可以。
\end{enumerate}

一整句：\texttt{next} 是指向函式的指標，該函式回傳 \texttt{char *
const}。 前面學過的 \texttt{static}、\texttt{extern}、\texttt{const}
若出現，一定寫在宣告最前面，先讀完儲存類別與限定符，再從識別字往外走
declarator。

\subsection{反過來：從描述寫出宣告}

讀懂之後，反向操作是更扎實的練習：給一句描述，自己組出宣告。 卡住的地
方，通常就是正著讀時會看漏的地方。

描述一：一個指向函式的指標陣列，這些函式不接收參數、回傳
\texttt{void}。

\begin{code}{c}{由描述寫出 handlers}
void (*handlers[10])();
\end{code}

先確定最外層是陣列（\texttt{handlers[10]}），每個元素是指向函式的指
標，所以元素外面包一層 \texttt{(*...)} 再接 \texttt{()}。

描述二：一個函式，接收一個 \texttt{int}，回傳指向長度 $5$ 的
\texttt{int} 陣列的指標。

\begin{code}{c}{由描述寫出 f}
int (*f(int n))[5];
\end{code}

先確定最外層是函式（接收 \texttt{int}），回傳值是指向陣列的指標，所以
整個 \texttt{(*f(int n))} 要再包一層才能接上 \texttt{[5]}。

\subsection{用 typedef 收短}

同一種複雜型別要出現好幾次時，先給它一個名字。

\begin{code}{c}{typedef 收短宣告}
typedef char * const *(*NextFn)();
NextFn next;
\end{code}

\texttt{typedef} 替既有型別取別名；讀法與變數宣告相同，最後得到的是型
別名。

\keypoint{由識別字往外讀，\texttt{[]} 與 \texttt{()} 優先於
\texttt{*}。 這套規則也叫順時針／螺旋法則（Expert C Programming 的說
法）：從識別字出發、順時針繞著讀，是同一條規則的另一種操作方式，讀不
確定的時候用它反過來驗算。}

\begin{drills}[label=drills:decl-complex-declarations]
\item \texttt{int *ap[10]} 跟 \texttt{int (*pa)[10]}，哪一個是「十個
指向 \texttt{int} 的指標」，哪一個是「指向一個長度 10 的 \texttt{int}
陣列的指標」？
\item \texttt{int f();} 跟 \texttt{int g(void);}，編譯器對呼叫端傳的引
數檢查力度有什麼不同？
\item 螺旋法則裡，\texttt{[]} 跟 \texttt{()} 的優先權比 \texttt{*} 高還
是低？
\item 給定「一個指向函式的指標陣列，這些函式不接收參數、回傳
\texttt{void}」，該怎麼宣告？
\end{drills}

\section{常見陷阱與最佳實踐}
\label{sec:decl-pitfalls}

共用檔案層迴圈變數的例子見 \texttt{static} 節開頭的星號程式。

\subsection{陷阱 1：混淆 printf 格式}

\begin{example}{c}{(bad-format.c) printf 格式不對}
#include <stdio.h>

int main(void) {
    long x = 10000000000L;
    printf("%d\n", x);
    printf("%ld\n", x);
}
\end{example}

\iosample{
}{
1410065408 \\
10000000000 \\
}

\keypoint{未定義行為：\texttt{printf} 格式指定符與引數型別不符。 這裡
印出的數字只是這次編譯的結果，換機器或最佳化層級不必相同。}

\texttt{\%d} 只把可變參數當 \texttt{int} 讀，在 \texttt{long} 比
\texttt{int} 寬的常見平台上，等於只讀走了 \texttt{x} 位元模式的低 $32$
位元。 這裡印出的 \texttt{1410065408} 剛好是 $10000000000$ 的低 $32$
位元當有號數解讀的結果，換一種呼叫慣例或架構未必是這個數字。 編譯器
不一定能在編譯期抓到，開 \texttt{-Wformat} 才會提醒。

\subsection{陷阱 2：帶符號與無符號混用}

\texttt{int} 與 \texttt{unsigned int} 出現在同一個比較裡，usual
arithmetic conversions 會把 \texttt{int} 轉成 \texttt{unsigned int}。
\texttt{-1} 的位元模式當成無號來看，是該寬度能表示的最大正數。

\begin{example}{c}{(signed-unsigned-cmp.c) 有號與無號比較}
#include <stdio.h>

int main(void) {
    int si = -1;
    unsigned int ui = 10;

    if (si < ui) {
        printf("si < ui\n");
    } else {
        printf("si >= ui\n");
    }
    printf("-1 as unsigned: %u\n", (unsigned int)si);
}
\end{example}

\iosample{
}{
si >= ui \\
-1 as unsigned: 4294967295 \\
}

直覺以為 $-1 < 10$，實際走的是 $4294967295 < 10$，條件為假。 編譯器開
\texttt{-Wsign-compare} 會對這一行發出警告。 前一節的提升規則解釋了為
什麼 \texttt{printf("\%d", short\_var)} 仍然正確：可變引數會先把
\texttt{short} 升成 \texttt{int}。

\keypoint{未定義行為：有號與無號混在同一個比較裡，卻仍把結果當成有號
去讀。 這裡印出的數字只是這次編譯的結果，換機器或最佳化層級不必相同。}

轉換本身（有號轉成無號）是有定義的：在 $32$ 位元 \texttt{unsigned
int} 上，\texttt{-1} 一定變成 $4294967295$。 學習要點針對的是把這個
無號值繼續當成 $-1$ 使用；寬度一變，印出來的十進位就會跟著變。

\subsection{陷阱 3：溢位無警告}

\begin{example}{c}{(silent-overflow.c) 迴圈裡累加，溢位沒有任何警告}
#include <stdio.h>

int main(void) {
    int total = 0;
    for (int i = 0; i < 1000000; i++) {
        total += 100000;
    }
    printf("total = %d\n", total);
}
\end{example}

\iosample{
}{
total = 1215752192 \\
}

\keypoint{未定義行為：迴圈裡帶符號累加溢位。 這裡印出的數字只是這次編
譯的結果，換機器或最佳化層級不必相同。}

正確答案是 $100000 \times 1000000 = 100000000000$，遠超過
\texttt{INT\_MAX}。 這裡的結果來自二補數環繞這個常見（但標準不保證）
的實作行為，剛好落在一個看起來完全合理的正整數。 程式沒有當掉、沒有
任何警告，只是安靜地給了一個錯得離譜的答案。 崩潰至少會被發現，這種
錯誤不會。

\subsection{陷阱 4：未初始化}

區塊裡的自動變數若沒寫初值，內容是垃圾；讀它是未定義行為，不要靠印出
來的數字猜。 檔案層或 \texttt{static} 變數沒寫初值時，保證是 $0$（指
標則是空指標）。

\begin{code}{c}{未初始化的 auto 與 static}
void f(void) {
    int junk;            // indeterminate
    static int zeros;    // 0
}
\end{code}

\stylepoint{自動變數與 \texttt{static} 都寫明初值。 後者預設為 $0$ 只
是語言保證，初值仍應寫明。}

\subsection{陷阱 5：巢狀 const 轉換不安全}

前面三種 \texttt{const} 指標，多一層 \texttt{*} 之後不能照直覺往外加
\texttt{const}。

\begin{code}{c}{危險：以為 const int ** 能收 int **}
void printAll(const int **arr, int n) {
    for (int i = 0; i < n; i++) {
        printf("%d\n", *arr[i]);
    }
}

int *data[3];
// printAll(data, 3);
// compiler rejects: int ** cannot convert to const int **
\end{code}

直覺以為多包一層 \texttt{const} 只會更安全，實際上 \texttt{int **} 轉
成 \texttt{const int **} 並不保證安全，C 標準因此不讓它隱式發生。 如果
放行，函式內部可以合法寫 \texttt{arr[0] = \&something\_non\_const;}，
把一個原本不該進去的位址塞進 \texttt{data}；之後再透過外面那份
non-const 的 \texttt{data} 去改到理論上受保護的東西，\texttt{const}
的保證就被繞過了。

要轉換就得寫明確的 cast。 強制轉型一旦放行，後面那次改寫就不再受標準
保護：

\keypoint{未定義行為：用強制轉型把 \texttt{int **} 塞進 \texttt{const
int **}，再從非 const 路徑寫入。 這裡若因此印出某個數字，那個數字只是
這次編譯的結果，換機器或最佳化層級不必相同。}

一開始就把資料結構設計成本來就是 \texttt{const}，可以避開這條路徑。

\subsection{最佳實踐}

\begin{itemize}
\item 跨平臺要固定寬度時，優先用 \texttt{int32\_t}、\texttt{uint64\_t}
\item 算術前用上下界檢查，尤其是帶符號整數
\item \texttt{\%d}／\texttt{\%ld}／\texttt{\%lld}／\texttt{\%u} 與型別對準
\item 有號與無號不要放進同一個比較
\item 編譯時打開 \texttt{-Wall -Wextra -Wconversion -Wsign-compare}
\end{itemize}

\begin{drills}[label=drills:decl-pitfalls]
\item \texttt{bad-format.c} 用 \texttt{\%d} 印一個 \texttt{long} 值，為
什麼會印出錯誤的數字？ 該改用哪個格式指定符？
\item \texttt{signed-unsigned-cmp.c} 裡，\texttt{si < ui}（\texttt{si =
-1}，\texttt{ui} 是 \texttt{unsigned}）為什麼判斷結果是假？
\item 為什麼迴圈裡帶符號整數累加溢位，比程式直接當掉更危險？
\item 為什麼 \texttt{int **} 不能隱式轉換成 \texttt{const int **}？
\end{drills}

\addcontentsline{toc}{section}{練習解答}
\section*{練習解答}

各節練習的答案收在這裡。 短題只給能一句話核對的結果。

\paragraph{第~\ref{sec:decl-int-sizes}~節}
\begin{enumerate}
\item 通常 2 個位元組，範圍約 $-32768$ 到 $32767$。
\item 印出 \texttt{-2147483648}；這只是這台機器、這個編譯器當下呈現的
一種可能，不是標準保證的行為（帶符號溢位是未定義行為）。
\item 標準保證環繞成 $0$（模 $2^n$）；帶符號溢位則是未定義行為，不保
證環繞。
\end{enumerate}

\paragraph{第~\ref{sec:decl-signed-unsigned}~節}
\begin{enumerate}
\item \texttt{unsigned int} 是 $0$ 到 $2^n - 1$；\texttt{signed int} 是
$-2^{n-1}$ 到 $2^{n-1} - 1$。
\item \texttt{\%d} 會把底下的位元當有號解讀，可能印出負數或錯誤的值；
該用 \texttt{\%u}。
\item 不會區分；一律把底下的位元當無號數排版，所以負數的十六進位看到
的是它的二補數表示法。
\end{enumerate}

\paragraph{第~\ref{sec:decl-stdint}~節}
\begin{enumerate}
\item \texttt{sizeof(a + b)} 是 4（\texttt{int}）；存回 \texttt{c} 之後
是 44。 因為運算當下 \texttt{a}、\texttt{b} 先被提升成 \texttt{int}，
$300$ 沒有溢位；溢位是最後窄化存回 \texttt{unsigned char} 時，依環繞
規則模 256 才發生的。
\item 提升成 \texttt{int}（裝不下才升成 \texttt{unsigned int}）。
\item 用減法／除法先算出還能加或乘多少，跟現有值比較，在真正執行加法
或乘法之前就先判斷會不會超過上界。
\end{enumerate}

\paragraph{第~\ref{sec:decl-const}~節}
\begin{enumerate}
\item \texttt{int * const p2} 是指標本身不能改；\texttt{const int *p1}
是指向的資料不能改。
\item 只是各檔案各自讀取、不需要通信，用 \texttt{static const} 放進標
頭讓每個檔案各自一份；要多檔案共用同一份、能通信，才用
\texttt{extern} 搭配單一定義。
\item 標記可能被程式以外的東西改掉的值，例如硬體暫存器或訊號處理函式
寫入的旗標，防止編譯器把讀取優化成暫存複本。
\end{enumerate}

\paragraph{第~\ref{sec:decl-static}~節}
\begin{enumerate}
\item 因為 \texttt{i} 是共用的檔案層變數，\texttt{print\_one\_row} 執行
完一輪後 \texttt{i} 已經是 11，外層 \texttt{for} 的 \texttt{i++} 與條件
測試接著失敗，迴圈只跑了一次。
\item 一般區域變數每次呼叫都重新配置、重新歸零；\texttt{static} 區域變
數只在程式啟動時初始化一次，之後每次呼叫都沿用同一格，數值會累積。
\item 改變的是 linkage（從 external 變成 internal），storage duration
本來就是 static，不需要靠 \texttt{static} 延長。
\item 不會；兩個檔案各自的 file-scope \texttt{static} 是各自獨立的內部
連結，不會在連結期衝突。
\end{enumerate}

\paragraph{第~\ref{sec:decl-complex-declarations}~節}
\begin{enumerate}
\item \texttt{int *ap[10]} 是十個指向 \texttt{int} 的指標（陣列優
先）；\texttt{int (*pa)[10]} 是指向長度 10 的 \texttt{int} 陣列的指標
（括號讓 \texttt{*} 先跟識別字結合）。
\item \texttt{int f();} 是舊式寫法，參數未指定，編譯器不會檢查呼叫端
傳的引數；\texttt{int g(void);} 明確表示不接收參數，傳引數會編譯錯
誤。
\item \texttt{[]} 與 \texttt{()} 的優先權比 \texttt{*} 高。
\item \texttt{void (*handlers[10])();}。
\end{enumerate}

\paragraph{第~\ref{sec:decl-pitfalls}~節}
\begin{enumerate}
\item \texttt{\%d} 只把可變引數當 \texttt{int} 讀，在 \texttt{long} 比
\texttt{int} 寬的平台上只會讀走低 32 位元；該用 \texttt{\%ld}。
\item \texttt{si} 與 \texttt{ui} 比較時，\texttt{int} 會被轉成
\texttt{unsigned int}，\texttt{-1} 的位元模式當無號數解讀是該寬度能表
示的最大正數，遠大於 10，所以 \texttt{si < ui} 為假。
\item 因為它不會產生任何警告或當機訊號，只是安靜地給出一個看起來合理
但錯得離譜的答案，比崩潰更難被發現。
\item 因為若放行，函式內部可以透過那個 \texttt{const int **} 把一個
non-const 的位址塞進原本的 \texttt{int **}，再經由外面那份 non-const
指標寫入，繞過 \texttt{const} 的保護，所以標準不讓它隱式發生。
\end{enumerate}

% 變更摘要
% - 修訂：'\textbackslash 0' 斷行、\{0\} 取代空初始化列、結構陣列
%   兌現開場伏筆、末尾 padding 理由、巢狀 Date、複合字面值改保守寫法
% - 螺旋法則／巢狀 const 陷阱不在本章（回第 13 章）
% TODO(future-chapter): union / enum / bit-field 細節
% TODO(future-chapter): 標頭檔保護、makefile

\chapter{結構}
\label{ch:structure}

一筆學生資料同時有學號、姓名、GPA，三種型別。 前面學過的陣列一次只能
裝同一種型別，沒辦法把這三樣自然地綁在同一筆紀錄裡。 結構（struct）
就是為這件事準備的：把不同型別的欄位捆成一個單位，之後傳遞、複製、
取址都針對這一整包。

這一章從為什麼需要它講起，再到定義、初始化、欄位、指標、參數傳遞。 寫
結構時若碰到 \texttt{*}、\texttt{[]}、\texttt{()} 疊在一起，讀法見第
13 章後段的螺旋法則。

\section{為什麼需要 struct}
\label{sec:struct-motivation}

假設班上有 $100$ 個學生，每個學生要記學號、姓名、GPA。 沒有結構的時
候，常見寫法是開三個平行陣列：

\begin{code}{c}{三個平行陣列硬湊一筆紀錄}
char ids[100][10];
char names[100][50];
double gpas[100];

strcpy(ids[0], "R123");
strcpy(names[0], "Ada");
gpas[0] = 3.8;
\end{code}

\texttt{ids[i]}、\texttt{names[i]}、\texttt{gpas[i]} 必須永遠對準同一
個人。 編譯器不會幫你核對第 $i$ 格對應同一個人：排序時若只動了
\texttt{gpas}、忘了跟著交換姓名，成績就會貼到別人身上。 想把一筆學生
當成函式的一個引數傳進去，也得一次傳三個陣列加一個索引，很容易漏。

結構把這三樣寫進同一個型別，一筆學生就是一個變數：

\begin{code}{c}{一筆學生綁成一個單位}
struct Student {
    char id[10];
    char name[50];
    double gpa;
};

struct Student s;
strcpy(s.id, "R123");
strcpy(s.name, "Ada");
s.gpa = 3.8;
\end{code}

之後排序、傳遞、放進陣列，動的都是這一整包。 核心就是：不同型別的資
料綁在一起、當成一個單位。

\begin{drills}[label=drills:struct-motivation]
\item 用三個平行陣列存學生資料時，排序或交換資料為什麼容易出錯？
\item \texttt{struct Student} 把哪三種型別的欄位綁成一個單位？
\item 把一筆學生資料當函式引數傳遞，平行陣列版要傳幾個東西？struct
版呢？
\end{drills}

\section{定義與宣告}
\label{sec:struct-declare}

\subsection{結構語法}

\texttt{struct} 後面先寫一個名字，叫做 \textbf{tag}，接著用大括號列出
成員。 tag 讓這份佈局有一個型別名字，之後宣告變數都寫
\texttt{struct} 加 tag。

\begin{code}{c}{帶 tag 的結構}
struct Point {
    int x;
    int y;
};

struct Student {
    char id[10];
    char name[50];
    double gpa;
};

struct Point p = {10, 20};
struct Student s = {"R123", "Ada", 3.8};
\end{code}

\texttt{struct Point} 有兩個 \texttt{int}；\texttt{struct Student} 有
字元陣列和一個 \texttt{double}。 宣告變數時，\texttt{struct Point p}
跟 \texttt{int n} 同一類事情：左邊是型別，右邊是變數名字。 大括號那
份是初值，成員依宣告順序對進去；下一節再把初始化的細節拆開。

\subsection{tag 與匿名結構}

tag 的作用是：這份佈局從此有一個可以反覆使用的型別名字。 省掉 tag、
直接在大括號後面宣告變數，叫做匿名結構。 兩個匿名結構即使成員列出來
完全一樣，仍是兩種型別，不能互相賦值。

\begin{code}{c}{長得一樣的匿名結構仍是兩種型別}
struct { int x; int y; } a = {1, 2};
struct { int x; int y; } b;
// b = a;   // compiler error: incompatible types
\end{code}

直覺以為版面相同就可以互相抄，C 比的是型別名字，不是成員長得像不像。
要讓兩個變數能互相賦值，兩邊都得寫同一個 tag，例如都是
\texttt{struct Point}。

\subsection{用 typedef 省略 struct}

每次都寫 \texttt{struct Point} 有點長。 \texttt{typedef} 可以替這個型
別取一個較短的名字，之後當普通型別用：

\begin{code}{c}{typedef 讓 tag 跟簡寫同一件事}
typedef struct Point Point;
typedef struct Student Student;

Point p = {10, 20};
Student s = {"R123", "Ada", 3.8};
\end{code}

\texttt{struct Point} 與 \texttt{Point} 指同一種型別，只是書寫長短不
同。 也可以在定義時一次寫完，大括號後面直接接別名：

\begin{code}{c}{定義時順便 typedef}
typedef struct {
    int x, y;
} Point;
\end{code}

這種寫法沒有對外的 tag，別名 \texttt{Point} 就是這個型別的名字。 後
文為了跟前面的 tag 討論對上，多數例子仍寫 \texttt{struct Point}；真
正寫程式時，\texttt{typedef} 過的短名比較常見。

\subsection{大小、對齊與填補}

\texttt{sizeof(struct Point)} 在常見平臺上是 $8$，剛好兩個
\texttt{int}。 成員型別一混，量出來的大小往往比各成員相加還大：編譯
器會在中間或末尾插入填補（padding）位元組，讓每個成員的起始位址符合
對齊（alignment）。 許多 CPU 讀 \texttt{int} 時希望位址是 $4$ 的倍
數，讀指標時希望是 $8$ 的倍數；對不齊就得額外填空。

下面這個結構混了指標、\texttt{short}、\texttt{int}。 在 $64$ 位元這
台機器上跑出來的位移如下。

\begin{example}{c}{(slot-layout.c) 用 offsetof 量位移}
#include <stdio.h>
#include <stddef.h>

struct Compact {
    char *name;
    short kind;
    int uses;
};

struct Spread {
    short kind;
    char *name;
    int uses;
};

int main(void) {
    printf("Compact sizeof=%zu\n", sizeof(struct Compact));
    printf("  name offset=%zu size=%zu\n",
           offsetof(struct Compact, name), sizeof(char *));
    printf("  kind offset=%zu size=%zu\n",
           offsetof(struct Compact, kind), sizeof(short));
    printf("  uses offset=%zu size=%zu\n",
           offsetof(struct Compact, uses), sizeof(int));

    printf("Spread sizeof=%zu\n", sizeof(struct Spread));
    printf("  kind offset=%zu size=%zu\n",
           offsetof(struct Spread, kind), sizeof(short));
    printf("  name offset=%zu size=%zu\n",
           offsetof(struct Spread, name), sizeof(char *));
    printf("  uses offset=%zu size=%zu\n",
           offsetof(struct Spread, uses), sizeof(int));
}
\end{example}

\iosample{
}{
Compact sizeof=16 \\
\ \ name offset=0 size=8 \\
\ \ kind offset=8 size=2 \\
\ \ uses offset=12 size=4 \\
Spread sizeof=24 \\
\ \ kind offset=0 size=2 \\
\ \ name offset=8 size=8 \\
\ \ uses offset=16 size=4 \\
}

\texttt{struct Compact} 的成員大小加起來是 $8 + 2 + 4 = 14$，整份卻
是 $16$。 中間空掉的 $2$ 位元組，是為了讓 \texttt{int uses} 從位移
$12$ 開始（$4$ 的倍數）：

\begin{center}
\begin{tabular}{|l|r|r|}
\hline
\textbf{成員} & \textbf{位移} & \textbf{大小} \\
\hline
\texttt{name} & $0$ & $8$ \\
\texttt{kind} & $8$ & $2$ \\
填補 & $10$ & $2$ \\
\texttt{uses} & $12$ & $4$ \\
\hline
合計 &  & $16$ \\
\hline
\end{tabular}
\end{center}

\texttt{struct Spread} 三個成員邏輯上存的資料量相同，只是把
\texttt{short} 搬到最前面。 指標要對齊到 $8$，\texttt{kind} 後面就得
先空 $6$ 格。 \texttt{uses} 到位移 $20$ 就結束了，但整份仍補到 $24$：
結構若放進陣列，下一個元素的起始位址必須同樣對齊到 $8$ 的倍數。 末尾
那 $4$ 格填補，就是為了這件事。

\begin{center}
\begin{tabular}{|l|r|r|}
\hline
\textbf{成員} & \textbf{位移} & \textbf{大小} \\
\hline
\texttt{kind} & $0$ & $2$ \\
填補 & $2$ & $6$ \\
\texttt{name} & $8$ & $8$ \\
\texttt{uses} & $16$ & $4$ \\
填補 & $20$ & $4$ \\
\hline
合計 &  & $24$ \\
\hline
\end{tabular}
\end{center}

\keypoint{換成員的型別或順序，整個 struct 的大小可能跟著變，即使邏輯
上存的資料量沒變。}

這些數字依架構而異；寫進課文的是這台 $64$ 位元機器的結果。 更細的對
齊規則，後面高等技術那一章會再講一次。

\begin{drills}[label=drills:struct-declare]
\item 兩個匿名結構即使成員列的型別、順序完全一樣，能不能互相賦值？
為什麼？
\item \texttt{typedef struct Point Point;} 之後，\texttt{struct Point}
跟 \texttt{Point} 是不是同一種型別？
\item \texttt{struct Compact} 三個成員的大小加起來是 $14$ 位元組，但
\texttt{sizeof(struct Compact)} 卻是 $16$，多出來的 $2$ 位元組在做
什麼？
\item 把 \texttt{short} 搬到結構最前面（\texttt{Spread} 版），為什麼
整個結構反而變大？
\end{drills}

\section{初始化}
\label{sec:struct-init}

\subsection{完整與部分初始化}

大括號依成員宣告的順序填值。 全部都給，叫做完整初始化。 只給前面幾
個，叫做部分初始化：大括號裡至少要有一個值，沒寫到的成員補 $0$。 想
讓每個成員都是 $0$，寫 \texttt{\{0\}} 即可，第一個成員得到 $0$，其餘
也補 $0$。 空的 \texttt{\{\}} 要到 C23 才是標準寫法，本書用
\texttt{\{0\}}。

\begin{code}{c}{完整與部分初始化}
struct Point p1 = {5, 10};   // x = 5, y = 10
struct Point p2 = {5};       // x = 5, y = 0
struct Point p3 = {0};       // x = 0, y = 0

struct Student s = {"R123"};
// id is "R123"; name every byte is 0, hence empty string
// gpa is 0.0
\end{code}

\texttt{char} 陣列補 $0$ 的意思要講清楚：每個元素都是字元
\texttt{'\textbackslash 0'}，所以 \texttt{s.name} 是空字串，
\texttt{printf("\%s", s.name)} 什麼都不會印。 這跟整數成員補 $0$ 是同
一條規則，只是套在字元陣列上剛好變成空字串。

\subsection{標籤化初始化}

C99 可以用 \texttt{.成員 = 值} 依名字對應，不必背宣告順序，也允許跳
過中間的成員。 這種寫法叫標籤化初始化（designated initializer）。

\begin{example}{c}{(designated-init.c) 用名字對應初值}
#include <stdio.h>

struct Student {
    char id[10];
    char name[50];
    double gpa;
};

int main(void) {
    struct Student s = {.gpa = 3.8, .id = "R123"};
    printf("id = %s\n", s.id);
    if (s.name[0] == '\0') {
        printf("name is empty\n");
    }
    printf("gpa = %.1f\n", s.gpa);
}
\end{example}

\iosample{
}{
id = R123 \\
name is empty \\
gpa = 3.8 \\
}

大括號裡先寫 GPA 再寫學號，照樣對得上；沒提到的 \texttt{name} 仍補
$0$，所以第一個字元是 \texttt{'\textbackslash 0'}，條件成立。 之後若
有人把 \texttt{gpa} 搬到結構的第一個成員，這份初值不用改；依位置寫的
\texttt{\{"R123", "Ada", 3.8\}} 就會全部錯位。

\stylepoint{成員一多、或以後可能調順序，優先用 \texttt{.成員 = 值}。}

\begin{drills}[label=drills:struct-init]
\item \texttt{struct Point p2 = \{5\};} 之後，\texttt{p2.y} 的值是多
少？
\item \texttt{struct Student s = \{"R123"\};} 之後，\texttt{s.name} 是
什麼樣的字串？ 為什麼不是垃圾值？
\item 用 \texttt{.gpa = 3.8, .id = "R123"} 這種標籤化初始化，跟依宣告
順序寫的 \texttt{\{"R123", "Ada", 3.8\}} 比，好處是什麼？
\end{drills}

\section{欄位存取}
\label{sec:struct-access}

\subsection{點運算子}

對結構變數用 \texttt{.} 取出或指定成員：

\begin{code}{c}{用點運算子讀寫欄位}
struct Point p;
p.x = 10;
p.y = 20;
printf("%d %d\n", p.x, p.y);
\end{code}

\texttt{p.x} 就是 \texttt{p} 裡面那格叫做 \texttt{x} 的 \texttt{int}，
可以出現在指定的左邊，也可以當普通整數拿去印。

\subsection{取址與優先權}

成員也佔記憶體，所以可以對它取址。 \texttt{\&p.y} 這種寫法要看運算子
優先權：\texttt{.} 比 \texttt{\&} 高，所以是先取出成員 \texttt{y}，再
對那一格取址。 整個結構的位址是 \texttt{\&p}。

\begin{example}{c}{(member-addr.c) 成員 y 相對整包的位移}
#include <stdio.h>

struct Point {
    int x, y;
};

int main(void) {
    struct Point p = {10, 20};
    printf("offset of y = %td\n", (char *)&p.y - (char *)&p);
}
\end{example}

\iosample{
}{
offset of y = 4 \\
}

\texttt{x} 在最前面，佔 $4$ 位元組，所以 \texttt{y} 從位移 $4$ 開始，
\texttt{\&p.y} 跟 \texttt{\&p} 差 $4$。 若把 \texttt{y} 寫成第一個成
員，兩個位址就會相同，對照就看不出差別，所以這裡用第二個成員。

把括號加錯邊，語意就斷了：

\begin{code}{c}{括號加在 p 外面編不過}
struct Point p = {10, 20};
printf("%d\n", p.y);       // 20
// printf("%d\n", (&p).y); // compiler error: (&p) is a pointer
\end{code}

\texttt{(\&p)} 的型別是 \texttt{struct Point *}，指標沒有點運算子可
用，要寫 \texttt{(\&p)->y} 或乾脆 \texttt{p.y}。

\subsection{整包賦值與不能比較}

結構可以用 \texttt{=} 一次複製所有成員：

\begin{example}{c}{(assign-copy.c) 整包賦值}
#include <stdio.h>

struct Point {
    int x, y;
};

int main(void) {
    struct Point a = {3, 4};
    struct Point b;
    b = a;
    printf("%d %d\n", b.x, b.y);
}
\end{example}

\iosample{
}{
3 4 \\
}

\texttt{b = a} 之後，\texttt{b.x}、\texttt{b.y} 各自等於 \texttt{a}
對應的成員。 這是逐成員複製，兩個變數仍是兩塊獨立的記憶體；改
\texttt{b.x} 不會動到 \texttt{a.x}。

相等比較沒有同樣的語法。 下面這行會被編譯器拒絕：

\begin{code}{c}{結構沒有 ==}
struct Point a = {1, 2};
struct Point b = {1, 2};
// if (a == b) {        // compiler error
//     printf("same\n");
// }
if (a.x == b.x && a.y == b.y) {
    printf("same\n");
}
\end{code}

直覺以為結構跟 \texttt{int} 一樣能用 \texttt{==}，C 沒有為結構定義這
個運算子，要自己比成員。 \texttt{!=} 同樣沒有。 成員裡若有陣列，還得
逐格比。 \texttt{memcmp} 比的是整塊位元組，含填補；填補裡的內容未經
定義，兩個成員完全相同的結構用 \texttt{memcmp} 比，仍可能得到不相等。

\subsection{記憶體佈局}

前面 \texttt{slot-layout.c} 已經用位移說明填補落在哪。 寫註解時，位移
要跟對齊規則對得上，不要只抄一個看起來合理的數字：

\begin{code}{c}{位移註解要對齊規則}
struct Compact {
    char *name;   // offset 0, 8 bytes
    short kind;   // offset 8, 2 bytes
    /* 2 bytes padding: next int wants a multiple of 4 */
    int uses;     // offset 12, 4 bytes
};
/* sizeof == 16 on this 64-bit machine */
\end{code}

\texttt{name} 是指標，對齊 $8$，所以結構整體也按 $8$ 對齊；
\texttt{kind} 後面那 $2$ 格空著，是為了 \texttt{uses}。 理由與上一節
的表相同。

\subsection{結構陣列}

開場那個班上 $100$ 個學生，現在可以真的放進陣列。 每個元素都是一整包
\texttt{struct Student}，交換兩個元素時，學號、姓名、GPA 一起走：

\begin{example}{c}{(class-swap.c) 整包交換，成績跟姓名對在一起}
#include <stdio.h>
#include <string.h>

struct Student {
    char id[10];
    char name[50];
    double gpa;
};

int main(void) {
    struct Student class[100];
    strcpy(class[0].id, "R123");
    strcpy(class[0].name, "Ada");
    class[0].gpa = 3.8;
    strcpy(class[1].id, "R200");
    strcpy(class[1].name, "Bob");
    class[1].gpa = 2.1;

    struct Student tmp = class[0];
    class[0] = class[1];
    class[1] = tmp;

    printf("%s %.1f\n", class[0].name, class[0].gpa);
    printf("%s %.1f\n", class[1].name, class[1].gpa);
}
\end{example}

\iosample{
}{
Bob 2.1 \\
Ada 3.8 \\
}

\texttt{class[0].gpa} 是第 $0$ 個學生的 GPA；\texttt{strcpy} 寫進
\texttt{class[1].name}。 用臨時變數把 \texttt{class[0]} 與
\texttt{class[1]} 對調之後，Bob 的 $2.1$ 仍貼在 Bob 身上，Ada 的
$3.8$ 仍貼在 Ada 身上。 開場那種只換 \texttt{gpas[i]}、忘了換
\texttt{names[i]} 的錯法，在這裡寫不出來。

\begin{drills}[label=drills:struct-access]
\item \texttt{(\&p).y} 這樣寫為什麼編譯不過？ 正確寫法該怎麼寫？
\item 兩個結構能不能直接用 \texttt{==} 比較？ 該怎麼做才對？
\item \texttt{struct Student class[100];} 整包交換 \texttt{class[0]}
和 \texttt{class[1]}，跟只交換 \texttt{gpa} 一個欄位比，為什麼前者
才不會讓姓名和成績錯位？
\end{drills}

上一節 \texttt{Spread} 末尾那 $4$ 格填補，就是為了 \texttt{class[1]}
的起始位址仍是 $8$ 的倍數。 下一節的指標會再碰到：\texttt{\&class[0]}
是第一包的位址，\texttt{ptr + 1} 一次跨過一整個結構，落到
\texttt{class[1]}。

\section{結構指標}
\label{sec:struct-pointer}

\subsection{宣告與箭頭}

結構變數在 stack 上時，可以取其位址；也可以用 \texttt{malloc} 在
heap 上要一塊同樣大小的空間，得到指向結構的指標。 第 9 章已經用過
\texttt{malloc}／\texttt{free}，這裡只是把型別換成結構。

\begin{code}{c}{指到既有變數，或自己配置}
struct Point p = {10, 20};
struct Point *q = &p;     // points at stack object

struct Point *r = malloc(sizeof(struct Point));
r->x = 10;
r->y = 20;
free(r);
r = NULL;
\end{code}

透過指標讀寫成員，有三種看起來不一樣、意思相同的寫法：

\begin{code}{c}{三種讀法並排}
struct Point p = {10, 20};
struct Point *ptr = &p;

printf("%d\n", p.x);        // the variable itself
printf("%d\n", (*ptr).x);   // dereference, then dot
printf("%d\n", ptr->x);     // arrow: same as (*ptr).x
\end{code}

\texttt{(*ptr).x} 的括號不能省：\texttt{.} 優先權高於 \texttt{*}，寫
成 \texttt{*ptr.x} 會先當成 \texttt{ptr.x} 再解參考，而 \texttt{ptr}
是指標、沒有成員 \texttt{x}。 \texttt{ptr->x} 就是為了少寫那對括號。

\keypoint{\texttt{ptr->field} 等同 \texttt{(*ptr).field}。 變數本身用
點，指標用箭頭。}

\subsection{指標成員的配置與釋放}

結構裡若有指標成員，外層那一次 \texttt{malloc} 只配置結構本身。 內部
的 \texttt{char *} 此時仍是垃圾，還沒有指向任何字串；要再為那串字另
外配置，結束時也要先放掉內層、再放掉外層。 只放外層，heap 上的字串會
留著，而指向它的唯一位址已經沒了。

下面用一筆短片紀錄當例子：標題是動態字串，時長是整數。 三個函式分別
負責配置、印出、釋放。

\begin{example}{c}{(clip-heap.c) 內層字串要自己配置與釋放}
#include <stdio.h>
#include <stdlib.h>
#include <string.h>

typedef struct {
    char *title;
    int seconds;
} Clip;

Clip *clip_create(const char *title, int seconds) {
    Clip *c = malloc(sizeof(Clip));
    if (c == NULL) {
        return NULL;
    }
    c->title = malloc(strlen(title) + 1);
    if (c->title == NULL) {
        free(c);
        return NULL;
    }
    strcpy(c->title, title);
    c->seconds = seconds;
    return c;
}

void clip_print(const Clip *c) {
    printf("%s (%d s)\n", c->title, c->seconds);
}

void clip_destroy(Clip *c) {
    if (c == NULL) {
        return;
    }
    free(c->title);
    free(c);
}

int main(void) {
    Clip *c = clip_create("sunrise", 12);
    clip_print(c);
    clip_destroy(c);
}
\end{example}

\iosample{
}{
sunrise (12 s) \\
}

\texttt{clip\_create} 先要一塊 \texttt{Clip}，再要
\texttt{strlen(title) + 1} 個字元（多那一格放
\texttt{'\textbackslash 0'}），最後才 \texttt{strcpy}。 若第二次
\texttt{malloc} 失敗，已經要來的外層要先 \texttt{free}，否則那塊結構
自己也會漏。 \texttt{clip\_destroy} 的順序相反：先放
\texttt{c->title}，再放 \texttt{c}；反過來的話，\texttt{c} 沒了就拿不
到內層位址。

關係可以想成兩層：\texttt{c} 這顆指標在 stack 上，指向 heap 上的
\texttt{Clip}；\texttt{Clip} 裡的 \texttt{title} 再指向 heap 上另一段
字元。 放掉外層等於拆掉中間那座橋，內層那串字還在 heap，只是程式再也
走不到。

兩種常見的錯法：

\begin{code}{c}{內層沒配置就解參考，或只放外層}
Clip *c = malloc(sizeof(Clip));
// printf("%s\n", c->title);  // title is garbage: undefined
// free(c);                   // leaks the string if title was malloc'd
\end{code}

第一行之後 \texttt{c->title} 沒有指向合法物件，拿去印是未定義行為。
第二種是標題其實有配置，卻只 \texttt{free(c)}，字串那塊從此消失在帳
本裡。 配置與釋放成對寫在 \texttt{clip\_create}／
\texttt{clip\_destroy} 裡，比較不容易漏。

若 \texttt{Clip} 本身在 stack 上、只有 \texttt{title} 在 heap，仍然要
\texttt{free(c.title)}，只是不用再 \texttt{free} 那份結構。

\begin{drills}[label=drills:struct-pointer]
\item \texttt{ptr->x} 等同於哪一個用 \texttt{(*ptr)} 寫的寫法？
\item \texttt{*ptr.x}（沒加括號）這樣寫會出什麼問題？
\item \texttt{clip\_destroy} 為什麼要先 \texttt{free(c->title)}、再
\texttt{free(c)}？ 順序反過來會發生什麼事？
\end{drills}

\section{參數傳遞}
\label{sec:struct-params}

結構當函式引數時，預設是傳值：呼叫的當下複製一整包進去。 函式也可以
改收指標，只複製一個位址。 這兩種寫法都能用，取捨要看要不要改到呼叫
端、以及複製成本接不接受。

\subsection{傳值：複數運算}

複數有實部、虛部，加法和乘法都是輸入兩個複數、算出一個新的再送回
來，很適合看傳值與回傳同時發生。

\begin{example}{c}{(cpx-arith.c) 傳兩個複數，回傳一個新的}
#include <stdio.h>

typedef struct {
    double re, im;
} Cpx;

Cpx add(Cpx a, Cpx b) {
    Cpx r;
    r.re = a.re + b.re;
    r.im = a.im + b.im;
    return r;
}

Cpx mul(Cpx a, Cpx b) {
    Cpx r;
    r.re = a.re * b.re - a.im * b.im;
    r.im = a.re * b.im + a.im * b.re;
    return r;
}

int main(void) {
    Cpx a = {1, 2};
    Cpx b = {3, 4};
    Cpx s = add(a, b);
    Cpx p = mul(a, b);
    printf("%.1f + %.1fi\n", s.re, s.im);
    printf("%.1f + %.1fi\n", p.re, p.im);
}
\end{example}

\iosample{
}{
4.0 + 6.0i \\
-5.0 + 10.0i \\
}

呼叫 \texttt{add(a, b)} 時，\texttt{a}、\texttt{b} 各複製一份進函式，
函式裡的 \texttt{a.re = 100} 就算真的寫了，也動不到 \texttt{main} 裡
的那份。 \texttt{return} 再複製一份出去，存進 \texttt{s}。 代入
$1 + 2i$ 與 $3 + 4i$：和是 $4 + 6i$，積是 $(1 \cdot 3 - 2 \cdot 4) +
(1 \cdot 4 + 2 \cdot 3)i = -5 + 10i$，跟印出來的兩行相符。
\texttt{main} 裡的 \texttt{a}、\texttt{b} 從頭到尾都還是 $1 + 2i$、
$3 + 4i$。

兩個 \texttt{double} 的結構很小，複製的成本可以忽略。 成員一多，每一
次呼叫都整包搬運就會開始有感。

\subsection{傳位址}

傳指標有兩個常見理由。 第一是效能：複製一個位址，比起複製整個結構便
宜，結構越大差越多。 第二是可修改性：函式要改到呼叫端的那一份，只能
拿到它的位址。

下面這個對照把兩種寫法放在同一個程式裡。 \texttt{move\_copy} 收結構
本身，\texttt{move} 收指標：

\begin{example}{c}{(move-point.c) 傳值改不到正本，傳位址可以}
#include <stdio.h>

struct Point {
    int x, y;
};

void move_copy(struct Point p, int x, int y) {
    p.x = x;
    p.y = y;
}

void move(struct Point *p, int x, int y) {
    p->x = x;
    p->y = y;
}

int main(void) {
    struct Point p = {0, 0};
    move_copy(p, 5, 5);
    printf("after copy: %d %d\n", p.x, p.y);
    move(&p, 5, 5);
    printf("after pointer: %d %d\n", p.x, p.y);
}
\end{example}

\iosample{
}{
after copy: 0 0 \\
after pointer: 5 5 \\
}

\texttt{move\_copy(p, 5, 5)} 把 $\{0, 0\}$ 複製進參數 \texttt{p}，函
式裡改的是複本，所以第一行仍是 $0\ 0$。 這就是傳值版改了
\texttt{p.x}、\texttt{p.y} 卻動不到 \texttt{main} 裡那份的原因。
\texttt{move(\&p, 5, 5)} 只傳位址，\texttt{p->x} 寫的是正本，第二行才
變成 $5\ 5$。

\subsection{怎麼選}

\begin{center}
\begin{tabular}{|l|l|l|}
\hline
 & \textbf{傳值} & \textbf{傳指標} \\
\hline
函式內能否改到正本 & 不能，改的是複本 & 能 \\
是否複製整個 struct & 是，成本隨大小增加 & 只複製一個位址 \\
\hline
\end{tabular}
\end{center}

只要讀、結構又小，傳值清楚，像上面的複數。 要改呼叫端的資料，或結構
已經大到不想每次搬運，傳指標。 只讀的指標參數可以標 \texttt{const}，
向呼叫者保證函式不會透過這個指標寫入，第 13 章的 \texttt{const} 參數
就是這種用法。

\begin{drills}[label=drills:struct-params]
\item \texttt{move\_copy(p, 5, 5)} 呼叫後，\texttt{main} 裡的
\texttt{p.x} 有沒有改變？ 為什麼？
\item 想讓函式真的改到呼叫端的結構內容，參數該傳值還是傳指標？
\item 結構很大、函式只需要讀取不需要修改時，指標參數該怎麼宣告比較
好？
\end{drills}

\section{結構應用}
\label{sec:struct-nested}

成員也可以是另一個結構。 學生除了學號、姓名、GPA，還有入學日期，就把
\texttt{Date} 嵌進 \texttt{Student}。 取用時從外往內走：
\texttt{s.enrolled.year}。

\begin{example}{c}{(enrolled.c) 學生裡嵌一筆日期}
#include <stdio.h>

struct Date {
    int year, month, day;
};

struct Student {
    char id[10];
    char name[50];
    double gpa;
    struct Date enrolled;
};

int main(void) {
    struct Student s = {
        .id = "R123",
        .name = "Ada",
        .gpa = 3.8,
        .enrolled = {.year = 2026, .month = 8, .day = 31}
    };
    printf("%s enrolled %d-%02d-%02d\n",
           s.name, s.enrolled.year,
           s.enrolled.month, s.enrolled.day);
    s.enrolled.day = 1;
    printf("day now %d\n", s.enrolled.day);
}
\end{example}

\iosample{
}{
Ada enrolled 2026-08-31 \\
day now 1 \\
}

\texttt{.enrolled = \{.year = 2026, \ldots\}} 是標籤化初始化套在內層
結構上。 印出來的第一行用 \texttt{s.enrolled.year} 這串點運算子往下
走；接著只改 \texttt{s.enrolled.day}，外層的姓名與 GPA 不受影響。 開
場那筆學生資料，到這裡學號、成績、入學日都在同一包裡。

\begin{drills}[label=drills:struct-nested]
\item 要存取內層 \texttt{Date} 結構的 \texttt{year} 成員，該怎麼寫？
\item 只改 \texttt{s.enrolled.day}，會不會影響到 \texttt{s.name} 或
\texttt{s.gpa}？
\end{drills}

\section{多檔案程式設計}
\label{sec:struct-multifile}

複數那種型別若會在好幾個 \texttt{.c} 裡用到，把定義和函式原型放進標
頭，實作放另一個檔，主程式只 \texttt{\#include} 標頭：

\begin{code}{c}{標頭與實作分開}
/* cpx.h */
typedef struct {
    double re, im;
} Cpx;

Cpx add(Cpx a, Cpx b);
Cpx mul(Cpx a, Cpx b);

/* cpx.c */
#include "cpx.h"
Cpx add(Cpx a, Cpx b) {
    Cpx r;
    r.re = a.re + b.re;
    r.im = a.im + b.im;
    return r;
}
/* ... mul ... */

/* main.c */
#include "cpx.h"
/* call add and mul */
\end{code}

標頭讓每個翻譯單元看見同一份型別與原型；函式本體只出現一次，避免重複
定義。 標頭檔保護、makefile 怎麼把這幾個檔編成一支程式，留給之後的
章節。

\keypoint{型別與函式原型放標頭，本體放 \texttt{.c}，主程式只透過標頭
使用它們。}

\begin{drills}[label=drills:struct-multifile]
\item 型別定義和函式原型該放在 \texttt{.h} 還是 \texttt{.c}？ 函式本
體呢？
\item 如果好幾個 \texttt{.c} 檔都要用同一個結構型別，為什麼要共用同
一份標頭，而不是各自定義一次？
\end{drills}

\addcontentsline{toc}{section}{練習解答}
\section*{練習解答}

各節練習的答案收在這裡。 短題只給能一句話核對的結果。

\paragraph{第~\ref{sec:struct-motivation}~節}
\begin{enumerate}
\item 因為索引 \texttt{i} 必須永遠對準同一個人，編譯器不會幫忙核對；
排序或交換時只要漏動其中一個陣列，資料就會錯位。
\item 學號（\texttt{char id[10]}）、姓名（\texttt{char name[50]}）、
GPA（\texttt{double gpa}）。
\item 平行陣列版要傳三個陣列加一個索引；struct 版只要傳一個
\texttt{struct Student} 或它的位址。
\end{enumerate}

\paragraph{第~\ref{sec:struct-declare}~節}
\begin{enumerate}
\item 不能。 C 比的是型別名字，不是成員長得像不像；兩個匿名結構即使
成員完全一樣，仍是兩種不相容的型別。
\item 是同一種型別，只是書寫長短不同。
\item 是填補（padding）位元組，用來讓 \texttt{int uses} 從位移 $12$
（$4$ 的倍數）開始，滿足對齊要求。
\item 因為 \texttt{char *name} 對齊到 $8$，把 \texttt{short} 放最前
面會讓後面的指標成員錯位，編譯器得在 \texttt{short} 後面多插入 $6$
位元組填補才能讓指標對齊，整體反而比把指標放最前面還大。
\end{enumerate}

\paragraph{第~\ref{sec:struct-init}~節}
\begin{enumerate}
\item $0$。 部分初始化時沒寫到的成員一律補 $0$。
\item 空字串。 \texttt{name} 沒被提到，整個字元陣列每個位元組都補
$0$，第一個字元就是 \texttt{'\textbackslash 0'}。
\item 依名字對應，不必背宣告順序、也不怕以後成員順序調整時初值全部
錯位。
\end{enumerate}

\paragraph{第~\ref{sec:struct-access}~節}
\begin{enumerate}
\item 因為 \texttt{(\&p)} 的型別是 \texttt{struct Point *}，是指標，
指標沒有 \texttt{.} 可用；要寫成 \texttt{(\&p)->y}，或乾脆用
\texttt{p.y}。
\item 不能，C 沒有為結構定義 \texttt{==}；要自己逐一比較每個成員
（例如 \texttt{a.x == b.x \&\& a.y == b.y}）。
\item 整包交換一次把學號、姓名、GPA 一起搬動，三個欄位永遠對齊同一
個人；只換 \texttt{gpa} 會讓成績跟錯的姓名配對在一起。
\end{enumerate}

\paragraph{第~\ref{sec:struct-pointer}~節}
\begin{enumerate}
\item \texttt{(*ptr).x}。
\item \texttt{.} 的優先權比 \texttt{*} 高，會先算成 \texttt{ptr.x}
再解參考；但 \texttt{ptr} 是指標，沒有成員 \texttt{x}，編譯不過。
\item 因為 \texttt{c} 一旦被 \texttt{free}，就再也拿不到
\texttt{c->title} 那個位址；順序反過來會漏放內層字串，造成記憶體洩
漏。
\end{enumerate}

\paragraph{第~\ref{sec:struct-params}~節}
\begin{enumerate}
\item 沒有改變，仍是 \texttt{0 0}。 \texttt{move\_copy} 收到的是
\texttt{p} 的複本，函式裡改的是複本，動不到 \texttt{main} 裡的正
本。
\item 要傳指標，讓函式透過位址寫回正本。
\item 宣告成 \texttt{const struct T *}，只複製一個位址、成本低，又
向呼叫者保證不會透過這個指標寫入。
\end{enumerate}

\paragraph{第~\ref{sec:struct-nested}~節}
\begin{enumerate}
\item \texttt{s.enrolled.year}，從外層結構的成員一路往內層走。
\item 不會，只改 \texttt{s.enrolled.day} 這一個欄位，其餘欄位不受
影響。
\end{enumerate}

\paragraph{第~\ref{sec:struct-multifile}~節}
\begin{enumerate}
\item 型別定義和函式原型放 \texttt{.h}；函式本體放 \texttt{.c}。
\item 這樣每個用到這個型別的翻譯單元看到的都是同一份定義，不必各自
重複定義，也避免不同檔案裡型別定義不一致。
\end{enumerate}

\chapter{遞迴}
\label{ch:recursion}

遞迴（recursion）是一種把問題交給同一個函式、但規模更小的輸入去解的寫
法。 一個問題能不能這樣拆，取決於它是否長得像自己的縮小版：空的情況可
以立刻回答，非空的情況則可以假設縮小版已經解完，再把那個答案組回現在
這層要的結果。

第~\ref{ch:function}~章的範例~\ref{ex:print-digits}（\texttt{printDigits}）
其實已經用過遞迴：先把數字除以 $10$ 交給下一層處理，這一層再印出個位。
當時我們把它當成 \texttt{void} 函式的例子。 這一章要把這種寫法變成一套
可以反覆套用的設計步驟，並一路看到它在整數、陣列、字串、鏈結結構上長
什麼樣子，以及在 C 裡會踩到哪些只屬於這門語言的陷阱。

這一章分四段推進。 先用自我相似把副本模型種下來，再用三個問題寫出第一
個遞迴，接著才畫呼叫堆疊、把回程上那些待完成的運算講清楚。 然後按資料
的形狀走一遍：整數、陣列區間、字串、鏈結與樹。 再處理多重遞迴、記憶化、
尾遞迴，以及把遞迴轉成迴圈。 最後才進入回溯法。 排列與 N 皇后會等到那
一節才出現，因為它們在遞迴之上又多了一層選擇與撤銷。

節末練習的答案收在章末。

\section{自我相似：不寫程式碼的暖身}
\label{sec:rec-self-sim}

資料夾裡面還可以再放資料夾。 打開最外層會看到下一層，一直打開到空資料
夾為止。 若要數整個目錄樹裡有幾個檔案，做法很固定：空資料夾貢獻 $0$；
非空資料夾的檔案數，等於自己這一層放的檔案，加上裡面每一個子資料夾的
檔案數。 子資料夾的問題跟原本的問題長得一模一樣，只是規模更小。 這種
自己裡面還有一個自己的結構，叫做\textbf{自我相似}。

走廊上站了一排同學，也可以看成同一件事。 最前面的人被問到 $1$ 加到
$5$ 是多少。 他自己只記得一件事：如果後面那個人能告訴他 $1$ 加到 $4$
的和，他再加 $5$ 就得到答案。 於是他把問題傳給後面的人，後面的人又如
法炮製，一直傳到最後一個被問 $1$ 加到 $0$ 的人，那個人立刻回答 $0$。
答案再一層一層傳回來，每一個人都把自己負責的那個數字加上去。

\begin{center}
\begin{tikzpicture}
    \foreach \i/\name/\q in {0/甲/\texttt{sum(3)}, 1/乙/\texttt{sum(2)}, 2/丙/\texttt{sum(1)}, 3/丁/\texttt{sum(0)}} {
        \node[cell, minimum width=1.7cm, minimum height=1.05cm] (p\i) at (2.35*\i, 0) {\name};
        \node[idxlbl, above=3pt of p\i] {\q};
    }
    \foreach \i in {0,1,2} {
        \pgfmathtruncatemacro{\j}{\i+1}
        \draw[diagarrow] ([yshift=4pt]p\i.east) -- ([yshift=4pt]p\j.west);
        \draw[diagarrow] ([yshift=-4pt]p\j.west) -- ([yshift=-4pt]p\i.east);
    }
    \node[idxlbl, above=2pt] at ($(p0.east)!0.5!(p1.west)+(0,0.22)$) {問};
    \node[idxlbl, below=2pt] at ($(p0.east)!0.5!(p1.west)+(0,-0.22)$) {答};
\end{tikzpicture}
\diagcaption{一排同學同時站在走廊上，每人手上拿著自己那一層的問題。
問題往後傳、答案往前傳；這四個人在同一時間都還在，各自獨立。}
\label{fig:rec-line-of-people}
\end{center}

這排同學在同一時間都還站在走廊上，每人手上拿著自己的問題，等後面那個
人回覆。 這就是這一章要先種下的圖像：同一份函式定義，同時存在好多份
各自獨立的執行。 讀過迴圈的人，很容易把遞迴想成同一個函式在原地把
$n$ 改小再跑一次；走廊上的同學告訴我們，每一層都是一個還活著的人，
有自己的問題、自己的等待。 函式那一章已經把這件事叫做 stack frame：
每一次呼叫都有自己的參數複本。 遞迴只是同一份函式被呼叫了很多次，所
以同一時間會有很多層長得很像、但參數不同的 frame 疊在堆疊上。

\keypoint{遞迴的每一層都是一次獨立的函式呼叫，有自己的參數與區域變數。
同一時間可以有許多層同時活著，各自等自己的子呼叫回來。}

\begin{drills}[label=drills:rec-self-sim]
\item 圖~\ref{fig:rec-line-of-people} 在問題傳到 \texttt{sum(0)} 的那一
刻，走廊上同時站著幾個人？ 每人手上的問題分別是什麼？
\item 若把遞迴想成同一個人把 $n$ 改小再算一次，這張圖裡哪一件事會對不
上？
\item 用三句話寫出數資料夾檔案的基礎情況、縮小方式、組合方式。
\end{drills}

\section{寫遞迴的三個問題}
\label{sec:rec-three-q}

要把一個問題寫成遞迴函式，我們每次都問同樣的三件事：

\begin{enumerate}
\item 最小的情況是什麼？ 它的答案是什麼？
\item 假設這個函式對更小的輸入已經是對的，我要怎麼用它的回傳值，組出
現在這層的答案？
\item 每次呼叫是否嚴格逼近最小情況？
\end{enumerate}

第一問決定\textbf{基礎情況}（base case）：小到可以立刻回答、不必再拆。
第二問決定\textbf{遞迴情況}：把問題縮小，再把子問題的答案組合回來。
第三問保證程式會停：如果每次呼叫都沒有比較小，基礎情況永遠走不到，堆
疊就會一直長，直到程式被作業系統殺掉。

第二問有一個寫的時候必須遵守的紀律。 我們在寫 $n$ 這一層的時候，不准
跟著子呼叫一路追下去看它內部怎麼算；我們直接相信 \texttt{sum(n - 1)}
已經回傳正確的 $1 + \cdots + (n-1)$，只負責把 $n$ 加上去。 這種把遞迴
呼叫當成已經算對的黑盒子的寫法，叫做\textbf{遞迴信仰之躍}
（recursive leap of faith）。 設計遞迴時靠的就是它；一層一層把堆疊畫
出來，是後面用來核對與除錯的工具，兩件事分開做。

\keypoint{寫遞迴時，把遞迴呼叫當成已經算對的黑盒子：只負責處理最小情
況，以及如何把子問題的回傳值組合成現在這層的答案。 這叫做遞迴信仰之躍。}

片語~\ref{idiom:three-q} 把三個問題寫成固定的程式骨架。 註解裡的 (1)(2)(3)
對應上面那三問，之後每一個例子都會再標一次。

\begin{code}[label=idiom:three-q]{text}{寫遞迴的三個問題}
回傳型別 f(較小的輸入) {
    if (已經是最小情況) {          // (1) 基礎情況
        return 最小情況的答案;
    }
    子答案 = f(更小的輸入);         // (2) 信仰之躍：假設這一步是對的
    return 用子答案組出這一層的答案; // (3) 組合；輸入必須比現在小
}
\end{code}

第一個完整例子用 $1$ 加到 $n$。 三個問題逐一填：

\begin{itemize}
\item 最小情況：$n = 0$，空的加總是 $0$。
\item 信仰之躍：假設 \texttt{sum(n - 1)} 已經是 $1 + \cdots + (n-1)$。
\item 組合：現在的答案是 $n + \texttt{sum(n - 1)}$。
\item 進度：每次把 $n$ 減 $1$，對非負整數一定會走到 $0$。
\end{itemize}

\begin{example}[label=ex:sum]{c}{(sum.c) 用遞迴求 1 加到 n}
#include <stdio.h>

int sum(int n) {
    if (n == 0) {              // (1) 最小情況：空的和是 0
        return 0;
    }
    return n + sum(n - 1);     // (2) 縮小成 n-1  (3) 用加法組合
}

int main(void) {
    int n;
    scanf("%d", &n);
    printf("%d\n", sum(n));
}
\end{example}

\iosample{
5 \\
}{
15 \\
}

用信仰之躍來看 \texttt{sum(5)}，我們只寫一步：
\texttt{sum(5) = 5 + sum(4)}。 相信 \texttt{sum(4) = 10}，所以答案是
$15$。 不必在這一層就把 \texttt{sum(4)}、\texttt{sum(3)}、
\texttt{sum(2)} 全部展開。 正是因為每一層都只做一件事，整條鏈才組得出
$5+4+3+2+1+0 = 15$。

基礎情況也可以寫成 $n = 1$ 時回傳 $1$，對正整數 $n$ 答案一樣。 差別在
$n = 0$：以 $0$ 為基礎情況的版本仍然正確，以 $1$ 為基礎的版本會再呼叫
\texttt{sum(0)}，然後 \texttt{sum(-1)}、\texttt{sum(-2)}，永遠走不到停
下的地方。 基礎情況選哪一個，取決於函式的定義域包含哪些輸入。

\stylepoint{先把三個問題用一句話寫在註解或草稿上，再填程式；不要一邊
寫 \texttt{return} 一邊想像整條呼叫鏈，那會把信仰之躍換成窮舉追蹤。}

\begin{drills}[label=drills:rec-three-q]
\item 只寫一步信仰之躍：\texttt{sum(4)} 等於什麼？ 不要把 \texttt{sum(3)}
再展開。
\item 把範例~\ref{ex:sum} 的基礎情況改成 \texttt{if (n == 1) return 1;}，
再呼叫 \texttt{sum(0)}。 程式會停嗎？ 停的話答案是多少？
\item 補基礎情況：求 $1^2 + 2^2 + \cdots + n^2$ 的遞迴，最小情況與它的
答案各是什麼？
\item 從零寫：用三個問題填表，寫出求 $1+3+5+\cdots+(2n-1)$ 的函式。
\end{drills}

\section{整數上的遞迴與溢位}
\label{sec:rec-int}

$n!$ 的定義本身就是遞迴的：$0! = 1$，且 $n! = n \times (n-1)!$。 三個
問題幾乎直接從定義抄下來。 片語~\ref{idiom:factorial} 是對應的骨架。

\begin{code}[label=idiom:factorial]{c}{階乘的遞迴骨架}
long long factorial(int n) {
    if (n <= 1) {                 // (1) 0! = 1! = 1
        return 1;
    }
    return n * factorial(n - 1);  // (2)(3) 乘上縮小後的階乘
}
\end{code}

這裡把回傳型別寫成 \texttt{long long}，是因為 $n!$ 長大的速度非常快。
32-bit 有號整數的上限 \texttt{INT\_MAX} 是 $2147483647$。
\[
12! = 479001600 < 2147483647, \qquad
13! = 6227020800 > 2147483647.
\]
若把回傳型別寫成 \texttt{int}，\texttt{factorial(13)} 在 C 裡是
\textbf{有號整數溢位}，屬於未定義行為：編譯器不保證印出哪個數字，也不
保證程式還能繼續正常跑。 64-bit 的 \texttt{long long} 可以撐到 $20!$
（$2432902008176640000$）；$21!$ 連 \texttt{long long} 都會溢位。

\begin{example}{c}{(factorial.c) 用 long long 計算階乘}
#include <stdio.h>

long long factorial(int n) {
    if (n <= 1) {
        return 1;
    }
    return n * factorial(n - 1);
}

int main(void) {
    int n;
    scanf("%d", &n);
    printf("%lld\n", factorial(n));
}
\end{example}

\iosample{
10 \\
}{
3628800 \\
}

輸入 $10$，程式先一路呼叫到 \texttt{factorial(1)} 回傳 $1$，再在回程
上依序乘 $2, 3, \ldots, 10$，得到 $3628800$。 乘法發生在遞迴呼叫的
\texttt{return} 裡、子呼叫回來之後；這一點下一節會畫出來。 若輸入改成
$20$，同一支程式印出 $2432902008176640000$；若仍用 \texttt{int}，連
$13$ 都已經超出範圍。

同樣的骨架可以套到乘冪：$a^0 = 1$，$a^n = a \times a^{n-1}$。

\begin{code}{c}{乘冪的遞迴骨架}
long long power(long long a, int n) {
    if (n == 0) {
        return 1;
    }
    return a * power(a, n - 1);
}
\end{code}

這個骨架只處理 $n \ge 0$。 若 $n$ 是負數，$n-1$ 會愈來愈小，永遠走不
到 $n == 0$，程式會一直呼叫直到堆疊溢位。 負指數要另外定義，例如改成
回傳分數，或一開始就拒絕 $n < 0$。

\keypoint{遞迴算階乘、費氏數列這類快速長大的整數時，先問結果會不會超
出 \texttt{int}。 $13!$ 與後面會遇到的 $\mathrm{fib}(47)$ 都已經超過
32-bit 有號整數；需要時改用 \texttt{long long}，並記得它也有上限。}

\begin{drills}[label=drills:rec-int]
\item 若 \texttt{factorial} 的回傳型別是 \texttt{int}，呼叫
\texttt{factorial(13)} 在 C 裡是什麼行為？ 能靠印出來的數字判斷它包了
幾圈嗎？
\item \texttt{power(2, -3)} 用本節的骨架會停嗎？ 為什麼？
\item $20!$ 用 \texttt{long long} 印得出來；$21!$ 呢？
\end{drills}

\section{遞迴呼叫時發生什麼事：堆疊與回程}
\label{sec:rec-stack}

信仰之躍讓我們能寫出 \texttt{factorial}。 若要核對它真的做對、或找出
為什麼會停不下來，就需要看執行期在記憶體裡發生什麼事。 第~\ref{ch:function}~章
已經介紹過呼叫堆疊：每次呼叫在堆疊頂端壓入一層 frame，參數與區域變數都
住在那一層裡，\texttt{return} 時整層彈掉。 遞迴只是同一份函式被壓了
很多次，所以同一時間會有好多層 \texttt{factorial}，各自帶著自己的
\texttt{n}。

關鍵在\textbf{回程}。 \texttt{return n * factorial(n - 1);} 這一行裡
的乘法，並沒有在往下呼叫的路上做完。 它被掛在這一層的 frame 裡，等子
呼叫把 $(n-1)!$ 帶回來才乘。 往基礎情況走的那一段叫做去程（active
flow），從基礎情況帶著回傳值往回走、把掛著的運算做完的那一段叫做回程
（passive flow / backward flow）。 學生最常漏掉的就是回程：以為工作
全部發生在往下走的路上，子呼叫一回來答案就已經好了。

圖~\ref{fig:rec-stack-down} 是 \texttt{factorial(3)} 走到基礎情況時
的堆疊。 因為基礎情況是 $n \le 1$，最深一層是 $n = 1$，正在執行、準備
回傳 $1$；下面兩層都暫停著，各自掛著一個還沒做的乘法。 這一層\emph{不
會}再呼叫 \texttt{factorial(0)}。

\begin{center}
\begin{tikzpicture}
    \node[hlcell, minimum width=5.6cm, minimum height=0.72cm] (f1)
        {\texttt{factorial} \quad $n=1$ \quad 回傳 $1$};
    \node[cell, minimum width=5.6cm, minimum height=0.72cm, below=0pt of f1] (f2)
        {\texttt{factorial} \quad $n=2$ \quad 待做：$2 \times \cdots$};
    \node[cell, minimum width=5.6cm, minimum height=0.72cm, below=0pt of f2] (f3)
        {\texttt{factorial} \quad $n=3$ \quad 待做：$3 \times \cdots$};
    \node[cell, minimum width=5.6cm, minimum height=0.72cm, below=0.55cm of f3] (mn)
        {\texttt{main} \quad 等 \texttt{factorial(3)} 回來};

    \node[idxlbl, anchor=east, xshift=-0.12cm] at (f1.west) {頂};
    \node[idxlbl, anchor=west, align=left, xshift=0.2cm] at (f1.east)
        {目前執行\\這一層};
    \node[idxlbl, anchor=west, align=left, xshift=0.2cm] at (f2.east)
        {暫停};
    \node[idxlbl, anchor=west, align=left, xshift=0.2cm] at (mn.east)
        {暫停};
\end{tikzpicture}
\diagcaption{\texttt{factorial(3)} 的去程結束時：三層 \texttt{factorial}
疊在 \texttt{main} 上面，最深停在 $n = 1$，下面兩層各自掛著一個還沒做
的乘法。}
\label{fig:rec-stack-down}
\end{center}

回程把掛著的乘法由上往下做完，如圖~\ref{fig:rec-stack-up}。
\texttt{n = 1} 回傳 $1$ 之後彈掉；\texttt{n = 2} 計算 $2 \times 1 = 2$；
最後 \texttt{n = 3} 計算 $3 \times 2 = 6$，交回 \texttt{main}。

\begin{center}
\begin{tikzpicture}
    \node[process, minimum width=3.4cm] (a) {$n=1$ 回傳 $1$};
    \node[process, minimum width=3.4cm, below=0.45cm of a] (c) {$2 \times 1 = 2$};
    \node[process, minimum width=3.4cm, below=0.45cm of c] (d) {$3 \times 2 = 6$};
    \node[startend, minimum width=3.4cm, below=0.45cm of d] (e) {\texttt{main} 印出 $6$};
    \draw[diagarrow] (a) -- (c);
    \draw[diagarrow] (c) -- (d);
    \draw[diagarrow] (d) -- (e);
    \node[idxlbl, right=8pt of c] {回程，由上往下};
\end{tikzpicture}
\diagcaption{\texttt{factorial(3)} 的回程：每一層把子呼叫的回傳值乘上
自己的 \texttt{n}，乘法是在回來的路上才發生的。}
\label{fig:rec-stack-up}
\end{center}

把 \texttt{sum(3)} 用同一套圖像走一遍，掛著的改成加法，而且基礎情況是
$n = 0$，所以會多一層：
$0$ 回傳 $0$，然後 $1+0=1$、$2+1=3$、$3+3=6$。 去程只是把問題越問越小，
真正把數字組起來的工作發生在回程。 第~\ref{ch:function}~章的
\texttt{printDigits}（範例~\ref{ex:print-digits}）
把 \texttt{printf} 放在遞迴呼叫後面，印出的順序是高位到低位，原因也是
回程：除到只剩最高位時先印，再一路回來印較低的位。

會把這張圖畫出來，跟會寫出遞迴，是兩件不同的技能。 設計時靠三個問題
與信仰之躍；把 frame 一層一層展開，是用來核對已經寫好的程式、找出基
礎情況沒寫到或組合步驟寫錯的工具。

\keypoint{\texttt{n * factorial(n - 1)} 裡的乘法掛在這一層的 frame 裡，
等子呼叫回來才做。 漏掉回程，就會以為遞迴只是把 $n$ 一路減到基礎情況
的另一種迴圈。}

\begin{drills}[label=drills:rec-stack]
\item 仿圖~\ref{fig:rec-stack-down} 畫 \texttt{factorial(4)} 去程結束時
的堆疊：最深一層的 $n$ 是多少？ 總共幾層 \texttt{factorial}？
\item 第~\ref{ch:function}~章的範例~\ref{ex:print-digits} 把
\texttt{printf} 放在遞迴呼叫後面。 若改放到呼叫前面，輸入 $123$ 會印出
什麼順序？
\item 乘法 \texttt{n * factorial(n - 1)} 發生在去程還是回程？ 用
\texttt{factorial(3)} 的一層說明。
\end{drills}

\section{遞迴過程 vs 迭代過程}
\label{sec:rec-process}

\texttt{sum} 與 \texttt{factorial} 每次呼叫都留下一個待完成的加法或乘
法，所以 $n$ 層呼叫就掛了 $n$ 個運算，概念上的空間是 $\Theta(n)$。 這
種執行過程叫做\textbf{線性遞迴過程}（linear recursive process）：程式
寫成遞迴函式，跑起來也真的在堆疊上留下一串待辦。

同一件工作可以改成每一層都不欠帳。 做法是多傳一個參數，專門用來累積
目前已經算好的結果，叫做\textbf{累加器}（accumulator）。 下面的片語
把兩種過程並排放：左邊每一層都欠一個加法，右邊把加法提前做完再往下走。

\begin{code}{c}{遞迴過程與帶累加器的迭代過程}
int sum(int n) {                 // 線性遞迴過程
    if (n == 0) {
        return 0;
    }
    return n + sum(n - 1);       // 加法掛著等回來
}

int sum_iter(int n, int acc) {   // 線性迭代過程
    if (n == 0) {
        return acc;
    }
    return sum_iter(n - 1, acc + n);  // 這一層什麼都不欠
}
\end{code}

\texttt{sum\_iter(5, 0)} 的每一層都把加法做完才呼叫下一層：
\texttt{sum\_iter(5, 0)} $\rightarrow$ \texttt{sum\_iter(4, 5)}
$\rightarrow$ \texttt{sum\_iter(3, 9)} $\rightarrow$
\texttt{sum\_iter(2, 12)} $\rightarrow$ \texttt{sum\_iter(1, 14)}
$\rightarrow$ \texttt{sum\_iter(0, 15)}，最後回傳 $15$。 堆疊上雖然仍
有一層 frame，但概念上這一層已經不欠任何運算，這種執行過程叫做
\textbf{線性迭代過程}（linear iterative process）。 空間在概念上是
$\Theta(1)$；至於 C 的編譯器會不會真的把這層 frame 重用，要等到尾遞迴
那一節才講，而且 C 標準並不保證。

\begin{example}{c}{(sum-acc.c) 兩種過程算同一個和}
#include <stdio.h>

int sum(int n) {
    if (n == 0) {
        return 0;
    }
    return n + sum(n - 1);
}

int sum_iter(int n, int acc) {
    if (n == 0) {
        return acc;
    }
    return sum_iter(n - 1, acc + n);
}

int main(void) {
    int n;
    scanf("%d", &n);
    printf("%d %d\n", sum(n), sum_iter(n, 0));
}
\end{example}

\iosample{
5 \\
}{
15 15 \\
}

兩個函式對同一個 $n$ 印出相同的 $15$，差別在過程。 之後若看到一個遞迴
在 \texttt{return} 裡還有加法、乘法、或其它必須等子呼叫回來才能做的事，
它跑的就是遞迴過程，深度為 $n$ 時堆疊也長到 $\Theta(n)$。 若遞迴呼叫本
身就是 \texttt{return} 的全部內容，它跑的就是迭代過程，這正是稍後尾遞
迴要利用的形狀。

\keypoint{函式寫成遞迴，跑起來不一定是遞迴過程。 差別在回程上有沒有掛
著還沒做的運算：有，就是線性遞迴過程；沒有、只靠累加器往下傳，就是線
性迭代過程。}

\begin{drills}[label=drills:rec-process]
\item \texttt{sum} 與 \texttt{sum\_iter} 哪一個在回程上欠著加法？ 概念
空間各是什麼？
\item 寫出 \texttt{sum\_iter(3, 0)} 每一層的兩個參數。
\item 把 \texttt{factorial} 改成帶累加器的版本（只寫函式，呼叫端傳
\texttt{acc = 1}）。
\end{drills}

\section{在陣列上遞迴：區間表示法}
\label{sec:rec-array}

C 沒有辦法把陣列當成一個可以拆開的值來傳。 傳進函式的是起始位址，長度
必須另外帶。 在陣列上寫遞迴，慣用的做法是把目前要處理的範圍表示成一個
\textbf{半開區間} $[lo, hi)$：從索引 \texttt{lo} 開始、到 \texttt{hi}
之前結束，長度是 \texttt{hi - lo}。 空區間就是 \texttt{lo >= hi}；整個
陣列就是 \texttt{lo = 0}、\texttt{hi = n}。

縮小問題的方式很固定：處理完 \texttt{a[lo]} 之後，剩下的是
$[lo+1, hi)$。 三個問題變成：空區間的答案是什麼、假設子區間已經算對、
怎麼把 \texttt{a[lo]} 組進去。

\begin{code}{text}{陣列區間上的遞迴}
回傳型別 f(int *a, int lo, int hi) {
    if (lo >= hi) {                    // (1) 空區間
        return 空區間的答案;
    }
    return 組合(a[lo], f(a, lo + 1, hi)); // (2)(3) 縮小成右半段
}
\end{code}

圖~\ref{fig:rec-array-range} 是 \texttt{array\_sum(a, 0, 5)} 第一層看
到的範圍：橘色那格是這一層要加進去的 \texttt{a[0]}，右邊 $[1, 5)$ 交給
子呼叫。

\begin{center}
\begin{tikzpicture}
    \node[hlcell] (c0) {1};
    \node[cell, right=0pt of c0] (c1) {4};
    \node[cell, right=0pt of c1] (c2) {8};
    \node[cell, right=0pt of c2] (c3) {9};
    \node[cell, right=0pt of c3] (c4) {10};
    \node[idxlbl, below=2pt of c0] {0};
    \node[idxlbl, below=2pt of c1] {1};
    \node[idxlbl, below=2pt of c2] {2};
    \node[idxlbl, below=2pt of c3] {3};
    \node[idxlbl, below=2pt of c4] {4};
    \node[idxlbl, above=3pt of c0] {\texttt{lo}};
    \node[idxlbl, above=3pt] at ($(c4.east)+(0.15,0)$) {\texttt{hi}};
    \draw[diagarrow] ($(c4.north east)+(0.15,0.35)$) -- ($(c4.north east)+(0.15,0.08)$);
\end{tikzpicture}
\diagcaption{半開區間 $[0, 5)$。 橘色是這一層的 \texttt{a[lo]}，右邊
$[1, 5)$ 是子問題；\texttt{hi} 指向最後一格的右邊，本身不是一個元素。}
\label{fig:rec-array-range}
\end{center}

\begin{example}{c}{(array-sum.c) 用區間遞迴求陣列元素和}
#include <stdio.h>

int array_sum(int *a, int lo, int hi) {
    if (lo >= hi) {                    // (1) 空區間的和是 0
        return 0;
    }
    return a[lo] + array_sum(a, lo + 1, hi);  // (2)(3)
}

int main(void) {
    int n, a[100];
    scanf("%d", &n);
    if (n < 0 || n > 100) {
        return 0;
    }
    for (int i = 0; i < n; i++) {
        scanf("%d", &a[i]);
    }
    printf("%d\n", array_sum(a, 0, n));
}
\end{example}

\iosample{
5 \\
1 4 8 9 10 \\
}{
32 \\
}

\texttt{array\_sum(a, 0, 5)} 依信仰之躍就是 $1 + \texttt{array\_sum(a, 1, 5)}$。
子區間 $[1, 5)$ 的和是 $4+8+9+10 = 31$，加上 $1$ 得到 $32$。 空區間回傳
$0$，與 \texttt{sum(0) = 0} 是同一件事，只是現在空的是一段索引範圍。

求最大值時，空區間沒有最大值，基礎情況改成\textbf{單元素}區間：
\texttt{lo + 1 == hi} 時答案就是 \texttt{a[lo]}。 這是基礎情況隨問題改
寫的典型例子，不要死記 $0$。 這個寫法有一個硬前提：呼叫時必須
\texttt{lo < hi}。 若傳進空區間（\texttt{lo == hi}），\texttt{lo + 1 == hi}
永遠不成立，$\texttt{lo}$ 會一路加過去、跳過 \texttt{hi}，變成無窮遞迴。
這正是後面陷阱那一節說的：基礎情況 off-by-one，重則堆疊溢位。

\begin{code}{c}{區間上的最大值，基礎情況是單元素}
int array_max(int *a, int lo, int hi) {
    /* 前提：lo < hi（非空區間）。空區間會無窮遞迴。 */
    if (lo + 1 == hi) {                // (1) 只有一格
        return a[lo];
    }
    int rest = array_max(a, lo + 1, hi); // (2)
    if (a[lo] > rest) {                // (3) 這一格跟其餘比
        return a[lo];
    }
    return rest;
}
\end{code}

陣列傳進函式時會退化成指標，\texttt{sizeof(a)} 在被呼叫的函式裡量到的
是指標的大小，不是陣列長度。 所以 \texttt{lo}、\texttt{hi} 或 $n$ 一定
要當參數傳，不能在 \texttt{array\_sum} 裡面用 \texttt{sizeof} 回頭推長
度。

區間也可以從兩頭往中間收。 判斷回文時，先比 \texttt{a[lo]} 與
\texttt{a[hi]}（這裡改用封閉區間比較好寫），相同就把範圍縮小成
$[lo+1, hi-1]$。 二分搜尋則每次把區間對半切，時間 $\Theta(\log n)$，堆
疊深度也是 $\Theta(\log n)$。

\begin{example}{c}{(binsearch.c) 在已排序陣列上遞迴二分搜尋}
#include <stdio.h>

int binsearch(int *a, int lo, int hi, int x) {
    if (lo >= hi) {                    // (1) 空區間：找不到
        return -1;
    }
    int mid = lo + (hi - lo) / 2;
    if (a[mid] == x) {
        return mid;
    }
    if (a[mid] > x) {
        return binsearch(a, lo, mid, x);      // 左半 [lo, mid)
    }
    return binsearch(a, mid + 1, hi, x);      // 右半 [mid+1, hi)
}

int main(void) {
    int n, x, a[100];
    scanf("%d %d", &n, &x);
    if (n < 0 || n > 100) {
        return 0;
    }
    for (int i = 0; i < n; i++) {
        scanf("%d", &a[i]);
    }
    printf("%d\n", binsearch(a, 0, n, x));
}
\end{example}

\iosample{
5 8 \\
1 4 8 9 10 \\
}{
2 \\
}

$n = 5$、$x = 8$，區間是 $[0, 5)$。 $mid = 2$，$\texttt{a[2]} = 8$，第
一次就找到，回傳索引 $2$。 若 $x = 7$：第一次 $mid = 2$、$8 > 7$，改搜
$[0, 2)$；第二次 $mid = 1$、$4 < 7$，改搜 $[2, 2)$，空區間回傳 $-1$。
\texttt{lo + (hi - lo) / 2} 寫成這樣是為了避免 \texttt{(lo + hi) / 2}
在 \texttt{lo}、\texttt{hi} 都很大時溢位；觀念上就是取中點。

\keypoint{C 在陣列上的遞迴，靠的是把範圍編碼成參數：$(a, lo, hi)$ 或
$(a, n)$。 陣列退化成指標之後，長度不會自己跟著走，一定要另外傳。}

\begin{drills}[label=drills:rec-array]
\item 抓蟲：對 \texttt{array\_max} 呼叫 \texttt{array\_max(a, 0, 0)}。
基礎情況會被碰到嗎？ 接下來會發生什麼事？
\item 補組合步驟：線性搜尋 $x$ 是否在 $[lo, hi)$ 裡，基礎情況是空區間
回傳 $0$。 \texttt{a[lo] == x} 時回傳 $1$，否則？
\item 已排序陣列 $\{1,4,8,9,10\}$，\texttt{binsearch} 找 $7$，依序經過
哪些 $[lo, hi)$？
\item 在 \texttt{array\_sum} 裡寫 \texttt{int n = sizeof(a) / sizeof(a[0]);}
得到的 $n$ 是陣列長度嗎？
\end{drills}

\section{在字串上遞迴}
\label{sec:rec-string}

C 的字串是以 \texttt{'\textbackslash 0'} 結尾的 \texttt{char} 陣列，所
以上一節的區間寫法完全適用。 另外還有一種更短的寫法：把指標往前移一格，
遞迴處理 \texttt{s + 1}，基礎情況就是看到結尾字元。 這等價於把區間的
\texttt{lo} 加一，只是現在移動的是指標本身。

\begin{code}{text}{以 s+1 推進的字串遞迴}
回傳型別 f(char *s) {
    if (*s == '\0') {          // (1) 空字串
        return 空字串的答案;
    }
    return 組合(*s, f(s + 1)); // (2)(3) 剩下的從下一格開始
}
\end{code}

\begin{example}{c}{(rec-strlen.c) 遞迴計算字串長度}
#include <stdio.h>

int rec_strlen(char *s) {
    if (*s == '\0') {          // (1)
        return 0;
    }
    return 1 + rec_strlen(s + 1);  // (2)(3)
}

int main(void) {
    char s[100];
    scanf("%99s", s);
    printf("%d\n", rec_strlen(s));
}
\end{example}

\iosample{
hello \\
}{
5 \\
}

對字串 hello 呼叫 \texttt{rec\_strlen}，依信仰之躍就是 $1$ 加上對
ello 的呼叫。
子字串長度 $4$，加上這一格的 $1$ 得到 $5$。 空字串的長度是 $0$，與空區
間的和是 $0$、空串列的長度是 $0$ 同一件事：最小情況貢獻單位元素。

判斷回文則適合用索引版，因為要同時看頭尾。 下面的範例讀入一個不含空白
的字串，是回文印 $1$，否則印 $0$。

\begin{example}{c}{(rec-palindrome.c) 遞迴判斷字串是否為回文}
#include <stdio.h>
#include <string.h>

int is_palindrome(char *s, int lo, int hi) {
    if (lo >= hi) {                // (1) 空或單字元都是回文
        return 1;
    }
    if (s[lo] != s[hi]) {
        return 0;
    }
    return is_palindrome(s, lo + 1, hi - 1);  // (2)(3) 往中間收
}

int main(void) {
    char s[100];
    scanf("%99s", s);
    int n = strlen(s);
    printf("%d\n", is_palindrome(s, 0, n - 1));
}
\end{example}

\iosample{
level \\
}{
1 \\
}

\texttt{level} 的頭尾都是 \texttt{l}，縮小成 \texttt{eve}；再比
\texttt{e} 與 \texttt{e}，縮小成 \texttt{v}；此時 \texttt{lo == hi}，
單字元是回文，回傳 $1$。 若輸入是 \texttt{hello}，第一層 \texttt{h} 與
\texttt{o} 就不同，立刻回傳 $0$，不必再往下走。

\texttt{s + 1} 這種推進方式以後在鏈結串列上會再見到一次：那裡移動的是
\texttt{p->next}。 兩者都是在資料自己的形狀上往前走一格，程式的形狀跟
資料的形狀對齊。

\begin{drills}[label=drills:rec-string]
\item 預測：\texttt{rec\_strlen} 接到空字串（第一個字元就是
\texttt{'\textbackslash 0'}）回傳多少？
\item Parsons：把 \texttt{is\_palindrome} 的四件事排成正確順序——比頭尾、
縮小成 $[lo+1,hi-1]$、\texttt{lo >= hi} 回傳 $1$、頭尾不同回傳 $0$。
\item 把 \texttt{rec\_strlen} 改成尾遞迴：加一個累加器參數，回程上不再
欠 $+1$。
\end{drills}

\section{多重遞迴與遞迴樹：費氏數列與河內塔}
\label{sec:rec-fib-hanoi}

到目前為止，每個函式在遞迴情況裡只呼叫自己一次。 若一次拆成兩個子問
題，執行過程就從一條鏈變成一棵樹，叫做\textbf{樹狀遞迴}
（tree recursion）或多重遞迴。 費氏數列是最常見的例子：
\[
F(0) = 0,\quad F(1) = 1,\quad F(n) = F(n-1) + F(n-2).
\]
兩個子問題都比 $n$ 小，基礎情況有兩個。 寫出來幾乎就是定義本身。

\begin{code}{c}{費氏數列的樸素遞迴}
long long fib(int n) {
    if (n <= 1) {                    // (1) F(0)=0, F(1)=1
        return n;
    }
    return fib(n - 1) + fib(n - 2);  // (2)(3) 兩個子問題
}
\end{code}

圖~\ref{fig:rec-fib-tree} 把 \texttt{fib(4)} 的呼叫全部畫出來。 同一
個 \texttt{fib(2)} 出現了兩次，\texttt{fib(1)} 出現了三次。 樹愈往下
長，重複的子問題愈多。

\begin{center}
\begin{tikzpicture}[
    recnode/.style={draw=black!55, rounded corners=2pt, inner sep=3pt,
                    font=\ttfamily\scriptsize, align=center, minimum width=1.15cm},
    recdup/.style={draw=memStroke, thick, fill=memFill, rounded corners=2pt,
                   inner sep=3pt, font=\ttfamily\scriptsize, align=center,
                   minimum width=1.15cm}
]
    \node[recnode] (f4) at (0,0) {fib(4)};
    \node[recnode] (f3) at (-2.4,-1.35) {fib(3)};
    \node[recdup]  (f2r) at (2.4,-1.35) {fib(2)};
    \node[recdup]  (f2l) at (-3.6,-2.7) {fib(2)};
    \node[recnode] (f1a) at (-1.2,-2.7) {fib(1)};
    \node[recnode] (f1b) at (1.2,-2.7) {fib(1)};
    \node[recnode] (f0a) at (3.6,-2.7) {fib(0)};
    \node[recnode] (f1c) at (-4.5,-4.05) {fib(1)};
    \node[recnode] (f0b) at (-2.7,-4.05) {fib(0)};
    \draw[diagarrow] (f4) -- (f3);
    \draw[diagarrow] (f4) -- (f2r);
    \draw[diagarrow] (f3) -- (f2l);
    \draw[diagarrow] (f3) -- (f1a);
    \draw[diagarrow] (f2r) -- (f1b);
    \draw[diagarrow] (f2r) -- (f0a);
    \draw[diagarrow] (f2l) -- (f1c);
    \draw[diagarrow] (f2l) -- (f0b);
\end{tikzpicture}
\diagcaption{\texttt{fib(4)} 的遞迴樹。 橘色的 \texttt{fib(2)} 被完整
算了兩遍；樹的形狀說明時間隨 $n$ 指數成長，深度仍只有 $\Theta(n)$。}
\label{fig:rec-fib-tree}
\end{center}

令 $T(n)$ 為 \texttt{fib(n)} 的呼叫次數，則 $T(0)=T(1)=1$，
$T(n)=T(n-1)+T(n-2)+1$。 可以驗證 $T(n) = 2F(n+1)-1$，所以
\texttt{fib(5)} 呼叫 $15$ 次，\texttt{fib(10)} 呼叫 $177$ 次，
\texttt{fib(40)} 大約三億三千萬次。 時間是 $\Theta(\varphi^n)$，其中
$\varphi = (1+\sqrt{5})/2 \approx 1.618$；空間則是去程最深的那一條鏈，
$\Theta(n)$。 兩個子問題讓學生特別容易只模擬其中一邊（一步模型），畫
遞迴樹就是為了把兩邊都攤開來看。

網格走格是同一個形狀的另一個例子：在 $r$ 列 $c$ 行的格子上，從左下那
一格走到右上那一格，每次只能向右或向上。 走法數等於
\[
\mathrm{paths}(r, c) =
\begin{cases}
1 & \text{若 } r=1 \text{ 或 } c=1,\\
\mathrm{paths}(r-1,c)+\mathrm{paths}(r,c-1) & \text{否則。}
\end{cases}
\]
第一列或第一行只剩一條直路，這就是基礎情況；一般格子拆成先走一步再解兩
個較小的格子。

\begin{example}{c}{(grid-paths.c) 網格走格的走法數}
#include <stdio.h>

int paths(int r, int c) {
    if (r == 1 || c == 1) {     // (1) 只剩一列或一行
        return 1;
    }
    return paths(r - 1, c) + paths(r, c - 1);  // (2)(3)
}

int main(void) {
    int r, c;
    scanf("%d %d", &r, &c);
    printf("%d\n", paths(r, c));
}
\end{example}

\iosample{
3 3 \\
}{
6 \\
}

$3 \times 3$ 的走法有 $6$ 種，對應 $\binom{4}{2} = 6$。 跟費氏數列一
樣，樸素寫法會把同一塊子網格算很多遍；格子稍大就該改用下一節的記憶化，
或直接用組合公式。

河內塔也是兩個遞迴呼叫，不過子問題是我們自己造出來的，不是輸入裡現成
的左右兩半。 $n$ 個盤子從 A 柱移到 C 柱、B 當輔助，步驟是：

\begin{enumerate}
\item 把上面 $n-1$ 個從 A 移到 B（C 當輔助）；
\item 把最底下面那個從 A 移到 C；
\item 再把 $n-1$ 個從 B 移到 C（A 當輔助）。
\end{enumerate}

$n = 0$ 時沒有盤子，什麼都不用搬，這就是基礎情況。 $n = 1$ 時會先對
$0$ 個盤子呼叫兩次（立刻回來），中間把那一個盤子印出來。 這樣寫的好處
是輸入 $0$ 也會停，不必在基礎情況裡再放一次 \texttt{printf}。

\begin{example}{c}{(hanoi.c) 河內塔}
#include <stdio.h>

void hanoi(int n, char from, char to, char aux) {
    if (n == 0) {              // (1) 沒有盤子，什麼都不用搬
        return;
    }
    hanoi(n - 1, from, aux, to);                         // 上面 n-1 移到輔助柱
    printf("Move disk %d from %c to %c\n", n, from, to); // 最大盤移到目標
    hanoi(n - 1, aux, to, from);                         // 上面 n-1 移到目標
}

int main(void) {
    int n;
    scanf("%d", &n);
    hanoi(n, 'A', 'C', 'B');
}
\end{example}

\iosample{
3 \\
}{
Move disk 1 from A to C \\
Move disk 2 from A to B \\
Move disk 1 from C to B \\
Move disk 3 from A to C \\
Move disk 1 from B to A \\
Move disk 2 from B to C \\
Move disk 1 from A to C \\
}

令 $T(n)$ 為移動 $n$ 個盤子的步數。 $T(0) = 0$，
$T(n) = 2T(n-1) + 1$。 展開得到
\[
T(n) = 2^{n-1} + 2^{n-2} + \cdots + 2 + 1 = 2^n - 1.
\]
$n = 3$ 時 $T(3) = 7$，與上面印出的七行相符；$n = 0$ 時 $2^0 - 1 = 0$，
什麼都不印。 遞迴關係式就是程式的形狀；把 $T(n)$ 解出來，才知道河內塔
為什麼不能靠多加幾行迴圈硬幹。

\keypoint{函式裡出現兩次遞迴呼叫時，執行過程是一棵樹。 先畫出遞迴樹、
再數重複的子問題，比只追其中一條鏈更能看出時間為什麼變成指數。}

\begin{drills}[label=drills:rec-fib-hanoi]
\item 畫出 \texttt{fib(3)} 的遞迴樹。 \texttt{fib(1)} 出現幾次？
\item 抓蟲：若河內塔的基礎情況寫成 \texttt{n == 1} 並在那裡
\texttt{printf}，輸入 $0$ 會怎樣？
\item 用 $T(n) = 2T(n-1)+1$、$T(0)=0$ 算出 $T(4)$。 跟 $2^4-1$ 一樣嗎？
\item $3 \times 3$ 格子從左下走到右上、只能向右或向上，走法數是多少？
先用遞迴關係、再用 $\binom{4}{2}$ 核對。
\end{drills}

\section{記憶化：從指數到線性}
\label{sec:rec-memo}

\texttt{fib(4)} 的樹裡，\texttt{fib(2)} 被完整算了兩遍。 $n$ 再大一點，
同一棵子樹會重複指數多次。 \textbf{記憶化}（memoization）就是開一個表
把算過的 $F(k)$ 存下來：要算 $F(n)$ 時先查表，有就直接回傳，沒有才真
的遞迴，算完再寫進表。 每個 $k$ 只會真正計算一次，時間從
$\Theta(\varphi^n)$ 降到 $\Theta(n)$。

表的初始值要能跟任何合法的費氏數區分開。 $F(n)$ 都是非負的，所以用
$-1$ 表示還沒算過。 範例裡用迴圈把表填成 $-1$，用意清楚，也適用於任意
哨兵值。

\begin{example}{c}{(fib-memo.c) 記憶化後的費氏數列}
#include <stdio.h>

long long memo[100];

long long fib(int n) {
    if (n <= 1) {
        return n;
    }
    if (memo[n] != -1) {       // 算過就直接用
        return memo[n];
    }
    memo[n] = fib(n - 1) + fib(n - 2);
    return memo[n];
}

int main(void) {
    int n;
    scanf("%d", &n);
    if (n < 0 || n >= 100) {
        return 0;
    }
    for (int i = 0; i < 100; i++) {
        memo[i] = -1;
    }
    printf("%lld\n", fib(n));
}
\end{example}

\iosample{
40 \\
}{
102334155 \\
}

\texttt{fib(40)} 用樸素寫法要三億多次呼叫，記憶化之後每個 $k$ 最多進
入函式本體一次，瞬間印出 $102334155$。 $\mathrm{fib}(47) = 2971215073$
已經超過 \texttt{INT\_MAX}，所以這裡繼續用 \texttt{long long}。
\texttt{long long} 自己也有頂：
\[
F(92) = 7540113804746346429 < \texttt{LLONG\_MAX},
\]
\[
F(93) = 12200160415121876738 > \texttt{LLONG\_MAX}.
\]
表開到 $100$ 綽綽有餘，因為還沒用到第 $99$ 格，$\texttt{long long}$ 已經
先到極限。 \texttt{main} 擋 $n < 0$ 或 $n \ge 100$，只是在防
\texttt{memo[n]} 越界。 $n$ 在 $93$ 到 $99$ 之間不會越界，但會溢位：表
的長度與型別上限是兩個獨立的界，擋越界擋不住未定義行為。

這張表放在檔案範圍，由 \texttt{fib} 直接改寫。 它存的是 $F(k)$ 的值，
對每一條呼叫路徑都一樣，寫進去就不再改，所以是快取，不是每一層各自的
狀態。 後面回溯法那節的棋盤才是路徑狀態：值取決於目前走
在哪一條分支，換一條分支就不同。 真要嚴謹，記憶表一樣可以由
\texttt{main} 配置再傳指標進來；這裡放檔案範圍，只是少傳一個參數。

有人會寫 \texttt{memset(memo, -1, sizeof(memo))}。 這招能用，是因為
\texttt{long}／\texttt{long long} 在二補數裡，每個位元組都填 \texttt{0xFF}
恰好等於 $-1$。 它只對全位元為 $1$ 的值（通常是 $-1$）以及全 $0$ 成立；
若想把表初始化成 $1$ 或 $42$，\texttt{memset} 會把每個位元組都設成同一個
數字，結果完全不是我們要的。 需要任意哨兵值時，用迴圈填最不容易誤導。

\stylepoint{記憶化的表要跟遞迴一起想：表的維度對應遞迴的參數，表的初
值對應還沒算過。 先寫對樸素遞迴，再加表，比較不容易把狀態漏掉。}

\keypoint{樹狀遞迴的指數來自重複子問題。 記憶化讓每個子問題只算一次，
時間降到子問題的個數；它沒有改遞迴的寫法，只是把信仰之躍的結果存下來
重用。}

\begin{drills}[label=drills:rec-memo]
\item 若把 \texttt{memset(memo, 1, sizeof(memo))} 拿來初始化費氏表，表
裡的每個 \texttt{long long} 會變成什麼？ 能不能當哨兵用？
\item 用 \texttt{long long} 算 $F(93)$ 會碰到哪一種上限：表的長度，還是
整數溢位？
\item 樸素 \texttt{fib(5)} 呼叫 $15$ 次；加上記憶化之後，每個 $k$ 最多
真正計算幾次？
\end{drills}

\section{在鏈狀結構上遞迴：串列與二元樹}
\label{sec:rec-list}

第~\ref{ch:structure}~章讓我們能把不同欄位捆成一個單位，第~\ref{ch:pointer}~章
讓一個欄位可以指向另一個同型的節點。 兩者合在一起，資料的定義本身就是
遞迴的：

\begin{code}{c}{鏈結串列的節點：定義裡就有自己}
struct Node {
    int v;
    struct Node *next;   // 要嘛是 NULL，要嘛指向另一個 Node
};
\end{code}

一份串列要嘛是空的（\texttt{NULL}），要嘛是一個節點接上一份更短的串
列。 這跟空區間、空字串是同一種形狀，所以程式也長得一樣：基礎情況處理
\texttt{NULL}，遞迴情況處理 \texttt{p->v} 與 \texttt{p->next}。

\begin{code}{c}{串列長度：形狀跟著資料走}
int length(struct Node *p) {
    if (p == NULL) {           // (1) 空串列長度 0
        return 0;
    }
    return 1 + length(p->next);  // (2)(3) 這一格加尾巴的長度
}
\end{code}

正序印與反序印只差 \texttt{printf} 放在遞迴呼叫的前面還是後面。 這是
回程最好認的一對例子：兩份程式的去程完全相同，印出的順序卻相反，差別
全在那一行輸出發生在去程還是回程。

\begin{example}{c}{(list-print.c) 正序印與反序印}
#include <stdio.h>

struct Node {
    int v;
    struct Node *next;
};

int length(struct Node *p) {
    if (p == NULL) {
        return 0;
    }
    return 1 + length(p->next);
}

void print_forward(struct Node *p) {
    if (p == NULL) {
        return;
    }
    printf("%d ", p->v);       // 去程：先印自己
    print_forward(p->next);
}

void print_reverse(struct Node *p) {
    if (p == NULL) {
        return;
    }
    print_reverse(p->next);    // 回程：先讓尾巴印完
    printf("%d ", p->v);
}

int main(void) {
    struct Node c = {3, NULL};
    struct Node b = {2, &c};
    struct Node a = {1, &b};

    printf("%d\n", length(&a));
    print_forward(&a);
    printf("\n");
    print_reverse(&a);
    printf("\n");
}
\end{example}

\iosample{}{
3 \\
1 2 3  \\
3 2 1  \\
}

\texttt{main} 用三個區域變數串成 $1 \rightarrow 2 \rightarrow 3$，不必
\texttt{malloc}。 \texttt{length} 回傳 $3$。 \texttt{print\_forward}
在往下走的路上依序印 $1$、$2$、$3$；\texttt{print\_reverse} 先走到
\texttt{NULL} 再一路回來印，所以是 $3$、$2$、$1$。 兩份函式只有
\texttt{printf} 與遞迴呼叫的先後不同。

二元樹的定義同樣是遞迴的：空樹，或一個節點加上左子樹與右子樹。 高度、
節點數、中序走訪都直接對應這個形狀。 空樹高度 $0$、節點數 $0$；非空樹
的高度是兩棵子樹高度的較大值再加 $1$，節點數是兩棵子樹的節點數再加
$1$。

\begin{example}{c}{(tree-height.c) 二元樹的高度、節點數與中序走訪}
#include <stdio.h>

struct TreeNode {
    int v;
    struct TreeNode *left;
    struct TreeNode *right;
};

int height(struct TreeNode *t) {
    if (t == NULL) {
        return 0;
    }
    int hl = height(t->left);
    int hr = height(t->right);
    if (hl > hr) {
        return 1 + hl;
    }
    return 1 + hr;
}

int count(struct TreeNode *t) {
    if (t == NULL) {
        return 0;
    }
    return 1 + count(t->left) + count(t->right);
}

void inorder(struct TreeNode *t) {
    if (t == NULL) {
        return;
    }
    inorder(t->left);
    printf("%d ", t->v);
    inorder(t->right);
}

int main(void) {
    struct TreeNode n4 = {4, NULL, NULL};
    struct TreeNode n5 = {5, NULL, NULL};
    struct TreeNode n2 = {2, &n4, &n5};
    struct TreeNode n3 = {3, NULL, NULL};
    struct TreeNode n1 = {1, &n2, &n3};

    printf("%d\n", height(&n1));
    printf("%d\n", count(&n1));
    inorder(&n1);
    printf("\n");
}
\end{example}

\iosample{}{
3 \\
5 \\
4 2 5 1 3  \\
}

這棵樹的根是 $1$，左子樹是 $2$（底下還有 $4$、$5$），右子樹是 $3$。 最
深的鏈有三層，高度 $3$；五個節點，\texttt{count} 印 $5$。 中序是左、根、
右，所以先把左子樹 $4, 2, 5$ 走完，再印根 $1$，再走右子樹 $3$。

\begin{center}
\begin{tikzpicture}[
    recnode/.style={draw=black!55, circle, inner sep=2pt,
                    font=\ttfamily\small, minimum size=0.7cm}
]
    \node[recnode] (n1) at (0,0) {1};
    \node[recnode] (n2) at (-1.6,-1.3) {2};
    \node[recnode] (n3) at (1.6,-1.3) {3};
    \node[recnode] (n4) at (-2.4,-2.6) {4};
    \node[recnode] (n5) at (-0.8,-2.6) {5};
    \draw[diagarrow] (n1) -- (n2);
    \draw[diagarrow] (n1) -- (n3);
    \draw[diagarrow] (n2) -- (n4);
    \draw[diagarrow] (n2) -- (n5);
\end{tikzpicture}
\diagcaption{範例中那棵樹。 \texttt{height}、\texttt{count}、
\texttt{inorder} 的程式碼形狀，就是這份資料定義的形狀。}
\label{fig:rec-tree}
\end{center}

\keypoint{資料的定義若本身是遞迴的（空，或一個單位接上更小的同類資
料），程式就照那個形狀寫：基礎情況對空，遞迴情況對這一個單位加上尾巴
或子樹。 這是結構遞迴，也是最不容易寫錯的一種。}

\begin{drills}[label=drills:rec-list]
\item 預測：對 $1 \rightarrow 2 \rightarrow 3$ 呼叫 \texttt{print\_reverse}，
印出的順序是什麼？ \texttt{printf} 發生在去程還是回程？
\item 把 \texttt{print\_reverse} 的 \texttt{printf} 改到呼叫前面。
印出順序變成什麼？ 和 \texttt{print\_forward} 還有沒有差別？
\item 只有根、左右都是 \texttt{NULL} 的樹，\texttt{height} 與
\texttt{count} 各回傳多少？
\item 一棵只有三個節點的樹：根 $1$，左子 $2$，右子 $3$。 中序走訪印出
什麼？
\end{drills}

\section{結構遞迴 vs 生成遞迴}
\label{sec:rec-struct-vs-gen}

上一節的 \texttt{length}、\texttt{height} 都在輸入裡已經存在的子結構
上遞迴：串列的尾巴、樹的左子樹。 子問題是原資料的一部分，每走一步資料
就嚴格變小，停機幾乎自動成立。 這種寫法叫做\textbf{結構遞迴}
（structural recursion）。 整數上的 \texttt{sum(n-1)}、陣列區間的
$[lo+1, hi)$、字串的 \texttt{s+1}，也都可以看成結構遞迴：我們走的是資
料自己的結構。

有些問題的子問題並不躺在輸入裡，而是我們另外造出來的。 河內塔要把
$n-1$ 個盤子移到輔助柱，那 $n-1$ 個盤子在原問題裡並不是一份現成的子結
構；輾轉相除的 \texttt{gcd(b, a \% b)} 更明顯，第二個引數是新算出來的
餘數。 這種寫法叫做\textbf{生成遞迴}（generative recursion）。 因為子
問題是生成的，停機必須另外論證：每次呼叫都要證明某個度量在下降，例如
盤子數 $n$、或 \texttt{gcd} 的第二個引數（餘數嚴格小於 $b$）。

\begin{example}{c}{(gcd.c) 輾轉相除：生成出來的子問題}
#include <stdio.h>

int gcd(int a, int b) {
    /* 前提：a, b >= 0。負數時終止論證不成立。 */
    if (b == 0) {              // (1) 餘數為 0，答案就是 a
        return a;
    }
    return gcd(b, a % b);      // (2)(3) 新的一對 (b, a%b)
}

int main(void) {
    int a, b;
    scanf("%d %d", &a, &b);
    printf("%d\n", gcd(a, b));
}
\end{example}

\iosample{
48 18 \\
}{
6 \\
}

\texttt{gcd(48, 18)} $\rightarrow$ \texttt{gcd(18, 12)} $\rightarrow$
\texttt{gcd(12, 6)} $\rightarrow$ \texttt{gcd(6, 0)}，回傳 $6$。 前提是
$a, b \ge 0$：在這個前提下，每一次的第二個引數都是非負整數且嚴格遞減，
所以一定走到 $0$。 這份終止論證要自己寫，編譯器不會幫我們檢查。 前提
一破，論證就失效。 C 的 $\%$ 對負運算元可能產生負餘數，
\texttt{gcd(-48, 18)} 仍會停
（$(-48, 18) \rightarrow (18, -12) \rightarrow (-12, 6) \rightarrow (6, 0)$），
但真正下降的是 $|b|$，不是 $b$ 本身。 生成遞迴的停機必須自己論證，論證
所依賴的前提被打破時，論證就不成立。

合併排序把陣列切成左右兩半再合併，切出來的兩半是新的區間，屬於生成遞
迴裡的分治。 區間長度每次減半，終止條件是長度 $0$ 或 $1$。 分治跟結構
遞迴一樣都問那三個問題，只是縮小的方式改成自己切，不再沿著 \texttt{next}
或 \texttt{n-1} 走。

\stylepoint{能寫成結構遞迴就先寫結構遞迴，停機幾乎免費。 只有子問題必
須另外構造時才用生成遞迴，那時把下降的度量明講出來，例如 $n$、非負的
$b$、或區間長度。}

\begin{drills}[label=drills:rec-struct-vs-gen]
\item \texttt{length(p->next)} 是結構遞迴還是生成遞迴？ \texttt{gcd(b, a \% b)}
呢？
\item 寫出 \texttt{gcd(48, 18)} 的呼叫鏈，直到基礎情況。
\item 河內塔每次呼叫的 $n$ 都比上一層小。 這份終止論證夠不夠？ 為什麼
它仍算生成遞迴？
\end{drills}

\section{尾遞迴與 C 的現實}
\label{sec:rec-tail}

若遞迴呼叫本身就是 \texttt{return} 的全部內容，這個呼叫位在\textbf{尾
位置}（tail position），這種函式叫做\textbf{尾遞迴}（tail recursion）。
\texttt{sum\_iter} 與下面的 \texttt{factorial\_acc} 都是：回程上不欠乘
法，累加器已經把帳結清。

\begin{code}{c}{階乘的尾遞迴形式}
long long factorial_acc(int n, long long acc) {
    if (n <= 1) {
        return acc;
    }
    return factorial_acc(n - 1, acc * n);  // 呼叫本身就是 return 的全部
}
\end{code}

\begin{example}{c}{(fact-tail.c) 用累加器寫尾遞迴階乘}
#include <stdio.h>

long long factorial_acc(int n, long long acc) {
    if (n <= 1) {
        return acc;
    }
    return factorial_acc(n - 1, acc * n);
}

int main(void) {
    int n;
    scanf("%d", &n);
    printf("%lld\n", factorial_acc(n, 1));
}
\end{example}

\iosample{
10 \\
}{
3628800 \\
}

有些語言保證把尾遞迴編譯成迴圈，frame 可以重用，深度多大都不佔堆疊。
C 標準\textbf{不保證}尾呼叫最佳化（tail-call optimization, TCO）。
gcc 通常要到 \texttt{-O2}／\texttt{-O3}（會打開
\texttt{-foptimize-sibling-calls}）才做；clang 有時從 \texttt{-O1} 就
開始做。 關掉最佳化（\texttt{-O0}）時，同一個尾遞迴函式一樣會把堆疊用
完。 鏈上只要有一個呼叫不在尾位置，整條最佳化就失效。

下面這個極短的例子可以自己編譯來看。 \texttt{foo} 是標準的尾遞迴，它
什麼工作都不做、永遠呼叫自己：

\begin{code}{c}{用來觀察 TCO 在不在的極短例子}
int foo(void) {
    return foo();
}

int main(void) {
    return foo();
}
\end{code}

在 Apple Clang 15（macOS, arm64）上實測：\texttt{-O0} 編譯後執行，程式
因堆疊溢位收到 segmentation fault（結束碼 $139$）；\texttt{-O2} 編譯後，
同一支程式變成一個不佔堆疊的無窮迴圈，不會爆掉，會一直跑到你把它停下。
這份編譯器從 \texttt{-O1} 就已經會做 TCO；gcc 在 Linux x86-64 上通常要
到 \texttt{-O2} 才打開 \texttt{-foptimize-sibling-calls}。 結果依編譯器
與版本而異。 這個對照只說明編譯器\emph{可以}把尾遞迴變成迴圈，不表示你
寫作業時可以依賴它。 可攜的做法是自己把尾遞迴改成 \texttt{while}，或像
下一節那樣用陣列模擬堆疊。

\keypoint{C 不保證尾呼叫最佳化。 就算函式寫成尾遞迴，\texttt{-O0} 或
換一個編譯器仍可能以 $\Theta(n)$ 的堆疊跑。 深度可能很大時，可攜的修
法是改成迴圈，或自己準備一個陣列當堆疊。}

\begin{drills}[label=drills:rec-tail]
\item \texttt{return n * factorial(n - 1);} 裡，哪一部分讓這個呼叫不在
尾位置？
\item 同一支 \texttt{foo}，\texttt{-O0} 與 \texttt{-O2} 各會看到什麼現
象？ 寫作業時能依賴哪一種？
\item 把 \texttt{rec\_strlen} 改成尾遞迴（累加器從 $0$ 加起），再改成
一個 \texttt{while} 迴圈。
\end{drills}

\section{把遞迴轉成迴圈：用顯式堆疊模擬}
\label{sec:rec-explicit-stack}

\texttt{print\_reverse} 在回程上才印，去程經過的節點必須先記住。 這些
節點本來住在一層層的 frame 裡；我們也可以自己開一個指標陣列當堆疊，去
程把節點壓進去，回程彈出來再印。 這樣深度受我們宣告的陣列大小限制，跟
系統給主執行緒的堆疊上限脫鉤。

順序在這裡不能交換。 若在去程經過一個節點時就印，得到的是正序，連堆疊
都不需要，那正是 \texttt{print\_forward}。 反序印出的恰好是壓入順序的
顛倒，所以非壓不可。 對照階乘：乘法可交換，去程邊走邊乘就得到 $n!$，
那種情況顯式堆疊是多餘的。

\begin{example}{c}{(print-rev-stack.c) 用陣列模擬反向列印的呼叫堆疊}
#include <stdio.h>

struct Node {
    int v;
    struct Node *next;
};

void print_reverse_stack(struct Node *p) {
    struct Node *pending[64];
    int top = 0;
    while (p != NULL) {            // 去程：把節點壓進去
        if (top >= 64) {           // 陣列滿了：拒絕，不是印出一部分
            return;
        }
        pending[top] = p;
        top++;
        p = p->next;
    }
    while (top > 0) {              // 回程：後進先出
        top--;
        printf("%d ", pending[top]->v);
    }
}

int main(void) {
    struct Node c = {3, NULL};
    struct Node b = {2, &c};
    struct Node a = {1, &b};
    print_reverse_stack(&a);
    printf("\n");
}
\end{example}

\iosample{}{
3 2 1  \\
}

以 $1 \rightarrow 2 \rightarrow 3$ 走一遍。 去程結束時 \texttt{pending}
裡由底到頂是指向 $1$、$2$、$3$ 的指標；回程依序彈出，印 $3$、$2$、$1$。
這跟圖~\ref{fig:rec-stack-up} 的回程是同一件事，只是堆疊從系統的
frame 換成我們自己的陣列，掛著的從待乘的整數換成待印的節點。 樹狀遞迴
（例如走訪二元樹）也能這樣轉：把節點指標壓進陣列，用迴圈模擬去程與回
程，只是每個 frame 要記的狀態比一個指標多。

\texttt{pending[64]} 限制了深度上限是 $64$。 串列若可能更長，陣列要開
到問題的最大深度，或改用 \texttt{malloc} 做動態堆疊。 這是用空間換對
系統堆疊的控制。 \texttt{top >= 64} 時直接 \texttt{return}，表示拒絕
這次列印，不是印出前 $64$ 個就算完成。

\begin{drills}[label=drills:rec-explicit-stack]
\item 對 $1 \rightarrow 2 \rightarrow 3$，去程結束後 \texttt{pending}
由底到頂對應哪些值？ 回程印出什麼？
\item 若在第一個迴圈裡、壓進去的當下就 \texttt{printf}，印出順序是什麼？
那還需要第二個迴圈嗎？
\item 階乘去程直接 \texttt{result *= n} 就能算。 為什麼本節的反向列印
不能同樣把陣列拿掉？
\item 拿掉 \texttt{top >= 64} 的檢查，串列有 $70$ 個節點時會發生哪一種
錯誤？
\end{drills}

\section{C 遞迴的陷阱}
\label{sec:rec-pitfalls}

這一節把前面散落的 C 細節收成一張清單。 Java 或 Python 的遞迴教材通常
不會碰到這些，因為那些語言要嘛有垃圾收集與例外，要嘛禁止回傳區域變數
的位址。

\textbf{無窮遞迴的結局是 segmentation fault。} Python 會拋
\texttt{RecursionError} 並印出一疊呼叫紀錄；C 把系統堆疊用完之後，下
一次寫入就踩到不能寫的記憶體，作業系統送出 \texttt{SIGSEGV}。 看到
\texttt{Segmentation fault} 時，若函式是遞迴的，第一個要查的就是基礎
情況有沒有寫到、進度有沒有真的讓輸入變小。

\textbf{堆疊大小是有限的、而且依平台而變。} 主執行緒的預設堆疊在 Linux
上常見約 $8$ MB（\texttt{ulimit -s} 常顯示 $8192$，單位是 KB），
Windows 常見 $1$ MB（連結時可用 \texttt{/STACK} 改）。 這些是典型預設
值，不是保證。 自己的環境用 \texttt{ulimit -s} 查看。 淺的 frame（只
有幾個 \texttt{int}）大約能到數十萬層；frame 裡若有大型區域陣列，深度
一下子就只剩幾百層。

\textbf{不要回傳指向區域變數的指標。} 區域陣列住在這一層 frame 裡，
\texttt{return} 之後那塊記憶體已經不屬於這個函式。 呼叫者拿到的是懸空
指標，再用它就是未定義行為。

\begin{code}{c}{回傳區域陣列的位址：frame 一彈就懸空}
char *label(void) {
    char buf[16];
    buf[0] = 'A';
    buf[1] = '\0';
    return buf;    // buf 隨這一層一起消失
}
\end{code}

若確實需要把字串交出去，讓呼叫者提供緩衝區當參數，或用 \texttt{malloc}
配置到堆積上並講清楚由誰 \texttt{free}。

\textbf{不要在遞迴 frame 裡放大型區域陣列。} 下面這種寫法每一層都再要
一塊 $10000$ 個 \texttt{int}（約 $40$ KB）。 以 $8$ MB 的堆疊來算，兩
百層就滿了。

\begin{code}{c}{每一層都帶一塊大陣列，深度很快就爆}
int rec(int n) {
    int buf[10000];    // 每一層約 40 KB
    buf[0] = n;
    if (n == 0) {
        return 0;
    }
    return rec(n - 1);
}
\end{code}

需要那塊陣列時，改成在 \texttt{main} 宣告一次再當指標傳進去，讓所有層
共用同一塊，或改用堆積。

\textbf{基礎情況的 off-by-one。} 空與單元素是兩種不同的最小情況。 求和、
長度對空有自然答案 $0$；求最大值、二元搜尋找到的那個元素，空區間沒有
意義，基礎情況要寫成單元素或找不到。 少寫一個基礎情況，輕則答案錯，重
則無窮遞迴。

\textbf{有號整數溢位。} $13!$ 與 $\mathrm{fib}(47)$ 都超出 32-bit
\texttt{int}。 溢位在 C 是未定義行為，不能靠印出來的數字回頭推算它
包了幾圈。

\textbf{用全域變數當遞迴的工作記憶。} 把棋盤、解的個數放在檔案範圍，讓
遞迴函式直接改寫，等於把每一層 frame 本來該各自擁有的狀態藏到同一個地
方。 這樣的函式不能重入、不好測試，也讓副本模型從圖上消失：你看不到每
一層各自拿著一份狀態。 下一節會把 N 皇后的全域寫法跟參數寫法並排放。
第~\ref{sec:rec-memo}~節的表也在檔案範圍，但那是快取：$F(k)$ 對所有路
徑都相同，寫進去就不再改。 這裡禁止的是路徑狀態，也就是棋盤這種每一條
分支各自不同的東西。

\keypoint{C 的遞迴錯誤通常以 segmentation fault、懸空指標或未定義行為
出現，不會給一則友善的遞迴太深訊息。 寫完先問三件事：基礎情況、進度、
這一層的區域變數會不會在 \texttt{return} 之後還被用到。}

\begin{drills}[label=drills:rec-pitfalls]
\item \texttt{label} 回傳 \texttt{buf} 之後，呼叫者立刻用那個指標印字串，
是哪一種錯誤？
\item 每一層放 \texttt{int buf[10000]}（約 $40$ KB），堆疊以 $8$ MB 計，
大概幾層就滿？
\item \texttt{array\_max} 傳空區間，對應本節哪一條陷阱？ 結果是答案錯，
還是無窮遞迴？
\end{drills}

\section{回溯法：選擇、探索、撤銷}
\label{sec:rec-backtrack}

前面的遞迴都在一條或兩條固定的縮小路上走。 \textbf{回溯}
（backtracking）多了一層：在每一層有若干個選擇，試一個、往下探索，若
這條路走不通或走完了，就把選擇撤銷，改試下一個。 結構是

\begin{code}{text}{回溯的三步：選擇、探索、撤銷}
void search(目前的狀態) {
    if (已經完成或無解) {         // (1)
        記錄或回傳;
        return;
    }
    for (每一個可以做的選擇) {
        做選擇;                    // choose
        search(新的狀態);          // explore
        撤銷選擇;                  // unchoose
    }
}
\end{code}

若狀態是用值傳進去的（例如子集和裡的剩餘容量 $k$、還剩幾個元素 $n$），
子呼叫結束後呼叫者手上的 $k$、$n$ 根本沒被改過，不必手動撤銷。 若狀態
是同一塊陣列被就地改寫（排列裡交換字元、棋盤上放皇后），子呼叫結束後
必須把改過的格子還原，否則下一條分支看到的是上一條分支留下的殘局。
C 的陣列傳的是位址，所以就地修改是預設行為，撤銷一定要寫出來。

先看不需要手動撤銷的那一種：子集和。 $n$ 個\textbf{正整數}，問有幾種
子集的和剛好是 $k$。 對最後一個數字 $a[n-1]$ 只有兩個選擇：不拿它，或
拿它。 兩個選擇都把 $n$ 減 $1$，拿的那條再把 $k$ 減去 $a[n-1]$。 參數
都是整數、傳值，子呼叫結束後這一層的 $n$、$k$ 還是原值。

\texttt{k == 0} 立刻回傳 $1$，以及 \texttt{k < 0} 剪枝，都只在元素為正
時才對。 若陣列裡有 $0$，$\texttt{k == 0}$ 會在湊滿的瞬間停止列舉剩下
的元素，於是只差幾個 $0$ 的子集被算成同一種，結果是少算。 例如
$a = \{0, 2\}$、$k = 2$ 只會回傳 $1$，但 $\{2\}$ 與 $\{0, 2\}$ 都成立。
若有負數，$\texttt{k < 0}$ 也不能當已經超過的停止條件。 下面的程式預設
輸入都是正整數。

\begin{example}{c}{(subset.c) 子集和的計數：拿或不拿}
#include <stdio.h>

int count_subset(int *a, int n, int k) {
    /* 前提：a[0..n) 皆為正整數。 */
    if (k == 0) {              // 剛好湊滿
        return 1;
    }
    if (n == 0 || k < 0) {     // 沒數字了，或已經超過
        return 0;
    }
    return count_subset(a, n - 1, k)                 // 不拿 a[n-1]
         + count_subset(a, n - 1, k - a[n - 1]);     // 拿 a[n-1]
}

int main(void) {
    int n, k, a[20];
    scanf("%d %d", &n, &k);
    if (n < 0 || n > 20) {
        return 0;
    }
    for (int i = 0; i < n; i++) {
        scanf("%d", &a[i]);
    }
    printf("%d\n", count_subset(a, n, k));
}
\end{example}

\iosample{
4 5 \\
1 2 3 4 \\
}{
2 \\
}

$\{1,2,3,4\}$ 裡和為 $5$ 的子集有 $\{1,4\}$ 與 $\{2,3\}$ 兩種。 每一層
把問題縮小成前 $n-1$ 個數字，兩條分支的 $k$ 不同。 $0$-$1$ 背包的拿或
不拿也是同一棵樹，只是回傳值從計數改成最大價值。 $n$ 到 $20$ 附近，
$2^n$ 次呼叫還勉強跑得完；$n$ 再大就要剪枝或改成之後會學的動態規劃。

N 皇后則必須手動撤銷。 在 $n \times n$ 棋盤放 $n$ 個皇后，彼此不同列、
不同行、不同對角線。 我們第 \texttt{row} 列決定放在哪一欄，\texttt{board[i]}
記錄第 $i$ 列的皇后在哪一欄。 棋盤是同一塊陣列，傳的是指標，所以放上
去之後一定要在子呼叫回來時拿掉。

常見的偷懶寫法是把棋盤和解的個數放在檔案範圍，讓 \texttt{solve} 直接
改：

\begin{code}{c}{反例：棋盤放在檔案範圍，副本從圖上消失}
int board[20];
int solutions;

int is_safe(int row, int col) {
    for (int i = 0; i < row; i++) {
        int dcol = board[i] - col;
        int drow = i - row;
        if (dcol == 0 || dcol == drow || dcol == -drow) {
            return 0;
        }
    }
    return 1;
}

void solve(int row, int n) {
    if (row == n) {
        solutions++;
        return;
    }
    for (int col = 0; col < n; col++) {
        if (is_safe(row, col)) {
            board[row] = col;      // 所有層改同一塊
            solve(row + 1, n);
            board[row] = -1;
        }
    }
}
\end{code}

\texttt{is\_safe} 沒有棋盤參數，因為它直接讀檔案範圍那一份
\texttt{board}。 兩次搜尋不能重疊進行：裡面那一次會把外面那一次的棋盤
覆寫掉。 堆疊圖上本來每一層該有自己的狀態，現在全部摺進兩個全域變數。
下面的範例改成由 \texttt{main} 提供陣列、當參數往下傳。

\begin{example}{c}{(nqueens.c) N 皇后：狀態當參數，撤銷寫出來}
#include <stdio.h>

int is_safe(int row, int col, int *board) {
    for (int i = 0; i < row; i++) {
        int dcol = board[i] - col;
        int drow = i - row;
        if (dcol == 0 || dcol == drow || dcol == -drow) {
            return 0;
        }
    }
    return 1;
}

void solve(int row, int n, int *board, int *solutions) {
    if (row == n) {            // (1) n 列都放好了，找到一解
        (*solutions)++;
        return;
    }
    for (int col = 0; col < n; col++) {
        if (is_safe(row, col, board)) {
            board[row] = col;                       // choose
            solve(row + 1, n, board, solutions);    // explore
            board[row] = -1;                        // unchoose
        }
    }
}

int main(void) {
    int n, board[20], solutions = 0;
    scanf("%d", &n);
    if (n < 0 || n > 20) {
        return 0;
    }
    for (int i = 0; i < n; i++) {
        board[i] = -1;
    }
    solve(0, n, board, &solutions);
    printf("%d\n", solutions);
}
\end{example}

\iosample{
8 \\
}{
92 \\
}

$n = 8$ 有 $92$ 組解，這是經典數字。 $n = 4$ 只有 $2$ 組，棋盤小，適
合自己把 \texttt{board} 的變化畫出來核對撤銷有沒有寫對。
\texttt{board} 與 \texttt{solutions} 都由 \texttt{main} 提供、當參數
往下傳，沒有檔案範圍的全域變數。 \texttt{board[row] = -1} 這一行在這
個問題裡即使省略，下一次寫入同一個 \texttt{board[row]} 也會蓋掉舊值，
答案碰巧仍對；那種寫法示範的是錯誤的心智模型。 換成要在走不通時印出棋
盤、或同一列可能被跳過不寫，缺撤銷就會看到上一條分支留下的皇后。 回溯
的紀律是：改了就還原，不要依賴下一次寫入剛好蓋到同一個位置。

$n$ 到 $13$ 附近就要等上一陣子；$n = 20$ 在這支樸素程式上實務跑不完。
緩衝區開 $20$ 是為了示範邊界檢查，不是因為 $n = 20$ 算得完。 這跟記
憶化那節一樣：表的長度與型別上限是兩個獨立的界。

排列生成用交換做出選擇，同樣要交換回來。 範例讀入一個字串，印出把字元
重排之後的所有結果。

\begin{example}{c}{(permute.c) 排列生成：交換、遞迴、交換回來}
#include <stdio.h>
#include <string.h>

void swap(char *a, char *b) {
    char t = *a;
    *a = *b;
    *b = t;
}

void permute(char *s, int l, int r) {
    if (l == r) {              // (1) 後段已經固定，印出目前的 s
        printf("%s\n", s);
        return;
    }
    for (int i = l; i <= r; i++) {
        swap(&s[l], &s[i]);          // choose：把 s[i] 換到位置 l
        permute(s, l + 1, r);        // explore：排列剩下的
        swap(&s[l], &s[i]);          // unchoose：換回來
    }
}

int main(void) {
    char s[20];
    scanf("%19s", s);
    permute(s, 0, (int)strlen(s) - 1);
}
\end{example}

\iosample{
ABC \\
}{
ABC \\
ACB \\
BAC \\
BCA \\
CBA \\
CAB \\
}

位置 $l$ 依序試 $s[l], s[l+1], \ldots, s[r]$ 每一個字元，試完就換回去，
所以下一輪 \texttt{for} 看到的仍是進入這一層時的字串。 輸入是
\texttt{ABC} 這種每個字元都不同的字串時，印出 $3! = 6$ 種排列。 若輸入
是 \texttt{AAB}，同一套程式會把 \texttt{AAB}、\texttt{ABA}、\texttt{BAA}
各印兩次，因為兩個 \texttt{A} 被當成不同的選擇。 要去掉重複，需要在選
擇前先跳過已經用過的相同字元，或改用計數表決定每個字元還能用幾次。 這
支範例故意不處理重複，是為了讓交換與撤銷本身先看清楚。

\keypoint{回溯 = 選擇 $\rightarrow$ 探索 $\rightarrow$ 撤銷。 狀態若以
值傳遞，回傳時自動還原；狀態若是同一塊被就地修改的陣列，撤銷一定要寫
出來。 C 的陣列屬於後者。}

\begin{drills}[label=drills:rec-backtrack]
\item 預測：把 \texttt{permute} 的第二次 \texttt{swap}（unchoose）拿掉，
輸入 \texttt{ABC} 還會印出六個不同的排列嗎？
\item 抓蟲：\texttt{count\_subset} 的陣列改成 $\{0, 2\}$、$k = 2$。
程式回傳幾？ 正解是幾？ $\texttt{k == 0}$ 那條為什麼會少算？
\item 缺 \texttt{board[row] = -1} 的 N 皇后，若在走不通時把棋盤印出來，
會看到什麼？ 只計數解的個數時，為什麼答案可能碰巧仍對？
\item 從零寫：對正整數陣列數有幾種子集的和為 $k$。 先填三個問題，再寫
拿／不拿兩條分支。
\end{drills}

\section{何時用遞迴、何時用迴圈}
\label{sec:rec-when}

現在可以用過程模型來選工具，而不只是優雅或高效這種口號。

線性掃描、只需要一份累加狀態的工作，遞迴過程會在堆疊上留下 $\Theta(n)$
個 frame，迭代過程（迴圈，或尾遞迴再自己改成迴圈）只需 $\Theta(1)$ 的
額外空間。 $1$ 加到 $n$、走訪一陣列求最高值，用 \texttt{for} 寫更直接，
也不會在 $n$ 到 $10^6$ 時把 $8$ MB 的堆疊用完。

資料本身是遞迴的（串列、樹、目錄、運算式），或問題的自然拆法是兩條以
上的分支（費氏、走格、分治、回溯），遞迴讓程式的形狀跟問題的形狀對齊。
這時候省下的是閱讀與除錯的成本。 若深度可能很大（鏈結串列有 $10^5$ 個
節點、或生成遞迴在最壞情況走到 $\Theta(n)$），再改用顯式堆疊把同一份
過程搬到堆積或自己的陣列上，遞迴這個想法仍然先保留。

C 額外要算進去的有兩筆。 第一，沒有 TCO 保證，尾遞迴在 \texttt{-O0}
或另一個編譯器上仍可能爆堆疊。 第二，frame 裡的大型區域陣列會把可用
深度砍得很短。 這兩件事成立時，就算問題看起來適合遞迴，實務上仍應改
成迴圈或顯式堆疊。

一個可以用的判斷順序：先問資料或問題是不是自我相似的。 若是，先用三
個問題把遞迴寫對，再看過程是線性遞迴、線性迭代還是樹狀。 線性迭代且
深度可觀，就轉迴圈；樹狀且有重複子問題，就加記憶化；樹狀且分支是選擇，
就用回溯並把撤銷寫清楚。 最後才用這台機器的堆疊上限、frame 大小，決
定要不要把同一份過程改成顯式堆疊。

\keypoint{遞迴適合自我相似的資料與多路拆分的問題；迴圈適合線性掃描。
選哪一種，看的是過程佔多少堆疊、以及 C 又不保證 TCO；跟哪一種看起來
比較短無關。}

\begin{drills}[label=drills:rec-when]
\item $1$ 加到 $10^7$，用 \texttt{sum} 還是 \texttt{for}？ 用本節的判
斷順序走一遍。
\item 二元樹高度呢？ 鏈結串列長度、而且節點可能有 $10^5$ 個呢？
\item 有重複子問題的走格、與沒有重複但有選擇的 N 皇后，各該在寫對遞迴
之後補上什麼？
\end{drills}

\phantomsection
\addcontentsline{toc}{section}{練習解答}
\section*{練習解答}

各節練習的答案收在這裡。 從零寫與改寫只給能一句話核對的結果；完整程式
見正文片語。

\paragraph{第~\ref{sec:rec-self-sim}~節}
\begin{enumerate}
\item 四個人。 甲拿 \texttt{sum(3)}，乙 \texttt{sum(2)}，丙
\texttt{sum(1)}，丁 \texttt{sum(0)}。
\item 走廊上同時站著四個人，各自等後面的回覆。 把遞迴想成同一個人把
$n$ 改小再算一次，解釋不了為什麼甲在等的時候乙也還在。
\item 空資料夾貢獻 $0$；非空的等於自己這一層的檔案加上每個子資料夾；子
資料夾的問題與原本相同、只是更小。
\end{enumerate}

\paragraph{第~\ref{sec:rec-three-q}~節}
\begin{enumerate}
\item $4 + \texttt{sum(3)}$。 不要再把 \texttt{sum(3)} 展開。
\item 不會停。 \texttt{n == 1} 這條永遠走不到，\texttt{sum(0)} 會呼叫
\texttt{sum(-1)}、\texttt{sum(-2)}，直到堆疊溢位。
\item $n = 0$ 時和為 $0$（若寫成 $n = 1$ 回傳 $1$，呼叫端就不得再傳
$0$）。
\item 基礎情況 $n = 0$ 回傳 $0$；遞迴情況是 $(2n-1) + f(n-1)$。
\end{enumerate}

\paragraph{第~\ref{sec:rec-int}~節}
\begin{enumerate}
\item 有號整數溢位，在 C 是未定義行為。 不能靠印出來的數字判斷它包了
幾圈。
\item 不停。 $n$ 愈來愈小，走不到 $n == 0$，直到堆疊溢位。
\item $20!$ 裝得進 \texttt{long long}；$21!$ 超過 \texttt{LLONG\_MAX}。
\end{enumerate}

\paragraph{第~\ref{sec:rec-stack}~節}
\begin{enumerate}
\item 最深一層 $n = 1$，一共四層 \texttt{factorial}。
\item $3$、$2$、$1$（先印個位）。
\item 回程。 $n = 3$ 那一層等 \texttt{factorial(2)} 把 $2$ 帶回來，才
做 $3 \times 2$。
\end{enumerate}

\paragraph{第~\ref{sec:rec-process}~節}
\begin{enumerate}
\item \texttt{sum} 欠加法，概念空間 $\Theta(n)$；\texttt{sum\_iter} 不
欠，概念空間 $\Theta(1)$。
\item $(\texttt{3},\texttt{0}) \rightarrow (\texttt{2},\texttt{3})
\rightarrow (\texttt{1},\texttt{5}) \rightarrow (\texttt{0},\texttt{6})$。
\item \texttt{return factorial\_acc(n - 1, acc * n);}，呼叫端傳
\texttt{acc = 1}。
\end{enumerate}

\paragraph{第~\ref{sec:rec-array}~節}
\begin{enumerate}
\item 基礎情況 \texttt{lo + 1 == hi} 不會成立，$\texttt{lo}$ 一路加過
去，無窮遞迴、堆疊溢位。
\item 回傳對 $[lo+1, hi)$ 的搜尋結果。
\item $[0, 5) \rightarrow [0, 2) \rightarrow [2, 2)$，空區間回傳 $-1$。
\item 不是。 \texttt{sizeof(a)} 量到的是指標大小。
\end{enumerate}

\paragraph{第~\ref{sec:rec-string}~節}
\begin{enumerate}
\item $0$。
\item 先 \texttt{lo >= hi} 回傳 $1$，再比頭尾，不同就回傳 $0$，相同才
縮小成 $[lo+1, hi-1]$。
\item 累加器版：空字串回傳 \texttt{acc}，否則
\texttt{rec\_strlen(s + 1, acc + 1)}。
\end{enumerate}

\paragraph{第~\ref{sec:rec-fib-hanoi}~節}
\begin{enumerate}
\item \texttt{fib(1)} 出現兩次：一次是 \texttt{fib(3)} 的右子呼叫，一
次在 \texttt{fib(2)} 底下。
\item 輸入 $0$ 會無窮遞迴，直到堆疊溢位。
\item $T(4) = 15$，與 $2^4 - 1$ 相同。
\item $6$，與 $\binom{4}{2} = 6$ 相符。
\end{enumerate}

\paragraph{第~\ref{sec:rec-memo}~節}
\begin{enumerate}
\item 每個 \texttt{long long} 的每個位元組都是 $1$，值是
\texttt{0x0101010101010101}，不是 $1$。 技術上可以當哨兵（那個值不會
跟任何要算的費氏數撞到），但比較式得寫成跟 \texttt{0x0101010101010101}
比，而不是跟 $1$ 比：你以為填進去的和實際填進去的不是同一個東西。 這
正是用迴圈填比較不容易誤導的原因。
\item 整數溢位。 表長 $100$，索引 $93$ 不越界。
\item 每個 $k$ 最多真正計算一次。
\end{enumerate}

\paragraph{第~\ref{sec:rec-list}~節}
\begin{enumerate}
\item $3\ 2\ 1$。 \texttt{printf} 發生在回程。
\item 變成 $1\ 2\ 3$，與 \texttt{print\_forward} 沒有差別。
\item \texttt{height} 回傳 $1$，\texttt{count} 回傳 $1$。
\item $2\ 1\ 3$。
\end{enumerate}

\paragraph{第~\ref{sec:rec-struct-vs-gen}~節}
\begin{enumerate}
\item \texttt{length(p->next)} 是結構遞迴；\texttt{gcd(b, a \% b)} 是
生成遞迴。
\item \texttt{gcd(48, 18)} $\rightarrow$ \texttt{gcd(18, 12)}
$\rightarrow$ \texttt{gcd(12, 6)} $\rightarrow$ \texttt{gcd(6, 0)}。
\item $n$ 變小足以保證會停，但它仍是生成遞迴：搬 $n-1$ 盤到輔助柱這件
事，並不是輸入裡現成的子結構。
\end{enumerate}

\paragraph{第~\ref{sec:rec-tail}~節}
\begin{enumerate}
\item 前面的 \texttt{n *}。 乘法讓這個呼叫不在尾位置。
\item \texttt{-O0} 堆疊溢位（segmentation fault）；\texttt{-O2} 變成
不佔堆疊的無窮迴圈。 寫作業時哪一種都不能當保證。
\item 空字串回傳 \texttt{acc}；否則 \texttt{f(s + 1, acc + 1)}。 迴圈
版就是 \texttt{while (*s) \{ acc++; s++; \}}。
\end{enumerate}

\paragraph{第~\ref{sec:rec-explicit-stack}~節}
\begin{enumerate}
\item 由底到頂對應 $1$、$2$、$3$；回程印 $3\ 2\ 1$。
\item 印 $1\ 2\ 3$。 不需要第二個迴圈，那就是 \texttt{print\_forward}。
\item 印出的順序不可交換。 拿掉陣列就只能在經過時立刻印，得到正序。
\item \texttt{pending} 越界寫入，未定義行為，常見結局是 segmentation
fault。
\end{enumerate}

\paragraph{第~\ref{sec:rec-pitfalls}~節}
\begin{enumerate}
\item 懸空指標，再用它是未定義行為。
\item 約 $200$ 層（$8$ MB $\div$ $40$ KB）。
\item 基礎情況的 off-by-one。 結果是無窮遞迴，不是只把答案算錯。
\end{enumerate}

\paragraph{第~\ref{sec:rec-backtrack}~節}
\begin{enumerate}
\item 不會。 仍印六行，但只有四種：\texttt{ABC}、\texttt{ACB}、
\texttt{CAB}、\texttt{CBA}，缺 \texttt{BAC} 與 \texttt{BCA}。
\item 程式回傳 $1$，正解是 $2$，也就是 $\{2\}$ 與 $\{0, 2\}$。
\texttt{k == 0} 擺在 \texttt{n == 0} 前面，取走 $2$ 之後立刻回傳，剩下
的 $0$ 沒被列舉，所以是少算。
\item 會看到上一條分支留下的皇后。 只計數時，下一次寫入會蓋掉同一個
\texttt{board[row]}，答案碰巧仍對。
\item 正整數；基礎情況 $k = 0$ 回傳 $1$、$n = 0$ 或 $k < 0$ 回傳 $0$；
兩條分支是不拿與拿。
\end{enumerate}

\paragraph{第~\ref{sec:rec-when}~節}
\begin{enumerate}
\item 用 \texttt{for}。 這是線性掃描，深度 $10^7$ 會把堆疊用完。
\item 樹的高度跟資料形狀走，用遞迴。 串列長度在 $10^5$ 個節點時改迴圈
或顯式堆疊，系統堆疊裝不下。
\item 走格加記憶化；N 皇后把撤銷寫清楚。
\end{enumerate}

\chapter{檔案處理}
\label{ch:fileio}

程式一結束，變數裡的數字就消失了。 要把資料留下來，得寫進檔案；要把
上次留下的資料拿回來，得從檔案讀。 這一章教這件事。

三條線貫穿全章。 第一，\texttt{printf} 其實一直在把資料寫進某個目的
地，只是到目前為止那個目的地都是螢幕。 第二，同一個來源會記得你上次
讀到哪，所以連續讀下去會拿到下一個位元組。 第三，每個讀寫函式都會告
訴你成功或失敗，讀懂那個回傳值，就不會把失敗當成一筆資料。

先用已經會的 \texttt{printf} 與 \texttt{getchar}，把三條預設的線接去
檔案。 接著看這條線自己記得哪些狀態。 再開檔、關檔，然後依字元、行、
格式化欄位、二進位紀錄往下選工具。 最後用一支小程式把前面收在一起。

節末練習的答案收在章末。

\section{三條預設串流與重新導向}
\label{sec:fio-stream}

C 的標準函式庫不直接把資料寫進磁碟。 程式面對的是串流（stream）：一
條可以讀或寫的位元組序列，底下可能接鍵盤、螢幕、檔案、或另一支程式。
啟動時已經開好三條：

\begin{code}{c}{三條預先開啟的串流}
stdin    /* 標準輸入，預設接鍵盤 */
stdout   /* 標準輸出，預設接螢幕 */
stderr   /* 標準錯誤，預設也接螢幕 */
\end{code}

\texttt{printf} 印出一行，就是把那一行寫進 \texttt{stdout}。
\texttt{getchar()} 就是 \texttt{getc(stdin)}。 第~\ref{ch:character}~章
逐字元讀到檔案結束，用的就是 \texttt{stdin}；當時輸入來自鍵盤，或由
評測系統把檔案接到 \texttt{stdin}。 這一節把接線這件事講清楚，還不必
學任何新函式。

作業系統在啟動程式時可以改接線，不必改程式碼。 把檔案接到
\texttt{stdin}／\texttt{stdout}／\texttt{stderr} 的寫法叫做\textbf{重新導向}；
把一支程式的輸出接到另一支的輸入，叫做\textbf{管線}。

\begin{code}{text}{重新導向與管線}
./a.out < in.txt          從 in.txt 讀，當成 stdin
./a.out > out.txt         把 stdout 寫進 out.txt（檔案先清空）
./a.out >> out.txt        把 stdout 接到 out.txt 結尾
./a.out 2> err.txt        把 stderr 寫進 err.txt
./a.out > out.txt 2> err.txt
./prog1 | ./prog2         管線：prog1 的 stdout 接到 prog2 的 stdin
\end{code}

\texttt{stdout} 與 \texttt{stderr} 預設都接到螢幕，看起來像同一條。 一
旦用 \texttt{>} 把 \texttt{stdout} 接到檔案，兩條就分開了：錯誤訊息若
走 \texttt{printf}，會被一起寫進那個檔案；走 \texttt{fprintf(stderr,
...)}，仍留在螢幕上。

\begin{example}[label=ex:fio-hello-err]{c}{(hello-err.c) 正常訊息與錯誤分道}
#include <stdio.h>

int main(void) {
    printf("ok\n");
    fprintf(stderr, "fail\n");
    return 0;
}
\end{example}

在終端機直接執行，兩行都出現。 改成 \texttt{./a.out > out.txt
2> err.txt} 之後，\texttt{out.txt} 只剩 \texttt{ok}，\texttt{err.txt}
只剩 \texttt{fail}。

\begin{center}
\begin{iotermbox}{\$ ./a.out > out.txt 2> err.txt}
\ttfamily\raggedright\noindent
out.txt: ok\\
err.txt: fail
\end{iotermbox}
\end{center}

配合重新導向，第~\ref{ch:character}~章那個逐字元讀 \texttt{stdin} 的
迴圈已經能處理檔案：\texttt{./a.out < in.txt} 就等於從檔案讀。 下一
節才打開這三條串流背後共同的型別。

\keypoint{錯誤與診斷走 \texttt{stderr}，真正的結果走 \texttt{stdout}。
這樣用重新導向存結果時，錯誤才不會被一起吃掉。}

\begin{drills}[label=drills:fio-stream]
\item \texttt{./a.out > log.txt} 之後，\texttt{printf} 與
\texttt{fprintf(stderr, ...)} 各寫到哪裡？
\item \texttt{./prog1 | ./prog2} 這條管線裡，\texttt{prog2} 的
\texttt{stdin} 接到哪裡？
\item 為什麼診斷訊息要走 \texttt{stderr}，即使它預設也接到螢幕？
\item 補空：把 \texttt{out.txt} 接到結尾、檔案先不清空，運算子是
\texttt{?}。
\end{drills}

\section{FILE 的狀態}
\label{sec:fio-state}

\texttt{stdin}、\texttt{stdout}、\texttt{stderr} 的型別都是
\texttt{FILE *}。 同一個指標背後是一個有狀態的物件，至少要記得四塊。
這一節把四塊當場講完；後面開檔、讀寫，都指回這張圖。

\begin{center}
\begin{tikzpicture}[
    box/.style={
        draw=black!55, thick, rounded corners=2pt,
        minimum height=1.45cm, align=center,
        font=\sffamily\small, inner sep=3pt
    },
]
    \node[box, fill=memFill, draw=memStroke, minimum width=3.6cm]
        (buf) {串流緩衝區\\[-1pt]{\ttfamily\small [H][e][l][l][o]\ldots}};
    \node[box, fill=coreFill, draw=coreStroke, minimum width=2.1cm,
          right=0.22cm of buf]
        (pos) {位置\\[-1pt]{\ttfamily 5}};
    \node[box, fill=procFill, draw=procStroke, minimum width=2.0cm,
          right=0.22cm of pos]
        (eof) {EOF 旗標\\[-1pt]{\ttfamily 0}};
    \node[box, fill=optFill, draw=optStroke, minimum width=2.1cm,
          right=0.22cm of eof]
        (err) {error 旗標\\[-1pt]{\ttfamily 0}};
    \node[draw=black!35, dashed, rounded corners=8pt, inner sep=9pt,
          fit=(buf)(pos)(eof)(err)] (fileobj) {};
    \node[font=\sffamily\small, above=3pt of fileobj.north west,
          anchor=south west] {\texttt{FILE}};
    \node[box, fill=greenFill, draw=greenStroke, minimum width=4.4cm,
          below=1.15cm of fileobj]
        (disk) {磁碟上的位元組};
    \draw[diagarrow] ([xshift=-1.4cm]fileobj.south) -- (disk.north);
\end{tikzpicture}
\diagcaption{\texttt{FILE} 四塊狀態：串流緩衝區暫存尚未送出或尚未交
給程式的位元組，位置記住下一次讀寫從哪裡開始，EOF 與 error 是兩面獨
立的旗標。}
\label{fig:fio-file-state}
\end{center}

\subsection{串流緩衝區}

資料寫出去時，常常先躺在 \texttt{FILE} 內部的串流緩衝區，等到滿了、
遇到換行、或呼叫 \texttt{fflush}/\texttt{fclose} 才真正送出。 常見三
種策略：全緩衝（滿了才送）、列緩衝（遇到換行才送）、無緩衝（每次都
送）。 \texttt{stderr} 通常無緩衝，所以錯誤訊息會立刻出現。
\texttt{stdout} 接到終端機時多半是列緩衝；接到檔案時多半是全緩衝。

因此 \texttt{printf("loading...")} 後面沒有換行時，字可能暫時留在串
流緩衝區裡，螢幕上看不到。 除錯時會誤以為程式卡在更前面。 要立刻送
出，呼叫 \texttt{fflush(stdout)}，或印一個換行讓列緩衝觸發。

\begin{code}{c}{強制把 stdout 的串流緩衝區送出}
printf("loading...");
fflush(stdout);
/* 長時間計算 */
printf(" done\n");
\end{code}

\texttt{fflush(fp)} 對輸出串流的意思是：把串流緩衝區裡尚未寫出的資料
交給作業系統。 對輸入串流呼叫 \texttt{fflush} 是未定義行為。 網路上
常見的 \texttt{fflush(stdin)} 用來丟掉鍵盤殘餘，在標準 C 裡不合法；
要清掉行尾，請自己用 \texttt{getchar} 讀到換行為止。

\begin{code}{c}{反例：fflush(stdin) 是未定義行為}
fflush(stdin);    /* 不要這樣寫 */
\end{code}

\begin{code}{c}{丟掉目前這一行剩下的字元}
int ch;
while ((ch = getchar()) != EOF && ch != '\n') {
}
\end{code}

\texttt{fclose} 會先沖洗串流緩衝區再關閉。 寫入錯誤有時要到那一刻才
浮現；下一節的片語因此要求檢查 \texttt{fclose} 的回傳值。

\subsection{位置}

每次成功讀或寫，位置都會往前推，所以連續兩次讀到的是相鄰的兩個位元
組。 這一節只需要這件事。 之後有函式能把位置搬到別處，到隨機存取那
一節再登場。

\subsection{EOF 旗標：永遠晚一步}

\texttt{feof(fp)} 在 EOF 旗標立起時回傳非 $0$。 它\textbf{不}負責執
行讀取，只回答上一次讀寫已經留下的狀態。

EOF 旗標在一次讀取試圖跨越檔案尾時才立起。 讀成功最後一個字元之後，
位置已經指在檔案尾，旗標仍是 $0$。 再讀一次，讀取失敗，旗標才變成
$1$。 這就是 \texttt{while (!feof(fp))} 永遠晚一步的原因。

下面借用 \texttt{fgets}：它一次讀一行到行陣列裡，成功回傳該陣列，失
敗回傳 \texttt{NULL}（細節見第~\ref{sec:fio-line}~節）。 這裡只看迴
圈條件。 下面這支程式是反例，不要當片語抄。

\begin{example}[label=ex:fio-feof-wrong]{c}{(feof-wrong.c) 反例：用 feof 當迴圈條件}
#include <stdio.h>

int main(void) {
    char line[64];
    while (!feof(stdin)) {
        fgets(line, sizeof line, stdin);
        printf("%s", line);
    }
    return 0;
}
\end{example}

\iosample{
alpha \\
beta \\
}{
alpha \\
beta \\
beta \\
}

開檔時 EOF 旗標是 $0$，所以 \texttt{!feof(stdin)} 為真，進入迴圈。 第
一次 \texttt{fgets} 讀到 \texttt{alpha\textbackslash n}，印出
\texttt{alpha}；旗標仍是 $0$。 第二次讀到 \texttt{beta\textbackslash n}，
印出 \texttt{beta}；位置已經在檔案尾，旗標\textbf{仍是} $0$，因為這一次
讀取成功了。 迴圈條件再看一次 \texttt{feof}，仍為假，於是第三趟
\texttt{fgets} 失敗、回傳 \texttt{NULL}、這才立起 EOF 旗標。
\texttt{printf} 不管回傳值，仍把行陣列裡上次留下的
\texttt{beta\textbackslash n} 再印一次。

正確順序永遠是：先讀，看回傳值，需要時再用 \texttt{feof} 與
\texttt{ferror} 分類失敗。

\begin{example}[label=ex:fio-feof-right]{c}{(feof-right.c) 先讀，再依回傳值行事}
#include <stdio.h>

int main(void) {
    char line[64];
    while (fgets(line, sizeof line, stdin) != NULL) {
        printf("%s", line);
    }
    return 0;
}
\end{example}

\iosample{
alpha \\
beta \\
}{
alpha \\
beta \\
}

\subsection{error 旗標}

\texttt{ferror(fp)} 在 error 旗標立起時回傳非 $0$，表示讀寫過程發生過
I/O 錯誤。 \texttt{clearerr(fp)} 把 EOF 與 error 兩面旗標都清掉。
\texttt{fgets} 失敗時可能是檔案尾，也可能是讀取錯誤；迴圈結束後問
\texttt{ferror}，才能把兩種分開。 寫入錯誤常常遲到，資料先躺在串流緩
衝區裡，直到 \texttt{fflush} 或 \texttt{fclose} 才真正失敗。

\keypoint{先做讀取，用回傳值決定要不要處理這筆資料。
\texttt{feof}/\texttt{ferror} 只回答已經發生的狀態，不拿來當迴圈的進
入條件。 \texttt{printf} 沒有換行時，資料可能還在串流緩衝區。}

\begin{drills}[label=drills:fio-state]
\item 輸入只有一行 \texttt{hi\textbackslash n}，\texttt{feof-wrong.c}
的輸出是什麼？
\item 讀成功最後一個字元之後、尚未再讀一次之前，\texttt{feof(stdin)}
是真還是假？
\item 指出下面這行的問題：\texttt{fflush(stdin);}。
\item Parsons：把這四件事排成正確順序——看回傳值、讀取、需要時問
\texttt{feof}/\texttt{ferror}、處理這一筆資料。
\item \texttt{printf("hi")} 沒有換行，接著進入一段很長的迴圈。 終端
機上可能暫時看不到 \texttt{hi}，原因與圖~\ref{fig:fio-file-state} 的
哪一塊有關？
\end{drills}

\section{開啟與關閉}
\label{sec:fio-open}

要讀寫磁碟上的檔案，先向函式庫要一個 \texttt{FILE *}。
\texttt{fopen} 成功回傳非 \texttt{NULL} 的指標，失敗回傳
\texttt{NULL}。 失敗時函式庫把原因寫進 \texttt{errno}，那是一個用來記
錄失敗原因的全域變數；\texttt{perror} 把它翻成人看得懂的訊息。 用完
一定要 \texttt{fclose}：它會先把串流緩衝區沖出去，失敗時回傳
\texttt{EOF}。 寫入錯誤有時要到關閉那一刻才浮現，所以
\texttt{fclose} 的回傳值也要看。

這一節把完整檢查寫成片語。 從下一節起，範例聚焦在讀寫本身，錯誤檢查
依這條片語處理：開檔失敗仍寫出 \texttt{NULL} 判斷（沒檔就沒後面），
其餘 \texttt{fclose} 與逐次寫入的檢查不再抄進每一支程式。

\begin{code}{c}{打開、確認非 NULL、用完看 fclose}
FILE *fp = fopen(path, mode);
if (fp == NULL) {
    perror(path);
    return 1;
}
/* 讀寫 */
if (fclose(fp) != 0) {
    perror("fclose");
    return 1;
}
\end{code}

\texttt{mode} 由三條軸組成，不必背十二種字串。 第一軸決定開檔時對既
有內容做什麼：

\begin{center}
\begin{tabular}{|l|l|l|l|l|}
\hline
字母 & 檔案不存在 & 既有內容 & 起始位置 & 預設方向 \\
\hline
\texttt{r} & 失敗 & 保留 & 開頭 & 讀 \\
\texttt{w} & 建立空檔 & 截成空 & 開頭 & 寫 \\
\texttt{a} & 建立空檔 & 保留 & 結尾 & 寫 \\
\hline
\end{tabular}
\end{center}

第二軸是 \texttt{+}：加上之後，讀與寫都允許（\texttt{r+}、
\texttt{w+}、\texttt{a+}）。 第三軸是 \texttt{b}：二進位檔一律加上它。
理由見第~\ref{sec:fio-bin}~節；在會把換行轉換成兩個位元組的系統上，
漏掉 \texttt{b} 會改掉檔案裡的位元組。 三軸可以組合，例如
\texttt{"r+b"} 是既有檔案、讀寫、二進位、從開頭開始。

\keypoint{\texttt{w} 會把既有檔案清空。 只是想讀，用 \texttt{r}；想
接在後面寫，用 \texttt{a}。}

檔名從哪裡來？ 第~\ref{ch:function}~章寫過 \texttt{main(int argc,
char *argv[])}，當時標成進階、先不詳述。 現在它有真實用途：
\texttt{argv[0]} 是程式自己的名字，\texttt{argv[1]} 起才是我們傳入的
引數。 \texttt{argc} 是包含 \texttt{argv[0]} 在內的個數，所以一個檔
名對應 \texttt{argc == 2}。

\begin{example}[label=ex:fio-cat-file]{c}{(cat-f{}ile.c) 從命令列取檔名並抄到 stdout}
#include <stdio.h>

int main(int argc, char *argv[]) {
    if (argc != 2) {
        fprintf(stderr, "usage: %s file\n", argv[0]);
        return 1;
    }

    FILE *fp = fopen(argv[1], "r");
    if (fp == NULL) {
        perror(argv[1]);
        return 1;
    }

    int ch;
    while ((ch = fgetc(fp)) != EOF) {
        fputc(ch, stdout);
    }

    if (fclose(fp) != 0) {
        perror("fclose");
        return 1;
    }
    return 0;
}
\end{example}

\texttt{sample.txt} 裡面兩行 \texttt{hello} 與 \texttt{world} 時，
\texttt{./a.out sample.txt} 的輸出就是那兩行。 開檔失敗走
\texttt{perror}，訊息出現在 \texttt{stderr}，即使有人把
\texttt{stdout} 用重新導向接到別的檔案，錯誤仍看得到。

\stylepoint{檔名交給 \texttt{argv}，讓同一支程式處理不同檔案。 把檔
名寫死在原始碼裡，示範可以、正式程式容易卡住。}

\texttt{remove} 刪檔，\texttt{rename} 改名；兩者成功回傳 $0$。 後面
重寫刪除紀錄時會用到。

進階、此處先不展開：\texttt{freopen} 把已經開啟的串流改接到另一個檔
案；\texttt{tmpfile} 開一個二進位暫存檔，程式結束時自動刪掉。 知道名
字即可。

\begin{drills}[label=drills:fio-open]
\item \texttt{fopen("a.txt", "w")} 時 \texttt{a.txt} 已經有內容，開
檔成功後檔案裡剩下什麼？
\item \texttt{argc} 為 $1$ 時，\texttt{argv[1]} 可不可以當檔名？
\item \texttt{fwrite} 當時回傳成功，\texttt{fclose} 卻回傳
\texttt{EOF}。 資料一定已經安全躺在磁碟上了嗎？
\item 要讀一個既有的二進位檔、稍後也要寫回同一檔，\texttt{mode} 寫
哪一種？
\end{drills}

\section{讀寫工具箱}
\label{sec:fio-toolbox}

開好檔之後，真正的問題通常是：我要讀檔，該用哪個函式？ 下面兩張表是
後面四節的共同參照。 選工具看你要的單位；看回傳值才知道這一趟有沒有
成功。

\begin{center}
\begin{tabular}{|l|l|}
\hline
你要做的事 & 用哪個 \\
\hline
一個一個字元判斷（空白、換行、註解） & \texttt{fgetc} \\
一次處理整行（含換行） & \texttt{fgets} \\
欄位整齊的文字（空格或指定分隔） & \texttt{fscanf}，或 \texttt{fgets}+\texttt{sscanf} \\
定長的二進位紀錄 & \texttt{fread}/\texttt{fwrite} \\
\hline
\end{tabular}
\end{center}

欄位可能缺漏、含空白、或一行壞掉時，走 \texttt{fgets} 再
\texttt{sscanf}：失敗就跳過這一行，檔案位置不會卡在半行中間。

\begin{center}
\small
\begin{tabular}{|l|l|l|p{4.6cm}|}
\hline
函式 & 成功回傳 & 失敗回傳 & 怎麼分辨 \\
\hline
\texttt{fopen} & 非 \texttt{NULL} & \texttt{NULL} & \texttt{perror}（看 \texttt{errno}） \\
\texttt{fclose} & $0$ & \texttt{EOF} & 沖洗或關閉失敗 \\
\texttt{fgetc} & 該位元組（$0$--$255$） & \texttt{EOF} & \texttt{feof}／\texttt{ferror} \\
\texttt{fgets} & 行陣列指標 & \texttt{NULL} & 同上 \\
\texttt{fscanf} & 轉換成功的個數 & 較小的數或 \texttt{EOF} &
  不足：格式不符；\texttt{EOF}：檔尾或讀取錯誤 \\
\texttt{fread}/\texttt{fwrite} & 項目數 & 小於請求的數 &
  短讀再問 \texttt{feof}／\texttt{ferror} \\
\texttt{fseek} & $0$ & 非 $0$ & 無法定位 \\
\texttt{ftell} & 從開頭算的位移 & \texttt{-1L} & 不要把失敗值當檔案大小 \\
\hline
\end{tabular}
\end{center}

\texttt{fseek} 與 \texttt{ftell} 到隨機存取那一節才用；先知道失敗長
什麼樣子。 後面各節只指回這張表，不再各講一遍先看回傳值。

\keypoint{先問要處理的單位，再選函式。 選好之後，每一次呼叫都先看回
傳值，再用 \texttt{feof}/\texttt{ferror} 分類已經發生的失敗。}

\begin{drills}[label=drills:fio-toolbox]
\item 要印出檔案裡所有以 \texttt{\#} 開頭的行，選表裡哪一個函式？
\item 檔案是一筆筆固定大小的學生成績，選哪一對？
\item \texttt{fscanf} 回傳 $1$、你請求的是兩個整數。 這是檔案尾，還
是格式不符？
\item 補空：\texttt{fread} 請求 $10$ 筆、回傳 $7$，且
\texttt{ferror} 為假。 這叫短讀，原因是 \texttt{?}。
\end{drills}

\section{以字元為單位}
\label{sec:fio-char}

錯誤檢查依第~\ref{sec:fio-open}~節的片語。 回傳值見
第~\ref{sec:fio-toolbox}~節的表。

\texttt{fgetc(fp)} 從 \texttt{fp} 讀一個位元組，成功時回傳該位元組的
值（當成 \texttt{unsigned char} 再升成 \texttt{int}），失敗時回傳
\texttt{EOF}。 \texttt{fputc(ch, fp)} 寫一個位元組，成功回傳寫出的
值，失敗回傳 \texttt{EOF}。

標準迴圈把賦值寫在條件裡，並用括號保住先賦值再比較的順序：

\begin{code}{c}{逐字元讀到 EOF}
int ch;
while ((ch = fgetc(fp)) != EOF) {
    /* 使用 ch */
}
\end{code}

第~\ref{ch:character}~章已經強調：接 \texttt{getchar} 的變數要是
\texttt{int}。 \texttt{fgetc} 同一條規則。 若寫成 \texttt{char ch =
fgetc(fp)}，在 \texttt{char} 有號的平臺上會出現兩種災難。 第一種
第~\ref{ch:character}~章見過：真正的 \texttt{EOF} 被轉成 \texttt{char}
之後再升回 \texttt{int}，可能變成 $255$，迴圈比不到 $-1$，於是停不
下來。 第二種是二進位檔裡的位元組 $\mathtt{0xFF}$：它升成
\texttt{int} 本該是 $255$，先塞進有號 \texttt{char} 卻變成 $-1$，迴圈
會以為檔案結束，後面的資料全部沒讀到。

\begin{code}{c}{兩種都會錯：用 char 接 fgetc}
char ch;
while ((ch = fgetc(fp)) != EOF) {
    /* 有號 char 時，EOF 比不到，或 0xFF 被當成 EOF */
}
\end{code}

\begin{example}[label=ex:fio-mini-cat]{c}{(mini-cat.c) 把 stdin 抄到 stdout}
#include <stdio.h>

int main(void) {
    int ch;
    while ((ch = fgetc(stdin)) != EOF) {
        fputc(ch, stdout);
    }
    return 0;
}
\end{example}

\iosample{
ab \\
}{
ab \\
}

這支程式沒有自己 \texttt{fopen}。 配合第~\ref{sec:fio-stream}~節的重
新導向，\texttt{./a.out < sample.txt} 就等於把檔案抄到螢幕。
\texttt{cat-file.c} 只是把同一個迴圈裡的 \texttt{stdin} 換成自己打開
的 \texttt{FILE *}。

統計字元、單字、行數，用的也是這道迴圈。 單字的邊界用
\texttt{isspace} 判斷：從空白進到非空白，就算一個新單字開始。

\begin{example}[label=ex:fio-count]{c}{(count-cwl.c) 列、單字、字元}
#include <stdio.h>
#include <ctype.h>

int main(void) {
    int ch;
    int chars = 0;
    int words = 0;
    int lines = 0;
    int in_word = 0;

    while ((ch = fgetc(stdin)) != EOF) {
        chars++;
        if (ch == '\n') {
            lines++;
        }
        if (isspace(ch)) {
            in_word = 0;
        } else if (!in_word) {
            in_word = 1;
            words++;
        }
    }

    printf("%d %d %d\n", lines, words, chars);
    return 0;
}
\end{example}

\iosample{
one two \\
three \\
}{
2 3 14 \\
}

輸入是 \texttt{one two\textbackslash nthree\textbackslash n}，一共
$14$ 個字元（含兩個換行）、三個單字、兩行。 最後一行若沒有換行，
\texttt{lines} 會少算 $1$，因為我們只在讀到 \texttt{\textbackslash n}
時才加行數。

進階、此處先不展開：\texttt{ungetc(ch, fp)} 把剛剛讀到的字元塞回串
流，下一次 \texttt{fgetc} 會再讀到它。 標準只保證可以退回一格。

\keypoint{\texttt{fgetc} 的回傳值用 \texttt{int} 接。 這條規則在文字
檔避免停不下來，在二進位檔避免把 $\mathtt{0xFF}$ 誤判成檔案結束。}

\begin{drills}[label=drills:fio-char]
\item \texttt{char ch;} 搭配 \texttt{while ((ch = fgetc(fp)) !=
EOF)}，在有號 \texttt{char} 的平臺上讀二進位檔，檔案中間有一個位元組
$\mathtt{0xFF}$。 迴圈在哪裡停？
\item \texttt{count-cwl.c} 讀 \texttt{a\textbackslash nb}（最後沒有換
行），印出的三個數字是什麼？
\item 把 \texttt{mini-cat.c} 的 \texttt{fgetc(stdin)} 改成
\texttt{getchar()}，行為差在哪裡？
\item 從零寫：從 \texttt{argv[1]} 打開檔案，用 \texttt{fgetc} 抄到
\texttt{stdout}（錯誤檢查依第~\ref{sec:fio-open}~節）。
\end{drills}

\section{以行為單位}
\label{sec:fio-line}

錯誤檢查依第~\ref{sec:fio-open}~節的片語。

\texttt{fgets(buf, size, fp)} 從 \texttt{fp} 讀一行，最多放進
\texttt{size - 1} 個字元，再補 \texttt{\textbackslash 0}。 成功回傳
\texttt{buf}（也就是那塊行陣列），讀到檔案尾且這一趟一個字元都沒讀到，
回傳 \texttt{NULL}。 換行字元若在這一趟的額度內，會被保留在行陣列裡。

\begin{code}{c}{逐行讀取}
char line[128];
while (fgets(line, sizeof line, fp) != NULL) {
    /* line 可能以 \n 結尾，也可能因太長被截斷 */
}
\end{code}

\texttt{size} 含結尾的空字元，所以 \texttt{char line[8]} 一次最多讀
$7$ 個字元。 若原始行是 \texttt{abcdefg\textbackslash n}（八個字元），
第一次 \texttt{fgets} 得到 \texttt{abcdefg}（沒有換行），第二次才得到
\texttt{\textbackslash n}。 判斷一行有沒有被截斷：成功讀到之後看行陣
列裡有沒有 \texttt{\textbackslash n}。

去掉行尾換行，用 \texttt{strcspn} 找第一個換行的位置，再寫入空字元：

\begin{code}{c}{去掉 fgets 留下的換行}
line[strcspn(line, "\n")] = '\0';
\end{code}

這行在沒有換行時也安全：\texttt{strcspn} 會走到字串結尾，等於沒改內
容。

\texttt{fputs(s, fp)} 把字串寫出，自己\textbf{不}補換行。 相對地，
\texttt{puts(s)} 寫到 \texttt{stdout} 而且\textbf{會}再寫一個換行。
從 \texttt{fgets} 讀來的行若仍帶著 \texttt{\textbackslash n}，用
\texttt{fputs} 抄出去就會維持原來的換行。

\texttt{gets} 沒有長度參數，行陣列要多大全憑運氣。 C11 已把它從標準
裡刪除；任何教材或舊程式若還在呼叫它，都改成 \texttt{fgets}。

\begin{example}[label=ex:fio-keep-hash]{c}{(keep-hash.c) 只留下以 \# 開頭的行}
#include <stdio.h>

int main(void) {
    char line[128];
    while (fgets(line, sizeof line, stdin) != NULL) {
        if (line[0] == '#') {
            fputs(line, stdout);
        }
    }
    return 0;
}
\end{example}

\iosample{
keep \\
\#note \\
skip \\
\#also \\
}{
\#note \\
\#also \\
}

一個跟換行有關的陷阱：使用者輸入 \texttt{quit} 結束。
\texttt{fgets} 讀到的是 \texttt{quit\textbackslash n}，
\texttt{strcmp(line, "quit")} 會判定兩字串不同。 先剝掉換行再比。

\keypoint{\texttt{fgets} 成功時回傳行陣列指標；用回傳值當迴圈條件，
不要用 \texttt{feof} 當迴圈條件。}

\begin{drills}[label=drills:fio-line]
\item \texttt{fgets} 之後立刻 \texttt{strcmp(line, "quit") == 0}，使
用者輸入 \texttt{quit} 並按 Enter，比得中嗎？
\item \texttt{char buf[4];} 讀一行 \texttt{abcd\textbackslash n}，第
一次 \texttt{fgets(buf, 4, fp)} 的內容是什麼（含有沒有換行）？
\item 補空：\texttt{line[strcspn(line, ?)] = '\textbackslash 0';}
問號填哪個字串？
\item \texttt{puts("ok")} 與 \texttt{fputs("ok", stdout)} 印到螢幕上
各是幾個字元？
\end{drills}

\section{格式化讀寫}
\label{sec:fio-fmt}

錯誤檢查依第~\ref{sec:fio-open}~節的片語。 回傳值見
第~\ref{sec:fio-toolbox}~節的表。

\texttt{fprintf(fp, fmt, ...)} 與 \texttt{printf} 相同，只是目的地改
成 \texttt{fp}。 \texttt{printf(...)} 就是
\texttt{fprintf(stdout, ...)}。 這件事有實際價值：把輸出目標做成參數，
同一支函式可以印到螢幕、寫進檔案、或送到診斷。

\begin{example}[label=ex:fio-report]{c}{(report.c) 同一個函式，三個目的地}
#include <stdio.h>

typedef struct {
    int id;
    char name[16];
    double gpa;
} Record;

static void report(FILE *out, const Record *r) {
    fprintf(out, "%d %s %.2f\n", r->id, r->name, r->gpa);
}

int main(void) {
    Record r = {1, "Alice", 3.80};
    report(stdout, &r);
    report(stderr, &r);
    return 0;
}
\end{example}

\iosample{
}{
1 Alice 3.80 \\
}

\texttt{stderr} 上也會出現同一行。 後面列出成績檔時，把 \texttt{out}
換成自己打開的 \texttt{FILE *} 即可，不必再寫一份 \texttt{printf}。

\texttt{fscanf} 的回傳值是成功完成的轉換個數。 三種結束方式見工具箱
那張表：個數相符、\texttt{EOF}、以及比較小的非負整數（格式不符）。

\begin{example}{c}{(fscanf-three.c) 用回傳值分開三種結局}
#include <stdio.h>

int main(void) {
    int a, b;
    int n = scanf("%d %d", &a, &b);
    printf("n=%d\n", n);
    if (n == 2) {
        printf("%d %d\n", a, b);
    } else if (n == EOF) {
        printf("eof or read error\n");
        if (feof(stdin)) {
            printf("feof set\n");
        }
    } else {
        printf("mismatch, converted %d\n", n);
    }
    return 0;
}
\end{example}

輸入 \texttt{3 4} 時 \texttt{n=2}。 輸入 \texttt{3 x} 時 \texttt{n=1}，
走格式不符。 輸入為空時 \texttt{n=-1}（也就是 \texttt{EOF}）。

\texttt{\%s} 會先跳過前導空白，再讀到下一個空白為止。 scanset
\texttt{\%[\^{},]} 是另一種轉換指定：方括號裡是\textbf{字元集合}，
\texttt{\^{}} 表示讀到集合裡的字元就停。 所以 \texttt{\%[\^{},]} 的
意思是一直讀、直到碰到逗號。 它跟 \texttt{\%s} 一樣要有寬度上限，否
則輸入比陣列長時會寫穿。 \texttt{char name[50]} 對應
\texttt{\%49[\^{},]}，留下一格給空字元。

直接用 \texttt{fscanf} 拆 CSV 很脆弱：某行缺欄，回傳值小於預期，檔案
位置停在失敗那一格。 比較穩的路徑是先 \texttt{fgets} 拿整行，再
\texttt{sscanf} 解析；失敗時丟掉這一行即可。

\begin{example}[label=ex:fio-csv]{c}{(csv-parse.c) 逐行解析 CSV，壞行跳過}
#include <stdio.h>
#include <string.h>

int main(void) {
    char line[128];
    int id;
    char name[50];
    double gpa;

    while (fgets(line, sizeof line, stdin) != NULL) {
        line[strcspn(line, "\n")] = '\0';
        if (sscanf(line, "%d,%49[^,],%lf", &id, name, &gpa) == 3) {
            printf("%d %s %.2f\n", id, name, gpa);
        } else {
            fprintf(stderr, "skip: %s\n", line);
        }
    }
    return 0;
}
\end{example}

\iosample{
1,Alice,3.80 \\
2,Bob,3.50 \\
badline \\
3,Carol,3.90 \\
}{
1 Alice 3.80 \\
2 Bob 3.50 \\
3 Carol 3.90 \\
}

寫入時若目標是字元陣列，用 \texttt{snprintf}：它接受陣列大小，保證寫
入空字元。

\stylepoint{欄位可能缺漏時，先 \texttt{fgets} 再 \texttt{sscanf}，失
敗就跳過或報錯。}

\begin{drills}[label=drills:fio-fmt]
\item \texttt{fscanf(fp, "\%d \%d", \&a, \&b)} 讀到 \texttt{3 x}，回
傳值是多少？
\item \texttt{char name[8];} 使用 \texttt{\%s} 而不寫寬度，輸入一個
$30$ 字元的名字，會發生什麼事？
\item 為什麼 \texttt{csv-parse.c} 的 scanset 要寫寬度 $49$？
\item 改寫：把只用 \texttt{fscanf}、迴圈條件寫 \texttt{== 3} 的 CSV
讀取，改成 \texttt{fgets}+\texttt{sscanf}。 遇到 \texttt{badline} 之
後還要讀得到下一行。
\end{drills}

\section{二進位讀寫}
\label{sec:fio-bin}

錯誤檢查依第~\ref{sec:fio-open}~節的片語。

同一筆數字，用文字存與用二進位存，磁碟上的位元組不一樣。 把
$32767$ 寫成十進位數字字串，檔案裡是五個 ASCII 字元。 若改成一個
$16$ 位元整數 \texttt{int16\_t}，檔案裡是兩個位元組。

\begin{example}[label=ex:fio-text-bin]{c}{(write-text-bin.c) 同一個數字的兩種存法}
#include <stdio.h>
#include <stdint.h>

int main(void) {
    FILE *txt = fopen("num.txt", "w");
    if (txt == NULL) {
        perror("num.txt");
        return 1;
    }
    fprintf(txt, "32767");
    fclose(txt);

    FILE *bin = fopen("num.bin", "wb");
    if (bin == NULL) {
        perror("num.bin");
        return 1;
    }
    int16_t n = 32767;
    fwrite(&n, sizeof n, 1, bin);
    fclose(bin);
    return 0;
}
\end{example}

用十六進位把每個位元組印出來（本機指令是 \texttt{xxd}，讀者不必會用，
後面每個二進位範例都會把結果貼上來）：

\begin{center}
\begin{verbatim}
$ xxd num.txt
00000000: 3332 3736 37                             32767

$ xxd num.bin
00000000: ff7f                                     ..
\end{verbatim}
\end{center}

\texttt{33}、\texttt{32}、\texttt{37}、\texttt{36}、\texttt{37} 就是
字元 \texttt{'3'} 到 \texttt{'7'}。 \texttt{num.bin} 只有
\texttt{ff 7f}。 多位元組整數在記憶體裡的排列順序叫做端序
（endianness）。 這台機器是小端序：低位元組排在前面，所以
$32767 = \mathtt{0x7FFF}$ 存成 \texttt{FF} 再 \texttt{7F}。 大端序機
器會存成 \texttt{7F FF}。 文字檔 \texttt{"32767"} 的字元順序與機器無
關；二進位檔換機器讀可能得到不同的數值。 這也是 mode 要加 \texttt{b}
的理由之一：文字模式可能改寫換行，二進位檔的每一個位元組都必須原樣
進出。

\texttt{fread} 與 \texttt{fwrite} 以記憶體裡的位元組為單位搬資料。 呼
叫形如 \texttt{fread(ptr, size, nmemb, fp)}，回傳\textbf{成功的項目數}，
型別是 \texttt{size\_t}。 檔案剩下的位元組不夠湊滿 \texttt{nmemb} 個
項目時，回傳值小於請求，叫做短讀。 再問 \texttt{feof} 與
\texttt{ferror} 分類。

\begin{code}{c}{fread 看回傳的項目數}
size_t n = fread(buf, sizeof buf[0], 5, fp);
if (n < 5) {
    if (ferror(fp)) {
        /* 讀取錯誤 */
    } else if (feof(fp)) {
        /* 檔案尾，前 n 個項目仍可用 */
    }
}
\end{code}

\begin{example}{c}{(short-read.c) 檔案只有兩筆，請求五筆}
#include <stdio.h>
#include <stdint.h>

int main(void) {
    int16_t vals[2] = {1, 2};
    FILE *fp = fopen("short.bin", "wb");
    if (fp == NULL) {
        perror("short.bin");
        return 1;
    }
    fwrite(vals, sizeof vals[0], 2, fp);
    fclose(fp);

    int16_t buf[5];
    fp = fopen("short.bin", "rb");
    if (fp == NULL) {
        perror("short.bin");
        return 1;
    }
    size_t n = fread(buf, sizeof buf[0], 5, fp);
    printf("got %zu items\n", n);
    if (n < 5 && feof(fp)) {
        printf("feof\n");
    }
    fclose(fp);
    return 0;
}
\end{example}

\iosample{
}{
got 2 items \\
feof \\
}

三個結構相關的陷阱，前兩個第~\ref{ch:structure}~章已經鋪過。

第一，填補。 \texttt{sizeof} 量到的是含填補的大小，\texttt{fwrite}
會把填補一起寫進檔案。

\begin{example}{c}{(padding-show.c) 五個位元組的欄位，寫出八個位元組}
#include <stdio.h>

struct Spread {
    char c;
    int n;
};

int main(void) {
    struct Spread s = {'A', 1};
    printf("sizeof=%zu field_sum=%zu\n",
           sizeof s, sizeof s.c + sizeof s.n);

    FILE *fp = fopen("pad.bin", "wb");
    if (fp == NULL) {
        perror("pad.bin");
        return 1;
    }
    fwrite(&s, sizeof s, 1, fp);
    fclose(fp);
    return 0;
}
\end{example}

\iosample{
}{
sizeof=8 field\_sum=5 \\
}

\begin{center}
\begin{verbatim}
$ xxd pad.bin
00000000: 4100 0000 0100 0000                      A.......
\end{verbatim}
\end{center}

\texttt{'A'} 是 \texttt{41}，後面三個 \texttt{00} 是填補，然後才是
整數 $1$。 宣告時用 \texttt{\{'A', 1\}} 初始化，填補也會被設成 $0$。

第二，端序。 同一個 \texttt{uint32\_t} 值 $\mathtt{0x01020304}$ 在這
台小端序機器上寫成：

\begin{center}
\begin{verbatim}
$ xxd endian.bin
00000000: 0403 0201                                ....
\end{verbatim}
\end{center}

\begin{example}{c}{(dump-endian.c) 同一個 uint32 在檔案裡的位元組}
#include <stdio.h>
#include <stdint.h>

int main(void) {
    FILE *fp = fopen("endian.bin", "wb");
    if (fp == NULL) {
        perror("endian.bin");
        return 1;
    }
    uint32_t n = 0x01020304u;
    fwrite(&n, sizeof n, 1, fp);
    fclose(fp);
    return 0;
}
\end{example}

第三，含指標的結構不能直接寫。 \texttt{char *name} 在結構裡只佔一個
指標那麼大；\texttt{fwrite} 寫出的是位址，字串內容不會出現在檔案裡。
下一次執行這個程式再讀回來，那個位址毫無意義。 姓名要能跟結構一起存
盤，請改成 \texttt{char name[16]} 這種內嵌陣列。

\begin{example}[label=ex:fio-recs]{c}{(dump-records.c) 內嵌陣列的定長紀錄}
#include <stdio.h>

typedef struct {
    int id;
    char name[16];
    double gpa;
} Record;

int main(void) {
    Record recs[2] = {
        {1, "Alice", 3.80},
        {2, "Bob", 3.50},
    };

    FILE *fp = fopen("recs.bin", "wb");
    if (fp == NULL) {
        perror("recs.bin");
        return 1;
    }
    fwrite(recs, sizeof recs[0], 2, fp);
    fclose(fp);
    printf("sizeof(Record)=%zu\n", sizeof(Record));
    return 0;
}
\end{example}

\iosample{
}{
sizeof(Record)=32 \\
}

\begin{center}
\begin{verbatim}
$ xxd recs.bin
00000000: 0100 0000 416c 6963 6500 0000 0000 0000  ....Alice.......
00000010: 0000 0000 0000 0000 6666 6666 6666 0e40  ........ffffff.@
00000020: 0200 0000 426f 6200 0000 0000 0000 0000  ....Bob.........
00000030: 0000 0000 0000 0000 0000 0000 0000 0c40  ...............@
\end{verbatim}
\end{center}

宣告時用 \texttt{\{1, "Alice", 3.80\}}，其餘位元組都是 $0$。 賦值那
邊不能寫同樣的大括號。

何時用文字、何時用二進位：需要人眼看、需要跨機器、欄位長度不固定，
走文字。 需要定長紀錄、檔案很大、只在同一種機器上讀寫，走二進位。

\keypoint{\texttt{fread}/\texttt{fwrite} 回傳的是項目數。 結構則連填
補、端序、指標欄位一起寫進檔案。}

\begin{drills}[label=drills:fio-bin]
\item \texttt{fread(buf, sizeof(int), 10, fp)} 回傳 $7$，且
\texttt{ferror(fp)} 為假。 陣列前幾格有資料？
\item 把含 \texttt{char *name} 的結構 \texttt{fwrite} 出去，檔案裡有
沒有字串的字元？
\item 同一份 \texttt{pad.bin} 拿到大端序機器上，用同一個
\texttt{struct Spread} 讀回來，\texttt{n} 還是 $1$ 嗎？
\item 文字檔存 \texttt{32767} 與二進位 \texttt{int16\_t} 存同一個數，
哪一種換機器讀比較不容易因端序出錯？
\end{drills}

\section{位置與隨機存取}
\label{sec:fio-seek}

錯誤檢查依第~\ref{sec:fio-open}~節的片語。

定長紀錄把檔案變成陣列：第 $i$ 筆從位元組
\texttt{i * sizeof(Record)} 開始。 \texttt{fseek(fp, offset, whence)}
移動位置，成功回傳 $0$。 \texttt{whence} 為 \texttt{SEEK\_SET}（從檔
案開頭）、\texttt{SEEK\_CUR}（從目前位置）、\texttt{SEEK\_END}（從檔
案尾）。 \texttt{ftell(fp)} 回傳從開頭算起的目前位置；失敗時回傳
\texttt{-1L}，所以結果要接在 \texttt{long} 裡再跟 \texttt{-1L} 比。
\texttt{rewind(fp)} 相當於把位置設回開頭並清錯誤旗標。

\begin{center}
\begin{tikzpicture}
    \node[cell, minimum width=1.6cm] (r0) {R0};
    \node[cell, minimum width=1.6cm, right=0cm of r0] (r1) {R1};
    \node[hlcell, minimum width=1.6cm, right=0cm of r1] (r2) {R2};
    \node[cell, minimum width=1.6cm, right=0cm of r2] (r3) {R3};
    \node[idxlbl, below=3pt of r0] {\texttt{0}};
    \node[idxlbl, below=3pt of r1] {\texttt{1}};
    \node[idxlbl, below=3pt of r2] {\texttt{2}};
    \node[idxlbl, below=3pt of r3] {\texttt{3}};
    \draw[diagarrow] ([yshift=0.55cm]r2.north) -- (r2.north);
    \node[font=\sffamily\scriptsize, above=8pt of r2]
        {\texttt{fseek(fp, 2 * sizeof(Record), SEEK\_SET)}};
\end{tikzpicture}
\diagcaption{定長紀錄在檔案裡連續排列，第 $i$ 筆的位移是
$i \times \texttt{sizeof(Record)}$。}
\label{fig:fio-seek}
\end{center}

\begin{code}{c}{跳到第 i 筆再讀}
long offset = (long)i * (long)sizeof(Record);
if (fseek(fp, offset, SEEK_SET) != 0) {
    perror("fseek");
}
Record r;
if (fread(&r, sizeof r, 1, fp) != 1) {
    /* 沒有這一筆，或讀取錯誤 */
}
\end{code}

更新某一筆時，通常用 \texttt{"r+b"} 打開。 標準規定，從讀切到寫、或
從寫切到讀，中間必須插入 \texttt{fflush} 或一次定位。 理由是：實作
常讓讀寫共用同一塊串流緩衝區，切方向前得先讓它同步，否則串流緩衝區
裡的內容與檔案位置會對不上。 讀完第 $i$ 筆之後位置已經指到第 $i+1$ 筆開
頭，若直接 \texttt{fwrite} 會寫錯地方，而且未先定位本身就違反這條規
定。 正確做法是改完欄位後，再 \texttt{fseek} 回同一筆的開頭，然後寫。

假設 \texttt{grades.bin} 已經有三筆，對應第~\ref{sec:fio-bin}~節那種
\texttt{Record}（Alice、Bob、Carol）。 下面只做原地更新。

\begin{example}[label=ex:fio-update]{c}{(seek-update.c) 原地改第 1 筆的 GPA}
#include <stdio.h>

typedef struct {
    int id;
    char name[16];
    double gpa;
} Record;

int main(void) {
    FILE *fp = fopen("grades.bin", "r+b");
    if (fp == NULL) {
        perror("grades.bin");
        return 1;
    }

    Record r;
    long offset = (long)1 * (long)sizeof(Record);
    fseek(fp, offset, SEEK_SET);
    if (fread(&r, sizeof r, 1, fp) != 1) {
        fclose(fp);
        return 1;
    }
    printf("before %d %s %.2f\n", r.id, r.name, r.gpa);
    r.gpa = 3.95;
    fseek(fp, offset, SEEK_SET);
    fwrite(&r, sizeof r, 1, fp);
    fclose(fp);
    return 0;
}
\end{example}

\iosample{
}{
before 2 Bob 3.50 \\
}

索引從 $0$ 起算，所以第 $1$ 筆是 Bob。 讀與寫之間那一次
\texttt{fseek} 同時滿足兩件事：位置回到 Bob 的開頭，以及標準要求的同
步。 文字模式下 \texttt{fseek} 的可攜性限制見附錄。

\keypoint{定長紀錄用 \texttt{index * sizeof} 當位移。 \texttt{r+b} 讀
寫切換時，中間必須 \texttt{fflush} 或呼叫定位函式，因為讀寫可能共用
同一塊串流緩衝區。}

\begin{drills}[label=drills:fio-seek]
\item 第 $0$ 筆、第 $3$ 筆的位移各是多少（以 \texttt{sizeof(Record)}
為單位寫出來）？
\item \texttt{r+b} 下 \texttt{fread} 成功後立刻 \texttt{fwrite} 同一
筆、中間沒有 \texttt{fseek}。 錯在哪兩件事？
\item \texttt{ftell} 回傳 \texttt{-1L} 時，能不能拿去當檔案大小？
\item 從零寫：打開既有的定長紀錄檔，印出最後一筆（先用
\texttt{fseek} 到結尾再往回一筆，或先 \texttt{ftell} 算筆數）。
\end{drills}

\section{綜合：定長紀錄與 CSV}
\label{sec:fio-combine}

錯誤檢查依第~\ref{sec:fio-open}~節的片語。

把前面的工具收成一支小程式：檔案裡放定長的 \texttt{Record}，用
\texttt{argv[1]} 當子命令。 \texttt{create} 建三筆（當成已知前置），
\texttt{list} 用第~\ref{sec:fio-fmt}~節的 \texttt{report} 依序印出，
\texttt{update} 原地改第 $n$ 筆的 GPA，\texttt{delete} 用重寫的方式
刪掉第 $n$ 筆。

\begin{example}[label=ex:fio-grades]{c}{(grades.c) 建檔、列出、原地更新、重寫刪除}
#include <stdio.h>
#include <stdlib.h>
#include <string.h>

typedef struct {
    int id;
    char name[16];
    double gpa;
} Record;

static void report(FILE *out, const Record *r) {
    fprintf(out, "%d %s %.2f\n", r->id, r->name, r->gpa);
}

static int create_file(const char *path) {
    Record recs[3] = {
        {1, "Alice", 3.80},
        {2, "Bob", 3.50},
        {3, "Carol", 3.90},
    };
    FILE *fp = fopen(path, "wb");
    if (fp == NULL) {
        perror(path);
        return 1;
    }
    fwrite(recs, sizeof recs[0], 3, fp);
    fclose(fp);
    return 0;
}

static int list_file(const char *path) {
    FILE *fp = fopen(path, "rb");
    if (fp == NULL) {
        perror(path);
        return 1;
    }
    Record r;
    while (fread(&r, sizeof r, 1, fp) == 1) {
        report(stdout, &r);
    }
    fclose(fp);
    return 0;
}

static int update_file(const char *path, int index, double gpa) {
    FILE *fp = fopen(path, "r+b");
    if (fp == NULL) {
        perror(path);
        return 1;
    }
    Record r;
    long offset = (long)index * (long)sizeof(Record);
    fseek(fp, offset, SEEK_SET);
    if (fread(&r, sizeof r, 1, fp) != 1) {
        fprintf(stderr, "no record at %d\n", index);
        fclose(fp);
        return 1;
    }
    r.gpa = gpa;
    fseek(fp, offset, SEEK_SET);
    fwrite(&r, sizeof r, 1, fp);
    fclose(fp);
    return 0;
}

static int delete_file(const char *path, int index) {
    FILE *in = fopen(path, "rb");
    if (in == NULL) {
        perror(path);
        return 1;
    }
    FILE *out = fopen("grades.tmp", "wb");
    if (out == NULL) {
        perror("grades.tmp");
        fclose(in);
        return 1;
    }
    Record r;
    int i = 0;
    int kept = 0;
    while (fread(&r, sizeof r, 1, in) == 1) {
        if (i != index) {
            fwrite(&r, sizeof r, 1, out);
            kept++;
        }
        i++;
    }
    int cin = fclose(in);
    int cout = fclose(out);
    if (cin != 0 || cout != 0) {
        perror("fclose");
        return 1;
    }
    remove(path);
    rename("grades.tmp", path);
    printf("kept %d\n", kept);
    return 0;
}

int main(int argc, char *argv[]) {
    const char *path = "grades.bin";
    if (argc < 2) {
        fprintf(stderr, "usage: %s create|list|update|delete ...\n",
                argv[0]);
        return 1;
    }
    if (strcmp(argv[1], "create") == 0) {
        return create_file(path);
    }
    if (strcmp(argv[1], "list") == 0) {
        return list_file(path);
    }
    if (strcmp(argv[1], "update") == 0) {
        if (argc != 4) {
            fprintf(stderr, "usage: %s update index gpa\n", argv[0]);
            return 1;
        }
        return update_file(path, atoi(argv[2]), atof(argv[3]));
    }
    if (strcmp(argv[1], "delete") == 0) {
        if (argc != 3) {
            fprintf(stderr, "usage: %s delete index\n", argv[0]);
            return 1;
        }
        return delete_file(path, atoi(argv[2]));
    }
    fprintf(stderr, "unknown command\n");
    return 1;
}
\end{example}

\texttt{fclose(in)} 與 \texttt{fclose(out)} 必須各呼叫一次再合併判斷。
若寫成兩個呼叫用 \texttt{||} 串在條件裡，左邊失敗時右邊不會執行，暫存
檔沒沖洗就去 \texttt{rename}。

文字與二進位之間的橋：從 CSV 讀進來，寫成同一種 \texttt{Record}。 每
一筆用 \texttt{Record r = \{0\}} 先清成零，避免名字欄留下上一筆的殘
餘字元。

\begin{example}[label=ex:fio-csv2bin]{c}{(csv-to-bin.c) CSV 轉定長二進位再讀回}
#include <stdio.h>
#include <string.h>

typedef struct {
    int id;
    char name[16];
    double gpa;
} Record;

int main(void) {
    char line[128];
    Record recs[8];
    int n = 0;

    while (fgets(line, sizeof line, stdin) != NULL && n < 8) {
        line[strcspn(line, "\n")] = '\0';
        Record r = {0};
        if (sscanf(line, "%d,%15[^,],%lf", &r.id, r.name, &r.gpa) != 3) {
            fprintf(stderr, "skip: %s\n", line);
            continue;
        }
        recs[n] = r;
        n++;
    }

    FILE *fp = fopen("fromcsv.bin", "wb");
    if (fp == NULL) {
        perror("fromcsv.bin");
        return 1;
    }
    if (n > 0) {
        fwrite(recs, sizeof recs[0], (size_t)n, fp);
    }
    fclose(fp);

    fp = fopen("fromcsv.bin", "rb");
    if (fp == NULL) {
        perror("fromcsv.bin");
        return 1;
    }
    Record r;
    while (fread(&r, sizeof r, 1, fp) == 1) {
        printf("%d %s %.2f\n", r.id, r.name, r.gpa);
    }
    fclose(fp);
    return 0;
}
\end{example}

\iosample{
1,Alice,3.80 \\
2,Bob,3.50 \\
bad \\
3,Carol,3.90 \\
}{
1 Alice 3.80 \\
2 Bob 3.50 \\
3 Carol 3.90 \\
}

\stylepoint{子命令放在 \texttt{argv}，比在程式裡做互動選單更利於重跑
與測試。 刪除定長紀錄用重寫，更新才用 \texttt{r+b} 原地改。}

\begin{drills}[label=drills:fio-combine]
\item 從零寫：印出 \texttt{stdin} 的最後一行（相當於 \texttt{tail -n
1} 的文字檔版本）。 行長不超過 $1023$ 字元。
\item \texttt{csv-to-bin.c} 若把 \texttt{Record r = \{0\}} 改成未初
始化的 \texttt{Record r;}，Alice 之後讀 Bob，名字欄可能出現什麼殘餘？
\item \texttt{list} 若改成 \texttt{while (!feof(fp)) \{ fread...
report...\}}，最後一筆會印幾次？
\item 改寫：迴圈條件若寫 \texttt{fscanf(...) == 3}，改成
\texttt{fgets} 加 \texttt{sscanf}。
\end{drills}

\section{錯誤清單}
\label{sec:fio-checklist}

下面每一條都對應前面某一節。

\begin{itemize}
\item 診斷走 \texttt{stderr}，結果走 \texttt{stdout}。
接線叫重新導向；\texttt{|} 叫管線。
（第~\ref{sec:fio-stream}~節）
\item 先讀、看回傳值，再用 \texttt{feof}/\texttt{ferror} 分類。
\texttt{while (!feof(fp))} 會讓最後一行重複。
（第~\ref{sec:fio-state}~節）
\item \texttt{printf} 無換行可能留在串流緩衝區。
\texttt{fflush(stdout)} 合法，\texttt{fflush(stdin)} 是未定義行為。
（第~\ref{sec:fio-state}~節）
\item \texttt{fopen} 失敗回傳 \texttt{NULL}，先 \texttt{perror} 再離
開；\texttt{fclose} 的回傳值也要看。 \texttt{w} 會清空既有檔案。 檔
名從 \texttt{argv} 取。
（第~\ref{sec:fio-open}~節）
\item 先問要處理的單位，再選函式。
（第~\ref{sec:fio-toolbox}~節）
\item \texttt{fgetc}/\texttt{getchar} 的回傳值用 \texttt{int} 接。
（第~\ref{sec:fio-char}~節）
\item 用 \texttt{fgets} 讀行時，換行會留在行陣列裡。 比對結束字串前
先剝換行。 標準已移除 \texttt{gets}。
（第~\ref{sec:fio-line}~節）
\item \texttt{fscanf} 看成功轉換數；scanset 是字元集合，要寫寬度。
CSV 用 \texttt{fgets} + \texttt{sscanf}。 輸出目標可以做成
\texttt{FILE *} 參數。
（第~\ref{sec:fio-fmt}~節）
\item \texttt{fread}/\texttt{fwrite} 回傳項目數。 填補、端序、指標欄
位都會進檔案。 結構在宣告時初始化。
（第~\ref{sec:fio-bin}~節）
\item 定長紀錄位移 \texttt{index * sizeof}。 \texttt{r+b} 讀寫之間必
須定位或 \texttt{fflush}。 \texttt{ftell} 失敗是 \texttt{-1L}。
（第~\ref{sec:fio-seek}~節）
\end{itemize}

\begin{drills}[label=drills:fio-checklist]
\item 下面這段有四個錯，請全部標出。
\begin{verbatim}
char ch;
char name[8];
struct S { int id; char *title; } s;
FILE *fp = fopen("d.bin", "r+b");
while (!feof(fp)) {
    ch = fgetc(fp);
    fscanf(fp, "%s", name);
    fwrite(&s, sizeof s, 1, fp);
}
\end{verbatim}
\item 指出錯誤：\texttt{fflush(stdin);}
\item 指出錯誤：結構裡有 \texttt{char *name}，直接
\texttt{fwrite(\&s, sizeof s, 1, fp)}。
\end{drills}

\phantomsection
\addcontentsline{toc}{section}{附錄：可攜性與底層}
\section*{附錄：可攜性與底層}
\label{sec:fio-advanced}

這些內容第一次學檔案用不到，需要時再回來。

\texttt{setvbuf(fp, buf, mode, size)} 可在開檔後、尚未讀寫前改串流緩
衝策略。 \texttt{mode} 為 \texttt{\_IOFBF}（全）、\texttt{\_IOLBF}（列）、
\texttt{\_IONBF}（無）。 多數程式用預設即可。

文字模式下，\texttt{fseek} 只保證能回到先前 \texttt{ftell} 回傳的值，
或回到開頭；隨便算一個位元組位移不一定合法，因為換行轉換會讓邏輯位
置與位元組位置脫鉤。 二進位模式下，ISO C 不保證 \texttt{SEEK\_END}
能用（POSIX 通常可以）。 用定位到檔案尾再問 \texttt{ftell} 來量大小，
在嚴格的標準意義上並非處處成立。

\texttt{fsetpos} 與 \texttt{fgetpos} 用一個不透明的位置型別記位置，
適合很大的檔案；這一章用 \texttt{fseek} 與 \texttt{ftell} 就夠。

\texttt{FILE *} 底下還有作業系統的檔案描述子，英文叫 file
descriptor。 POSIX 的 \texttt{read}/\texttt{write} 走那一層；這一章
停在 \texttt{FILE}。

\phantomsection
\addcontentsline{toc}{section}{練習解答}
\section*{練習解答}

各節練習的答案收在這裡。 從零寫與改寫只給能一句話核對的結果；完整程式
見正文片語。

\paragraph{第~\ref{sec:fio-stream}~節}
\begin{enumerate}
\item \texttt{printf} 進 \texttt{log.txt}，\texttt{fprintf(stderr,
...)} 仍在螢幕。
\item \texttt{prog1} 的 \texttt{stdout}。
\item 一旦 \texttt{stdout} 被重新導向接到檔案，走 \texttt{printf} 的
診斷會跟著消失；\texttt{stderr} 仍留在螢幕。
\item \texttt{>>}。
\end{enumerate}

\paragraph{第~\ref{sec:fio-state}~節}
\begin{enumerate}
\item 兩行 \texttt{hi}（最後一行重複）。
\item 假。 旗標尚未立起。
\item 對輸入串流 \texttt{fflush} 是未定義行為。
\item 讀取 $\rightarrow$ 看回傳值 $\rightarrow$ 處理這一筆（僅當成功）
$\rightarrow$ 需要時問 \texttt{feof}/\texttt{ferror}。
\item 串流緩衝區。 列緩衝還沒遇到換行。
\end{enumerate}

\paragraph{第~\ref{sec:fio-open}~節}
\begin{enumerate}
\item 空的。 \texttt{w} 會截斷。
\item 不行。 \texttt{argc == 1} 表示只有 \texttt{argv[0]}。
\item 不一定。 寫入錯誤可能到 \texttt{fclose} 沖洗串流緩衝區時才發生。
\item \texttt{"r+b"}。
\end{enumerate}

\paragraph{第~\ref{sec:fio-toolbox}~節}
\begin{enumerate}
\item \texttt{fgets}。
\item \texttt{fread}/\texttt{fwrite}。
\item 格式不符。
\item 檔案尾（EOF）。
\end{enumerate}

\paragraph{第~\ref{sec:fio-char}~節}
\begin{enumerate}
\item 讀到 $\mathtt{0xFF}$ 那一格就停，後面沒讀。
\item $1$ $2$ $3$。
\item 沒有差別。 \texttt{getchar()} 就是 \texttt{getc(stdin)}。
\item 同 \texttt{cat-file.c} 的迴圈，檔名來自 \texttt{argv[1]}。
\end{enumerate}

\paragraph{第~\ref{sec:fio-line}~節}
\begin{enumerate}
\item 比不中。 行陣列裡是 \texttt{quit\textbackslash n}。
\item \texttt{abc} 加上空字元，沒有換行。
\item \texttt{\textbackslash n}（字串 \texttt{"\textbackslash n"}）。
\item \texttt{puts} 三個字元；\texttt{fputs} 兩個字元，沒有換行。
\end{enumerate}

\paragraph{第~\ref{sec:fio-fmt}~節}
\begin{enumerate}
\item 回傳 $1$。
\item 寫穿 \texttt{name}，未定義行為。
\item \texttt{name[50]} 只能再容 $49$ 個字元加一個空字元。
\item 迴圈條件改 \texttt{fgets != NULL}，每一行用 \texttt{sscanf}，
失敗就 \texttt{continue}。
\end{enumerate}

\paragraph{第~\ref{sec:fio-bin}~節}
\begin{enumerate}
\item 前 $7$ 個 \texttt{int} 有資料。
\item 沒有。 檔案裡是指標值。
\item 不一定。 \texttt{int} 的四個位元組端序相反，\texttt{n} 的數值
會變。
\item 文字。
\end{enumerate}

\paragraph{第~\ref{sec:fio-seek}~節}
\begin{enumerate}
\item $0$ 與 $3 \times \texttt{sizeof(Record)}$。
\item 位置已經移到下一筆開頭；而且讀後直接寫違反共用串流緩衝區的同
步規定。
\item 不能。 那是失敗值。
\item \texttt{fseek(fp, -sizeof(Record), SEEK\_END)} 再 \texttt{fread}
一筆（附錄：\texttt{SEEK\_END} 在 ISO C 不保證）。 或先量大小再定位
最後一筆。
\end{enumerate}

\paragraph{第~\ref{sec:fio-combine}~節}
\begin{enumerate}
\item 用行陣列逐行 \texttt{fgets}，每次成功就覆寫；迴圈結束後印出
（先確認至少成功讀過一行）。
\item Alice 的殘餘字元留在 Bob 的空字元後面。
\item 兩次。 與第~\ref{sec:fio-state}~節同一種晚一步。
\item 同第~\ref{sec:fio-fmt}~節改寫。
\end{enumerate}

\paragraph{第~\ref{sec:fio-checklist}~節}
\begin{enumerate}
\item 四個錯：\texttt{while (!feof(fp))}；\texttt{char ch} 接
\texttt{fgetc}；\texttt{\%s} 無寬度；讀之後直接 \texttt{fwrite}、中間
沒有定位。 （\texttt{char *title} 再另算一個，寫出的是位址。）
\item 未定義行為。
\item 寫出的是位址，字串內容不在檔案裡。
\end{enumerate}

\chapter{前置處理}
\label{ch:preprocessing}

前置處理器（preprocessor）在編譯最前面工作，做的事情只有一種：把原始碼裡符
合特定樣式的文字，換成另一段文字。 它不管型別、不管運算子優先度、不管變數
的作用域——那些是編譯器的工作，前置處理器根本看不到。

這一章從頭到尾只靠三句話撐起來：

\begin{itemize}
\item \textbf{前置處理器不懂 C。} 它認得的只有一行一行的指令、以及一個一
個的 token；後面每一個陷阱，都是這句話的推論。
\item \textbf{展開是可觀察的。} \texttt{gcc -E} 讓前置處理器的輸出直接被
看見。 這一章每個坑都遵守同一個順序：先自己手寫展開結果、再用
\texttt{gcc -E} 核對——「精熟」跟「背過規則」的分界，就在這裡。
\item \textbf{巨集是最後手段。} 常數優先用 \texttt{enum}，邏輯優先用
\texttt{static inline} 函式；巨集只留給它們做不到的事。
\end{itemize}

章節安排跟著這三句話走：先弄懂前置處理器實際做的事、指令的通用語法規則，
再接上第一個真正用得到的指令 \texttt{\#include}；接著進物件式與函式式巨
集，中段插入「巨集還是函式」的判準，再教兩個進階運算子與可變參數巨集；然
後是條件編譯與命令列開關；最後把所有零件收攏成一組讀者自己寫得出來的
\texttt{LOG} 與 \texttt{assert}。 節末練習的答案收在章末。

\section{翻譯的第一階段}
\label{sec:pp-what}

一支 \texttt{.c} 檔要變成執行檔，中間經過好幾個階段：前置處理、編譯（翻譯
成組合語言）、組譯（組合語言變成目的檔）、連結（把目的檔和函式庫接起來）。
每個階段都有自己的輸出：前置處理器的輸出叫\textbf{翻譯單元}
（translation unit）——它仍然是純文字的 C 原始碼，不是二進位，也不是什麼
特殊格式，會直接餵給下一站的編譯器；編譯器把翻譯單元變成組合語言；組譯器
把組合語言變成\textbf{目的檔}（object file，\texttt{.o}）——這時候才第一
次出現二進位；連結器把一個或多個目的檔、加上函式庫，接成最終的執行檔。

\begin{center}
\begin{tikzpicture}[node distance=1.2cm and 1.0cm]
    \node[startend, minimum width=1.9cm, font=\small\sffamily] (src) {原始碼\\\texttt{.c}};
    \node[process, right=of src, minimum width=2.6cm, font=\small\sffamily] (cpp) {前置處理器};
    \node[startend, right=of cpp, minimum width=2.4cm, font=\small\sffamily] (tu) {翻譯單元\\（仍是文字）};
    \node[process, right=of tu, minimum width=3.3cm, font=\small\sffamily] (exe) {編譯器／組譯器\\／連結器};
    \draw[flowarrow] (src) -- (cpp);
    \draw[flowarrow] (cpp) -- (tu);
    \draw[flowarrow] (tu) -- (exe);
\end{tikzpicture}
\diagcaption{前置處理是編譯流程的第一站，輸出（翻譯單元）還是文字形式的
C 原始碼，不是二進位。 這一章只管第一站；後面三站在函式那幾章已經摸過邊。}
\label{fig:pp-pipeline}
\end{center}

眼見為憑最有效。 下面這支程式定義了一個常數字串巨集：

\begin{example}{c}{(hello.c) 讓前置處理器做一次文字替換}
#include <stdio.h>

#define GREETING "Hello, preprocessor!"

int main(void) {
    printf("%s\n", GREETING);
    return 0;
}
\end{example}

用 \texttt{-E} 旗標讓 \texttt{gcc} 只跑到前置處理這一站就停，把結果印到標
準輸出：

\begin{code}{bash}{gcc -E}
gcc -E hello.c | wc -l      # 幾百行，因平台與函式庫版本而異
gcc -E hello.c | tail -6    # 看最後幾行
\end{code}

\begin{code}{text}{gcc -E 輸出（節錄最後幾行）}
int main(void) {
    printf("%s\n", "Hello, preprocessor!");
    return 0;
}
\end{code}

前面幾百行是 \texttt{\#include <stdio.h>} 展開出來的宣告——\texttt{printf}
的原型、一堆型別定義、其他標頭檔又 \texttt{\#include} 進來的東西，全部攤平
變成一份巨大的純文字檔，這正是翻譯單元的樣子。 最後幾行才是我們自己寫的
\texttt{main}，而且 \texttt{GREETING} 已經被換成它展開後的樣子：字面上的
字串常數，跟我們直接寫 \texttt{"Hello, preprocessor!"} 沒有任何差別。 前
置處理器不知道 \texttt{GREETING} 「是」字串——它只知道看到這個 token 就換
成另一串 token。

\keypoint{前置處理器的輸出（翻譯單元）仍是合法的 C 原始碼，會原封不動送進
編譯器。 \texttt{\#define} 只是純文字替換，替換之後，替換前的名字就不存在
了。}

\begin{drills}[label=drills:pp-what]
\item \texttt{gcc -E hello.c} 印出的行數，跟 \texttt{printf(...)} 那一行
比起來，誰佔的行數多？多出來的內容從哪裡來？
\item 前置處理器的輸出（翻譯單元）還是不是合法的 C 程式？會不會被拿去繼
續編譯？
\item 如果把 \texttt{\#include <stdio.h>} 整行刪掉，只留
\texttt{\#define GREETING ...} 和呼叫 \texttt{printf} 的那一行，會在前置
處理、編譯、連結哪一個階段出錯？
\item 手寫：\texttt{\#define TWICE(x) x+x}，呼叫 \texttt{TWICE(3)} 展開成
什麼？先自己寫，再用 \texttt{gcc -E} 核對。
\end{drills}

\section{指令的通用規則}
\label{sec:pp-syntax}

在教第一個真正用得到的指令之前，先把所有前置處理指令共用的語法規則講清
楚，之後每一節都會直接用到，不再重複解釋。

\subsection{一行一個指令，到換行為止}

前置處理指令一定以 \texttt{\#} 開頭。 \texttt{\#} 前面只能有空白字元（用
來縮排），不能有其他程式碼；指令從 \texttt{\#} 開始，一路到這一行的換行字
元為止，中間不用分號結尾（這跟一般 C 陳述式不一樣）：

\begin{example}{c}{(continuation-demo.c) \# 前面可以有空白}
#include <stdio.h>

#define BIG_NUMBER \
    123456

int main(void) {
    #define LOCAL_VALUE 5   /* 縮排的 #，合法 */
    printf("%d\n", BIG_NUMBER);
    printf("%d\n", LOCAL_VALUE);
    return 0;
}
\end{example}

\iosample{}{123456\\5}

\texttt{\#define LOCAL\_VALUE 5} 前面有四個空白，這是合法的——縮排跟指令
本身無關，前置處理器只在乎 \texttt{\#} 前面除了空白之外沒有別的東西。

\subsection{反斜線續行}

一個指令原則上要寫在同一個實體行裡，但如果這一行的最後一個字元是反斜線
\texttt{\textbackslash}（後面緊接著換行字元，中間不能有空白，否則續行失
敗），前置處理器會把這一行和下一行接成邏輯上的同一行，再當成同一個指令處
理。 上面範例裡的 \texttt{BIG\_NUMBER} 就是這樣定義的：

\begin{code}{c}{反斜線續行}
#define BIG_NUMBER \
    123456
\end{code}

這等價於把兩行接成一行寫：\texttt{\#define BIG\_NUMBER 123456}。 反斜線續
行不是巨集專屬的語法——它是前置處理器對「一行」的定義，任何指令都能用；後
面函式式巨集展開成好幾個陳述式時，會大量用到這個語法把巨集寫成好幾個實體
行、但邏輯上仍是同一個 \texttt{\#define}。

\subsection{token：前置處理器真正比對的單位}

前置處理器做替換時，比對的單位是\textbf{token}（識別字、關鍵字、數字常
數、字串常數、運算子這些不可再分割的最小單元），不是子字串：

\begin{example}{c}{(token-demo.c) N100 與 "N" 不會被 \#define N 換掉}
#include <stdio.h>

#define N 100

int main(void) {
    int N100 = 1;       /* N100 是單一個 token，不是 N 後面接 100 */
    char *s = "N";       /* 整個 "N" 是一個字串常數 token，不會被拆開比對 */
    printf("%d\n", N100);
    printf("%s\n", s);
    printf("%d\n", N);
    return 0;
}
\end{example}

\iosample{}{1\\N\\100}

\texttt{N100} 這五個字元對編譯器和前置處理器來說是同一個識別字 token，跟
\texttt{N} 完全無關，所以不會被換成 \texttt{100100}；\texttt{"N"} 整體是
一個字串常數 token，前置處理器不會鑽進字串內部去比對裡面的字元，所以引號
裡的 \texttt{N} 不會被換掉。 只有真正單獨出現、跟 \texttt{N} 這個識別字
token 完全相等的地方，才會被替換——最後一行的 \texttt{N} 就是這種情況，換
成了 \texttt{100}。

\texttt{\#define} 的名字後面如果緊接著 \texttt{(}（中間沒有空白），前置
處理器會把它當成函式式巨集；名字後面如果先隔了空白才出現
\texttt{(}，反而會被當成物件式巨集，這個空白引發的陷阱留到第
\ref{sec:pp-funcmacro}~節函式式巨集一開始示範。

\keypoint{指令從 \texttt{\#} 開始，到這一行的換行字元結束；反斜線
\texttt{\textbackslash} 後面緊接換行，能讓指令跨實體行接續。 巨集替換比
對的單位是 token，不是子字串——\texttt{N100} 不是 \texttt{N} 加
\texttt{100}，字串常數裡的文字也不會被鑽進去替換。}

\begin{drills}[label=drills:pp-syntax]
\item 手寫：把 \texttt{BIG\_NUMBER} 的定義拿掉反斜線、直接分成兩行寫
（不續行），\texttt{BIG\_NUMBER} 會變成什麼內容？下一行的
\texttt{123456} 會發生什麼事？
\item 為什麼 \texttt{\#} 前面可以有空白（像縮排在 \texttt{if} 區塊裡的
\texttt{\#define}），但 \texttt{\#} 後面不能先出現其他程式碼？
\item \texttt{\#define N 100} 之後，\texttt{N100} 和 \texttt{"N"} 會不會
被換成別的東西？為什麼？
\item 補空：一個前置處理指令從 \texttt{\#} 開始，到 \underline{\hspace{2em}}
為止；除非用 \underline{\hspace{2em}} 續行。
\end{drills}

\section{\#include}
\label{sec:pp-include}

\texttt{\#include} 是最簡單、也最先用到的前置處理指令：把某個檔案的內容原
封不動貼進來，取代掉這一行。 兩種寫法找檔案的位置不同：

\begin{code}{c}{include 的兩種形式}
#include <stdio.h>     // 先找系統／編譯器內建的標頭檔路徑
#include "myheader.h"  // 先找目前檔案所在的目錄，找不到才退到系統路徑
\end{code}

\texttt{<>} 用在標準函式庫或第三方函式庫的標頭檔；\texttt{""} 用在自己專案
裡的標頭檔。

\subsection{標頭檔放宣告，不放定義}

標頭檔該放的是\textbf{宣告}——「這個函式存在、長這個樣子」；函式的本體（\textbf{定義}）
留在對應的 \texttt{.c} 檔裡。 原因不是風格問題，是連結器的規則：同一個函式
的本體在整個程式裡只能出現一次。 如果標頭檔裡放了定義，兩個 \texttt{.c} 檔
各自 \texttt{\#include} 它，各自編譯出的目的檔裡都會有一份函式本體，連結時
就會撞名：

\begin{code}{c}{baddefs.h（錯誤示範：標頭檔裡放了定義）}
int square(int x) {   /* 不該放在標頭檔 */
    return x * x;
}
\end{code}

\texttt{a.c} 與 \texttt{b.c} 都 \texttt{\#include "baddefs.h"}，個別編譯都
不會出錯——前置處理器只是把這三行貼進兩個檔案，編譯器看到的也都是合法的
函式定義。 錯誤要到連結那一站才會現形：

\begin{code}{text}{連結器的錯誤}
duplicate symbol '_square' in:
    a.o
    b.o
ld: 1 duplicate symbols
\end{code}

（不同平台的連結器訊息不一樣，Linux 上常見的是「multiple definition of
\texttt{square}」，意思相同。）

改成標頭檔只放宣告、實作放進各自的 \texttt{.c} 檔，才是正確的分工：

\begin{example}{c}{(mathutils.h / mathutils.c / main.c) 標頭檔放宣告，實作放 .c 檔}
/* mathutils.h：只有宣告 */
int max_int(int a, int b);
int min_int(int a, int b);

/* mathutils.c：真正的定義 */
#include "mathutils.h"

int max_int(int a, int b) {
    return a > b ? a : b;
}

int min_int(int a, int b) {
    return a < b ? a : b;
}

/* main.c：只需要知道函式的樣子，不需要看到本體 */
#include <stdio.h>
#include "mathutils.h"

int main(void) {
    printf("%d %d\n", max_int(3, 7), min_int(3, 7));
    return 0;
}
\end{example}

\iosample{}{7 3}

編譯時把兩個 \texttt{.c} 檔一起交給編譯器：\texttt{gcc main.c mathutils.c
-o main}。 \texttt{main.c} 完全不需要知道 \texttt{max\_int} 怎麼實作，只
要看過標頭檔裡的宣告，就知道怎麼呼叫它。

\subsection[重複引入的問題]{重複引入的問題（先看問題，第~\ref{sec:pp-cond}~節見真章）}

標頭檔常常互相 \texttt{\#include}——例如 \texttt{A.h} 和 \texttt{B.h} 都
\texttt{\#include} 了同一個 \texttt{C.h}，而某個 \texttt{.c} 檔同時用到
\texttt{A.h} 與 \texttt{B.h}，那麼 \texttt{C.h} 的內容就會被貼進同一個翻
譯單元兩次。 如果 \texttt{C.h} 裡有 \texttt{struct} 或型別定義，重複貼上
會變成重複定義，編不過。 解法要用到條件編譯——\texttt{\#ifndef} +
\texttt{\#define} 的 include guard、或非標準但幾乎所有編譯器都支援的
\texttt{\#pragma once}，留到第~\ref{sec:pp-cond}~節一起講。

\keypoint{\texttt{\#include} 是檔案層級的文字貼上。 標頭檔放宣告不放定
義，否則同一個函式本體出現在多個目的檔裡，連結期會撞名；同一個標頭檔被間
接引入超過一次，也可能因為重複定義而編不過。}

\begin{drills}[label=drills:pp-include]
\item \texttt{\#include <stdio.h>} 與 \texttt{\#include "myheader.h"} 找
檔案的路徑順序有什麼不同？
\item \texttt{baddefs.h} 那個例子錯在哪個階段被抓到：前置處理、編譯，還是
連結？為什麼編譯階段抓不到？
\item \texttt{mathutils.h} 這一節的例子只被 \texttt{\#include} 一次，沒有
出任何問題。 猜猜看：同一個標頭檔被同一個 \texttt{.c} 檔案間接
\texttt{\#include} 兩次，會是什麼情況造成的？（第~\ref{sec:pp-cond}~節見
真章）
\item 補空：標頭檔裡對函式該寫 \underline{\hspace{2em}}，\texttt{.c} 檔裡
才寫函式的 \underline{\hspace{2em}}。
\end{drills}

\section{物件式巨集 \#define}
\label{sec:pp-define}

\texttt{\#define NAME 替換內容} 定義一個\textbf{物件式巨集}（object-like
macro）：往後原始碼裡出現的 \texttt{NAME} 這個 token，都會被換成後面那串
內容。 最常見的用法是取代程式裡散落的魔術數字：

\begin{code}{c}{define 常數}
#define MAX_SIZE 100
#define PI 3.14159

int arr[MAX_SIZE];
double area = PI * r * r;
\end{code}

\subsection{重新定義與 \texttt{\#undef}}

同一個名字重複 \texttt{\#define} 成完全一樣的內容，合法、不會有任何警告；
定義成不一樣的內容，會得到一個「重新定義」的警告：

\begin{code}{text}{未先 \#undef 就重新定義}
warning: 'DOUBLE' redefined
note: this is the location of the previous definition
\end{code}

要合法地換掉一個巨集的定義，中間要先 \texttt{\#undef}：

\begin{code}{c}{先 \#undef 再重新 \#define}
#define DOUBLE(x) 2 * x
#undef DOUBLE
#define DOUBLE(x) (2 * (x))
\end{code}

這樣寫編譯乾淨，沒有警告。 下一節會回來解釋 \texttt{DOUBLE} 這個例子本身
的問題。

\subsection[\#define vs const vs enum]{\texttt{\#define} 與 \texttt{const} 與 \texttt{enum}：陣列大小}

想要一個「編譯期常數」，C 裡至少有三個工具，行為並不相同：

\begin{example}{c}{(const-vs-define.c) 三種常數在陣列大小上的差別}
#include <stdio.h>

#define MAX_SIZE 100          /* 巨集：純文字替換，隨處可當常數 */
enum { ENUM_SIZE = 100 };     /* enum 常數：真正的整數常數表達式 */
const int array_size = 100;   /* const 變數：有型別、有作用域，
                                  但不是常數表達式 */

int global_by_macro[MAX_SIZE]; /* OK：展開後就是字面常數 100 */
int global_by_enum[ENUM_SIZE]; /* OK：enum 常數也是編譯期常數 */
/* int global_by_const[array_size];  錯誤：見下方 */

int main(void) {
    int local_by_const[array_size]; /* 這裡合法，但它是 VLA */
    printf("%zu\n", sizeof local_by_const / sizeof(int));
    return 0;
}
\end{example}

\iosample{}{100}

把註解拿掉、讓 \texttt{global\_by\_const} 真的出現在檔案作用域，會編不過：

\begin{code}{text}{檔案作用域用 const 當陣列大小}
error: variably modified 'global_by_const' at file scope
\end{code}

\texttt{array\_size} 是一個有型別、有位址的變數，即使它是 \texttt{const}、
即使它的值一眼就能看出來是 \texttt{100}，C 標準也不承認它是「整數常數表達
式」——那是 C++ 的規則，C 沒有這樣放寬。 檔案作用域的陣列大小必須是真正的
常數表達式，所以編不過。 在函式內部，\texttt{array\_size} 可以用來宣告陣
列，但那是 C99 引入的\textbf{可變長度陣列}（VLA）：大小在執行期才決定，不
是編譯期常數，不能有初始化列表、不能是 \texttt{static}、不能當結構成員。
\texttt{\#define} 與 \texttt{enum} 常數兩個都保留了「隨處可用的編譯期常
數」這個特性；\texttt{enum} 額外多了一點型別性質（它是一種列舉型別的值），
\texttt{\#define} 則完全沒有型別、沒有作用域，也不會出現在偵錯器裡。

三種工具怎麼選，用下面這個表格判斷：

\begin{center}
\begin{tabular}{ll}
\toprule
需求 & 選擇 \\
\midrule
只是取代程式裡的魔術數字，不需要型別 & \texttt{enum}（偵錯器看得到名字） \\
陣列大小、\texttt{switch case} 等需要整數常數表達式 & \texttt{\#define} 或 \texttt{enum} \\
有型別、有值、但不需要出現在常數表達式裡 & \texttt{const} 變數 \\
型別本身要取別名 & \texttt{typedef}（見下方，\textbf{不要}用 \texttt{\#define}）\\
值超出 \texttt{int} 範圍（如 \texttt{1u << 31}） & \texttt{\#define}（見下方說明）\\
\bottomrule
\end{tabular}
\end{center}

最後一列值得展開講：C 標準規定\textbf{列舉常數的型別是 \texttt{int}}，像
\texttt{1u << 31}、\texttt{0xFFFFFFFFUL} 這種超出 \texttt{int} 範圍的值放
進 \texttt{enum}，嚴格來說不可攜——GCC／Clang 在預設選項下會用擴充功能悄
悄接受，但加上 \texttt{-pedantic} 就會看到

\begin{code}{text}{-pedantic 抓到 enum 常數超出 int 範圍}
warning: ISO C restricts enumerator values to range of 'int'
    (2147483648 is too large)
\end{code}

這種需要塞下大範圍或位元遮罩的場合（跟位元操作那一章的 flag 常數常常一
起出現），\texttt{enum} 裝不下，只能回頭用 \texttt{\#define}。

\subsection{危險示範：\texttt{\#define} 當型別別名}

\texttt{\#define} 也能拿來幫型別取別名，但這是它最不該被用在的地方：

\begin{example}{c}{(intptr-demo.c) \#define 當型別別名的陷阱}
#include <stdio.h>

#define INT_PTR int * /* 危險：#define 版本 */

int main(void) {
    INT_PTR p1, p2; /* 展開成 int * p1, p2; */
    printf("sizeof p1 = %zu (pointer)\n", sizeof(p1));
    printf("sizeof p2 = %zu (int, not pointer!)\n", sizeof(p2));
    return 0;
}
\end{example}

\iosample{}{sizeof p1 = 8 (pointer)\\sizeof p2 = 4 (int, not pointer!)}

（\texttt{8} 與 \texttt{4} 是常見 64 位元平臺——macOS／Linux 的
x86-64、ARM64——上指標與 \texttt{int} 的大小，不是 C 標準保證的數字；換
一個資料模型的平臺，這兩個數字可能不同，但「星號只黏 \texttt{p1}」這個文
字替換的結論不會變。）

\texttt{INT\_PTR p1, p2;} 展開後是 \texttt{int * p1, p2;}——星號只黏在
\texttt{p1} 前面，\texttt{p2} 只是一個普通的 \texttt{int}。 這正是文字替換
的後果：\texttt{\#define} 不知道自己「代表一個型別」，它只是把 \texttt{INT\_PTR}
換成 \texttt{int *} 這三個 token，剩下的宣告語法照舊解讀，而 C 的宣告語
法本來就是「星號跟著變數名，不是跟著型別」。 \texttt{typedef} 不會有這個問
題，因為 \texttt{typedef} 是編譯器真正認得的型別別名，不是文字替換：

\begin{code}{c}{typedef 沒有這個陷阱}
typedef int *IntPtr;
IntPtr p1, p2;  /* 兩個都是 int*，因為型別是 IntPtr 本身 */
\end{code}

\keypoint{\texttt{\#define} 沒有型別、沒有作用域，重複定義成不同內容要先
\texttt{\#undef}。 陣列大小要編譯期常數，選 \texttt{\#define} 或
\texttt{enum}，不要用 \texttt{const} 變數。 型別別名一律用
\texttt{typedef}，不要用 \texttt{\#define}。}

\begin{drills}[label=drills:pp-define]
\item 手寫：\texttt{\#define MAX\_SIZE 100} 之後又寫一次一模一樣的
\texttt{\#define MAX\_SIZE 100}，編譯器會抗議嗎？
\item \texttt{int global\_by\_const[array\_size];} 在檔案作用域為什麼編不
過？錯誤訊息在說什麼？
\item \texttt{IntPtr} 那個例子裡，\texttt{sizeof(p1)} 和 \texttt{sizeof(p2)}
為什麼不一樣？
\item 補空：要在檔案最外層宣告陣列大小，除了字面常數，能用的是
\texttt{\#define} 與 \underline{\hspace{2em}}。
\end{drills}

\section{函式式巨集}
\label{sec:pp-funcmacro}

\texttt{\#define} 也能定義帶參數的\textbf{函式式巨集}（function-like
macro），語法看起來很像函式呼叫，但骨子裡還是文字替換：

\begin{code}{c}{函式式巨集}
#define SQUARE(x) ((x) * (x))
#define ABS(x) ((x) < 0 ? (-(x)) : (x))
\end{code}

\subsection{名字後面那個空白，差很多}

前置處理器判斷一個 \texttt{\#define} 是不是函式式巨集，規則很死板：巨集
名字後面要\textit{緊接著}左括號，中間不能有空白。 只要多了一個空白，整個
就變成物件式巨集，連小括號帶後面的內容全部原封不動當成替換文字：

\begin{code}{c}{名字跟左括號之間多了一個空白}
#define TWICE (x) x+x   /* 危險：TWICE 後面有空白，這是物件式巨集 */
\end{code}

\begin{code}{text}{gcc -E 的展開結果}
int result = (x) x+x(3);
\end{code}

\texttt{TWICE} 被定義成一個物件式巨集，替換內容是\textit{整段}
\texttt{(x) x+x}；\texttt{TWICE(3)} 裡的 \texttt{TWICE} 被換成
\texttt{(x) x+x}，後面呼叫時多打的 \texttt{(3)} 完全沒被巨集吃進去，只是
原封不動貼在展開結果後面，變成 \texttt{(x) x+x(3)}——一段語法上不知所云的
文字，\texttt{x} 也從沒被定義過。 這個坑純粹是空白造成的：本章其餘函式式
巨集的定義，名字跟左括號之間都刻意不留空白。

\begin{example}{c}{(macro-basic.c) 呼叫函式式巨集}
#include <stdio.h>

#define SQUARE(x) ((x) * (x))
#define ABS(x) ((x) < 0 ? (-(x)) : (x))

int main(void) {
    int a = 5;
    printf("%d\n", SQUARE(a));     /* 25 */
    printf("%d\n", ABS(-10));      /* 10 */
    printf("%d\n", SQUARE(1 + 2)); /* 括號救了這一行：9，不是 5 */
    return 0;
}
\end{example}

\iosample{}{25\\10\\9}

\subsection{括號規則}

函式式巨集的參數，以及整個展開結果，都要用括號包起來，理由是運算子優先
度。 沒括號的巨集會被展開位置的上下文吃掉優先度：

\begin{code}{c}{沒括號的巨集}
#define DOUBLE(x) 2 * x   /* 危險：沒括號 */

DOUBLE(3 + 4)
\end{code}

\texttt{DOUBLE(3 + 4)} 展開後是 \texttt{2 * 3 + 4}——乘法先做，結果是
\texttt{10}，不是預期的 \texttt{14}。 加上兩層括號才安全：

\begin{code}{c}{加上括號}
#define DOUBLE(x) (2 * (x))

DOUBLE(3 + 4)   /* 展開成 (2 * (3 + 4))，得到 14 */
\end{code}

參數的括號防止參數本身是一串運算式時被拆散；最外層的括號防止展開結果被
呼叫端的其他運算子拆散。 兩層都不能省。

\subsection{副作用：參數被求值超過一次}

文字替換還有一個後果：巨集的參數在展開結果裡出現幾次，就被求值幾次。
如果參數帶有副作用（例如 \texttt{++}），麻煩就來了：

\begin{example}{c}{(macro-side-effect.c) SQUARE(i++) 的未定義行為}
#include <stdio.h>

#define SQUARE(x) ((x) * (x))

static inline int square(int x) {
    return x * x;
}

int main(void) {
    int i = 5;
    printf("%d\n", SQUARE(i++)); /* i++ 被展開了兩次 */
    printf("i = %d\n", i);

    i = 5;
    printf("%d\n", square(i++)); /* i++ 只呼叫一次 */
    printf("i = %d\n", i);
    return 0;
}
\end{example}

\iosample{}{30\\i = 7\\25\\i = 6}

\texttt{SQUARE(i++)} 展開成 \texttt{((i++) * (i++))}——對同一個變數 \texttt{i}
做了兩次修改，中間沒有一個標準保證的\textbf{序列點}，這在 C 裡是\textbf{未
定義行為}，不是「巨集固定會算出某個錯誤答案」。 這台機器、這個編譯器算出
\texttt{30} 且 \texttt{i} 變成 \texttt{7}（編譯器甚至會主動警告
\texttt{-Wsequence-point} 或 \texttt{-Wunsequenced}），換一顆編譯器或換一
個優化等級，結果都可能不同。 唯一可靠的說法是：\textbf{不要把有副作用的運
算式傳進函式式巨集}。 真正的解法是 \texttt{static inline} 函式：
\texttt{square(i++)} 裡的 \texttt{i++} 是一個真正的函式呼叫參數，只求值一
次，行為和一般函式呼叫完全相同、可預期。

\subsection{多行展開：先埋個伏筆}

函式式巨集的展開結果不一定只有一個運算式；如果需要好幾個陳述式，就要靠
第~\ref{sec:pp-syntax}~節教過的反斜線續行把它們接成同一個 \texttt{\#define}。
但這樣寫出來的多行巨集，直接拿去接在 \texttt{if}／\texttt{else} 後面會出
問題——第~\ref{sec:pp-assert}~節做 \texttt{LOG} 與 \texttt{assert} 時會示
範完整的坑與修法（\texttt{do \{ ... \} while (0)}）。

\keypoint{函式式巨集是文字替換，不是函式呼叫。 參數與整個展開式都要加括
號；參數若有副作用會被重複求值，是未定義行為，不是「巨集特性」。 需要保
證只求值一次、又想要保留一點效能，用 \texttt{static inline} 函式。}

\begin{drills}[label=drills:pp-funcmacro]
\item \texttt{DOUBLE(3+4)} 用沒有括號的版本展開成什麼？印出來是多少？
\item 加上括號的版本怎麼寫？驗證 \texttt{DOUBLE(3+4)} 變成 \texttt{14}。
\item \texttt{SQUARE(i++)} 印出 \texttt{30}、\texttt{i} 變成 \texttt{7}
——這是「巨集固定會這樣算」還是「未定義行為，這台機器剛好這樣算」？
\item 補空：函式式巨集碰到有副作用的參數（例如 \texttt{i++}）時，最安全
的替代方案是 \underline{\hspace{2em}}。
\item 手寫：\texttt{\#define TWICE (x) x+x}（名字後面有空白），
\texttt{TWICE(3)} 會展開成什麼？先自己寫，再用 \texttt{gcc -E} 核對。
\end{drills}

\section{巨集還是函式}
\label{sec:pp-guidelines}

巨集不是壞東西，但它比函式危險，原因前面兩節已經看到兩個：沒有型別檢
查（\texttt{IntPtr}）、參數可能被重複求值（\texttt{SQUARE(i++)}）。 第三
個理由是巨集會\textbf{隱藏控制流程}：

\begin{example}{c}{(hidden-control-flow.c) 巨集裡藏著 return}
#include <stdio.h>

#define CHECK_POSITIVE(x) if ((x) <= 0) return -1

int check(int a, int b) {
    CHECK_POSITIVE(a);
    printf("checked a\n");
    CHECK_POSITIVE(b);
    printf("checked b\n");
    return 0;
}

int main(void) {
    printf("%d\n", check(1, 2));
    printf("%d\n", check(1, -2));
    return 0;
}
\end{example}

\iosample{}{checked a\\checked b\\0\\checked a\\-1}

\texttt{check(1, -2)} 只印出 \texttt{checked a}，接著就直接印
\texttt{-1}——\texttt{CHECK\_POSITIVE(b)} 那一行看起來像個普通的函式呼叫，
但它展開出來的 \texttt{return -1;} 直接讓 \texttt{check} 提前返回，呼叫端
完全看不出這一行「可能會讓函式結束」。 如果改成回傳 \texttt{bool} 的函
式，呼叫端得自己寫 \texttt{if (!check\_positive(b)) return -1;}，控制流程
一眼就看得到。

\texttt{static inline} 函式速度上有沒有損失？在現代編譯器（開了最基本的
優化）底下，\texttt{static inline} 函式幾乎一定會被展開成跟巨集一樣的內
嵌程式碼，速度差異可以忽略；犧牲的那一點編譯期間，換來型別檢查、只求值
一次、偵錯器可以單步進去——這筆帳幾乎總是划算。

跨檔案共用的 \texttt{static inline} 函式該放哪？第~\ref{sec:pp-include}~
節的鐵律是「標頭檔放宣告不放定義」，原因是同一個函式本體出現在多個目的
檔裡會撞名。 但 \texttt{static inline} 是這條鐵律的正當例外：
\texttt{static} 讓這個函式只有\textbf{內部連結}，每個 \texttt{\#include}
它的翻譯單元各自拿到一份獨立的本體，彼此之間根本不會撞名，連結器也就不
會抱怨。 所以跨檔案共用的 \texttt{static inline} 函式，直接整份（宣告加
定義）寫在標頭檔裡就是正確用法，不用像一般函式那樣把定義拆到 \texttt{.c}
檔。

歸納整章的判準：

\begin{itemize}
\item 純數字常數，選 \texttt{enum} 優先於 \texttt{\#define}——多一點型別
資訊，偵錯器也看得到名字（第~\ref{sec:pp-define}~節）。
\item 型別別名一律用 \texttt{typedef}，不要用 \texttt{\#define}
（第~\ref{sec:pp-define}~節）。
\item 邏輯稍微複雜、或參數可能有副作用，優先用 \texttt{static inline}
函式（第~\ref{sec:pp-funcmacro}~節）。
\item 巨集會讓陳述消失或提前返回時（像 \texttt{CHECK\_POSITIVE}），寧可
多寫幾個字，換成看得到控制流程的函式。
\end{itemize}

真正只有巨集能做到、函式完全取代不了的，只有四種場合：

\begin{enumerate}
\item \textbf{Token 操作}：字串化 \texttt{\#} 與 token 貼接 \texttt{\#\#}
（第~\ref{sec:pp-stringize}~節）——函式的參數是值，沒辦法把參數名字變成
字串，也沒辦法用參數拼出新的識別字。
\item \textbf{接受不定數量的參數}：可變參數巨集 \texttt{\_\_VA\_ARGS\_\_}
（第~\ref{sec:pp-variadic}~節）。
\item \textbf{整段程式碼要在編譯期就消失或選擇性存在}：條件編譯
（第~\ref{sec:pp-cond}~節）——函式呼叫至少會被編譯出符號，巨集能讓程式碼
連編譯器都看不到。
\item \textbf{需要呼叫端當下的原始碼位置、或需要懶惰求值}：
\texttt{\_\_FILE\_\_}／\texttt{\_\_LINE\_\_}（第~\ref{sec:pp-cmdline}~節）、
或條件不成立時完全不求值某個運算式——函式呼叫前，所有引數一定會先被求
值，巨集不用。
\end{enumerate}

\keypoint{巨集沒有自己的函式呼叫框架，debugger 只看得到展開後的結果，看
不到巨集名字本身，也不能單步跳進去。 能用 \texttt{enum}、
\texttt{typedef}、\texttt{static inline} 解決的，就不要用巨集；真正該用
巨集的只有 token 操作、可變參數、條件編譯、與 \texttt{\_\_FILE\_\_}／
\texttt{\_\_LINE\_\_} 這類懶惰求值的場合。}

\begin{drills}[label=drills:pp-guidelines]
\item \texttt{CHECK\_POSITIVE(b)} 那一行看起來像函式呼叫，但實際藏了什麼
陳述？\texttt{check(1, -2)} 為什麼在 \texttt{checked a} 之後就沒有
\texttt{checked b}？
\item 如果把 \texttt{CHECK\_POSITIVE} 換成一個回傳 \texttt{bool} 的函
式，呼叫端要多寫哪一行，換來的可讀性是什麼？
\item 這一章出現過三種「巨集比函式危險」的理由，各自配一個本章出現過的
例子：重複求值、型別不安全、隱藏控制流程。
\item 補空：巨集真正不可取代的四種場合是 token 操作、可變參數、
\underline{\hspace{2em}}、與需要懶惰求值或呼叫端原始碼位置的
\underline{\hspace{2em}}。
\end{drills}

\section{\texttt{\#} 與 \texttt{\#\#}}
\label{sec:pp-stringize}

函式式巨集還有兩個特殊運算子，只能在巨集定義裡使用。 \texttt{\#} 把參數變
成字串常數（\textbf{字串化}）；\texttt{\#\#} 把兩個 token 黏成一個
（\textbf{token 貼接}）。

\begin{example}{c}{(stringize-paste.c) \# 與 \#\# 的效果}
#include <stdio.h>

#define STR(x) #x
#define PRINT_VAR(v) printf(#v " = %d\n", v)
#define CONCAT(a, b) a##b

int main(void) {
    int count = 42;
    printf("%s\n", STR(hello));     /* "hello" */
    printf("%s\n", STR(1 + 2));     /* "1 + 2"，不是 "3" */
    PRINT_VAR(count);               /* count = 42 */

    int CONCAT(my, Var) = 10;       /* 變成 int myVar = 10; */
    printf("%d\n", myVar);
    return 0;
}
\end{example}

\iosample{}{hello\\1 + 2\\count = 42\\10}

\texttt{STR(1 + 2)} 展開成字串常數 \texttt{"1 + 2"}——\texttt{\#} 只是把參
數的原始文字包上引號，並不會先把 \texttt{1 + 2} 算出來變成 \texttt{"3"}，
再次證明前置處理器不懂 C 的運算規則。 \texttt{PRINT\_VAR(v)} 同時用了
\texttt{\#v}（把變數名字變成字串印出來）和 \texttt{v}（真正的變數，拿去
求值），兩件事各用各的方式表達，缺一不可。 \texttt{CONCAT(a, b)} 用
\texttt{\#\#} 把 \texttt{my} 和 \texttt{Var} 黏成一個新的識別字
\texttt{myVar}——這跟只用空白隔開完全不同：

\begin{code}{c}{沒有 \#\# 只用空白}
#define CONCAT_SPACE(a, b) a b
int CONCAT_SPACE(my, Var) = 5;  /* 展開成 int my Var = 5;，兩個 token，編不過 */
\end{code}

沒有 \texttt{\#\#} 時，\texttt{a b} 展開後仍是兩個各自獨立的 token
\texttt{my} 和 \texttt{Var}，中間隔著一個空白，語法上變成宣告了
\texttt{my}、後面又冒出一個不知道哪來的 \texttt{Var}，編譯器會直接報錯。
\texttt{\#\#} 才是真正把兩個 token 熔接成一個新 token 的運算子。

\subsection{陷阱：\texttt{\#}／\texttt{\#\#} 不會先展開參數}

一般情況下，函式式巨集的參數在替換進去之前，如果自己也是一個巨集，會先
被展開；但參數只要被 \texttt{\#} 或 \texttt{\#\#} 用到，就是唯一的例外
——這兩個運算子作用到的參數\textbf{不會}先展開，是原始傳進來的文字直接
字串化或貼接：

\begin{example}{c}{(str-xstr.c) 兩層間接才能先展開再字串化}
#include <stdio.h>

#define VERSION 3
#define STR(x) #x
#define XSTR(x) STR(x)

int main(void) {
    printf("%s\n", STR(VERSION));   /* "VERSION" */
    printf("%s\n", XSTR(VERSION));  /* "3" */
    return 0;
}
\end{example}

\iosample{}{VERSION\\3}

\texttt{STR(VERSION)} 裡的 \texttt{VERSION} 直接被 \texttt{\#} 字串化，
前置處理器不管它本身也是一個巨集，結果是字面上的 \texttt{"VERSION"}；
\texttt{XSTR(VERSION)} 多繞一層：\texttt{XSTR} 的參數 \texttt{x} 在它的
替換內容（\texttt{STR(x)}）裡沒有直接被 \texttt{\#} 或 \texttt{\#\#}
用到，所以會先正常展開——\texttt{VERSION} 先變成 \texttt{3}，再把展開後
的 \texttt{3} 傳給 \texttt{STR}，才字串化成 \texttt{"3"}。 這個「多一層
間接讓參數先展開」的寫法（\texttt{XSTR} 包一層 \texttt{STR}）是標準慣
用手法，任何時候想對一個「本身可能是巨集」的東西做字串化或貼接，都要先
繞這一層。

\keypoint{\texttt{\#} 把參數字串化，不會先求值；\texttt{\#\#} 把兩個 token
黏成一個新的識別字，跟單純用空白隔開是兩件事。 這兩個運算子只在巨集定義
裡有效，第~\ref{sec:pp-assert}~節自製 \texttt{assert} 會用到 \texttt{\#}。
被 \texttt{\#}／\texttt{\#\#} 用到的參數不會先展開，要先展開再處理就得
像 \texttt{XSTR} 包 \texttt{STR} 那樣多繞一層間接。}

\begin{drills}[label=drills:pp-stringize]
\item \texttt{STR(hello)} 與 \texttt{STR(1 + 2)} 分別展開成什麼字串常數？
\item \texttt{PRINT\_VAR(count)} 為什麼可以同時印出變數名字又印出它的
值？拆開巨集的兩個部分說明。
\item \texttt{CONCAT(my, Var)} 展開後變成什麼識別字？
\item 手寫：把 \texttt{CONCAT} 換成只用空白隔開（不用 \texttt{\#\#}）的
版本，\texttt{CONCAT\_SPACE(my, Var)} 會展開成什麼？跟 \texttt{a\#\#b} 的
版本差在哪裡？
\item \texttt{STR(VERSION)} 跟 \texttt{XSTR(VERSION)} 為什麼展開結果不
一樣？\texttt{XSTR} 多做了哪一步？
\end{drills}

\section{可變參數巨集}
\label{sec:pp-variadic}

函式式巨集的參數列跟函式一樣，本來個數是固定的。 但很多時候我們想要一個
巨集接受「不定數量」的參數——最典型的例子就是包一層 \texttt{printf}：訊息
格式化字串後面要接幾個值，事先不知道。 C99 引入 \texttt{...} 與
\texttt{\_\_VA\_ARGS\_\_} 解決這件事：

\begin{code}{c}{可變參數巨集的語法}
#define NAME(固定參數, ...) 替換內容，用 __VA_ARGS__ 代表 ... 的部分
\end{code}

\texttt{...} 放在參數列最後，代表「後面接任意個數的參數」；替換內容裡用
\texttt{\_\_VA\_ARGS\_\_} 這個特殊識別字，代表呼叫端實際傳進來、對應到
\texttt{...} 的那些引數，原封不動貼在這裡：

\begin{example}{c}{(log-variadic.c) 把 LOG 改成 printf 風格的可變參數版本}
#include <stdio.h>

#define LOG(fmt, ...) \
    fprintf(stderr, "[%s:%d] " fmt "\n", __FILE__, __LINE__, __VA_ARGS__)

int main(void) {
    int x = 10;
    LOG("x = %d", x);
    LOG("x doubled = %d", x * 2);
    return 0;
}
\end{example}

\iosample{}{[log-variadic.c:8] x = 10\\{[log-variadic.c:9] x doubled = 20}}

\texttt{LOG("x = \%d", x)} 展開時，\texttt{fmt} 對應到 \texttt{"x = \%d"}、
\texttt{\_\_VA\_ARGS\_\_} 對應到 \texttt{x}，整行變成
\texttt{fprintf(stderr, "[\%s:\%d] " "x = \%d" "\textbackslash n", \_\_FILE\_\_,
\_\_LINE\_\_, x)}——注意 \texttt{"[\%s:\%d] "} 跟 \texttt{fmt} 展開出來的
\texttt{"x = \%d"} 是兩個相鄰的字串常數，C 會自動把相鄰字串接成一個，等
價於一路寫好的格式字串。 這跟第~\ref{sec:pp-funcmacro}~節的一般參數行為
完全一致，只是 \texttt{\_\_VA\_ARGS\_\_} 一次代表好幾個引數，而不是一
個。

\subsection{呼叫端至少要傳一個變動參數}

嚴格的 C99／C11 規定 \texttt{...} 至少要對應到一個引數——也就是說
\texttt{LOG("just a message")}（後面完全不接任何變動參數）在標準模式下編
不過，因為 \texttt{fmt} 後面的逗號沒有東西可以接。 GCC 與 Clang 都提供一
個常見的擴充寫法 \texttt{, \#\#\_\_VA\_ARGS\_\_}：當變動參數是空的，
\texttt{\#\#} 會連同前面那個逗號一起吞掉，讓零個變動參數也能編過（C23 用
\texttt{\_\_VA\_OPT\_\_} 標準化了同一件事）。 這一章的範例都保證至少傳一
個變動參數，不需要用到這個擴充。

\keypoint{\texttt{...} 放在參數列最後，\texttt{\_\_VA\_ARGS\_\_} 在展開時
被換成呼叫端實際傳入、對應 \texttt{...} 的所有引數。 標準要求至少傳一個
變動參數；想支援零個，要用編譯器擴充 \texttt{\#\#\_\_VA\_ARGS\_\_} 或
C23 的 \texttt{\_\_VA\_OPT\_\_}。}

\begin{drills}[label=drills:pp-variadic]
\item 手寫：\texttt{LOG("x = \%d", x)} 展開後大致長什麼樣子？
\texttt{\_\_VA\_ARGS\_\_} 被換成了什麼？
\item 如果呼叫 \texttt{LOG("just a message")}（後面沒有任何變動參數），
在嚴格的 C11 底下會發生什麼事？
\item \texttt{\_\_VA\_ARGS\_\_} 跟一般巨集參數（像 \texttt{SQUARE} 的
\texttt{x}）比起來，差在哪裡？
\item 補空：想讓巨集像 \texttt{printf} 一樣接受不定數量的參數，參數列最
後要寫 \underline{\hspace{2em}}，展開內容裡用 \underline{\hspace{2em}}
代表呼叫端實際傳入的東西。
\end{drills}

\section{條件編譯}
\label{sec:pp-cond}

條件編譯讓前置處理器根據某個巨集有沒有被定義、或某個常數表達式是不是為
真，選擇性地保留或丟掉一段程式碼。 這一節最重要的應用，是解決第
\ref{sec:pp-include} 節留下的問題：同一個標頭檔被間接 \texttt{\#include}
兩次以上，會發生什麼事。

\subsection{Include guard}

先看沒有防護的情況。 標頭檔裡放了一個 \texttt{struct} 定義：

\begin{code}{c}{guard\_test2.h（沒有 include guard）}
struct Point { int x, y; };
\end{code}

如果它被同一個 \texttt{.c} 檔 \texttt{\#include} 兩次（實務上常常是因為
\texttt{A.h} 和 \texttt{B.h} 都 \texttt{\#include} 了同一個
\texttt{C.h}，而某個 \texttt{.c} 檔同時用到 \texttt{A.h} 與
\texttt{B.h}），結果是這樣：

\begin{code}{text}{struct 被定義兩次}
error: redefinition of 'struct Point'
note: originally defined here
\end{code}

修法是在標頭檔最外層包一層條件編譯，讓內容最多只被前置處理器「看到」一
次：

\begin{example}{c}{(guard-test1.h) include guard}
#ifndef GUARD_TEST1_H
#define GUARD_TEST1_H

struct Point { int x, y; };

#endif
\end{example}

第一次 \texttt{\#include} 這個檔案時，\texttt{GUARD\_TEST1\_H} 還沒被定
義，\texttt{\#ifndef} 為真，於是定義 \texttt{GUARD\_TEST1\_H}、往下保留
\texttt{struct} 的定義。 第二次 \texttt{\#include} 同一個檔案時，
\texttt{GUARD\_TEST1\_H} 已經定義過了，\texttt{\#ifndef} 變成假，中間整段
（包括 \texttt{struct} 的定義）被前置處理器直接跳過，不會出現在展開結果
裡，也就不會有重複定義的問題。 這是本章到目前為止唯一一段值得背下來的固
定格式：每一個標頭檔都該包這一層。

網路上常見的 include guard 命名還有 \texttt{\_GUARD\_H\_}、
\texttt{\_\_GUARD\_H\_\_} 這種前後加底線的寫法，這裡刻意不用——之後
《C 標準程式庫》那一章會看到，底線開頭的識別字整批保留給編譯器與函式庫實
作使用，自己寫的 guard 巨集踩到這條規則不算違規（不會編不過），但已經算
是在跟實作可能用到的名字打架，不值得冒這個險。

大部分主流編譯器也支援一句非標準但幾乎通用的替代寫法：

\begin{code}{c}{\#pragma once}
#pragma once

struct Point { int x, y; };
\end{code}

\texttt{\#pragma once} 讓編譯器自己記住「這個檔案已經看過了」，效果跟
include guard 一樣，只是少打三行字；缺點是它不是 C 標準的一部分（雖然實
務上幾乎所有編譯器都支援）。 兩種寫法擇一即可，這本書後面統一用 include
guard，因為它是任何合格 C 編譯器都保證支援的寫法。

\keypoint{Include guard 是 \texttt{\#ifndef} + \texttt{\#define} +
\texttt{\#endif} 的固定寫法，保護標頭檔內容最多只被前置處理器展開一次。
每一個自己寫的標頭檔都該包上這一層，不是只有出問題才補；\texttt{\#pragma
once} 是效果相同的非標準捷徑。}

\subsection{常見的條件編譯指令}

\begin{code}{c}{條件編譯指令}
#ifdef SYMBOL      // 如果 SYMBOL 有被 #define 過
#ifndef SYMBOL     // 如果 SYMBOL 沒有被 #define 過
#if CONDITION      // 如果常數表達式 CONDITION 為真
#elif CONDITION    // 前面都不成立、且這個為真
#else              // 前面都不成立
#endif             // 結束這一段條件編譯
\end{code}

\texttt{\#ifdef X} 等價於 \texttt{\#if defined(X)}；把 \texttt{defined()}
寫成函式呼叫的樣子，是為了能跟其他條件用 \texttt{\&\&}／\texttt{||} 組合，
例如 \texttt{\#if defined(DEBUG) \&\& !defined(NDEBUG)}。

\subsection[\#if 只算整數常數]{\texttt{\#if} 只做整數常數運算，未定義的識別字算 0}

\texttt{\#if} 後面的條件必須是一個\textbf{整數常數表達式}——前置處理器不
懂變數、不懂函式呼叫，只能算加減乘除、比較、邏輯運算這些純整數運算。 更
危險的規則是：如果條件裡出現一個\textbf{根本沒被 \texttt{\#define} 過}的
識別字，前置處理器不會報錯，而是把它當成 \texttt{0}：

\begin{example}{c}{(if-undefined-demo.c) 誤打錯字或忘記定義，靜靜地算成 0}
#include <stdio.h>

int main(void) {
#if UNDEFINED_SYMBOL > 0
    printf("branch A\n");
#else
    printf("branch B\n");
#endif
    return 0;
}
\end{example}

\iosample{}{branch B}

\texttt{UNDEFINED\_SYMBOL} 從頭到尾沒有被 \texttt{\#define} 過，
\texttt{\#if UNDEFINED\_SYMBOL > 0} 被當成 \texttt{\#if 0 > 0}，結果是假，
程式走進 \texttt{\#else}。 一般編譯選項（連 \texttt{-Wall -Wextra} 都不
會）\textbf{不會警告這件事}——如果你把巨集名字打錯字，或者忘了在某個標
頭檔裡 \texttt{\#define} 它，程式照樣默默編過、默默走錯分支，除錯起來很難
察覺。 想讓編譯器主動抓這個地雷，要另外加上 \texttt{-Wundef}：

\begin{code}{bash}{用 -Wundef 抓未定義識別字}
gcc -Wundef if-undefined-demo.c -o if-undefined-demo
\end{code}

\begin{code}{text}{-Wundef 的警告}
if-undefined-demo.c:4:5: warning: 'UNDEFINED_SYMBOL' is not defined,
evaluates to 0 [-Wundef]
\end{code}

\subsection{平台分支}

另一個常見用途是根據編譯器內建定義的平台巨集，切換不同平台的程式碼：

\begin{example}{c}{(platform.c) 平台判斷}
#include <stdio.h>

int main(void) {
#ifdef __linux__
    printf("Running on Linux\n");
#elif defined(__APPLE__)
    printf("Running on macOS\n");
#elif defined(_WIN32)
    printf("Running on Windows\n");
#else
    printf("Running on other platform\n");
#endif
    return 0;
}
\end{example}

\iosample{}{Running on macOS}

注意巨集名字是 \texttt{\_\_linux\_\_}，前後各兩條底線，全部小寫；沒有這個
巨集叫 \texttt{\_\_LINUX\_\_}。 這些平台巨集是編譯器自己內建定義的，不需
要（也不應該）自己 \texttt{\#define} 它們；如果打錯成 \texttt{\_\_LINUX\_\_}
這種不存在的名字，就正好踩到上一小節的地雷——它會被靜靜當成 \texttt{0}，
永遠走到 \texttt{\#else}，而且不會有任何警告（除非加了 \texttt{-Wundef}）。
另外 \texttt{\_WIN32} 這個名字容易讓人誤會——它在 32 位元\textit{與} 64
位元 Windows 上都會被定義（是「Windows API」的意思，不是「32 位元」），
不要以為它只代表 32 位元平臺就另外去找一個不存在的「64 位元判斷」。

\subsection{除錯開關}

條件編譯也常用來整段打開或關閉除錯輸出：

\begin{code}{c}{除錯開關}
#define DEBUG 1

#if DEBUG
    printf("Debug mode enabled\n");
#endif
\end{code}

跟 \texttt{if (DEBUG) printf(...)} 不同：\texttt{\#if} 是前置處理期就決定
要不要保留這段程式碼，\texttt{DEBUG} 為假的分支根本不會被編譯器看到；一般
的 \texttt{if} 是執行期才判斷，兩個分支都會被編譯進執行檔。 想要完全排除
一段程式碼（連編譯都不編）、或整段程式碼在某些平台上根本不合語法，只能用
條件編譯，不能用 \texttt{if}。 同樣的手法還能反過來用：把一大段暫時不想
編譯、但還不想刪掉的程式碼包在 \texttt{\#if 0 ... \#endif} 裡，等於一種
「編譯期註解」。

\keypoint{\texttt{\#if}／\texttt{\#ifdef} 是前置處理期的分支，不合條件的
分支連編譯器都看不到；一般的 \texttt{if} 是執行期分支，兩邊都會被編譯。
\texttt{\#if} 只做整數常數運算，\textbf{未定義的識別字會被當成 0}，預設
不會有警告——這是本章最容易被忽略的地雷，養成打錯字或漏
\texttt{\#define} 就用 \texttt{-Wundef} 檢查的習慣。}

\begin{drills}[label=drills:pp-cond]
\item \texttt{guard-test1.h} 被同一個 \texttt{.c} 檔 \texttt{\#include}
兩次，為什麼沒有出錯？guard 裡的三行分別做了什麼？
\item \texttt{guard-test2.h}（沒有 guard）被 \texttt{\#include} 兩次會出
現什麼錯誤，是編譯期還是連結期？
\item \texttt{\#if UNDEFINED\_SYMBOL > 0} 在沒有先 \texttt{\#define
UNDEFINED\_SYMBOL} 的情況下，會被前置處理器當成什麼值？預設編譯會不會
警告？
\item 想讓編譯器主動警告「用到未定義的識別字」，要加上哪個旗標？
\item 平台判斷那段程式碼在 macOS 上為什麼走 \texttt{\_\_APPLE\_\_} 那個分
支？
\item 補空：include guard 巨集的名字通常取標頭檔名的大寫加底線，例如
\texttt{mathutils.h} 對應 \underline{\hspace{2em}}。
\end{drills}

\section{命令列與預定義巨集}
\label{sec:pp-cmdline}

\subsection{命令列 \texttt{-D} 與 \texttt{-U}}

\texttt{-D} 讓你在編譯時從命令列定義巨集，不用改原始碼；\texttt{-U} 則是
取消某個巨集的定義。 這是條件編譯真正發揮威力的地方——同一份原始碼，靠
編譯旗標切換出不同版本。

\begin{example}{c}{(cmdline-define.c) 命令列決定要不要優化、要哪個版本}
#include <stdio.h>

int main(void) {
#ifdef OPTIMIZE
    printf("Optimized version\n");
#else
    printf("Debug version\n");
#endif

#ifdef VERSION
    printf("Version: %d\n", VERSION);
#endif
    return 0;
}
\end{example}

\begin{code}{bash}{不同的編譯旗標}
gcc cmdline-define.c -o prog                     # 不加旗標
gcc -DOPTIMIZE -DVERSION=2 cmdline-define.c -o prog2
\end{code}

\iosample{}{Debug version}

加了 \texttt{-DOPTIMIZE -DVERSION=2} 之後，輸出變成 \texttt{Optimized
version} 加 \texttt{Version: 2}。

最常見的真實用法是 \texttt{NDEBUG}：C 標準函式庫的 \texttt{assert} 只要
偵測到 \texttt{NDEBUG} 被定義，就會把自己展開成幾乎什麼都不做，整批關掉斷
言檢查——開發時開著斷言抓錯，發布版用 \texttt{-DNDEBUG} 一次關掉，不用手
動刪程式碼裡的 \texttt{assert}。 第~\ref{sec:pp-assert}~節會自己寫一個一
模一樣行為的版本。

還有一個常跟 \texttt{-D}／\texttt{-U} 一起用到的旗標：\texttt{-I目錄}，
讓 \texttt{\#include "..."} 除了目前檔案所在目錄，額外搜尋這個目錄。
第~\ref{sec:pp-include}~節的 \texttt{mathutils.h} 例子把標頭檔跟原始碼放
在同一層，不需要 \texttt{-I}；一旦專案變大、標頭檔搬進獨立的
\texttt{include/} 資料夾，就要靠 \texttt{-I} 告訴編譯器去哪裡找：

\begin{code}{bash}{-I 指定額外的標頭檔搜尋路徑}
gcc main.c mathutils.c -Iinclude -o main
\end{code}

\keypoint{\texttt{-D名字}／\texttt{-D名字=值} 從命令列定義巨集，
\texttt{-U名字} 取消定義，都不用改原始碼。 \texttt{-DNDEBUG} 是關閉
\texttt{assert} 的標準做法；\texttt{-I目錄} 讓 \texttt{\#include "..."}
額外搜尋指定的目錄。}

\subsection{預定義巨集}

編譯器會自動定義一批巨集，不用自己 \texttt{\#define}，直接可以用：

\begin{code}{c}{常見的預定義巨集}
__FILE__            // 目前檔案名字（字串）
__LINE__            // 目前行號（整數）
__DATE__            // 編譯日期（字串）
__TIME__            // 編譯時間（字串）
__STDC_VERSION__    // 這個編譯器遵循的 C 標準版本（整數）
\end{code}

\begin{example}{c}{(predefined.c) 印出所有預定義巨集}
#include <stdio.h>

int main(void) {
    printf("File: %s\n", __FILE__);
    printf("Line: %d\n", __LINE__);
    printf("Date: %s\n", __DATE__);
    printf("Time: %s\n", __TIME__);
    printf("STDC_VERSION: %ld\n", __STDC_VERSION__);
    printf("func: %s\n", __func__);
    return 0;
}
\end{example}

\iosample{}{File: predefined.c\\Line: 5\\Date: Sep  1 2026\\Time: 11:30:15\\STDC\_VERSION: 201112\\func: main}

\texttt{\_\_STDC\_VERSION\_\_} 印出 \texttt{201112} 代表這個編譯器目前用
C11（\texttt{201710} 是 C17，\texttt{202311} 是 C23，依實際編譯選項而
定）。 \texttt{\_\_DATE\_\_}／\texttt{\_\_TIME\_\_} 是編譯當下的日期時間，
每次重新編譯都會不一樣，跟原始碼本身沒有關係。 \texttt{\_\_LINE\_\_} 算的
是它\textit{被展開的那一行}，如果它出現在另一個巨集（例如下一節的
\texttt{LOG}）裡面，算的是\textbf{呼叫端}那一行，不是巨集定義那一行——這
正是它能用來做日誌定位的原因。

\subsection{\texttt{\_\_func\_\_} 不是巨集}

\texttt{\_\_func\_\_} 看起來跟其他雙底線名字很像，但它其實是 C99 引入的
一個\textbf{識別字}，由編譯器在每個函式一進來就自動填好一個字串，不是前
置處理器展開出來的：

\begin{code}{c}{驗證 \_\_func\_\_ 不是巨集}
#ifdef __func__
    // 不會走到這裡：__func__ 對前置處理器來說根本不存在
#else
    // 會走到這裡：它是識別字，不是巨集
#endif
\end{code}

差別在於：巨集在前置處理階段就被整段換掉，\texttt{\#ifdef} 能看到它；
\texttt{\_\_func\_\_} 要等到編譯階段才被賦值，前置處理器完全不知道有這個
東西存在，所以 \texttt{\#ifdef \_\_func\_\_} 永遠是假。

\subsection{其他兩個指令：\texttt{\#error} 與 \texttt{\#line}}

\texttt{\#error 訊息} 讓前置處理器主動中止編譯、印出這則訊息，常用來檢查
前提條件（例如「這支程式需要 C99 以上」）；\texttt{\#line 數字 "檔名"} 能
覆寫接下來 \texttt{\_\_LINE\_\_}／\texttt{\_\_FILE\_\_} 回報的行號與檔名，
主要給程式碼產生器（例如 yacc／bison）使用，手寫程式碼很少直接用到。

\keypoint{\texttt{\_\_FILE\_\_}／\texttt{\_\_LINE\_\_} 在巨集展開鏈裡永遠
指向呼叫端，不是巨集定義處，這是它們能用來定位錯誤的原因。
\texttt{\_\_func\_\_} 是識別字不是巨集，\texttt{\#ifdef} 抓不到它；
\texttt{\#error} 用來在編譯期主動中止並回報訊息。}

\begin{drills}[label=drills:pp-cmdline]
\item 不加任何 \texttt{-D} 旗標編譯 \texttt{cmdline-define.c}，輸出是什
麼？加 \texttt{-DOPTIMIZE -DVERSION=2} 之後呢？
\item \texttt{-DNDEBUG} 通常搭配哪個標準函式庫巨集使用？它做的事是什麼？
\item 補空：想要在編譯時「取消」某個巨集的定義，命令列旗標是
\underline{\hspace{2em}}。
\item \texttt{\_\_LINE\_\_} 如果出現在另一個巨集（例如 \texttt{LOG}）裡
面，印出來的行號是巨集定義那一行，還是呼叫端那一行？
\item \texttt{\_\_func\_\_} 是巨集還是識別字？這代表 \texttt{\#ifdef
\_\_func\_\_} 會是真還是假？
\item 補空：想知道原始碼檔名，用 \underline{\hspace{2em}}；想在編譯期主
動中止並印出一則錯誤訊息，用 \underline{\hspace{2em}}。
\end{drills}

\section{綜合：自己寫 \texttt{LOG} 與 \texttt{assert}}
\label{sec:pp-assert}

前面十節的零件——反斜線續行、函式式巨集、巨集還是函式的判準、\texttt{\#}、
可變參數巨集、條件編譯、\texttt{-DNDEBUG}、\texttt{\_\_FILE\_\_}／
\texttt{\_\_LINE\_\_}——湊在一起，剛好可以自己寫出兩個每個 C 專案幾乎都會
用到的工具：一個帶檔名行號、支援任意參數的 \texttt{LOG}，和一個仿造標準
函式庫 \texttt{assert} 行為的斷言巨集。 \texttt{LOG} 第~\ref{sec:pp-variadic}~
節已經寫過，這裡直接拿來用：

\begin{code}{c}{回顧：可變參數版的 LOG}
#define LOG(fmt, ...) \
    fprintf(stderr, "[%s:%d] " fmt "\n", __FILE__, __LINE__, __VA_ARGS__)
\end{code}

\subsection{裸 \texttt{if} 巨集踩到的 dangling else}

\texttt{LOG} 只展開成一個 \texttt{fprintf} 呼叫，語法上是單一個運算式陳
述，接在 \texttt{if}／\texttt{else} 後面很安全。 但接下來要寫的
\texttt{assert} 展開成好幾個陳述式（判斷、印訊息、中止程式），如果直接用
裸的 \texttt{if} 包起來，會踩到一個經典的坑：

\begin{example}{c}{(dangling-else.c) 裸 if 巨集接不上 else}
#include <stdio.h>

#define MY_ASSERT_BAD(cond) \
    if (!(cond)) { \
        printf("assertion failed\n"); \
    }

int main(void) {
    int x = 5;
    if (x > 0)
        MY_ASSERT_BAD(x > 10);
    else
        printf("x <= 0\n");
    return 0;
}
\end{example}

展開之後，\texttt{if (x > 0) MY\_ASSERT\_BAD(x > 10);} 那一行變成
\texttt{if (x > 0) if (!(x > 10)) \{...\} ; else ...}——真正的問題出在呼
叫端那個分號：\texttt{if (!(x > 10)) \{...\}} 本身已經是一句完整的陳述式
（大括號結尾不用再接分號），呼叫端寫的那個 \texttt{;} 變成緊接在後面的一
個\textbf{獨立空陳述式}。 於是外層 \texttt{if (x > 0)} 的陳述部分，到這個
空陳述式就結束了；後面的 \texttt{else} 往前找最近的 \texttt{if}，卻只找到
一個空陳述式擋在中間，沒有 \texttt{if} 可以接，編不過（不同編譯器的錯誤
訊息用字不盡相同，GCC 上是 \texttt{error: 'else' without a previous
'if'}）。

更危險的是相反的情形：如果呼叫端\textbf{不寫}分號——
\texttt{MY\_ASSERT\_BAD(x > 10)} 後面直接接 \texttt{else}——反而編得過：
外層 \texttt{if (x > 0)} 的陳述部分就是巨集展開出的
\texttt{if (!(x > 10)) \{...\}}，而 \texttt{else} 順理成章接在\textit{這
個內層} \texttt{if} 後面，不是接在外層 \texttt{if} 後面。 這台機器上編譯
器通常會用 \texttt{-Wdangling-else} 提醒「加括號避免歧義」，但這只是警
告，程式照樣編出來、照樣執行，而且執行結果完全不對——\texttt{x > 0} 卻走
進「\texttt{x <= 0}」的分支，是看不出來的\textbf{靜默語意錯誤}，比編不過
危險得多。 \texttt{do-while(0)} 真正要防的就是這個：不管呼叫端寫不寫分
號，都不該讓 \texttt{else} 有機會誤接到巨集內部的 \texttt{if}。

改成 \texttt{do \{ ... \} while (0)} 包起來，整段就變成語法上「單獨一
句」，\texttt{else} 就能接回外層的 \texttt{if}：

\begin{example}{c}{(dangling-else-fixed.c) do-while(0) 接得回 else}
#include <stdio.h>

#define MY_ASSERT_OK(cond) \
    do { \
        if (!(cond)) { \
            printf("assertion failed\n"); \
        } \
    } while (0)

int main(void) {
    int x = 5;
    if (x > 0)
        MY_ASSERT_OK(x > 10);
    else
        printf("x <= 0\n");
    return 0;
}
\end{example}

\iosample{}{assertion failed}

\subsection{自製 \texttt{assert}}

零件都到位了：\texttt{\#} 把條件字串化印出來、\texttt{\_\_FILE\_\_}／
\texttt{\_\_LINE\_\_} 定位、\texttt{do-while(0)} 保證語法安全、
\texttt{\#ifdef NDEBUG} 讓它在發布版整批消失：

\begin{example}{c}{(my-assert.c) 完整的自製 assert}
#include <stdio.h>
#include <stdlib.h>

#ifdef NDEBUG
    #define MY_ASSERT(cond) ((void)0)
#else
    #define MY_ASSERT(cond) \
        do { \
            if (!(cond)) { \
                fprintf(stderr, "Assertion failed: %s\n", #cond); \
                fprintf(stderr, "  at %s:%d\n", __FILE__, __LINE__); \
                abort(); \
            } \
        } while (0)
#endif

int main(void) {
    int x = 5;

    if (x > 0)
        MY_ASSERT(x < 100);
    else
        MY_ASSERT(x >= 0);

    MY_ASSERT(x > 0);
    printf("all assertions passed\n");

    MY_ASSERT(x > 10); /* 這一行會讓程式中止 */
    printf("unreachable\n");
    return 0;
}
\end{example}

一般編譯、直接執行：

\begin{code}{bash}{偵錯版本}
gcc my-assert.c -o my-assert && ./my-assert
\end{code}

\iosample{}{all assertions passed\\Assertion failed: x > 10\\\ \ at my-assert.c:28}

加 \texttt{-DNDEBUG} 重新編譯，\texttt{MY\_ASSERT} 整批變成
\texttt{((void)0)}，條件根本不會被求值，程式一路跑到底：

\begin{code}{bash}{發布版本}
gcc -DNDEBUG my-assert.c -o my-assert-release && ./my-assert-release
\end{code}

\iosample{}{all assertions passed\\unreachable}

跟標準函式庫真正的 \texttt{assert} 一模一樣的行為：\texttt{NDEBUG} 定義前
就 \texttt{\#include <assert.h>} 或改變 \texttt{NDEBUG}，行為都會跟著變，
因為 \texttt{assert.h} 本身也是用同一套條件編譯手法寫的——所以實務上直接
用標準的 \texttt{assert} 就好，這裡自己動手寫一次是為了把整章的零件串起
來看。

\subsection{回到起點：用 \texttt{gcc -E} 看整段展開}

把 \texttt{LOG} 和 \texttt{MY\_ASSERT} 放進同一支程式，用第
\ref{sec:pp-what} 節學過的 \texttt{gcc -E}，可以直接看到前置處理器真正產
生的翻譯單元：

\begin{example}{c}{(debug-demo.c) LOG 與 MY\_ASSERT 放進同一支程式}
#include <stdio.h>
#include <stdlib.h>

#define LOG(fmt, ...) \
    fprintf(stderr, "[%s:%d] " fmt "\n", __FILE__, __LINE__, __VA_ARGS__)

#ifdef NDEBUG
    #define MY_ASSERT(cond) ((void)0)
#else
    #define MY_ASSERT(cond) \
        do { \
            if (!(cond)) { \
                fprintf(stderr, "Assertion failed: %s\n", #cond); \
                fprintf(stderr, "  at %s:%d\n", __FILE__, __LINE__); \
                abort(); \
            } \
        } while (0)
#endif

int main(void) {
    int x = 10;
    LOG("x starts at %d", x);

    MY_ASSERT(x > 0);
    x = x * 2;
    LOG("x doubled to %d", x);

    MY_ASSERT(x < 15); /* 這一行會讓程式中止 */
    printf("unreachable\n");
    return 0;
}
\end{example}

\iosample{}{[debug-demo.c:22] x starts at 10\\{[debug-demo.c:26] x doubled to 20}\\Assertion failed: x < 15\\\ \ at debug-demo.c:28}

\begin{code}{text}{gcc -E 輸出（節錄 main，示意用的原始 C 陳述式）}
int main(void) {
    int x = 10;
    fprintf(stderr, "[%s:%d] " "x starts at %d" "\n", "debug-demo.c", 22, x);

    do { if (!(x > 0)) { fprintf(stderr, "Assertion failed: %s\n", "x > 0");
        fprintf(stderr, "  at %s:%d\n", "debug-demo.c", 24); abort(); } } while (0);
    x = x * 2;
    fprintf(stderr, "[%s:%d] " "x doubled to %d" "\n", "debug-demo.c", 26, x);

    do { if (!(x < 15)) { fprintf(stderr, "Assertion failed: %s\n", "x < 15");
        fprintf(stderr, "  at %s:%d\n", "debug-demo.c", 28); abort(); } } while (0);
    printf("unreachable\n");
    return 0;
}
\end{code}

每一個看起來像函式呼叫的巨集，展開後都變成一長串沒有巨集名字的原始
C 陳述式——\texttt{\_\_FILE\_\_}、\texttt{\_\_LINE\_\_}、\texttt{\#cond}、
\texttt{\_\_VA\_ARGS\_\_} 全部被換成具體的文字。 （這裡的 \texttt{stderr}
是刻意保留的識別字寫法：有些平臺的 C 函式庫把 \texttt{stderr} 本身也定義
成巨集，實際在你自己的機器上跑 \texttt{gcc -E}，看到的可能是那個平臺內部
的替代名字，不影響這裡要示範的重點。）這正好呼應第
\ref{sec:pp-what} 節的第一句話：前置處理器只做文字替換，替換完之後，
\texttt{LOG} 和 \texttt{MY\_ASSERT} 這兩個名字就徹底從翻譯單元裡消失了。

\keypoint{\texttt{\#cond} 把條件字串化印出來、\texttt{do-while(0)} 讓多行
巨集在語法上安全、\texttt{\#ifdef NDEBUG} 讓斷言可以整批關閉、
\texttt{\_\_VA\_ARGS\_\_} 讓 \texttt{LOG} 接受任意參數——這四件事合起來
就是標準函式庫 \texttt{assert} 與常見 log 巨集的完整實作套路。}

\begin{drills}[label=drills:pp-assert]
\item \texttt{MY\_ASSERT(x > 10)} 失敗時印出的兩行，分別用到巨集裡的哪個
運算子／哪個預定義巨集？
\item 用 \texttt{-DNDEBUG} 重新編譯 \texttt{my-assert.c}，同一行
\texttt{MY\_ASSERT(x > 10)} 還會中止程式嗎？為什麼？
\item \texttt{dangling-else.c} 為什麼會出現「\texttt{else} 沒有對應的
\texttt{if}」這個錯誤？
\item 手寫：把 \texttt{MY\_ASSERT\_BAD} 換成 \texttt{do-while(0)} 版本，
為什麼 \texttt{else} 就接得回去了？
\item 如果呼叫端\textbf{不寫}分號（\texttt{MY\_ASSERT\_BAD(x > 10)} 後
面直接接 \texttt{else}），這段程式編不編得過？如果編得過，\texttt{else}
實際上接到哪個 \texttt{if}？
\end{drills}

\section{常見錯誤清單}
\label{sec:pp-checklist}

下面每一條都對應前面某一節。

\begin{itemize}
\item 前置處理器輸出的翻譯單元還是合法 C 原始碼；\texttt{\#define} 只是
文字替換，替換前的名字之後就不存在了。
（第~\ref{sec:pp-what}~節）
\item 指令從 \texttt{\#} 到換行結束，反斜線續行後面不能有多餘空白，否則
續行失敗；巨集替換比對的是 token，不是子字串。
（第~\ref{sec:pp-syntax}~節）
\item 標頭檔放宣告不放定義，否則連結期會撞名。
（第~\ref{sec:pp-include}~節）
\item 陣列大小要編譯期常數，選 \texttt{\#define} 或 \texttt{enum}；\\
型別別名用 \texttt{typedef}，不要用 \texttt{\#define}。\\
（第~\ref{sec:pp-define}~節）
\item 函式式巨集的參數與整個展開式都要加括號；參數有副作用會被重複求
值，是未定義行為。
（第~\ref{sec:pp-funcmacro}~節）
\item 能用 \texttt{enum}／\texttt{typedef}／\texttt{static inline} 解決
的，不要用巨集；巨集只留給 token 操作、可變參數、條件編譯、與懶惰求值。\\
（第~\ref{sec:pp-guidelines}~節）
\item \texttt{\#} 字串化不會先求值；\texttt{\#\#} 才是真正的 token 貼
接，跟空白隔開不同。
（第~\ref{sec:pp-stringize}~節）
\item 呼叫可變參數巨集時，標準要求至少傳一個變動參數，\\
零個變動參數要靠 \texttt{\#\#\_\_VA\_ARGS\_\_} 這類編譯器擴充才能編過。
（第~\ref{sec:pp-variadic}~節）
\item 每個自己寫的標頭檔都要包 include guard；平台巨集是
\texttt{\_\_linux\_\_}，全小寫，不是 \texttt{\_\_LINUX\_\_}；\texttt{\#if}
對未定義的識別字視為 \texttt{0}，預設不會警告，要加 \texttt{-Wundef} 才
抓得到。\\
（第~\ref{sec:pp-cond}~節）
\item \texttt{-D名字}／\texttt{-D名字=值} 定義巨集，\texttt{-U名字} 取消
定義；\texttt{\_\_LINE\_\_}／\texttt{\_\_FILE\_\_} 在巨集裡指向呼叫端；
\texttt{\_\_func\_\_} 是識別字不是巨集。
（第~\ref{sec:pp-cmdline}~節）
\item 多行函式式巨集要包 \texttt{do \{ ... \} while (0)}，\\
否則裸 \texttt{if} 會踩到 dangling else。
（第~\ref{sec:pp-assert}~節）
\end{itemize}

\begin{drills}[label=drills:pp-checklist]
\item 下面這段有四個錯，請全部標出。
\begin{verbatim}
#define SQUARE(x) x * x
#ifdef __LINUX__
    #define PLATFORM "linux"
#endif
#if BUILD_MODE
    #define MODE "release"
#endif
#define LOG_ON(msg) if (1) printf(msg)
\end{verbatim}
\item 指出錯誤：\texttt{\#define IntPtr int *} 後面接
\texttt{IntPtr p1, p2;}。
\item 指出錯誤：標頭檔裡直接寫函式本體（不是宣告），被兩個 \texttt{.c}
檔一起 \texttt{\#include}。
\item 圈出哪幾行會真的進入翻譯單元：
\begin{verbatim}
#define LEVEL 2
#if LEVEL >= 3
    printf("high\n");
#elif LEVEL >= 1
    printf("mid\n");
#else
    printf("low\n");
#endif
\end{verbatim}
\end{drills}

\phantomsection
\addcontentsline{toc}{section}{練習解答}
\section*{練習解答}

各節練習的答案收在這裡。 手寫展開類的題目只給展開後的結果；完整程式見
正文片語／範例。

\paragraph{第~\ref{sec:pp-what}~節}
\begin{enumerate}
\item \texttt{gcc -E hello.c} 印出幾百行，遠比一行 \texttt{printf} 多；
多出來的內容是 \texttt{\#include <stdio.h>} 展開出的一大串宣告。
\item 是。 前置處理器的輸出（翻譯單元）仍是合法的 C 原始碼，會直接餵給
編譯器繼續處理。
\item 編譯階段：\texttt{printf} 沒有先被宣告。 較新的編譯器（GCC 14／
Clang 16 這個等級以上）預設直接把「隱含宣告函式」（\texttt{implicit
declaration of function}）當成錯誤；較舊的編譯器只給警告、照樣編過連結成
功。 但不管是哪一種，呼叫一個沒有原型的函式，本身都是未定義行為。
\item \texttt{3+3}，也就是 \texttt{TWICE(3)} 展開成 \texttt{3+3}（不是
\texttt{6}，前置處理器不會幫你先算出來，但 \texttt{3+3} 這個運算式的值是
\texttt{6}）。
\end{enumerate}

\paragraph{第~\ref{sec:pp-syntax}~節}
\begin{enumerate}
\item 兩個錯誤，第二個更早被編譯器抓到：\texttt{BIG\_NUMBER} 會變成沒有
替換內容的空巨集，下一行單獨的 \texttt{123456} 不再屬於任何指令，變成檔
案作用域裡一段孤立的數字，不是合法的頂層宣告，編不過；更直接的是使用
\texttt{BIG\_NUMBER} 的那一行——\texttt{printf("\%d\textbackslash n",
BIG\_NUMBER);} 因為 \texttt{BIG\_NUMBER} 展開成空，變成
\texttt{printf("\%d\textbackslash n", );}，多出一個逗號後面沒有引數，這
是更直接、通常也更早被編譯器指出來的語法錯誤。
\item \texttt{\#} 前面的空白只是縮排，跟指令本身無關；\texttt{\#} 後面如
果先出現其他程式碼，那一行就不再符合「以 \texttt{\#} 開頭的指令」這個語
法規則，會被當成普通程式碼解讀。
\item 不會。 \texttt{N100} 是單一個識別字 token，跟 \texttt{N} 無關；
\texttt{"N"} 裡的 \texttt{N} 是字串常數的一部分，前置處理器不會鑽進字串
內部替換。
\item 換行字元；反斜線 \texttt{\textbackslash}。
\end{enumerate}

\paragraph{第~\ref{sec:pp-include}~節}
\begin{enumerate}
\item \texttt{<>} 先找系統／編譯器內建的標頭檔搜尋路徑；\texttt{""} 先從
目前檔案所在目錄找，找不到才退到系統路徑。
\item 連結期。 個別編譯 \texttt{a.c}、\texttt{b.c} 都不會出錯，因為前置
處理器只是把定義貼進兩個檔案、編譯器看到的也都是合法的函式定義；要到連
結器要把兩個目的檔接起來、發現兩份 \texttt{square} 本體撞名時才會出錯。
\item 標頭檔互相 \texttt{\#include}、或多個標頭檔被同一個 \texttt{.c}
檔一起用到時，同一個標頭檔可能被間接展開超過一次。
\item 宣告；定義。
\end{enumerate}

\paragraph{第~\ref{sec:pp-define}~節}
\begin{enumerate}
\item 不會。 內容一字不差的重新定義是合法的，不會有警告。
\item \texttt{array\_size} 是 \texttt{const} 變數，不是整數常數表達式；
檔案作用域的陣列大小必須是真正的常數表達式，錯誤訊息說它是「檔案作用域
的可變長度陣列」（variably modified at file scope）。
\item \texttt{INT\_PTR p1, p2;} 展開成 \texttt{int * p1, p2;}，星號只黏
著 \texttt{p1}；\texttt{p2} 只是宣告了一個 \texttt{int}。
\item \texttt{enum}（\texttt{enum \{ NAME = value \};}）。
\end{enumerate}

\paragraph{第~\ref{sec:pp-funcmacro}~節}
\begin{enumerate}
\item 展開成 \texttt{2 * 3+4}，也就是 \texttt{2 * 3 + 4 = 10}。
\item \texttt{\#define DOUBLE(x) (2 * (x))}；\texttt{DOUBLE(3+4)} 展開成
\texttt{(2 * (3+4))}，等於 \texttt{14}。
\item 未定義行為：\texttt{(i++) * (i++)} 對同一個變數做了兩次沒有序列點
隔開的修改，這台機器剛好算出 \texttt{30} 且 \texttt{i} 變成 \texttt{7}，
換一個編譯器結果可能不同。
\item \texttt{static inline} 函式。
\item \texttt{TWICE} 名字後面有空白，是物件式巨集，替換內容是整段
\texttt{(x) x+x}；\texttt{TWICE(3)} 展開成 \texttt{(x) x+x(3)}。
\end{enumerate}

\paragraph{第~\ref{sec:pp-guidelines}~節}
\begin{enumerate}
\item 藏了 \texttt{return -1;}；\texttt{b = -2} 讓
\texttt{CHECK\_POSITIVE(b)} 展開出的 \texttt{if} 成立，\texttt{check}
直接返回，後面的 \texttt{printf("checked b")} 與 \texttt{return 0} 都沒
機會執行。
\item 呼叫端要寫 \texttt{if (!check\_positive(b)) return -1;}，多一行，
但換來的是一眼就看得到「這裡可能會提前返回」，不用去翻巨集定義。
\item 重複求值配 \texttt{SQUARE(i++)}；型別不安全配 \texttt{IntPtr}（星
號只黏 \texttt{p1}，型別悄悄跟著出錯）；隱藏控制流程配
\texttt{CHECK\_POSITIVE}。
\item 條件編譯；\texttt{\_\_FILE\_\_}／\texttt{\_\_LINE\_\_}（懶惰求值／
呼叫端位置）。
\end{enumerate}

\paragraph{第~\ref{sec:pp-stringize}~節}
\begin{enumerate}
\item \texttt{STR(hello)} 展開成 \texttt{"hello"}；\texttt{STR(1 + 2)}
展開成 \texttt{"1 + 2"}（保留原始文字，不會先算出 \texttt{3}）。
\item \texttt{\#v} 把參數名字字串化印出來（顯示 \texttt{"count"}），另
一個獨立的 \texttt{v} 則是真正拿去求值的變數，兩者各做各的事。
\item \texttt{myVar}。
\item 展開成 \texttt{my Var}（中間隔著空白，是兩個 token），拿去宣告變
數會編不過；\texttt{a\#\#b} 才會把兩個 token 真正黏成一個新識別字
\texttt{myVar}。
\item \texttt{STR(VERSION)} 的參數被 \texttt{\#} 直接用到，不會先展開，
所以是字面上的 \texttt{"VERSION"}；\texttt{XSTR(VERSION)} 的參數在
\texttt{XSTR} 自己的替換內容裡沒有直接接 \texttt{\#}，會先正常展開成
\texttt{3}，再交給 \texttt{STR} 字串化，多做了「先讓 \texttt{VERSION}
展開成 \texttt{3}」這一步。
\end{enumerate}

\paragraph{第~\ref{sec:pp-variadic}~節}
\begin{enumerate}
\item 展開成 \texttt{fprintf(stderr, "[\%s:\%d] " "x = \%d" "\textbackslash n",
\_\_FILE\_\_, \_\_LINE\_\_, x)}；\texttt{\_\_VA\_ARGS\_\_} 被換成
\texttt{x}。
\item 標準 C11 底下編不過，因為 \texttt{...} 後面沒有對應到任何引數；要
用 GCC／Clang 的 \texttt{\#\#\_\_VA\_ARGS\_\_} 擴充或 C23 的
\texttt{\_\_VA\_OPT\_\_} 才能支援零個變動參數。
\item \texttt{\_\_VA\_ARGS\_\_} 一次代表任意個數的引數（照原樣整段貼
上），一般參數（像 \texttt{x}）固定只代表一個引數。
\item \texttt{...}；\texttt{\_\_VA\_ARGS\_\_}。
\end{enumerate}

\paragraph{第~\ref{sec:pp-cond}~節}
\begin{enumerate}
\item 第一次 \texttt{\#include} 時 \texttt{GUARD\_TEST1\_H} 還沒定義，
\texttt{\#ifndef} 為真，定義它並保留內容；第二次 \texttt{\#include} 時
\texttt{GUARD\_TEST1\_H} 已經定義過，\texttt{\#ifndef} 為假，中間內容被
整段跳過。
\item \texttt{redefinition of 'struct Point'}，編譯期錯誤（兩次
\texttt{\#include} 發生在同一次編譯裡，不是連結期）。
\item 被當成 \texttt{0}；預設不會警告。
\item \texttt{-Wundef}。
\item 因為 \texttt{\_\_APPLE\_\_} 是 macOS 上的編譯器內建定義的巨集，
\texttt{\_\_linux\_\_} 沒有被定義，\texttt{\#ifdef \_\_linux\_\_} 為假，
往下走到 \texttt{\#elif defined(\_\_APPLE\_\_)}。
\item \texttt{MATHUTILS\_H}。
\end{enumerate}

\paragraph{第~\ref{sec:pp-cmdline}~節}
\begin{enumerate}
\item 不加旗標輸出 \texttt{Debug version}；加 \texttt{-DOPTIMIZE
-DVERSION=2} 之後輸出 \texttt{Optimized version} 加 \texttt{Version: 2}。
\item \texttt{assert}；\texttt{NDEBUG} 定義後，\texttt{assert(...)} 會展
開成幾乎什麼都不做，關掉所有斷言檢查。
\item \texttt{-U名字}（例如 \texttt{gcc -UNDEBUG}）。
\item 呼叫端那一行，不是巨集定義那一行。
\item 識別字，不是巨集；\texttt{\#ifdef \_\_func\_\_} 永遠是假，因為前置
處理器根本不知道有這個東西。
\item \texttt{\_\_FILE\_\_}；\texttt{\#error}。
\end{enumerate}

\paragraph{第~\ref{sec:pp-assert}~節}
\begin{enumerate}
\item 第一行用 \texttt{\#cond} 把條件字串化；第二行用
\texttt{\_\_FILE\_\_} 與 \texttt{\_\_LINE\_\_} 定位。
\item 不會。 \texttt{NDEBUG} 定義後 \texttt{MY\_ASSERT} 被換成
\texttt{((void)0)}，條件根本不會被求值，也不會呼叫 \texttt{abort}。
\item 因為呼叫端那個分號變成了一個獨立的空陳述式：巨集展開出的
\texttt{if (!(x > 10)) \{...\}} 本身已經結束，後面的 \texttt{;} 是另一句
空陳述式，外層 \texttt{if (x > 0)} 的陳述部分到這裡就結束了，\texttt{else}
往前找不到可以接的 \texttt{if}。
\item \texttt{do \{ ... \} while (0)} 把整段包成語法上單獨一句，
\texttt{if (x > 0) MY\_ASSERT\_OK(...);} 只有一句陳述，後面的
\texttt{else} 能正確接回外層的 \texttt{if}。
\item 編得過。 這時候外層 \texttt{if (x > 0)} 的陳述部分就是巨集展開出
的 \texttt{if (!(x > 10)) \{...\}}，\texttt{else} 會接到\textit{這個內
層} \texttt{if}，不是原本想要的外層 \texttt{if}——程式照樣編過、照樣執
行，但語意悄悄錯了，是比編譯錯誤更危險的靜默 bug。
\end{enumerate}

\paragraph{第~\ref{sec:pp-checklist}~節}
\begin{enumerate}
\item 四個錯：\texttt{SQUARE(x)} 沒有括號（該寫成
\texttt{((x) * (x))}）；\texttt{\_\_LINUX\_\_} 應該是全小寫的
\texttt{\_\_linux\_\_}；\texttt{BUILD\_MODE} 沒有先 \texttt{\#define}，
\texttt{\#if BUILD\_MODE} 被當成 \texttt{\#if 0}，\texttt{MODE} 永遠不
會被定義；\\
\texttt{LOG\_ON} 展開成裸的 \texttt{if}，沒有
\texttt{do-while(0)}，會踩 dangling else。
\item \texttt{p2} 不是指標，是 \texttt{int}：\texttt{\#define} 展開後星
號只黏著 \texttt{p1}。
\item 個別編譯都不會出錯，連結時會因為兩份相同的函式本體撞名而失敗
（\texttt{duplicate symbol} 或 \texttt{multiple definition}）。
\item \texttt{LEVEL} 是 \texttt{2}，\texttt{LEVEL >= 3} 為假、
\texttt{LEVEL >= 1} 為真，所以只有 \texttt{\#elif} 分支裡的\\
\texttt{printf("mid\textbackslash n");} 那一行會真的進入翻譯單元，其餘
分支都被前置處理器整段跳過。
\end{enumerate}

\chapter{C 標準程式庫}
\label{ch:stdlib}

「精熟標準程式庫」不是背過很多函式名字，是三件可遷移的能力：看得懂一份函
式規格、認得失敗長什麼樣子、會用函式指標把資料型別和演算法拆開。 把這三件
事當骨架，本章出現的每一個函式都掛得上去；讀者遇到書上沒教的函式，也能靠
同一套骨架自己查手冊。

這一章分三段走。 第一段練「讀規格」的能力：標準程式庫的邊界在哪裡、怎麼拆
解一份函式原型、C 的錯誤回報有幾種家族、程式怎麼開始與結束。 第二段是全章
重心：用函式指標把演算法從資料型別裡拆出來，\texttt{qsort} 與
\texttt{bsearch} 是唯一的教學載體。 第三段快速巡覽字元／字串、數學、時
間、可變參數、國際化幾個分領域。 章末有一個把所有主題串起來的整合專案。
節末練習的答案收在章末。

\section{標準程式庫是什麼}
\label{sec:sl-what}

標準程式庫不是 C 語言本身的一部分，是隨編譯器附帶的一套函式與型別，C11
標準把它切成 29 個標頭檔（\texttt{<stdio.h>}、\texttt{<stdlib.h>}、
\texttt{<string.h>}⋯，比 C99 多了 \texttt{<stdalign.h>}、
\texttt{<stdatomic.h>}、\texttt{<stdnoreturn.h>}、\texttt{<threads.h>}、
\texttt{<uchar.h>} 這五個——本書其他地方用到的 \texttt{\_Static\_assert}、
\texttt{\_\_STDC\_VERSION\_\_} 站的都是 C11 的立場）。 標頭檔提供的是
\textbf{宣告}，函式真正的實作藏在
別處的函式庫檔案裡——這也是為什麼有些函式庫要在編譯時額外加旗標才連結得
起來：

\begin{code}{bash}{數學函式庫要額外連結}
gcc prog.c -o prog       # 用到 sqrt、pow 等函式，缺 -lm 可能連結失敗
gcc prog.c -o prog -lm   # -lm 明確連結數學函式庫
\end{code}

\texttt{<math.h>} 只是宣告，\texttt{sqrt}／\texttt{pow} 的本體在
\texttt{libm} 裡；\texttt{-lm} 告訴連結器去哪裡找這些本體。

\subsection{保留名稱}

C 標準替未來的標準程式庫版本、還有編譯器自己的內部用途，保留了一批名字
——這些名字讀者不該自己拿來用：

\begin{itemize}
\item 底線開頭的識別字（\texttt{\_foo}、\texttt{\_\_bar}）整批保留給實作
（編譯器與函式庫）使用。
\item \texttt{str} 加小寫字母開頭的名字（\texttt{strfoo}、
\texttt{strxyz}）保留給未來的標準函式庫擴充，即使 C 標準目前還沒有這個
函式。
\end{itemize}

自己寫的函式、變數、巨集名字，避開這兩類就好，不用背完整清單。

\subsection{巨集會遮蔽函式}

第~\ref{sec:pp-define}~節提過巨集是純文字替換。 標準程式庫允許某些函式
\textit{同時}有一個效果相同的巨集版本——常見的是 \texttt{getchar}：

\begin{example}{c}{(macro-hide.c) 用括號躲開巨集}
#include <stdio.h>

int main(void) {
    int c1 = getchar();       /* 可能是巨集展開，也可能是函式呼叫 */
    int c2 = (getchar)();     /* 一定是函式呼叫 */
    printf("c1=%d c2=%d\n", c1, c2);
    return 0;
}
\end{example}

巨集要展開，名字後面必須\textit{直接}跟著左括號；\texttt{(getchar)} 中間
被括號擋住，前置處理器找不到「\texttt{getchar} 緊接左括號」這個樣式，就
不會展開，於是編譯器看到的是一個普通的函式呼叫。 這一招（用括號把名字架
起來）在你懷疑某個標準函式庫名字被巨集蓋住、想強迫呼叫真正的函式時很有
用；另一招是 \texttt{\#undef getchar} 之後那個名字就只剩函式版本。

\keypoint{標頭檔只給宣告，函式庫檔案才有實作，缺少對應的連結旗標
（如 \texttt{-lm}）會連結失敗。 底線開頭與 \texttt{str+小寫} 是保留名
稱，不要自己用。 \texttt{(name)()} 或 \texttt{\#undef name} 能繞開遮蔽
函式的巨集版本。}

\begin{drills}[label=drills:sl-what]
\item 為什麼用到 \texttt{sqrt} 的程式，有些平台不加 \texttt{-lm} 會連結
失敗，而不是編譯失敗？
\item 自己想定義一個叫 \texttt{\_config} 的全域變數，為什麼不建議？
\item \texttt{(getchar)()} 為什麼一定是函式呼叫，不會被巨集展開？
\item 補空：想強迫呼叫某個可能被巨集遮蔽的標準函式庫函式，除了加括號，
還可以先 \underline{\hspace{2em}} 那個名字。
\end{drills}

\section{怎麼讀一份函式規格}
\label{sec:sl-spec}

標準程式庫的函式手冊看起來都長一個樣子：一行函式原型，加上幾句話說明參
數與回傳值。 讀熟一份規格，比背下所有函式更有用。 用 \texttt{strtol}
逐字拆解：

\begin{code}{c}{strtol 的原型}
long strtol(const char *restrict str, char **restrict endptr, int base);
\end{code}

\begin{itemize}
\item 第一個參數 \texttt{const char *}：只讀，函式不會修改這個字串。
\item 第二個參數 \texttt{char **}：一個「指標的指標」，用來讓函式\textit{回
傳第二個東西}——除了函式本身的回傳值，還要告訴呼叫端「解析到哪裡停下
來」。 呼叫端傳一個 \texttt{char *} 變數的位址進去，函式內部把這個變數
設成指向沒被解析到的第一個字元。
\item 兩個 \texttt{restrict} 各自限定它所在的那個指標參數，不是一句合起
來的承諾。 \texttt{str} 上的 \texttt{restrict} 承諾：函式執行期間，
\texttt{str} 指到的字串內容不會透過其他指標被修改。 \texttt{endptr} 上
的 \texttt{restrict} 承諾：\texttt{endptr} 指到的那個 \texttt{char *} 變
數（呼叫端拿來接住「解析到哪裡」的那個變數本身），不會透過其他指標被存
取。 兩個承諾各自獨立，讓編譯器可以放心優化這兩條各自的存取路徑
（\texttt{restrict} 這個關鍵字在很多標準函式庫原型裡都會出現，先認得它
的意思）。
\item 回傳型別是 \texttt{long}，不是 \texttt{int}：轉換出來的數字可能超
過 \texttt{int} 能裝的範圍，用更大的型別接住，呼叫端自己決定要不要縮
小、要不要檢查縮小會不會溢位。
\end{itemize}

\subsection{回傳值協定的四個家族}

C 沒有例外（exception）機制，函式回報失敗只能靠回傳值或修改某個全域狀
態。 整個標準程式庫的錯誤回報，可以分成四個家族，後面每一節出現的函式
都能對到其中一個：

\begin{center}
\begin{tabular}{p{2.6cm}p{4.0cm}p{5.2cm}}
\toprule
家族 & 代表 & 失敗長什麼樣 \\
\midrule
哨兵值 & \texttt{getchar}、\texttt{fgetc} & 回傳一個資料範圍外的特殊值
（\texttt{EOF}） \\
指標或 \texttt{NULL} & \texttt{fopen}、\texttt{bsearch}、\texttt{malloc}
& 回傳 \texttt{NULL} \\
計數 & \texttt{fread}、\texttt{printf} & 回傳值小於預期，或是負數 \\
設定 \texttt{errno} & \texttt{strtol}、\texttt{math.h} 的函式 & 回傳值
本身合法，但 \texttt{errno} 被設成非零 \\
\bottomrule
\end{tabular}
\end{center}

遇到一個沒教過的函式，先看它的回傳型別落在哪個家族，就知道該檢查什麼：
回傳型別是可能無效的指標，就檢查 \texttt{NULL}；回傳型別是計數，就跟預
期的數字比對；回傳值本身看不出失敗（例如 \texttt{strtol} 的
\texttt{0} 既可能是真的轉出 \texttt{0}，也可能是失敗），就要檢查
\texttt{errno}。 下一節細講 \texttt{errno} 這個機制。

\keypoint{看到一個新函式，先讀完整的參數型別（尤其是 \texttt{char **}、
\texttt{void *} 這類）與回傳型別，回傳型別會告訴你失敗長什麼樣，四個家
族（哨兵值／NULL／計數／errno）幾乎涵蓋全部標準函式庫。}

\begin{drills}[label=drills:sl-spec]
\item \texttt{strtol} 的第二個參數為什麼是 \texttt{char **} 而不是
\texttt{char *}？
\item \texttt{strtol} 回傳 \texttt{long} 而不是 \texttt{int}，用意是什
麼？
\item \texttt{fopen} 屬於哪一個回傳值家族？\texttt{fread} 呢？
\item 補空：\texttt{strtol("abc", \&endptr, 10)} 回傳 \texttt{0}，這個
\texttt{0} 有兩種可能：一種是字串裡真的解析出 \texttt{0}，另一種是
\underline{\hspace{2em}}。 光看回傳值分不出來，還要看
\underline{\hspace{2em}}。
\end{drills}

\section{錯誤處理}
\label{sec:sl-errors}

上一節提到的 \texttt{errno} 是一個標準函式庫定義的整數（技術上是一個
執行緒區域的左值），某些函式偵測到錯誤時會把它設成一個非零的代碼。 用
法有兩個規則容易忽略：

\begin{itemize}
\item \textbf{呼叫前先清零}：\texttt{errno} 不會在函式成功時被自動清成
\texttt{0}，只在某些函式偵測到錯誤時才被寫入。 上一次呼叫留下的舊值可
能還在，所以要檢查 \texttt{errno} 之前，先手動 \texttt{errno = 0;}。
\item \textbf{只在規格明確說會設的情況下才檢查}：不是每個函式失敗都會
設 \texttt{errno}，規格沒提到 \texttt{errno} 的函式，檢查它毫無意義。
\end{itemize}

\begin{code}{c}{errno 的正確用法}
errno = 0;
long v = strtol(s, &endptr, 10);
if (errno == ERANGE) {
    /* 真的溢位了 */
}
\end{code}

\texttt{perror(prefix)} 把 \texttt{errno} 目前的值翻成一句人看得懂的英
文訊息，印到 \texttt{stderr}；\\
\texttt{strerror(errno)} 則是把同一句訊息變成字串，自己決定要印在哪裡：

\begin{code}{c}{perror 與 strerror}
FILE *fp = fopen("missing.txt", "r");
if (!fp) {
    perror("fopen");                          /* fopen: No such file or directory */
    fprintf(stderr, "detail: %s\n", strerror(errno)); /* 同一句訊息 */
}
\end{code}

\subsection{\texttt{assert} 與 \texttt{NDEBUG}}

第~\ref{sec:pp-assert}~節已經自己寫過一個功能對等的版本：\texttt{assert(cond)}
在 \texttt{cond} 為假時印出訊息並呼叫 \texttt{abort}；只要在
\texttt{\#include <assert.h>} 之前定義了 \texttt{NDEBUG}，整個
\texttt{assert} 就變成幾乎不做事的空殼。 \texttt{assert} 檢查的是「這件
事按照程式邏輯不該發生」（程式設計者的錯誤），不是「使用者輸入可能有
誤」（那種情況要用回傳值或 \texttt{errno}，因為使用者輸入錯誤在發布版
也需要處理，不能被 \texttt{NDEBUG} 關掉）。

C11 加入 \texttt{\_Static\_assert(常數表達式, "訊息")}，在\textit{編譯
期}檢查一個常數條件，不像 \texttt{assert} 要等到執行期：

\begin{code}{c}{\_Static\_assert}
_Static_assert(sizeof(int) == 4, "this code assumes 32-bit int");
\end{code}

條件在編譯期就能算出來、又想在寫錯假設時直接編不過，用
\texttt{\_Static\_assert}；條件要等執行期才知道，用 \texttt{assert}。

\texttt{<signal.h>} 的 \texttt{signal} 用來註冊訊號處理函式，
\texttt{<setjmp.h>} 的 \texttt{setjmp}／\texttt{longjmp} 可以做非局部跳
轉——這兩個工具能處理的情境比本書涵蓋的範圍複雜（訊號處理函式裡能安全
做的事非常有限），這裡只記名字，之後真的需要再深入。

\keypoint{檢查 \texttt{errno} 前先清零，只在規格明確提到的情況下檢查。
\texttt{assert} 抓程式設計者的錯誤、可以被 \texttt{NDEBUG} 關掉；使用者
輸入錯誤要用回傳值處理，不能靠 \texttt{assert}。 \texttt{\_Static\_assert}
是編譯期版本。}

\begin{drills}[label=drills:sl-errors]
\item 為什麼檢查 \texttt{errno} 之前要先手動清零？
\item \texttt{perror} 跟 \texttt{strerror} 印出的內容有什麼關係？差別
在哪裡？
\item 使用者在 CLI 程式輸入了一個不存在的檔名，該用 \texttt{assert} 還
是回傳值處理？為什麼？
\item 補空：\texttt{\_Static\_assert} 在 \underline{\hspace{2em}} 期檢
查條件，\texttt{assert} 在 \underline{\hspace{2em}} 期檢查。
\end{drills}

\section{程式的開始與結束}
\label{sec:sl-startend}

\texttt{main} 的回傳值就是程式的結束狀態碼，\texttt{0}（或
\texttt{EXIT\_SUCCESS}）表示成功，非零（或 \texttt{EXIT\_FAILURE}）表示
失敗，殼層（shell）或呼叫這支程式的其他程式可以讀到這個值。
\texttt{exit(status)} 讓程式從\textit{任何地方}立刻結束，效果跟
\texttt{main} 執行 \texttt{return status;} 一樣。

\texttt{atexit(func)} 註冊一個「程式結束前要跑的函式」，可以註冊多個，
呼叫順序跟註冊順序\textbf{相反}：

\begin{example}{c}{(exit-demo.c) atexit 以反序呼叫}
#include <stdio.h>
#include <stdlib.h>

void cleanup1(void) { printf("cleanup1\n"); }
void cleanup2(void) { printf("cleanup2\n"); }

int main(void) {
    atexit(cleanup1);
    atexit(cleanup2);
    printf("Program running\n");
    exit(0);
}
\end{example}

\iosample{}{Program running\\cleanup2\\cleanup1}

\subsection{\texttt{exit} 與 \texttt{\_Exit} 的真實差別}

\texttt{exit} 在結束前會依序做三件事：呼叫所有 \texttt{atexit} 註冊的
函式、沖出（flush）所有開啟的串流緩衝區、關閉所有開啟的檔案。
\texttt{\_Exit} 跳過這三件事，直接結束程式。 差別不是敘述能說清楚的，
用會不會印出東西來看最直接：

\begin{example}{c}{(exit-vs-underscore-exit.c) 有沒有沖出緩衝區}
#include <stdio.h>
#include <stdlib.h>

int main(int argc, char **argv) {
    printf("about to exit"); /* 故意不加換行 */
    if (argc > 1 && argv[1][0] == 'E') {
        _Exit(0); /* 不沖出緩衝區、不呼叫 atexit 函式 */
    } else {
        exit(0); /* 沖出緩衝區、呼叫 atexit 函式 */
    }
    return 0;
}
\end{example}

\begin{code}{bash}{把標準輸出重新導向到檔案，逼緩衝區真的積著}
./prog > out1.txt; cat out1.txt      # about to exit
./prog E > out2.txt; cat out2.txt    # 空的
\end{code}

\texttt{stdout} 接到檔案時走全緩衝，\texttt{printf} 那一行沒有換行，內
容留在緩衝區裡還沒真正寫進檔案。 走 \texttt{exit} 的那一份，緩衝區在程
式結束前被沖出，\texttt{out1.txt} 裡看得到 \texttt{about to exit}；走
\texttt{\_Exit} 的那一份，緩衝區內容直接被丟棄，\texttt{out2.txt} 是空
的。 \texttt{\_Exit} 通常只用在確定不需要清理、想儘快結束的場合（例如
\texttt{fork} 出來的子行程裡發生無法處理的錯誤）。

\subsection{\texttt{abort}：不清理，直接中止}

\texttt{abort()} 送出 \texttt{SIGABRT} 訊號，讓程式異常終止，\textbf{不
會}呼叫 \texttt{atexit} 註冊的函式：

\begin{example}{c}{(abort-demo.c) abort 跳過 atexit}
#include <stdio.h>
#include <stdlib.h>

void cleanup(void) { printf("cleanup ran\n"); }

int main(void) {
    atexit(cleanup);
    printf("before abort\n");
    abort();
    printf("unreachable\n");
    return 0;
}
\end{example}

\iosample{}{before abort}

只印出 \texttt{before abort}，\texttt{cleanup ran} 沒有機會執行——第
\ref{sec:pp-assert}~節自製的 \texttt{assert} 失敗時就是呼叫 \texttt{abort}，
現在知道原因了：斷言失敗代表程式已經處在不可信任的狀態，\texttt{abort}
選擇「立刻停下」而不是「按正常程序清理後再走」。

\texttt{atexit} 不被呼叫這件事，C 標準說得很清楚（\texttt{abort} 屬於
「異常終止」，跟 \texttt{exit} 的「正常終止」不同）。 但\textbf{已開啟串
流裡還沒沖出的緩衝內容會不會保留下來}，標準只說是
\textit{implementation-defined}——不保證會、也不保證不會，換一個平台可
能不一樣。 沿用第~\ref{sec:sl-startend}~節 \texttt{exit} vs
\texttt{\_Exit} 的手法：故意不加換行、把輸出重新導向到檔案，實際看這台機
器怎麼選：

\begin{example}{c}{(abort-flush.c) abort 有沒有沖出緩衝區，是 implementation-defined}
#include <stdio.h>
#include <stdlib.h>

int main(void) {
    printf("before abort"); /* 故意不加換行 */
    abort();
    return 0;
}
\end{example}

\begin{code}{bash}{重新導向到檔案，逼緩衝區真的積著}
./abort-flush > out.txt; cat out.txt
\end{code}

\iosample{}{before abort}

在這台機器（macOS）上，\texttt{out.txt} 裡確實看得到
\texttt{before abort}——代表這個平台的 \texttt{abort} 實作會沖出已開啟串
流的緩衝內容。 但這不是所有平台都保證的行為：C 標準把它列為
implementation-defined，不能寫程式時假設緩衝內容一定會（或一定不會）留
下來；真的要保證重要資料寫進磁碟，該在 \texttt{abort} 之前自己
\texttt{fflush}，不要依賴這個沒有保證的副作用。

\keypoint{\texttt{atexit} 註冊的函式以反序呼叫。 \texttt{exit} 會沖出緩
衝區並呼叫 \texttt{atexit} 函式，\texttt{\_Exit} 兩者都跳過，
\texttt{abort} 也跳過 \texttt{atexit} 且直接送出中止訊號。 三個「結束程
式」的函式清理程度完全不同，不要當成同義詞。}

\begin{drills}[label=drills:sl-startend]
\item \texttt{atexit(cleanup1); atexit(cleanup2);} 之後程式結束，呼叫順
序是什麼？
\item \texttt{exit-vs-underscore-exit.c} 用 \texttt{exit} 跑一次、用
\texttt{\_Exit} 跑一次，都重新導向到檔案，兩個檔案內容為什麼不一樣？
\item \texttt{abort()} 之後，用 \texttt{atexit} 註冊的清理函式會不會執
行？
\item 補空：\texttt{main} 執行到最後的 \texttt{return status;}，效果等
同呼叫 \underline{\hspace{2em}}。
\end{drills}

\section{字串與數字之間：轉換函式}
\label{sec:sl-strtol}

\texttt{atoi}／\texttt{atof} 把字串轉成數字，用起來像下面這樣：

\begin{code}{c}{atoi 基本用法}
int num = atoi("12345");
double d = atof("3.14159");
\end{code}

但 \texttt{atoi} 有三個缺陷，湊在一起讓它幾乎不能用在真正需要檢查輸入的
場合：

\begin{example}{c}{(atoi-flaws.c) atoi 的三個缺陷}
#include <stdio.h>
#include <stdlib.h>

int main(void) {
    printf("%d\n", atoi("0"));       /* 合法的 0 */
    printf("%d\n", atoi("abc"));     /* 也是 0：轉換失敗還是真的是 0？ */
    printf("%d\n", atoi("42abc"));   /* 42：吃了幾個字元？不知道 */
    printf("%d\n", atoi("99999999999999999999")); /* 溢位：未定義行為 */
    return 0;
}
\end{example}

\iosample{}{0\\0\\42\\-1（此平台觀察到的值，技術上是未定義行為）}

\texttt{atoi("abc")} 和 \texttt{atoi("0")} 回傳一模一樣的 \texttt{0}，
呼叫端沒有任何辦法分辨轉換失敗與真的轉出 \texttt{0}；\texttt{"42abc"}
轉出 \texttt{42}，但呼叫端不知道字串裡還有沒被吃掉的 \texttt{"abc"}；超
出 \texttt{int} 範圍時，C 標準沒有規定會發生什麼事。

\subsection[strtol 錯誤檢查包裝函式]{用 \texttt{strtol} 寫一個真正檢查錯誤的包裝函式}

\texttt{strtol(str, \&endptr, base)} 靠 \texttt{endptr} 告訴你解析到哪
裡停下來，靠 \texttt{errno} 告訴你有沒有溢位，把這兩個資訊用起來，才能
做到 \texttt{atoi} 做不到的事：

\begin{example}{c}{(parse-int.c) 完整檢查錯誤的 strtol 包裝}
#include <stdio.h>
#include <stdlib.h>
#include <errno.h>
#include <limits.h>
#include <stdbool.h>

bool parse_int(const char *s, int *out) {
    errno = 0;
    char *endptr;
    long v = strtol(s, &endptr, 10);

    if (endptr == s) return false;       /* 完全沒吃到數字 */
    if (*endptr != '\0') return false;   /* 有尾巴垃圾 */
    if (errno == ERANGE) return false;   /* long 溢位 */
    if (v < INT_MIN || v > INT_MAX) return false; /* long 縮到 int 溢位 */

    *out = (int)v;
    return true;
}

int main(void) {
    const char *tests[] = {
        "42", "  42", "42abc", "abc", "",
        "2147483647", "2147483648", "-17"
    };
    for (int i = 0; i < 8; i++) {
        int out;
        if (parse_int(tests[i], &out))
            printf("\"%s\" -> %d\n", tests[i], out);
        else
            printf("\"%s\" -> parse failed\n", tests[i]);
    }
    return 0;
}
\end{example}

\iosample{}{"42" -> 42\\"  42" -> 42\\"42abc" -> parse failed\\"abc" -> parse failed\\"" -> parse failed\\"2147483647" -> 2147483647\\"2147483648" -> parse failed\\"-17" -> -17}

四個檢查缺一不可：\texttt{endptr == s} 抓「完全沒有數字」（例如空字串
或純字母）；\texttt{*endptr != '\textbackslash 0'} 抓「有數字，但後面還
有殘留的垃圾字元」；\texttt{errno == ERANGE} 抓「連 \texttt{long} 都裝不
下」；範圍檢查抓「\texttt{long} 裝得下、但縮到 \texttt{int} 會溢位」。
另外，\texttt{"  42"} 前面帶空白也能成功解析——\texttt{strtol} 會先跳過
開頭的空白字元，這是規格保證的行為，不是巧合。

\subsection{\texttt{base} 參數}

第三個參數決定用哪個進位制解析，\texttt{2} 到 \texttt{36} 都合法；傳
\texttt{0} 讓 \texttt{strtol} 自動偵測——\texttt{0x}／\texttt{0X} 開頭當
十六進位，\texttt{0} 開頭當八進位，其餘當十進位，跟 C 原始碼裡數字常數
的規則一致。 浮點數的等價工具是 \texttt{strtod}，用法跟 \texttt{strtol}
同一套（\texttt{endptr}、\texttt{errno}），只是沒有 \texttt{base} 參
數。

\subsection{跟 \texttt{sscanf} 比較}

只要解析一個獨立的數字，\texttt{strtol}／\texttt{strtod} 檢查得比
\texttt{sscanf} 更細（分得出「沒吃到」跟「有殘留」跟「溢位」）；一行裡
要同時解析好幾個不同型別的欄位（像 CSV 的一行），\texttt{sscanf} 或
\texttt{fgets}+\texttt{sscanf} 的組合更省事，細節見第
\ref{sec:fio-fmt}~節。 兩者不是互斥，各自適合的場合不一樣。

\keypoint{\texttt{atoi} 分不出「失敗」與「轉出 0」，也不知道吃了幾個字
元，溢位是未定義行為。 \texttt{strtol} 用 \texttt{endptr} 定位解析終
點、用 \texttt{errno} 回報溢位，四個檢查（沒吃到、有殘留、long 溢位、縮
小溢位）湊齊才算真正安全。}

\begin{drills}[label=drills:sl-strtol]
\item \texttt{atoi("abc")} 跟 \texttt{atoi("0")} 都回傳 \texttt{0}，呼叫
端要怎麼分辨轉換有沒有成功？（提示：用 \texttt{atoi} 分辨不出來）
\item \texttt{parse\_int("42abc", \&out)} 為什麼回傳 \texttt{false}？是
哪一個檢查抓到的？
\item \texttt{parse\_int("2147483648", \&out)}（\texttt{INT\_MAX + 1}）
為什麼失敗？是 \texttt{strtol} 本身失敗，還是縮小成 \texttt{int} 才失
敗？
\item 補空：\texttt{base} 傳 \texttt{0} 讓 \texttt{strtol} 自動偵測進位
制，\texttt{"0x1A"} 會被當成 \underline{\hspace{2em}} 進位。
\end{drills}

\section{函式指標與 \texttt{qsort}}
\label{sec:sl-qsort}

C 沒有泛型（generic），要寫一個能排序\textit{任何}型別陣列的函式，只剩
一條路：把「怎麼比較兩個元素」這件事，從排序演算法本身抽出來，變成一個
外部傳進來的函式。 這正是 \texttt{qsort} 的原型長這樣的原因：

\begin{code}{c}{qsort 的原型}
void qsort(void *base, size_t n, size_t size,
           int (*compar)(const void *, const void *));
\end{code}

\texttt{void *base} 是陣列開頭的位址——\texttt{void *} 能指向任何型別，
但沒有型別資訊，\texttt{qsort} 不知道每個元素多大；所以要多給
\texttt{size\_t size} 告訴它每個元素占多少位元組，\texttt{qsort} 內部就
能靠「位址 + 位移量」自己搬動任意型別的元素，完全不需要知道那個型別叫
什麼。 最後一個參數 \texttt{compar} 是函式指標：呼叫端提供一個「怎麼比
較兩個元素」的函式，\texttt{qsort} 排序過程中需要比較時就呼叫它。 位址
＋元素大小＋比較函式，三件事合起來取代了「泛型」。

比較函式的契約：接兩個 \texttt{const void *}（各指向陣列裡的一個元
素），回傳負值表示第一個排在前面、正值表示第二個排在前面、\texttt{0}
表示兩者相等。 這個規則要對陣列裡\textit{所有}可能的配對都成立，形成一
個一致的全序，否則排序結果不可靠。

\subsection{整數：小心減法溢位}

\begin{code}{c}{危險與安全的整數比較函式}
/* 危險：極端值相減會溢位 */
int compare_int_bad(const void *a, const void *b) {
    return *(const int *)a - *(const int *)b;
}

/* 安全：不做減法 */
int compare_int(const void *a, const void *b) {
    int x = *(const int *)a, y = *(const int *)b;
    return (x > y) - (x < y);
}
\end{code}

\begin{example}{c}{(qsort-int.c) 減法版在極端值上真的排錯}
#include <stdio.h>
#include <stdlib.h>

int compare_int_bad(const void *a, const void *b) {
    return *(const int *)a - *(const int *)b;
}
int compare_int(const void *a, const void *b) {
    int x = *(const int *)a, y = *(const int *)b;
    return (x > y) - (x < y);
}

int main(void) {
    int extreme[] = {2147483647, -2147483648, 0};
    qsort(extreme, 3, sizeof(int), compare_int_bad);
    for (int i = 0; i < 3; i++) printf("%d ", extreme[i]);
    printf("\n");

    int extreme2[] = {2147483647, -2147483648, 0};
    qsort(extreme2, 3, sizeof(int), compare_int);
    for (int i = 0; i < 3; i++) printf("%d ", extreme2[i]);
    printf("\n");
    return 0;
}
\end{example}

\iosample{}{2147483647 -2147483648 0\\-2147483648 0 2147483647}

第一行完全沒排好：\texttt{INT\_MAX - INT\_MIN} 本身就溢位（有號整數溢位
是未定義行為），比較函式回傳的正負號變得不可靠，\texttt{qsort} 靠這個
不可靠的訊號做出的排序自然是錯的。 第二行用
\texttt{(x > y) - (x < y)} 完全不做減法，兩個布林值相減永遠落在
\texttt{-1}／\texttt{0}／\texttt{1}，不會溢位，排序正確。 這個寫法值得
記下來，之後看到任何「比較函式」都優先考慮它，不要下意識寫
\texttt{a - b}。

\subsection{字串：兩種佈局，兩種比較函式}

字串陣列常見的兩種記憶體佈局，比較函式\textbf{不一樣}，這是
\texttt{qsort} 最高頻出錯的地方：

\begin{example}{c}{(qsort-string.c) char *arr[] 與 char arr[][N] 的比較函式}
#include <stdio.h>
#include <stdlib.h>
#include <string.h>

/* 佈局一：char *arr[]，每個元素是「指標」 */
int compare_string_ptr(const void *a, const void *b) {
    return strcmp(*(const char **)a, *(const char **)b);
}

/* 佈局二：char arr[][N]，每個元素本身就是字元陣列 */
int compare_string_arr(const void *a, const void *b) {
    return strcmp((const char *)a, (const char *)b);
}

int main(void) {
    char *words[] = {"zebra", "apple", "mango"};
    qsort(words, 3, sizeof(char *), compare_string_ptr);
    for (int i = 0; i < 3; i++) printf("%s ", words[i]);
    printf("\n");

    char fixed[3][10] = {"zebra", "apple", "mango"};
    qsort(fixed, 3, sizeof(fixed[0]), compare_string_arr);
    for (int i = 0; i < 3; i++) printf("%s ", fixed[i]);
    printf("\n");
    return 0;
}
\end{example}

\iosample{}{apple mango zebra\\apple mango zebra}

差別在 \texttt{qsort} 傳進比較函式的 \texttt{a}、\texttt{b} 永遠是「指
向陣列元素的位址」。 佈局一的元素本身就是 \texttt{char *}，所以
\texttt{a} 是「指向 \texttt{char *} 的位址」，也就是 \texttt{char **}，
要先解一次參考拿到真正的字串指標，才能傳給 \texttt{strcmp}——這就是
\texttt{*(const char **)a} 的由來。 佈局二的元素本身就是字元陣列（也就
是字串本身），\texttt{a} 已經是「指向這個字串第一個字元」的位址，直接
當 \texttt{const char *} 用即可，不用多解一次參考。 兩種比較函式對調著
用，程式通常不會馬上當掉，但比較結果毫無意義（把一段記憶體位址的位元
組硬套進 \texttt{strcmp}），是未定義行為，結果無法預期。

\subsection{結構：多鍵比較}

\begin{example}{c}{(qsort-date.c) 日期的多鍵比較}
#include <stdio.h>
#include <stdlib.h>

typedef struct { int year, month, day; } Date;

int compare_date(const void *a, const void *b) {
    const Date *d1 = (const Date *)a;
    const Date *d2 = (const Date *)b;
    if (d1->year != d2->year) return d1->year - d2->year;
    if (d1->month != d2->month) return d1->month - d2->month;
    return d1->day - d2->day;
}

int main(void) {
    Date dates[] = {{2021,7,15}, {2020,3,8}, {2020,1,5}, {2020,3,1}};
    qsort(dates, 4, sizeof(Date), compare_date);
    for (int i = 0; i < 4; i++)
        printf("%04d-%02d-%02d\n", dates[i].year, dates[i].month, dates[i].day);
    return 0;
}
\end{example}

\iosample{}{2020-01-05\\2020-03-01\\2020-03-08\\2021-07-15}

先比年，年不同就直接用年的差分出勝負，不再看月和日；年相同才比月，月
也相同才比日——這是\textbf{短路}的寫法：只要某一鍵已經分出勝負，後面
的鍵完全不用看。 這裡的年月日範圍都小，直接減法不會溢位，可以簡化寫
成 \texttt{a - b}；如果鍵的範圍可能碰到整數邊界，還是要換成
\texttt{(x > y) - (x < y)}。 順便一提：如果把日期編碼成單一整數
\texttt{YYYYMMDD}（例如 \texttt{20200308}），比較函式會退化成上一小節
的整數版本——這說明多鍵比較本質上就是「把好幾個鍵，依優先度合併成一個
排序訊號」，編碼成單一整數只是這件事的極端特例，代價是遇到日期範圍超
出整數容量、或月/日需要補零比較之類的情況會比較不直觀。

\subsection{一個位元的鍵：穩定性}

\texttt{qsort} 的規格\textbf{不保證穩定}——鍵相等的兩個元素，排序後的
相對順序可能被打散。 鍵的可能值只有兩種（例如某個位元是 \texttt{0} 還
是 \texttt{1}）時，這個問題特別容易被看到：

\begin{example}{c}{(qsort-stability.c) 位元鍵的不穩定與修復}
#include <stdio.h>
#include <stdlib.h>

typedef struct { int bit, orig_index; } Item;

int compare_bit_unstable(const void *a, const void *b) {
    return ((const Item *)a)->bit - ((const Item *)b)->bit;
}

int compare_bit_stable(const void *a, const void *b) {
    const Item *x = (const Item *)a, *y = (const Item *)b;
    if (x->bit != y->bit) return x->bit - y->bit;
    return x->orig_index - y->orig_index; /* tie-break */
}

int main(void) {
    enum { N = 40 };
    Item items[N];
    for (int i = 0; i < N; i++) {
        items[i].bit = (i < N / 2) ? 1 : 0;
        items[i].orig_index = i;
    }
    for (int i = N - 1; i > 0; i--) { /* 洗牌 */
        int j = (i * 37 + 11) % (i + 1);
        Item t = items[i]; items[i] = items[j]; items[j] = t;
    }

    Item copy[N];
    for (int i = 0; i < N; i++) copy[i] = items[i];
    qsort(copy, N, sizeof(Item), compare_bit_unstable);
    printf("unstable, bit=0 group order: ");
    for (int i = 0; i < N; i++) if (copy[i].bit == 0) printf("%d ", copy[i].orig_index);
    printf("\n");

    for (int i = 0; i < N; i++) copy[i] = items[i];
    qsort(copy, N, sizeof(Item), compare_bit_stable);
    printf("stable, bit=0 group order:   ");
    for (int i = 0; i < N; i++) if (copy[i].bit == 0) printf("%d ", copy[i].orig_index);
    printf("\n");
    return 0;
}
\end{example}

\iosample{}{unstable, bit=0 group order: 28 34 20 22 24 38 21 26 23 32 36 25 27 29 31 33 35 37 39 30\\stable, bit=0 group order:   20 21 22 23 24 25 26 27 28 29 30 31 32 33 34 35 36 37 38 39}

不加 tie-break 的版本，\texttt{bit == 0} 那一群元素排完序之後，
\texttt{orig\_index} 的順序被打亂了（不是原來的 \texttt{20, 21, 22,
...}），這就是不穩定的真面目；加了 tie-break（鍵相等時再比
\texttt{orig\_index}）的版本，\texttt{bit} 相同的元素之間永遠按照原始索
引排列，行為等同一個穩定排序。 只要在結構裡多塞一個「原始位置」欄位、
在比較函式裡當最後一個 tie-break 鍵，就能讓不保證穩定的 \texttt{qsort}
表現得像穩定排序——這個技巧跟位元操作那一章的取位元運算銜接得上：位元
鍵這種只有兩三種可能值、大量並列（tie）的情境，正是最容易踩到不穩定問
題的場合。

上面「不穩定」那組輸出是這台機器、這個 \texttt{qsort} 實作跑出來的結
果，\textbf{不是每個平台都保證會亂}：有些 libc 的 \texttt{qsort} 在陣列
夠小時會退化成插入排序，剛好又是穩定的，讓你在自己機器上怎麼測都「看
起來穩定」。 這正是不能依賴 \texttt{qsort} 穩定性的意思——沒有任何一段
標準文字保證它穩定，你的機器沒亂不代表換一個函式庫版本、換一個陣列大
小不會亂。

\keypoint{\texttt{void *} + \texttt{size\_t} 讓 \texttt{qsort} 不需要知
道元素型別；比較函式的契約是回傳負／零／正、且對所有配對一致。 整數比
較用 \texttt{(x>y)-(x<y)} 避免溢位；\texttt{char *arr[]} 要解兩次參考、
\texttt{char arr[][N]} 只需解一次；\texttt{qsort} 不保證穩定，鍵可能相
等時要自己加 tie-break。}

\begin{drills}[label=drills:sl-qsort]
\item \texttt{compare\_int\_bad} 在哪兩個值上會出錯？換成
\texttt{(x>y)-(x<y)} 為什麼能避開這個問題？
\item \texttt{char *words[]} 跟 \texttt{char fixed[3][10]} 的比較函式為
什麼不能互換？互換之後編譯器會不會報錯？
\item \texttt{compare\_date} 為什麼可以只在「年不同」時才回傳，年相同
時完全不看月和日的差？
\item 補空：\texttt{qsort} 不保證 \underline{\hspace{2em}}，鍵相等的元
素要自己在結構裡加一個欄位當 tie-break。
\end{drills}

\section{\texttt{bsearch}}
\label{sec:sl-bsearch}

\texttt{bsearch} 在已排序的陣列裡做二分搜尋，時間複雜度 $O(\log n)$，但
它有三條前提，缺一不可：

\begin{enumerate}
\item 陣列必須已經排序。
\item 搜尋時用的比較函式，必須跟排序時用的是同一套語意（同一個鍵、同
一個順序）。
\item 陣列裡不能有兩個「比較函式判定為相等」但其實不是同一筆資料、卻
期待搜到特定一筆的假設——\texttt{bsearch} 找到任何一個相等的元素就會停
下來，不保證是哪一個。
\end{enumerate}

\begin{code}{c}{bsearch 的原型}
void *bsearch(const void *key, const void *base, size_t n, size_t size,
              int (*compar)(const void *, const void *));
\end{code}

\subsection{key 與元素型別不對稱}

\texttt{bsearch} 跟 \texttt{qsort} 共用同一種比較函式格式，但呼叫時第
一個參數 \texttt{key} 不是陣列裡的元素，是呼叫端自己準備的「要找的
值」。 比較函式的兩個參數，第一個永遠是 \texttt{key} 的位址，第二個才
是陣列裡某個元素的位址——這個不對稱在同型別搜尋時看不出差別，但概念上
要分清楚。

\begin{example}{c}{(bsearch-demo.c) 重用 qsort 的三個比較函式}
#include <stdio.h>
#include <stdlib.h>

int compare_int(const void *a, const void *b) {
    int x = *(const int *)a, y = *(const int *)b;
    return (x > y) - (x < y);
}

typedef struct { int year, month, day; } Date;
int compare_date(const void *a, const void *b) {
    const Date *d1 = (const Date *)a, *d2 = (const Date *)b;
    if (d1->year != d2->year) return d1->year - d2->year;
    if (d1->month != d2->month) return d1->month - d2->month;
    return d1->day - d2->day;
}

int main(void) {
    int arr[] = {1, 2, 5, 8, 9, 10};
    int target = 5;
    int *r1 = (int *)bsearch(&target, arr, 6, sizeof(int), compare_int);
    printf("find 5: %s\n", r1 ? "found" : "not found");

    int missing = 6;
    int *r2 = (int *)bsearch(&missing, arr, 6, sizeof(int), compare_int);
    printf("find 6: %s\n", r2 ? "found" : "not found");

    Date dates[] = {{2020,1,5}, {2020,3,8}, {2021,7,15}};
    Date key = {2020, 3, 8};
    Date *r3 = (Date *)bsearch(&key, dates, 3, sizeof(Date), compare_date);
    if (r3) printf("find date: %04d-%02d-%02d\n", r3->year, r3->month, r3->day);
    return 0;
}
\end{example}

\iosample{}{find 5: found\\find 6: not found\\find date: 2020-03-08}

三個呼叫直接重用第~\ref{sec:sl-qsort}~節寫過的 \texttt{compare\_int}、
\texttt{compare\_date}——同一個比較函式，先給 \texttt{qsort} 決定排序順
序，再給 \texttt{bsearch} 決定搜尋順序，兩邊的「順序」必須是同一個定
義，這正是「泛型」最有說服力的展示：演算法（排序、搜尋）跟資料的比較
邏輯完全分離，寫一次比較函式，兩種演算法都能用。 找不到時回傳
\texttt{NULL}，跟其他「回傳指標或 NULL」家族的函式一致（第
\ref{sec:sl-spec}~節）。

\subsection{\texttt{bsearch} 做不到的事}

\texttt{bsearch} 只回答「陣列裡有沒有一個等於 \texttt{key} 的元素」，
回傳的是「某一個」相等的元素，不是「第一個」；它也做不到常見的
lower\_bound／upper\_bound 語意（找第一個不小於 \texttt{key} 的位置），
更沒辦法在「沒找到」的時候告訴你應該插入哪個位置。 這些需求只能自己手
寫二分搜尋，用半開區間 $[lo, hi)$ 的不變量寫法最不容易寫錯：

\begin{example}{c}{(lower-bound.c) 手寫 lower\_bound}
#include <stdio.h>

/* 半開區間 [lo, hi)：回傳第一個 >= target 的位置；
   找不到時回傳 n（代表該插入的位置） */
int lower_bound(const int *arr, int n, int target) {
    int lo = 0, hi = n;
    while (lo < hi) {
        int mid = lo + (hi - lo) / 2;
        if (arr[mid] < target) lo = mid + 1;
        else hi = mid;
    }
    return lo;
}

int main(void) {
    int arr[] = {1, 2, 5, 5, 5, 8, 9, 10};
    printf("lower_bound(5) = %d\n", lower_bound(arr, 8, 5));
    printf("lower_bound(6) = %d\n", lower_bound(arr, 8, 6));
    printf("lower_bound(0) = %d\n", lower_bound(arr, 8, 0));
    printf("lower_bound(11) = %d\n", lower_bound(arr, 8, 11));
    return 0;
}
\end{example}

\iosample{}{lower\_bound(5) = 2\\lower\_bound(6) = 5\\lower\_bound(0) = 0\\lower\_bound(11) = 8}

\texttt{lower\_bound(5)} 回傳 \texttt{2}：陣列裡三個 \texttt{5} 當中第
一個的位置；\texttt{lower\_bound(6)} 回傳 \texttt{5}：陣列裡沒有
\texttt{6}，但如果要插入一個 \texttt{6}、還想保持排序，該插進索引
\texttt{5}（也就是原本 \texttt{5} 那一串之後、\texttt{8} 之前）。 這正
是遞迴那一章講過的區間表示法——用半開區間表示「還沒排除的搜尋範圍」，
迴圈不變量永遠是「答案在 \texttt{[lo, hi)} 裡面」，寫法上跟遞迴版的二分
搜尋是同一個思路，只是這裡用迴圈實作。

\keypoint{\texttt{bsearch} 的前提是陣列已排序、比較函式跟排序時一致；
它回傳「某一個」相等元素、找不到時回傳 \texttt{NULL}，做不到
lower\_bound／upper\_bound。 需要插入位置或第一個/最後一個相等元素，要
自己手寫半開區間二分搜尋。}

\begin{drills}[label=drills:sl-bsearch]
\item \texttt{bsearch} 的三條前提各自缺一個會發生什麼事？各舉一個後果。
\item 比較函式的第一個參數是 \texttt{key} 的位址還是陣列元素的位址？跟
第二個參數比起來有什麼不同？
\item 為什麼 \texttt{bsearch-demo.c} 可以直接重用
\texttt{qsort-date.c} 裡的 \texttt{compare\_date}？這件事說明了什麼？
\item 補空：\texttt{bsearch} 做不到「找不到時回傳插入位置」，這種需求
要改用 \underline{\hspace{2em}}。
\end{drills}

\section{隨機數}
\label{sec:sl-rand}

\texttt{rand()} 回傳一個 \texttt{0} 到 \texttt{RAND\_MAX} 之間的偽隨機
整數，\texttt{srand(seed)} 設定產生序列的種子——同一個種子永遠產生同一
串序列（方便寫可重現的測試），\texttt{srand(time(NULL))} 用目前時間當
種子，讓每次執行的序列不同。

\subsection{模擬骰子}

\begin{example}{c}{(dice.c) 模擬 60000 次骰子，看均勻性}
#include <stdio.h>
#include <stdlib.h>
#include <time.h>

int main(void) {
    srand((unsigned)time(NULL));
    int counts[7] = {0};
    int trials = 60000;
    for (int i = 0; i < trials; i++) {
        int die = rand() % 6 + 1;
        counts[die]++;
    }
    for (int face = 1; face <= 6; face++) {
        printf("%d: %5d (%.2f%%)\n", face, counts[face], 100.0 * counts[face] / trials);
    }
    return 0;
}
\end{example}

\iosample{}{1: 10092 (16.82\%)\\2: 10153 (16.92\%)\\3:  9904 (16.51\%)\\4:  9872 (16.45\%)\\5:  9964 (16.61\%)\\6: 10015 (16.69\%)}

六個面的比例都貼近理論值 $1/6 \approx 16.67\%$，眼睛看均勻性比背公式直
覺。

\subsection{模偏差與拒絕取樣}

\texttt{rand() \% n} 是最常見的寫法，但有兩個問題：低位元的隨機品質通
常比高位元差（某些演算法的實作，低位元的週期性比高位元明顯）；如果
\texttt{RAND\_MAX + 1} 不是 \texttt{n} 的整數倍，\texttt{\% n} 的結果不
會完全均勻，某些餘數會比另一些稍微更常出現，這叫\textbf{模偏差}。
\texttt{n} 很小、\texttt{RAND\_MAX} 很大時偏差通常小到量不出來，但公平
性要求高的場合（例如洗牌、抽獎），正確做法是\textbf{拒絕取樣}：

\begin{example}{c}{(rand-reject.c) 拒絕取樣避免模偏差}
#include <stdio.h>
#include <stdlib.h>

int fair_rand(int n) {
    int limit = (RAND_MAX / n) * n; /* 只保留能整除 n 的範圍 */
    int r;
    do {
        r = rand();
    } while (r >= limit); /* 超出範圍的值直接丟掉，重新抽 */
    return r % n;
}

int main(void) {
    srand(1);
    int counts[3] = {0};
    for (int i = 0; i < 30000; i++) counts[fair_rand(3)]++;
    for (int i = 0; i < 3; i++) printf("%d: %d\n", i, counts[i]);
    return 0;
}
\end{example}

\iosample{}{0: 9980\\1: 10079\\2: 9941}

\texttt{limit} 是小於或等於 \texttt{RAND\_MAX} 的最大的 \texttt{n} 的倍
數（\texttt{(RAND\_MAX / n) * n} 用整數除法自動算出這個值）；抽到的值大
於或等於 \texttt{limit} 就直接丟掉重抽，只在保留下來的均勻範圍裡取餘數，
這樣每個餘數出現的機率才是真正相等。 有個邊界情況值得一提：如果
\texttt{RAND\_MAX + 1} 剛好是 \texttt{n} 的倍數，這個算法會比理論上必要
的範圍再少留一個 \texttt{n}，多丟掉一整段其實還算均勻的值——不影響正確
性，只是稍微浪費了一點點抽樣。 最後要誠實說明一件
事：\texttt{rand} 的演算法是公開的，序列可預測，\textbf{不能用在任何安
全相關的場合}（密碼、token、亂數金鑰），那些場合要用作業系統提供的密
碼學安全亂數來源，不在本書範圍內。

\subsection{蒙地卡羅法估算 $\pi$}

\begin{example}{c}{(monte-carlo.c) 隨機取樣估算圓周率}
#include <stdio.h>
#include <stdlib.h>
#include <time.h>
#include <math.h>

double estimate_pi(long samples) {
    long inside = 0;
    for (long i = 0; i < samples; i++) {
        double x = (double)rand() / RAND_MAX;
        double y = (double)rand() / RAND_MAX;
        if (x * x + y * y <= 1.0) inside++;
    }
    return 4.0 * inside / samples;
}

int main(void) {
    srand((unsigned)time(NULL));
    const double pi_exact = acos(-1.0); /* M_PI 不是標準 C，acos(-1.0) 才可攜 */
    for (long n = 1000; n <= 10000000; n *= 100) {
        double pi_est = estimate_pi(n);
        printf("N=%8ld  pi~=%.6f  error=%.6f\n", n, pi_est, fabs(pi_est - pi_exact));
    }
    return 0;
}
\end{example}

\texttt{<math.h>} 定義的 \texttt{M\_PI} 是 POSIX／編譯器擴充，不是 ISO C
標準的一部分，用 \texttt{-std=c11} 這種嚴格模式編譯會找不到這個名字。
\texttt{acos(-1.0)} 完全靠標準函式算出同一個 $\pi$，不依賴任何非標準巨
集——這是本章少數幾個「非標準擴充」的例子，看到 \texttt{M\_} 開頭的巨集
（\texttt{M\_E}、\texttt{M\_SQRT2}⋯）都要留意它們不保證每個編譯器都有。

\iosample{}{N=    1000  pi~=3.136000  error=0.005593\\N=  100000  pi~=3.143880  error=0.002287\\N=10000000  pi~=3.142468  error=0.000876}

在單位正方形裡隨機取點，落在四分之一圓內的比例乘上 \texttt{4} 就是
$\pi$ 的估計值——圓面積跟正方形面積的比例是 $\pi/4$。 這種蒙地卡羅法
的誤差量級是 $O(1/\sqrt{N})$：樣本數要乘 \texttt{100}，誤差才平均縮小
到十分之一（多一位有效數字），不是乘 \texttt{10} 就有效果。 上面這組
實測數字（因為是隨機的，每次執行都會有一點出入）大致印證了這個趨勢：
$N$ 每乘 \texttt{100}，誤差量級確實往下掉了一截，但比起「精確」的數學
方法收斂得慢得多——這正是隨機演算法的現實：簡單好寫，但要更精確就得付
出跟樣本數平方成正比的代價。

\keypoint{\texttt{srand(seed)} 決定 \texttt{rand()} 的整條序列，同 seed
可重現，測試很需要這個特性。 \texttt{rand() \% n} 有模偏差，公平性要求
高時用拒絕取樣；\texttt{rand} 絕對不能用在安全相關的場合。 蒙地卡羅法
的誤差是 $O(1/\sqrt{N})$，樣本數要乘 100 才多一位精度。}

\begin{drills}[label=drills:sl-rand]
\item 同一個 \texttt{seed} 呼叫 \texttt{srand} 兩次、各接著呼叫十次
\texttt{rand()}，兩串結果會一樣還是不一樣？這個特性對寫測試有什麼用？
\item \texttt{rand() \% 3} 在 \texttt{RAND\_MAX} 不是 \texttt{3} 的倍數
時，為什麼不是完全均勻？
\item 拒絕取樣裡，抽到的值大於或等於 \texttt{limit} 時為什麼要丟掉重
抽，而不是直接對 \texttt{limit} 取餘數？
\item 補空：蒙地卡羅法的誤差量級是 \underline{\hspace{2em}}，樣本數要
乘 \underline{\hspace{2em}} 倍，誤差才平均縮小到十分之一。
\end{drills}

\section{字元與字串}
\label{sec:sl-charstring}

\texttt{<ctype.h>} 的分類函式（\texttt{isalpha}、\texttt{isdigit}、
\texttt{isspace}⋯）有一條幾乎沒有中文教材強調過的規則：\textbf{參數必
須先轉成 \texttt{unsigned char}}（或是 \texttt{EOF}），否則是未定義行
為。

\begin{example}{c}{(ctype-ub.c) 有號 char 傳進 ctype 函式的陷阱}
#include <stdio.h>
#include <ctype.h>

int main(void) {
    char c = -1; /* 在有號 char 的平台上，這是負值 */
    printf("isalpha(c) [dangerous] = %d\n", isalpha(c));
    printf("isalpha((unsigned char)c) [safe] = %d\n", isalpha((unsigned char)c));
    return 0;
}
\end{example}

\iosample{}{isalpha(c) [dangerous] = 0\\isalpha((unsigned char)c) [safe] = 0}

兩行印出來的結果一樣，這正是這個陷阱最陰險的地方：這台機器上兩種寫法看
起來「都對」，程式也沒有當掉，很容易讓人誤以為 \texttt{isalpha(c)} 這樣
寫沒問題。 但這只是這個平台的查表剛好夠寬容，不是保證——原因是這些函式的
典型實作是一張查表：用參數的數值當索引，去一張長度
\texttt{256}（或更長）的表裡查分類結果。 \texttt{char} 在很多平台是有
號的，值 \texttt{-1} 當索引會變成負數，往表格陣列的邊界之前存取，是不
是造成問題完全看實作——這台機器的表格剛好多留了緩衝、沒有壞掉，但 C
標準沒有做這個保證，換一個平台可能直接讀到不相關的記憶體甚至當掉。
規則很簡單：呼叫 \texttt{<ctype.h>} 的任何函式之前，先把參數轉成
\texttt{unsigned char}（如果來源可能是 \texttt{EOF}，直接傳
\texttt{int} 型別的 \texttt{EOF} 值即可，不用轉型，因為 \texttt{EOF}
本身就在允許的輸入範圍裡）。

\subsection{\texttt{<string.h>} 的搜尋函式群}

\begin{example}{c}{(string-search.c) strchr／strstr／strspn／strcspn}
#include <stdio.h>
#include <string.h>

int main(void) {
    const char *s = "hello, world";
    char *p = strchr(s, 'w');
    printf("strchr: %s\n", p ? p : "(null)");

    char *q = strstr(s, "wor");
    printf("strstr: %s\n", q ? q : "(null)");

    const char *digits = "12345abc";
    size_t n = strspn(digits, "0123456789");
    printf("strspn: %zu digits at the front\n", n);

    char line[] = "hello\n";
    line[strcspn(line, "\n")] = '\0';
    printf("strcspn: \"%s\" (length %zu)\n", line, strlen(line));
    return 0;
}
\end{example}

\iosample{}{strchr: world\\strstr: world\\strspn: 5 digits at the front\\strcspn: "hello" (length 5)}

\texttt{strchr} 找一個字元，\texttt{strstr} 找一個子字串，兩者都回傳
\texttt{NULL} 表示沒找到（回傳值家族：指標或 NULL）；\texttt{strspn}
回傳「從開頭算起，連續屬於某個字元集合的長度」，常用來檢查一段前綴是
不是全部由某些字元組成。 \texttt{strcspn} 剛好相反：回傳「從開頭算起，連
續\textbf{不屬於}某個字元集合的長度」，也就是第一個「有出現在集合裡」的
字元位置。

\keypoint{\texttt{line[strcspn(line, "\textbackslash n")] = '\textbackslash 0';}
是去掉 \texttt{fgets} 讀進來那個換行字元的標準慣用法：\texttt{strcspn}
找到第一個 \texttt{\textbackslash n} 的位置，直接把那個位置寫成字串結尾，
比自己找換行、比對長度都省事，找不到換行時也剛好等於字串長度，不會出
錯。}

\subsection{\texttt{strtok} 的靜態狀態}

\texttt{strtok(str, delim)} 第一次呼叫傳真正的字串，之後每次呼叫傳
\texttt{NULL} 就會接著上次停下的地方繼續切：

\begin{example}{c}{(strtok-reentrant.c) strtok 的靜態狀態不可重入}
#include <stdio.h>
#include <string.h>

void process_line(char *line) {
    char *tok = strtok(line, " ");
    while (tok != NULL) {
        printf("  word: %s\n", tok);
        tok = strtok(NULL, " ");
    }
}

int main(void) {
    char lines[] = "aa bb cc\nxx yy zz";
    char *line = strtok(lines, "\n");
    while (line != NULL) {
        printf("line: %s\n", line);
        char copy[32];
        snprintf(copy, sizeof copy, "%s", line); /* 有邊界檢查，不像 strcpy 會溢位 */
        process_line(copy); /* 內部又呼叫了一次 strtok！ */
        line = strtok(NULL, "\n");
    }
    return 0;
}
\end{example}

\iosample{}{line: aa bb cc\\\ \ word: aa\\\ \ word: bb\\\ \ word: cc}

第二行 \texttt{xx yy zz} 完全沒有出現。 \texttt{strtok} 靠一個函式內部
的\textbf{靜態變數}記住「上次切到哪裡」，這個狀態是全域共用的，不分是
哪一段字串在切。 外層 \texttt{strtok(lines, "\textbackslash n")} 切完
第一行後，靜態狀態記著「還沒切到 \texttt{xx yy zz} 那一段」；但
\texttt{process\_line} 內部又呼叫了 \texttt{strtok(line, " ")}，把同一
份靜態狀態改成「這一份 \texttt{copy} 已經切完」——外層迴圈接著呼叫
\texttt{strtok(NULL, "\textbackslash n")} 時，讀到的是被內層弄亂的狀
態，於是直接回傳 \texttt{NULL}，第二行就這樣不見了。 這正是「不可重
入」的意思：一個函式如果同一時間（或巢狀呼叫）被用兩次，會互相干擾。
需要巢狀切字串，要用 POSIX 的 \texttt{strtok\_r} 或 C11 的
\texttt{strtok\_s}，把狀態變成呼叫端自己拿著的變數，不共用。

\keypoint{\texttt{<ctype.h>} 系列函式的參數必須先轉成
\texttt{unsigned char}，否則有號 \texttt{char} 的負值是未定義行為。
\texttt{strtok} 用共用的靜態狀態記位置，巢狀呼叫（或多執行緒）會互相
干擾，這種情況要換 \texttt{strtok\_r}。}

\begin{drills}[label=drills:sl-charstring]
\item 為什麼 \texttt{isalpha} 這類函式的參數要先轉成
\texttt{unsigned char}？這台機器沒出事，代表這樣寫是安全的嗎？
\item \texttt{strchr} 與 \texttt{strstr} 找不到時各回傳什麼？
\item \texttt{strtok-reentrant.c} 為什麼第二行 \texttt{xx yy zz} 沒有被
印出來？
\item 補空：需要在巢狀呼叫裡各自維護獨立的切字串狀態，要換用
\underline{\hspace{2em}}。
\end{drills}

\section{數學}
\label{sec:sl-math}

\texttt{<math.h>} 的函式碰到定義域外的輸入（例如 \texttt{sqrt(-1)}）或
結果超出可表示範圍（例如 \texttt{log(0)}）時，標準允許兩種回報機制：
設定 \texttt{errno}（\texttt{EDOM} 表示定義域錯誤、\texttt{ERANGE} 表示
值域錯誤），或是設定浮點例外旗標，兩種機制哪個生效由
\texttt{math\_errhandling} 這個巨集決定，\textbf{不是所有平台都用
errno}：

\begin{example}{c}{(math-errors.c) 定義域與值域錯誤}
#include <stdio.h>
#include <math.h>
#include <errno.h>

int main(void) {
    errno = 0;
    double r1 = sqrt(-1.0);
    printf("sqrt(-1) = %f, isnan = %d, errno==EDOM? %d\n",
           r1, isnan(r1), errno == EDOM);

    errno = 0;
    double r2 = log(0.0);
    printf("log(0) = %f, isinf = %d, errno==ERANGE? %d\n",
           r2, isinf(r2), errno == ERANGE);
    return 0;
}
\end{example}

\iosample{}{sqrt(-1) = nan, isnan = 1, errno==EDOM? 0\\log(0) = -inf, isinf = 1, errno==ERANGE? 0}

在這台機器上兩個 \texttt{errno} 檢查都是假的——不是程式寫錯，是這個平
臺的 \texttt{math\_errhandling} 只用浮點例外（\texttt{MATH\_ERREXCEPT}），
沒有用 \texttt{errno}（\texttt{MATH\_ERRNO}）。 這也是為什麼
\texttt{isnan}／\texttt{isinf} 比 \texttt{errno} 更值得信賴：不管平台選
哪種回報機制，算出來的結果只要是「非數」或「無限大」，這兩個分類巨集
一定能檢查出來；\texttt{errno} 檢查則要先確認自己的平台真的支援
\texttt{MATH\_ERRNO}。 想知道自己平台用哪一種，直接印
\texttt{math\_errhandling} 跟標頭裡定義的 \texttt{MATH\_ERRNO}／
\texttt{MATH\_ERREXCEPT} 比較就知道。

\texttt{fabs} 跟 \texttt{abs} 不能混用：\texttt{abs} 是
\texttt{<stdlib.h>} 裡處理 \texttt{int} 的版本，\texttt{fabs} 是
\texttt{<math.h>} 裡處理 \texttt{double} 的版本，把浮點數傳進
\texttt{abs} 會先被截斷成整數，得到完全不是本意的結果。

\texttt{<float.h>} 跟 \texttt{<limits.h>} 不是函式，是一批巨集常數，用
來查詢目前平台各種型別的極限值，例如 \texttt{FLT\_EPSILON}（\texttt{float}
能分辨的最小差距）、\texttt{DBL\_MAX}、\texttt{INT\_MAX}——想寫出不假設
特定平台的可攜程式碼，該查這兩個標頭，而不是自己猜一個數字寫死。

\keypoint{\texttt{math.h} 的錯誤可能走 \texttt{errno}、也可能只走浮點
例外，由 \texttt{math\_errhandling} 決定，別預設 \texttt{errno} 一定有
效；\texttt{isnan}／\texttt{isinf} 更可靠。 \texttt{fabs} 處理
\texttt{double}，\texttt{abs} 處理 \texttt{int}，不能混用。
\texttt{<float.h>}／\texttt{<limits.h>} 查平台極限，不要自己猜數字。}

\begin{drills}[label=drills:sl-math]
\item 為什麼這台機器上檢查 \texttt{errno == EDOM} 抓不到 \texttt{sqrt(-1)}
的錯誤？該檢查哪個函式代替？
\item 把一個 \texttt{double} 傳進 \texttt{abs} 會發生什麼事？
\item \texttt{FLT\_EPSILON} 跟 \texttt{DBL\_MAX} 分別在說什麼？
\item 補空：判斷 \texttt{math.h} 函式的計算結果是不是「非數」，應該用
\underline{\hspace{2em}}，而不是只靠 \texttt{errno}。
\end{drills}

\section{時間}
\label{sec:sl-time}

\texttt{time\_t} 是一個代表「某個時間點」的型別（實務上通常是從某個固
定時刻起算的秒數，但標準沒有規定具體編碼）；\texttt{struct tm} 把時間
拆成年、月、日、時、分、秒等欄位，方便讀取或組裝。 兩個型別之間要靠
\texttt{mktime} 轉換：

\begin{example}{c}{(time-demo.c) struct tm、mktime、difftime、strftime}
#include <stdio.h>
#include <time.h>

int main(void) {
    struct tm t1 = {0};
    t1.tm_year = 2020 - 1900; /* tm_year 是「距 1900 年」 */
    t1.tm_mon = 3 - 1;        /* tm_mon 是 0-11 */
    t1.tm_mday = 8;
    t1.tm_isdst = -1;         /* 不確定有沒有夏令時間，交給 mktime 判斷 */

    struct tm t2 = {0};
    t2.tm_year = 2021 - 1900;
    t2.tm_mon = 7 - 1;
    t2.tm_mday = 15;
    t2.tm_isdst = -1;

    time_t time1 = mktime(&t1);
    time_t time2 = mktime(&t2);

    double diff_days = difftime(time2, time1) / (60 * 60 * 24);
    printf("days between: %.0f\n", diff_days);

    char buf[64];
    strftime(buf, sizeof buf, "%Y-%m-%d (%A)", &t1);
    printf("t1 formatted: %s\n", buf);
    return 0;
}
\end{example}

\iosample{}{days between: 494\\t1 formatted: 2020-03-08 (Sunday)}

兩個欄位最容易搞錯：\texttt{tm\_year} 存的是\textit{距離 1900 年}的差
值，不是西元年份本身，\texttt{2020} 年要存 \texttt{120}；
\texttt{tm\_mon} 是 \texttt{0} 到 \texttt{11}，一月是 \texttt{0}，不是
\texttt{1}。 還有一個容易漏掉的欄位：\texttt{tm\_isdst}。 用
\texttt{\{0\}} 初始化會讓它是 \texttt{0}，等於明確告訴 \texttt{mktime}
「這個時間點不是夏令時間」；如果實際所在時區當時真的在實施夏令時間，
\texttt{mktime} 算出來的結果會差一小時。 傳 \texttt{-1} 才是正確的預設
值——請 \texttt{mktime} 自己依據 \texttt{tm\_year}／\texttt{tm\_mon}／
\texttt{tm\_mday} 判斷該不該套用夏令時間，不用呼叫端自己猜。
\texttt{mktime} 除了轉換型別，還會把 \texttt{struct tm}
「正規化」（例如 \texttt{tm\_mday} 給 \texttt{32} 時自動推算成下個月的
第幾天，順便算出正確的星期），這也是上面 \texttt{t1 formatted} 那一行
輸出裡星期幾（\texttt{Sunday}）是對的的原因——我們一開始沒有填
\texttt{tm\_wday}，是 \texttt{mktime} 幫忙補上的。 \texttt{difftime}
回傳兩個 \texttt{time\_t} 之間的秒數差（\texttt{double}），除以
\texttt{86400} 就是天數差。 \texttt{strftime} 依格式字串把
\texttt{struct tm} 印成字串，比自己拼字串安全，也會自動套用當地語言的
月份／星期名稱（如果設定了對應的 locale，見第
\ref{sec:sl-i18n}~節）。

回頭看第~\ref{sec:sl-qsort}~節的日期比較函式：那裡用的
\texttt{Date} 結構是簡化版，只有年月日三個整數；真正處理日期時，換成
\texttt{struct tm}、用 \texttt{mktime} 轉成 \texttt{time\_t} 再直接用
數值比較（或用 \texttt{difftime} 的正負號），可以少考慮「日期進位」
（例如二月有幾天、閏年）的邊界問題，因為 \texttt{mktime} 已經幫忙正規
化過。

\keypoint{\texttt{tm\_year} 是距 1900 年的差值，\texttt{tm\_mon} 是
\texttt{0}-\texttt{11}，兩個都不是字面上的年份／月份數字；
\texttt{tm\_isdst} 該設成 \texttt{-1} 讓 \texttt{mktime} 自己判斷夏令時
間，不要留著 \texttt{\{0\}} 初始化出來的 \texttt{0}。 \texttt{mktime}
會正規化欄位並補上星期；\texttt{difftime} 給秒數差；\texttt{strftime}
是印時間字串的安全做法。}

\begin{drills}[label=drills:sl-time]
\item \texttt{struct tm} 存 \texttt{2024} 年 \texttt{tm\_year} 該填多
少？存三月 \texttt{tm\_mon} 該填多少？
\item \texttt{time-demo.c} 沒有填 \texttt{t1.tm\_wday}，為什麼
\texttt{strftime} 印出的星期還是對的？
\item \texttt{difftime} 回傳的是什麼單位？換算成天要除以多少？
\item 補空：把 \texttt{Date} 結構換成 \texttt{struct tm} 再用
\texttt{mktime} 轉換的好處是，不用自己處理 \underline{\hspace{2em}} 之
類的日期進位問題。
\end{drills}

\section{可變參數}
\label{sec:sl-variadic}

\texttt{printf} 能接受不固定數量的參數，靠的是 \texttt{<stdarg.h>} 提供
的一組工具。 自己寫一個「像 \texttt{printf} 一樣能吃任意參數」的函式，
需要 \texttt{va\_list}（存取參數列的狀態）、\texttt{va\_start}（開始存
取）、\texttt{va\_end}（結束存取），加上 \texttt{<stdio.h>} 裡專門吃
\texttt{va\_list} 的 \texttt{vfprintf} 家族：

\begin{example}{c}{(log-printf.c) 寫一個吃任意參數的 log 函式}
#include <stdio.h>
#include <stdarg.h>

void log_printf(const char *fmt, ...) {
    va_list args;
    va_start(args, fmt);
    fprintf(stderr, "[LOG] ");
    vfprintf(stderr, fmt, args);
    fprintf(stderr, "\n");
    va_end(args);
}

int main(void) {
    int x = 42;
    log_printf("x = %d", x);
    log_printf("y = %d, name = %s", 7, "test");
    return 0;
}
\end{example}

\iosample{}{[LOG] x = 42\\{[LOG] y = 7, name = test}}

\texttt{va\_start(args, fmt)} 的第二個參數必須是「最後一個有名字的參
數」（這裡是 \texttt{fmt}），這樣 \texttt{va\_list} 才知道從哪裡開始接
著找可變的那一批參數。 \texttt{vfprintf} 跟 \texttt{fprintf} 做的事一
樣，差別只在參數列已經被包成一個 \texttt{va\_list}，不是攤開的
\texttt{...}——標準函式庫裡幾乎每一個 \texttt{printf} 家族的函式（
\texttt{printf}、\texttt{sprintf}、\texttt{fprintf}⋯）都有一個對應的
\texttt{v} 開頭版本，就是為了讓像 \texttt{log\_printf} 這樣「包一層」的
函式能把收到的可變參數直接轉手傳下去，不需要自己一個一個攤開解析格式
字串。

\subsection{巨集層級 vs 函式層級的可變參數}

第~\ref{sec:pp-variadic}~節已經示範過另一種可變參數：巨集層級的
\texttt{\_\_VA\_ARGS\_\_}，那裡的 \texttt{LOG(fmt, ...)} 在\textit{前置
處理期}就直接展開成一整行 \texttt{fprintf} 呼叫——展開完之後跟手寫的
\texttt{fprintf} 沒有任何差別，前置處理器不管參數個數，只是照樣把
\texttt{\_\_VA\_ARGS\_\_} 換成呼叫端傳的那一串東西。 這一節的
\texttt{log\_printf} 是另一回事：它是一個真正在\textit{執行期}用
\texttt{va\_list} 走訪任意個數參數的\textbf{函式}，不是文字替換出來
的。

兩種機制可以合作，各司其職：巨集負責取得呼叫端當下的
\texttt{\_\_FILE\_\_}／\texttt{\_\_LINE\_\_}（這件事只有巨集辦得到——
函式呼叫發生在巨集展開之後，函式內部已經看不到呼叫端的原始碼位置），
格式化與輸出的部分轉手交給函式：

\begin{code}{c}{巨集轉呼叫函式，各司其職}
#define LOG(fmt, ...) \
    (fprintf(stderr, "[%s:%d] ", __FILE__, __LINE__), \
     log_printf(fmt "\n", __VA_ARGS__))
\end{code}

跟第~\ref{sec:pp-assert}~節直接在巨集裡塞好一整行 \texttt{fprintf} 比起
來，這裡把「要不要多加時間戳記、要不要同時寫進檔案」這些以後可能變複
雜的邏輯，都收進 \texttt{log\_printf} 這一個函式——以後要改 log 的行
為，只要改這個函式本體；巨集永遠只做它唯一辦得到的事：取得呼叫端的原
始碼位置，其餘全部轉手交給函式。

\keypoint{\texttt{va\_start}／\texttt{va\_end} 搭配 \texttt{va\_list} 讓
函式在\textbf{執行期}真正走訪任意數量的參數；\texttt{vfprintf} 家族專門
吃已經包好的 \texttt{va\_list}。 巨集的 \texttt{\_\_VA\_ARGS\_\_}（第
\ref{sec:pp-variadic}~節）在\textbf{前置處理期}就展開成固定的呼叫。 兩
者可以合作：巨集取得呼叫端的原始碼位置，函式負責真正的格式化邏輯。}

\begin{drills}[label=drills:sl-variadic]
\item \texttt{va\_start(args, fmt)} 的第二個參數為什麼一定要是最後一個
有名字的參數？
\item \texttt{vfprintf} 跟 \texttt{fprintf} 的差別是什麼？為什麼
\texttt{log\_printf} 要用 \texttt{vfprintf} 而不是直接呼叫
\texttt{fprintf(stderr, fmt, args)}？（提示：\texttt{fprintf} 的參數列
長什麼樣）
\item 為什麼 \texttt{\_\_FILE\_\_}／\texttt{\_\_LINE\_\_} 一定要留在
\texttt{LOG} 巨集裡取得，不能搬進 \texttt{log\_printf} 函式內部用同樣的
方式取得？
\item 補空：巨集層級的可變參數用 \underline{\hspace{2em}} 代表整包參
數，在\underline{\hspace{2em}}期展開；函式層級的可變參數要靠
\underline{\hspace{2em}} 走訪，在\underline{\hspace{2em}}期真正運作。
\end{drills}

\section{國際化（只講一件事）}
\label{sec:sl-i18n}

國際化的完整規則對初學者太重，這裡只回答一個問題：為什麼
\texttt{strlen("中文")} 不是 \texttt{2}？

\begin{example}{c}{(utf8-len.c) 位元組數不等於字元數}
#include <stdio.h>
#include <string.h>
#include <stdlib.h>
#include <locale.h>
#include <wchar.h>

int main(void) {
    const char *s = "中文";
    printf("strlen(\"%s\") = %zu\n", s, strlen(s));

    setlocale(LC_ALL, "");
    wchar_t wbuf[16];
    size_t wn = mbstowcs(wbuf, s, 16);
    if (wn == (size_t)-1) {
        /* 遇到不合法的多位元組序列（跟目前 locale 對不上）就會這樣回報 */
        fprintf(stderr, "decode failed: invalid multibyte sequence\n");
        return 1;
    }
    printf("characters decoded: %zu\n", wn);
    return 0;
}
\end{example}

\iosample{}{strlen("中文") = 6\\characters decoded: 2}

原始碼裡的字串常量在編譯器眼中就是一串 \texttt{char}（位元組），
\texttt{strlen} 數的是位元組數，跟人類眼中的「字元數」完全是兩件事。
這份原始碼用 UTF-8 編碼儲存，「中」跟「文」在 UTF-8 裡各占 \texttt{3}
個位元組，兩個字合計 \texttt{6} 個位元組——這就是 \texttt{strlen} 印出
\texttt{6} 的原因。 要拿到人類直覺上的「字元數」，得先知道字串的編碼規
則，把位元組序列\textit{解碼}成一串「寬字元」（\texttt{wchar\_t}），
\texttt{mbstowcs}（multi-byte string to wide-character string）就是做
這件事的函式，靠目前的 \texttt{locale} 設定知道要用哪種編碼規則解碼；
\texttt{setlocale(LC\_ALL, "")} 讓程式改用執行環境（例如殼層的
\texttt{LANG} 環境變數）指定的 locale，而不是 C 語言預設、只認得
ASCII 的 \texttt{"C"} locale。

一句話總結三個層次：\textbf{位元組}是儲存與傳輸的單位（\texttt{char}）；
\textbf{編碼字元}是人類認知的一個符號，可能占一到多個位元組（UTF-8 下
的中文字通常是三個）；\textbf{寬字元}是程式內部用來統一處理各種語言字
元的型別（\texttt{wchar\_t}），跟位元組數脫鉤。 更完整的
locale／編碼規則，以及 \texttt{wchar\_t} 系列函式的細節，超出本書範
圍，需要時查手冊或 cppreference。

\keypoint{\texttt{strlen} 數的是位元組數，不是人類眼中的字元數；UTF-8
編碼下常見的中文字元占三個位元組。 想要正確的「字元數」，要先靠
\texttt{setlocale} 加 \texttt{mbstowcs} 之類的函式解碼成寬字元。}

\begin{drills}[label=drills:sl-i18n]
\item \texttt{strlen("中文")} 印出 \texttt{6}，這個 \texttt{6} 數的是什
麼單位？
\item 如果沒有先呼叫 \texttt{setlocale(LC\_ALL, "")}，程式預設用的是哪
個 locale？這個 locale 認得中文字元嗎？
\item \texttt{mbstowcs} 這個函式名字裡的 \texttt{mbs} 跟 \texttt{wcs}
分別代表什麼？
\item 補空：想知道一個字串「人類眼中有幾個字」，不能只看
\texttt{strlen}，要先用 \underline{\hspace{2em}} 把位元組序列解碼成寬
字元。
\end{drills}

\phantomsection
\addcontentsline{toc}{section}{整合專案：CSV 資料表 CLI}
\section*{整合專案：CSV 資料表 CLI}
\label{sec:sl-capstone}

單一個專案，比十三節各自的小範例更能看出這些工具怎麼配合。 專案是一個
處理 CSV 資料表的命令列工具，每個子命令對應前面某幾節：

\begin{itemize}
\item 讀檔用第~\ref{sec:fio-line}~節的 \texttt{fgets}，每個欄位用
\texttt{strtol}／\texttt{strtod} 解析（第~\ref{sec:sl-strtol}~節），解
析失敗印出行號、用 \texttt{EXIT\_FAILURE} 結束（第
\ref{sec:sl-startend}~節）。
\item \texttt{sort} 子命令用多鍵比較函式（日期主鍵、姓名次鍵）交給
\texttt{qsort}（第~\ref{sec:sl-qsort}~節）。
\item \texttt{--gen N} 子命令用 \texttt{rand} 產生 \texttt{N} 筆隨機測
資（第~\ref{sec:sl-rand}~節）。
\item \texttt{--stats} 子命令排序後用 \texttt{strftime} 印出資料的日期
範圍（第~\ref{sec:sl-time}~節）。
\item 用 \texttt{atexit} 註冊一個清理函式，程式結束時報告處理了幾筆資
料（第~\ref{sec:sl-startend}~節）。
\end{itemize}

\begin{example}{c}{(csvtool.c) CSV 資料表 CLI（節錄核心邏輯）}
typedef struct {
    int year, month, day;
    char name[32];
    double score;
} Row;

static int row_count = 0;

static void report_cleanup(void) {
    fprintf(stderr, "[cleanup] processed %d rows\n", row_count);
}

/* 解析 "YYYY-MM-DD" 格式的日期，失敗回傳 0 */
static int parse_date(const char *s, int *y, int *m, int *d) {
    char *endptr;
    errno = 0;
    long yy = strtol(s, &endptr, 10);
    if (errno != 0 || *endptr != '-') return 0;
    long mm = strtol(endptr + 1, &endptr, 10);
    if (errno != 0 || *endptr != '-') return 0;
    long dd = strtol(endptr + 1, &endptr, 10);
    if (errno != 0 || *endptr != '\0') return 0;
    *y = (int)yy; *m = (int)mm; *d = (int)dd;
    return 1;
}

/* 逐行讀取，任何一欄解析失敗就印出行號並 exit(EXIT_FAILURE) */
static Row *load_csv(const char *path, int *out_n) {
    FILE *fp = fopen(path, "r");
    if (!fp) { perror(path); exit(EXIT_FAILURE); }

    int capacity = 16, n = 0, lineno = 0;
    Row *rows = malloc((size_t)capacity * sizeof(Row));
    if (!rows) { fprintf(stderr, "out of memory\n"); exit(EXIT_FAILURE); }
    char line[256];

    while (fgets(line, sizeof line, fp)) {
        lineno++;
        line[strcspn(line, "\n")] = '\0';
        if (line[0] == '\0') continue;

        char date_str[32], name_str[32], score_str[32];
        if (sscanf(line, "%31[^,],%31[^,],%31s",
                   date_str, name_str, score_str) != 3) {
            fprintf(stderr, "line %d: bad format: %s\n", lineno, line);
            exit(EXIT_FAILURE);
        }

        Row r;
        if (!parse_date(date_str, &r.year, &r.month, &r.day)) {
            fprintf(stderr, "line %d: bad date: %s\n", lineno, date_str);
            exit(EXIT_FAILURE);
        }
        strncpy(r.name, name_str, sizeof r.name - 1);
        r.name[sizeof r.name - 1] = '\0';

        errno = 0;
        char *endptr;
        r.score = strtod(score_str, &endptr);
        if (errno != 0 || *endptr != '\0') {
            fprintf(stderr, "line %d: bad score: %s\n", lineno, score_str);
            exit(EXIT_FAILURE);
        }

        if (n == capacity) {
            capacity *= 2;
            Row *tmp = realloc(rows, (size_t)capacity * sizeof(Row));
            if (!tmp) { fprintf(stderr, "out of memory\n"); exit(EXIT_FAILURE); }
            rows = tmp;
        }
        rows[n++] = r;
    }
    fclose(fp);
    *out_n = n;
    return rows;
}

/* 多鍵比較：日期主鍵，姓名次鍵 */
static int compare_row(const void *a, const void *b) {
    const Row *r1 = (const Row *)a, *r2 = (const Row *)b;
    if (r1->year != r2->year) return r1->year - r2->year;
    if (r1->month != r2->month) return r1->month - r2->month;
    if (r1->day != r2->day) return r1->day - r2->day;
    return strcmp(r1->name, r2->name);
}

int main(int argc, char **argv) {
    if (argc < 3) { fprintf(stderr, "usage: %s sort <file>\n", argv[0]); return 1; }
    atexit(report_cleanup);

    int n;
    Row *rows = load_csv(argv[2], &n);
    row_count = n; /* atexit 的 report_cleanup 靠這個變數報告筆數 */

    qsort(rows, (size_t)n, sizeof(Row), compare_row);
    for (int i = 0; i < n; i++) {
        printf("%04d-%02d-%02d,%s,%.1f\n",
               rows[i].year, rows[i].month, rows[i].day, rows[i].name, rows[i].score);
    }
    free(rows);
    return 0;
}
\end{example}

\begin{code}{bash}{執行}
gcc csvtool.c -o csvtool
./csvtool sort data.csv
./csvtool --stats data.csv
./csvtool --gen 5
\end{code}

\begin{code}{text}{data.csv}
2020-03-08,bob,88.5
2020-01-05,alice,92.0
2021-07-15,carol,75.0
2020-03-08,alice,81.0
\end{code}

\iosample{}{2020-01-05,alice,92.0\\2020-03-08,alice,81.0\\2020-03-08,bob,88.5\\2021-07-15,carol,75.0\\{[cleanup] processed 4 rows}}

日期相同的兩筆（\texttt{2020-03-08} 的 \texttt{bob} 跟 \texttt{alice}）
靠次鍵姓名分出順序——這是第~\ref{sec:sl-qsort}~節多鍵比較的直接應用。
值得留意一個容易忽略的執行細節：\texttt{atexit} 註冊的
\texttt{report\_cleanup} 用 \texttt{fprintf(stderr, ...)}（不緩衝，馬上
輸出），而 \texttt{sort} 印出的資料是透過 \texttt{printf} 寫進
\texttt{stdout}。 在互動式終端機上，\texttt{stdout} 是\textbf{行緩衝}，
每一行排序結果印完（帶著換行）就立刻沖出去，所以上面
\texttt{\textbackslash iosample} 呈現的順序（排序結果在前、
\texttt{cleanup} 訊息最後）就是終端機上會看到的正常順序。 但如果把 \texttt{stdout} 重新導向到檔案（或者像
\texttt{./csvtool sort data.csv > out.txt 2>\&1} 這樣連 \texttt{stderr}
也一起導向、合併觀察），緩衝模式會切成全緩衝——這正是第
\ref{sec:sl-startend}~節 \texttt{exit-vs-underscore-exit.c} 那個例子的設
定。 \texttt{sort} 印出的內容會先留在 \texttt{stdout} 的緩衝區裡，直到
\texttt{exit} 收尾時才真正沖出；\texttt{cleanup} 走的 \texttt{stderr}
不緩衝，當場就寫出去。 於是在合併觀察的情況下，反而會先看到
\texttt{cleanup} 那一行、排序結果最後才出現——順序完全由緩衝模式決定，不
是「終端機上就會這樣」。

\section{常見錯誤清單}
\label{sec:sl-checklist}

下面每一條都對應前面某一節。

\begin{itemize}
\item 底線開頭、\texttt{str+小寫} 是保留名稱；巨集可能遮蔽同名函式，用
\texttt{(name)()} 或 \texttt{\#undef} 繞開。
（第~\ref{sec:sl-what}~節）
\item 回傳值協定分四個家族：哨兵值、指標或 NULL、計數、errno。 看新函
式先看回傳型別落在哪個家族。
（第~\ref{sec:sl-spec}~節）
\item 檢查 \texttt{errno} 前先清零，只在規格明確提到的情況下檢查；
\texttt{assert} 抓程式錯誤、可被 \texttt{NDEBUG} 關掉，使用者輸入錯誤
不能靠 \texttt{assert}。
（第~\ref{sec:sl-errors}~節）
\item \texttt{exit}／\texttt{\_Exit}／\texttt{abort} 清理程度不同：
\texttt{exit} 沖緩衝區＋呼叫 \texttt{atexit}，\texttt{\_Exit} 都跳過，
\texttt{abort} 也跳過 \texttt{atexit} 且直接送出中止訊號。
（第~\ref{sec:sl-startend}~節）
\item \texttt{atoi} 分不出失敗與轉出 0；\texttt{strtol} 要四個檢查（沒
吃到、有殘留、long 溢位、縮小溢位）才算安全。
（第~\ref{sec:sl-strtol}~節）
\item 整數比較函式用 \texttt{(x>y)-(x<y)} 避免溢位；\texttt{char *arr[]}
與 \texttt{char arr[][N]} 的比較函式不能互換；\texttt{qsort} 不保證穩
定。
（第~\ref{sec:sl-qsort}~節）
\item \texttt{bsearch} 前提是已排序、比較函式一致；做不到插入位置，要
自己手寫二分搜尋。
（第~\ref{sec:sl-bsearch}~節）
\item \texttt{rand() \% n} 有模偏差；\texttt{rand} 不能用在安全相關場
合。
（第~\ref{sec:sl-rand}~節）
\item \texttt{ctype.h} 參數要先轉 \texttt{unsigned char}；\\
\texttt{strtok} 的靜態狀態巢狀呼叫會互相干擾。
（第~\ref{sec:sl-charstring}~節）
\item \texttt{math.h} 錯誤可能只走浮點例外，不一定設 \texttt{errno}；
\texttt{fabs} 處理 \texttt{double}，\texttt{abs} 處理 \texttt{int}。
（第~\ref{sec:sl-math}~節）
\item \texttt{tm\_year} 是距 1900 年差值，\texttt{tm\_mon} 是
\texttt{0}-\texttt{11}。
（第~\ref{sec:sl-time}~節）
\item \texttt{strlen} 數位元組不是字元數，UTF-8 中文字通常占 3 個位元組。\\
（第~\ref{sec:sl-i18n}~節）
\end{itemize}

\begin{drills}[label=drills:sl-checklist]
\item 下面這段有三個錯，請全部標出。
\begin{verbatim}
int compare_int_bad(const void *a, const void *b) {
    return *(const int *)a - *(const int *)b;
}
char *words[] = {"b", "a"};
qsort(words, 2, sizeof(char *), compare_string_arr); /* 用了 char arr[][N] 版 */
char c = -1;
isalpha(c);
\end{verbatim}
\item 指出錯誤：檢查 \texttt{errno} 之前沒有先清零。
\item 指出錯誤：對使用者輸入的檔名用 \texttt{assert(file\_exists(name))}。
\end{drills}

\phantomsection
\addcontentsline{toc}{section}{練習解答}
\section*{練習解答}

各節練習的答案收在這裡。 完整程式見正文片語／範例。

\paragraph{第~\ref{sec:sl-what}~節}
\begin{enumerate}
\item 標頭檔只有宣告，前置處理與編譯階段只需要知道函式的樣子就能通
過；\texttt{sqrt} 的本體在 \texttt{libm} 裡，要到連結階段把目的檔跟函
式庫接起來時，缺少 \texttt{-lm} 才會找不到本體。
\item \texttt{\_} 開頭的識別字整批保留給編譯器與函式庫實作使用，跟未
來標準或編譯器內部機制撞名的風險較高。
\item 巨集展開要求名字後面直接跟著左括號，\texttt{(getchar)} 中間被括
號隔開，不符合展開樣式，所以一定被當成普通函式呼叫。
\item \texttt{\#undef}。
\end{enumerate}

\paragraph{第~\ref{sec:sl-spec}~節}
\begin{enumerate}
\item 因為函式要透過這個參數「回傳」另一個資訊——解析停在哪個字元，
呼叫端傳一個 \texttt{char *} 變數的位址進去，函式才能修改那個變數本
身，讓它指向正確的位置。
\item 因為轉換出來的數字可能超過 \texttt{int} 能裝的範圍，用更大的型
別接住，讓呼叫端自己決定要不要檢查縮小到 \texttt{int} 會不會溢位。
\item \texttt{fopen} 屬於「指標或 NULL」家族；\texttt{fread} 屬於「計
數」家族。
\item 轉換失敗（完全沒吃到合法的數字）；\texttt{endptr} 有沒有前進（是
不是還等於原始的字串起點）。
\end{enumerate}

\paragraph{第~\ref{sec:sl-errors}~節}
\begin{enumerate}
\item 因為 \texttt{errno} 不會在函式成功時自動清成 \texttt{0}，上一次
呼叫留下的非零舊值可能還在，沒有先清零就檢查會誤判。
\item 兩者印出的內容是同一句話；\texttt{perror} 直接把訊息印到
\texttt{stderr}，\texttt{strerror} 只回傳字串，印在哪裡自己決定。
\item 回傳值處理。 使用者輸入錯誤在發布版也要能正確反應給使用者，不
能被 \texttt{NDEBUG} 關掉檢查。
\item 編譯；執行。
\end{enumerate}

\paragraph{第~\ref{sec:sl-startend}~節}
\begin{enumerate}
\item \texttt{cleanup2} 先呼叫，\texttt{cleanup1} 後呼叫（反序）。
\item \texttt{exit} 版本結束前會沖出緩衝區，沒換行的內容也會寫進檔
案；\texttt{\_Exit} 版本直接跳過沖出這一步，緩衝區內容被丟棄，檔案是
空的。
\item 不會。 \texttt{abort} 直接送出中止訊號，跳過 \texttt{atexit}。
\item \texttt{exit(status)}。
\end{enumerate}

\paragraph{第~\ref{sec:sl-strtol}~節}
\begin{enumerate}
\item 用 \texttt{atoi} 分辨不出來；要換成 \texttt{strtol} 搭配
\texttt{endptr} 檢查（\texttt{endptr == s} 代表完全沒轉成功）。
\item \texttt{*endptr != '\textbackslash 0'} 抓到的：\texttt{42} 後面
還留著沒被解析的 \texttt{"abc"}。
\item 是縮小成 \texttt{int} 才失敗：\texttt{strtol} 本身用
\texttt{long} 接，\texttt{2147483648} 沒有超過 \texttt{long} 的範圍
（在常見的 64 位元平台上），但超過了 \texttt{INT\_MAX}，被範圍檢查那
一行擋下來。
\item 十六（16）。
\end{enumerate}

\paragraph{第~\ref{sec:sl-qsort}~節}
\begin{enumerate}
\item 在 \texttt{INT\_MAX} 與 \texttt{INT\_MIN} 這類極端值上，兩者相減
會超出 \texttt{int} 能表示的範圍，是未定義行為；
\texttt{(x>y)-(x<y)} 完全不做減法，兩個布林值相減永遠落在
\texttt{-1}/\texttt{0}/\texttt{1}，不會溢位。
\item 因為 \texttt{qsort} 傳進來的位址型別意義不同：\texttt{char *arr[]}
每個元素是指標，比較函式收到的是「指向指標的位址」，要解兩次參考；
\texttt{char arr[][N]} 每個元素是字元陣列本身，只需解一次。 互換不會
讓編譯器報錯（兩邊參數型別都是 \texttt{const void *}），但比較結果是
未定義行為。
\item 因為年不同就已經足以分出兩個日期的先後順序，後面的月、日完全
不影響這個判斷結果，所以可以直接回傳、不用再往下比。
\item 穩定（stable）。
\end{enumerate}

\paragraph{第~\ref{sec:sl-bsearch}~節}
\begin{enumerate}
\item 陣列沒排序：二分搜尋的前提被破壞，結果不可靠、可能漏找。 比較
函式跟排序時不一致：語意錯位，一樣找不到本該找到的元素。
\item 第一個參數是 \texttt{key} 的位址；第二個才是陣列裡某個元素的位
址。 兩者角色不同：一個是外部提供的搜尋目標，一個是陣列裡實際存在的
資料。
\item 因為兩者的排序（也就是「順序」的定義）是同一套語意，寫一次比
較函式，排序跟搜尋都能用；這正是把演算法跟資料型別拆開的好處。
\item 手寫二分搜尋（例如 lower\_bound 的半開區間寫法）。
\end{enumerate}

\paragraph{第~\ref{sec:sl-rand}~節}
\begin{enumerate}
\item 兩串結果一樣；這個特性讓測試可以用固定 seed 重現同一組「隨機」
輸入，方便除錯與比對。
\item 因為 \texttt{RAND\_MAX + 1} 除以 \texttt{3} 如果除不盡，靠近
\texttt{RAND\_MAX} 的那一小段餘數對應的原始值比其他餘數少，取餘數後
這些餘數出現的機率會略低於其他餘數。
\item 因為如果直接對 \texttt{limit} 取餘數，被丟棄範圍原本對應的那些
餘數會被硬塞進去，一樣造成不均勻；丟掉重抽才能保證只在真正均勻的範
圍裡取餘數。
\item $O(1/\sqrt{N})$；\texttt{100}。
\end{enumerate}

\paragraph{第~\ref{sec:sl-charstring}~節}
\begin{enumerate}
\item 因為這些函式常見的實作是查表，用參數值當索引；有號 \texttt{char}
的負值當索引會存取到表格邊界之前，標準沒保證這樣安全。 這台機器沒出
事只是這個平台的表格剛好多留了緩衝，換一個平台可能真的出事，不能當
成「這樣寫沒問題」的證據。
\item 都回傳 \texttt{NULL}。
\item 因為外層 \texttt{strtok(lines, "\textbackslash n")} 跟內層
\texttt{process\_line} 裡的 \texttt{strtok(line, " ")} 共用同一份靜態
狀態，內層呼叫把狀態改成「已經切完」，外層接著呼叫時讀到被弄亂的狀
態，直接回傳 \texttt{NULL}。
\item \texttt{strtok\_r}（或 C11 的 \texttt{strtok\_s}）。
\end{enumerate}

\paragraph{第~\ref{sec:sl-math}~節}
\begin{enumerate}
\item 因為這個平台的 \texttt{math\_errhandling} 只用浮點例外
（\texttt{MATH\_ERREXCEPT}），沒有用 \texttt{errno}；該檢查
\texttt{isnan}。
\item 會先被截斷成整數，得到完全不是本意的結果（例如小數部分整個消
失）。
\item \texttt{FLT\_EPSILON} 是 \texttt{float} 能分辨的最小差距；
\texttt{DBL\_MAX} 是 \texttt{double} 能表示的最大值。
\item \texttt{isnan}（配合 \texttt{isinf} 檢查無限大）。
\end{enumerate}

\paragraph{第~\ref{sec:sl-time}~節}
\begin{enumerate}
\item \texttt{124}（\texttt{2024 - 1900}）；三月的 \texttt{tm\_mon}
填 \texttt{2}（\texttt{0} 起算，一月是 \texttt{0}）。
\item 因為呼叫 \texttt{mktime} 時，它會依其他欄位重新計算並正規化整
個 \texttt{struct tm}，順便補上正確的 \texttt{tm\_wday}。
\item 秒；除以 \texttt{86400}（\texttt{60*60*24}）。
\item 閏年、每個月天數不同。
\end{enumerate}

\paragraph{第~\ref{sec:sl-variadic}~節}
\begin{enumerate}
\item 因為 \texttt{va\_start} 要知道從哪個位置開始接著找可變參數，那
個位置就是最後一個有名字的參數之後。
\item \texttt{fprintf} 直接接可變參數列（\texttt{...}）；
\texttt{vfprintf} 接的是已經包好的 \texttt{va\_list}。
\texttt{log\_printf} 收到的 \texttt{args} 已經是 \texttt{va\_list}，沒
辦法直接攤開成 \texttt{...} 再交給 \texttt{fprintf}，只能用
\texttt{vfprintf} 轉手傳遞。
\item 因為 \texttt{\_\_FILE\_\_}／\texttt{\_\_LINE\_\_} 是預定義巨集，
永遠展開成「目前這一行程式碼」的檔名與行號；如果搬進
\texttt{log\_printf} 函式內部使用，展開出來的會是
\texttt{log\_printf} 定義所在那一行的檔名行號，不是呼叫端真正呼叫
\texttt{LOG} 的那一行。
\item \texttt{\_\_VA\_ARGS\_\_}；前置處理；\texttt{va\_list}；執行。
\end{enumerate}

\paragraph{第~\ref{sec:sl-i18n}~節}
\begin{enumerate}
\item 位元組數。
\item \texttt{"C"} locale；不認得，只認得 ASCII 範圍的字元。
\item \texttt{mbs} 是 multi-byte string（多位元組字串，也就是一般的
\texttt{char} 字串）；\texttt{wcs} 是 wide-character string（寬字元字
串）。
\item \texttt{setlocale} 加 \texttt{mbstowcs}。
\end{enumerate}

\paragraph{第~\ref{sec:sl-checklist}~節}
\begin{enumerate}
\item 三個錯：\texttt{compare\_int\_bad} 用減法比較整數，極端值會溢
位；\texttt{words} 是 \texttt{char *arr[]} 佈局卻用了
\texttt{char arr[][N]} 版的比較函式；\texttt{isalpha(c)} 傳進去的是有
號 \texttt{char}，沒有先轉成 \texttt{unsigned char}。
\item 上一次呼叫留下的非零 \texttt{errno} 可能還在，沒清零就檢查會誤
判成「這次呼叫失敗」。
\item 錯：使用者輸入的檔名不存在，是執行期會發生、需要處理的正常情
況，不是程式設計者的邏輯錯誤；用 \texttt{assert} 會導致這個檢查在
\texttt{-DNDEBUG} 的發布版裡直接消失，應該用回傳值／錯誤訊息處理。
\end{enumerate}

\chapter{資料結構}
\label{ch:datastructure}

這一章從頭到尾只靠三句話撐起來：

\begin{itemize}
\item \textbf{heap 上的東西不會自己消失。} 陣列與一般區域變數的生命週期
由編譯器自動管理；\texttt{malloc} 配置出來的記憶體不會，忘記處理就是
洩漏，處理錯了就是懸置指標或雙重釋放。
\item \textbf{每個資料結構都有不變式。} 「從 \texttt{head} 沿 \texttt{next}
走有限步會到 \texttt{NULL}」「左子樹都小於根、右子樹都大於根」——每個
操作的本質都是「暫時打破這句話，再想辦法恢復」。 bug 幾乎都是「打破
了，卻忘記恢復」。
\item \textbf{遞迴結構要用遞迴思考。} 鏈結串列與二元樹的定義本身就是遞
迴的（「一個節點，加上一個更小的同類結構」），走訪、搜尋、重建的程式
碼形狀，會跟著資料的形狀走。
\end{itemize}

章節安排跟著這三句話走：先弄懂 heap 記憶體怎麼一回事、三種會壞掉的方
式、怎麼用工具親眼看見壞掉；再進鏈結串列的節點與不變式、插入為什麼要
按順序、為什麼要用「指標的指標」修改呼叫端的指標；最後用同一套不變式
語言講二元搜尋樹的插入、走訪、重建與集合運算。

\section{為什麼陣列不夠}
\label{sec:ds-why-not-array}

想像要讀入一次考試的成績，題目沒告訴你會有幾筆——可能 5 筆，可能
5000 筆。 陣列宣告時就要決定大小：宣告太小，資料一多就裝不下；宣告
「夠大」（例如 \texttt{int scores[100000];}），大多數時候都在浪費記
憶體。 陣列還有第二個限制：要在\textbf{中間}插入或刪除一筆，得把後面
整段元素搬家，成本跟「後面還剩幾格」成正比。

\begin{center}
\begin{tikzpicture}
    \node[cell] (a0) at (0,0) {1};
    \node[cell] (a1) at (0.9,0) {4};
    \node[cell] (a2) at (1.8,0) {8};
    \node[idxlbl, anchor=east] at (-0.55,0) {陣列};

    \node[cell] (n0) at (4.2,0) {1};
    \node[cell] (n1) at (6.4,0) {4};
    \node[cell] (n2) at (8.6,0) {8};
    \node[idxlbl, anchor=east] at (3.65,0) {鏈結};

    \draw[diagarrow] (n0.east) -- (n1.west);
    \draw[diagarrow] (n1.east) -- (n2.west);
    \node[idxlbl, anchor=west] at (9.35,0) {NULL};
\end{tikzpicture}
\diagcaption{左：陣列元素必須連續存放，大小固定；右：鏈結串列的節點可
以散落在記憶體各處，靠指標維持順序，可以一個一個動態長大。}
\label{fig:ds-array-vs-linked}
\end{center}

鏈結串列（linked list）換一種思路：每筆資料各自包成一個\textbf{節
點}，節點裡除了資料，還有一個指標指出「下一筆是誰」。 要多少節點就配
置多少，不用事先猜上限；插入或刪除一筆，只要改幾個指標，不用搬動其他
元素。 代價是：要透過索引直接跳到第 $i$ 筆，只能從頭一路數過去，做不
到陣列的 $O(1)$ 隨機存取。

\keypoint{陣列擅長「用索引直接定位」，大小固定、隨處可用；鏈結串列擅
長「大小隨輸入動態增長」與「已知節點旁插入刪除」，代價是失去隨機存
取。 選哪一種，看題目最頻繁的操作是什麼。}

\section{heap 與 stack：變數活多久}
\label{sec:ds-heap-stack}

一支執行中的程式，記憶體大致分成幾個區域：程式碼本身、全域變數、
\textbf{stack}（存區域變數與函式呼叫的紀錄）、\textbf{heap}（存動態配
置出來的資料）。 這一節先不出現任何資料結構，只講兩件事的差別：區域
變數活多久、\texttt{malloc} 出來的東西活多久。

\begin{center}
\begin{tikzpicture}[node distance=0mm]
    \node[cell, minimum width=3.6cm, minimum height=0.8cm] (code) at (0,3) {程式碼（唯讀）};
    \node[cell, minimum width=3.6cm, minimum height=0.8cm] (glob) at (0,2.1) {全域變數};
    \node[cell, minimum width=3.6cm, minimum height=0.8cm] (heap) at (0,1.2) {heap（往上長）};
    \node[cell, minimum width=3.6cm, minimum height=0.8cm] (stack) at (0,0.3) {stack（往下長）};
    \node[idxlbl, anchor=west] at (2.0,1.2) {malloc/free};
    \node[idxlbl, anchor=west] at (2.0,0.3) {函式呼叫};
\end{tikzpicture}
\diagcaption{stack 存區域變數，函式一返回就自動收回；heap 存
\texttt{malloc} 配置的資料，要自己 \texttt{free} 才會收回，兩者互不干
涉、各自往中間長大。}
\label{fig:ds-memory-layout}
\end{center}

區域變數的生命週期跟著它所在的函式：函式一返回，區域變數占用的 stack
空間就可以被下一次函式呼叫覆寫。 回傳一個指向區域變數的位址，是危險
動作：

\begin{example}{c}{(lifetime-bad.c) 回傳區域變數的位址}
#include <stdio.h>

int *make_local(void) {
    int x = 42;
    return &x; /* 危險：x 是區域變數，函式一返回就沒了 */
}

int main(void) {
    int *p = make_local();
    printf("%d\n", *p);
    return 0;
}
\end{example}

\iosample{}{42}

編譯器已經主動警告了（\texttt{-Wreturn-stack-address}），但程式依然編
得過、跑得出 \texttt{42}——因為 \texttt{x} 那塊 stack 空間還沒被別的呼
叫覆寫，剛好看起來沒事。 這正是「未定義行為」的典型長相：不是保證壞
掉，是\textbf{不保證好}，換一次呼叫順序、換一顆編譯器，結果可能完全不
同。 \texttt{malloc} 配置的記憶體剛好相反：它活在 heap 上，不受函式返
回影響，一直存在到你自己呼叫 \texttt{free}（或程式結束）。 這正是鏈結
串列、樹這些「跨函式呼叫依然要存在」的結構，非用 \texttt{malloc} 不可
的原因。

\keypoint{區域變數活在 stack 上，函式一返回就消失；\texttt{malloc} 配
置的資料活在 heap 上，activity 不跟著任何函式的返回結束，要自己
\texttt{free}。 回傳指向區域變數的位址，是這一章第一個要牢記的地雷。}

\section{malloc 與 free}
\label{sec:ds-malloc-free}

\texttt{malloc(n)} 向系統要一塊至少 \texttt{n} 個位元組的記憶體，回傳
一個 \texttt{void *}；配置失敗時回傳 \texttt{NULL}。 完整的使用模式永
遠是三步驟：配置、使用、釋放。

\begin{code}{c}{malloc 的三步驟}
int *p = malloc(sizeof(int));  /* 配置 */
if (p == NULL) {
    /* 配置失敗，通常是系統記憶體不足 */
}
*p = 5;                        /* 使用 */
free(p);                       /* 釋放 */
p = NULL;                      /* 習慣：free 之後立刻讓指標指向 NULL */
\end{code}

\subsection{不要轉型 \texttt{malloc} 的回傳值}

很多教材會寫 \texttt{(int *)malloc(...)}，但 C 裡\textbf{不建議}轉
型：\texttt{void *} 在 C 裡會自動轉成任何指標型別，轉型只是多餘的文
字。 更重要的是：轉型從前有個真正的風險——舊式 K\&R 編譯器在沒有
\texttt{\#include <stdlib.h>} 時，會假設 \texttt{malloc} 回傳
\texttt{int}，而顯式轉型會把這個型別不合的警告直接蓋掉，在 32 位元
\texttt{int}／64 位元指標的平臺上悄悄截斷位址。 現代編譯器（本書使用
的這一顆也一樣）已經把「呼叫沒有原型的函式」列為編譯期錯誤，不管有沒
有轉型都逃不掉，這個歷史風險已經不成立；但轉型本身仍然是純粹的雜訊，
不轉型才是目前的慣例。

\subsection{\texttt{sizeof(*p)} 優於 \texttt{sizeof(型別名)}}

\begin{code}{c}{sizeof 用誰}
Node *p = malloc(sizeof(Node));   /* 能動，但型別名要跟著 p 的宣告手動同步 */
Node *p = malloc(sizeof(*p));     /* 更好：p 的型別變了，這裡自動跟著變 */
\end{code}

寫 \texttt{sizeof(*p)} 時，配置的大小永遠跟著 \texttt{p} 實際的型別
走；哪天把 \texttt{Node *p} 改成 \texttt{TreeNode *p} 卻忘記同步改
\texttt{sizeof(Node)}，會配置出\textit{大小不對}的記憶體——這種錯誤編
譯器抓不到，執行期也不一定馬上出事，是很難除的臭蟲。 \texttt{sizeof(*p)}
從結構上避開了這個問題。

\subsection{\texttt{calloc} 與 \texttt{realloc}}

\texttt{calloc(n, size)} 跟 \texttt{malloc(n * size)} 類似，但會把整塊
記憶體清成 \texttt{0}，且對 \texttt{n * size} 的乘法溢位有內建檢查；
\texttt{realloc(p, newsize)} 試著把 \texttt{p} 指向的區塊改成新大小，
可能原地擴充、也可能搬到別處（並把舊內容複製過去），回傳新位址，失敗
時回傳 \texttt{NULL}\textbf{且原本的 \texttt{p} 仍然有效}——所以絕對不
能寫 \texttt{p = realloc(p, ...)}，失敗時會把唯一的位址蓋成
\texttt{NULL}，原本那塊記憶體就找不到了，變成洩漏。 正確寫法是接一個
暫存指標：

\begin{code}{c}{realloc 要用暫存指標接}
int *tmp = realloc(p, new_n * sizeof(int));
if (tmp == NULL) {
    /* p 仍然指向原本的記憶體，可以在這裡善後 */
} else {
    p = tmp;
}
\end{code}

\keypoint{三步驟：配置、使用、釋放；\texttt{malloc} 不用轉型；用
\texttt{sizeof(*p)} 而不是 \texttt{sizeof(型別名)}；\texttt{realloc}
一定要接暫存指標，不能直接蓋掉原本的指標。}

\section[三種記憶體錯誤]{三種會壞掉的方式：洩漏、懸置指標、雙重釋放}
\label{sec:ds-memory-bugs}

「忘了處理 heap 記憶體」不是一種錯誤，是\textbf{三種不同}的錯誤，症狀
與成因都不一樣。

\subsection{記憶體洩漏（memory leak）}

配置出去的記憶體，沒有任何指標還指得到它，也沒被 \texttt{free}——它就
永遠回不去了，直到程式結束。

\begin{example}{c}{(leak-demo.c) 迴圈裡配置卻忘記釋放}
#include <stdio.h>
#include <stdlib.h>

int main(void) {
    for (int i = 0; i < 5; i++) {
        int *p = malloc(sizeof(int));
        *p = i;
        printf("%d\n", *p); /* 用起來完全正常，看不出任何異狀 */
    } /* 每次迴圈結束，p 這個區域變數消失，但它指向的記憶體沒有被 free */
    return 0;
}
\end{example}

\iosample{}{0\\1\\2\\3\\4}

這支程式輸出完全正常，短時間執行也感覺不到任何異狀——這正是洩漏最陰
險的地方：\textbf{不會馬上爆，這才可怕}。 長時間執行的伺服器或反覆呼
叫的函式，一次洩漏幾十個位元組看不出問題，跑幾天、幾百萬次呼叫後，記
憶體可能被慢慢吃光。

\subsection{懸置指標（dangling pointer）}

\texttt{free} 之後，指標本身的\textit{值}沒有改變，還是指向那塊記憶
體——但那塊記憶體已經被系統收回，可能隨時被別的配置要走、內容被覆寫。
繼續用這個指標，就是懸置指標。

\begin{example}{c}{(dangling-demo.c) free 之後繼續使用}
#include <stdio.h>
#include <stdlib.h>

int main(void) {
    int *p = malloc(sizeof(int));
    *p = 99;
    free(p);
    printf("%d\n", *p); /* 懸置指標：p 指向的記憶體已經還給系統了 */
    return 0;
}
\end{example}

\iosample{}{3}

注意輸出不是 \texttt{99}，而是別的數字（這台機器上是 \texttt{3}，換一
臺機器、換一次執行都可能不一樣）——\texttt{free} 之後那塊記憶體的內容
不再屬於你，讀到什麼是未定義行為，讀到「看起來正常的舊值」有時反而比
讀到亂碼更危險，因為程式表面上還是「正常」執行下去。

\subsection{雙重釋放（double free）}

同一個指標 \texttt{free} 兩次：

\begin{example}{c}{(doublefree-demo.c) 同一塊記憶體 free 兩次}
#include <stdlib.h>

int main(void) {
    int *p = malloc(sizeof(int));
    free(p);
    free(p); /* 危險：同一塊記憶體被釋放兩次 */
    return 0;
}
\end{example}

跟前兩種不一樣，雙重釋放在很多平臺（包括這台示範機器）會被配置器自己
的安全檢查抓到，程式直接異常中止——這是三種錯誤裡\textbf{唯一}通常會
馬上有動靜的一種，但這不是 C 標準保證的行為，只是常見的實作選擇；換一
個平臺或關掉安全檢查，雙重釋放一樣可能像懸置指標那樣，安靜地損壞資料
而不中止。

\keypoint{洩漏（配置了、沒人指得到、沒 free）、懸置指標（free 之後還
在用）、雙重釋放（同一塊 free 兩次）是三種不同的錯誤。 洩漏與懸置指標
通常\textbf{不會}馬上讓程式壞掉，這正是它們難抓的原因；free 之後立刻
把指標設成 \texttt{NULL} 能讓「不小心又用了一次」直接變成容易發現的
\texttt{NULL} 解參考，而不是安靜的懸置指標。}

\section{工具：親眼看見記憶體壞在哪}
\label{sec:ds-tools}

上一節的三種錯誤，光靠讀程式碼不一定看得出來——這正是專門的記憶體檢
查工具存在的理由。 Linux 上最常見的是 \texttt{valgrind}；本書的示範機
器是 Apple Silicon 版 macOS，\texttt{valgrind} 目前不支援這個平臺，改
用 macOS 內建（隨 Xcode 命令列工具安裝）的 \texttt{leaks} 指令示範同一
件事：

\begin{code}{bash}{用 leaks 抓 leak-demo.c 的洩漏}
gcc leak-demo.c -o leak-demo
leaks --atExit -- ./leak-demo
\end{code}

\begin{code}{text}{leaks 輸出（節錄摘要）}
Process 12345: 193 nodes malloced for 17 KB
Process 12345: 5 leaks for 80 total leaked bytes.
\end{code}

\texttt{leaks} 精準地數出「5 leaks」——正好對應迴圈跑了 5 次、每次都
配置了一個 \texttt{int} 卻沒 \texttt{free}。 讀程式碼要花心思才看得出
的洩漏，工具一眼就能量化，這正是它的教學回報率這麼高的原因：\textbf{記憶體
管理從一個抽象概念，變成一個可以親眼觀察、可以計數的現象}。

另一個常見工具是編譯器內建的 AddressSanitizer（\texttt{-fsanitize=address}），
在 Linux 與多數 macOS 環境下可以直接偵測懸置指標、雙重釋放、緩衝區溢
位，一發現就印出出錯的原始碼行號並中止程式：

\begin{code}{bash}{用 AddressSanitizer 編譯}
gcc -fsanitize=address -g dangling-demo.c -o dangling-demo
./dangling-demo
\end{code}

\texttt{-fsanitize=address} 會讓程式變慢、變大，只適合開發／測試階段
使用，不要用在正式發布的版本；\\
\texttt{-g} 讓錯誤訊息附上原始碼行號。 兩種工具（\texttt{leaks} 或
\texttt{valgrind}／\texttt{-fsanitize=address}）挑一種，在寫鏈結串列與
樹之前先養成「寫完就跑一次檢查工具」的習慣，比事後除錯有效率得多。

\keypoint{記憶體錯誤光看程式碼不一定抓得出來，要靠工具實際跑一次。\\
Linux 常用 \texttt{valgrind}，macOS 可以用內建的 \texttt{leaks}，兩個
平臺都能用編譯器內建的 \texttt{-fsanitize=address}。 開發階段就該養成
用工具檢查的習慣，不要等到題目「怪怪的」才想到記憶體可能出問題。}

\section{assert：抓的是「不該發生」的事}
\label{sec:ds-assert}

\texttt{assert(cond)} 在 \texttt{cond} 為假時印出訊息並中止程式，是開
發階段快速抓出「按照邏輯不該發生」的狀況的工具；\texttt{-DNDEBUG}
編譯後，所有 \texttt{assert} 會整批消失，變成完全不檢查。

\begin{example}{c}{(assert-malloc-bad.c) 反面教材：用 assert 檢查 malloc}
#include <assert.h>
#include <stdlib.h>

int *make_array(int n) {
    int *arr = malloc(n * sizeof(int));
    assert(arr != NULL); /* 危險示範，見下方說明 */
    return arr;
}

int main(void) {
    int *a = make_array(1000000000); /* 刻意要一大塊，示意配置可能失敗 */
    return a != NULL;
}
\end{example}

這是本章要糾正的一個常見錯誤：\textbf{不要}用 \texttt{assert} 檢查
\texttt{malloc} 的回傳值。 理由有兩個。 第一，\texttt{malloc} 失敗
（系統記憶體不足）是執行期真的可能發生的事，不是「程式設計者的邏輯錯
誤」——這正是 \texttt{assert} 該抓的對象的定義相反。 第二，一旦發布版
加上 \texttt{-DNDEBUG}，這行 \texttt{assert} 會整句消失，變成完全沒有
檢查，配置失敗時 \texttt{arr} 是 \texttt{NULL}，函式若接著寫入
\texttt{arr[0] = ...}，會是解參考 \texttt{NULL} 的未定義行為——比「沒
檢查」更難查，因為程式在偵錯版本裡表面上工作正常（因為那裡 assert 有
效，失敗時直接中止並指出行號），一到發布版就無聲地壞掉。 正確作法是
用一般的 \texttt{if} 檢查，讓配置失敗這個「執行期可能發生的狀況」在任
何編譯模式下都有一致的處理方式（印出錯誤訊息、回傳失敗、或
\texttt{exit}，依場合而定）。

\texttt{assert} 真正適合的場合，是像「這個函式的參數不該是 \texttt{NULL}」
「走到這裡表示前面的邏輯有 bug」這種呼叫端如果守規矩就一定不會發生的
狀況——這一章後面在鏈結串列與樹的操作裡，會用 \texttt{assert} 檢查
「傳進來的指標不該是 \texttt{NULL}」這類前提，而不是檢查
\texttt{malloc} 本身。

\keypoint{\texttt{assert} 抓的是「按邏輯不該發生」的程式設計者錯誤，可
以被 \texttt{-DNDEBUG} 整批關掉。 \texttt{malloc} 失敗是執行期真的可
能發生的事，要用一般的 \texttt{if} 處理，不要用 \texttt{assert}。}

\section{節點與不變式}
\label{sec:ds-node-invariant}

單向鏈結串列（singly linked list）的每個節點包含要存的\textbf{資
料}，以及指向下一個節點的 \texttt{next} 指標。 串列的「開頭」用一個
\texttt{head} 指標記住；空串列時 \texttt{head == NULL}。

\begin{code}{c}{節點型別}
typedef struct Node {
    int data;
    struct Node *next;
} Node;

Node *head = NULL;  /* 空串列 */
\end{code}

這一章要建立的核心詞彙是\textbf{不變式}（invariant）：一句話，描述這
個結構「任何時刻都該成立」的性質。 鏈結串列的不變式是：

\begin{center}
\itshape 從 \texttt{head} 出發，沿著 \texttt{next} 走，經過有限步一定
會走到 \texttt{NULL}；串列裡的每個節點，恰好被前一個節點的
\texttt{next}（或 \texttt{head} 本身）指到一次。
\end{center}

之後每一個操作——插入、刪除、合併——本質上都是同一件事：\textbf{暫時
打破這句話，再想辦法恢復}。 插入時，新節點在被正確接上之前，串列有一
瞬間是「斷開」或「有兩個節點指向同一個地方」的；刪除時，被刪的節點在
真正 \texttt{free} 之前，有一瞬間已經沒有人指向它。 幾乎所有鏈結串列
的 bug，都是「打破了不變式，卻忘記在操作結束前恢復」。 記住這句話，後
面每一節的程式碼都會回來檢查它有沒有被滿足。

\keypoint{鏈結串列的不變式：從 \texttt{head} 沿 \texttt{next} 走有限步
到 \texttt{NULL}，每個節點恰好被指到一次。 每個操作都是「暫時破壞→立
刻恢復」，寫完程式碼後回頭檢查這句話還成不成立，是抓 bug 最快的方
法。}

\section{走訪：迴圈版與遞迴版}
\label{sec:ds-traverse}

走訪整條串列的模式很固定：用一個指標從 \texttt{head} 出發，處理目前
節點，再往下一個走，直到 \texttt{NULL}。

\begin{code}{c}{走訪鏈結串列}
void print_list(Node *head) {
    for (Node *cur = head; cur != NULL; cur = cur->next) {
        printf("%d ", cur->data);
    }
    printf("\n");
}

/* 遞迴版：base case 是空串列，其餘處理目前節點再遞迴處理剩下的部分 */
void print_list_rec(Node *p) {
    if (p == NULL) return;
    printf("%d ", p->data);
    print_list_rec(p->next);
}
\end{code}

遞迴版的形狀，直接對應鏈結串列「一個節點接一個更短的同類結構」這個定
義本身——這跟遞迴那一章講過的「base case ＋ 縮小問題」是同一個思路。
遞迴版每往下一層都會多用一點 stack 空間，串列非常長（例如幾十萬個節
點）時可能耗盡 stack；迴圈版沒有這個限制，實務上處理長串列優先用迴圈
版。

\keypoint{走訪的迴圈版與遞迴版做的事完全一樣，只是實作機制不同：遞迴
版形狀跟串列的遞迴定義一致，但長串列可能耗盡 stack；迴圈版沒有這個風
險，是實務上的預設選擇。}

\section{頭端插入——順序為什麼重要}
\label{sec:ds-head-insert}

在 \texttt{head} 前插入一個新節點，只需要兩個步驟，但\textbf{順序不能
顛倒}：

\begin{center}
\begin{tikzpicture}
    \node[idxlbl] at (0,1.3) {\textbf{1. 初始}};
    \node[hlcell2] (h1) at (0,0.5) {head};
    \node[cell] (b1) at (1.8,0) {5};
    \draw[diagarrow] (h1.south) -- (b1.north);

    \node[idxlbl] at (4.2,1.3) {\textbf{2. p->next = head}};
    \node[hlcell2] (h2) at (4.2,0.5) {head};
    \node[hlcell] (p2) at (4.2,-0.7) {3};
    \node[cell] (b2) at (6.0,0) {5};
    \draw[diagarrow] (h2.south) -- (b2.north);
    \draw[diagarrow] (p2.north) -- (b2.south west);

    \node[idxlbl] at (8.6,1.3) {\textbf{3. head = p}};
    \node[hlcell2] (h3) at (8.6,0.5) {head};
    \node[hlcell] (p3) at (8.6,-0.7) {3};
    \node[cell] (b3) at (10.4,0) {5};
    \draw[diagarrow] (h3.south) -- (p3.north);
    \draw[diagarrow] (p3.east) -- (b3.west);
\end{tikzpicture}
\diagcaption{先讓新節點的 \texttt{next} 指向原本的 \texttt{head}（第 2
格），\textbf{再}把 \texttt{head} 改指向新節點（第 3 格）。 順序寫反的
話，第 2 步做完 \texttt{head} 就已經是新節點自己，舊串列會整條斷開。}
\label{fig:ds-head-insert-order}
\end{center}

\begin{code}{c}{頭端插入}
void push_front(Node **head, int x) {
    Node *p = malloc(sizeof(Node));
    p->data = x;
    p->next = *head; /* 第一步：新節點先接住舊的 head */
    *head = p;        /* 第二步：head 才改指向新節點 */
}
\end{code}

（這裡先用 \texttt{Node **head} 讓函式能真正改到呼叫端的指標，下一節
會完整解釋原因。）如果把這兩行順序寫反：

\begin{example}{c}{(wrong-order-insert.c) 順序寫反的後果}
#include <stdio.h>
#include <stdlib.h>

typedef struct Node {
    int data;
    struct Node *next;
} Node;

int main(void) {
    Node *head = NULL;
    Node *a = malloc(sizeof(Node)); a->data = 1; a->next = NULL;
    head = a;
    Node *b = malloc(sizeof(Node)); b->data = 2; b->next = NULL;
    head = b;        /* 順序寫反：先改了 head */
    b->next = head;  /* head 已經是 b 自己了！ */

    printf("b->next == b ? %s\n", (b->next == b) ? "yes" : "no");
    return 0;
}
\end{example}

\iosample{}{b->next == b ? yes}

先做 \texttt{head = b}，\texttt{head} 已經變成 \texttt{b} 了；再執行
\texttt{b->next = head}，等於是 \texttt{b->next = b}——新節點指向自己，
形成一個獨立的自我循環，原本 \texttt{a} 所在的那條串列從此沒有任何指
標指得到，直接洩漏。 這正是第~\ref{sec:ds-node-invariant}~節說的「暫
時破壞不變式」出了差錯的具體例子：兩行程式碼之間，不變式必須被小心地
維持住，順序寫反就是最常見的維持失敗。

\keypoint{頭端插入永遠是「新節點先接住舊 head，head 才改指向新節點」；
順序寫反會讓新節點指向自己、舊串列整條洩漏。 更新鏈結指標前，先在紙上
畫出操作前、操作中、操作後三格，確認每一步「誰的 next 該指向誰」，再
動手寫程式。}

\section{指標的指標}
\label{sec:ds-pointer-to-pointer}

上一節的 \texttt{push\_front} 用了 \texttt{Node **head}，這一節解釋原
因。 先看一個「看起來合理，其實沒用」的版本：

\begin{example}{c}{(why-ptrptr.c) 為什麼 Node *list 沒用}
#include <stdio.h>
#include <stdlib.h>

typedef struct Node {
    int data;
    struct Node *next;
} Node;

/* 錯誤示範：list 只是呼叫端指標的「複本」 */
void add_bad(Node *list, int x) {
    Node *p = malloc(sizeof(Node));
    p->data = x;
    p->next = list;
    list = p; /* 只改了區域變數 list，呼叫端的指標完全沒被動到 */
}

void add_good(Node **list, int x) {
    Node *p = malloc(sizeof(Node));
    p->data = x;
    p->next = *list;
    *list = p; /* 透過位址，真正改到呼叫端的指標 */
}

int main(void) {
    Node *head = NULL;

    add_bad(head, 1);
    printf("after add_bad:  head is %s\n", head == NULL ? "still NULL" : "not NULL");

    add_good(&head, 1);
    printf("after add_good: head is %s\n", head == NULL ? "still NULL" : "not NULL");

    free(head);
    return 0;
}
\end{example}

\iosample{}{after add\_bad:  head is still NULL\\after add\_good: head is not NULL}

C 的函式參數永遠是\textbf{傳值}——指標也不例外。 呼叫 \texttt{add\_bad(head, 1)}
時，\texttt{list} 這個參數拿到的是 \texttt{head} 這個位址值的\textit{複
本}，跟呼叫端的 \texttt{head} 是兩個不同的變數，只是剛好存著一樣的
值。 函式內部寫 \texttt{list = p;}，只是把這個複本改成指向新節點，呼
叫端的 \texttt{head} 完全沒被動到——函式一返回，複本連同新配置的節點
一起消失，那個節點還洩漏了。 想要真正修改呼叫端的指標，得傳「指標的
位址」，也就是 \texttt{Node **}：\texttt{*list} 解一次參考，拿到的正
是呼叫端 \texttt{head} 這個變數本身，對它賦值才會真正生效。

這正是本章跟「回傳新 head」兩種寫法並存的原因：

\begin{code}{c}{兩種同樣正確的風格}
Node *push_front_v1(Node *head, int x) {   /* 回傳新 head，呼叫端要接回去 */
    Node *p = malloc(sizeof(Node));
    p->data = x;
    p->next = head;
    return p;
}
head = push_front_v1(head, x);

void push_front_v2(Node **head, int x) {   /* 傳位址，函式內部直接改 */
    Node *p = malloc(sizeof(Node));
    p->data = x;
    p->next = *head;
    *head = p;
}
push_front_v2(&head, x);
\end{code}

兩種寫法都正確，差別只在「誰負責把新 head 接回去」。 本章接下來統一
用\texttt{Node **} 風格（下一節的刪除操作會看到它真正的威力），「回傳
新 head」風格留給章末練習對照。

\keypoint{函式參數是傳值，指標也一樣——函式內部改變參數本身（讓它指向
別的地方），呼叫端的指標不會被影響。 想讓函式修改呼叫端的指標，要傳
\texttt{Node **}，透過 \texttt{*pp} 才能碰到那個指標變數本身。}

\section{搜尋與刪除}
\label{sec:ds-search-delete}

搜尋跟走訪幾乎一樣，只是多了「找到就提前返回」：

\begin{code}{c}{搜尋}
Node *search(Node *head, int target) {
    for (Node *p = head; p != NULL; p = p->next) {
        if (p->data == target) return p;
    }
    return NULL;
}
\end{code}

刪除是鏈結串列最容易寫錯的操作——原因是「刪頭」「刪中間」「刪尾」看
起來像三種不同的情況，很多教材也真的分三段程式碼處理。 用
\texttt{Node **} 可以把三種情況統一成同一份程式碼：

\begin{example}{c}{(ptrptr-delete.c) 用 Node ** 統一刪頭／刪中／刪尾}
#include <stdio.h>
#include <stdlib.h>

typedef struct Node {
    int data;
    struct Node *next;
} Node;

void push_front(Node **head, int x) {
    Node *p = malloc(sizeof(Node));
    p->data = x;
    p->next = *head;
    *head = p;
}

/* pp 一路指向「目前正在檢查的那個 next 欄位」（或 head 本身） */
void delete_value(Node **pp, int target) {
    while (*pp != NULL) {
        if ((*pp)->data == target) {
            Node *dead = *pp;
            *pp = dead->next;
            free(dead);
            return;
        }
        pp = &(*pp)->next;
    }
}

void print_list(Node *head) {
    for (Node *cur = head; cur != NULL; cur = cur->next) {
        printf("%d", cur->data);
        if (cur->next != NULL) printf(" ");
    }
    printf("\n");
}

void free_list(Node *head) {
    while (head != NULL) {
        Node *tmp = head;
        head = head->next;
        free(tmp);
    }
}

int main(void) {
    Node *head = NULL;
    int vals[] = {5, 4, 3, 2, 1};
    for (int i = 0; i < 5; i++) push_front(&head, vals[i]);
    print_list(head);

    delete_value(&head, 1); print_list(head); /* 刪頭 */
    delete_value(&head, 3); print_list(head); /* 刪中間 */
    delete_value(&head, 5); print_list(head); /* 刪尾 */
    delete_value(&head, 99); print_list(head); /* 不存在，什麼都不做 */

    free_list(head);
    return 0;
}
\end{example}

\iosample{}{1 2 3 4 5\\2 3 4 5\\2 4 5\\2 4\\2 4}

\texttt{pp} 一開始指向 \texttt{head} 本身；每往下走一格，就讓
\texttt{pp} 指向「目前節點的 \texttt{next} 欄位」的位址——\texttt{pp}
永遠指向「某個東西的 next 欄位，或 head 本身」，這兩種情況在
\texttt{Node **} 的視角下完全一樣，不需要另外判斷「是不是頭」。 找到
目標後，\texttt{*pp = dead->next;} 這一行同時涵蓋了「頭被繞過」跟「前
一個節點的 next 被接到下一個節點」兩種效果，接著 \texttt{free(dead)}
——三種情況、一份程式碼，這正是「指標的指標」真正的威力。

\keypoint{把 \texttt{pp} 想成「一路指向某個 next 欄位（或 head）的位
址」，刪頭、刪中間、刪尾就是同一件事，不需要分三段程式碼特殊處理。}

\section{釋放整條串列}
\label{sec:ds-free-list}

\begin{code}{c}{釋放整條串列}
void free_list(Node *head) {
    while (head != NULL) {
        Node *tmp = head;
        head = head->next;  /* 先移到下一個 */
        free(tmp);           /* 才釋放剛剛那個 */
    }
}

/* 遞迴版 */
void free_list_rec(Node *p) {
    if (p == NULL) return;
    free_list_rec(p->next); /* 先遞迴處理剩下的部分 */
    free(p);                 /* 再釋放自己 */
}
\end{code}

順序同樣不能顛倒：一定要\textbf{先存下 \texttt{next}}，才能
\texttt{free} 目前節點；先 \texttt{free} 再讀 \texttt{->next}，讀到的
是已經被釋放的記憶體，是懸置指標：

\begin{code}{c}{危險：free 之後才讀 next}
while (head != NULL) {
    free(head);
    head = head->next;  /* 危險：head 已經被釋放了，這裡是懸置指標 */
}
\end{code}

這正好呼應第~\ref{sec:ds-memory-bugs}~節的懸置指標——這裡的差別是懸置
的時間極短（只有一行的距離），但一樣是未定義行為，不能靠「反正很快就
換掉了」自我安慰。

\keypoint{釋放整條串列，永遠要先把 \texttt{next} 存到暫存變數，再
\texttt{free} 目前節點；順序反過來就是懸置指標。}

\section{串接（concat）與有序合併（merge）是兩件事}
\label{sec:ds-merge}

「合併兩條串列」常被混成一件事，其實至少有兩種完全不同的操作。

\subsection{串接：直接把一條接到另一條後面}

\begin{code}{c}{concat：串接兩條串列}
Node *concat(Node *list1, Node *list2) {
    if (list1 == NULL) return list2;
    Node *p = list1;
    while (p->next != NULL) p = p->next;
    p->next = list2; /* list1 的尾巴直接接上 list2 */
    return list1;
}
\end{code}

串接之後，\texttt{list1} 跟 \texttt{list2}\textbf{共用}同一段節點
——\texttt{list2} 原本的節點現在也是 \texttt{list1} 的一部分。 這件事
不寫程式很難感覺到危險，直到你 \texttt{free} 兩次：

\begin{example}{c}{(concat-doublefree.c) 串接後分別 free 兩條串列}
#include <stdio.h>
#include <stdlib.h>

typedef struct Node {
    int data;
    struct Node *next;
} Node;

Node *make_list(int *vals, int n) {
    Node *head = NULL, *tail = NULL;
    for (int i = 0; i < n; i++) {
        Node *p = malloc(sizeof(Node));
        p->data = vals[i];
        p->next = NULL;
        if (head == NULL) head = tail = p;
        else { tail->next = p; tail = p; }
    }
    return head;
}

Node *concat(Node *list1, Node *list2) {
    if (list1 == NULL) return list2;
    Node *p = list1;
    while (p->next != NULL) p = p->next;
    p->next = list2;
    return list1;
}

void free_list(Node *head) {
    while (head != NULL) {
        Node *tmp = head;
        head = head->next;
        free(tmp);
    }
}

int main(void) {
    int a[] = {1, 2, 3};
    int b[] = {4, 5};
    Node *list1 = make_list(a, 3);
    Node *list2 = make_list(b, 2);

    concat(list1, list2); /* list1 現在包含了 list2 的所有節點 */

    free_list(list1);  /* 把 1,2,3,4,5 全部 free 掉 */
    free_list(list2);  /* 危險：4,5 已經被上面那行 free 過了！ */
    return 0;
}
\end{example}

這支程式在這台機器上會直接異常中止（雙重釋放被配置器的安全檢查抓
到）。 串接之後，\texttt{list2} 這個名字雖然還在，指向的節點卻已經被
併進 \texttt{list1} 的責任範圍——串接之後只該用\textit{一條}串列的名
字去釋放，\texttt{list2} 這個變數繼續存在只是留給你誤用的陷阱。

\subsection[有序合併：兩條遞增串列合併]{有序合併：兩條各自遞增的串列，合併成一條遞增的串列}

\begin{code}{c}{merge：有序合併，用 dummy node 骨架}
Node *merge_sorted(Node *a, Node *b) {
    Node dummy;
    Node *tail = &dummy;
    dummy.next = NULL;

    while (a != NULL && b != NULL) {
        if (a->data <= b->data) {
            tail->next = a;
            a = a->next;
        } else {
            tail->next = b;
            b = b->next;
        }
        tail = tail->next;
    }
    tail->next = (a != NULL) ? a : b;
    return dummy.next;
}
\end{code}

\texttt{dummy} 是一個「假的頭節點」，宣告在 stack 上（不是
\texttt{malloc} 出來的，函式結束前也不需要、也不能 \texttt{free}
它），純粹用來讓「插入第一個節點」跟「插入後面的節點」變成同一種操
作，不用另外判斷「這是不是第一次接」。 \texttt{merge\_sorted} 直接重
新利用 \texttt{a}、\texttt{b} 原本的節點（不額外配置新節點），結果串
列的每個節點原本屬於 \texttt{a} 或 \texttt{b} 其中一條——這是它跟串接
的共同點：兩者都不新建節點，都會導致原本兩條串列的名字不能再各自獨立
使用。

\keypoint{串接（concat）把一條接到另一條後面，不管順序；有序合併
（merge）把兩條遞增的串列交錯接成一條遞增的串列，靠 dummy node 讓第一
次插入跟後續插入用同一份程式碼。 兩者都會讓原本兩條串列共用節點，合
併後只能用\textit{一個}名字去釋放。}

\section{二元搜尋樹：不變式是對整棵子樹}
\label{sec:ds-bst-define}

二元樹的每個節點最多有兩個子節點：\texttt{left}、\texttt{right}。

\begin{code}{c}{樹節點型別}
typedef struct TreeNode {
    int data;
    struct TreeNode *left, *right;
} TreeNode;
\end{code}

\textbf{二元搜尋樹}（BST）的不變式是：對\textbf{每一個}節點，左子樹裡
\textit{所有}的值都小於它、右子樹裡\textit{所有}的值都大於它。 關鍵字
是「所有」——不是只跟直接子節點比較。 一個常見的誤解是只檢查直接父子
關係就以為驗證完了：

\begin{example}{c}{(bst-invariant-check.c) 只檢查直接子節點是不夠的}
#include <stdio.h>
#include <stdlib.h>
#include <limits.h>

typedef struct TreeNode {
    int data;
    struct TreeNode *left, *right;
} TreeNode;

TreeNode *make(int v, TreeNode *l, TreeNode *r) {
    TreeNode *p = malloc(sizeof(TreeNode));
    p->data = v; p->left = l; p->right = r;
    return p;
}

/* 錯的檢查：只比較「直接子節點」，看起來很合理，其實不夠 */
int is_bst_wrong(TreeNode *root) {
    if (root == NULL) return 1;
    if (root->left && root->left->data >= root->data) return 0;
    if (root->right && root->right->data <= root->data) return 0;
    return is_bst_wrong(root->left) && is_bst_wrong(root->right);
}

/* 對的檢查：對「整棵子樹」成立，用範圍往下傳 */
int is_bst_range(TreeNode *root, long lo, long hi) {
    if (root == NULL) return 1;
    if (root->data <= lo || root->data >= hi) return 0;
    return is_bst_range(root->left, lo, root->data)
        && is_bst_range(root->right, root->data, hi);
}

int main(void) {
    /*        5
             / \
            3   8
               / \
              4   9   <- 4 在 5 的右子樹裡，卻比 5 小：違反 BST 性質 */
    TreeNode *root = make(5,
        make(3, NULL, NULL),
        make(8, make(4, NULL, NULL), make(9, NULL, NULL)));

    printf("is_bst_wrong: %d\n", is_bst_wrong(root));
    printf("is_bst_range: %d\n", is_bst_range(root, LONG_MIN, LONG_MAX));
    return 0;
}
\end{example}

\iosample{}{is\_bst\_wrong: 1\\is\_bst\_range: 0}

節點 \texttt{4} 是 \texttt{8} 的左子節點，\texttt{4 < 8} 沒有錯，
\texttt{is\_bst\_wrong} 因此判定為合法；但 \texttt{4} 同時也在根節點
\texttt{5} 的\textit{右}子樹裡，而右子樹的每個值都該大於 \texttt{5}
——\texttt{4} 違反了這一點。 \texttt{is\_bst\_range} 把「這個節點允許
的範圍」一路往下傳（一開始是 $(-\infty, \infty)$，往左子樹傳
$(lo, root)$、往右子樹傳 $(root, hi)$），才是對「整棵子樹」成立的正確
檢查。

\keypoint{BST 的不變式是對\textit{整棵子樹}成立，不是只比較直接子節
點。 驗證一棵樹是否為合法 BST，要用「允許範圍往下傳」的寫法，不能只
看每個節點跟它兩個孩子的大小關係。}

\section{建立 BST：插入}
\label{sec:ds-bst-insert}

\begin{code}{c}{BST 插入（遞迴寫法）}
TreeNode *insert(TreeNode *root, int x) {
    if (root == NULL) {
        TreeNode *p = malloc(sizeof(TreeNode));
        p->data = x;
        p->left = p->right = NULL;
        return p;
    }
    if (x < root->data) {
        root->left = insert(root->left, x);
    } else if (x > root->data) {
        root->right = insert(root->right, x);
    }
    /* x == root->data：重複鍵，本書統一選擇忽略，不重複插入 */
    return root;
}
\end{code}

從根往下比較，比根小往左走、比根大往右走，直到走到 \texttt{NULL} 才
真正配置新節點；每一層回傳的都是「這個子樹更新後的根」，所以每一層都
要用 \texttt{root->left = insert(root->left, x)} 把新根接回去，不能只
呼叫 \texttt{insert} 而不接回傳值——這是本節唯一需要牢記的細節，漏寫
賦值，插入的節點會憑空消失（呼叫端拿不到更新後的子樹根）。 \textbf{重
複鍵的處理，全書統一選擇「忽略」}：\texttt{x == root->data} 時，兩個
分支都不成立，直接回傳原本的 \texttt{root}，不會插入新節點、也不會有
記憶體洩漏。

\keypoint{插入時，每一層都要把子樹的新根接回
\texttt{root->left}／\texttt{root->right}，漏接會讓插入的節點消失。 重
複鍵一律忽略——這是本書統一的選擇，其他書籍可能選擇往右塞或允許重
複，遇到時留意題目自己的規定。}

\section{三種遍歷：只差 printf 放哪一行}
\label{sec:ds-traversals}

前序、中序、後序三種走訪，遞迴結構完全一樣，唯一的差別是「拜訪」（這
裡是 \texttt{printf}）放在遞迴呼叫的哪個位置：

\begin{code}{c}{前序／中序／後序}
void preorder(TreeNode *root) {
    if (root == NULL) return;
    printf("%d ", root->data);  /* 拜訪放最前面 */
    preorder(root->left);
    preorder(root->right);
}

void inorder(TreeNode *root) {
    if (root == NULL) return;
    inorder(root->left);
    printf("%d ", root->data);  /* 拜訪放中間 */
    inorder(root->right);
}

void postorder(TreeNode *root) {
    if (root == NULL) return;
    postorder(root->left);
    postorder(root->right);
    printf("%d ", root->data);  /* 拜訪放最後面 */
}
\end{code}

\begin{center}
\begin{tikzpicture}
    \node[hlcell2] (r) at (2.0,0) {5};
    \node[cell] (l) at (0.8,-1.0) {3};
    \node[cell] (r2) at (3.2,-1.0) {8};
    \node[idxlbl, anchor=west] at (4.0,0) {前序拜訪 5 的時機：一進來就印};
    \node[idxlbl, anchor=west] at (4.0,-0.5) {中序拜訪 5 的時機：左子樹處理完才印};
    \node[idxlbl, anchor=west] at (4.0,-1.0) {後序拜訪 5 的時機：兩子樹都處理完才印};
    \draw[diagarrow] (r.south west) -- (l.north);
    \draw[diagarrow] (r.south east) -- (r2.north);
\end{tikzpicture}
\diagcaption{同一個節點 5，三種走訪只差「印出它」的時機在遞迴呼叫的前
面、中間、還是後面；左右子樹本身也遞迴地遵守同一個順序。}
\label{fig:ds-traversal-timing}
\end{center}

三者裡最重要的性質是：\textbf{對 BST 做中序走訪，會依遞增順序印出所
有的鍵}。 這不是巧合——BST 的不變式保證每個節點左邊都比它小、右邊都
比它大，中序走訪「先左、後自己、後右」剛好完全對應「先印小的、再印自
己、再印大的」。 這個性質後面重建與集合運算都會用到。

\keypoint{三種走訪的遞迴結構完全一樣，只差 \texttt{printf} 放在遞迴呼
叫的前／中／後面。 中序走訪 BST 會得到遞增序列，這是「用結構隱含排
序」的核心性質，記住它比背三段程式碼更有用。}

\section{二元樹重建：從遍歷結果反推樹}
\label{sec:ds-tree-rebuild}

走訪是把樹「壓平」成一串數字；反過來，能不能從一串（或兩串）數字重建
出原本的樹？ 答案是：\textbf{前序 加 中序}可以，\textbf{前序 加 後
序}不行（後面會說明原因）。

前序的第一個元素，一定是這整棵（子）樹的根——這是前序「先印自己」的
定義保證的。 找到這個根在中序序列裡的位置，中序裡它\textit{左邊}的部
分就是左子樹的中序、右邊就是右子樹的中序；左子樹有幾個節點，前序裡緊
接著根之後的那幾個，就是左子樹的前序。 用這兩條資訊切開兩個子序列，
分別遞迴，重建左右子樹。

\begin{example}{c}{(rebuild-tree.c) 從前序＋中序重建二元樹}
#include <stdio.h>
#include <stdlib.h>

typedef struct TreeNode {
    int data;
    struct TreeNode *left, *right;
} TreeNode;

/* preorder[preStart..preEnd] 與 inorder[inStart..inEnd] 都是閉區間 */
TreeNode *rebuild(int *preorder, int preStart, int preEnd,
                   int *inorder, int inStart, int inEnd) {
    if (preStart > preEnd) return NULL;

    int root_val = preorder[preStart]; /* 前序的第一個一定是這棵子樹的根 */
    TreeNode *root = malloc(sizeof(TreeNode));
    root->data = root_val;

    int mid = inStart;
    while (inorder[mid] != root_val) mid++; /* 在中序裡找到根的位置 */
    int left_size = mid - inStart;

    root->left = rebuild(preorder, preStart + 1, preStart + left_size,
                          inorder, inStart, mid - 1);
    root->right = rebuild(preorder, preStart + left_size + 1, preEnd,
                           inorder, mid + 1, inEnd);
    return root;
}

void postorder_print(TreeNode *root) {
    if (root == NULL) return;
    postorder_print(root->left);
    postorder_print(root->right);
    printf("%d ", root->data);
}

void free_tree(TreeNode *root) {
    if (root == NULL) return;
    free_tree(root->left);
    free_tree(root->right);
    free(root);
}

int main(void) {
    int preorder[] = {3, 1, 2, 5, 4, 6};
    int inorder[]  = {1, 2, 3, 4, 5, 6};
    int n = 6;

    TreeNode *root = rebuild(preorder, 0, n - 1, inorder, 0, n - 1);
    postorder_print(root); /* 用另一種走訪印出來，驗證重建的樹是對的 */
    printf("\n");
    free_tree(root);
    return 0;
}
\end{example}

\iosample{}{2 1 4 6 5 3}

用後序印出重建好的樹，而不是直接印前序或中序（那樣只是把輸入原封不動
印回去，證明不了重建對不對）——後序的結果只有樹的形狀真的對，才會是
\texttt{2 1 4 6 5 3}，是驗證重建演算法正確性很自然的方式。

\textbf{為什麼前序＋後序不夠？} 前序告訴你「根是誰」，但不能告訴你根
的\textit{哪些}子節點屬於左子樹、哪些屬於右子樹——如果某個節點只有一
個孩子，前序＋後序沒辦法分辨這個孩子是接在 \texttt{left} 還是
\texttt{right}。 中序之所以夠，是因為它把「比根小」跟「比根大」明確地
分在根的兩側，這條資訊前序跟後序都沒有單獨提供。

\keypoint{前序的第一個元素是根；在中序裡定位根，左邊是左子樹的中序、
右邊是右子樹的中序，左子樹的節點數決定前序要在哪裡切開。 前序＋中序
才能唯一重建整棵樹，前序＋後序做不到，因為缺少「哪些節點屬於左子樹」
這條資訊。}

\section{BST 的交集與聯集}
\label{sec:ds-bst-set-ops}

第~\ref{sec:ds-traversals}~節已經看到「中序走訪 BST 得到遞增序列」這
個性質，這裡直接派上用場：把兩棵 BST 各自中序走訪成排序好的陣列，剩
下就是排序陣列常見的\textbf{雙指標合併}，時間複雜度 $O(n+m)$——不需要
對每個節點各自搜尋一次（那樣是 $O(nm)$，而且順序不會是排序好的）。

\begin{example}{c}{(bst-intersect-union.c) 中序轉陣列後雙指標求交集聯集}
#include <stdio.h>
#include <stdlib.h>

typedef struct TreeNode {
    int data;
    struct TreeNode *left, *right;
} TreeNode;

TreeNode *insert(TreeNode *root, int x) {
    if (root == NULL) {
        TreeNode *p = malloc(sizeof(TreeNode));
        p->data = x;
        p->left = p->right = NULL;
        return p;
    }
    if (x < root->data) root->left = insert(root->left, x);
    else if (x > root->data) root->right = insert(root->right, x);
    return root;
}

void inorder_to_array(TreeNode *root, int *arr, int *idx) {
    if (root == NULL) return;
    inorder_to_array(root->left, arr, idx);
    arr[(*idx)++] = root->data;
    inorder_to_array(root->right, arr, idx);
}

void free_tree(TreeNode *root) {
    if (root == NULL) return;
    free_tree(root->left);
    free_tree(root->right);
    free(root);
}

int main(void) {
    int a_vals[] = {5, 2, 8, 1, 9, 3};
    int b_vals[] = {2, 3, 4, 8, 10};
    TreeNode *t1 = NULL, *t2 = NULL;
    for (int i = 0; i < 6; i++) t1 = insert(t1, a_vals[i]);
    for (int i = 0; i < 5; i++) t2 = insert(t2, b_vals[i]);

    int arr1[6], arr2[5], n1 = 0, n2 = 0;
    inorder_to_array(t1, arr1, &n1); /* 中序 = 排序好的陣列 */
    inorder_to_array(t2, arr2, &n2);

    printf("intersection:");
    int i = 0, j = 0;
    while (i < n1 && j < n2) {
        if (arr1[i] < arr2[j]) i++;
        else if (arr1[i] > arr2[j]) j++;
        else { printf(" %d", arr1[i]); i++; j++; }
    }
    printf("\n");

    printf("union:");
    i = 0; j = 0;
    while (i < n1 || j < n2) {
        int v;
        if (i >= n1) v = arr2[j++];
        else if (j >= n2) v = arr1[i++];
        else if (arr1[i] < arr2[j]) v = arr1[i++];
        else if (arr1[i] > arr2[j]) v = arr2[j++];
        else { v = arr1[i]; i++; j++; }
        printf(" %d", v);
    }
    printf("\n");

    free_tree(t1);
    free_tree(t2);
    return 0;
}
\end{example}

\iosample{}{intersection: 2 3 8\\union: 1 2 3 4 5 8 9 10}

交集：兩個指標各自指向一個陣列，指到的值較小的那邊指標前進；相等就是
交集的成員，兩邊指標一起前進。 聯集：一樣雙指標往前走，較小的那個先
印出來，相等時只印一次、兩邊指標一起前進。 兩段程式碼的骨架幾乎一
樣，都是排序陣列合併的標準寫法——這正是「中序＝排序」這個性質真正有
用的地方：把樹的問題，轉換成陣列的問題，陣列上成熟的技巧就都能直接
用。

\keypoint{BST 求交集或聯集，不要對每個節點各自搜尋一次（$O(nm)$，順序
還亂）；先各自中序走訪成排序陣列，再用雙指標合併，是 $O(n+m)$ 且結果
自然有序的正確做法。}

\section{常見錯誤清單}
\label{sec:ds-checklist}

下面每一條都對應前面某一節。

\begin{itemize}
\item 回傳指向區域變數的位址：函式一返回，那塊 stack 空間就可能被覆
寫；編譯器通常會警告，但程式仍可能「看起來正常」。
（第~\ref{sec:ds-heap-stack}~節）
\item \texttt{malloc} 不需要轉型；用 \texttt{sizeof(*p)} 而不是
\texttt{sizeof(型別名)}；\texttt{realloc} 要接暫存指標，不能直接蓋掉
原指標。
（第~\ref{sec:ds-malloc-free}~節）
\item 洩漏、懸置指標、雙重釋放是三種不同的錯誤；前兩種通常不會馬上讓
程式壞掉，這正是它們難抓的原因。
（第~\ref{sec:ds-memory-bugs}~節）
\item 不要用 \texttt{assert} 檢查 \texttt{malloc} 的回傳值——配置失敗
是執行期可能發生的事，不是邏輯錯誤，而且 \texttt{-DNDEBUG} 後檢查會消
失。
（第~\ref{sec:ds-assert}~節）
\item 頭端插入要先接住舊 \texttt{head} 再改 \texttt{head}，順序寫反會
讓新節點指向自己、舊串列洩漏。
（第~\ref{sec:ds-head-insert}~節）
\item 函式參數是傳值；要修改呼叫端的指標，得傳 \texttt{Node **}。
（第~\ref{sec:ds-pointer-to-pointer}~節）
\item 釋放串列要先存下 \texttt{next} 再 \texttt{free}，順序反過來是懸
置指標。
（第~\ref{sec:ds-free-list}~節）
\item 串接（concat）與有序合併（merge）都會讓兩條串列共用節點，合併
後只能用一個名字釋放，兩邊都釋放會雙重釋放。
（第~\ref{sec:ds-merge}~節）
\item BST 的不變式是對整棵子樹成立，只檢查直接子節點不夠。
（第~\ref{sec:ds-bst-define}~節）
\item BST 插入每一層都要把子樹新根接回
\texttt{root->left}／\texttt{root->right}，漏接插入的節點會消失。
（第~\ref{sec:ds-bst-insert}~節）
\item 前序＋中序才能唯一重建二元樹，前序＋後序不夠，因為分不出哪些節
點屬於左子樹。
（第~\ref{sec:ds-tree-rebuild}~節）
\item BST 交集聯集用中序轉陣列＋雙指標合併，不要對每個節點各自搜尋。
（第~\ref{sec:ds-bst-set-ops}~節）
\end{itemize}

\phantomsection
\addcontentsline{toc}{section}{練習解答}
\section*{練習解答}

\paragraph{第~\ref{sec:ds-pointer-to-pointer}~節}
\begin{enumerate}
\item 把本節的 \texttt{push\_front\_v1}（回傳新 head）風格，改寫成刪
除操作：寫一個 \texttt{Node *delete\_v1(Node *head, int target)}，回
傳更新後的新 head。 提示：跟 \texttt{Node **} 版比起來，這個版本要不
要特別處理「刪的剛好是 head」這個情況？

\textit{答}：需要。 因為沒有 \texttt{Node **} 可以直接改到呼叫端的
\texttt{head}，函式只能「回傳」新的 head，所以必須明確判斷
\texttt{head->data == target} 這個特例，跟刪除中間節點分開處理——這正
是 \texttt{Node **} 版能把三種情況合併成一份程式碼的原因所在。
\end{enumerate}

\paragraph{第~\ref{sec:ds-bst-insert}~節}
\begin{enumerate}
\item 如果 BST 插入時把 \texttt{root->left = insert(root->left, x);}\\
漏寫成只呼叫 \texttt{insert(root->left, x);}（不接回傳值），對一棵原
本是空樹的 BST 插入第一筆資料，會發生什麼事？

\textit{答}：\texttt{insert} 內部確實配置了新節點並回傳它，但呼叫端沒
有接住這個回傳值，\texttt{root} 本身沒有被更新，函式結束後
\texttt{root} 仍然是 \texttt{NULL}——插入的節點形同消失（而且沒有任何
指標指向它，直接洩漏）。
\end{enumerate}

\paragraph{第~\ref{sec:ds-tree-rebuild}~節}
\begin{enumerate}
\item 承接 \texttt{rebuild-tree.c} 的例子（前序
\texttt{3 1 2 5 4 6}、中序 \texttt{1 2 3 4 5 6}），手算這棵樹的中序走
訪結果，並說明為什麼答案一定跟輸入的中序陣列一模一樣。

\textit{答}：一模一樣，就是 \texttt{1 2 3 4 5 6}——因為
\texttt{rebuild} 本來就是拿這份中序陣列切出左右子樹，重建出來的樹結
構，中序走訪回去，必然照原樣把這份陣列走一遍。 這也是為什麼驗證重建
正不正確要用\textit{另一種}走訪（例如後序）而不是中序：中序無論樹的
形狀對不對，走出來都會跟輸入相同，驗證不出形狀是否正確。
\end{enumerate}

\chapter{物件}
\label{ch:objects}

上一章回答的是「資料怎麼\emph{串起來}」；這一章回答的是另一個問題：
「程式碼怎麼\emph{分開}」。 一支程式只要稍微大一點，把所有東西塞進一
個 \texttt{main} 就會開始難以維護——狀態變數散落各處、同樣的邏輯被複
製貼上好幾份、換一種內部實作就要把用到它的地方全部重寫一遍。 C 沒有
\texttt{class} 關鍵字，但只要善用\textbf{標頭檔}、\textbf{struct}、
\textbf{分別編譯}，一樣可以做到「物件」真正重要的那些好處：把狀態包
起來、把介面和實作分開、換實作不影響呼叫端。

這一章只用一個例子貫穿到底：一個堆疊（stack）模組。 不是為了教堆疊
這個資料結構本身（上一章已經打好指標與 heap 的基礎），而是拿堆疊當
\textbf{載體}，一步一步示範「從寫死在 \texttt{main} 裡的陣列，走到一
個可以換掉內部實作、呼叫端完全不用改的模組」這整個過程。 全章的高潮
發生在中段：同一份 \texttt{main.c}、同一份 \texttt{stack.h}，只換掉
\texttt{stack.c}，重新編譯就能跑——那一刻你會真正感覺到「物件」帶來
的好處是什麼，而不是背一句抽象定義。

\section{從一個具體的痛開始}
\label{sec:obj-pain}

括號平衡是堆疊最經典的應用：從左到右掃過字串，遇到 \texttt{(} 就推
進堆疊，遇到 \texttt{)} 就彈出一個；掃完後堆疊該剛好清空。 最直接的
寫法，把堆疊整個寫死在 \texttt{main} 裡：

\begin{example}{c}{(paren1.c) 堆疊直接寫死在 main 裡}
#include <stdio.h>

int main(void) {
    char line[1000];
    if (!fgets(line, sizeof line, stdin)) return 0;

    int stack[1000];
    int top = 0;
    int ok = 1;

    for (int i = 0; line[i] != '\0' && line[i] != '\n'; i++) {
        char c = line[i];
        if (c == '(') {
            stack[top++] = c;
        } else if (c == ')') {
            if (top == 0) { ok = 0; break; }
            top--;
        }
    }
    if (top != 0) ok = 0;

    printf("%s\n", ok ? "balanced" : "not balanced");
    return 0;
}
\end{example}

\iosample{(a(b)c)}{balanced}

能跑，也不難懂。 但現在假設題目多了一個要求：字串裡同時要平衡
\texttt{()} 跟 \texttt{[]} 兩種括號，各自獨立判斷。 最直覺的做法是再
開一組 \texttt{stack2}、\texttt{top2}，把上面判斷括號的邏輯複製貼上一
份，只把變數名字換掉。 再多一種括號，再複製一份。 這份程式碼裡「堆
疊」這件事——怎麼推進、怎麼彈出、怎麼判斷空——其實完全沒變，變的只
是「用哪個陣列、哪個索引」。 每次要多一個堆疊，就要重新複製貼上整套
邏輯，這正是這一章要解決的痛點：\textbf{堆疊的操作邏輯，應該只寫一
次，寫成「可以重複開新的堆疊來用」的樣子}。

\section{把狀態包成 struct}
\label{sec:obj-struct}

第一步很小：把 \texttt{stack} 陣列跟 \texttt{top} 這兩個「總是綁在一
起使用」的變數，包進一個 \texttt{struct}：

\begin{code}{c}{把陣列與 top 包成一個 Stack}
#define STACK_CAP 1000

typedef struct {
    int data[STACK_CAP];
    int top;
} Stack;
\end{code}

原本各自散落的操作函式，通通改成吃一個 \texttt{Stack *} 參數：

\begin{code}{c}{操作函式改吃 Stack *}
void stack_init(Stack *s) {
    s->top = 0;
}

void stack_push(Stack *s, int x) {
    s->data[s->top] = x;
    s->top++;
}

int stack_pop(Stack *s) {
    s->top--;
    return s->data[s->top];
}
\end{code}

現在要開第二個堆疊，只要再宣告一個 \texttt{Stack} 變數，呼叫同一組函
式、傳不同的位址進去就好——不用複製貼上任何邏輯：

\begin{code}{c}{開兩個獨立的堆疊}
Stack s1, s2;
stack_init(&s1);
stack_init(&s2);
stack_push(&s1, 'a');
stack_push(&s2, 'b'); /* s1、s2 互不影響 */
\end{code}

\keypoint{把「總是綁在一起用」的變數包進一個 \texttt{struct}，操作函式
改吃這個 struct 的指標——這是把重複邏輯變成「可以開多份」的第一步，也
是本章所有後續技巧的地基。}

\section{介面：函式宣告先於實作}
\label{sec:obj-interface}

上一節的 \texttt{stack\_init}／\texttt{stack\_push}／\texttt{stack\_pop}
還需要再補兩個常用操作：判斷是否為空、判斷是否已滿。 五個函式的\textbf{原
型}湊在一起，就是這個堆疊模組的\textbf{介面}（interface）——只描述「能
做什麼」，不涉及「怎麼做到」：

\begin{code}{c}{堆疊模組的介面：五個函式原型}
void stack_init(Stack *s);
void stack_push(Stack *s, int x);
int  stack_pop(Stack *s);
int  stack_is_empty(const Stack *s);
int  stack_is_full(const Stack *s);
\end{code}

這一節要做一件聽起來奇怪的事：\textbf{只靠這五個原型，先把括號平衡的
\texttt{main.c} 完整寫出來}——這幾個函式的本體到現在都還沒寫。

\begin{code}{c}{main.c：只靠原型，函式本體還不存在}
Stack s;
stack_init(&s);
int ok = 1;

for (int i = 0; line[i] != '\0' && line[i] != '\n'; i++) {
    char c = line[i];
    if (c == '(') {
        stack_push(&s, c);
    } else if (c == ')') {
        if (stack_is_empty(&s)) { ok = 0; break; }
        stack_pop(&s);
    }
}
if (!stack_is_empty(&s)) ok = 0;
\end{code}

這支 \texttt{main.c} 能不能\textit{編譯}？ 能——編譯器只要看過原型，
就知道每個函式接受什麼參數、回傳什麼型別，足夠產生呼叫它們的機器碼。
能不能\textit{連結成執行檔}？ 不能，因為函式本體還不存在，連結器找不
到它們真正的位址。 這個「能編譯、不能連結」的中間狀態，正是這一節要
凸顯的重點：\textbf{寫使用端程式碼的人，完全不需要知道實作怎麼寫，只
需要一份介面}。 這正是\textbf{抽象}（abstraction）這個詞在這裡的具體
意思——不是什麼玄妙的哲學，就是「先講清楚能做什麼，再決定怎麼做到」。

\keypoint{只靠函式原型（介面），就能把使用端的程式碼完整寫出來、甚至
編譯出目的檔——不需要等實作寫完。 這件事本身就是「抽象」：把「能做什
麼」跟「怎麼做到」拆成兩個可以分開處理的問題。}

\section{標頭檔與 include guard}
\label{sec:obj-header}

上一節那五個原型，該放進一個 \texttt{.h} 檔，讓每個要用堆疊的
\texttt{.c} 檔都能 \texttt{\#include} 它：

\begin{example}{c}{(stack.h) 堆疊模組的標頭檔}
#ifndef STACK_H
#define STACK_H

#define STACK_CAP 1000

typedef struct {
    int data[STACK_CAP];
    int top;
} Stack;

void stack_init(Stack *s);
void stack_push(Stack *s, int x);
int  stack_pop(Stack *s);
int  stack_is_empty(const Stack *s);
int  stack_is_full(const Stack *s);

#endif
\end{example}

\texttt{\#ifndef}／\texttt{\#define}／\texttt{\#endif} 這層 include
guard 前置處理那一章已經講過原理，這裡不重複；記住每個自己寫的標頭檔
都該包這一層就好。 \texttt{Stack} 這個型別定義也放進標頭檔，因為呼叫
端要宣告 \texttt{Stack s;} 這種區域變數，編譯器得知道它確切多大。

真正的實作放進對應的 \texttt{stack.c}：

\begin{example}{c}{(stack.c) 堆疊模組的實作}
#include "stack.h"

void stack_init(Stack *s) {
    s->top = 0;
}

void stack_push(Stack *s, int x) {
    s->data[s->top] = x;
    s->top++;
}

int stack_pop(Stack *s) {
    s->top--;
    return s->data[s->top];
}

int stack_is_empty(const Stack *s) {
    return s->top == 0;
}

int stack_is_full(const Stack *s) {
    return s->top == STACK_CAP;
}
\end{example}

如果 \texttt{stack.c} 裡還想要一個只給自己內部用、不想讓其他檔案呼叫
的輔助函式，在它前面加上 \texttt{static}：

\begin{code}{c}{static：只在這個檔案裡看得見}
static void log_internal_state(const Stack *s) {
    /* 只有 stack.c 自己會呼叫這個函式，其他 .c 檔看不到它、也連結不到它 */
}
\end{code}

檔案範圍的 \texttt{static} 給的是\textbf{內部連結}（internal linkage）
——這個名字只在目前這個翻譯單元裡看得到，就算另一個 \texttt{.c} 檔剛
好也定義了一個同名函式，也不會撞名。 這是模組「只暴露該暴露的東西」
的第一層工具，之後第~\ref{sec:obj-opaque}~節會看到更徹底的做法。

\keypoint{標頭檔放「呼叫端需要知道的東西」（型別定義、函式原型），實
作放對應的 \texttt{.c} 檔；每個標頭檔都要包 include guard。 檔案範圍的
\texttt{static} 讓一個名字只在該檔案內部可見，不會跟其他檔案撞名。}

\section{分別編譯}
\label{sec:obj-compile}

\texttt{stack.h}、\texttt{stack.c}、\texttt{main.c} 三個檔案，怎麼變
成一支執行檔？ 前置處理那一章已經看過「編譯→組譯→連結」這條流水線；
分別編譯的意思是：\texttt{stack.c} 跟 \texttt{main.c} \textbf{各自獨
立}走過前面兩站，變成各自的目的檔，最後才由連結器把它們接起來。

\begin{center}
\begin{tikzpicture}[node distance=1.0cm and 0.8cm]
    \node[startend, minimum width=2.0cm, font=\small\sffamily] (s1) {stack.c};
    \node[process, right=of s1, minimum width=2.6cm, font=\small\sffamily] (c1) {編譯／組譯};
    \node[startend, right=of c1, minimum width=2.0cm, font=\small\sffamily] (o1) {stack.o};

    \node[startend, below=of s1, minimum width=2.0cm, font=\small\sffamily] (s2) {main.c};
    \node[process, right=of s2, minimum width=2.6cm, font=\small\sffamily] (c2) {編譯／組譯};
    \node[startend, right=of c2, minimum width=2.0cm, font=\small\sffamily] (o2) {main.o};

    \node[process, right=3.2cm of o1, minimum width=2.0cm, font=\small\sffamily, yshift=-0.65cm] (link) {連結};
    \node[startend, right=of link, minimum width=2.0cm, font=\small\sffamily] (exe) {執行檔};

    \draw[flowarrow] (s1) -- (c1);
    \draw[flowarrow] (c1) -- (o1);
    \draw[flowarrow] (s2) -- (c2);
    \draw[flowarrow] (c2) -- (o2);
    \draw[flowarrow] (o1) -- (link);
    \draw[flowarrow] (o2) -- (link);
    \draw[flowarrow] (link) -- (exe);
\end{tikzpicture}
\diagcaption{\texttt{stack.c} 與 \texttt{main.c} 各自獨立編譯成目的
檔，互不干涉；只有連結這一步才把兩份目的檔接成一支執行檔。 只改了
\texttt{stack.c}，就只需要重新編譯 \texttt{stack.o} 再重新連結，不用
動 \texttt{main.o}。}
\label{fig:obj-separate-compile}
\end{center}

\begin{code}{bash}{手動分別編譯再連結}
gcc -c stack.c -o stack.o    # 只編譯，不連結，產生目的檔
gcc -c main.c -o main.o
gcc stack.o main.o -o paren  # 連結兩個目的檔成執行檔
\end{code}

\texttt{-c} 讓 \texttt{gcc} 只做到目的檔這一步就停下來。 大型專案通
常不會手動打這幾行，而是寫一份 \texttt{Makefile}：

\begin{code}{text}{Makefile（最小示意）}
paren: stack.o main.o
	gcc stack.o main.o -o paren

stack.o: stack.c stack.h
	gcc -c stack.c -o stack.o

main.o: main.c stack.h
	gcc -c main.c -o main.o
\end{code}

\texttt{make} 會比對檔案的修改時間，只重新編譯真的改過的 \texttt{.c}
檔，沒改的直接沿用舊的 \texttt{.o}——這也是分別編譯除了「介面與實作分
開」之外，另一個實務上的好處：專案越大，省下的重複編譯時間越可觀。

\keypoint{分別編譯讓每個 \texttt{.c} 檔各自變成目的檔，連結才把它們接
起來；只改了其中一個檔案，只需要重新編譯那一個，不用整個專案重編。
\texttt{Makefile} 把這件事自動化。}

\section{應用一：括號平衡}
\label{sec:obj-app-paren}

有了 \texttt{stack.h}，第~\ref{sec:obj-pain}~節那支括號平衡程式可以完
整改寫，而且從頭到尾\textbf{只 \texttt{\#include "stack.h"}}，不需要
知道 \texttt{stack.c} 內部怎麼實作：

\begin{example}{c}{(main\_paren.c) 只碰 stack.h 的括號平衡}
#include <stdio.h>
#include "stack.h"

int main(void) {
    char line[1000];
    if (!fgets(line, sizeof line, stdin)) return 0;

    Stack s;
    stack_init(&s);
    int ok = 1;

    for (int i = 0; line[i] != '\0' && line[i] != '\n'; i++) {
        char c = line[i];
        if (c == '(') {
            stack_push(&s, c);
        } else if (c == ')') {
            if (stack_is_empty(&s)) { ok = 0; break; }
            stack_pop(&s);
        }
    }
    if (!stack_is_empty(&s)) ok = 0;

    printf("%s\n", ok ? "balanced" : "not balanced");
    return 0;
}
\end{example}

\iosample{(a(b)c)}{balanced}

\section{應用二：對稱字串}
\label{sec:obj-app-palindrome}

第二個應用一樣只碰 \texttt{stack.h}：把字串整個推進堆疊，再一個一個
彈出來跟原字串從頭比對——彈出的順序剛好是原字串的反序，兩者完全相同
就是回文（palindrome）：

\begin{example}{c}{(main\_palindrome.c) 只碰 stack.h 的回文判斷}
#include <stdio.h>
#include <string.h>
#include "stack.h"

int main(void) {
    char line[1000];
    if (!fgets(line, sizeof line, stdin)) return 0;
    int len = (int)strlen(line);
    if (len > 0 && line[len - 1] == '\n') line[--len] = '\0';

    Stack s;
    stack_init(&s);
    for (int i = 0; i < len; i++) stack_push(&s, line[i]);

    int ok = 1;
    for (int i = 0; i < len; i++) {
        int c = stack_pop(&s);
        if (c != line[i]) { ok = 0; break; }
    }

    printf("%s\n", ok ? "palindrome" : "not palindrome");
    return 0;
}
\end{example}

\iosample{racecar}{palindrome}

兩個應用程式合計沒有出現任何 \texttt{s.data[...]} 或 \texttt{s.top}
這種直接碰觸內部欄位的寫法，一律透過 \texttt{stack\_push}／
\texttt{stack\_pop}／\texttt{stack\_is\_empty} 操作。 這不是巧合，是
下一節「換實作」能夠成立的前提。

\keypoint{只要應用程式全程只透過介面（五個函式）操作
\texttt{Stack}，不直接碰觸內部欄位，介面背後的實作要怎麼換都不會影響
到它——這正是下一節要展示的事。}

\section{換一個實作：換掉 stack.c，main.c 不用改}
\label{sec:obj-swap}

這一節是全章的高潮：把 \texttt{stack.c} 整份換掉，\texttt{stack.h} 一
個字都不改，\texttt{main.c} 也一個字都不改，重新編譯就能跑。 換掉的
方向，是把固定大小的陣列，換成\textbf{用得完才長大}的動態陣列——原本
的 \texttt{Stack} 得先把 \texttt{data} 從「固定大小的陣列欄位」改成
「指標＋容量」，這一步標頭檔要跟著微調：

\begin{example}{c}{(stack2.h) 用指標＋容量取代固定陣列}
#ifndef STACK_H
#define STACK_H

typedef struct {
    int *data;
    int top;
    int capacity;
} Stack;

void stack_init(Stack *s);
void stack_push(Stack *s, int x);
int  stack_pop(Stack *s);
int  stack_is_empty(const Stack *s);
int  stack_is_full(const Stack *s);
void stack_free(Stack *s);

#endif
\end{example}

多了一個 \texttt{stack\_free}（動態配置的記憶體要記得釋放），介面其
餘五個函式的原型完全沒變。 現在準備兩份 \texttt{stack.c}，都實作同一
份介面：

\begin{example}{c}{(stack\_fixed.c) 固定容量版}
#include <stdlib.h>
#include "stack2.h"

#define FIXED_CAP 1000

void stack_init(Stack *s) {
    s->data = malloc(FIXED_CAP * sizeof(int));
    s->top = 0;
    s->capacity = FIXED_CAP;
}

void stack_push(Stack *s, int x) {
    s->data[s->top] = x;
    s->top++;
}

int stack_pop(Stack *s) {
    s->top--;
    return s->data[s->top];
}

int stack_is_empty(const Stack *s) { return s->top == 0; }
int stack_is_full(const Stack *s)  { return s->top == s->capacity; }

void stack_free(Stack *s) {
    free(s->data);
    s->data = NULL;
}
\end{example}

\begin{example}{c}{(stack\_dynamic.c) 用完就長大，永遠不會滿}
#include <stdlib.h>
#include "stack2.h"

#define INITIAL_CAP 4

void stack_init(Stack *s) {
    s->data = malloc(INITIAL_CAP * sizeof(int));
    s->top = 0;
    s->capacity = INITIAL_CAP;
}

void stack_push(Stack *s, int x) {
    if (s->top == s->capacity) {
        s->capacity *= 2;
        s->data = realloc(s->data, s->capacity * sizeof(int));
    }
    s->data[s->top++] = x;
}

int stack_pop(Stack *s) {
    s->top--;
    return s->data[s->top];
}

int stack_is_empty(const Stack *s) { return s->top == 0; }
int stack_is_full(const Stack *s)  { (void)s; return 0; }

void stack_free(Stack *s) {
    free(s->data);
    s->data = NULL;
}
\end{example}

\texttt{main\_paren.c}（只需要在原本的版本結尾補上一行
\texttt{stack\_free(\&s);}）分別跟這兩份 \texttt{stack.c} 連結：

\begin{code}{bash}{同一份 main.c，換掉 stack.c 重新編譯}
gcc -c stack_fixed.c -o stack.o
gcc -c main_paren.c -o main.o
gcc stack.o main.o -o paren   # 固定容量版

gcc -c stack_dynamic.c -o stack.o   # 換掉這一步
gcc stack.o main.o -o paren         # main.o 完全不用重新編譯！
\end{code}

兩次連結出來的 \texttt{paren} 跑起來結果完全一樣。 差別只有：固定容
量版遇到超過 1000 層括號巢狀就會寫壞記憶體，動態版永遠不會滿，代價是
每次擴充要重新配置、搬移資料。 這個抉擇——要不要為了容量彈性多付一
點效能代價——完全被藏在 \texttt{stack.c} 內部，\texttt{main.c} 完全不
需要知道，也不需要改。

\keypoint{只要應用程式全程只透過介面操作，換掉背後的實作
（\texttt{stack.c}），呼叫端的原始碼（\texttt{main.c}）完全不用改，
重新編譯（甚至只需要重新連結）就能跑。 這是「介面與實作分開」最直接
的回報。}

\section{不完整型別：更徹底的資訊隱藏}
\label{sec:obj-opaque}

上一節換了容量的管理策略，但 \texttt{Stack} 的欄位形狀
（\texttt{data}／\texttt{top}／\texttt{capacity}）沒有變；如果想換成
完全不同的內部表示——例如上一章講過的\textbf{鏈結串列}——這些欄位名
字就完全用不上了，而且只要 \texttt{Stack} 的定義還寫在標頭檔裡，任何
看得到標頭檔的人也看得到內部欄位，可以繞過介面直接寫
\texttt{s.top = 0}，這種「偷看」在小程式裡看不出問題，專案一大、內部
表示一換，所有偷看內部欄位的呼叫端都會一起壞掉。

\textbf{不完整型別}（incomplete type）能徹底擋掉這件事：標頭檔只
\texttt{typedef} 這個型別的\textit{名字}，不寫出它的內容：

\begin{example}{c}{(stack\_opaque.h) 只宣告型別名字，不定義內容}
#ifndef STACK_H
#define STACK_H

typedef struct Stack Stack; /* 不完整型別：只宣告，不定義內部長什麼樣 */

Stack *stack_create(void);
void   stack_destroy(Stack *s);
void   stack_push(Stack *s, int x);
int    stack_pop(Stack *s);
int    stack_is_empty(const Stack *s);

#endif
\end{example}

看到這份標頭檔的人，只知道「\texttt{Stack} 是某種東西」，完全不知道
它裡面有什麼欄位——編譯器也一樣不知道，所以呼叫端不能再宣告
\texttt{Stack s;} 這種區域變數（編譯器不知道要留多少空間），改成一律
透過 \texttt{stack\_create} 拿一個 \texttt{Stack *}、用完呼叫
\texttt{stack\_destroy} 釋放。 這正是這個技巧的代價：\textbf{呼叫端從
此只能透過指標使用這個型別}，換來的是內部表示可以隨便換，標頭檔完全
不用動。 真正的 \texttt{struct Stack} 定義，只出現在實作的 \texttt{.c}
檔裡，而且兩種完全不同的內部表示可以並存：

\begin{example}{c}{(stack\_opaque\_array.c) 動態陣列版，struct 定義只有這裡看得到}
#include <stdlib.h>
#include "stack_opaque.h"

struct Stack {
    int *data;
    int top;
    int capacity;
};

Stack *stack_create(void) {
    Stack *s = malloc(sizeof(Stack));
    s->capacity = 4;
    s->data = malloc(s->capacity * sizeof(int));
    s->top = 0;
    return s;
}

void stack_destroy(Stack *s) {
    free(s->data);
    free(s);
}

void stack_push(Stack *s, int x) {
    if (s->top == s->capacity) {
        s->capacity *= 2;
        s->data = realloc(s->data, s->capacity * sizeof(int));
    }
    s->data[s->top++] = x;
}

int stack_pop(Stack *s) {
    s->top--;
    return s->data[s->top];
}

int stack_is_empty(const Stack *s) { return s->top == 0; }
\end{example}

\begin{example}{c}{(stack\_opaque\_linked.c) 鏈結串列版，跟上面的 struct Stack 完全不同}
#include <stdlib.h>
#include "stack_opaque.h"

typedef struct StackNode {
    int data;
    struct StackNode *next;
} StackNode;

struct Stack {
    StackNode *top;
};

Stack *stack_create(void) {
    Stack *s = malloc(sizeof(Stack));
    s->top = NULL;
    return s;
}

void stack_destroy(Stack *s) {
    while (s->top != NULL) {
        StackNode *dead = s->top;
        s->top = s->top->next;
        free(dead);
    }
    free(s);
}

void stack_push(Stack *s, int x) {
    StackNode *p = malloc(sizeof(StackNode));
    p->data = x;
    p->next = s->top;
    s->top = p;
}

int stack_pop(Stack *s) {
    StackNode *dead = s->top;
    int v = dead->data;
    s->top = dead->next;
    free(dead);
    return v;
}

int stack_is_empty(const Stack *s) { return s->top == NULL; }
\end{example}

兩份 \texttt{struct Stack} 定義的欄位完全不同（一個是指標＋容量、一個
是鏈結節點），但因為呼叫端只透過 \texttt{Stack *} 這個不透明的指標操
作，\texttt{main.c} 連 \texttt{.o} 檔都不需要重新編譯，只要換一顆
\texttt{stack\_*.o} 重新連結就能跑：

\begin{code}{bash}{main.o 完全不用重新編譯，只重新連結}
gcc -c main_opaque.c -o main.o          # main.o 只編譯這一次

gcc -c stack_opaque_array.c -o stack.o
gcc stack.o main.o -o paren             # 陣列版

gcc -c stack_opaque_linked.c -o stack.o # 換成完全不同的資料結構
gcc stack.o main.o -o paren             # main.o 沒有被動過
\end{code}

Richard Reese 的 \textit{Understanding and Using C Pointers}（見書末
第~\ref{ch:references}~章）對這個技巧有一句很精準的說法：只把使用一
個資料結構所需要的資訊暴露給使用者；暴露太多內部細節，使用者遲早會
依賴這些細節，之後想改內部實作就會弄壞使用者的程式碼。 不完整型別把
這句話變成編譯器強制的規則，不只是君子協定。

\keypoint{不完整型別（\texttt{typedef struct Stack Stack;} 只宣告不定
義）讓標頭檔完全不透漏內部欄位，呼叫端只能透過指標使用；代價是要改
用 create／destroy 動態配置，換來的是內部表示可以換成任何形狀（陣列、
鏈結串列⋯）都不影響呼叫端，連目的檔都不用重新編譯。}

\section{錯誤處理該由誰負責}
\label{sec:obj-errors}

\texttt{stack\_pop} 如果被呼叫在一個空堆疊上，該怎麼辦？ 這個問題沒有
單一「正確」答案，是設計時要做的取捨，常見選項有四種：

\begin{center}
\begin{tabular}{p{2.6cm}p{5.0cm}p{4.2cm}}
\toprule
做法 & 說明 & 取捨 \\
\midrule
哨兵值 & 回傳一個資料範圍外的特殊值（例如 \texttt{-1}），呼叫端自己檢
查 & 簡單，但如果 \texttt{-1} 剛好是合法資料，會混淆 \\
設 \texttt{errno} & 回傳值照舊，另外設一個全域錯誤碼 & 呼叫端很容易忘
記檢查，跟標準函式庫的慣例一致 \\
\texttt{assert} & 直接假設呼叫端不會犯這個錯，錯了就中止 & 只適合「呼
叫端的邏輯錯誤」，不適合處理使用者輸入 \\
輸出參數＋\texttt{bool} & 回傳 \texttt{bool} 表示成功與否，真正的值透
過指標參數帶出來 & 呼叫端一定要檢查回傳值才拿得到資料，介面稍微複雜
一點點 \\
\bottomrule
\end{tabular}
\end{center}

前三種在第~\ref{ch:preprocessing}、\ref{ch:stdlib}~章都看過對應的例
子。 這裡示範第四種——\textbf{輸出參數＋bool}——因為它最適合「呼叫端
一定要能區分成功與失敗，而且不能靠特殊值蒙混過去」的場合：

\begin{code}{c}{輸出參數 + bool}
bool stack_pop(Stack *s, int *out) {
    if (stack_is_empty(s)) return false;
    s->top--;
    *out = s->data[s->top];
    return true;
}

/* 呼叫端一定要檢查回傳值，才知道 *out 有沒有意義 */
int value;
if (stack_pop(&s, &value)) {
    printf("%d\n", value);
} else {
    printf("stack is empty\n");
}
\end{code}

這種寫法的好處是編譯器幫不上忙、但至少介面本身「誠實」：拿不到
\texttt{bool} 的回傳值，就沒有正當理由去用 \texttt{*out}。 本章前面
的例子選用最簡單的「呼叫端自己先用 \texttt{stack\_is\_empty} 檢查」
風格，是為了保持範例精簡；實務上寫一個要給別人用的模組，輸出參數＋
\texttt{bool} 通常是更穩妥的選擇。

\keypoint{錯誤處理沒有單一正確答案，是設計取捨：哨兵值簡單但可能跟合
法值混淆、\texttt{errno} 容易被忘記檢查、\texttt{assert} 只適合邏輯錯
誤、輸出參數＋\texttt{bool} 最誠實但介面稍微複雜。 選一種在整個模組裡
貫徹，不要同一個模組裡的函式各用一套。}

\section{物件的優點：回頭看做過的事}
\label{sec:obj-summary}

「物件」常常先給一串抽象名詞：抽象、封裝、資訊隱藏、多型。 這一節反
過來，把每個詞回指到前面已經動手做過的具體段落：

\begin{itemize}
\item \textbf{抽象}（abstraction）：第~\ref{sec:obj-interface}~節只靠
函式原型就把 \texttt{main.c} 寫完——先講清楚「能做什麼」，不用等「怎
麼做到」寫完。
\item \textbf{封裝}（encapsulation）：第~\ref{sec:obj-header}~節把狀
態（\texttt{Stack} 的欄位）跟操作它的函式綁在一起，透過
\texttt{.h}／\texttt{.c} 分工，標頭檔只放呼叫端需要的東西。
\item \textbf{資訊隱藏}（information hiding）：第~\ref{sec:obj-opaque}~
節用不完整型別，讓呼叫端連內部欄位長什麼樣都看不到，只能透過介面操
作。
\item \textbf{多型}（polymorphism）：下一節（選讀）會示範怎麼用函式
指標，讓同一段呼叫端程式碼，依實際傳進來的資料型態，執行不同的實
作。
\end{itemize}

四個詞背後都是同一件事的不同切面：\textbf{把「呼叫端需要知道的東
西」跟「實作者需要自由改動的東西」，用一條清楚的界線分開}。 界線畫
在哪裡（介面該包含幾個函式、要不要用不完整型別、錯誤處理選哪一
種），是每次設計模組時真正要做決定的地方，沒有一套公式能自動算出正
確答案。

\keypoint{抽象、封裝、資訊隱藏、多型不是四個要背的定義，是「介面與實
作之間那條界線」在不同情境下的四種說法。 回頭看本章走過的每一步，每
個詞都對應到一段你已經親手寫過、編譯過、跑過的程式碼。}

\section{進階：用函式指標做多型（選讀）}
\label{sec:obj-polymorphism}

前面的堆疊只處理 \texttt{int}，呼叫端永遠知道自己在跟哪一種
\texttt{Stack} 打交道。 有些情境需要反過來：呼叫端寫\textit{同一段}
程式碼，但依實際傳進來的資料型態，執行不同的操作——這正是「多型」在
C 裡最常見的實現方式：把一張\textbf{函式指標表}放進 struct 的第一個
欄位。

\begin{example}{c}{(shape-poly.c) 函式指標模擬 vtable}
#include <stdio.h>

typedef struct Shape {
    double (*area)(const struct Shape *self);
    double (*perimeter)(const struct Shape *self);
} Shape;

typedef struct {
    Shape base; /* 一定要是第一個欄位 */
    double width, height;
} Rectangle;

double rect_area(const Shape *self) {
    const Rectangle *r = (const Rectangle *)self;
    return r->width * r->height;
}
double rect_perimeter(const Shape *self) {
    const Rectangle *r = (const Rectangle *)self;
    return 2 * (r->width + r->height);
}
Rectangle make_rectangle(double w, double h) {
    Rectangle r;
    r.base.area = rect_area;
    r.base.perimeter = rect_perimeter;
    r.width = w;
    r.height = h;
    return r;
}

typedef struct {
    Shape base;
    double radius;
} Circle;

double circle_area(const Shape *self) {
    const Circle *c = (const Circle *)self;
    return 3.14159265358979 * c->radius * c->radius;
}
double circle_perimeter(const Shape *self) {
    const Circle *c = (const Circle *)self;
    return 2 * 3.14159265358979 * c->radius;
}
Circle make_circle(double r) {
    Circle c;
    c.base.area = circle_area;
    c.base.perimeter = circle_perimeter;
    c.radius = r;
    return c;
}

void print_shape_info(const Shape *s) {
    printf("area=%.2f perimeter=%.2f\n", s->area(s), s->perimeter(s));
}

int main(void) {
    Rectangle r = make_rectangle(3, 4);
    Circle c = make_circle(2);
    print_shape_info((const Shape *)&r);
    print_shape_info((const Shape *)&c);
    return 0;
}
\end{example}

\iosample{}{area=12.00 perimeter=14.00\\area=12.57 perimeter=12.57}

關鍵在 \texttt{Rectangle} 與 \texttt{Circle} 都把 \texttt{Shape base}
放在 struct 的\textbf{第一個}欄位——C 標準保證一個 struct 的位址跟它
第一個欄位的位址相同，所以 \texttt{(const Shape *)\&r} 這個轉型是安全
的：透過 \texttt{Shape *} 存取，看到的正是 \texttt{r.base} 那兩個函式
指標。 \texttt{print\_shape\_info} 完全不知道傳進來的是矩形還是圓
形，只透過 \texttt{s->area(s)} 呼叫「這個形狀自己填好的那個函式」
——同一行呼叫，依實際型態跑出不同的計算，這正是多型的核心。 在物件導
向語言裡，編譯器會自動幫你產生並維護這張表（通常稱為 vtable）；C 沒
有這個語法糖，需要自己手動填。

\keypoint{把函式指標表放進 struct 的第一個欄位，讓不同的具體型別
（\texttt{Rectangle}、\texttt{Circle}）都能安全地被轉型成基底型別
（\texttt{Shape}）操作——這是 C 裡手動模擬「多型」最常見的寫法，物件
導向語言的 vtable 做的正是同一件事，只是編譯器自動幫你維護。}

\section{常見錯誤清單}
\label{sec:obj-checklist}

下面每一條都對應前面某一節。

\begin{itemize}
\item 同樣的邏輯散落在多份複製貼上的程式碼裡，是這一章要解決的起點症
狀；把狀態包進 struct、操作改吃這個 struct 的指標，是第一步。
（第~\ref{sec:obj-struct}~節）
\item 只靠函式原型（介面）就能把使用端程式碼寫完、編譯出目的檔，不需
要等實作寫完——這是抽象的具體意思，不是需要死背的名詞。
（第~\ref{sec:obj-interface}~節）
\item 每個自己寫的標頭檔都要包 include guard；檔案範圍的
\texttt{static} 讓一個名字只在該檔案內部可見。
（第~\ref{sec:obj-header}~節）
\item 只改了一個 \texttt{.c} 檔，只需要重新編譯那一個再重新連結，不用
整個專案重編；\texttt{Makefile} 把這件事自動化。
（第~\ref{sec:obj-compile}~節）
\item 應用程式全程只透過介面操作，不直接碰觸內部欄位，是「換實作不影
響呼叫端」能夠成立的前提。
（第~\ref{sec:obj-app-paren}、\ref{sec:obj-app-palindrome}~節）
\item 換掉 \texttt{stack.c} 的內部實作，只要介面（\texttt{stack.h}）
沒變，\texttt{main.c} 完全不用改。
（第~\ref{sec:obj-swap}~節）
\item 標頭檔裡把 struct 定義完整寫出來，呼叫端就能繞過介面直接碰觸內
部欄位；不完整型別把這條路徹底擋掉，代價是呼叫端只能透過指標使用。
（第~\ref{sec:obj-opaque}~節）
\item 錯誤處理沒有單一正確答案，是設計取捨，要在整個模組裡選一種貫
徹，不要同一個模組各用一套。
（第~\ref{sec:obj-errors}~節）
\end{itemize}

\phantomsection
\addcontentsline{toc}{section}{練習解答}
\section*{練習解答}

\paragraph{第~\ref{sec:obj-swap}~節}
\begin{enumerate}
\item \texttt{stack\_dynamic.c} 版本裡，如果把 \texttt{stack\_push}
裡「容量不夠就 \texttt{realloc}」那段判斷拿掉，只留
\texttt{s->data[s->top++] = x;}，會發生什麼事？

\textit{答}：一開始（容量 4 以內）完全正常，一旦推進的元素超過目前
的 \texttt{capacity}，就會寫到配置範圍之外——這是陣列越界寫入，未定
義行為，可能悄悄損壞其他資料，也可能讓程式當掉，取決於那塊記憶體後
面剛好放了什麼。
\end{enumerate}

\paragraph{第~\ref{sec:obj-opaque}~節}
\begin{enumerate}
\item 為什麼不完整型別的版本裡，呼叫端不能再寫
\texttt{Stack s;} 這種區域變數宣告，一定要改成
\texttt{Stack *s = stack\_create();}？

\textit{答}：宣告一個區域變數，編譯器要在編譯期就知道它確切占多少位
元組，才能在 stack 上留出對應的空間；不完整型別只宣告了名字，沒有欄
位資訊，編譯器不知道 \texttt{Stack} 有多大，所以不允許宣告這種變
數。 透過 \texttt{stack\_create} 回傳的指標則沒有這個問題，因為配置
多少記憶體是實作 \texttt{.c} 檔自己知道、自己處理的事，呼叫端只需要
拿到一個位址。
\end{enumerate}

\chapter{高等技術}
\label{ch:extendedtechniques}

到目前為止，我們用的型別大致上是 \texttt{int}、\texttt{char}、\texttt{double}
這些基本型別，加上陣列與函式。 實務上還有幾種 C 語言提供的「進階型別
工具」。 \texttt{typedef}、\texttt{enum}、\texttt{union}。 以及跟「變數能活
多久、誰能看到、在記憶體裡怎麼排」有關的 storage class 與對齊規則。 這些
觀念在題庫裡常常跟位元操作一起出現（例如用 \texttt{typedef} 包一個集合型
別、再用 \texttt{unsigned long long} 存 bitset），但本身並不是位元技巧；
這一章先把型別與記憶體布局的基礎打好，下一章再專心處理位元運算。

\section{typedef：替型別取別名}
\label{sec:ext-typedef}

\texttt{typedef} 讓我們替既有的型別取一個新名字。 編譯器看見
\texttt{Score} 時，仍然把它當成 \texttt{int} 來處理。 不會多占記憶體、也不
會改變運算規則；差別只在於程式碼讀起來更像在描述「這個 \texttt{int} 代表
什麼意義」。 例如考試分數、座標、計數器，都可以各自 typedef 成不同的名
字，讓函式簽名一眼看出參數的語意。

\begin{code}{c}{替 int 取別名 Score}
typedef int Score;

Score a, b;
scanf("%d %d", &a, &b);
printf("%d\n", a + b);
\end{code}

範例 \thechapter.2 把兩個分數讀進來相加。 如果之後題目改成「平均分要印
小數」，我們只需要把 \texttt{typedef int Score} 改成
\texttt{typedef double Score}，其餘變數名稱 \texttt{a}、\texttt{b} 不必
全部重命名。 這正是 typedef 在維護大型程式時的價值：型別的「物理表示」可以
換，但「語意上的名字」可以留著。

\begin{example}{c}{(type-alias.c) 用 Score 表示整數分數}
#include <stdio.h>

typedef int Score;

int main(void) {
    Score a, b;
    scanf("%d %d", &a, &b);
    printf("%d\n", a + b);
}
\end{example}

\iosample{
3 4 \\
}{
7 \\
}

\keypoint{\texttt{typedef} 不創造新型別，只是替既有型別取別名；編譯器仍
依底層型別決定大小與運算，但程式碼裡的名字可以表達設計意圖。}

題庫裡較後段的集合 API 會用 \texttt{typedef struct set \{ \ldots \} Set;}
把「內部怎麼存 bitset」藏起來，對外只暴露 \texttt{Set} 這個名字。  typedef
在這裡扮演的角色，跟本節的 \texttt{Score} 完全一樣，只是包裝的底層型別
從 \texttt{int} 換成了 \texttt{struct}。

\begin{drills}[label=drills:ext-typedef]
\item \texttt{typedef int Score;} 之後寫 \texttt{Score a;}，\texttt{a}
實際的底層型別是什麼？ 編譯器看到 \texttt{Score} 時怎麼處理它？
\item 若要把 \texttt{Score} 底層型別從 \texttt{int} 換成 \texttt{double}，
需要改幾處程式碼？ 這是 \texttt{typedef} 在維護程式時的什麼價值？
\item 題庫裡 \texttt{typedef struct set \{ ... \} Set;} 跟本節的
\texttt{typedef int Score;}，兩者裡 \texttt{typedef} 扮演的角色一樣嗎？
\end{drills}

\section{enum：具名整數常數}
\label{sec:ext-enum}

有時候程式需要在少數幾個固定選項之間做選擇。 例如兩位候選人誰得票較
多、星期幾、方向上下左右。 我們當然可以直接用 \texttt{0}、\texttt{1}、
\texttt{2} 這些魔術數字，但讀程式的人得靠註解或變數名才能猜出
\texttt{1} 到底代表什麼。 \texttt{enum}（enumeration，列舉）讓編譯器自
動替這些選項分配整數常數，並給每個常數一個名字。

\begin{code}{c}{宣告列舉型別與列舉常數}
typedef enum { CAND_A, CAND_B } Candidate;

Candidate winner;
if (votesA > votesB) {
    winner = CAND_A;
} else {
    winner = CAND_B;
}
\end{code}

沒有特別指定值的時候，第一個列舉常數是 \texttt{0}，之後依序遞增：
\texttt{CAND\_A} 是 \texttt{0}、\texttt{CAND\_B} 是 \texttt{1}。 範例
\thechapter.4 模擬一場只有兩位候選人的公投：讀入兩個得票數，決定勝者，
再印出勝者代號（\texttt{A} 或 \texttt{B}）與總票數。 邏輯跟「比大小」
沒什麼兩樣，差別在於 \texttt{winner} 的型別是 \texttt{Candidate} 而
不是裸的 \texttt{int}，我們比較時用的是 \texttt{CAND\_A} 這種有名字的
常數，而不是硬記「0 代表 A、1 代表 B」。

\begin{example}{c}{(referendum-winner.c) 用 enum 表示候選人並決定勝者}
#include <stdio.h>

typedef enum { CAND_A, CAND_B } Candidate;

int main(void) {
    int votesA, votesB;
    scanf("%d %d", &votesA, &votesB);

    Candidate winner;
    if (votesA > votesB) {
        winner = CAND_A;
    } else {
        winner = CAND_B;
    }

    char label = (winner == CAND_A) ? 'A' : 'B';
    printf("%c %d\n", label, votesA + votesB);
}
\end{example}

\iosample{
100 200 \\
}{
B 300 \\
}

以範例的輸入為例：\texttt{votesA = 100}、\texttt{votesB = 200}，條件
\texttt{votesA > votesB} 不成立，所以 \texttt{winner = CAND\_B}，對應
的字元代號是 \texttt{'B'}，總票數 \texttt{100 + 200 = 300}，輸出
\texttt{B 300}。 如果改成 \texttt{500 123}，則 \texttt{winner} 會是
\texttt{CAND\_A}，輸出 \texttt{A 623}。

\stylepoint{需要把 enum 常數轉成字元或字串給 \texttt{printf} 時，通常
還得另外寫一個 \texttt{if} 或 \texttt{switch}；enum 解決的是「程式內部用
具名常數做判斷」，不是自動變成可印的字串。}

\begin{drills}[label=drills:ext-enum]
\item \texttt{typedef enum \{ CAND\_A, CAND\_B \} Candidate;} 沒有特別
指定值時，\texttt{CAND\_A} 和 \texttt{CAND\_B} 分別是多少？
\item 範例讀入 \texttt{500 123} 時，\texttt{winner} 會是哪個常數？ 輸出
是什麼？
\item 能不能直接寫 \texttt{printf("\%s", winner)} 印出 \texttt{"CAND\_A"}
這樣的名字？ 為什麼？
\end{drills}

\section{union：同一段記憶體的不同解讀}
\label{sec:ext-union}

\texttt{struct} 把多個欄位\emph{各自}占一段記憶體排在一起；\texttt{union}
則讓多個欄位\emph{共用}同一段記憶體。 同一時刻只有其中一個欄位「有效」，
但你可以用不同的型別去\emph{解讀}同一段位元組。 這在 low-level 程式裡很
常見：例如把一個 \texttt{int} 拆成四個 \texttt{unsigned char} 去看每一
個位元組，或把浮點數的位元模式當整數讀出來（這在浮點數那一章我們刻意
沒做，因為牽涉 IEEE 754 細節，但 union 正是標準做法）。

\begin{code}{c}{union 讓 int 與 char 陣列共用記憶體}
union IntBytes {
    int value;
    unsigned char bytes[4];
};

union IntBytes x;
scanf("%d", &x.value);
printf("%u\n", x.bytes[0]);
\end{code}

union 的大小，等於它所有成員中\emph{最大}的那一個所占的空間（這裡是
\texttt{int} 的 $4$ 位元組）。 寫入 \texttt{x.value} 之後，再讀
\texttt{x.bytes[0]} 到 \texttt{x.bytes[3]}，看到的其實是同一個
\texttt{int} 在記憶體裡的逐位元組分解。 順序取決於 CPU 是 little-endian
還是 big-endian，這也是為什麼 union 常跟「記憶體布局」一起教。

\begin{center}
\begin{tikzpicture}
    \node[hlcell] (b0) {b0};
    \node[hlcell, right=0pt of b0] (b1) {b1};
    \node[hlcell, right=0pt of b1] (b2) {b2};
    \node[hlcell, right=0pt of b2] (b3) {b3};
    \coordinate (bMid) at ($(b0.south)!0.5!(b3.south)$);
    \node[idxlbl, below=0.2cm of bMid] {\texttt{x.value}（同一個 int）};
    \node[idxlbl, anchor=west, align=left, right=0.35cm of b3] {也可寫成\\ \texttt{x.bytes[0..3]}};
\end{tikzpicture}
\diagcaption{\texttt{union IntBytes} 的四個位元組與單一 \texttt{int}
共用同一段記憶體；改 \texttt{value} 會連帶改變 \texttt{bytes} 的內容。}
\label{fig:ch21-union-bytes}
\end{center}

範例 \thechapter.6 讀入一個整數，依序印出它的四個位元組（由低位元組到
高位元組）。 在常見的 little-endian 機器上，輸入 \texttt{256} 會得到
\texttt{0 1 0 0}：因為 $256 = 0 \times 1 + 1 \times 256 + 0 + 0$，最低
位元組是 \texttt{0}，次低位元組是 \texttt{1}。

\begin{example}{c}{(union-bytes.c) 用 union 檢視 int 的各個位元組}
#include <stdio.h>

union IntBytes {
    int value;
    unsigned char bytes[4];
};

int main(void) {
    union IntBytes x;
    scanf("%d", &x.value);

    for (int i = 0; i < 4; i++) {
        if (i) {
            printf(" ");
        }
        printf("%u", x.bytes[i]);
    }
    printf("\n");
}
\end{example}

\iosample{
256 \\
}{
0 1 0 0 \\
}

\keypoint{union 的所有成員共享同一段記憶體；寫入一個成員後再讀另一個成
員，讀到的是同一位元模式的\emph{不同解讀}，不是複製出來的第二份資料。}

\begin{drills}[label=drills:ext-union]
\item \texttt{union IntBytes \{ int value; unsigned char bytes[4]; \}}
這個 union 的 \texttt{sizeof} 是多少？
\item 在 little-endian 機器上，輸入 \texttt{256} 給 \texttt{x.value}，
\texttt{x.bytes[0]} 到 \texttt{x.bytes[3]} 分別是多少？
\item \texttt{struct} 跟 \texttt{union} 的差別是什麼——各自的成員各占一
段記憶體，還是共用同一段？
\end{drills}

\section{storage class：變數能活多久、誰能看到}
\label{sec:ext-storage-class}

每個變數除了型別之外，還有一個 storage class（儲存類別），決定它「在
哪裡分配」「能活多久」「哪些函式看得到」。 初學時最常碰到的是
\texttt{auto}（區域變數的預設）跟 \texttt{static}；之後寫多檔案程式會再
遇到 \texttt{extern}。

\begin{code}{c}{區域 auto 變數 vs. static 區域變數}
void foo() {
    int autoVar = 0;      // 每次呼叫都重新建立，離開函式就消失
    autoVar++;

    static int persist = 0;  // 只初始化一次，值在多次呼叫之間保留
    persist++;
}
\end{code}

\texttt{auto} 區域變數存在每次函式呼叫的 stack frame 裡：函式一返回，這
一層堆疊就被整個丟棄，變數也跟著消失。 \texttt{static} 區域變數雖然
「寫在函式裡面」，但記憶體不在 stack 上反覆建立／銷毀。 編譯器只會在
程式啟動時初始化一次，之後每次呼叫都沿用同一塊儲存空間。 這跟函式那一
章講的 call stack 可以對照：一般區域變數是跟著 stack frame 來來去去，
\texttt{static} 區域變數則像躲在函式後面的「永久抽屜」。

範例 \thechapter.8 連續呼叫 \texttt{countCalls} 三次，每次印出累計次
數 \texttt{1}、\texttt{2}、\texttt{3}。 如果把 \texttt{static int count}
改成一般的 \texttt{int count}，每次呼叫都會從 \texttt{0} 重新開始，三
次都會印 \texttt{1}。

\begin{example}{c}{(static-counter.c) static 區域變數保留呼叫次數}
#include <stdio.h>

void countCalls(void) {
    static int count = 0;
    count++;
    printf("%d\n", count);
}

int main(void) {
    int n;
    scanf("%d", &n);
    for (int i = 0; i < n; i++) {
        countCalls();
    }
}
\end{example}

\iosample{
3 \\
}{
1 \\
2 \\
3 \\
}

\texttt{extern} 則是「這個變數的\emph{定義}在別的檔案，這裡只是宣告我
要用它」。 片語 \thechapter.9 示範典型寫法：在標頭檔用 \texttt{extern}
宣告，在某個 \texttt{.c} 檔真正定義並初始化。 這讓多個原始檔可以共享
同一個全域變數，又不會在連結時因為重複定義而報錯。

\begin{code}{c}{extern 宣告 vs. 另一個檔案中的定義}
// shared.h
extern int globalCounter;

// fileA.c
#include "shared.h"
int globalCounter = 0;   // 唯一的定義

// fileB.c
#include "shared.h"
void bump() { globalCounter++; }  // 使用 extern 宣告的全域變數
\end{code}

\stylepoint{\texttt{static} 也可以修飾「檔案層級的全域變數或函式」，意
思變成「只在這個 \texttt{.c} 檔內可見」。 跟函式內的 \texttt{static}
區域變數是不同用法，但關鍵字相同，讀程式時要看它修飾的是區域變數還是
全域名稱。}

\begin{drills}[label=drills:ext-storage-class]
\item \texttt{countCalls} 裡的 \texttt{static int count} 若改成一般的
\texttt{int count}，連續呼叫三次會印出什麼？ 跟原本的 \texttt{1 2 3}
差在哪裡？
\item \texttt{auto} 區域變數存在哪裡？ 函式返回後它會發生什麼事？
\item \texttt{extern} 宣告和它真正的定義，通常分別放在什麼樣的檔案裡？
\end{drills}

\section{動態記憶體管理簡介}
\label{sec:ext-dynamic-memory}

到此為止，我們的變數都是在\textbf{編譯時}決定大小、在\textbf{執行時}建立
和銷毀。 但有時候程式需要「根據使用者輸入決定需要多少記憶體」，或「在
執行期間新增、刪除資料元素」。 這時就需要\textbf{動態配置}（dynamic
allocation）——在程式執行時，用 \texttt{malloc} 從 \emph{heap} 上要求一
塊記憶體，用完後用 \texttt{free} 歸還。 指標是動態配置的核心應用：沒
有指標，就無法把記憶體中四散各地、執行期間才建立的區塊連接起來，這正
是鏈結串列、樹這類\emph{資料結構}的基礎。

\texttt{malloc}／\texttt{free} 的完整規則——三種會壞掉的方式（記憶體
洩漏、懸置指標、雙重釋放）、怎麼用 AddressSanitizer 之類的工具把它們
抓出來、\texttt{calloc}／\texttt{realloc} 怎麼用——已經在
第~\ref{ch:datastructure}~章（資料結構）系統講過，這裡不重複；還沒讀
過的話，建議先讀那一章再回來看後面的案例。

\begin{drills}[label=drills:ext-dynamic-memory]
\item \texttt{malloc}／\texttt{free} 的完整規則（記憶體洩漏、懸置指
標、雙重釋放三種錯誤，以及用 AddressSanitizer 之類的工具找出它們）寫
在哪一章？
\item 為什麼說「指標是動態配置的核心應用」？
\end{drills}

\section{struct 的記憶體布局與對齊}
\label{sec:ext-struct-padding}

到這裡，我們已經知道 \texttt{int} 占 $4$ 位元組、\texttt{char} 占
$1$ 位元組。 但把不同型別的欄位\emph{放進同一個 struct} 時，總大小
往往\emph{不是}各欄位大小單純相加。 編譯器會在欄位之間插入 padding
（填充位元組），讓每個欄位的起始位址符合 CPU 的\emph{對齊}
（alignment）要求。 對齊不是 C 語言隨意加的累贅：許多 CPU 從「不是
$4$ 的倍數」的位址讀取 \texttt{int} 會比較慢，有些架構甚至直接當
成硬體錯誤。

先看一個只有三個 \texttt{char} 的 struct，沒有 padding 的餘地：

\begin{code}{c}{連續三個 char，大小就是 3}
struct ThreeChars {
    char a, b, c;
};
// sizeof(struct ThreeChars) == 3
\end{code}

如果在兩個 \texttt{char} 中間插入一個 \texttt{int}，布局就會變複雜：

\begin{center}
\begin{tikzpicture}
    \node[hlcell] (structA) {a};
    \node[cell, minimum width=0.45cm, right=0pt of structA] (pad1a) {};
    \node[cell, minimum width=0.45cm, right=0pt of pad1a] (pad1b) {};
    \node[cell, minimum width=0.45cm, right=0pt of pad1b] (pad1c) {};
    \node[hlcell2, right=0pt of pad1c] (x1) {x};
    \node[hlcell2, right=0pt of x1] (x2) {x};
    \node[hlcell2, right=0pt of x2] (x3) {x};
    \node[hlcell2, right=0pt of x3] (x4) {x};
    \node[hlcell, right=0pt of x4] (structC) {c};
    \node[cell, minimum width=0.45cm, right=0pt of structC] (pad2a) {};
    \node[cell, minimum width=0.45cm, right=0pt of pad2a] (pad2b) {};
    \node[cell, minimum width=0.45cm, right=0pt of pad2b] (pad2c) {};

    \node[idxlbl, below=0.2cm of structA] {a};
    \node[idxlbl, below=0.2cm of pad1b] {pad};
    \coordinate (xMid) at ($(x1.south)!0.5!(x4.south)$);
    \node[idxlbl, below=0.2cm of xMid] {x (int)};
    \node[idxlbl, below=0.2cm of structC] {c};
    \node[idxlbl, below=0.2cm of pad2b] {pad};
\end{tikzpicture}
\diagcaption{\texttt{struct CharIntChar \{ char a; int x; char c; \}} 在
常見 $32$-bit 平台上的布局：\texttt{int x} 必須從 $4$ 的倍數位址開始，
所以 \texttt{a} 後面插入 $3$ 個 padding 位元組，末尾也可能再 pad 到
整個 struct 大小是 $4$ 的倍數。}
\label{fig:ch21-struct-padding}
\end{center}

範例 \thechapter.11 用 \texttt{sizeof} 直接量出兩個 struct 的大小：
\texttt{struct ThreeChars} 是 \texttt{3}，\texttt{struct CharIntChar} 在
同一台機器上是 \texttt{12}。 中間那個 \texttt{int} 加上前後 padding，
把三個「邏輯上只有 $1 + 4 + 1 = 6$ 位元組」的欄位撐成 $12$ 位元組。

\begin{example}{c}{(struct-padding.c) 用 sizeof 觀察 padding 造成的差異}
#include <stdio.h>

struct ThreeChars {
    char a, b, c;
};

struct CharIntChar {
    char a;
    int x;
    char c;
};

int main(void) {
    printf("%zu %zu\n", sizeof(struct ThreeChars), sizeof(struct CharIntChar));
}
\end{example}

\iosample{}{
3 12 \\
}

\keypoint{\texttt{sizeof(struct ...)} 量的是「編譯器實際分配的大小」，包
含 padding；不能假設 struct 大小等於各欄位大小之和。 這在之後用
\texttt{memcpy} 複製 raw 位元組、或跟硬體暫存器對接時特別重要。}

如果真的很在意每個位元組都不能浪費（例如嵌入式協定封包），可以請編譯
器用 \texttt{\_\_attribute\_\_((packed))} 或 \texttt{\#pragma pack}
關閉 padding。 但這可能導致未對齊存取變慢，甚至觸發硬體例外，所以一般
應用程式很少這樣做，比較常見的是\emph{調整欄位順序}（把相同型別或
較大的欄位放一起）來減少 padding，而不是完全關掉對齊。

\begin{drills}[label=drills:ext-struct-padding]
\item \texttt{struct ThreeChars \{ char a, b, c; \}} 的 \texttt{sizeof}
是多少？ 為什麼沒有 padding？
\item \texttt{struct CharIntChar \{ char a; int x; char c; \}} 在範例機
器上 \texttt{sizeof} 是 \texttt{12}，不是邏輯上的 \texttt{6}，多出來的
$6$ 個位元組是什麼？
\item 為什麼 \texttt{int} 欄位要從 $4$ 的倍數位址開始，編譯器才要插入
padding？
\end{drills}

\section{整數在記憶體裡長什麼樣}
\label{sec:ext-int-bytes}

把前面幾節串起來：一個 \texttt{int} 在記憶體裡就是連續的 $4$ 個位元
組；union 讓我們用 \texttt{bytes[0..3]} 逐一檢視；struct padding 則告訴
我們「型別大小」跟「布局規則」會一起決定位元組怎麼排。 下一章的位元
遮罩、移位、popcount，全都建立在「整數就是一串固定長度的位元，從最低
位（LSB）到最高位（MSB）有明確編號」這個圖像上。 如果這裡的 union 範
例 \thechapter.6 還記得 \texttt{256} 拆成 \texttt{0 1 0 0}，下一章讀
\texttt{x \& 1}、\texttt{x >> 1} 時就不會覺得那些運算子從天而降。

\begin{center}
\begin{tikzpicture}
    \foreach \v/\lbl [count=\i from 0] in {0/b0 (LSB), 1/b1, 0/b2, 0/b3 (MSB)} {
        \ifnum\i=0
            \node[cell] (byte\i) {\v};
        \else
            \pgfmathtruncatemacro{\prev}{\i-1}
            \node[cell, right=0pt of byte\prev] (byte\i) {\v};
        \fi
        \node[idxlbl, below=0.2cm of byte\i] {\lbl};
    }
    \node[idxlbl, right=0.65cm of byte3] {值 $= 256$};
\end{tikzpicture}
\diagcaption{little-endian 機器上 \texttt{int} 值 \texttt{256} 的四個位
元組；位元運算章節會改從「位元 0、位元 1、\ldots」的角度看同一串
資料。}
\label{fig:ch21-little-endian-bytes}
\end{center}

\begin{drills}[label=drills:ext-int-bytes]
\item 這一節說「整數就是一串固定長度的位元，從 LSB 到 MSB 有明確編
號」，這個圖像跟哪一節的 union 範例直接對應？
\item little-endian 機器上，\texttt{256} 的四個位元組由低到高分別是
什麼？
\end{drills}

\section[位元操作（見第 12 章）]{位元操作詳解}
\label{sec:ext-bitops-pointer}

位元操作的完整內容在獨立的第~\ref{ch:bitoperations}~章（位元操作）：
進制轉換、操作符定義，到打包、解包與演算法應用都在那裡。 本章只在後
面的案例裡把位元操作跟其他高等型別結合，示範實戰應用，不重複講一次
基礎，詳見第~\ref{ch:bitoperations}~章。

\begin{drills}[label=drills:ext-bitops-pointer]
\item 位元操作的完整內容（進制轉換、操作符定義、打包解包與演算法應
用）在哪一章？
\item 這一節為什麼不重新講一次位元操作的基礎？
\end{drills}

\section{高等技術的綜合應用}
\label{sec:ext-synthesis}

前面幾節的觀念——型別、記憶體布局、位元運算、動態配置——在實務中往往
一起用。 這一節用三個具體例子展示「高等技術」怎麼組織出有用的東西。

\subsection{案例 1：集合型別（bitset）}

題庫裡常見「集合」題目：給定元素上限（例如最多 64 個元素），用位元表示
某元素是否在集合裡。 一個 \texttt{unsigned long long} 可以表示 64 個
元素的集合。

\begin{code}{c}{集合型別的設計}
typedef unsigned long long Set;

// 包含元素 i
Set insert(Set s, int i) {
    return s | (1ULL << i);
}

// 移除元素 i（不叫 remove：那是 <stdio.h> 刪檔案函式的名字，撞名要避免）
Set set_remove(Set s, int i) {
    return s & ~(1ULL << i);
}

// 檢查元素 i 是否存在
int contains(Set s, int i) {
    return (s >> i) & 1;
}

// 聯集
Set union_set(Set a, Set b) {
    return a | b;
}

// 交集
Set intersect(Set a, Set b) {
    return a & b;
}
\end{code}

這個設計結合了 \texttt{typedef}（提供有語意的名字）、位元運算（高效
實作）、以及位元字段的思想（每位代表一個元素）。

\subsection{案例 2：結構序列化}

讀寫二進位檔案或網路協定時，常常要把 \texttt{struct} 的位元模式逐位元
組寫出去。 這時會結合 \texttt{union} 和 \texttt{memcpy}：

\begin{code}{c}{用 union 進行結構序列化}
struct Color {
    unsigned char r, g, b;
};

union ColorBytes {
    struct Color rgb;
    unsigned int packed;  // 32-bit integer
};

// 把 RGB 顏色轉成 32-bit 整數（little-endian）
unsigned int colorToInt(struct Color c) {
    union ColorBytes cb;
    cb.rgb = c;
    return cb.packed;
}
\end{code}

這裡結合了 \texttt{struct} 的語意清晰度、\texttt{union} 的位元重新解讀，
以及 \texttt{sizeof} 對記憶體布局的理解。

\subsection{案例 3：鏈結串列節點}

鏈結串列是資料結構的基礎，每個節點是用 \texttt{malloc} 動態建立的
\texttt{struct}，內含指向下一節點的指標。 這同時用上了 \texttt{typedef}、
\texttt{struct}、指標，以及動態配置：

\begin{code}{c}{鏈結串列節點的典型設計}
typedef struct Node {
    int data;
    struct Node *next;  // 自指型別
} Node;

Node *head = NULL;

// 新增節點到鏈結串列前端
void insert(int val) {
    Node *new_node = (Node *)malloc(sizeof(Node));
    if (new_node == NULL) return;
    new_node->data = val;
    new_node->next = head;
    head = new_node;
}

// 走訪鏈結串列
void print_all(void) {
    for (Node *p = head; p != NULL; p = p->next) {
        printf("%d ", p->data);
    }
}
\end{code}

這三個案例展示了高等技術不是為了複雜而複雜，而是精確地用編譯器和 CPU
的能力去解決問題：集合用位元的高效性，序列化用 union 的位元重新解讀，
鏈結串列用指標和動態配置的靈活性。 掌握這些工具，就能寫出既高效又清晰
的程式碼。

\begin{drills}[label=drills:ext-synthesis]
\item 案例 1 的 bitset 為什麼把「移除元素」的函式命名為
\texttt{set\_remove} 而不是直接叫 \texttt{remove}？
\item 案例 2 的 \texttt{union ColorBytes} 讓 \texttt{struct Color} 跟
\texttt{unsigned int packed} 共用記憶體，這裡 union 的用途，跟前面
\texttt{union IntBytes} 那節的用途是同一件事嗎？
\item 案例 3 的 \texttt{Node} 結構裡，\texttt{struct Node *next;} 為什
麼可以在自己的定義裡出現指向自己型別的指標，卻不會造成結構大小是無
窮大？
\end{drills}

\addcontentsline{toc}{section}{練習解答}
\section*{練習解答}

各節練習的答案收在這裡。 短題只給能一句話核對的結果。

\paragraph{第~\ref{sec:ext-typedef}~節}
\begin{enumerate}
\item 底層型別仍是 \texttt{int}；編譯器把 \texttt{Score} 當成
\texttt{int} 處理，不會多占記憶體或改變運算規則。
\item 只需要改 \texttt{typedef} 那一行本身；這正是 typedef 讓「物理表
示可以換、語意名字留著」的價值。
\item 一樣，兩者都是替既有型別（一個是 \texttt{int}、一個是
\texttt{struct}）取一個表達語意的別名。
\end{enumerate}

\paragraph{第~\ref{sec:ext-enum}~節}
\begin{enumerate}
\item \texttt{CAND\_A} 是 \texttt{0}，\texttt{CAND\_B} 是 \texttt{1}。
\item \texttt{winner} 是 \texttt{CAND\_A}，輸出 \texttt{A 623}。
\item 不行；enum 常數只是有名字的整數，本身不會自動變成可印的字串，
要另外寫 \texttt{if} 或 \texttt{switch} 轉換。
\end{enumerate}

\paragraph{第~\ref{sec:ext-union}~節}
\begin{enumerate}
\item 等於最大成員 \texttt{int value} 的大小，也就是 \texttt{4}。
\item \texttt{0 1 0 0}（\texttt{256} 的最低位元組是 \texttt{0}，次低位
元組是 \texttt{1}）。
\item \texttt{struct} 的成員各自占一段記憶體；\texttt{union} 的成員共
用同一段記憶體，同一時刻只有其中一個有效。
\end{enumerate}

\paragraph{第~\ref{sec:ext-storage-class}~節}
\begin{enumerate}
\item 會連續印出 \texttt{1 1 1}，因為每次呼叫都重新從 \texttt{0} 開
始；原本的 \texttt{static} 版本則保留累計值，印出 \texttt{1 2 3}。
\item 存在函式呼叫的 stack frame 裡；函式返回後這層堆疊被丟棄，變數
跟著消失。
\item \texttt{extern} 宣告通常放在標頭檔（\texttt{.h}），真正的定義放
在某一個 \texttt{.c} 檔裡。
\end{enumerate}

\paragraph{第~\ref{sec:ext-dynamic-memory}~節}
\begin{enumerate}
\item 第~\ref{ch:datastructure}~章（資料結構）。
\item 因為執行期間才建立、分散在記憶體各處的區塊，必須靠指標才能連
接起來，這正是鏈結串列、樹這類資料結構的基礎。
\end{enumerate}

\paragraph{第~\ref{sec:ext-struct-padding}~節}
\begin{enumerate}
\item \texttt{3}；三個 \texttt{char} 都是 $1$ 位元組，彼此對齊要求相
同，不需要 padding。
\item padding：\texttt{a} 後面補 $3$ 個位元組讓 \texttt{int x} 對齊到
$4$ 的倍數位址，結尾可能再補 $3$ 個位元組讓整個 struct 大小是 $4$ 的
倍數。
\item 因為許多 CPU 從不是 $4$ 的倍數的位址讀取 \texttt{int} 會比較慢，
甚至可能直接觸發硬體錯誤，編譯器插入 padding 讓每個欄位起始位址符合
對齊要求。
\end{enumerate}

\paragraph{第~\ref{sec:ext-int-bytes}~節}
\begin{enumerate}
\item 跟 \texttt{union IntBytes} 那節、把 \texttt{256} 拆成
\texttt{0 1 0 0} 四個位元組的範例對應。
\item \texttt{0}、\texttt{1}、\texttt{0}、\texttt{0}（由低到高）。
\end{enumerate}

\paragraph{第~\ref{sec:ext-bitops-pointer}~節}
\begin{enumerate}
\item 第~\ref{ch:bitoperations}~章（位元操作）。
\item 因為完整的基礎已經在第~\ref{ch:bitoperations}~章系統講過，這一
節只在案例裡把位元操作跟其他高等型別結合示範應用，重複講一次沒有必
要。
\end{enumerate}

\paragraph{第~\ref{sec:ext-synthesis}~節}
\begin{enumerate}
\item 因為 \texttt{remove} 是 \texttt{<stdio.h>} 裡刪除檔案的函式名
字，跟它撞名容易造成混淆或連結衝突，所以改用 \texttt{set\_remove}。
\item 是同一件事：兩處的 union 都是讓一段記憶體同時用兩種型別解讀，
藉此把資料在不同表示法之間轉換或拆解。
\item 不會，因為 \texttt{struct Node} 裡放的是指向 \texttt{struct Node}
的\emph{指標}，指標本身大小固定（不隨指向的型別而變），所以不會造成
無窮遞迴展開。
\end{enumerate}

\phantomsection
\addcontentsline{toc}{chapter}{參考資料}
\chapter*{參考資料}
\label{ch:references}

本書寫作過程中，幾本教材反覆被引用、對照，或直接影響了某幾節的節序安
排。 這裡列出來，方便讀者需要更完整、更深入的說明時查閱原書。

\begin{description}
\item[K. N. King, \textit{C Programming: A Modern Approach}, 2nd
edition.] 全書節序參照最多的一本，尤其是〈前置處理〉一章的
\texttt{\#if} 與條件編譯、〈物件〉一章「介面先於實作」的設計次序，都
直接照這本書 ch.19「Program Design」的骨架安排。

\item[Stephen Prata, \textit{C Primer Plus}, 6th edition.] 〈前置處理〉
與〈C 標準程式庫〉兩章裡，字串化（\texttt{\#}）、token 貼接
（\texttt{\#\#}）、可變參數巨集的範例編排，參照了這本書 ch.16 的作
法。

\item[Peter van der Linden, \textit{Expert C Programming: Deep C
Secrets}.] 〈宣告修飾詞〉一章「順時針／螺旋法則」的說法出自這本書；
〈前置處理〉一章對編譯／連結流程的描述，也參照了這本書對建置流程的
講解。

\item[Richard Reese, \textit{Understanding and Using C Pointers}.]
（O'Reilly, 2013）〈物件〉一章「不完整型別」那一節，關於「只暴露使用
者需要的資訊」這個資訊隱藏原則的講法，直接引用自這本書。

\item[劉邦鋒，《由片語學習 C 程式設計》，第二版，國立臺灣大學。] 本書
最初的骨架設計——用「片語＋範例」交替、搭配學習要點與風格要點——延續
自這本書的編排慣例，是本書體例最直接的源頭。
\end{description}

這五本之外，各章節末的「常見錯誤清單」與練習，部分取材自標準函式庫
（\texttt{<assert.h>}、\texttt{<stdlib.h>} 等）的官方文件與 ISO C 標
準本身逐條核對過的行為，沒有另外列出版本資訊；行文中提到「C 標準規
定」的地方，指的都是 C11（ISO/IEC 9899:2011）。


\end{document}
